% !TeX program = pdflatex
% Geometria algebraiczna — audyt szczegółowy całości 2026-11-18, poprawki matematyczno-redakcyjne i układ rysunków.
% Rozdziały 0--10: od teorii pierścieni do topologii odmian algebraicznych.
% Dokument obejmuje także rozdział 3 wskazany przez użytkownika.
\documentclass[11pt,a4paper,openany]{report}
\usepackage[T1]{fontenc}
\usepackage[utf8]{inputenc}
\usepackage[polish]{babel}
\usepackage{lmodern}
\usepackage{amsmath,amssymb,amsthm,mathtools}
\usepackage{microtype}
\usepackage{xcolor}
\usepackage{geometry}
\usepackage{enumitem}
\usepackage{tikz}
\usetikzlibrary{arrows.meta,positioning,calc,decorations.pathreplacing}
\usepackage[unicode]{hyperref}
\hypersetup{colorlinks=true,linkcolor=blue!45!black,urlcolor=blue!45!black}
\geometry{margin=2.5cm}
\setlength{\parindent}{0pt}
\setlength{\parskip}{3pt}
\setlength{\emergencystretch}{2em}
\setlist[enumerate]{topsep=4pt,itemsep=2pt}
\setlist[itemize]{topsep=4pt,itemsep=2pt}
\numberwithin{equation}{section}
\theoremstyle{plain}
\newtheorem{twierdzenie}{Twierdzenie}[section]
\newtheorem{lemat}[twierdzenie]{Lemat}
\newtheorem{wniosek}[twierdzenie]{Wniosek}
\newtheorem{stwierdzenie}[twierdzenie]{Stwierdzenie}
\theoremstyle{definition}
\newtheorem{definicja}[twierdzenie]{Definicja}
\newtheorem{przyklad}[twierdzenie]{Przykład}
\theoremstyle{remark}
\newtheorem{uwaga}[twierdzenie]{Uwaga}
\newcommand{\N}{\mathbb N}
\newcommand{\Z}{\mathbb Z}
\newcommand{\Q}{\mathbb Q}
\newcommand{\R}{\mathbb R}
\newcommand{\C}{\mathbb C}
\newcommand{\F}{\mathbb F}
\newcommand{\A}{\mathbb A}
\newcommand{\PP}{\mathbb P}
\newcommand{\kk}{\Bbbk}
\newcommand{\V}{\operatorname V}
\newcommand{\I}{\operatorname I}
\newcommand{\id}{\operatorname{id}}
\newcommand{\im}{\operatorname{im}}
\newcommand{\kerf}{\operatorname{ker}}
\newcommand{\Spec}{\operatorname{Spec}}
\newcommand{\rad}{\sqrt}
\newcommand{\ev}{\operatorname{ev}}
\newcommand{\Frac}{\operatorname{Frac}}
\newcommand{\trdeg}{\operatorname{trdeg}}
\title{Geometria algebraiczna\\[3pt]\large Pierścienie, krzywe, powierzchnie, snopy, geometria biwymierna i topologia}
\author{}
\date{}

% BEGIN GLOBAL CODE LISTINGS (generated from latex/pl/code-listings.tex)
\usepackage{xcolor}
\usepackage{listings}
\definecolor{codebg}{RGB}{248,249,251}
\definecolor{codeborder}{RGB}{180,191,203}
\definecolor{codecomment}{RGB}{61,113,79}
\definecolor{codekeyword}{RGB}{41,78,132}
\lstdefinelanguage{Julia}{
  morekeywords={abstract,baremodule,begin,break,catch,const,continue,do,else,elseif,
    end,export,false,finally,for,function,global,if,import,in,isa,let,local,macro,
    module,mutable,primitive,quote,return,struct,true,try,type,using,where,while,
    missing,nothing,println,print},
  sensitive=true,morecomment=[l]{\#},morecomment=[n]{\#=}{=\#},
  morestring=[b]",
}
\lstdefinestyle{skrypt}{
  backgroundcolor=\color{codebg},frame=single,rulecolor=\color{codeborder},
  framesep=5pt,frameround=tttt,framerule=.5pt,
  basicstyle=\ttfamily\footnotesize,keywordstyle=\color{codekeyword}\bfseries,
  commentstyle=\color{codecomment}\itshape,stringstyle=\color{purple!65!black},
  numbers=left,numberstyle=\tiny\color{gray!75!black},numbersep=8pt,
  stepnumber=1,breaklines=true,breakatwhitespace=false,
  columns=fullflexible,keepspaces=true,showstringspaces=false,tabsize=4,
  linewidth=.90\linewidth,captionpos=b,aboveskip=11pt,belowskip=12pt,
  xleftmargin=12pt,xrightmargin=2pt,
}
\lstdefinestyle{pythonpde}{style=skrypt,language=Python}
\lstdefinestyle{python}{style=skrypt,language=Python}
\lstdefinestyle{julia}{style=skrypt,language=Julia}
\lstset{style=skrypt}
\renewcommand{\lstlistingname}{Kod}
\renewcommand{\lstlistlistingname}{Spis kodów}
% END GLOBAL CODE LISTINGS
\begin{document}
\hypersetup{pageanchor=false}
\maketitle
\hypersetup{pageanchor=true}
\pagenumbering{roman}
\tableofcontents
\clearpage
\pagenumbering{arabic}
\setcounter{chapter}{-1}

\chapter{Teoria pierścieni: od arytmetyki do geometrii równań}
\label{ch:pierscienie}
\section{Po co nam pierścienie?}

Przed badaniem zbiorów opisanych równaniami wielomianowymi potrzebujemy języka,
który pozwoli rozumieć same wielomiany i relacje między nimi. Nie chodzi tylko
 o wykonywanie rachunków. Równania $f=0$ oraz $g=0$ mogą opisywać ten sam zbiór,
różne zbiory, albo różne \emph{struktury algebraiczne} ukryte pod tym samym
zbiorem punktów. Te różnice najlepiej badać za pomocą pierścieni i ideałów.

\begin{przyklad}[Trzy rodzaje arytmetyki]
Rozważmy wyrażenie $x^2-1=(x-1)(x+1)$.
\begin{enumerate}[label=\arabic*)]
\item W $\R$ z równania $x^2=1$ wynika $x=1$ albo $x=-1$.
\item W $\Z/8\Z$ równanie $x^2=1$ ma aż cztery rozwiązania:
$\overline1,\overline3,\overline5,\overline7$.
Rzeczywiście, jeśli $x=2m+1$, to $x^2=4m(m+1)+1$, a ponieważ $m(m+1)$
jest parzyste, mamy $x^2\equiv1\pmod8$.
Jeżeli natomiast $x$ jest parzyste, to $x^2$ jest parzyste,
a więc nie może być przystające do $1$ modulo $8$.
Wszystkie cztery klasy nieparzyste i tylko one są zatem rozwiązaniami.
\item W pierścieniu $\R[t]/(t^2)$ element $t$ jest niezerowy, lecz
$t^2=0$. Nie wolno zatem wnioskować z $ab=0$, że $a=0$ lub $b=0$.
\end{enumerate}
Te różnice nie wynikają z odmiennego zapisu symboli, lecz z własności
struktury, w której wykonujemy działania.
\end{przyklad}

\begin{przyklad}[Równania i funkcje]
Wielomian $y-x^2\in\R[x,y]$ opisuje parabolę. Gdy na tej paraboli
pracujemy z wielomianami, nie potrzebujemy osobno pamiętać wszystkich
wielomianów przyjmujących na niej te same wartości: dwa z nich różnią się
o wielokrotność $y-x^2$ dokładnie wtedy, gdy wyznaczają na \emph{rzeczywistych}
punktach paraboli tę samą funkcję wielomianową. Wykorzystujemy tu, że
niezerowy wielomian jednej zmiennej nad $\R$ nie może mieć nieskończenie
wielu pierwiastków; nad ciałem skończonym analogiczne twierdzenie o funkcjach
nie jest prawdziwe. Wyrazimy to później izomorfizmem
$\R[x,y]/(y-x^2)\cong\R[x]$ i udowodnimy go, zamiast jedynie odwoływać się
do rysunku.
\end{przyklad}

W tym rozdziale budujemy pojęcia w następującej kolejności:
\emph{pierścień $\to$ elementy szczególne $\to$ homomorfizm $\to$
ideał $\to$ pierścień ilorazowy $\to$ ideał pierwszy i maksymalny $\to$
pierścienie wielomianów}. Na końcu połączymy je z pierwszymi przykładami
geometrycznymi. Nie zakładamy znajomości teorii pierścieni; potrzebna jest
jedynie elementarna algebra i pojęcia grupy oraz odwzorowania.

\section{Definicja pierścienia i jej konsekwencje}

\begin{definicja}[Pierścień]
\label{def:pierscien}
Pierścieniem nazywamy zbiór $A$ z działaniami
$+\colon A\times A\to A$ i $\cdot\colon A\times A\to A$ oraz wyróżnionymi
elementami $0,1\in A$, które spełniają warunki:
\begin{enumerate}[label=(P\arabic*)]
\item $(A,+)$ jest grupą abelową: dodawanie jest łączne i przemienne,
$0$ jest elementem neutralnym, a każdy $a\in A$ ma element przeciwny $-a$;
\item mnożenie jest łączne: $(ab)c=a(bc)$ dla wszystkich $a,b,c\in A$;
\item $1$ jest obustronnym elementem neutralnym mnożenia:
$1a=a1=a$ dla każdego $a\in A$;
\item zachodzą obie rozdzielności:
$a(b+c)=ab+ac$ oraz $(a+b)c=ac+bc$ dla wszystkich $a,b,c\in A$.
\end{enumerate}
Pierścień jest \emph{przemienny}, gdy ponadto $ab=ba$ dla dowolnych $a,b\in A$.
\end{definicja}

Pozwalamy na $1=0$; wtedy, jak zaraz sprawdzimy, pierścień ma tylko jeden
element i nazywa się \emph{pierścieniem zerowym}. Jest to wygodne, ponieważ
iloraz $A/A$ będzie pierścieniem. \textbf{Jeżeli nie powiedziano inaczej,
od tego miejsca wszystkie pierścienie są przemienne i mają jedynkę.}
Gdy mowa o dziedzinie lub ciele, wymagamy dodatkowo $1\ne0$.

\begin{uwaga}[Dlaczego rozdzielamy aksjomaty?]
W definicji nie ma założenia, że wolno dzielić przez element różny od zera.
Nie zakładamy nawet, że z $ab=ac$ i $a\ne0$ wynika $b=c$.
Te własności odróżnią później ciała i dziedziny od ogólnych pierścieni.
Również przemienność mnożenia nie wynika z pozostałych aksjomatów:
na przykład macierze kwadratowe tworzą zwykle pierścień nieprzemienny.
\end{uwaga}

\begin{lemat}[Podstawowe reguły rachunku]
\label{lem:reguly}
Niech $A$ będzie dowolnym pierścieniem, niekoniecznie przemiennym.
Dla wszystkich $a,b\in A$ zachodzą równości
$0a=a0=0$, $(-a)b=-(ab)=a(-b)$ oraz $(-a)(-b)=ab$.
Jeżeli $1=0$, to $A=\{0\}$.
\end{lemat}
\begin{proof}
Z rozdzielności mamy $0a=(0+0)a=0a+0a$.
Dodając do obu stron element przeciwny do $0a$ w grupie $(A,+)$,
otrzymujemy $0=0a$. Analogicznie $a0=0$.
Ponieważ $(-a)b+ab=((-a)+a)b=0b=0$, element $(-a)b$ jest przeciwny
do $ab$, czyli $(-a)b=-(ab)$. Tak samo $a(-b)=-(ab)$.
Wreszcie $(-a)(-b)=-\bigl(a(-b)\bigr)=-(-(ab))=ab$.
Jeżeli $1=0$, to dla każdego $a\in A$ mamy $a=1a=0a=0$.
\end{proof}

\begin{definicja}[Podpierścień]
Niech $A$ będzie pierścieniem. Podzbiór $B\subseteq A$ jest
\emph{podpierścieniem}, jeżeli $0,1\in B$, dla wszystkich $a,b\in B$
zachodzi $a-b\in B$ i $ab\in B$. Działania na $B$ są ograniczeniami
działań w $A$.
\end{definicja}

Warunek $1\in B$ jest częścią naszej konwencji. Dlatego $2\Z$ nie jest
podpierścieniem $\Z$ (chociaż jest zamknięty na dodawanie, odejmowanie
i mnożenie). Później okaże się, że $2\Z$ jest \emph{ideałem}.

\begin{przyklad}[Łańcuch podpierścieni]
Mamy $\Z\subseteq\Q\subseteq\R\subseteq\C$.
Każdy z tych zbiorów zawiera tę samą jedynkę $1$ i jest zamknięty
na odpowiednie działania. Podpierścienie mogą się znacznie różnić:
w $\Z$ liczba $2$ nie ma odwrotności względem mnożenia,
a w $\Q$ taką odwrotnością jest $1/2$.
\end{przyklad}

\section{Najważniejsze przykłady pierścieni}

\subsection{Liczby całkowite i arytmetyka modulo \texorpdfstring{$n$}{n}}

W $\Z$ dodajemy i mnożymy w zwykły sposób. Aby otrzymać nowy przykład,
ustalamy $n\ge2$ i mówimy, że $a,b\in\Z$ są przystające modulo $n$,
jeżeli $n\mid(a-b)$. Zapisujemy to $a\equiv b\pmod n$.
Relacja ta jest relacją równoważności: zwrotność wynika z $n\mid0$,
symetria z $n\mid(a-b)\Rightarrow n\mid(b-a)$,
a przechodniość z faktu, że suma dwóch wielokrotności $n$
również jest wielokrotnością $n$.

Klasę liczby $a$ oznaczamy $\overline a$; zbiór klas to $\Z/n\Z$.
Chcemy zdefiniować $\overline a+\overline b=\overline{a+b}$ i
$\overline a\,\overline b=\overline{ab}$. Trzeba sprawdzić,
że wynik nie zależy od wybranych reprezentantów.

\begin{stwierdzenie}
\label{st:modulo}
Dla każdego $n\ge2$ powyższe wzory definiują działania na $\Z/n\Z$,
a z tymi działaniami $\Z/n\Z$ jest pierścieniem przemiennym.
\end{stwierdzenie}
\begin{proof}
Załóżmy, że $a'=a+nr$ oraz $b'=b+ns$ dla pewnych $r,s\in\Z$.
Wtedy $(a'+b')-(a+b)=n(r+s)$, zatem $a'+b'\equiv a+b\pmod n$.
Ponadto
$a'b'-ab=(a+nr)(b+ns)-ab=n(as+br+nrs)$,
więc $a'b'\equiv ab\pmod n$. Działania są zatem dobrze określone.
Wszystkie aksjomaty pierścienia przechodzą z $\Z$ na klasy:
przykładowo $\overline a(\overline b+\overline c)
=\overline{a(b+c)}=\overline{ab+ac}
=\overline a\,\overline b+\overline a\,\overline c$.
Pozostałe aksjomaty sprawdza się w ten sam sposób, używając
odpowiednich własności liczb całkowitych.
\end{proof}

\begin{center}
\begin{tikzpicture}[x=0.79cm,y=1cm,>=Stealth]
\draw[->,gray!75] (-6.7,0)--(7.1,0) node[right] {$\Z$};
\foreach \i in {-6,...,6} {\draw (\i,0.10)--(\i,-0.10);}
\foreach \i in {-6,...,6} {
  \pgfmathtruncatemacro{\r}{mod(\i+99,3)}
  \ifnum\r=0\def\col{blue!70!black}\else
    \ifnum\r=1\def\col{red!75!black}\else\def\col{teal!70!black}\fi
  \fi
  \fill[\col] (\i,0) circle (2.7pt);
  \node[font=\scriptsize,below=5pt,\col] at (\i,0) {$\i$};
}
\draw[blue!70!black,thick,rounded corners] (-5.35,0.48)--(-2.65,0.48)
  node[midway,above,font=\small] {$\overline0$};
\draw[red!75!black,thick,rounded corners] (-4.35,1.03)--(-1.65,1.03)
  node[midway,above,font=\small] {$\overline1$};
\draw[teal!70!black,thick,rounded corners] (-3.35,1.58)--(-0.65,1.58)
  node[midway,above,font=\small] {$\overline2$};
\node[font=\small,align=center] at (3.6,1.25)
  {Kolor oznacza klasę\\reszt modulo $3$.};
\end{tikzpicture}
\end{center}

Na rysunku punkty tego samego koloru leżą w jednej klasie równoważności.
Klasa nie jest jedną liczbą: na przykład $\overline1$ zawiera
$\dots,-5,-2,1,4,7,\dots$. W ilorazie traktujemy je jako jeden element.

\subsection{Wielomiany, funkcje i iloczyny}

\begin{przyklad}[Pierścień wielomianów]
Jeżeli $A$ jest pierścieniem, to $A[x]$ oznacza zbiór formalnych sum
$f(x)=a_0+a_1x+\dots+a_dx^d$, gdzie $d\ge0$ i $a_j\in A$, oraz wielomian
zerowy. Dwa wielomiany są równe wtedy i tylko wtedy, gdy mają identyczne
współczynniki przy każdej potędze $x$. Dodajemy je współczynnikami,
a mnożenie określa rozdzielność i reguła $x^ix^j=x^{i+j}$.
Przykładowo w $\Z[x]$ mamy
$(2+x)(1+3x)=2+7x+3x^2$.
Nie utożsamiamy na tym etapie formalnego wielomianu z funkcją:
na $\F_2=\Z/2\Z$ wielomian $x^2-x$ nie jest wielomianem zerowym,
choć w obu punktach ciała $\F_2$ przyjmuje wartość zero.
\end{przyklad}

\begin{uwaga}[Wielomian formalny a funkcja wielomianowa]
\label{uw:formalna-funkcja}
Dla dowolnego pierścienia przemiennego $A$ podstawienie określa
odwzorowanie $A[x]\to A^A$, które wielomianowi $f$ przyporządkowuje
funkcję $a\mapsto f(a)$. Nie zawsze jest ono różnowartościowe.
Przykładowo w $\mathbb F_2[x]$ wielomiany $x^2-x$ oraz $0$ są różne,
ale wyznaczają tę samą funkcję na $\mathbb F_2$, gdyż
$0^2-0=0$ i $1^2-1=0$.

Jeśli natomiast $k$ jest ciałem nieskończonym, to odwzorowanie
$k[x]\to k^k$ jest różnowartościowe: gdy dwa wielomiany wyznaczają tę samą
funkcję, ich różnica znika w każdym punkcie $k$, a niezerowy wielomian
ma tylko skończenie wiele pierwiastków (twierdzenie~\ref{tw:korzenie}).
W przypadku $n$ zmiennych wniosek jest analogiczny. Dowód prowadzimy
indukcyjnie: dla wielomianu
$f(x_1,\dots,x_n)=\sum_{j=0}^d a_j(x_1,\dots,x_{n-1})x_n^j$
zerującego się na całym $k^n$ ustalamy pierwsze $n-1$ współrzędnych.
Otrzymany wielomian jednej zmiennej jest zerowy, więc wszystkie
$a_j$ znikają na $k^{n-1}$; indukcja pokazuje, że wszystkie są
wielomianami zerowymi.

W geometrii algebraicznej rozróżniamy więc 
\emph{pierścień formalnych wielomianów}, 
\emph{funkcje na punktach nad $k$} oraz 
\emph{pierścień zadany relacjami} (iloraz przez ideał).
Te obiekty nie są zamienne bez dodatkowych założeń.
\end{uwaga}

\begin{przyklad}[Pierścień funkcji]
Niech $X$ będzie niepustym zbiorem, a $A$ pierścieniem.
Na zbiorze $A^X$ wszystkich odwzorowań $f\colon X\to A$ definiujemy
$(f+g)(x)=f(x)+g(x)$ oraz $(fg)(x)=f(x)g(x)$.
Zero i jedynka są funkcjami stałymi równymi odpowiednio $0$ i $1$.
Sprawdzenie rozdzielności sprowadza się do równości
$(f(g+h))(x)=f(x)(g(x)+h(x))=f(x)g(x)+f(x)h(x)$
dla każdego $x\in X$. Otrzymujemy pierścień.
\end{przyklad}

\begin{przyklad}[Iloczyn pierścieni]
Dla pierścieni $A,B$ zbiór $A\times B$ jest pierścieniem z działaniami
$(a,b)+(a',b')=(a+a',b+b')$ i
$(a,b)(a',b')=(aa',bb')$. Jedynką jest $(1_A,1_B)$.
W pierścieniu $\R\times\R$ zachodzi $(1,0)(0,1)=(0,0)$,
choć żaden z dwóch czynników nie jest zerowy.
\end{przyklad}

\begin{przyklad}[Pierścień macierzy — przykład ostrzegawczy]
W zbiorze $M_2(\R)$ macierzy $2\times2$ mamy naturalne dodawanie,
mnożenie i macierz jednostkową. Jest to pierścień, ale nieprzemienny:
przyjmując
$P=\begin{psmallmatrix}1&0\\0&0\end{psmallmatrix}$ oraz
$Q=\begin{psmallmatrix}0&1\\0&0\end{psmallmatrix}$,
otrzymujemy $PQ=Q$, podczas gdy $QP=0$.
Od tej chwili będziemy głównie pracować w pierścieniach przemiennych,
ponieważ to właśnie ich spektra stanowią podstawę klasycznej geometrii
algebraicznej.
\end{przyklad}

\section{Elementy odwracalne, dzielniki zera i elementy nilpotentne}

\begin{definicja}[Jednostka, dzielnik zera, element nilpotentny]
Niech $A$ będzie pierścieniem przemiennym.
\begin{enumerate}[label=\arabic*)]
\item Element $a\in A$ jest \emph{odwracalny} (jest \emph{jednostką}),
jeżeli istnieje $b\in A$ takie, że $ab=1$. Wtedy $b$ oznaczamy $a^{-1}$.
Zbiór jednostek oznaczamy $A^\times$.
\item Niezerowy $a\in A$ jest \emph{dzielnikiem zera}, jeżeli istnieje
niezerowy $b\in A$ taki, że $ab=0$.
\item Element $a\in A$ jest \emph{nilpotentny}, jeżeli $a^m=0$
dla pewnego całkowitego $m\ge1$.
\item Element $e\in A$ jest \emph{idempotentny}, jeżeli $e^2=e$.
\end{enumerate}
\end{definicja}

\begin{stwierdzenie}
Niech $A$ będzie pierścieniem przemiennym. Jednostki tworzą grupę
abelową $(A^\times,\cdot)$. Jeżeli $a$ jest jednostką i $ab=ac$,
to $b=c$. Żadna jednostka w niezerowym pierścieniu nie jest dzielnikiem zera.
\end{stwierdzenie}
\begin{proof}
Jedynka jest jednostką, a odwrotność jednostki również jest jednostką.
Jeżeli $ab=1$ i $cd=1$, to $(ac)(bd)=(ab)(cd)=1$, więc iloczyn jednostek
jest jednostką. Łączność i przemienność dziedziczymy z $A$.
Z $ab=ac$ po pomnożeniu przez $a^{-1}$ dostajemy
$(a^{-1}a)b=(a^{-1}a)c$, czyli $b=c$.
Jeśli $ab=0$ i $a$ jest jednostką, to $b=a^{-1}0=0$.
\end{proof}

\begin{przyklad}[Które klasy modulo $n$ są odwracalne?]
W $\Z/12\Z$ jednostkami są dokładnie
$\overline1,\overline5,\overline7,\overline{11}$.
Na przykład $\overline5^2=\overline{25}=\overline1$.
Element $\overline6$ nie jest jednostką, ponieważ
$\overline6\,\overline2=\overline0$ i $\overline2\ne\overline0$.
Poniższe twierdzenie daje ogólną odpowiedź.
\end{przyklad}

\begin{lemat}[Tożsamość Bézouta dla liczb całkowitych]
\label{lem:bezout-int}
Niech $a,b\in\Z$ nie będą jednocześnie zerowe. Wówczas istnieją
$u,v\in\Z$ takie, że $ua+vb=\gcd(a,b)$.
\end{lemat}
\begin{proof}
Rozważmy zbiór $S=\{ua+vb:u,v\in\Z\}$.
Jest on zamknięty na dodawanie i odejmowanie. Ponieważ przynajmniej
jedna z liczb $a,b$ jest niezerowa, $S$ zawiera pewną liczbę dodatnią
(w razie potrzeby zmieniamy znak). Z zasady dobrego uporządkowania
liczb naturalnych istnieje najmniejsza dodatnia liczba $d\in S$.
Z definicji $S$ mamy $d=ua+vb$ dla pewnych $u,v\in\Z$.

Dzieląc $a$ przez $d$ z resztą, otrzymujemy $a=qd+r$,
gdzie $0\le r<d$. Ponieważ $a,d\in S$, także $r=a-qd\in S$.
Minimalność $d$ wymusza $r=0$, a zatem $d\mid a$.
Identycznie dowodzimy $d\mid b$.

Jeśli $c$ jest dowolnym wspólnym dzielnikiem $a$ i $b$,
to $c$ dzieli również ich kombinację $ua+vb=d$.
Wykazaliśmy, że $d$ jest dodatnim wspólnym dzielnikiem,
przez który dzieli się każdy wspólny dzielnik $a,b$.
Jest to właśnie największy wspólny dzielnik: $d=\gcd(a,b)$.
\end{proof}

\begin{twierdzenie}[Jednostki modulo $n$]
\label{tw:jednostki-mod}
Niech $n\ge2$ oraz $a\in\Z$. Klasa $\overline a\in\Z/n\Z$
jest jednostką wtedy i tylko wtedy, gdy $\gcd(a,n)=1$.
\end{twierdzenie}
\begin{proof}
Jeśli $\overline a$ jest jednostką, to istnieje $b\in\Z$ takie, że
$ab\equiv1\pmod n$. Oznacza to, że $ab-1=kn$ dla pewnego $k\in\Z$,
czyli $ab-kn=1$. Każdy wspólny dzielnik $a$ i $n$ dzieli lewą stronę,
a więc dzieli $1$. Stąd $\gcd(a,n)=1$.

Na odwrót, dla względnie pierwszych $a,n$ lemat~\ref{lem:bezout-int} daje $u,v\in\Z$ takie, że $au+nv=1$.
Przechodząc do klas modulo $n$, otrzymujemy
$\overline a\,\overline u=\overline1$. Zatem $\overline a$ jest jednostką.
\end{proof}

\begin{uwaga}[Skąd bierze się tożsamość Bézouta?]
Można ją otrzymać algorytmem Euklidesa. Dla $a=5,n=12$ mamy
$12=2\cdot5+2$, $5=2\cdot2+1$, a z podstawienia
$1=5-2(12-2\cdot5)=5\cdot5-2\cdot12$.
Stąd odwrotnością $\overline5$ w $\Z/12\Z$ jest $\overline5$.
Analogiczny algorytm dla wielomianów pojawi się dalej.
\end{uwaga}

\begin{definicja}[Dziedzina całkowitości i ciało]
\label{def:dziedzina-cialo}
\emph{Dziedziną całkowitości} nazywamy niezerowy pierścień przemienny,
w którym z $ab=0$ wynika $a=0$ lub $b=0$.
\emph{Ciałem} nazywamy niezerowy pierścień przemienny,
w którym każdy element różny od zera jest jednostką.
\end{definicja}

Każde ciało jest dziedziną: jeśli $ab=0$ i $a\ne0$, możemy pomnożyć
przez $a^{-1}$ i otrzymać $b=0$. Przykładami ciał są $\Q,\R,\C$,
a przykładem dziedziny niebędącej ciałem jest $\Z$.

\begin{stwierdzenie}[Skracanie w dziedzinie]
Dla niezerowego pierścienia przemiennego $A$ następujące warunki są
równoważne: (i) $A$ jest dziedziną całkowitości;
(ii) dla wszystkich $a,b,c\in A$ z $a\ne0$ i $ab=ac$ wynika $b=c$.
\end{stwierdzenie}
\begin{proof}
Jeśli $A$ jest dziedziną, to $ab=ac$ daje $a(b-c)=0$,
a $a\ne0$ wymusza $b-c=0$. Jeśli zachodzi (ii) oraz $ab=0$ przy $a\ne0$,
to $ab=a0$, więc $b=0$. Jest to definicja dziedziny.
\end{proof}

\begin{twierdzenie}[Skończona dziedzina jest ciałem]
\label{tw:skonczona}
Każda skończona dziedzina całkowitości $A$ jest ciałem.
\end{twierdzenie}
\begin{proof}
Ustalmy $a\in A$ z $a\ne0$ i rozważmy odwzorowanie
$m_a\colon A\to A$ zadane przez $m_a(x)=ax$.
Jeżeli $m_a(x)=m_a(y)$, to $ax=ay$. W dziedzinie możemy skrócić
niezerowy czynnik $a$, więc $x=y$; odwzorowanie $m_a$ jest różnowartościowe.
Ponieważ $A$ jest skończony, każde różnowartościowe odwzorowanie $A\to A$
jest też na. Zatem $1=m_a(b)=ab$ dla pewnego $b\in A$.
Dowolny element niezerowy $a$ ma odwrotność, a więc $A$ jest ciałem.
\end{proof}

\begin{wniosek}
Jeżeli $p$ jest liczbą pierwszą, to $\Z/p\Z$ jest ciałem.
Jeżeli $n\ge2$ jest złożone, to $\Z/n\Z$ nie jest dziedziną.
\end{wniosek}
\begin{proof}
Dla pierwszego $p$ każda liczba $1\le a\le p-1$ spełnia
$\gcd(a,p)=1$, więc każda niezerowa klasa jest odwracalna
na mocy twierdzenia~\ref{tw:jednostki-mod}.
Jeśli $n=rs$ z $1<r,s<n$, to
$\overline r\ne0$ i $\overline s\ne0$, ale
$\overline r\,\overline s=\overline n=0$.
\end{proof}

\begin{przyklad}[Nilpotentność nie jest tym samym co bycie dzielnikiem zera]
W $\Z/12\Z$ element $\overline6$ jest nilpotentny,
gdyż $\overline6^2=\overline{36}=0$.
Natomiast $\overline4$ jest idempotentny:
$\overline4^2=\overline{16}=\overline4$.
Element $\overline4$ jest też dzielnikiem zera, bo
$\overline4\,\overline3=0$, ale nie jest nilpotentny:
z $e^2=e$ wynika indukcyjnie $e^m=e$ dla każdego $m\ge1$.
W niezerowej dziedzinie nie ma niezerowych elementów nilpotentnych
ani idempotentów innych niż $0$ i $1$.
\end{przyklad}

\begin{lemat}[Nilpotentność i odwracalność]
Niech $A$ będzie pierścieniem przemiennym i niech $u\in A$ spełnia
$u^m=0$ dla pewnego $m\ge1$. Wtedy $1-u$ jest jednostką, a jej
odwrotnością jest $1+u+\dots+u^{m-1}$.
\end{lemat}
\begin{proof}
Korzystamy ze wzoru na sumę geometryczną, który działa w każdym
pierścieniu przemiennym:
$(1-u)(1+u+\dots+u^{m-1})=1-u^m=1$.
\end{proof}


\subsection{Charakterystyka pierścienia}

\begin{definicja}[Charakterystyka]
Niech $A$ będzie niezerowym pierścieniem.
Jeżeli istnieje dodatnia liczba całkowita $n$ spełniająca
$\underbrace{1+\dots+1}_{n\ \text{składników}}=0$,
to najmniejszą taką liczbę nazywamy \emph{charakterystyką} $A$
i oznaczamy $\operatorname{char}A$.
Jeżeli nie ma takiej liczby, przyjmujemy $\operatorname{char}A=0$.
\end{definicja}

\begin{przyklad}
Mamy $\operatorname{char}\Z=\operatorname{char}\Q
=\operatorname{char}\R=\operatorname{char}\C=0$,
a dla $n\ge2$ zachodzi $\operatorname{char}(\Z/n\Z)=n$.
Dla dowolnego niezerowego pierścienia $A$ charakterystyka $A[x]$ jest równa
charakterystyce $A$, ponieważ $1$ w $A[x]$ jest wielomianem stałym $1$.
\end{przyklad}

\begin{stwierdzenie}
Niech $A$ będzie dziedziną całkowitości. Wówczas
$\operatorname{char}A=0$ albo $\operatorname{char}A=p$
dla pewnej liczby pierwszej $p$.
\end{stwierdzenie}
\begin{proof}
Przypuśćmy, że charakterystyka jest dodatnia i wynosi $n$.
Gdyby $n=rs$ dla $1<r,s<n$, to
$(\underbrace{1+\dots+1}_{r})(\underbrace{1+\dots+1}_{s})
=\underbrace{1+\dots+1}_{rs}=0$.
Z minimalności $n$ oba czynniki są niezerowe, co przeczy
brakowi dzielników zera w dziedzinie. Zatem $n$ jest liczbą pierwszą.
\end{proof}

\section{Homomorfizmy: kiedy dwa opisy zachowują rachunek?}

\begin{definicja}[Homomorfizm pierścieni]
Niech $A,B$ będą pierścieniami. Odwzorowanie
$\varphi\colon A\to B$ jest \emph{homomorfizmem pierścieni}, jeżeli dla
wszystkich $a,b\in A$ zachodzą równości
$\varphi(a+b)=\varphi(a)+\varphi(b)$,
$\varphi(ab)=\varphi(a)\varphi(b)$ oraz $\varphi(1_A)=1_B$.
Bijektywny homomorfizm nazywamy \emph{izomorfizmem}.
\end{definicja}

Warunku zachowywania jedynki nie pomijamy. Jest on istotny:
zerowe odwzorowanie między niezerowymi pierścieniami zachowuje dodawanie
i mnożenie, ale nie wysyła jedynki na jedynkę.

\begin{stwierdzenie}
Jeżeli $\varphi\colon A\to B$ jest homomorfizmem pierścieni, to
$\varphi(0_A)=0_B$ i $\varphi(-a)=-\varphi(a)$ dla każdego $a\in A$.
Ponadto obraz jednostki $a\in A^\times$ jest jednostką w $B$,
a $\varphi(a^{-1})=\varphi(a)^{-1}$.
\end{stwierdzenie}
\begin{proof}
Ponieważ $\varphi(0_A)=\varphi(0_A+0_A)
=\varphi(0_A)+\varphi(0_A)$, możemy odjąć $\varphi(0_A)$
i otrzymać $\varphi(0_A)=0_B$.
Następnie $\varphi(a)+\varphi(-a)=\varphi(a-a)=0_B$,
więc $\varphi(-a)=-\varphi(a)$.
Gdy $ab=1_A$, to $\varphi(a)\varphi(b)=\varphi(1_A)=1_B$,
co dowodzi ostatniej części.
\end{proof}

\begin{definicja}[Jądro i obraz]
Dla homomorfizmu $\varphi\colon A\to B$ definiujemy
$\kerf\varphi=\{a\in A:\varphi(a)=0_B\}$ oraz
$\im\varphi=\{\varphi(a):a\in A\}\subseteq B$.
\end{definicja}

\begin{przyklad}[Redukcja modulo $n$]
Dla $n\ge2$ odwzorowanie
$\pi_n\colon\Z\to\Z/n\Z$, $a\mapsto\overline a$, jest
surjektywnym homomorfizmem. Jego jądro składa się dokładnie z liczb
podzielnych przez $n$: $\kerf\pi_n=n\Z$.
Odwzorowanie to \emph{zapomina}, którą konkretną liczbą całkowitą
reprezentowana jest klasa, ale zachowuje dodawanie i mnożenie.
\end{przyklad}

\begin{przyklad}[Podstawienie do wielomianu]
Niech $A$ będzie pierścieniem i $a\in A$. Odwzorowanie
$\ev_a\colon A[x]\to A$ dane przez $\ev_a(f)=f(a)$
jest homomorfizmem pierścieni. Dla $f=\sum_i f_ix^i$ i
$g=\sum_jg_jx^j$ mamy
$\ev_a(fg)=\sum_{i,j}f_ig_ja^{i+j}
=(\sum_if_ia^i)(\sum_jg_ja^j)=f(a)g(a)$.
Dodawanie i zachowywanie jedynki sprawdzamy bezpośrednio.
Surjektywność wynika z tego, że każda stała $b\in A$ jest wartością
stałego wielomianu $b$.
\end{przyklad}

\begin{uwaga}[Dlaczego jądro jest ważniejsze od samego wzoru?]
Dwa elementy $a,b\in A$ mają ten sam obraz wtedy i tylko wtedy, gdy
$\varphi(a-b)=0$, czyli $a-b\in\kerf\varphi$.
Jądro dokładnie opisuje więc \emph{jakie elementy stają się
nierozróżnialne} po zastosowaniu homomorfizmu. Aby z takiego
utożsamiania otrzymać nowy pierścień, potrzebujemy pojęcia ideału.
\end{uwaga}

\section{Ideały: podzbiory nadające się do ilorazowania}

\begin{definicja}[Ideał]
Niech $A$ będzie pierścieniem przemiennym. Podzbiór $I\subseteq A$
nazywamy \emph{ideałem}, jeżeli
\begin{enumerate}[label=(I\arabic*)]
\item $0\in I$ i $a-b\in I$ dla wszystkich $a,b\in I$;
\item $ra\in I$ dla wszystkich $r\in A$ i $a\in I$.
\end{enumerate}
Ideał jest \emph{właściwy}, gdy $I\ne A$.
\end{definicja}

Pierwszy warunek znaczy, że $(I,+)$ jest podgrupą $(A,+)$.
Drugi nazywa się \emph{pochłanianiem przez mnożenie}.
Można wyobrażać sobie ideał jako zbiór elementów, które chcemy
\emph{uznać za zero}: jeśli $a$ ma stać się zerem, to również
$ra$ musi stać się zerem, aby mnożenie pozostawało sensowne.

\begin{przyklad}[Dlaczego $2\Z$ jest ideałem]
Suma i różnica dwóch liczb parzystych są parzyste,
a iloczyn dowolnej liczby całkowitej z parzystą jest parzysty.
Zatem $2\Z$ jest ideałem $\Z$. Nie jest jednak podpierścieniem
w naszym znaczeniu, ponieważ $1\notin2\Z$.
\end{przyklad}

\begin{stwierdzenie}[Test właściwości ideału]
Niech $I$ będzie ideałem pierścienia $A$. Wówczas
$I=A$ wtedy i tylko wtedy, gdy $1\in I$.
Jeśli $I$ zawiera jednostkę, to $I=A$.
\end{stwierdzenie}
\begin{proof}
Jeżeli $I=A$, to oczywiście $1\in I$. Jeśli zaś $1\in I$,
to dla każdego $a\in A$ mamy $a=a\cdot1\in I$, więc $I=A$.
Gdy $u\in I$ jest jednostką, pochłanianie daje
$1=u^{-1}u\in I$, więc z pierwszej części $I=A$.
\end{proof}

\begin{definicja}[Ideał generowany przez zbiór]
Dla podzbioru $S\subseteq A$ przez $(S)$ oznaczamy najmniejszy ideał
pierścienia $A$ zawierający $S$, to znaczy przecięcie wszystkich ideałów
zawierających $S$. Jeśli $S=\{a_1,\dots,a_r\}$, piszemy
$(a_1,\dots,a_r)$. Ideał $(a)$ nazywamy \emph{głównym}.
\end{definicja}

\begin{stwierdzenie}[Jawny opis generowania]
Dla $S\subseteq A$ ideał $(S)$ składa się z wszystkich skończonych sum
$\sum_{j=1}^m r_js_j$, gdzie $m\ge0$, $r_j\in A$ i $s_j\in S$.
W szczególności $(a_1,\dots,a_r)=\{r_1a_1+\dots+r_ra_r:r_j\in A\}$
oraz $(a)=\{ra:r\in A\}$.
\end{stwierdzenie}
\begin{proof}
Nazwijmy zbiór wszystkich wskazanych sum $J$.
Suma pusta daje $0\in J$. Suma oraz różnica dwóch takich wyrażeń
też jest sumą odpowiedniej postaci (minus przenosimy na współczynniki).
Dla $t\in A$ mamy $t\sum_jr_js_j=\sum_j(tr_j)s_j\in J$.
Zatem $J$ jest ideałem. Ponadto każdy $s\in S$ należy do $J$,
bo $s=1\cdot s$.
Jeśli $I$ jest dowolnym ideałem zawierającym $S$, to pochłanianie
daje $r_js_j\in I$ dla wszystkich $j$, a zamknięcie na dodawanie
daje $\sum_jr_js_j\in I$. Stąd $J\subseteq I$.
Zbiór $J$ jest więc najmniejszym ideałem zawierającym $S$, czyli $J=(S)$.
\end{proof}

\begin{przyklad}[Dwa generatory mogą być potrzebne]
W $\Z$ mamy $(6,10)=(2)$: każda kombinacja $6a+10b$ jest parzysta,
a z drugiej strony $2=2\cdot6-1\cdot10$.
Natomiast w $k[x,y]$, gdzie $k$ jest ciałem, ideał $(x,y)$
nie jest generowany przez jeden wielomian.
Dowód tego drugiego faktu podamy po omówieniu dzielenia wielomianów;
ważne jest na razie rozróżnienie ideałów jedno- i wielogeneratorowych.
\end{przyklad}

\begin{stwierdzenie}[Suma, przecięcie i iloczyn ideałów]
Dla ideałów $I,J$ pierścienia $A$ zbiory
$I\cap J$ oraz $I+J=\{a+b:a\in I,\ b\in J\}$ są ideałami.
Suma $I+J$ jest najmniejszym ideałem zawierającym $I\cup J$.
Zbiór
$IJ=\{\sum_{r=1}^m a_rb_r:m\ge0,\ a_r\in I,\ b_r\in J\}$
także jest ideałem, a ponadto $IJ\subseteq I\cap J$.
\end{stwierdzenie}
\begin{proof}
Przecięcie dwóch ideałów zawiera zero i jest zamknięte na odejmowanie
oraz mnożenie przez elementy $A$, bo każda z tych własności zachodzi
w obu ideałach. Dla sumy: $(a+b)-(a'+b')=(a-a')+(b-b')\in I+J$,
a $r(a+b)=ra+rb\in I+J$. Ideał $I+J$ zawiera $I$ i $J$,
bo $0$ należy do obu; dowolny ideał zawierający $I$ i $J$ zawiera
również wszystkie sumy $a+b$.

Zbiór $IJ$ jest ideałem: można łączyć i odejmować skończone sumy,
a mnożenie przez $r\in A$ daje
$r\sum a_rb_r=\sum(ra_r)b_r$ z $ra_r\in I$.
Każdy składnik $a_rb_r$ należy do $I$ (bo $a_r\in I$)
i do $J$ (bo $b_r\in J$), a więc $IJ\subseteq I\cap J$.
\end{proof}

\begin{przyklad}[Przecięcie nie musi równać się iloczynowi]
W $\Z$ mamy $(2)\cap(4)=(4)$, natomiast $(2)(4)=(8)$.
Inaczej jest dla $(2)$ i $(3)$: wtedy
$(2)\cap(3)=(6)=(2)(3)$. Równość w drugim przypadku wiąże się
z faktem, że $2$ i $3$ są względnie pierwsze.
\end{przyklad}

\begin{definicja}[Ideały współmaksymalne]
Ideały $I,J\subseteq A$ nazywamy \emph{współmaksymalnymi}, jeśli
$I+J=A$. Równoważnie istnieją $a\in I$, $b\in J$ takie, że $a+b=1$.
\end{definicja}

\begin{stwierdzenie}
Niech $A$ będzie pierścieniem przemiennym z jedynką, a $I,J\subseteq A$ ideałami.
Jeżeli $I+J=A$, to $I\cap J=IJ$.
\end{stwierdzenie}
\begin{proof}
Wiemy już, że $IJ\subseteq I\cap J$.
Wybierzmy $a\in I,b\in J$ z $a+b=1$. Jeśli $c\in I\cap J$,
to $c=c(a+b)=ca+cb$. Ponieważ $c\in J$ i $a\in I$,
mamy $ca\in IJ$; ponieważ $c\in I$ i $b\in J$,
także $cb\in IJ$. Zatem $c\in IJ$, co daje drugą inkluzję.
\end{proof}

\section{Pierścienie ilorazowe: co oznacza utożsamienie elementów?}

\begin{definicja}[Pierścień ilorazowy]
Niech $I$ będzie ideałem pierścienia $A$. Definiujemy relację
$a\sim_Ib$ wtedy i tylko wtedy, gdy $a-b\in I$.
Klasę elementu $a$ oznaczamy $a+I=\{a+i:i\in I\}$.
Zbiór tych klas zapisujemy $A/I$.
Działania na klasach określamy wzorami
$(a+I)+(b+I)=(a+b)+I$ oraz $(a+I)(b+I)=ab+I$.
\end{definicja}

\begin{twierdzenie}[Poprawność konstrukcji ilorazowej]
\label{tw:iloraz}
Dla każdego ideału $I\subseteq A$ relacja $\sim_I$ jest relacją
równoważności, powyższe działania są dobrze określone,
a $A/I$ jest pierścieniem przemiennym z zerem $0+I$
i jedynką $1+I$.
\end{twierdzenie}
\begin{proof}
Zwrotność: $a-a=0\in I$.
Symetria: jeśli $a-b\in I$, to $b-a=-(a-b)\in I$.
Przechodniość: jeśli $a-b,b-c\in I$, to
$a-c=(a-b)+(b-c)\in I$.

Załóżmy teraz, że $a+I=a'+I$ i $b+I=b'+I$.
Oznacza to $a-a'\in I$ i $b-b'\in I$.
Dla dodawania mamy
$(a+b)-(a'+b')=(a-a')+(b-b')\in I$.
Dla mnożenia korzystamy z rozkładu
$ab-a'b'=(a-a')b+a'(b-b')$.
Oba składniki po prawej należą do $I$ na mocy własności pochłaniania,
więc $ab+I=a'b'+I$. To jest kluczowy punkt dowodu:
\emph{bez pochłaniania mnożenie klas nie musiałoby być dobrze określone}.

Aksjomaty pierścienia wynikają z aksjomatów w $A$;
na przykład łączność mnożenia to
$((a+I)(b+I))(c+I)=(ab)c+I=a(bc)+I
=(a+I)((b+I)(c+I))$.
Analogicznie sprawdzamy rozdzielność, przemienność i istnienie
odwrotności addytywnej $(-a)+I$.
\end{proof}

\begin{przyklad}[Dlaczego podgrupa addytywna nie wystarczy]
W $A=\Z[x]$ zbiór $H=\Z$ jest podgrupą grupy addytywnej $(A,+)$,
ale nie jest ideałem $A$, ponieważ $1\in H$ i $x\cdot1=x\notin H$.
Gdybyśmy w klasach addytywnych modulo $H$ chcieli mnożyć,
to $0$ i $1$ byłyby w tej samej klasie, ale ich iloczyny przez $x$
różniłyby się o $x$, który nie należy do $H$.
Wynik mnożenia zależałby zatem od reprezentanta.
\end{przyklad}

\begin{definicja}[Rzutowanie kanoniczne]
Dla ideału $I\subseteq A$ odwzorowanie
$\pi_I\colon A\to A/I$, $\pi_I(a)=a+I$,
nazywamy \emph{rzutowaniem kanonicznym}.
\end{definicja}

\begin{stwierdzenie}
Niech $A$ będzie pierścieniem przemiennym z jedynką, a $I\subseteq A$
ideałem. Rzutowanie kanoniczne $\pi_I:A\to A/I$ jest
surjektywnym homomorfizmem pierścieni i $\kerf\pi_I=I$.
\end{stwierdzenie}
\begin{proof}
Każda klasa w $A/I$ jest obrazem swojego reprezentanta,
więc $\pi_I$ jest na. Zachowywanie działań wynika z ich definicji.
Ponadto $\pi_I(a)=0+I$ jest równoważne $a-0\in I$,
a więc $a\in I$.
\end{proof}

\begin{twierdzenie}[Pierwsze twierdzenie o izomorfizmie]
\label{tw:izo}
Dla każdego homomorfizmu pierścieni $\varphi\colon A\to B$
istnieje dokładnie jeden izomorfizm
$\widetilde\varphi\colon A/\kerf\varphi\to\im\varphi$
spełniający $\widetilde\varphi(a+\kerf\varphi)=\varphi(a)$
dla wszystkich $a\in A$.
\end{twierdzenie}
\begin{proof}
Najpierw sprawdzamy, że $\kerf\varphi$ jest ideałem:
$\varphi(0)=0$; jeśli $a,b$ należą do jądra, to
$\varphi(a-b)=\varphi(a)-\varphi(b)=0$;
jeśli $r\in A$ i $a\in\kerf\varphi$, to
$\varphi(ra)=\varphi(r)\varphi(a)=0$.

Definiujemy $\widetilde\varphi(a+\kerf\varphi)=\varphi(a)$.
Jeśli $a+\kerf\varphi=b+\kerf\varphi$, to
$a-b\in\kerf\varphi$, zatem $\varphi(a)=\varphi(b)$.
Wzór nie zależy więc od reprezentanta. Homomorfizm zachowuje działania,
bo na przykład
$\widetilde\varphi((a+K)(b+K))=\varphi(ab)
=\varphi(a)\varphi(b)$ dla $K=\kerf\varphi$.
Z definicji obrazu jest surjektywny na $\im\varphi$.
Jeżeli $\widetilde\varphi(a+K)=\widetilde\varphi(b+K)$,
to $\varphi(a-b)=0$, a więc $a-b\in K$, czyli $a+K=b+K$.
Jest zatem różnowartościowy. Jedyność wynika z tego, że każda klasa
w ilorazie jest postaci $a+K$, na której żądany wzór ustala wartość.
\end{proof}

\begin{center}
\begin{tikzpicture}[>=Stealth,node distance=2.15cm]
\node (a) {$A$};
\node[right=3.0cm of a] (im) {$\im\varphi$};
\node[below=1.6cm of a] (q) {$A/\kerf\varphi$};
\draw[->] (a)--node[above] {$\varphi$} (im);
\draw[->] (a)--node[left] {$\pi$} (q);
\draw[->] (q)--node[below right] {$\widetilde\varphi\ \cong$} (im);
\end{tikzpicture}
\end{center}
W diagramie obie drogi z $A$ do $\im\varphi$ dają ten sam wynik.
Mówimy, że diagram jest \emph{przemienny}.

\begin{przyklad}[Wartość wielomianu i iloraz]
Dla $a\in A$ homomorfizm $\ev_a\colon A[x]\to A$ jest surjektywny.
Udowodnimy teraz, że $\kerf\ev_a=(x-a)$.
Dla dowolnego $f(x)=\sum_{j=0}^m c_jx^j$ zapisujemy
\[
 f(x)-f(a)=\sum_{j=1}^m c_j(x^j-a^j)
 =(x-a)\sum_{j=1}^m c_j\sum_{r=0}^{j-1}x^{j-1-r}a^r.
\]
Jeżeli $f(a)=0$, to $f(x)$ jest wielokrotnością $x-a$.
Jeśli odwrotnie $f(x)=(x-a)g(x)$, to podstawienie $x=a$ daje $f(a)=0$.
Wykorzystując twierdzenie~\ref{tw:izo}, otrzymujemy
$A[x]/(x-a)\cong A$.
To elementarna wersja zasady: \emph{równanie $x=a$ pozwala usunąć
jedną zmienną}.
\end{przyklad}

\subsection{Ideały w ilorazie i ich odpowiedniki w pierścieniu}

\begin{twierdzenie}[Odpowiedniość ideałów]
\label{tw:odp-idealow}
Niech $A$ będzie pierścieniem i $I$ jego ideałem.
Przyporządkowanie $J\mapsto J/I=\{a+I:a\in J\}$ jest bijekcją
między ideałami $J\subseteq A$ spełniającymi $I\subseteq J$
a ideałami pierścienia $A/I$. Odwrotne przyporządkowanie ma postać
$K\mapsto\pi_I^{-1}(K)$. Bijekcja zachowuje inkluzje.
\end{twierdzenie}
\begin{proof}
Dla $I\subseteq J$ zbiór $J/I$ jest ideałem $A/I$:
zerowa klasa do niego należy, różnice klas elementów $J$ pozostają
w $J/I$, a dla $r+I\in A/I$ i $a+I\in J/I$ iloczyn
$ra+I$ leży w $J/I$ dzięki $ra\in J$.

Jeśli $K$ jest ideałem $A/I$, to przeciwobraz $\pi_I^{-1}(K)$ jest
ideałem $A$ (sprawdzamy warunki ideału, przenosząc je przez homomorfizm).
Zawiera $I$, ponieważ $I=\kerf\pi_I$, a zero należy do $K$.

Dla $I\subseteq J$ zachodzi
$\pi_I^{-1}(J/I)=J$: jeśli $a+I=j+I$ dla pewnego $j\in J$,
to $a-j\in I\subseteq J$, a więc $a\in J$; odwrotna inkluzja
jest oczywista. Z kolei dla ideału $K\subseteq A/I$ mamy
$\pi_I(\pi_I^{-1}(K))=K$ dzięki surjektywności $\pi_I$.
To dowodzi, że przyporządkowania są odwrotne.
Jeżeli $J_1\subseteq J_2$, to oczywiście $J_1/I\subseteq J_2/I$;
analogicznie dla przeciwobrazów.
\end{proof}

\begin{przyklad}[Wszystkie ideały $\Z/12\Z$]
Za chwilę pokażemy, że wszystkie ideały $\Z$ mają postać $d\Z$.
Te spośród nich, które zawierają $12\Z$, są generowane przez
dzielniki $12$ (z dokładnością do znaku):
$12\Z,6\Z,4\Z,3\Z,2\Z,\Z$.
Z twierdzenia~\ref{tw:odp-idealow} otrzymujemy sześć ideałów
pierścienia $\Z/12\Z$, oznaczanych odpowiednio
$(0),(\overline6),(\overline4),(\overline3),(\overline2),(1)$.
Ich uporządkowanie przez zawieranie można narysować następująco:
\end{przyklad}

\begin{center}
\begin{tikzpicture}[>=Stealth, every node/.style={font=\small,inner sep=3pt}]
\node (top) at (0,3.5) {$(1)$};
\node (two) at (-1.2,2.55) {$(\overline2)$};
\node (three) at (1.2,2.55) {$(\overline3)$};
\node (four) at (-1.2,1.55) {$(\overline4)$};
\node (six) at (1.2,1.55) {$(\overline6)$};
\node (zero) at (0,0.55) {$(0)$};
\draw (top)--(two)--(four)--(zero);
\draw (top)--(three)--(six)--(zero);
\draw (two)--(six);
\end{tikzpicture}
\end{center}
Każdy odcinek oznacza relację bezpośredniego zawierania:
ideał na dole jest zawarty w ideale na górze. Na przykład
$(\overline6)\subseteq(\overline2)$ i
$(\overline6)\subseteq(\overline3)$, ponieważ $6$ jest wielokrotnością
zarówno $2$, jak i $3$. Rysunek nie jest wykresem funkcji,
lecz diagramem porządku częściowego.


\subsection{Chińskie twierdzenie o resztach}

Istnieje ważna sytuacja, w której pierścień ilorazowy można rozłożyć
na iloczyn dwóch prostszych pierścieni. Pozwala to rozwiązywać układy
kongruencji, a zarazem tłumaczy obecność idempotentów w pierścieniach.

\begin{twierdzenie}[Chińskie twierdzenie o resztach]
\label{tw:crt}
Niech $A$ będzie pierścieniem przemiennym, a $I,J$ jego ideałami
takimi, że $I+J=A$. Wtedy odwzorowanie
\[
 \Theta\colon A/(I\cap J)\longrightarrow A/I\times A/J,\qquad
 a+(I\cap J)\longmapsto(a+I,a+J),
\]
jest izomorfizmem pierścieni. Ponieważ dla ideałów współmaksymalnych
$I\cap J=IJ$, można równoważnie zapisać $A/IJ\cong A/I\times A/J$.
\end{twierdzenie}
\begin{proof}
Rozważmy najpierw odwzorowanie
$\varphi\colon A\to A/I\times A/J$ zadane wzorem
$\varphi(a)=(a+I,a+J)$. Jest ono homomorfizmem, ponieważ każde
z rzutowań jest homomorfizmem i działania w iloczynie określamy
współrzędnie. Jego jądro składa się z elementów spełniających
jednocześnie $a+I=0+I$ oraz $a+J=0+J$; jest więc równe $I\cap J$.

Pozostaje wykazać surjektywność. Z założenia $I+J=A$ możemy wybrać
$i\in I$ oraz $j\in J$ z $i+j=1$.
Niech $(u+I,v+J)$ będzie dowolnym elementem $A/I\times A/J$.
Weźmy $a=uj+vi$. Ponieważ $j=1-i$, mamy
$a-u=u(j-1)+vi=-ui+vi\in I$.
Analogicznie $i=1-j$, więc
$a-v=uj+v(i-1)=uj-vj\in J$.
Stąd $\varphi(a)=(u+I,v+J)$, co dowodzi, że $\varphi$ jest na.
Pierwsze twierdzenie o izomorfizmie daje
$A/(I\cap J)\cong A/I\times A/J$ i dokładnie wskazany wzór na izomorfizm.
\end{proof}

\begin{przyklad}[Rozkład $\Z/12\Z$]
Ideały $3\Z$ oraz $4\Z$ są współmaksymalne, gdyż
$1=4-3\in3\Z+4\Z$. Ich przecięcie jest równe $12\Z$:
liczba całkowita jest podzielna jednocześnie przez $3$ i przez $4$
dokładnie wtedy, gdy jest podzielna przez $12$.
Chińskie twierdzenie o resztach daje
$\Z/12\Z\cong\Z/3\Z\times\Z/4\Z$.
Izomorfizm wysyła $\overline a$ na
$(a\bmod 3,a\bmod 4)$.
W szczególności $\overline4\mapsto(\overline1,\overline0)$
oraz $\overline9\mapsto(\overline0,\overline1)$.
Tłumaczy to rachunki $\overline4^2=\overline4$ i
$\overline9^2=\overline9$: odpowiadają one oczywistym równościom
$(1,0)^2=(1,0)$ oraz $(0,1)^2=(0,1)$ w iloczynie pierścieni.
Zauważmy także, że $\overline4+\overline9=\overline1$.
\end{przyklad}

\begin{uwaga}[Idea geometryczna]
Idempotenty w pierścieniu pozwalają rozdzielać jego strukturę na
części. W dalszej geometrii algebraicznej odpowiada to rozkładowi
obiektów na pewne rozłączne składowe. Nie należy jednak mylić tego
z rozkładem na składowe nierozkładalne: dwie proste przecinające się
w jednym punkcie tworzą zbiór rozkładalny, ale nierozłączny.
\end{uwaga}

\section{Ideały pierwsze i maksymalne}

\begin{definicja}[Ideał pierwszy]
Właściwy ideał $\mathfrak p\subsetneq A$ jest \emph{pierwszy}, gdy
dla wszystkich $a,b\in A$ zachodzi
$ab\in\mathfrak p\Rightarrow a\in\mathfrak p\ \text{lub}\ b\in\mathfrak p$.
\end{definicja}

\begin{definicja}[Ideał maksymalny]
Właściwy ideał $\mathfrak m\subsetneq A$ jest \emph{maksymalny},
jeżeli dla każdego ideału $J\subseteq A$ z
$\mathfrak m\subseteq J\subseteq A$ zachodzi
$J=\mathfrak m$ lub $J=A$.
\end{definicja}

\begin{uwaga}[Dwa różne znaczenia słowa „pierwszy”]
Element pierwszy i ideał pierwszy to pojęcia pokrewne, ale nie tożsame.
Najpierw definiujemy \emph{ideały pierwsze}, ponieważ działają one
w każdym pierścieniu przemiennym. Elementy pierwsze i nierozkładalne
pojawią się później wraz z teorią podzielności.
\end{uwaga}

\begin{twierdzenie}[Kryterium ilorazowe]
\label{tw:kryterium}
Niech $I$ będzie właściwym ideałem pierścienia $A$.
Wówczas:
\begin{enumerate}[label=\arabic*)]
\item $I$ jest pierwszy wtedy i tylko wtedy, gdy $A/I$ jest dziedziną całkowitości;
\item $I$ jest maksymalny wtedy i tylko wtedy, gdy $A/I$ jest ciałem.
\end{enumerate}
\end{twierdzenie}
\begin{proof}
Ponieważ $I$ jest właściwy, mamy $1+I\ne0+I$.

(1) Jeśli $I$ jest pierwszy i $(a+I)(b+I)=0+I$,
to $ab\in I$, więc $a\in I$ lub $b\in I$.
W takim razie $a+I=0+I$ lub $b+I=0+I$, czyli iloraz jest dziedziną.
Odwrotnie, jeśli $A/I$ jest dziedziną i $ab\in I$, to
$(a+I)(b+I)=0+I$, więc jeden z czynników jest zerową klasą,
co znaczy $a\in I$ lub $b\in I$.

(2) Załóżmy, że $I$ jest maksymalny. Weźmy niezerową klasę $a+I$,
czyli $a\notin I$. Ideał $I+(a)$ zawiera $I$ ściśle,
zatem z maksymalności jest równy $A$.
Istnieją więc $i\in I$ i $r\in A$ takie, że $i+ra=1$.
Przechodząc do ilorazu, dostajemy
$(r+I)(a+I)=1+I$, czyli każda niezerowa klasa ma odwrotność.
Zatem $A/I$ jest ciałem.

Odwrotnie, załóżmy, że $A/I$ jest ciałem.
Niech $J$ będzie ideałem zawierającym $I$ ściśle.
Wybierzmy $a\in J\setminus I$. Klasa $a+I$ jest niezerowa,
więc ma odwrotność $r+I$. Zatem $ra+I=1+I$, czyli $1-ra\in I\subseteq J$.
Ponieważ także $ra\in J$, otrzymujemy $1\in J$, a więc $J=A$.
Właśnie to oznacza maksymalność $I$.
\end{proof}

\begin{wniosek}
Każdy ideał maksymalny pierścienia przemiennego jest pierwszy.
\end{wniosek}
\begin{proof}
Jeśli $\mathfrak m$ jest maksymalny, to $A/\mathfrak m$ jest ciałem,
a więc dziedziną. Z pierwszej części twierdzenia~\ref{tw:kryterium}
wynika, że $\mathfrak m$ jest pierwszy.
\end{proof}

\begin{przyklad}[Ideały w $\Z$]
Pokażemy, że każdy ideał $I\subseteq\Z$ ma postać $n\Z$ dla pewnego
$n\ge0$. Gdy $I=\{0\}$, bierzemy $n=0$.
W przeciwnym przypadku $I$ zawiera dodatnią liczbę całkowitą
(jeżeli $a\in I$ jest ujemne, to $-a\in I$). Niech $n$ będzie
najmniejszą dodatnią liczbą w $I$.
Oczywiście $n\Z\subseteq I$. Dla dowolnego $a\in I$ wykonujemy
dzielenie z resztą $a=qn+r$, gdzie $0\le r<n$.
Wtedy $r=a-qn\in I$. Minimalność $n$ wymusza $r=0$,
więc $a\in n\Z$. Otrzymaliśmy $I=n\Z$.

Dla $n\ge2$ ideał $(n)$ jest pierwszy (równoważnie: maksymalny)
dokładnie wtedy, gdy $n$ jest liczbą pierwszą.
Istotnie, dla liczby pierwszej $p$ iloraz $\Z/(p)$ jest ciałem,
a dla liczby złożonej $n$ ma dzielniki zera.
Ideał $(0)$ jest pierwszy, ponieważ $\Z$ jest dziedziną,
ale nie jest maksymalny, ponieważ $\Z$ nie jest ciałem.
\end{przyklad}

\begin{przyklad}[Różnica między jednym a dwiema zmiennymi]
Niech $k$ będzie ciałem. W $k[x,y]$ ideał $(x)$ jest pierwszy,
ponieważ po podstawieniu $x=0$ uzyskujemy
$k[x,y]/(x)\cong k[y]$, a $k[y]$ jest dziedziną.
Nie jest jednak maksymalny, ponieważ $k[y]$ nie jest ciałem:
wielomian $y$ nie ma odwrotności w $k[y]$.
Z kolei $(x,y)$ jest maksymalny, ponieważ
$k[x,y]/(x,y)\cong k$ (w ilorazie $x$ i $y$ stają się zerem).
Dla pierwszego izomorfizmu wystarczy rozważyć homomorfizm
$\psi\colon k[x,y]\to k[y]$ dany przez $\psi(f)=f(0,y)$.
Jest on surjektywny (każdy wielomian w $y$ jest obrazem samego siebie).
Jeśli $f(x,y)=\sum_{j=0}^m a_j(y)x^j$, to $\psi(f)=a_0(y)$,
zatem $\psi(f)=0$ wtedy i tylko wtedy, gdy $a_0(y)=0$,
czyli gdy $f$ jest wielokrotnością $x$. Mamy $\kerf\psi=(x)$,
więc $k[x,y]/(x)\cong k[y]$. Drugi izomorfizm zostanie
uzasadniony ogólnym twierdzeniem o ideale punktu.
\end{przyklad}

\subsection{Pierwiastnik ideału i informacja, której nie widać na punktach}

\begin{definicja}[Pierwiastek ideału]
Dla ideału $I\subseteq A$ definiujemy jego \emph{pierwiastek}
$\sqrt I=\{a\in A:\text{istnieje }m\ge1\text{ takie, że }a^m\in I\}$.
Ideał $I$ nazywamy \emph{radykalnym}, gdy $\sqrt I=I$.
Pierścień $A$ jest \emph{zredukowany}, gdy jego jedynym
nilpotentnym elementem jest $0$.
\end{definicja}

\begin{twierdzenie}
\label{tw:pierwiastek}
Dla każdego ideału $I$ pierścienia przemiennego $A$
zbiór $\sqrt I$ jest ideałem zawierającym $I$.
Ponadto $A/I$ jest zredukowany wtedy i tylko wtedy, gdy $I$ jest radykalny.
\end{twierdzenie}
\begin{proof}
Jeżeli $a\in I$, to $a^1\in I$, więc $I\subseteq\sqrt I$.
Jeśli $a^m\in I$ oraz $r\in A$, to $(ra)^m=r^ma^m\in I$,
a więc $ra\in\sqrt I$. Podobnie $(-a)^m=(-1)^ma^m\in I$,
więc $-a\in\sqrt I$.

Niech teraz $a,b\in\sqrt I$, czyli $a^m\in I$ i $b^n\in I$
dla pewnych $m,n\ge1$. Rozwińmy $(a+b)^{m+n-1}$
wzorem dwumianowym. W każdym składniku
$\binom{m+n-1}{j}a^jb^{m+n-1-j}$
mamy $j\ge m$ lub $m+n-1-j\ge n$:
w przeciwnym razie $j\le m-1$ i $m+n-1-j\le n-1$,
co po dodaniu dałoby $m+n-1\le m+n-2$, sprzeczność.
Każdy składnik zawiera zatem czynnik $a^m$ lub $b^n$
i należy do $I$. Suma także należy do $I$.
Stąd $a+b\in\sqrt I$, co kończy dowód, że $\sqrt I$ jest ideałem.

Element $a+I$ jest nilpotentny w $A/I$ dokładnie wtedy, gdy
$(a+I)^m=0+I$ dla pewnego $m\ge1$, a to znaczy $a^m\in I$,
czyli $a\in\sqrt I$. W ilorazie nie ma niezerowych nilpotentów
dokładnie wtedy, gdy każdy taki $a$ należy już do $I$.
Warunek ten jest równoważny $\sqrt I=I$.
\end{proof}

\begin{przyklad}[Podwójny punkt]
Niech $k$ będzie ciałem, $A=k[t]$, $I=(t^2)$.
Pierścień $A/I$ zawiera niezerowy element $\overline t=t+(t^2)$,
ale $\overline t^2=0$. Element jest niezerowy,
bo wielomian $t$ nie może być wielokrotnością $t^2$
(jego stopień jest mniejszy).
Mamy $\sqrt{(t^2)}=(t)$: jeśli $f(t)^m$ jest podzielny przez $t^2$,
to $f(0)^m=0$, więc $f(0)=0$ i $t\mid f$;
odwrotnie, dla $f=tg$ mamy $f^2=t^2g^2\in(t^2)$.

Zatem $k[t]/(t)$ i $k[t]/(t^2)$ opisują ten sam warunek
na \emph{zwykłych punktach nad ciałem}: $t=0$.
Nie są jednak izomorficzne jako pierścienie, ponieważ w drugim
występuje niezerowy nilpotent, a w pierwszym (izomorficznym z $k$) nie.
\end{przyklad}

\section{Wielomiany dokładniej: stopień, dzielenie i pierwiastki}

\begin{definicja}[Stopień wielomianu]
Niech $A$ będzie pierścieniem. Jeśli
$f=a_0+a_1x+\dots+a_dx^d\in A[x]$ i $a_d\ne0$,
to $\deg f=d$, a $a_d$ jest \emph{współczynnikiem wiodącym}.
Dla wielomianu zerowego przyjmujemy formalnie $\deg0=-\infty$.
Wielomian jest \emph{moniczny}, gdy jego współczynnik wiodący jest równy $1$.
\end{definicja}

\begin{stwierdzenie}[Stopień iloczynu]
Jeżeli $A$ jest dziedziną całkowitości, a $f,g\in A[x]$ są niezerowe,
to $\deg(fg)=\deg f+\deg g$. W szczególności $A[x]$ jest dziedziną.
\end{stwierdzenie}
\begin{proof}
Niech $f=a_mx^m+\dots$ i $g=b_nx^n+\dots$, gdzie $a_m,b_n\ne0$.
Współczynnik przy $x^{m+n}$ w iloczynie wynosi $a_mb_n$:
aby uzyskać wykładnik $m+n$, trzeba pomnożyć wyrazy najwyższych stopni.
W dziedzinie $a_mb_n\ne0$. Nie ma wyrazów wyższego stopnia,
więc $\deg(fg)=m+n$. W szczególności iloczyn dwóch niezerowych
wielomianów nie jest zerowy, a zatem $A[x]$ jest dziedziną.
\end{proof}

\begin{uwaga}
Założenia o dziedzinie nie wolno pomijać. W $(\Z/6\Z)[x]$
mamy $(\overline2x)(\overline3x)=0$, choć oba czynniki są
niezerowymi wielomianami stopnia $1$.
\end{uwaga}

\begin{twierdzenie}[Dzielenie z resztą]
\label{tw:dzielenie}
Niech $k$ będzie ciałem, a $f,g\in k[x]$, przy czym $g\ne0$.
Istnieją dokładnie jedne wielomiany $q,r\in k[x]$ takie, że
$f=qg+r$ oraz $r=0$ lub $\deg r<\deg g$.
\end{twierdzenie}
\begin{proof}
Najpierw konstruujemy $q$ i $r$.
Jeśli $f=0$ albo $\deg f<\deg g$, przyjmujemy $q=0,r=f$.
W przeciwnym przypadku zapiszmy
$f=ax^m+\text{wyrazy niższe}$,
$g=bx^n+\text{wyrazy niższe}$, gdzie $m\ge n$ i $a,b\ne0$.
Ponieważ $k$ jest ciałem, element $b$ ma odwrotność.
Wielomian $h=f-ab^{-1}x^{m-n}g$ ma stopień mniejszy od $m$:
wyrazy wiodące obu składników dokładnie się znoszą.
Powtarzając tę operację na $h$, po skończonej liczbie kroków
(za każdym razem stopień maleje) otrzymujemy resztę $r$ o stopniu
mniejszym od $\deg g$ lub zero, a zebrane wielomiany
$ab^{-1}x^{m-n}$ tworzą iloraz $q$. To dowodzi istnienia.

Dla jedyności załóżmy, że $f=qg+r=q'g+r'$ i obie reszty
mają stopień mniejszy od $\deg g$ albo są zerowe.
Mamy $(q-q')g=r'-r$. Jeśli $q-q'\ne0$, to z własności stopnia
w dziedzinie $k[x]$ dostajemy
$\deg((q-q')g)=\deg(q-q')+\deg g\ge\deg g$.
Z drugiej strony $\deg(r'-r)<\deg g$ albo $r'-r=0$,
co jest niemożliwe. Zatem $q=q'$, a w konsekwencji $r=r'$.
\end{proof}

\begin{przyklad}[Dzielenie w praktyce]
W $\Q[x]$ podzielmy $f=x^3-2x+3$ przez $g=x-2$.
Najpierw odejmujemy $x^2(x-2)=x^3-2x^2$,
co zostawia $2x^2-2x+3$.
Następnie odejmujemy $2x(x-2)=2x^2-4x$,
co zostawia $2x+3$. Odejmujemy $2(x-2)=2x-4$
i otrzymujemy resztę $7$. A zatem
$x^3-2x+3=(x^2+2x+2)(x-2)+7$.
Warto sprawdzić wynik przez podstawienie $x=2$:
$f(2)=8-4+3=7$, zgodnie z resztą.
\end{przyklad}

\begin{wniosek}[Twierdzenie o reszcie i czynniku]
Dla ciała $k$, wielomianu $f\in k[x]$ i elementu $a\in k$
reszta z dzielenia $f$ przez $x-a$ jest stałą $f(a)$.
W szczególności $f(a)=0$ wtedy i tylko wtedy, gdy $(x-a)\mid f$.
\end{wniosek}
\begin{proof}
Z twierdzenia o dzieleniu mamy $f=(x-a)q+r$,
gdzie $r$ jest stałą lub zerem. Podstawiając $x=a$,
dostajemy $f(a)=r$. Stąd od razu wynika równoważność.
\end{proof}

\begin{twierdzenie}[Ograniczenie liczby pierwiastków]
\label{tw:korzenie}
Jeżeli $k$ jest ciałem, a $f\in k[x]$ jest niezerowym wielomianem
stopnia $d$, to $f$ ma w $k$ co najwyżej $d$ różnych pierwiastków.
\end{twierdzenie}
\begin{proof}
Dowodzimy indukcyjnie względem $d$.
Wielomian stopnia $0$ jest niezerową stałą i nie ma pierwiastków.
Jeśli $f$ stopnia $d\ge1$ nie ma pierwiastków, teza jest oczywista.
W przeciwnym razie wybierzmy pierwiastek $a$.
Na mocy twierdzenia o czynniku $f=(x-a)g$, gdzie $\deg g=d-1$.
Dla każdego innego pierwiastka $b\ne a$ zachodzi
$0=f(b)=(b-a)g(b)$.
Ponieważ $k$ jest ciałem, $b-a\ne0$ i możemy skrócić ten czynnik,
więc $g(b)=0$. Pozostałe pierwiastki $f$ są więc pierwiastkami $g$,
których jest najwyżej $d-1$ z założenia indukcyjnego.
Łącznie $f$ ma najwyżej $d$ różnych pierwiastków.
\end{proof}

\begin{uwaga}
Nie jest to prawda dla wielomianów nad dowolnymi pierścieniami:
w $\Z/8\Z$ wielomian $x^2-1$ ma cztery różne pierwiastki,
co widzieliśmy na początku. To dlatego warunek „$k$ jest ciałem”
nie jest kosmetyczny.
\end{uwaga}

\subsection{Algorytm Euklidesa i ideały w \texorpdfstring{$k[x]$}{k[x]}}

\begin{definicja}[Podzielność i element nierozkładalny]
Niech $A$ będzie dziedziną. Dla $a,b\in A$ piszemy $a\mid b$,
gdy istnieje $c\in A$ takie, że $b=ac$.
Element $u\in A$ jest \emph{nierozkładalny}, jeśli jest niezerowy,
nie jest jednostką i z $u=ab$ wynika, że $a$ lub $b$ jest jednostką.
Dwa elementy $a,b$ nazywamy \emph{stowarzyszonymi}, jeśli
$a=ub$ dla pewnej jednostki $u\in A^\times$.
\end{definicja}

\begin{stwierdzenie}[Ideały w {$k[x]$} są główne]
\label{st:kx-glowne}
Jeżeli $k$ jest ciałem, to każdy ideał $I\subseteq k[x]$
ma postać $(d)$ dla pewnego $d\in k[x]$.
Dla niezerowego $I$ wielomian $d$ można wybrać moniczny,
a wtedy jest on jednoznacznie wyznaczony.
\end{stwierdzenie}
\begin{proof}
Dla $I=(0)$ bierzemy $d=0$.
Jeżeli $I\ne(0)$, wybieramy niezerowy wielomian $d\in I$
o najmniejszym możliwym stopniu. Mnożąc przez odwrotność
współczynnika wiodącego, możemy założyć, że $d$ jest moniczny.
Dla dowolnego $f\in I$ dzielimy przez $d$:
$f=qd+r$, gdzie $r=0$ lub $\deg r<\deg d$.
Ponieważ $f,qd\in I$, mamy $r=f-qd\in I$.
Minimalność stopnia $d$ wymusza $r=0$, czyli $d\mid f$.
Stąd $I\subseteq(d)$. Druga inkluzja wynika z $d\in I$.

Jeśli moniczne $d,e$ generują ten sam niezerowy ideał,
to $d\mid e$ i $e\mid d$. Z wzoru na stopień otrzymujemy
$\deg d=\deg e$, więc $e=cd$ dla stałej $c\in k^\times$.
Moniczność obu wielomianów wymusza $c=1$, czyli $d=e$.
\end{proof}

\begin{twierdzenie}[Wielomianowy Bézout]
\label{tw:bezout-pol}
Niech $k$ będzie ciałem, a $f,g\in k[x]$ nie będą równocześnie zerowe.
Istnieje jedyny moniczny wielomian $d\in k[x]$ taki, że
$(f,g)=(d)$. Ponadto istnieją $u,v\in k[x]$ takie, że
$d=uf+vg$. Wielomian $d$ jest największym wspólnym dzielnikiem
$f$ i $g$ w sensie podzielności, znormalizowanym do postaci monicznej.
\end{twierdzenie}
\begin{proof}
Z twierdzenia~\ref{st:kx-glowne} ideał $(f,g)$ ma jedyny moniczny
generator $d$. Ponieważ $f,g\in(d)$, mamy $d\mid f$ oraz $d\mid g$.
Każdy wspólny dzielnik $c$ wielomianów $f,g$ dzieli wszystkie
kombinacje $uf+vg$; w szczególności dzieli $d$.
Ponieważ $d\in(f,g)$, na mocy jawnego opisu generowania
istnieją $u,v\in k[x]$ takie, że $d=uf+vg$.
\end{proof}

\begin{przyklad}[Obliczenie NWD wielomianów]
W $\Q[x]$ rozważmy $f=x^3-1$ i $g=x^2-1$.
Mamy $f=xg+(x-1)$ oraz $g=(x+1)(x-1)$,
więc $\gcd(f,g)=x-1$ (w postaci monicznej).
Co więcej, bezpośrednio z pierwszego dzielenia dostajemy
tożsamość Bézouta $x-1=f-xg$.
To świadectwo mówi, że ideał $(f,g)$ jest równy $(x-1)$.
\end{przyklad}

\begin{twierdzenie}[Ideały maksymalne w {$k[x]$}]
\label{tw:irreducible}
Niech $k$ będzie ciałem, a $f\in k[x]$ wielomianem niezerowym
i niestałym. Ideał $(f)$ jest maksymalny wtedy i tylko wtedy,
gdy $f$ jest nierozkładalny w $k[x]$. W takim przypadku
$k[x]/(f)$ jest ciałem.
\end{twierdzenie}
\begin{proof}
Załóżmy, że $f$ jest nierozkładalny.
Niech $g+(f)$ będzie niezerowym elementem ilorazu.
Wtedy $f\nmid g$. Każdy wspólny dzielnik $f$ i $g$
jest albo jednostką, albo jest stowarzyszony z $f$,
gdyż $f$ jest nierozkładalny; druga możliwość dawałaby $f\mid g$.
Zatem $\gcd(f,g)=1$. Z twierdzenia~\ref{tw:bezout-pol}
istnieją $u,v\in k[x]$ takie, że $uf+vg=1$.
Po przejściu do ilorazu dostajemy
$(v+(f))(g+(f))=1+(f)$. Każda niezerowa klasa ma więc odwrotność,
zatem $k[x]/(f)$ jest ciałem, a $(f)$ maksymalny.

Odwrotnie, jeśli $f=gh$ dla niejednostek $g,h\in k[x]$,
to $\deg g,\deg h\ge1$ i oba stopnie są mniejsze od $\deg f$.
Mamy $(f)\subseteq(g)$. Inkluzja jest ścisła:
gdyby $g\in(f)$, to $g=qf$ dla niezerowego $q$,
a więc $\deg g\ge\deg f$, sprzeczność.
Ideał $(g)$ jest właściwy, bo $g$ nie jest jednostką.
W takim razie $(f)$ nie jest maksymalny.
\end{proof}

\begin{przyklad}[Konstrukcja ciała liczb zespolonych]
Wielomian $x^2+1\in\R[x]$ jest nierozkładalny nad $\R$:
gdyby rozkładał się na dwa wielomiany liniowe,
miałby rzeczywisty pierwiastek, a $a^2+1>0$ dla każdego $a\in\R$.
Dlatego $\R[x]/(x^2+1)$ jest ciałem.
Jeśli $i=x+(x^2+1)$, to $i^2=-1$, a każdą klasę można zapisać
jednoznacznie jako $a+bi$ dla $a,b\in\R$,
ponieważ dzielenie przez $x^2+1$ zostawia resztę stopnia co najwyżej $1$.
Działania zgadzają się z działaniami w $\C$,
więc $\R[x]/(x^2+1)\cong\C$.
\end{przyklad}

\section{Wiele zmiennych i pierwsze pierścienie współrzędnych}

\subsection{Podstawianie, ideały punktów i usuwanie zmiennych}

Niech $k$ będzie ciałem. Pierścień $k[x_1,\dots,x_n]$
można skonstruować iteracyjnie:
$k[x_1,\dots,x_n]=k[x_1,\dots,x_{n-1}][x_n]$.
Elementami są skończone sumy jednomianów
$c_{\alpha}x_1^{\alpha_1}\cdots x_n^{\alpha_n}$,
gdzie $\alpha\in\Z_{\ge0}^n$ i $c_\alpha\in k$.
Równość wielomianów oznacza równość współczynników.

\begin{stwierdzenie}[Podstawienie do punktu]
\label{st:eval-point}
Dla $a=(a_1,\dots,a_n)\in k^n$ homomorfizm
$\ev_a\colon k[x_1,\dots,x_n]\to k$ zadany przez
$\ev_a(f)=f(a_1,\dots,a_n)$ jest surjektywny, a jego jądro wynosi
$\mathfrak m_a=(x_1-a_1,\dots,x_n-a_n)$.
W szczególności $\mathfrak m_a$ jest ideałem maksymalnym.
\end{stwierdzenie}
\begin{proof}
Homomorfizm jest surjektywny, bo każda stała $c\in k$
jest obrazem stałego wielomianu $c$. Każdy generator $x_j-a_j$
przechodzi na zero, więc
$\mathfrak m_a\subseteq\kerf\ev_a$.

Dla odwrotnej inkluzji weźmy wielomian $f$ spełniający $f(a)=0$.
Kolejno zastępujmy zmienne ich wartościami:
\begin{align*}
 f(x_1,\dots,x_n)-f(a_1,\dots,a_n)
 &=\bigl[f(x_1,\dots,x_n)-f(a_1,x_2,\dots,x_n)\bigr]\\
 &\quad+\bigl[f(a_1,x_2,\dots,x_n)-f(a_1,a_2,x_3,\dots,x_n)\bigr]
 +\dots\\
 &\quad+\bigl[f(a_1,\dots,a_{n-1},x_n)-f(a_1,\dots,a_n)\bigr].
\end{align*}
Każdy nawias jest podzielny przez odpowiedni $x_j-a_j$:
to wzór $x_j^m-a_j^m=(x_j-a_j)\sum_{r=0}^{m-1}x_j^{m-1-r}a_j^r$
zastosowany współczynnik po współczynniku w pozostałych zmiennych.
Zatem $f-f(a)\in\mathfrak m_a$. Skoro $f(a)=0$,
mamy $f\in\mathfrak m_a$. Otrzymaliśmy równość jądra i ideału.
Ponieważ iloraz przez jądro jest izomorficzny z ciałem $k$,
ideał $\mathfrak m_a$ jest maksymalny na mocy
 twierdzenia~\ref{tw:kryterium}.
\end{proof}

\begin{uwaga}[Co wynika, a co jeszcze nie wynika?]
Udowodniliśmy, że każdemu punktowi $a\in k^n$ odpowiada ideał
maksymalny $\mathfrak m_a$. Nie udowodniliśmy odwrotnego zdania,
że każdy ideał maksymalny ma tę postać. Jest ono fałszywe dla
niektórych ciał, np. nad $\R$ ideał $(x^2+1)\subseteq\R[x]$
jest maksymalny, lecz nie pochodzi od żadnego punktu $\R$.
Dla ciała algebraicznie domkniętego odwrotne zdanie jest częścią
słabego twierdzenia Hilberta o zerach, którym zajmiemy się
w rozdziale poświęconym zbiorom algebraicznym.
\end{uwaga}

\begin{przyklad}[Parabola i jej pierścień]
\label{pr:parabola}
Rozważmy homomorfizm
$\Phi\colon k[x,y]\to k[t]$ zadany przez
$\Phi(f(x,y))=f(t,t^2)$.
Jest on surjektywny: $t=\Phi(x)$, więc każdy wielomian w $t$
jest obrazem odpowiedniego wielomianu w $x$.
Ponadto $\Phi(y-x^2)=t^2-t^2=0$, czyli
$(y-x^2)\subseteq\kerf\Phi$.

Aby udowodnić inkluzję odwrotną, dla dowolnego
$f(x,y)=\sum_{j=0}^m a_j(x)y^j$ zapiszmy
\[
 f(x,y)-f(x,x^2)
 =\sum_{j=1}^m a_j(x)(y^j-x^{2j})
 =(y-x^2)\sum_{j=1}^m a_j(x)
   \sum_{r=0}^{j-1}y^{j-1-r}x^{2r}.
\]
Jeżeli $\Phi(f)=0$, to $f(t,t^2)=0$ jako wielomian w $k[t]$,
a więc również $f(x,x^2)=0$ jako wielomian w $k[x]$
(oznaczenie zmiennej nie ma znaczenia). Z powyższego wzoru
wynika $f\in(y-x^2)$. Zatem $\kerf\Phi=(y-x^2)$.
Pierwsze twierdzenie o izomorfizmie daje
\[
 k[x,y]/(y-x^2)\cong k[t].
\]
Algebraicznie oznacza to, że po narzuceniu relacji $y=x^2$
współrzędna $y$ jest już określona przez $x$: każda klasa wielomianów
ma reprezentanta zależnego tylko od $x$.
Nad nieskończonym ciałem $k$ izomorfizm można też interpretować
bezpośrednio jako opis funkcji wielomianowych na punktach paraboli.
Nad ciałem skończonym trzeba zachować rozróżnienie z
uwagi~\ref{uw:formalna-funkcja}.
\end{przyklad}

\begin{center}
\begin{tikzpicture}[scale=0.86,>=Stealth]
\draw[->,gray!80] (-2.6,0)--(2.7,0) node[right] {$x$};
\draw[->,gray!80] (0,-0.25)--(0,4.9) node[above] {$y$};
\draw[blue!70!black,very thick,domain=-2.12:2.12,samples=100]
  plot (\x,{\x*\x});
\fill[red!75!black] (1.35,1.8225) circle (2pt);
\draw[densely dashed,red!75!black] (1.35,0)--(1.35,1.8225)--(0,1.8225);
\node[below,red!75!black] at (1.35,0) {$t$};
\node[left,red!75!black] at (0,1.8225) {$t^2$};
\node[blue!70!black,anchor=west] at (1.25,4.1) {$y=x^2$};
\end{tikzpicture}
\end{center}
Punkt paraboli opisujemy jedną liczbą $t$ poprzez parametryzację
$t\mapsto(t,t^2)$. Obliczenie jądra $\Phi$ jest algebraicznym
odpowiednikiem tego rysunku, ale działa nad \emph{dowolnym} ciałem $k$,
również wtedy, gdy nie możemy rysować jego punktów na płaszczyźnie.

\begin{przyklad}[Dwie osie i iloraz przez $(xy)$]
W $k[x,y]$ rozważmy ideał $(xy)$.
Równanie $xy=0$ opisuje sumę dwóch osi: $x=0$ lub $y=0$.
Pierścień $k[x,y]/(xy)$ pamięta, że iloczyn współrzędnych $x$ i $y$
ma być zerem, choć żadna z tych współrzędnych nie musi być zerem
jako element tego pierścienia.
W istocie ich klasy są niezerowe. Gdyby $x=xyh$ dla pewnego
$h\in k[x,y]$, to po skróceniu niezerowego $x$ w dziedzinie
$k[x,y]$ otrzymalibyśmy $1=yh$, co po podstawieniu $y=0$
dawałoby sprzeczność $1=0$ w ciele $k$. Zatem $x\notin(xy)$;
analogicznie $y\notin(xy)$. Jednocześnie
$(x+(xy))(y+(xy))=xy+(xy)=0+(xy)$.
Iloraz ma zatem dzielniki zera. Jest to algebraiczny ślad faktu,
że zbiór geometryczny składa się z dwóch składowych.
\end{przyklad}

\begin{center}
\begin{tikzpicture}[scale=0.87,>=Stealth]
\draw[gray!65,->] (-2.4,0)--(2.5,0) node[right] {$x$};
\draw[gray!65,->] (0,-2.2)--(0,2.5) node[above] {$y$};
\draw[blue!70!black,line width=2.2pt] (-2.1,0)--(2.1,0);
\draw[red!75!black,line width=2.2pt] (0,-1.95)--(0,2.05);
\fill[black] (0,0) circle (2.3pt);
\node[blue!70!black,below right] at (1.25,0) {$y=0$};
\node[red!75!black,above right] at (0,1.2) {$x=0$};
\node[align=left,font=\small] at (3.9,-1.3)
  {Przekrój rzeczywisty\\zbioru $V(xy)$.};
\end{tikzpicture}
\end{center}

\begin{przyklad}[Pierścień wielomianów w dwóch zmiennych nie jest pierścieniem ideałów głównych]
Przypuśćmy, że w $k[x,y]$ ideał $(x,y)$ jest równy $(f)$
dla pewnego $f\in k[x,y]$. Wtedy $f\mid x$ i $f\mid y$.
Pokażemy, że każdy wspólny dzielnik $x$ i $y$ jest jednostką.
Jeżeli $x=fg$ w $k[x,y]$, traktujemy wszystkie wielomiany
jako wielomiany w $x$ o współczynnikach w dziedzinie $k[y]$.
Stopnie względem $x$ sumują się, więc jeden z $f,g$ ma stopień
względem $x$ równy zero. Jeżeli $f\in k[y]$, to z $x=fg$
współczynnik przy $x$ pokazuje, że $f$ dzieli $1$ w $k[y]$,
więc $f$ jest jednostką. W przeciwnym razie $g\in k[y]$
i ten sam argument pokazuje, że $g$ jest jednostką;
wtedy $f$ jest stowarzyszony z $x$.
W rezultacie każdy dzielnik $x$ jest jednostką albo jest
stowarzyszony z $x$.

Gdyby $f$ nie był jednostką, musiałby być stowarzyszony z $x$,
ale z $f\mid y$ wynikałoby wtedy $x\mid y$, co jest niemożliwe:
po podstawieniu $x=0$ równanie $y=xh(x,y)$ dawałoby $y=0$
w $k[y]$. Zatem $f$ jest jednostką, skąd $(f)=k[x,y]$.
Tymczasem $(x,y)$ jest ideałem właściwym, ponieważ
wszystkie jego elementy przyjmują wartość zero w $(0,0)$,
a wielomian $1$ nie. Sprzeczność dowodzi, że $(x,y)$ nie jest główny.
\end{przyklad}

\section{Czego uczą nas przykłady geometryczne?}

Przyjrzyjmy się trzem ilorazom pierścienia $k[x,y]$:
\begin{enumerate}[label=\arabic*)]
\item $k[x,y]/(y-x^2)$ odpowiada paraboli. Jest izomorficzny z $k[t]$,
więc jest dziedziną: ideał $(y-x^2)$ jest pierwszy.
Gdy $k$ jest algebraicznie domknięte, oznacza to też
nierozkładalność paraboli w topologii Zariskiego; dokładne uzasadnienie
podamy w następnym rozdziale. Nie wyciągamy tego wniosku bez założeń
o ciele: nad ciałem skończonym topologia Zariskiego na zbiorze
punktów $k^n$ jest dyskretna.
\item $k[x,y]/(xy)$ odpowiada sumie osi. Ma dzielniki zera,
bo klasy $x$ i $y$ są niezerowe, ale ich iloczyn jest zerem.
\item $k[t]/(t^2)$ ma niezerowy nilpotent. W odróżnieniu od
$k[t]/(t)$ przechowuje informację o „pogrubieniu” punktu $t=0$.
\end{enumerate}
Te obserwacje prowadzą do trzech idei, które rozwinie kolejny rozdział.

\textbf{Pierwsza idea:} badamy nie tylko rozwiązania $f=0$, ale cały
ideał równań, ponieważ kombinacje równań również muszą znikać.

\textbf{Druga idea:} iloraz pierścienia wielomianów przez ideał równań
koduje algebraiczne funkcje na rozpatrywanym obiekcie.

\textbf{Trzecia idea:} własności ideału odzwierciedlają geometrię.
Ideał pierwszy daje iloraz bez dzielników zera,
a pierwiastek ideału kontroluje, czy iloraz zawiera nilpotenty.
Dopiero nad ciałem algebraicznie domkniętym połączymy pierwszość
ideału z nierozkładalnością geometryczną odpowiadającego mu zbioru.

\begin{uwaga}[Granica dotychczasowej teorii]
W tym rozdziale nie dowodziliśmy jeszcze, że dla dowolnego ideału
$I\subseteq k[x_1,\dots,x_n]$ odpowiednie punkty geometryczne
pozwalają odzyskać $\sqrt I$. Takie twierdzenie wymaga
założenia, że $k$ jest algebraicznie domknięte, i stanowi treść
silnego twierdzenia Hilberta o zerach. Będzie to główny cel
pierwszego rozdziału właściwej geometrii algebraicznej.
\end{uwaga}

\section{Podsumowanie logicznych zależności}

\begin{center}
\begin{tikzpicture}[>=Stealth,
  box/.style={draw=gray!65,rounded corners=3pt,align=center,inner sep=6pt,fill=gray!5},
  node distance=9mm and 14mm]
\node[box] (r) {pierścień $A$};
\node[box,below left=of r] (i) {ideał $I\subseteq A$};
\node[box,below right=of r] (h) {homomorfizm\\$\varphi\colon A\to B$};
\node[box,below=of i] (q) {iloraz $A/I$};
\node[box,below=of h] (k) {jądro $\kerf\varphi$};
\node[box,below=of q] (p) {ideał pierwszy\\$\Longleftrightarrow A/I$ dziedzina};
\node[box,below=of k] (m) {ideał maksymalny\\$\Longleftrightarrow A/I$ ciało};
\draw[->] (r)--(i);
\draw[->] (r)--(h);
\draw[->] (i)--(q);
\draw[->] (h)--(k);
\draw[->] (k)--(q);
\draw[->] (q)--(p);
\draw[->] (q)--(m);
\end{tikzpicture}
\end{center}

Ten diagram nie oznacza, że każdy ideał jest pierwszy lub maksymalny.
Dwie dolne ramki podają \emph{kryteria} dla ideałów, które mają
daną własność. Diagram przypomina natomiast przepływ pojęć:
ideał pozwala utworzyć iloraz, a typ ilorazu rozstrzyga o typie ideału.

W dalszej części skryptu przejdziemy do zbiorów algebraicznych
w $k^n$, topologii Zariskiego, korespondencji między zbiorami
a ideałami oraz twierdzenia Hilberta o zerach. Wyjaśnimy też,
czym różni się zbiór punktów nad $k$ od jego pierścienia współrzędnych. Następnie
omówimy geometrię rzutową, osobliwości, normalizację i krzywe
algebraiczne wraz z rodzajem oraz twierdzeniem Riemanna--Rocha.

\section{Jednoznaczność rozkładu wielomianów na czynniki}
\label{sec:ufd-preparation}

W dalszej geometrii będziemy używać informacji silniejszej niż
algorytm Euklidesa dla $k[t]$: wielomiany w \emph{wielu} zmiennych
również mają jednoznaczny rozkład na czynniki nierozkładalne.
Trzeba to uzasadnić; nie wynika to automatycznie z faktu,
że $k[x,y]$ jest dziedziną.

\begin{definicja}[Pierścień z jednoznacznym rozkładem]
Niech $A$ będzie dziedziną całkowitości. Elementy $a,b\in A$
są \emph{stowarzyszone}, gdy $a=ub$ dla pewnej jednostki
$u\in A^\times$. Dziedzinę $A$ nazywamy \emph{pierścieniem
z jednoznacznym rozkładem na czynniki} (ang. \emph{UFD}), jeżeli
każdy niezerowy nieodwracalny element jest skończonym iloczynem
elementów nierozkładalnych, a dwa takie rozkłady różnią się
tylko kolejnością czynników i ich stowarzyszeniem.
\end{definicja}

\begin{lemat}[Lemat Gaussa]
\label{lem:gauss-ufd}
Niech $A$ będzie pierścieniem z jednoznacznym rozkładem,
$K=\Frac(A)$, a $f,g\in A[T]$. Wielomian jest
\emph{pierwotny}, gdy żaden nieodwracalny element $A$
nie dzieli wszystkich jego współczynników.
Wówczas: (i) iloczyn dwóch pierwotnych wielomianów jest
pierwotny; (ii) jeśli pierwotny $f$ rozkłada się
w $K[T]$ na czynniki dodatnich stopni, to rozkłada się
także w $A[T]$ na czynniki dodatnich stopni.
\end{lemat}
\begin{proof}
(i) Gdyby iloczyn $fg$ nie był pierwotny, to wszystkie
jego współczynniki dzieliłby pewien element nierozkładalny
$p\in A$. W pierścieniu z jednoznacznym rozkładem każdy
element nierozkładalny jest pierwszy: jeśli $p\mid ab$,
porównanie rozkładów na czynniki wymusza $p\mid a$ albo
$p\mid b$. Zatem $A/(p)$ jest dziedziną.
Po redukcji współczynników modulo $p$ niezerowe
wielomiany $\bar f,\bar g\in(A/(p))[T]$ miałyby
iloczyn zero, co przeczy temu, że pierścień
wielomianów nad dziedziną jest dziedziną.

(ii) Załóżmy $f=uv$ dla $u,v\in K[T]$ o dodatnich stopniach.
Usuwając mianowniki i wyciągając wspólne czynniki
współczynników, możemy przedstawić
$u=\alpha u_0$, $v=\beta v_0$, gdzie
$\alpha,\beta\in K^\times$, a $u_0,v_0\in A[T]$
są pierwotne. Z (i) iloczyn $u_0v_0$ jest pierwotny.
Mamy więc $f=(\alpha\beta)u_0v_0$.
Jeśli iloczyn dwóch pierwotnych wielomianów jest mnożony
przez ułamek $a/b\in K$ w postaci nieskracalnej,
a wynik ma współczynniki w $A$, to $b$ musi być jednostką:
w przeciwnym razie pewien czynnik nierozkładalny $p\mid b$
dzieliłby wszystkie współczynniki $a(u_0v_0)$;
przez $p\nmid a$ dzieliłby wszystkie współczynniki
$u_0v_0$, sprzeczność. Po wykazaniu, że $b$ jest jednostką, skalar $a/b$ należy do $A$.
Gdyby nie był jednostką, miałby w $A$ dzielnik nierozkładalny
$p$, który dzieliłby wszystkie współczynniki $f=(a/b)u_0v_0$.
Przeczyłoby to pierwotności $f$. Zatem $a/b\in A^\times$.
W konsekwencji $f$ jest iloczynem dwóch wielomianów
w $A[T]$ dodatnich stopni.
\end{proof}

\begin{twierdzenie}[Pierścienie wielomianów są UFD]
\label{tw:polynomial-UFD}
Jeśli $k$ jest ciałem i $n\ge1$, to
$k[x_1,\dots,x_n]$ jest pierścieniem z jednoznacznym
rozkładem na czynniki. Bardziej ogólnie, jeśli $A$
jest UFD, to $A[T]$ jest UFD.
\end{twierdzenie}
\begin{proof}
Najpierw udowodnimy zdanie ogólne. Dla niezerowego
$f\in A[T]$ wyznaczmy \emph{zawartość} $c(f)$, czyli
największy wspólny dzielnik jego niezerowych współczynników,
określony z dokładnością do jednostki. Dzięki istnieniu
i jednoznaczności rozkładu współczynników w $A$ taki
element istnieje. Mamy $f=c(f)f_0$ z $f_0$ pierwotnym.
Lemat Gaussa pokazuje, że iloczyn dwóch pierwotnych
wielomianów jest pierwotny, a zatem
$c(fg)$ jest stowarzyszony z $c(f)c(g)$.

W ciele ułamków $K$ pierścień $K[T]$ ma jednoznaczny rozkład
na czynniki nierozkładalne. Uzasadnijmy to na podstawie
wcześniejszych wyników. Algorytm dzielenia z resztą
(twierdzenie~\ref{tw:dzielenie}) pokazuje, że każdy ideał
$K[T]$ jest główny (stwierdzenie~\ref{st:kx-glowne}).
Jeśli nierozkładalny $p$ dzieli $ab$, lecz nie dzieli $a$,
to największy wspólny dzielnik $p,a$ jest równy $1$
z dokładnością do jednostki. Tożsamość Bézouta
(twierdzenie~\ref{tw:bezout-pol}) daje $up+va=1$.
Po pomnożeniu przez $b$ otrzymujemy $upb+vab=b$;
oba składniki po lewej są podzielne przez $p$, zatem $p\mid b$.
Każdy nierozkładalny jest więc pierwszy. Rozkład niezerowego
wielomianu niebędącego jednostką istnieje przez indukcję względem
stopnia: wielomian nierozkładalny kończy postępowanie, a rozkładalny
jest iloczynem dwóch wielomianów dodatnich, ściśle mniejszych stopni.
Jednoznaczność wynika z pierwszości czynników: pierwszy czynnik
jednego rozkładu dzieli któryś czynnik drugiego, jest z nim
stowarzyszony, po czym można go skrócić i powtórzyć rozumowanie.
Jeżeli $f_0$ jest stały, to z pierwotności wynika,
że jest jednostką $A$, więc nie wymaga dalszego rozkładu.
W przeciwnym razie niech $f_0=q_1\cdots q_r$ będzie
rozkładem pierwotnego wielomianu $f_0$ na nierozkładalne
czynniki w $K[T]$
(z ewentualnym czynnikiem stałym włączonym do $q_1$).
Dla każdego $i$ usuwamy mianowniki współczynników $q_i$,
a następnie dzielimy przez ich największy wspólny dzielnik.
Otrzymujemy $q_i=\alpha_i h_i$, gdzie
$\alpha_i\in K^\times$, a $h_i\in A[T]$ jest pierwotny.
Z pierwszej części lematu Gaussa iloczyn $H=h_1\cdots h_r$
jest pierwotny. Ponieważ $f_0=(\prod_i\alpha_i)H$
i oba wielomiany należą do $A[T]$ oraz są pierwotne,
argument o nieskracalnym mianowniku z dowodu lematu Gaussa
pokazuje, że $\prod_i\alpha_i$ jest jednostką $A$.
Włączamy tę jednostkę do $h_1$ i uzyskujemy rozkład $f_0$
w $A[T]$. Każdy $h_i$ jest nierozkładalny w $A[T]$:
gdyby $h_i=ab$ z niejednostkami $a,b\in A[T]$, to
pierwotność $h_i$ i wzór na zawartość wymuszałyby
pierwotność obu czynników. Żaden z nich nie może być
nieodwracalną stałą, a więc oba miałyby dodatni stopień.
Dawałoby to rozkład $q_i$ w $K[T]$, sprzeczny z jego
nierozkładalnością. Istnienie rozkładu w $A[T]$ wynika
teraz z rozkładu zawartości $c(f)$ w $A$ i powyższego
rozkładu pierwotnej części $f_0$.

Dla jednoznaczności porównajmy dwa rozkłady $f$.
Po przejściu do $K[T]$ elementy stałe z $A$ stają się
jednostkami, a pierwotne czynniki nierozkładalne
pozostają nierozkładalne na mocy lematu Gaussa.
Jednoznaczność w $K[T]$ pozwala dopasować czynniki
wielomianowe parami z dokładnością do skalarów z $K$.
Uzasadnijmy dokładnie, że dwa pierwotne wielomiany
$h,h'\in A[T]$, dla których $h'=\lambda h$ z
$\lambda\in K^\times$, różnią się jednostką $A$.
Zapiszmy $\lambda=a/b$, gdzie $a,b\in A$ są względnie
pierwsze (nie mają wspólnego czynnika nierozkładalnego).
Z równości $bh'=ah$ wynika, że gdyby nierozkładalny
$p$ dzielił $b$, to przez pierwszość $p$ i $p\nmid a$
dzieliłby każdy współczynnik $h$, wbrew pierwotności.
Zatem $b$ jest jednostką, czyli $\lambda\in A$.
Gdyby $\lambda$ nie był jednostką, pewien nierozkładalny
dzielnik $\lambda$ dzieliłby wszystkie współczynniki
$h'$, znów wbrew pierwotności. Stąd
$\lambda\in A^\times$.

W każdym rozkładzie $f$ czynniki stałe nierozkładalne
pochodzą, z dokładnością do jednostek, z rozkładu
zawartości $c(f)$ w $A$, gdyż zawartość iloczynu
pierwotnych czynników jest jednostką.
Jednoznaczność rozkładu $c(f)$ w $A$ oraz wykazane
dopasowanie pierwotnych czynników wielomianowych
kończą dowód jednoznaczności w $A[T]$.

Ciało $k$ jest UFD w sposób trywialny, bo każdy jego
niezerowy element jest jednostką. Stosując wykazane
zdanie kolejno do $k[x_1]$, $k[x_1][x_2]$ itd.,
otrzymujemy tezę dla dowolnego $n$.
\end{proof}

\begin{uwaga}[Dlaczego potrzebujemy tego twierdzenia?]
W pierścieniu $k[X,Y,Z]$ będziemy korzystać z tego,
że jeśli wielomiany $F,G$ nie mają wspólnego czynnika
nierozkładalnego, a $F\mid Gh$, to $F\mid h$.
Rzeczywiście, gdy $F\ne0$, rozkładamy $F$ na czynniki
nierozkładalne. Każdy z nich jest pierwszy, ponieważ
pierścień jest UFD. Żaden nie dzieli $G$, więc z podzielności
$F\mid Gh$ wynika, że dzieli $h$ z co najmniej taką samą
krotnością, z jaką występuje w $F$: usuwamy kolejno jego
wystąpienia z iloczynu $Gh$, korzystając za każdym razem
z pierwszości i z faktu, że nie dzieli $G$.
Jednoznaczność rozkładu daje $F\mid h$.
Dla $F=0$ założenie $0\mid Gh$ oznacza $Gh=0$;
ponieważ $F$ i $G$ nie mają wspólnego czynnika
nierozkładalnego, $G\ne0$, zatem $h=0$ i także $0\mid h$.
Podobne rozumowania są podstawą dowodu twierdzenia Bézouta.
\end{uwaga}


\chapter{Geometria afiniczna: równania, ideały i twierdzenie Hilberta o zerach}
\label{ch:afiniczna}
\section{Od układów równań do geometrii}

W geometrii elementarnej opisujemy prostą równaniem $y=mx+b$, parabolę
równaniem $y=x^2$, a okrąg równaniem $x^2+y^2=1$. Geometria algebraiczna
bada analogiczne zbiory w przestrzeniach o dowolnej skończonej liczbie
współrzędnych. Różnica polega na tym, że interesują nas nie tylko
rysunki rozwiązań, ale także wszystkie wielomiany, które na nich znikają,
i funkcje, które można na tych zbiorach określić algebraicznie.

Odpowiemy kolejno na pytania:
\begin{enumerate}[label=\arabic*)]
\item Kiedy dwa układy równań wielomianowych mają ten sam zbiór rozwiązań?
\item Jaką topologię otrzymamy, uznając zbiory rozwiązań za zbiory domknięte?
\item Jak odzyskać informacje o równaniach ze zbioru ich wspólnych zer?
\item Dlaczego odpowiedź zmienia się wraz z ciałem współczynników?
\item Jak przekształcenia geometryczne przekładają się na homomorfizmy pierścieni?
\end{enumerate}

\begin{przyklad}[Te same punkty, różne równania]
W $\C^2$ układy $x=0$ oraz $x^2=0$ mają ten sam zbiór
rozwiązań: oś $y$. Również układy $x=0$ oraz $x^2=0,\ xy=0$
mają ten sam zbiór rozwiązań. Jednak ideały $(x)$, $(x^2)$ oraz
$(x^2,xy)$ nie są jednakowe. Przykładowo w ilorazie
$\C[x,y]/(x^2)$ klasa $x$ jest niezerowa, ale jej kwadrat jest zerowy.
Zatem sam zbiór rozwiązań nie odróżnia wszystkich relacji algebraicznych.
To pierwsza wskazówka, że potrzebujemy jednocześnie zbiorów i ideałów.
\end{przyklad}

\begin{uwaga}[Założenia o ciele]
Przez większość rozdziału $k$ oznacza dowolne \emph{ciało}; jego
charakterystyka może być dodatnia. Przy twierdzeniu Hilberta o zerach
przyjmiemy dodatkowo, że $k$ jest \emph{algebraicznie domknięte}.
Za każdym razem wyraźnie wskażemy, które założenie jest potrzebne.
Wszystkie pierścienie są przemienne z jedynką, a homomorfizmy zachowują jedynkę.
\end{uwaga}

\section{Przestrzeń afiniczna i zbiory algebraiczne}

\begin{definicja}[Przestrzeń afiniczna]
Niech $k$ będzie ciałem, a $n\ge1$ liczbą całkowitą. \emph{Przestrzenią
afiniczną wymiaru $n$ nad $k$}, na poziomie zbioru jej punktów nad $k$,
nazywamy $\A^n(k):=k^n$. Punkt zapisujemy jako
$a=(a_1,\dots,a_n)$, gdzie $a_i\in k$.
\end{definicja}

Symbol $\A^n(k)$ nie oznacza tu jeszcze schematu ani przestrzeni z nadaną
topologią. W tej chwili jest to zbiór uporządkowanych $n$-tek liczb z $k$.
Topologię wprowadzimy dopiero za pomocą równań.

\begin{definicja}[Zbiór zerowy]
Niech $k$ będzie ciałem, $n\ge1$ i $S\subseteq k[x_1,\dots,x_n]$.
Definiujemy \emph{zbiór wspólnych zer wielomianów z $S$} przez
\[
\V_k(S):=\{a\in k^n: f(a)=0\text{ dla każdego }f\in S\}.
\]
Gdy $S=\{f_1,\dots,f_r\}$, piszemy
$\V_k(f_1,\dots,f_r)$, a dla pojedynczego $f$ piszemy $\V_k(f)$.
Podzbiór $X\subseteq k^n$ nazywamy \emph{zbiorem algebraicznym afinicznym},
jeśli istnieje $S\subseteq k[x_1,\dots,x_n]$ takie, że $X=\V_k(S)$.
\end{definicja}

Dopuszczamy $S=\varnothing$. Wtedy każdy punkt spełnia pusty zestaw
równań i $\V_k(\varnothing)=k^n$. Z drugiej strony
$\V_k(1)=\varnothing$, gdyż żadna liczba nie spełnia $1=0$.
Możemy już sprawdzić, dlaczego od listy równań ważniejszy jest ideał.

\begin{stwierdzenie}[Zbiór zerowy zależy tylko od generowanego ideału]
\label{st:V-ideal}
Dla dowolnego ciała $k$, liczby $n\ge1$ i zbioru
$S\subseteq k[x_1,\dots,x_n]$ zachodzi $\V_k(S)=\V_k((S))$,
gdzie $(S)$ jest ideałem generowanym przez $S$.
\end{stwierdzenie}
\begin{proof}
Ponieważ $S\subseteq(S)$, każdy punkt zerujący wszystkie elementy
$(S)$ zeruje też elementy $S$, zatem $\V_k((S))\subseteq\V_k(S)$.
Dla inkluzji odwrotnej ustalmy $a\in\V_k(S)$ i weźmy dowolny
$h\in(S)$. Z definicji ideału generowanego
$h=\sum_{j=1}^r g_js_j$, gdzie $g_j\in k[x_1,\dots,x_n]$,
$s_j\in S$; suma jest skończona. Po podstawieniu punktu $a$ mamy
$h(a)=\sum_jg_j(a)s_j(a)=0$, ponieważ każde $s_j(a)=0$.
Stąd $a\in\V_k((S))$.
\end{proof}

\begin{przyklad}[Zbiór zerowy może być pusty mimo nietrywialnych równań]
W $\R^2$ zbiór $\V_{\R}(x^2+y^2+1)$ jest pusty, bo suma
$x^2+y^2+1$ jest dodatnia w każdym punkcie rzeczywistym.
Tymczasem w $\C^2$ ten sam wielomian ma zera, np. $(i,0)$.
Zmiana ciała współczynników może więc zmienić geometrię zbioru zer.
\end{przyklad}

\begin{przyklad}[Przecięcie dwóch krzywych]
W $\R^2$ układ $y=x^2$ i $y=2-x$ jest równoważny
$y=x^2$, $x^2+x-2=0$. Ponieważ
$x^2+x-2=(x-1)(x+2)$, otrzymujemy dokładnie dwa punkty:
$(1,1)$ i $(-2,4)$. Zapisujemy to jako
$\V_{\R}(y-x^2,y+x-2)=\{(1,1),(-2,4)\}$.
Ten przykład pokazuje, że \emph{przecięcie} zbiorów rozwiązań odpowiada
\emph{połączeniu} ich równań.
\end{przyklad}

\begin{center}
\begin{tikzpicture}[scale=.85,>=Stealth]
\draw[->,gray!75] (-2.9,0)--(2.9,0) node[right] {$x$};
\draw[->,gray!75] (0,-.35)--(0,4.9) node[above] {$y$};
\draw[blue!75!black,line width=1.35pt,domain=-2.15:2.15,samples=100] plot (\x,{\x*\x});
\draw[red!75!black,line width=1.25pt,domain=-2.7:2.35] plot (\x,{2-\x});
\fill (1,1) circle(2pt);
\node[anchor=west,fill=white,inner sep=1pt] at (1.37,1.2) {$(1,1)$};
\fill (-2,4) circle(2pt) node[above right,fill=white,inner sep=1pt] {$(-2,4)$};
\node[blue!75!black] at (1.75,4.65) {$y=x^2$};
\node[red!75!black,fill=white,inner sep=1pt] at (-1.05,3.4) {$y=2-x$};
\end{tikzpicture}
\end{center}
\begin{center}\small Rysunek 1.1. Przecięcie paraboli i prostej nad $\R$; zaznaczono wspólne zera ich równań.\end{center}

\section{Topologia Zariskiego: domknięte są rozwiązania równań}

\subsection{Dlaczego zbiory algebraiczne tworzą rodzinę zbiorów domkniętych?}

Aby zadać topologię na zbiorze $X$, można wskazać rodzinę
$\mathcal C$ jego zbiorów \emph{domkniętych}, pod warunkiem że
$\varnothing,X\in\mathcal C$, dowolne przecięcie zbiorów z
$\mathcal C$ należy do $\mathcal C$, a suma dwóch zbiorów z
$\mathcal C$ także należy do $\mathcal C$. Zbiory otwarte to dopełnienia
tych domkniętych. Sprawdzimy te warunki dla zbiorów zerowych.

\begin{twierdzenie}[Istnienie topologii Zariskiego]
\label{tw:zariski-top}
Niech $k$ będzie dowolnym ciałem i $n\ge1$.
Rodzina $\mathcal C=\{\V_k(S):S\subseteq k[x_1,\dots,x_n]\}$
jest rodziną zbiorów domkniętych pewnej topologii na $k^n$.
Topologię tę nazywamy \emph{topologią Zariskiego}.
Dla ideałów $I,J\subseteq k[x_1,\dots,x_n]$ zachodzą wzory
\[
\V_k(I+J)=\V_k(I)\cap\V_k(J),\qquad
\V_k(IJ)=\V_k(I)\cup\V_k(J).
\]
\end{twierdzenie}
\begin{proof}
Całe $k^n$ jest zbiorem $\V_k(0)$, a zbiór pusty jest równy
$\V_k(1)$.

Niech $\{S_\alpha\}_{\alpha\in T}$ będzie dowolną rodziną zbiorów
wielomianów. Punkt $a$ należy do wszystkich $\V_k(S_\alpha)$
dokładnie wtedy, gdy zeruje wszystkie wielomiany należące do
$\bigcup_{\alpha\in T}S_\alpha$. Dlatego
$\bigcap_{\alpha\in T}\V_k(S_\alpha)=
\V_k(\bigcup_{\alpha\in T}S_\alpha)$.
Z twierdzenia~\ref{st:V-ideal} wynika w szczególności wzór
$\V_k(I+J)=\V_k(I)\cap\V_k(J)$: suma ideałów jest ideałem
generowanym przez zbiór $I\cup J$.

Pozostaje suma dwóch zbiorów. Jeśli $a\in\V_k(I)$, to dla
każdych $f\in I$ i $g\in J$ mamy $f(a)g(a)=0$; każda
skończona suma iloczynów elementów z $I$ i $J$ także zeruje się w $a$.
To daje $\V_k(I)\cup\V_k(J)\subseteq\V_k(IJ)$.

Odwrotnie, niech $a\in\V_k(IJ)$ i załóżmy, że
$a\notin\V_k(I)$. Istnieje wtedy $f\in I$ takie, że
$f(a)\ne0$. Dla dowolnego $g\in J$ wielomian $fg$ należy do
$IJ$, a więc $0=(fg)(a)=f(a)g(a)$. Ponieważ $k$ jest ciałem
i $f(a)\ne0$, dostajemy $g(a)=0$. Dotyczy to każdego $g\in J$,
zatem $a\in\V_k(J)$. Dowiedliśmy odwrotnej inkluzji.
Dzięki domknięciu na dowolne przecięcia i skończone sumy
rodzina $\mathcal C$ jest rodziną zbiorów domkniętych topologii.
\end{proof}

\begin{uwaga}[Dlaczego użyliśmy iloczynu ideałów?]
Równość $\V_k(IJ)=\V_k(I)\cup\V_k(J)$ bierze się z prostego
faktu o liczbach: w ciele $uv=0$ oznacza $u=0$ lub $v=0$.
\emph{Suma} ideałów wymaga jednoczesnego spełnienia dwóch rodzin
równań, więc daje \emph{przecięcie} zbiorów. \emph{Iloczyn} ideałów
wymaga spełnienia co najmniej jednej rodziny, więc daje \emph{sumę}
zbiorów. To odwrócenie intuicji będzie powtarzać się w całym skrypcie.
\end{uwaga}

\begin{przyklad}[Dwie osie i ich punkt wspólny]
W $k[x,y]$ weźmy $I=(x)$ oraz $J=(y)$.
Wtedy $\V_k(I)$ to oś $y$, $\V_k(J)$ to oś $x$,
$\V_k(I+J)=\V_k(x,y)=\{(0,0)\}$, a
$\V_k(IJ)=\V_k(xy)$ jest sumą zbiorów obu osi.
\end{przyklad}

\begin{center}
\begin{tikzpicture}[scale=.82,>=Stealth]
\draw[gray!70,->] (-2.45,0)--(2.45,0) node[right] {$x$};
\draw[gray!70,->] (0,-2.2)--(0,2.4) node[above] {$y$};
\draw[blue!70!black,line width=2pt] (-2.05,0)--(2.05,0);
\draw[red!75!black,line width=2pt] (0,-1.95)--(0,2.0);
\fill[black] (0,0) circle (2.7pt);
\node[blue!70!black,below right,fill=white,inner sep=1pt] at (1.6,-.05) {$\V_k(y)$};
\node[red!75!black,above left,fill=white,inner sep=1pt] at (-.05,1.4) {$\V_k(x)$};
\end{tikzpicture}
\end{center}
\begin{center}\small Rysunek 1.2. Zbiór $\V_{\R}(xy)$: suma zbiorów dwóch przecinających się osi rzeczywistych.\end{center}

\subsection{Najprostsza topologia Zariskiego: prosta afiniczna}

\begin{stwierdzenie}[Zbiory domknięte na prostej nad ciałem nieskończonym]
\label{st:A1-top}
Jeśli $k$ jest nieskończonym ciałem, to domkniętymi zbiorami
Zariskiego w $\A^1(k)$ są dokładnie: zbiór pusty, całe $k$
oraz wszystkie skończone podzbiory $k$.
\end{stwierdzenie}
\begin{proof}
Niech $X=\V_k(S)$ będzie domknięty i różny od całego $k$.
Istnieje wtedy punkt $a\in k\setminus X$. Z definicji $X$
istnieje wielomian $f\in S$ taki, że $f(a)\ne0$;
w szczególności $f$ nie jest wielomianem zerowym.
Każdy punkt $X$ jest pierwiastkiem $f$, zatem $X$ jest skończony:
niezerowy wielomian jednej zmiennej nad ciałem ma nie więcej
pierwiastków niż wynosi jego stopień.

Odwrotnie, dla dowolnego skończonego zbioru
$X=\{a_1,\dots,a_r\}$ mamy
$X=\V_k(\prod_{j=1}^r(x-a_j))$, a dla $X=\varnothing$
można użyć wielomianu stałego $1$. Całe $k$ jest $\V_k(0)$.
\end{proof}

Zbiory otwarte niepuste na tej prostej są więc dopełnieniami
zbiorów skończonych. Dwa dowolne niepuste zbiory otwarte przecinają się,
ponieważ suma dwóch skończonych zbiorów jest skończona, a $k$ jest
nieskończone. W szczególności \emph{topologia Zariskiego na $\A^1(k)$
nie jest Hausdorffa}: nie potrafimy rozdzielić dwóch różnych punktów
rozłącznymi otoczeniami otwartymi.

\begin{przyklad}[Każdy punkt jest domknięty]
Dla dowolnego ciała $k$, liczby całkowitej $n\ge1$
i $a=(a_1,\dots,a_n)\in k^n$
zachodzi $\{a\}=\V_k(x_1-a_1,\dots,x_n-a_n)$.
Topologia Zariskiego spełnia zatem aksjomat $T_1$: wszystkie
jednopunktowe podzbiory są domknięte. Nie przeczy to temu, że
nad ciałem nieskończonym może nie spełniać aksjomatu Hausdorffa.
\end{przyklad}

\begin{przyklad}[Nad ciałem skończonym topologia jest dyskretna]
Niech $n\ge1$. Jeśli $k=\F_q$ ma $q$ elementów, to $k^n$ jest zbiorem skończonym.
Ponieważ każdy punkt jest domknięty, każdy podzbiór $k^n$
stanowi skończoną sumę punktów domkniętych, a zatem też jest domknięty.
W konsekwencji \emph{każdy podzbiór jest również otwarty}.
Topologia Zariskiego na $\A^n(\F_q)$ jest więc dyskretna.
To ostrzeżenie przed przenoszeniem bez zmian intuicji geometrycznej
z $\C$ na ciała skończone.
\end{przyklad}

\subsection{Zbiory podstawowo otwarte}

\begin{definicja}[Zbiór podstawowo otwarty]
Dla ciała $k$, liczby $n\ge1$ oraz $f\in k[x_1,\dots,x_n]$
definiujemy $D_k(f)=\{a\in k^n:f(a)\ne0\}=k^n\setminus\V_k(f)$.
\end{definicja}

\begin{stwierdzenie}[Baza topologii Zariskiego]
Niech $k$ będzie ciałem, $n\ge1$ liczbą całkowitą,
a $f,g\in k[x_1,\dots,x_n]$ wielomianami.
Rodzina $\{D_k(h):h\in k[x_1,\dots,x_n]\}$ jest bazą topologii
Zariskiego na $k^n$. Ponadto $D_k(1)=k^n$, $D_k(0)=\varnothing$
oraz $D_k(f)\cap D_k(g)=D_k(fg)$.
\end{stwierdzenie}
\begin{proof}
Równości $D_k(1)=k^n$ i $D_k(0)=\varnothing$ wynikają z definicji.
Punkt należy do $D_k(f)\cap D_k(g)$ dokładnie wtedy,
gdy $f(a)\ne0$ i $g(a)\ne0$. W ciele jest to równoważne
$(fg)(a)\ne0$, czyli $a\in D_k(fg)$.

Niech $U$ będzie dowolnym zbiorem otwartym.
Jego dopełnienie ma postać $\V_k(S)$ dla pewnego $S$.
Punkt $a$ należy do $U$ wtedy i tylko wtedy, gdy nie spełnia
wszystkich równań z $S$, czyli istnieje $f\in S$ z $f(a)\ne0$.
Zatem $U=\bigcup_{f\in S}D_k(f)$, jak wymaga definicja bazy.
\end{proof}

\begin{uwaga}[Funkcja $1/x$ i otwarta część prostej]
Zbiór $D_k(x)=k\setminus\{0\}$ jest otwarty w $\A^1(k)$.
Na tym zbiorze funkcja $x\mapsto1/x$ ma sens, mimo że nie jest
ograniczeniem wielomianu na całe $k$, gdy $k$ jest nieskończone.
To zapowiada pojęcie \emph{funkcji regularnej na zbiorze otwartym}:
lokalnie będziemy dopuszczać ilorazy wielomianów o nieznikającym
mianowniku. Formalny opis otrzymuje się przez lokalizację pierścienia;
przedstawimy prosty przypadek w dalszej części rozdziału.
\end{uwaga}

\section{Ideał znikania: od geometrii z powrotem do algebry}

\begin{definicja}[Ideał znikania]
Niech $k$ będzie ciałem, $n\ge1$ i $X\subseteq k^n$.
\emph{Ideałem znikania} $X$ nazywamy
\[
\I_k(X)=\{f\in k[x_1,\dots,x_n]:f(a)=0\text{ dla każdego }a\in X\}.
\]
Jeżeli $X=\varnothing$, warunek jest pusty, więc
$\I_k(\varnothing)=k[x_1,\dots,x_n]$.
\end{definicja}

\begin{twierdzenie}[Podstawowe własności ideału znikania]
\label{tw:I-basic}
Niech $k$ będzie ciałem, $n\ge1$ i $X\subseteq k^n$.
Wtedy $\I_k(X)$ jest ideałem radykalnym pierścienia
$k[x_1,\dots,x_n]$. Dla dowolnej rodziny podzbiorów
$\{X_\alpha\}_{\alpha\in T}$ zachodzi
$\I_k(\bigcup_\alpha X_\alpha)=\bigcap_\alpha\I_k(X_\alpha)$.
\end{twierdzenie}
\begin{proof}
Wielomian zerowy znika na $X$, zatem $0\in\I_k(X)$.
Jeżeli $f,g\in\I_k(X)$, to dla dowolnego $a\in X$ mamy
$(f-g)(a)=f(a)-g(a)=0$, więc $f-g\in\I_k(X)$.
Jeśli $h\in k[x_1,\dots,x_n]$, to $(hf)(a)=h(a)f(a)=0$,
a więc $hf\in\I_k(X)$. To dowodzi, że $\I_k(X)$ jest ideałem.

Przypuśćmy, że $f^m\in\I_k(X)$ dla pewnego $m\ge1$.
Dla każdego $a\in X$ zachodzi $f(a)^m=0$ w ciele $k$.
W ciele nie ma niezerowych elementów nilpotentnych, stąd
$f(a)=0$ dla każdego $a\in X$, czyli $f\in\I_k(X)$.
Zatem ideał jest radykalny.

Wreszcie $f$ znika na sumie $\bigcup_\alpha X_\alpha$
dokładnie wtedy, gdy znika na każdym $X_\alpha$ z osobna.
Jest to równoważne $f\in\bigcap_\alpha\I_k(X_\alpha)$.
\end{proof}

\begin{przyklad}[Ideał jednego punktu]
Niech $a=(a_1,\dots,a_n)\in k^n$. W rozdziale o pierścieniach
udowodniliśmy, że jądro homomorfizmu
$\ev_a:k[x_1,\dots,x_n]\to k$, $f\mapsto f(a)$,
wynosi $(x_1-a_1,\dots,x_n-a_n)$.
Ponieważ jądro to dokładnie $\I_k(\{a\})$, mamy
$\I_k(\{a\})=(x_1-a_1,\dots,x_n-a_n)$.
Jest to ideał maksymalny, bo iloraz jest izomorficzny z $k$.
\end{przyklad}

\begin{przyklad}[Rozwiązania układu równań i jego ideał]
Niech $X=\V_k(x,y)\subseteq k^2$. Wtedy $X=\{(0,0)\}$,
a z poprzedniego przykładu $\I_k(X)=(x,y)$.
Natomiast dla $J=(x^2,y^2)$ mamy również
$\V_k(J)=\{(0,0)\}$, lecz $J\ne(x,y)$.
W rzeczywistości $\sqrt J=(x,y)$: elementy $x$ i $y$
należą do $\sqrt J$; zatem $(x,y)\subseteq\sqrt J$.
W drugą stronę, jeśli $f^m\in J$, to $f(0,0)^m=0$,
więc $f(0,0)=0$, a stąd $f\in(x,y)$.
Widać, dlaczego z punktów możemy co najwyżej oczekiwać
odzyskania \emph{pierwiastnika} ideału.
\end{przyklad}

\subsection{Odwracanie inkluzji i operatory domknięcia}

\begin{twierdzenie}[Korespondencja porządku odwracająca inkluzje]
\label{tw:galois}
Niech $k$ będzie ciałem, $n\ge1$, a $R=k[x_1,\dots,x_n]$.
Dla zbiorów $S,T\subseteq R$ i podzbiorów $X,Y\subseteq k^n$
zachodzą implikacje
\begin{align*}
S\subseteq T&\ \Longrightarrow\ \V_k(T)\subseteq\V_k(S),\\
X\subseteq Y&\ \Longrightarrow\ \I_k(Y)\subseteq\I_k(X).
\end{align*}
Ponadto $S\subseteq\I_k(\V_k(S))$ oraz
$X\subseteq\V_k(\I_k(X))$.
W szczególności dla ideału $I\subseteq R$ i zbioru $X\subseteq k^n$
\[
\V_k(\I_k(\V_k(I)))=\V_k(I),\qquad
\I_k(\V_k(\I_k(X)))=\I_k(X).
\]
\end{twierdzenie}
\begin{proof}
Jeżeli $S\subseteq T$, to spełnienie wszystkich równań
z większego zbioru $T$ wymusza spełnienie równań z $S$.
Daje to pierwszą inkluzję. Jeśli $X\subseteq Y$, wielomian
zerujący się na większym zbiorze $Y$ zeruje się również na $X$;
tak otrzymujemy drugą inkluzję.

Każdy wielomian $f\in S$ znika na $\V_k(S)$ z definicji,
czyli $f\in\I_k(\V_k(S))$. Analogicznie każdy punkt $a\in X$
zeruje wszystkie wielomiany, które znikają na $X$, stąd
$a\in\V_k(\I_k(X))$.

Aby dowieść pierwszej z dwóch ostatnich równości, oznaczmy
$Z=\V_k(I)$ i zauważmy, że $I\subseteq\I_k(Z)$, więc
$\V_k(\I_k(Z))\subseteq\V_k(I)=Z$. Inkluzja odwrotna
$Z\subseteq\V_k(\I_k(Z))$ wynika z już udowodnionej własności.
Dowód drugiej równości jest analogiczny: z
$X\subseteq\V_k(\I_k(X))$ wynika
$\I_k(\V_k(\I_k(X)))\subseteq\I_k(X)$,
a odwrotna inkluzja wynika z własności dla $S=\I_k(X)$.
\end{proof}

\begin{center}
\begin{tikzpicture}[>=Stealth,node distance=4cm]
\node[draw,rounded corners,align=center,minimum width=4.7cm,minimum height=1.15cm] (left)
  {rodziny równań\\$I\subseteq k[x_1,\dots,x_n]$};
\node[draw,rounded corners,align=center,right=of left,minimum width=4.7cm,minimum height=1.15cm] (right)
  {zbiory punktów\\$X\subseteq k^n$};
\draw[->,thick,blue!70!black] (left.north east) to[bend left=20]
 node[above] {$I\mapsto\V_k(I)$} (right.north west);
\draw[->,thick,red!70!black] (right.south west) to[bend left=20]
 node[below] {$X\mapsto\I_k(X)$} (left.south east);
\node[below=1.4cm of left,align=center,font=\small] {więcej równań\\$\Longrightarrow$ mniej punktów};
\node[below=1.4cm of right,align=center,font=\small] {więcej punktów\\$\Longrightarrow$ mniej równań};
\end{tikzpicture}
\end{center}
\begin{center}\small Rysunek 1.3. Dwa przejścia odwracające inkluzję: $\V_k$ i $\I_k$.\end{center}

\begin{stwierdzenie}[Domknięcie Zariskiego dowolnego podzbioru]
\label{st:closure}
Dla dowolnego ciała $k$, liczby $n\ge1$ i podzbioru
$X\subseteq k^n$ jego domknięcie w topologii Zariskiego jest równe
$\overline X^{\mathrm{Zar}}=\V_k(\I_k(X))$.
\end{stwierdzenie}
\begin{proof}
Wiemy, że $X\subseteq\V_k(\I_k(X))$, a zbiór po prawej
jest domknięty w topologii Zariskiego.
Niech $Y=\V_k(S)$ będzie dowolnym zbiorem domkniętym zawierającym $X$.
Jeśli $f\in S$, to $f$ znika na $Y$, a więc także na $X$.
Stąd $S\subseteq\I_k(X)$, co odwracając inkluzję daje
$\V_k(\I_k(X))\subseteq\V_k(S)=Y$.
Zatem $\V_k(\I_k(X))$ jest najmniejszym zbiorem domkniętym
zawierającym $X$, czyli jego domknięciem.
\end{proof}

\begin{przyklad}[Nieskończony zbiór jest gęsty na prostej]
Jeżeli $k$ jest nieskończone, a $X\subseteq k$ jest nieskończony,
to $\I_k(X)=(0)$: niezerowy wielomian jednej zmiennej nie może
zerować się na nieskończenie wielu punktach.
Wobec tego $\overline X^{\mathrm{Zar}}=\V_k(0)=k$.
Na przykład $\Z\subseteq\C$ jest gęste w prostej afinicznej
$\A^1(\C)$ w topologii Zariskiego, mimo że w zwykłej topologii
zespolonej $\Z$ jest zbiorem domkniętym i dyskretnym.
\end{przyklad}

\section{Pierścień współrzędnych i algebraiczne funkcje na zbiorze}

\begin{definicja}[Pierścień współrzędnych]
Niech $k$ będzie ciałem, $n\ge1$ i $X\subseteq k^n$ zbiorem
algebraicznym. \emph{Pierścieniem współrzędnych} $X$ nazywamy
$k[X]:=k[x_1,\dots,x_n]/\I_k(X)$.
Klasa $x_i+\I_k(X)$ jest $i$-tą \emph{funkcją współrzędną na $X$}.
\end{definicja}

Konstrukcja jest motywowana następującym rachunkiem. Dwa wielomiany
$f,g$ są równe jako funkcje na $X$ wtedy i tylko wtedy, gdy
$f(a)=g(a)$ dla każdego $a\in X$. To zachodzi dokładnie wtedy,
gdy $(f-g)(a)=0$ dla wszystkich $a\in X$, a więc wtedy i tylko wtedy,
gdy $f-g\in\I_k(X)$. Zatem \emph{elementy $k[X]$ odpowiadają dokładnie
funkcjom wielomianowym na zbiorze punktów $X$}.

\begin{uwaga}[Uwaga o wyjściu od dowolnego ideału]
Jeżeli zaczynamy nie od zbioru $X$, lecz od ideału $J$ i konstruujemy
$k[x_1,\dots,x_n]/J$, to iloraz może zawierać niezerowe nilpotenty.
Natomiast pierścień współrzędnych zbioru punktów $X$ jest zawsze
\emph{zredukowany}, ponieważ $\I_k(X)$ jest radykalny.
Nad ciałem algebraicznie domkniętym okaże się, że
$k[\V_k(J)]\cong k[x_1,\dots,x_n]/\sqrt J$.
Bez tego założenia równość może zawieść.
\end{uwaga}

\begin{przyklad}[Parabola]
Niech $X=\V_k(y-x^2)\subseteq k^2$.
Dla dowolnego ciała $k$ mamy izomorfizm
$k[x,y]/(y-x^2)\cong k[t]$, wynikający z podstawienia
$x=t$, $y=t^2$. Jeśli $k$ jest nieskończone, to
$\I_k(X)=(y-x^2)$: gdy wielomian $f$ znika na wszystkich
punktach $(t,t^2)$, to $f(t,t^2)$ jest wielomianem jednej
zmiennej z nieskończenie wieloma pierwiastkami, czyli jest zerowy;
z poprzedniego rozdziału wiemy, że wówczas $f\in(y-x^2)$.
W takim przypadku $k[X]\cong k[t]$.

Nad $k=\F_q$ mamy natomiast $x^q-x\in\I_k(X)$,
choć $x^q-x\notin(y-x^2)$: gdyby należał, jego obraz w
$k[x,y]/(y-x^2)\cong k[x]$ byłby zerowy, a nie jest.
Nie wolno zatem nad skończonym ciałem utożsamiać obu ilorazów.
\end{przyklad}

\begin{przyklad}[Dwie osie]
Niech $k$ będzie ciałem nieskończonym i
$X=\V_k(xy)\subseteq k^2$. Wtedy $\I_k(X)=(xy)$.
Rzeczywiście, jeśli $f$ znika na osi $y=0$, to
$f(x,0)=0$ jako wielomian w $x$ (bo ma nieskończenie wiele zer),
a zatem $y\mid f$. Analogicznie znikanie na osi $x=0$
daje $x\mid f$. Zapiszmy $f=yg$; ponieważ $f(0,y)=yg(0,y)=0$
jako wielomian, z faktu, że $k[y]$ jest dziedziną, wynika
$g(0,y)=0$, czyli $x\mid g$. Mamy więc $xy\mid f$.
Odwrotna inkluzja $(xy)\subseteq\I_k(X)$ jest oczywista.
Stąd $k[X]\cong k[x,y]/(xy)$. Klasy $x$ oraz $y$
są niezerowe, lecz ich iloczyn jest zerowy: algebraicznie
widzimy dwie składowe przechodzące przez wspólny punkt.
\end{przyklad}

\begin{przyklad}[Krzywa hiperboliczna i odwracalna współrzędna]
Niech $k$ będzie nieskończonym ciałem i
$H=\V_k(xy-1)\subseteq k^2$. Każdy punkt $H$ ma postać
$(t,t^{-1})$ dla pewnego $t\in k^\times:=k\setminus\{0\}$.
W ilorazie $k[x,y]/(xy-1)$ klasy $x$ i $y$ są wzajemnie
odwrotne. Rozważmy \emph{pierścień wielomianów Laurenta}
$k[t,t^{-1}]=\{\sum_{j=-m}^r a_jt^j:a_j\in k\}$.
Homomorfizm $\Phi:k[x,y]\to k[t,t^{-1}]$, zadany przez
$\Phi(x)=t$, $\Phi(y)=t^{-1}$, jest surjektywny:
każda potęga $t^j$ o $j\ge0$ jest obrazem $x^j$, a każda
potęga $t^{-j}$ dla $j>0$ jest obrazem $y^j$.
Ponadto $xy-1\in\kerf\Phi$, więc otrzymujemy homomorfizm
$\overline\Phi:k[x,y]/(xy-1)\to k[t,t^{-1}]$.
Aby wykazać, że jest on izomorfizmem, skonstruujmy odwrotność:
wielomianowi Laurenta $\sum_j a_jt^j$ przyporządkujmy
$\sum_{j\ge0}a_j[x]^j+\sum_{j<0}a_j[y]^{-j}$,
gdzie $[x]$ i $[y]$ oznaczają klasy w ilorazie.
Jest to homomorfizm, ponieważ $[x][y]=1$, zatem
odpowiedniki dodatnich i ujemnych potęg mnożą się zgodnie
z regułą $t^it^j=t^{i+j}$. Złożenia w obu kierunkach
ustalają generatory $t,t^{-1}$ oraz $[x],[y]$, więc są identycznościami.

Obliczmy jeszcze $\I_k(H)$. Jeśli wielomian $f(x,y)$ znika na
wszystkich punktach $(t,t^{-1})$, to wielomian Laurenta
$\Phi(f)$ znika dla wszystkich $t\in k^\times$.
Po pomnożeniu przez dostatecznie wysoką potęgę $t$
otrzymujemy zwykły wielomian jednej zmiennej, który znika
na nieskończonym zbiorze $k^\times$; jest więc wielomianem zerowym.
Stąd $f\in\kerf\Phi=(xy-1)$. Inkluzja odwrotna jest
oczywista, ponieważ $xy-1$ zeruje się na $H$. W rezultacie
$k[H]\cong k[t,t^{-1}]$.
\end{przyklad}

\section{Dlaczego nie każdą nieskończoną rodzinę równań trzeba zapisywać?}

W definicji $\V_k(S)$ dopuściliśmy dowolny, także nieskończony
zbiór wielomianów. Można zapytać, czy taki opis jest rzeczywiście
potrzebny. Odpowiedź daje \emph{twierdzenie Hilberta o bazie}.
Jest to zarazem ważny moment: po raz pierwszy użyjemy własności
pierścienia wielomianów do stwierdzenia czegoś o \emph{wszystkich}
układach równań jednocześnie.

\begin{definicja}[Pierścień noetherowski]
Pierścień przemienny $A$ nazywamy \emph{noetherowskim}, jeśli każdy
jego ideał jest skończenie generowany.
To znaczy: dla każdego ideału $I\subseteq A$ istnieją
$a_1,\dots,a_r\in I$ takie, że $I=(a_1,\dots,a_r)$.
\end{definicja}

\begin{lemat}[Warunek łańcuchów wstępujących]
\label{lem:acc}
Dla pierścienia przemiennego $A$ następujące warunki są równoważne:
\begin{enumerate}[label=\arabic*)]
\item każdy ideał $A$ jest skończenie generowany;
\item każdy wstępujący ciąg ideałów
$I_1\subseteq I_2\subseteq I_3\subseteq\cdots$
jest od pewnego miejsca stały.
\end{enumerate}
\end{lemat}
\begin{proof}
Załóżmy (1) i rozważmy ciąg jak w (2).
Suma mnogościowa $I:=\bigcup_{n\ge1}I_n$ jest ideałem:
jeśli $a\in I_r$ i $b\in I_s$, to dla $m=\max(r,s)$
mamy $a,b\in I_m$, więc $a-b\in I_m\subseteq I$;
analogicznie $ca\in I$ dla $c\in A$.
Jeżeli $I=(0)$, to wszystkie $I_n=(0)$ i ciąg jest stały.
W przeciwnym przypadku z (1) wybieramy niepusty skończony
układ generatorów $a_1,\dots,a_t$ ideału $I$.
Każdy $a_j$ należy do pewnego
$I_{n_j}$. Dla $N=\max_j n_j$ wszystkie generatory należą
do $I_N$, więc $I\subseteq I_N$. Ponieważ $I_N\subseteq I$,
mamy $I=I_N$, a dla $n\ge N$ zachodzi $I_n=I_N$.

Załóżmy teraz (2) i niech $I$ będzie ideałem.
Gdyby $I$ nie był skończenie generowany, moglibyśmy wybrać
$a_1\in I$ i następnie kolejno $a_{r+1}\in I\setminus(a_1,\dots,a_r)$: taki element istnieje, ponieważ żaden
skończony zbiór nie generuje $I$.
Powstałby ściśle rosnący ciąg
$(a_1)\subsetneq(a_1,a_2)\subsetneq\cdots$,
co przeczy (2). Zatem (1) zachodzi.
\end{proof}

\begin{twierdzenie}[Twierdzenie Hilberta o bazie]
\label{tw:hilbert-basis}
Jeśli $A$ jest pierścieniem noetherowskim, to pierścień wielomianów
$A[t]$ również jest noetherowski.
\end{twierdzenie}
\begin{proof}
Niech $I\subseteq A[t]$ będzie dowolnym ideałem.
Dla każdej liczby całkowitej $d\ge0$ zdefiniujmy zbiór
$L_d\subseteq A$ złożony z tych $a\in A$, dla których istnieje
wielomian $f\in I$ o stopniu nie większym niż $d$,
którego współczynnik przy $t^d$ jest równy $a$.
Dopuszczamy więc zerowy współczynnik wiodący (a także wielomian
zerowy). Intuicyjnie $L_d$ zbiera współczynniki, które możemy
\emph{otrzymać na pozycji $t^d$} bez używania wyższych potęg.

Każdy $L_d$ jest ideałem: jeśli $a,b\in L_d$ mają świadków
$f,g\in I$ stopnia co najwyżej $d$, to $f-g\in I$ ma stopień
co najwyżej $d$ i współczynnik $a-b$ przy $t^d$.
Dla $c\in A$ wielomian $cf\in I$ ma taki współczynnik równy $ca$.
Co więcej $L_d\subseteq L_{d+1}$: jeżeli $f$ świadczy, że
$a\in L_d$, to $tf$ ma stopień co najwyżej $d+1$ i współczynnik
$a$ przy $t^{d+1}$. Ponieważ $A$ jest noetherowski,
lemat~\ref{lem:acc} daje liczbę $N\ge0$ taką, że
$L_N=L_{N+1}=L_{N+2}=\cdots$.

Dla każdego $d=0,\dots,N$ wybierzmy skończony układ generatorów
$a_{d,1},\dots,a_{d,r_d}$ ideału $L_d$. Z definicji $L_d$
możemy wybrać wielomiany $f_{d,j}\in I$ stopnia co najwyżej $d$
i ze współczynnikiem $a_{d,j}$ przy $t^d$.
Niech $J$ będzie ideałem $A[t]$ generowanym przez \emph{wszystkie}
$f_{d,j}$, gdzie $0\le d\le N$ i $1\le j\le r_d$.
To skończony zbiór generatorów; oczywiście $J\subseteq I$.

Wykażemy przez indukcję względem stopnia, że każdy $f\in I$
należy do $J$. Dla $f=0$ nie ma czego dowodzić.
Niech $m=\deg f$ i $a$ będzie współczynnikiem przy $t^m$.
Wprowadźmy $d=\min(m,N)$.
Skoro $a\in L_m$, a $L_m=L_N$ gdy $m\ge N$,
to $a\in L_d$. Możemy zatem zapisać
$a=\sum_{j=1}^{r_d}c_ja_{d,j}$ dla pewnych $c_j\in A$.
Wielomian
$h:=\sum_{j=1}^{r_d}c_jt^{m-d}f_{d,j}$ należy do $J$,
ma stopień co najwyżej $m$ i współczynnik $a$ przy $t^m$.
Zatem $f-h\in I$ ma stopień mniejszy od $m$ albo jest zerowy.
Na mocy założenia indukcyjnego $f-h\in J$, a więc $f\in J$.
W rezultacie $I=J$ jest skończenie generowany.
\end{proof}

\begin{wniosek}[Wielomiany w wielu zmiennych]
\label{wn:noetherian-polys}
Jeżeli $k$ jest ciałem i $n\ge1$, to $k[x_1,\dots,x_n]$
jest pierścieniem noetherowskim. W szczególności każdy zbiór
algebraiczny $X\subseteq k^n$ można zapisać w postaci
$X=\V_k(f_1,\dots,f_r)$ dla pewnej skończonej liczby wielomianów.
\end{wniosek}
\begin{proof}
Każdy ideał ciała $k$ jest równy $(0)$ lub $k$:
niezerowy element ideału jest odwracalny, więc taki ideał jest całym
ciałem. Zatem $k$ jest noetherowski.
Stosując twierdzenie~\ref{tw:hilbert-basis} kolejno do
$k[x_1]$, $k[x_1,x_2]=k[x_1][x_2]$ itd., otrzymujemy
noetherowskość $k[x_1,\dots,x_n]$.
Jeśli $X=\V_k(S)$, to ideał $(S)$ ma skończony zbiór
generatorów $f_1,\dots,f_r$. Każdy $f_i$ jest skończoną
sumą wielokrotności wielomianów pochodzących z $S$.
Zbierając wszystkie wielomiany z $S$ występujące w tych
sumach, uzyskujemy skończony podzbiór $S_0\subseteq S$
z własnością $(S_0)=(S)$. Istotnie, $f_i\in(S_0)$ daje
$(S)=(f_1,\dots,f_r)\subseteq(S_0)$, a z $S_0\subseteq S$
wynika inkluzja odwrotna. Zatem
$X=\V_k((S))=\V_k(S_0)$: wystarcza skończenie wiele
\emph{pierwotnych} równań z rodziny $S$.
\end{proof}

\begin{uwaga}[Twierdzenie o istnieniu, nie o algorytmie]
Twierdzenie Hilberta o bazie gwarantuje, że skończony opis istnieje,
ale nie podaje efektywnej metody znalezienia generatorów dla
arbitralnie zadanej rodziny równań. W obliczeniowej geometrii
algebraicznej rolę efektywnego narzędzia pełnią przede wszystkim
\emph{bazy Gröbnera}, które omówimy później tylko wtedy,
gdy okażą się potrzebne do konkretnych rachunków.
\end{uwaga}

\section{Ciała algebraicznie domknięte i problem zer}

\subsection{Dlaczego potrzebujemy domkniętości algebraicznej?}

\begin{definicja}[Ciało algebraicznie domknięte]
Ciało $k$ nazywamy \emph{algebraicznie domkniętym}, jeśli każdy
niezerowy wielomian $f\in k[t]$ dodatniego stopnia ma pierwiastek w $k$.
\end{definicja}

\begin{stwierdzenie}[Rozkład wielomianu na czynniki liniowe]
Niech $k$ będzie ciałem. Następujące warunki są równoważne:
\begin{enumerate}[label=\arabic*)]
\item $k$ jest algebraicznie domknięte;
\item każdy niezerowy wielomian $f\in k[t]$ dodatniego stopnia
rozpada się w $k[t]$ na iloczyn czynników liniowych, z uwzględnieniem
krotności i niezerowego współczynnika stałego;
\item każdy wielomian nierozkładalny w $k[t]$ ma stopień $1$.
\end{enumerate}
\end{stwierdzenie}
\begin{proof}
Załóżmy (1). Dla niezerowego $f$ dodatniego stopnia wybierzmy
pierwiastek $a\in k$. Na mocy twierdzenia o czynniku
$f=(t-a)g$, gdzie $\deg g=\deg f-1$. Powtarzamy tę operację,
a indukcja kończy się na wielomianie stopnia $0$. Otrzymujemy (2).
Jeśli (2), to nierozkładalny wielomian dodatniego stopnia nie może
mieć co najmniej dwóch czynników liniowych, zatem ma stopień $1$.
Tak otrzymujemy (3). Wreszcie każdy wielomian dodatniego stopnia
posiada nierozkładalny dzielnik dodatniego stopnia:
można go znaleźć, rozkładając wielomian, jeśli sam jest rozkładalny,
a następnie postępując indukcyjnie względem stopnia.
Z (3) wynika, że dzielnik ten ma postać $c(t-a)$,
więc $a$ jest pierwiastkiem wyjściowego wielomianu. To daje (1).
\end{proof}

\begin{przyklad}[Przykłady ciał]
Ciało $\C$ jest algebraicznie domknięte (podstawowe twierdzenie algebry).
Ciało $\R$ nie jest: wielomian $t^2+1$ nie ma rzeczywistych pierwiastków.
Ciało $\Q$ też nie jest: wielomian $t^2-2$ nie ma wymiernych
pierwiastków. Ciało skończone $\F_q$ o $q$ elementach nie jest algebraicznie
domknięte, ponieważ wielomian $t^q-t+1$ przyjmuje wartość $1$
na wszystkich jego elementach. Istotnie, dla $a\ne0$
mnożenie przez $a$ permutuje $q-1$ niezerowych elementów ciała.
Mnożąc elementy obu stron tej permutacji, mamy
$\prod_{b\in\F_q^\times}b
=a^{q-1}\prod_{b\in\F_q^\times}b$.
Iloczyn po prawej jest niezerowy, więc $a^{q-1}=1$,
a zatem $a^q=a$. Dla $a=0$ ta równość jest oczywista.
\end{przyklad}

\begin{uwaga}[Algebraiczne domknięcie a nieskończoność]
Każde algebraicznie domknięte ciało jest nieskończone.
Rzeczywiście, gdyby $k=\{a_1,\dots,a_q\}$ było skończone,
wielomian $f(t)=\prod_{i=1}^q(t-a_i)+1$ miałby dodatni stopień,
a dla każdego $a_i$ przyjmowałby wartość $1$.
Nie miałby więc żadnego pierwiastka w $k$.
\end{uwaga}

\subsection{Dwie wersje twierdzenia Hilberta o zerach}

Poniżej udowodnimy dwie fundamentalne wersje
\emph{Nullstellensatz} (twierdzenia Hilberta o zerach).
Pierwsza zapewnia istnienie rozwiązania układu równań;
druga dokładnie identyfikuje wszystkie wielomiany, które
zerują się na zbiorze tych rozwiązań.

\begin{twierdzenie}[Słabe twierdzenie Hilberta o zerach]
\label{tw:weak-NS}
Niech $k$ będzie ciałem algebraicznie domkniętym,
$n\ge1$, a $I\subsetneq k[x_1,\dots,x_n]$ ideałem właściwym.
Wtedy istnieje $a\in k^n$ taki, że $f(a)=0$ dla każdego $f\in I$,
czyli $\V_k(I)\ne\varnothing$.
\end{twierdzenie}

Równoważna kontrapozycja głosi: jeśli układ wielomianowych
równań nad ciałem algebraicznie domkniętym nie ma wspólnego
rozwiązania, to wielomiany generują ideał jednostkowy, czyli
istnieją $f_1,\dots,f_r\in I$ oraz
$h_1,\dots,h_r\in k[x_1,\dots,x_n]$ takie, że
$h_1f_1+\cdots+h_rf_r=1$.
Taką równość można rozumieć jako \emph{algebraiczne świadectwo}
niemożliwości rozwiązania układu.

\begin{twierdzenie}[Silne twierdzenie Hilberta o zerach]
\label{tw:strong-NS}
Niech $k$ będzie ciałem algebraicznie domkniętym,
$n\ge1$ i $I\subseteq k[x_1,\dots,x_n]$ ideałem.
Wtedy
\[
\I_k(\V_k(I))=\sqrt I
=\{f\in k[x_1,\dots,x_n]:f^m\in I\text{ dla pewnego }m\ge1\}.
\]
\end{twierdzenie}

Zauważmy, że tezy dotyczą \emph{każdego} ideału,
a nie tylko ideału głównego czy radykalnego. Najtrudniejszym
składnikiem dowodu słabej wersji będzie lemat Zariskiego.
Udowodnimy go od podstaw, aby nie ukrywać istotnego przejścia
od algebry do geometrii.

\section{Lemat Zariskiego: dlaczego ciało skończenie generowane jako algebra jest algebraiczne?}

\begin{definicja}[Rozszerzenie algebraiczne i element całkowity]
Niech $k\subseteq L$ będzie rozszerzeniem ciał. Element $u\in L$
jest \emph{algebraiczny nad $k$}, jeśli spełnia pewne niezerowe
równanie wielomianowe o współczynnikach w $k$.
Jeżeli $R\subseteq B$ jest zawieraniem pierścieni, to $u\in B$
nazywamy \emph{całkowitym nad $R$}, jeśli istnieją
$a_0,\dots,a_{m-1}\in R$ takie, że
$u^m+a_{m-1}u^{m-1}+\cdots+a_0=0$ dla pewnego $m\ge1$.
W tym drugim przypadku wymagamy, by wielomian był \emph{moniczny},
czyli miał współczynnik wiodący równy $1$.
\end{definicja}

Nad ciałem każda zależność wielomianowa z niezerowym
współczynnikiem wiodącym może zostać podzielona przez ten współczynnik.
Dlatego dla rozszerzenia ciał \emph{algebraiczny} i
\emph{całkowity} oznaczają to samo. Nad ogólnym pierścieniem jest
inaczej: np. $1/2$ jest algebraiczny nad $\Q$, lecz nie jest
całkowity nad $\Z$.

\begin{definicja}[Algebra skończenie generowana]
Niech $k$ będzie ciałem. Pierścień $B$ zawierający $k$
nazywamy \emph{skończenie generowaną $k$-algebrą}, jeśli
istnieją $u_1,\dots,u_s\in B$ takie, że
$B=k[u_1,\dots,u_s]$, czyli każdy element $B$ jest
wartością pewnego wielomianu nad $k$ w punktach $u_1,\dots,u_s$.
Nie należy tego mylić z warunkiem, że $B$ ma skończony wymiar
jako przestrzeń wektorowa nad $k$.
\end{definicja}

\begin{przyklad}[Algebra nie musi mieć skończonego wymiaru]
Pierścień $k[t]$ jest generowany przez jeden element $t$
jako $k$-algebra, ale jego wymiar jako przestrzeni wektorowej
nad $k$ jest nieskończony: wielomiany $1,t,t^2,\dots$
są liniowo niezależne. Lemat Zariskiego mówi, że sytuacja
radykalnie zmienia się, gdy skończenie generowana algebra
jest \emph{ciałem}.
\end{przyklad}

\subsection{Ciało ułamków i algebraiczna niezależność}

Zanim przejdziemy do dowodu lematu Zariskiego, potrzebujemy jeszcze
dwóch pojęć. W pierwszym przechodzimy od dziedziny do ciała,
w którym można dzielić przez każdy element niezerowy.
W drugim sprawdzamy, czy pewne wielkości spełniają jakąkolwiek
relację wielomianową.

\begin{definicja}[Ciało ułamków dziedziny]
Niech $R$ będzie dziedziną całkowitości.
Rozważmy pary $(a,b)$, gdzie $a,b\in R$ i $b\ne0$,
oraz relację $(a,b)\sim(c,d)$ wtedy i tylko wtedy,
gdy $ad=bc$. Klasę pary $(a,b)$ oznaczamy $a/b$.
Zbiór klas nazywamy \emph{ciałem ułamków} $R$
i oznaczamy $\Frac(R)$. Działania definiujemy wzorami
\[
\frac ab+\frac cd=\frac{ad+bc}{bd},\qquad
\frac ab\cdot\frac cd=\frac{ac}{bd}.
\]
\end{definicja}

\begin{lemat}[Poprawność konstrukcji ciała ułamków]
\label{lem:frac-construction}
Dla każdej dziedziny $R$ powyższa relacja jest relacją
równoważności, działania są dobrze określone, a $\Frac(R)$
jest ciałem. Odwzorowanie $R\to\Frac(R)$,
$a\mapsto a/1$, jest różnowartościowym homomorfizmem pierścieni.
\end{lemat}
\begin{proof}
Zwrotność wynika z $ab=ba$, symetria z symetrii równości.
Dla przechodniości załóżmy $a/b=c/d$ i $c/d=e/f$.
Mamy $ad=bc$ i $cf=de$. Mnożąc pierwszą równość przez
$f$, otrzymujemy $adf=bcf=bde$. Ponieważ $d\ne0$
i $R$ jest dziedziną, możemy skrócić $d$; stąd $af=be$,
czyli $a/b=e/f$.

Sprawdźmy niezależność dodawania od reprezentantów.
Jeśli $a/b=a'/b'$ oraz $c/d=c'/d'$, to
$ab'=a'b$ i $cd'=c'd$. Z tych równości otrzymujemy
\begin{align*}
(ad+bc)b'd'
 &=ab'dd'+bb'cd'\\
 &=a'bdd'+bb'c'd\\
 &=(a'd'+b'c')bd.
\end{align*}
Dlatego sumy wyznaczają tę samą klasę. Dla mnożenia
wystarczy pomnożyć równości $ab'=a'b$ oraz $cd'=c'd'$:
$acb'd'=a'c'bd$, czyli również iloczyn nie zależy od
reprezentantów.

Łączność, przemienność i rozdzielność sprowadzają się
po podstawieniu wzorów do odpowiednich własności $R$.
Przykładowo $(a/b)(c/d+e/f)=a(cf+de)/(bdf)$,
podczas gdy $(a/b)(c/d)+(a/b)(e/f)=ac/(bd)+ae/(bf)$;
po sprowadzeniu do wspólnego mianownika $bdf$ otrzymujemy
tę samą klasę $a(cf+de)/(bdf)$.
Elementy neutralne to $0/1$ i $1/1$.
Przeciwnym do $a/b$ jest $(-a)/b$.
Jeśli $a/b\ne0/1$, to $a\ne0$ i $b/a$ jest dozwoloną klasą,
a $(a/b)(b/a)=1/1$. Zatem mamy ciało.
Wreszcie $a/1=0/1$ oznacza $a=0$, więc odwzorowanie
$R\to\Frac(R)$ jest różnowartościowe; zachowanie działań
i jedynki wynika z ich definicji.
\end{proof}

\begin{uwaga}[Ułamki jako funkcje wymierne]
Dla $R=k[t_1,\dots,t_r]$ otrzymujemy ciało
$k(t_1,\dots,t_r)=\Frac(k[t_1,\dots,t_r])$.
Jego elementy mają postać $p(t)/q(t)$, gdzie $p,q$ są
wielomianami, a $q$ jest \emph{niezerowym wielomianem},
niekoniecznie wielomianem, który nie znika w każdym punkcie $k^r$.
To ważne: ułamek jest formalnym elementem ciała ułamków,
nie zaś funkcją określoną we wszystkich punktach przestrzeni afinicznej.
\end{uwaga}

\begin{definicja}[Algebraiczna niezależność]
Niech $k\subseteq L$ będzie rozszerzeniem ciał i
$u_1,\dots,u_r\in L$. Elementy te są
\emph{algebraicznie niezależne nad $k$}, jeśli
$P(u_1,\dots,u_r)\ne0$ dla każdego niezerowego
$P\in k[t_1,\dots,t_r]$. W przeciwnym razie
nazywamy je \emph{algebraicznie zależnymi}.
Pustą rodzinę traktujemy jako algebraicznie niezależną.
\end{definicja}

\begin{przyklad}[Zależność i niezależność]
W ciele $k(t)$ elementy $t$ i $t^2$ są algebraicznie zależne,
ponieważ niezerowy wielomian $P(u,v)=v-u^2$ spełnia
$P(t,t^2)=0$. Natomiast sam element $t$ jest algebraicznie
niezależny nad $k$: gdyby niezerowy wielomian $Q\in k[u]$
spełniał $Q(t)=0$ jako element ciała $k(t)$,
byłby równy zeru jako formalny wielomian, co jest sprzecznością.
\end{przyklad}

\begin{lemat}[Dołączenie elementu do rodziny niezależnej]
\label{lem:independence-algebraic}
Niech $k\subseteq L$ będzie rozszerzeniem ciał,
a $u_1,\dots,u_r\in L$ będą algebraicznie niezależne nad $k$.
Dla $v\in L$ następujące warunki są równoważne:
\begin{enumerate}[label=\arabic*)]
\item rodzina $u_1,\dots,u_r,v$ jest algebraicznie zależna nad $k$;
\item element $v$ jest algebraiczny nad ciałem $k(u_1,\dots,u_r)$.
\end{enumerate}
\end{lemat}
\begin{proof}
Jeśli (1), istnieje niezerowy wielomian
$P(t_1,\dots,t_r,z)=\sum_{j=0}^d P_j(t_1,\dots,t_r)z^j$
z $P(u_1,\dots,u_r,v)=0$.
Przynajmniej jeden $P_j$ jest niezerowym wielomianem.
Algebraiczna niezależność $u_i$ oznacza, że
$P_j(u_1,\dots,u_r)\ne0$, więc otrzymujemy
niezerowy wielomian jednej zmiennej $z$ o współczynnikach
w ciele $k(u_1,\dots,u_r)$, który zeruje się dla $z=v$.
To znaczy, że zachodzi (2).

Odwrotnie, jeśli $v$ jest algebraiczny nad
$F=k(u_1,\dots,u_r)$, to
$\sum_{j=0}^d c_j v^j=0$ dla pewnych $c_j\in F$,
nie wszystkich zerowych. Wszystkie $c_j$ są ilorazami
wielomianów w $u_1,\dots,u_r$. Po pomnożeniu przez wspólny
niezerowy mianownik otrzymujemy niezerową zależność
$\sum_{j=0}^d Q_j(u_1,\dots,u_r)v^j=0$
z $Q_j\in k[t_1,\dots,t_r]$, nie wszystkimi zerowymi.
Zatem rodzina $u_1,\dots,u_r,v$ jest algebraicznie zależna.
\end{proof}

\begin{lemat}[Skończona liczba elementów całkowitych]
\label{lem:integral-finite}
Niech $R\subseteq B$ będzie zawieraniem pierścieni i niech
elementy $u_1,\dots,u_s\in B$ będą całkowite nad $R$.
Wtedy $R[u_1,\dots,u_s]$ jest skończenie generowanym
$R$-modułem: istnieją elementy $b_1,\dots,b_m$
tej algebry takie, że każdy jej element można zapisać jako
$\sum_{j=1}^m r_jb_j$ z $r_j\in R$.
\end{lemat}
\begin{proof}
Najpierw rozważmy jeden element $u$, który spełnia
$u^d+a_{d-1}u^{d-1}+\cdots+a_0=0$.
Przekształcamy tę równość do postaci
$u^d=-a_{d-1}u^{d-1}-\cdots-a_0$.
Każdą wyższą potęgę $u^N$ możemy kolejno zredukować,
mnożąc tę relację przez $u^{N-d}$. Indukcja względem $N$
daje $R[u]=R\cdot1+R\cdot u+\cdots+R\cdot u^{d-1}$.

Dla kilku elementów $u_j$ wybierzmy dla każdego moniczne
równanie stopnia $d_j$. Każdy iloczyn potęg tych elementów,
a zatem dowolny element $R[u_1,\dots,u_s]$,
można zredukować do kombinacji monomów
$u_1^{e_1}\cdots u_s^{e_s}$ z $0\le e_j<d_j$.
Jest tylko $d_1\cdots d_s$ takich monomów.
\end{proof}

\begin{lemat}[Dlaczego elementy skończonej algebry spełniają równania moniczne]
\label{lem:determinant-integral}
Niech $R\subseteq B$ będą pierścieniami i załóżmy, że $B$
jest skończenie generowanym $R$-modułem. Wtedy każdy element
$z\in B$ jest całkowity nad $R$.
\end{lemat}
\begin{proof}
Wybierzmy generatory $b_1,\dots,b_m$ modułu $B$ nad $R$.
Mnożenie przez $z$ przenosi każdy generator w element $B$,
możemy więc znaleźć współczynniki $a_{ij}\in R$ takie, że
$zb_i=\sum_{j=1}^m a_{ij}b_j$.
Niech $M=(a_{ij})\in M_m(R)$ i niech
$b=(b_1,\dots,b_m)^T$. Otrzymujemy równanie
$(zE-M)b=0$, gdzie $E$ jest macierzą jednostkową.

Dla macierzy kwadratowej $C$ przez $\operatorname{adj}(C)$
oznaczamy transponowaną macierz dopełnień algebraicznych.
Rozwinięcie wyznacznika względem wiersza lub kolumny daje
$\operatorname{adj}(C)C=\det(C)E$: na przekątnej otrzymujemy
wyznacznik $C$, a poza przekątną wyznacznik macierzy
z dwoma jednakowymi wierszami, czyli zero.
Tożsamość obowiązuje też nad dowolnym pierścieniem przemiennym.
Stosujemy ją do $C=zE-M$ i dostajemy
$\operatorname{adj}(zE-M)(zE-M)=\det(zE-M)E$.
Po pomnożeniu równania przez tę macierz dostajemy
$\det(zE-M)b_i=0$ dla każdego $i$.
Ponieważ generatory $b_i$ generują cały $B$ jako $R$-moduł,
a w szczególności generują element $1\in B$,
wynika stąd $\det(zE-M)\cdot1=0$ w $B$.
Wielomian $P(T):=\det(TE-M)\in R[T]$ jest moniczny
stopnia $m$. Otrzymaliśmy $P(z)=0$, czyli $z$ jest całkowity nad $R$.
\end{proof}

\begin{lemat}[Dziedzina całkowitości nad pierścieniem]
\label{lem:field-integral}
Niech $R\subseteq L$ będą pierścieniami, gdzie $L$ jest ciałem.
Jeżeli każdy element $L$ jest całkowity nad $R$, to $R$ jest ciałem.
\end{lemat}
\begin{proof}
Ustalmy $0\ne r\in R$. W ciele $L$ istnieje $r^{-1}$.
Na mocy założenia spełnia ono równanie
$(r^{-1})^m+a_{m-1}(r^{-1})^{m-1}+\cdots+a_0=0$
z $a_j\in R$. Mnożąc przez $r^{m-1}$, otrzymujemy
$ r^{-1}=-(a_{m-1}+a_{m-2}r+\cdots+a_0r^{m-1})\in R$.
Zatem każdy niezerowy element $R$ ma odwrotność w $R$.
\end{proof}

\begin{lemat}[Lokalizacja pierścienia wielomianów nie staje się ciałem]
\label{lem:localization-not-field}
Niech $k$ będzie ciałem, $r\ge1$ i
$0\ne d\in k[t_1,\dots,t_r]$.
Pierścień $k[t_1,\dots,t_r,1/d]$, złożony z wyrażeń
$p/d^m$ z $p\in k[t_1,\dots,t_r]$ i $m\ge0$,
nie jest ciałem.
\end{lemat}
\begin{proof}
Wprowadźmy $P=k[t_1,\dots,t_r]$ oraz $s=1+t_1d\in P$.
Wielomian $s$ jest niezerowy i niestały: stopień iloczynu
$t_1d$ względem $t_1$ jest o jeden większy od stopnia $d$,
więc dodanie stałej $1$ nie może usunąć wyrazu najwyższego stopnia.
Wielomian niestały nie jest jednostką pierścienia $P$,
ponieważ stopnie niezerowych iloczynów w zmiennej $t_1$
sumują się, a dla zmiennych pozostałych współczynniki tworzą
dziedzinę $k[t_2,\dots,t_r]$.

Załóżmy dla sprzeczności, że $s$ jest odwracalny w
$P[1/d]$. Wtedy istnieją $p\in P$ i $m\ge0$ takie, że
$s(p/d^m)=1$, czyli $sp=d^m$ w dziedzinie $P$.
W ilorazie $P/(s)$ zachodzi $t_1d=-1$.
Oznacza to, że klasa $d$ ma odwrotność $-t_1$.
Jednocześnie z $sp=d^m$ wynika, że $d^m=0$ w tym ilorazie.
Jeśli $m\ge1$, odwracalność $d$ daje $1=0$ w $P/(s)$.
Jeśli $m=0$, z $sp=1$ otrzymujemy to samo.
W obu przypadkach $1\in(s)$, więc $s$ jest jednostką $P$,
sprzecznie z poprzednim akapitem.
\end{proof}

\begin{twierdzenie}[Lemat Zariskiego]
\label{tw:zariski-lemma}
Niech $k$ będzie ciałem, a $L$ ciałem zawierającym $k$.
Jeśli istnieją $a_1,\dots,a_n\in L$ takie, że
$L=k[a_1,\dots,a_n]$ jako $k$-algebra, to każdy element $L$
jest algebraiczny nad $k$. W rzeczywistości $L$ jest
skończenie wymiarową przestrzenią wektorową nad $k$.
\end{twierdzenie}
\begin{proof}
Spośród generatorów $a_1,\dots,a_n$ wybierzmy maksymalną
(ze względu na liczbę elementów) podrodzinę algebraicznie
niezależną nad $k$. Po zmianie numeracji możemy przyjąć,
że jest nią $a_1,\dots,a_r$; być może $r=0$.
Algebraiczna niezależność oznacza, że dla każdego
niezerowego $P\in k[t_1,\dots,t_r]$ mamy
$P(a_1,\dots,a_r)\ne0$.
Dlatego $k[a_1,\dots,a_r]$ jest izomorficzny z
pierścieniem wielomianów $P=k[t_1,\dots,t_r]$,
a wewnątrz ciała $L$ możemy rozważyć jego ciało ułamków
$F=k(a_1,\dots,a_r)$.

Z maksymalności wybranej rodziny i
lematu~\ref{lem:independence-algebraic} wynika, że każdy
pozostały generator $a_j$, $j>r$, jest algebraiczny nad $F$.
Gdyby któryś nie był, moglibyśmy go dołączyć do rodziny,
zachowując algebraiczną niezależność, co przeczyłoby jej
maksymalności. Wobec tego dla każdego $j>r$ istnieje
równanie moniczne
$ a_j^{m_j}+c_{j,m_j-1}a_j^{m_j-1}+\cdots+c_{j,0}=0$
o współczynnikach $c_{j,l}\in F$.
Każdy taki współczynnik jest ilorazem dwóch elementów $P$.
Wybierzmy jeden niezerowy $d\in P$ podzielny przez wszystkie
mianowniki występujące w skończonej liczbie tych równań.
Wtedy $c_{j,l}\in R:=P[1/d]$, więc wszystkie $a_j$ dla $j>r$
są całkowite nad $R$.

Zauważmy, że $R\subseteq L$: elementy $P$ są w $L$,
a $1/d$ również należy do $L$, bo $L$ jest ciałem i $d\ne0$.
Co więcej,
$L=k[a_1,\dots,a_n]=R[a_{r+1},\dots,a_n]$:
lewa strona zawiera prawą, zaś prawa zawiera wszystkie generatory
lewej. Na mocy lematu~\ref{lem:integral-finite}
$L$ jest skończenie generowanym $R$-modułem.
Z lematu~\ref{lem:determinant-integral} każdy element $L$
jest całkowity nad $R$, a z lematu~\ref{lem:field-integral}
wynika, że $R$ musi być ciałem.

Gdyby $r\ge1$, przeczyłoby to lematowi
\ref{lem:localization-not-field}, ponieważ
$R=k[t_1,\dots,t_r,1/d]$. Zatem $r=0$.
Oznacza to, że wszystkie generatory $a_j$ są algebraiczne nad
samym $k$. Możemy przyjąć $R=k$ i zastosować ponownie
lemat~\ref{lem:integral-finite}: otrzymujemy, że
$L=k[a_1,\dots,a_n]$ jest skończenie wymiarowe nad $k$.
Każdy element takiego rozszerzenia jest algebraiczny
(np. na mocy lematu~\ref{lem:determinant-integral}), co kończy dowód.
\end{proof}

\begin{uwaga}[Gdzie tkwi idea dowodu?]
Gdyby $L$ zawierało element transcendentalny $t$, to konieczne
byłyby odwrotności wielu wielomianów w $t$ (na przykład
$1/(1+t)$), ponieważ $L$ jest ciałem.
Tymczasem skończona liczba generatorów algebry pozwala
odwrócić tylko mianowniki pochodzące z pewnego skończonego
zbioru wielomianów. Element $1+t d$ z dowodu pokazuje,
że zawsze znajdziemy nowy nieodwracalny mianownik.
To jest geometryczno-algebraiczny mechanizm stojący za lematem Zariskiego.
\end{uwaga}

\section{Dowód twierdzenia Hilberta o zerach}

\subsection{Od maksymalnego ideału do punktu}

\begin{lemat}[Istnienie ideału maksymalnego w pierścieniu noetherowskim]
\label{lem:maximal-exists}
Niech $R$ będzie niezerowym pierścieniem noetherowskim.
Wtedy istnieje ideał maksymalny $\mathfrak m\subsetneq R$.
Co więcej, każdy ideał właściwy $I\subsetneq R$ jest zawarty
w pewnym ideale maksymalnym.
\end{lemat}
\begin{proof}
Rozważmy rodzinę wszystkich ideałów właściwych zawierających $I$.
Nie jest ona pusta, ponieważ zawiera sam $I$.
Gdyby żaden element tej rodziny nie był maksymalny ze względu na
inkluzję, moglibyśmy rozpocząć od $I_0=I$ i wybierać kolejno
ideały właściwe $I_0\subsetneq I_1\subsetneq I_2\subsetneq\cdots$.
Jest to sprzeczne z warunkiem łańcuchów wstępujących
z lematu~\ref{lem:acc}. Zatem istnieje ideał właściwy
$\mathfrak m$ zawierający $I$ i nieposiadający większego
ideału właściwego. Z definicji jest maksymalny.
Dla pierwszej części wystarczy przyjąć $I=(0)$;
jest to ideał właściwy, ponieważ $R\ne(0)$.
\end{proof}

\begin{lemat}[Ciało skończenie generowane nad ciałem algebraicznie domkniętym]
\label{lem:extension-trivial}
Niech $k$ będzie algebraicznie domkniętym ciałem,
a $L$ ciałem zawierającym $k$.
Jeśli $L$ jest skończenie generowaną $k$-algebrą,
to $L=k$.
\end{lemat}
\begin{proof}
Na mocy lematu Zariskiego każdy $a\in L$ jest algebraiczny
nad $k$. Wybierzmy niezerowy wielomian
$f\in k[t]$ o najmniejszym stopniu spośród wielomianów
spełniających $f(a)=0$. Ponieważ $k$ jest algebraicznie
domknięte, $f$ rozpada się na czynniki liniowe:
$f(t)=c\prod_{j=1}^r(t-b_j)$, gdzie $b_j\in k$ i $c\ne0$.
Z równości
$0=f(a)=c\prod_{j=1}^r(a-b_j)$ w ciele $L$
wynika, że $a=b_j$ dla pewnego $j$, a więc $a\in k$.
Dowolny element $L$ należy do $k$, stąd $L=k$.
\end{proof}

\begin{proof}[Dowód słabego twierdzenia Hilberta o zerach~\ref{tw:weak-NS}]
Oznaczmy $R=k[x_1,\dots,x_n]$ i niech $I\subsetneq R$.
Na mocy wniosku~\ref{wn:noetherian-polys} pierścień $R$
jest noetherowski; pierścień ilorazowy $R/I$ też jest
noetherowski, ponieważ jego ideały odpowiadają ideałom $R$
zawierającym $I$, a ilorazy generatorów tych ideałów generują
odpowiednie ideały w $R/I$.
Ponieważ $I$ jest właściwy, $R/I$ jest niezerowy.
Z lematu~\ref{lem:maximal-exists} istnieje ideał maksymalny
$\overline{\mathfrak m}$ w $R/I$.
Jego przeciwobraz względem rzutowania $R\to R/I$
jest ideałem maksymalnym $\mathfrak m$ w $R$ zawierającym $I$:
rzeczywiście $R/\mathfrak m\cong(R/I)/\overline{\mathfrak m}$,
a ostatni iloraz jest ciałem.

Rozważmy $L:=R/\mathfrak m$.
Jest to ciało oraz skończenie generowana $k$-algebra:
$L=k[\overline x_1,\dots,\overline x_n]$,
gdzie $\overline x_i=x_i+\mathfrak m$.
Włączenie $k\to L$ jest różnowartościowe: jego jądro jest
ideałem ciała $k$, a ponieważ $1\mapsto1\ne0$,
jądro jest zerowe. Możemy więc traktować $k$ jako podciało $L$.
Na mocy lematu~\ref{lem:extension-trivial} mamy $L=k$.
Każda klasa $\overline x_i$ odpowiada zatem pewnemu
$a_i\in k$. Niech $a=(a_1,\dots,a_n)\in k^n$.

Dla każdego $f\in R$ jego klasa w $L$ jest równa
$f(\overline x_1,\dots,\overline x_n)$, a po identyfikacji
$L=k$ jest to $f(a_1,\dots,a_n)$.
Jeżeli $f\in I\subseteq\mathfrak m$, jego klasa w $L$
jest zerowa, więc $f(a)=0$.
Punkt $a$ jest wspólnym zerem wszystkich wielomianów z $I$.
\end{proof}

\begin{wniosek}[Maksymalne ideały są ideałami punktów]
\label{wn:maximal-point}
Jeżeli $k$ jest algebraicznie domkniętym ciałem i $n\ge1$,
to każdy ideał maksymalny $\mathfrak m\subseteq k[x_1,\dots,x_n]$
ma dokładnie jedną postać
$\mathfrak m=(x_1-a_1,\dots,x_n-a_n)$ z $a\in k^n$.
\end{wniosek}
\begin{proof}
Zastosujmy dowód słabego twierdzenia do samego ideału
maksymalnego $\mathfrak m$.
Otrzymamy punkt $a$ taki, że wszystkie wielomiany z
$\mathfrak m$ znikają w $a$, czyli
$\mathfrak m\subseteq\I_k(\{a\})=(x_1-a_1,\dots,x_n-a_n)$.
Ideał po prawej jest właściwy, bo wielomian $1$ nie znika w $a$.
Maksymalność $\mathfrak m$ wymusza równość.
Dla jednoznaczności: jeśli ideały punktów $a$ i $b$ są równe,
to wielomian $x_i-a_i$ należy do ideału punktu $b$,
więc $0=(x_i-a_i)(b)=b_i-a_i$ dla każdego $i$.
Stąd $a=b$.
\end{proof}

\begin{uwaga}[Gdzie użyliśmy domkniętości algebraicznej?]
Jedyny moment w dowodzie, gdy potrzebowaliśmy algebraicznej
domkniętości $k$, to przejście od skończenie generowanego
\emph{ciała} $L=R/\mathfrak m$ do równości $L=k$.
Nad $\R$ możemy mieć $L\cong\C$, jak w przykładzie
$\R[t]/(t^2+1)$. Wówczas ideał maksymalny nie musi być
ideałem rzeczywistego punktu.
\end{uwaga}

\subsection{Trik Rabinowitscha i silne twierdzenie Hilberta o zerach}

\begin{lemat}[Trik Rabinowitscha]
\label{lem:rabinowitsch}
Niech $k$ będzie ciałem algebraicznie domkniętym, $n\ge1$,
$I\subseteq k[x_1,\dots,x_n]$ ideałem i
$0\ne f\in k[x_1,\dots,x_n]$ wielomianem, który znika
we wszystkich punktach $\V_k(I)$.
Wtedy istnieje $N\ge1$ takie, że $f^N\in I$.
\end{lemat}
\begin{proof}
Niech $R=k[x_1,\dots,x_n]$.
Z twierdzenia Hilberta o bazie możemy wybrać generatory
$I=(g_1,\dots,g_r)$. Dodajmy nową zmienną $t$ i
rozważmy ideał
$J=(g_1,\dots,g_r,1-tf)\subseteq R[t]$.
Załóżmy, że punkt $(a,b)\in k^{n+1}$ zeruje wszystkie
generatory $J$. Równości $g_j(a)=0$ oznaczają,
że $a\in\V_k(I)$, a więc z założenia $f(a)=0$.
Ostatnie równanie $1-bf(a)=0$ staje się jednak $1=0$,
co jest niemożliwe. Zatem $\V_k(J)=\varnothing$.

Słabe twierdzenie Hilberta o zerach, zastosowane w $n+1$
zmiennych, daje $J=R[t]$.
Zatem istnieją $u_1,\dots,u_r,v\in R[t]$ takie, że
\[
1=\sum_{j=1}^r u_j(x,t)g_j(x)+v(x,t)(1-tf(x)).
\]
Ponieważ $R$ jest dziedziną, możemy przejść do ciała jego
ułamków $K=\Frac(R)$, w którym niezerowy wielomian $f$
ma odwrotność. Podstawmy $t=1/f$ do tożsamości w $K$.
Ostatni składnik znika, ponieważ $1-(1/f)f=0$, i otrzymujemy
$1=\sum_{j=1}^r u_j(x,1/f)g_j(x)$.

Każde $u_j$ jest wielomianem w $t$ o współczynnikach z $R$.
Po podstawieniu $t=1/f$ jego współczynniki są ilorazami
wielomianów przez potęgi $f$. Wybieramy $N\ge1$
nie mniejsze od wszystkich występujących potęg mianownika.
Po pomnożeniu równości przez $f^N$ otrzymujemy
$f^N=\sum_{j=1}^r h_j(x)g_j(x)$ z $h_j\in R$.
A więc $f^N\in I$, jak należało wykazać.
\end{proof}

\begin{proof}[Dowód silnego twierdzenia Hilberta o zerach~\ref{tw:strong-NS}]
Oznaczmy $R=k[x_1,\dots,x_n]$.
Najpierw wykażmy $\sqrt I\subseteq\I_k(\V_k(I))$.
Jeśli $f\in\sqrt I$, to istnieje $m\ge1$ takie,
że $f^m\in I$. Dla dowolnego $a\in\V_k(I)$ mamy
$f(a)^m=f^m(a)=0$. Ponieważ $k$ jest ciałem,
$f(a)=0$, czyli $f\in\I_k(\V_k(I))$.
Ta inkluzja zachodzi nawet bez założenia algebraicznej domkniętości.

Odwrotnie, niech $f\in\I_k(\V_k(I))$.
Jeżeli $f=0$, to $f\in\sqrt I$ bez dalszych obliczeń.
Jeśli $f\ne0$, to lemat~\ref{lem:rabinowitsch}
daje $N\ge1$ z $f^N\in I$, czyli $f\in\sqrt I$.
To kończy dowód równości.
\end{proof}

\begin{przyklad}[Algebraiczne świadectwo braku punktów]
Rozważmy w $\C[x,y]$ wielomiany $f_1=x$ oraz
$f_2=xy-1$. Wspólne zera nie istnieją: z $x=0$
wynika $xy-1=-1\ne0$.
Równość $1=yx-(xy-1)$ pokazuje bezpośrednio,
że $(x,xy-1)=(1)$.
Słabe twierdzenie Hilberta o zerach gwarantuje,
że dla \emph{dowolnego} układu wielomianów bez wspólnych
zer nad $\C$ istnieje podobne świadectwo wielomianowe.
\end{przyklad}

\begin{przyklad}[Dlaczego powstaje potęga?]
Niech $I=(x^3,x^2y)\subseteq\C[x,y]$.
Wówczas $\V_{\C}(I)=\V_{\C}(x)$, bo $x^3=0$
w ciele daje $x=0$. Wielomian $x$ znika na tym zbiorze,
ale $x\notin I$: każda niezerowa suma wielokrotności
generatorów ma wszystkie monomy podzielne przez $x^2$.
Jednak $x^3\in I$, więc $x\in\sqrt I$.
Silne twierdzenie mówi dokładnie, że zawsze wystarczy
pewna dodatnia potęga znikającego wielomianu.
\end{przyklad}

\section{Co dokładnie daje Nullstellensatz?}

\begin{twierdzenie}[Korespondencja ideałów radykalnych i zbiorów algebraicznych]
\label{tw:radical-correspondence}
Niech $k$ będzie ciałem algebraicznie domkniętym, $n\ge1$
i $R=k[x_1,\dots,x_n]$.
Odwzorowania $I\mapsto\V_k(I)$ i $X\mapsto\I_k(X)$
ustanawiają wzajemnie odwrotne bijekcje pomiędzy rodziną
ideałów radykalnych $I\subseteq R$ a rodziną
zbiorów algebraicznych $X\subseteq k^n$.
Bijekcje odwracają inkluzje.
\end{twierdzenie}
\begin{proof}
Dla ideału radykalnego $I$ silne twierdzenie Hilberta o zerach
daje $\I_k(\V_k(I))=\sqrt I=I$.
Jeśli $X$ jest zbiorem algebraicznym, to $X=\V_k(J)$
dla pewnego ideału $J$. Z twierdzenia~\ref{tw:galois}
wiemy, że $\V_k(\I_k(\V_k(J)))=\V_k(J)$,
a więc $\V_k(\I_k(X))=X$.
Każdy $\I_k(X)$ jest radykalny na mocy
tego, co wykazaliśmy w twierdzeniu~\ref{tw:I-basic}.
Otrzymujemy obie potrzebne identyczności.
Odwracanie inkluzji było dowiedzione w twierdzeniu~\ref{tw:galois}.
\end{proof}

\begin{center}
\begin{tikzpicture}[>=Stealth]
\node[draw,rounded corners,fill=gray!5,align=center,minimum width=4.1cm,minimum height=1cm] (a) at (0,0)
 {ideały radykalne\\$I=\sqrt I$};
\node[draw,rounded corners,fill=gray!5,align=center,minimum width=4.1cm,minimum height=1cm] (b) at (7,0)
 {zbiory algebraiczne\\$X=\V_k(\I_k(X))$};
\draw[->,blue!70!black,thick] (a.north east) to[bend left=19]
 node[above] {$\V_k$} (b.north west);
\draw[->,red!75!black,thick] (b.south west) to[bend left=19]
 node[below] {$\I_k$} (a.south east);
\node[font=\small] at (3.5,-1.35)
 {$k$ algebraicznie domknięte};
\end{tikzpicture}
\end{center}
\begin{center}\small Rysunek 1.4. Korespondencja ideałów i zbiorów punktów po przejściu do ideałów radykalnych.\end{center}

\begin{wniosek}[Pierścień współrzędnych odczytany z równań]
\label{wn:coordinate-ring}
Niech $k$ będzie ciałem algebraicznie domkniętym,
$n\ge1$ i $I\subseteq k[x_1,\dots,x_n]$ ideałem.
Dla $X=\V_k(I)$ istnieje naturalny izomorfizm
$k[X]\cong k[x_1,\dots,x_n]/\sqrt I$.
Jeśli $I$ jest radykalny, to $k[X]\cong k[x_1,\dots,x_n]/I$.
\end{wniosek}
\begin{proof}
Z definicji $k[X]=k[x_1,\dots,x_n]/\I_k(X)$.
Z silnego twierdzenia Hilberta o zerach mamy
$\I_k(X)=\I_k(\V_k(I))=\sqrt I$. Podstawiamy tę równość.
\end{proof}

\begin{przyklad}[Nierówność obiektów nad $\R$]
Niech $I=(x^2+1)\subseteq\R[x]$.
Iloraz $\R[x]/I$ jest ciałem izomorficznym z $\C$, więc
$I$ jest ideałem maksymalnym, a zatem radykalnym.
Mimo to $\V_{\R}(I)=\varnothing$, skąd
$\I_{\R}(\V_{\R}(I))=\R[x]$, nie zaś $I$.
Jest to kontrprzykład do silnego Nullstellensatz nad
ciałem, które nie jest algebraicznie domknięte.
\end{przyklad}

\begin{przyklad}[Nierówność obiektów nad $\F_q$]
Niech $k=\F_q$, a $I=(0)\subseteq k[x]$.
Pierścień $k[x]$ jest dziedziną, więc $I$ jest radykalny.
Jednak $\V_k(I)=k$ i niezerowy wielomian $x^q-x$
znika we wszystkich punktach $k$. Stąd
$x^q-x\in\I_k(\V_k(I))$, mimo że $x^q-x\notin\sqrt I=(0)$.
Samo założenie, że $k$ jest ciałem, zdecydowanie nie wystarcza.
\end{przyklad}

\section{Nierozkładalność zbioru a pierwszość ideału}

\begin{definicja}[Zbiór nierozkładalny]
Niech $T$ będzie przestrzenią topologiczną, a $X\subseteq T$
niepustym zbiorem domkniętym.
Mówimy, że $X$ jest \emph{nierozkładalny}, jeśli dla każdych
dwóch podzbiorów domkniętych $Y,Z\subseteq X$ z
$X=Y\cup Z$ zachodzi $Y=X$ lub $Z=X$.
Jeśli istnieje rozkład na dwa właściwe zbiory domknięte,
to $X$ jest \emph{rozkładalny}.
\end{definicja}

Nierozkładalność nie oznacza tego samego co spójność topologiczna.
Zbiór nierozkładalny jest spójny, ponieważ rozkład na dwa
rozłączne niepuste zbiory jednocześnie otwarte i domknięte
dawałby w szczególności rozkład na dwa właściwe zbiory domknięte.
Odwrotna implikacja jest na ogół fałszywa.

\begin{twierdzenie}[Nierozkładalność a ideały pierwsze]
\label{tw:irreducible-prime}
Niech $k$ będzie dowolnym ciałem, $n\ge1$
i $X\subseteq k^n$ niepustym zbiorem algebraicznym.
Wtedy $X$ jest nierozkładalny w topologii Zariskiego
wtedy i tylko wtedy, gdy $\I_k(X)$ jest ideałem pierwszym.
Równoważnie: $X$ jest nierozkładalny wtedy i tylko wtedy,
gdy jego pierścień współrzędnych $k[X]$ jest dziedziną całkowitości.
\end{twierdzenie}
\begin{proof}
Najpierw załóżmy, że $X$ jest nierozkładalny.
Niech $f,g\in k[x_1,\dots,x_n]$ oraz $fg\in\I_k(X)$.
Każdy punkt $X$ spełnia $f(a)g(a)=0$,
a więc należy do $\V_k(f)$ albo do $\V_k(g)$.
Otrzymujemy rozkład
$X=(X\cap\V_k(f))\cup(X\cap\V_k(g))$
na dwa podzbiory domknięte w $X$. Nierozkładalność oznacza,
że jeden z nich jest całym $X$.
Jeśli $X\subseteq\V_k(f)$, to $f\in\I_k(X)$;
w przeciwnym razie $g\in\I_k(X)$. Ponieważ $X\ne\varnothing$,
ideał $\I_k(X)$ jest właściwy: wielomian $1$ nigdzie nie znika.
Wykazaliśmy więc, że $\I_k(X)$ jest pierwszy.
Ta implikacja nie wymaga algebraicznej domkniętości $k$.

Odwrotnie, załóżmy, że $\mathfrak p:=\I_k(X)$ jest pierwszy,
a $X=Y\cup Z$ z $Y,Z$ domkniętymi w $X$.
Ponieważ $X$ jest domknięty w $k^n$, również $Y$ i $Z$
są domknięte w $k^n$. Możemy zatem zapisać
$Y=\V_k(J)$ i $Z=\V_k(K)$ dla pewnych ideałów $J,K$.
Wtedy $X=\V_k(JK)$.
Każdy element $JK$ znika na $X$, więc $JK\subseteq\I_k(X)=\mathfrak p$.
Jeśli $J\not\subseteq\mathfrak p$, wybierzmy $f\in J\setminus\mathfrak p$.
Dla dowolnego $g\in K$ iloczyn $fg\in JK\subseteq\mathfrak p$,
zatem pierwszość $\mathfrak p$ wymusza $g\in\mathfrak p$.
Stąd $K\subseteq\mathfrak p$ i
$Z=\V_k(K)\supseteq\V_k(\mathfrak p)$.
Ponieważ $X$ jest zbiorem algebraicznym, mamy
$\V_k(\mathfrak p)=\V_k(\I_k(X))=X$
na mocy twierdzenia~\ref{tw:galois}, bez dodatkowych
założeń o ciele $k$. Zatem $Z=X$.
Jeśli zaś $J\subseteq\mathfrak p$, otrzymujemy
$Y\supseteq\V_k(\mathfrak p)=X$, więc $Y=X$.
W obu przypadkach rozkład jest trywialny.

Ostatnia równoważność wynika z kryterium ilorazowego:
ideał $\I_k(X)$ jest pierwszy wtedy i tylko wtedy,
gdy $k[x_1,\dots,x_n]/\I_k(X)=k[X]$ jest dziedziną.
\end{proof}

\begin{uwaga}[Dlaczego rozróżniamy dwa kryteria?]
Dla każdego ciała $k$ mamy równoważność:
$X$ niepusty algebraiczny nierozkładalny wtedy i tylko wtedy,
gdy \emph{jego ideał znikania} $\I_k(X)$ jest pierwszy.
Inne, silniejsze zdanie brzmi: dla dowolnego ideału pierwszego
$\mathfrak p\subseteq k[x_1,\dots,x_n]$ zbiór
$\V_k(\mathfrak p)$ jest niepusty i nierozkładalny.
To drugie zdanie wymaga dodatkowych założeń. Nad $\R$
ideał pierwszy $(x^2+1)$ ma pusty zbiór zer; nad ciałem
skończonym $\F_q$ ideał $(0)\subseteq\F_q[x]$
jest pierwszy, ale jego zbiór zer $\F_q$ ma dyskretną
topologię i jest rozkładalny, jeśli $q>1$.
\end{uwaga}

\begin{przyklad}[Parabola nierozkładalna nad $\C$]
Niech $X=\V_{\C}(y-x^2)$.
Wtedy $\I_{\C}(X)=(y-x^2)$, a iloraz
$\C[x,y]/(y-x^2)\cong\C[t]$ jest dziedziną.
Wobec tego $X$ jest nierozkładalny.
Jest to ścisłe uzasadnienie geometrycznej interpretacji
izomorfizmu, który pojawił się w rozdziale o pierścieniach.
\end{przyklad}

\begin{przyklad}[Suma dwóch osi jest rozkładalna]
Niech $X=\V_{\C}(xy)$.
Mamy $X=\V_{\C}(x)\cup\V_{\C}(y)$,
a oba zbiory po prawej są właściwymi podzbiorami domkniętymi
$X$. Ponadto $\C[X]\cong\C[x,y]/(xy)$ ma dzielniki zera:
klasy $x$ i $y$ są niezerowe, a ich iloczyn jest zerem.
Widzimy ten sam fenomen z perspektywy topologii i algebry.
\end{przyklad}

\begin{twierdzenie}[Skończony rozkład na składowe nierozkładalne]
\label{tw:finite-components}
Niech $k$ będzie dowolnym ciałem, $n\ge1$ i
$X\subseteq k^n$ niepustym zbiorem algebraicznym.
Wtedy istnieją niepuste nierozkładalne podzbiory
algebraiczne $X_1,\dots,X_r\subseteq X$ takie, że
$X=X_1\cup\cdots\cup X_r$.
Można wybrać je jako \emph{maksymalne nierozkładalne podzbiory}
$X$; po usunięciu powtórzeń wybór ten jest jednoznaczny.
\end{twierdzenie}
\begin{proof}
Najpierw pokażemy, że w topologii Zariskiego nie istnieje
nieskończony ściśle malejący ciąg zbiorów algebraicznych
$Y_1\supsetneq Y_2\supsetneq\cdots$.
Gdyby taki istniał, ideały znikania tworzyłyby ciąg
$\I_k(Y_1)\subseteq\I_k(Y_2)\subseteq\cdots$.
Inkluzje byłyby ściśle: jeśli $\I_k(Y_i)=\I_k(Y_{i+1})$,
to po zastosowaniu $\V_k$ dostalibyśmy
$Y_i=\V_k(\I_k(Y_i))=\V_k(\I_k(Y_{i+1}))=Y_{i+1}$,
sprzeczność. Ale pierścień $k[x_1,\dots,x_n]$ jest noetherowski,
więc taki nieskończenie rosnący ciąg ideałów nie istnieje.

Załóżmy teraz, że istnieje niepusty zbiór algebraiczny,
który \emph{nie} daje się zapisać jako skończona suma zbiorów
nierozkładalnych. Spośród wszystkich takich zbiorów wybierzmy
minimalny względem inkluzji: wybór jest możliwy dzięki
właśnie udowodnionemu warunkowi braku nieskończonych ściśle
malejących ciągów. Niech ten zbiór nazywa się $Y$.
Z definicji kontrprzykładu $Y$ nie jest nierozkładalny,
a więc $Y=Y'\cup Y''$ dla właściwych domkniętych
podzbiorów $Y',Y''\subsetneq Y$.
Z minimalności $Y$ oba zbiory $Y',Y''$ mają skończone
rozkłady nierozkładalne. Łącząc te rozkłady, otrzymujemy
taki rozkład również dla $Y$, sprzeczność.

Mamy więc skończony rozkład $X=Z_1\cup\cdots\cup Z_s$.
Jeśli któryś $Z_i$ zawiera się w innym $Z_j$, możemy go usunąć.
Powtarzając tę czynność, otrzymujemy rozkład bez zawierań.
Każdy pozostały $Z_i$ jest maksymalnym nierozkładalnym
podzbiorem $X$: gdyby istniał nierozkładalny
$Z_i\subsetneq W\subseteq X$, to z
$W=\bigcup_j(W\cap Z_j)$ i nierozkładalności $W$
wynikałoby $W\subseteq Z_j$ dla pewnego $j$.
Zatem $Z_i\subsetneq W\subseteq Z_j$, co przeczy brakowi zawierań.

Aby uzasadnić jednoznaczność, zauważmy najpierw, że jeśli
nierozkładalny zbiór $W$ zawiera się w skończonej sumie
domkniętych zbiorów $Z_1\cup\cdots\cup Z_s$,
to $W\subseteq Z_j$ dla pewnego $j$: dla $s=2$ jest to definicja,
a dla większego $s$ wynika przez indukcję po $s$.
Niech $X=W_1\cup\cdots\cup W_t$ będzie drugim rozkładem
na maksymalne nierozkładalne zbiory. Każdy $Z_i$ zawiera się
w pewnym $W_j$ na mocy poprzedniego faktu, a maksymalność
$Z_i$ wymusza $Z_i=W_j$. Zamieniając role rozkładów,
otrzymujemy tę samą rodzinę podzbiorów z dokładnością
do kolejności.
\end{proof}

\begin{przyklad}[Trzy proste]
Zbiór $X=\V_{\C}(xy(x-y))\subseteq\C^2$
jest sumą zbiorów trzech prostych: $x=0$, $y=0$ oraz $x=y$.
Każda jest nierozkładalna, bo jej pierścień współrzędnych
jest izomorficzny z $\C[t]$. Żadna nie zawiera innej,
więc są to dokładnie trzy składowe nierozkładalne $X$.
\end{przyklad}

\section{Morfizmy afiniczne i odwrócenie kierunku odwzorowań}

\subsection{Wielomianowe odwzorowania zbiorów algebraicznych}

\begin{definicja}[Odwzorowanie regularne zbiorów afinicznych]
Niech $k$ będzie ciałem algebraicznie domkniętym, $n,m\ge1$,
a $X\subseteq k^n$ i $Y\subseteq k^m$ niepustymi zbiorami algebraicznymi.
Odwzorowanie $F:X\to Y$ nazywamy \emph{regularnym} (w sensie
afinicznym), jeśli istnieją wielomiany
$f_1,\dots,f_m\in k[x_1,\dots,x_n]$ takie, że
$F(a)=(f_1(a),\dots,f_m(a))$ dla każdego $a\in X$.
Wielomiany $f_i$ nie muszą być jednoznaczne, ale ich klasy w $k[X]$
są jednoznacznie określone przez $F$.
\end{definicja}

\begin{przyklad}[Parametryzacja paraboli]
Niech $k$ będzie ciałem algebraicznie domkniętym i $X=\A^1(k)=k$ i
$Y=\V_k(y-x^2)\subseteq k^2$.
Odwzorowanie $F:k\to Y$, $t\mapsto(t,t^2)$, jest regularne,
bo jego współrzędne są wielomianami. Odwzorowanie odwrotne
$G:Y\to k$, $(x,y)\mapsto x$, jest również regularne,
a złożenia $G\circ F$ i $F\circ G$ są identycznościami.
Parabola jest więc \emph{izomorficzna jako zbiór algebraiczny}
z prostą afiniczną. Sam kształt paraboli na rysunku nie jest
tutaj niezmiennikiem geometrycznym: ważna jest struktura funkcji
wielomianowych.
\end{przyklad}

\begin{definicja}[Homomorfizm $k$-algebr]
Niech $A,B$ będą pierścieniami zawierającymi ciało $k$ jako podpierścień.
Homomorfizm pierścieni $\varphi:A\to B$ nazywamy
\emph{homomorfizmem $k$-algebr}, gdy $\varphi(c)=c$
dla każdego $c\in k$.
\end{definicja}

\begin{twierdzenie}[Kontrawariantny opis odwzorowań regularnych]
\label{tw:affine-morphism}
Niech $k$ będzie ciałem algebraicznie domkniętym, $n,m\ge1$,
a $X\subseteq k^n$ i $Y\subseteq k^m$ niepustymi zbiorami
algebraicznymi. Wtedy każdemu odwzorowaniu regularnemu
$F:X\to Y$ odpowiada homomorfizm $k$-algebr
\[
F^*:k[Y]\longrightarrow k[X],\qquad
F^*([h])=[h\circ F],
\]
gdzie $[h]$ oznacza klasę funkcji wielomianowej.
To przyporządkowanie jest bijekcją między
odwzorowaniami regularnymi $X\to Y$ a homomorfizmami
$k$-algebr $k[Y]\to k[X]$.
Jeżeli ponadto $q\ge1$, $Z\subseteq k^q$ jest niepustym
zbiorem algebraicznym, a $G:Y\to Z$ odwzorowaniem regularnym,
to $(G\circ F)^*=F^*\circ G^*$.
\end{twierdzenie}
\begin{proof}
Niech $F=(f_1,\dots,f_m)$.
Jeśli $[h]=[h']$ w $k[Y]$, to wielomiany $h,h'$ mają
te same wartości w każdym punkcie $Y$. Ponieważ $F(a)\in Y$
dla każdego $a\in X$, otrzymujemy
$h(F(a))=h'(F(a))$ dla wszystkich $a\in X$.
Zatem $[h\circ F]=[h'\circ F]$ w $k[X]$ i wzór na $F^*$
nie zależy od reprezentanta.
Podstawienie wielomianów zachowuje sumę, iloczyn i stałe,
więc $F^*$ jest homomorfizmem $k$-algebr.

W przeciwnym kierunku niech
$\varphi:k[Y]\to k[X]$ będzie homomorfizmem $k$-algebr.
Oznaczmy przez $[y_1],\dots,[y_m]$ klasy funkcji współrzędnych
na $Y$. Dla każdego $j$ wybierzmy wielomian
$f_j\in k[x_1,\dots,x_n]$ reprezentujący
$\varphi([y_j])\in k[X]$.
Zdefiniujmy $F:X\to k^m$ przez
$F(a)=(f_1(a),\dots,f_m(a))$.
Musimy sprawdzić, że faktycznie $F(X)\subseteq Y$.
Dla dowolnego $h\in\I_k(Y)$ klasa $[h]$ w $k[Y]$ jest zerowa.
Stąd $\varphi([h])=0$ w $k[X]$.
Ponieważ $h$ jest wielomianem w zmiennych $y_1,\dots,y_m$,
mamy $\varphi([h])=[h(f_1,\dots,f_m)]$.
Ta klasa jest zerowa dokładnie wtedy, gdy
$h(f_1(a),\dots,f_m(a))=0$ dla wszystkich $a\in X$.
Oznacza to, że każdy punkt $F(a)$ zeruje wszystkie wielomiany
ideału $\I_k(Y)$. Ponieważ $Y$ jest algebraiczny,
$Y=\V_k(\I_k(Y))$, a więc $F(a)\in Y$.
Otrzymaliśmy odwzorowanie regularne $F:X\to Y$.

Skonstruowany homomorfizm $F^*$ jest równy $\varphi$,
ponieważ oba mają te same wartości na stałych z $k$
i na klasach $[y_j]$, które generują $k[Y]$ jako $k$-algebrę.
Z drugiej strony odwzorowanie regularne $F$ jest jednoznacznie
wyznaczone przez jego $m$ funkcji współrzędnych na $X$.
Obie konstrukcje są więc odwrotne.

Dla złożenia i dowolnego $h\in k[Z]$ mamy
$((G\circ F)^*h)(a)=h(G(F(a)))$
oraz $(F^*(G^*h))(a)=(G^*h)(F(a))=h(G(F(a)))$.
Zatem $(G\circ F)^*=F^*\circ G^*$.
\end{proof}

\begin{center}
\begin{tikzpicture}[>=Stealth]
\node (x) at (0,1.3) {$X$};
\node (y) at (5.5,1.3) {$Y$};
\draw[->,thick,blue!70!black] (x)--node[above] {$F$}(y);
\node (kx) at (0,-.55) {$k[X]$};
\node (ky) at (5.5,-.55) {$k[Y]$};
\draw[->,thick,red!70!black] (ky)--node[below] {$F^*$}(kx);
\draw[densely dashed,gray!70] (0,.95)--(0,-.16);
\draw[densely dashed,gray!70] (5.5,.95)--(5.5,-.16);
\end{tikzpicture}
\end{center}
\begin{center}\small Rysunek 1.5. Odwzorowanie punktów $F:X\to Y$ indukuje homomorfizm pierścieni w kierunku przeciwnym.\end{center}

\begin{uwaga}[Dlaczego strzałka biegnie w drugą stronę?]
Jeśli mamy funkcję $h$ określoną na $Y$ i odwzorowanie
$F:X\to Y$, możemy złożyć je jako $h\circ F$,
otrzymując funkcję na $X$. Innymi słowy, funkcje na \emph{celu}
można cofać do funkcji na \emph{źródle}.
Wiele fundamentalnych konstrukcji geometrii algebraicznej
ma właśnie taki kontrawariantny charakter.
\end{uwaga}

\subsection{Zaskakujący przykład: bijekcja nie musi być izomorfizmem}

\begin{przyklad}[Krzywa z ostrzem i jej parametryzacja]
\label{pr:cusp}
Niech $k$ będzie ciałem algebraicznie domkniętym.
Rozważmy krzywą
$C=\V_k(y^2-x^3)\subseteq k^2$ oraz odwzorowanie
$\nu:k\to C$, $t\mapsto(t^2,t^3)$.
Jest ono regularne i rzeczywiście przyjmuje wartości w $C$,
bo $(t^3)^2-(t^2)^3=0$.

Pokażemy, że $\nu$ jest bijekcją na zbiorach punktów.
Jeżeli $(x,y)\in C$ i $x=0$, to $y^2=0$, a więc $y=0$;
wówczas jedynym parametrem jest $t=0$.
Jeśli $x\ne0$, definiujemy $t=y/x$. Z równania $y^2=x^3$
wynika $t^2=y^2/x^2=x$, a następnie
$t^3=t\,t^2=(y/x)x=y$. Zatem każdy punkt krzywej ma postać
$(t^2,t^3)$. Jednoznaczność także wynika z rachunku:
dla punktu różnego od $(0,0)$ parametr musi wynosić $y/x$.

Mimo to $\nu$ \emph{nie jest izomorfizmem algebraicznym}.
Aby tego dowieść, obliczymy homomorfizm pierścieni
$\Phi:k[x,y]\to k[t]$, $\Phi(x)=t^2$, $\Phi(y)=t^3$.
Ma on obraz $k[t^2,t^3]$.
Pokażemy, że jego jądro jest równe $(y^2-x^3)$.
Dowolny wielomian $f(x,y)$ możemy zredukować modulo
$y^2-x^3$ do postaci $A(x)+yB(x)$:
za każdym razem zastępujemy $y^m$, $m\ge2$, przez
$x^3y^{m-2}$, używając równania $y^2=x^3$.
Jeśli $\Phi(f)=0$, to po podstawieniu mamy
$A(t^2)+t^3B(t^2)=0$.
Pierwsza suma zawiera tylko parzyste potęgi $t$,
a druga tylko nieparzyste potęgi $t$ stopnia co najmniej $3$.
Nie mogą się więc wzajemnie znosić. Wszystkie współczynniki
obu wielomianów są zerowe, czyli $A=B=0$.
To znaczy $f\in(y^2-x^3)$.
Ponieważ ciało algebraicznie domknięte $k$ jest nieskończone,
dla dowolnego wielomianu $f$ warunek, że $f$ znika na
całym $C$, jest (dzięki surjektywności $\nu$) równoważny temu,
że wielomian $\Phi(f)(t)=f(t^2,t^3)$ znika dla każdego
$t\in k$, czyli jest zerowy jako wielomian.
Stąd $\I_k(C)=\kerf\Phi=(y^2-x^3)$.
W konsekwencji pierwsze twierdzenie o izomorfizmie daje
\[
k[C]\cong k[x,y]/(y^2-x^3)\cong k[t^2,t^3]\subsetneq k[t].
\]
Ostatnia inkluzja jest ścisła, ponieważ $t\notin k[t^2,t^3]$:
każdy monom o dodatnim stopniu w generatorach $t^2,t^3$
ma stopień co najmniej $2$, zatem z ich skończonej sumy
nie otrzymamy wielomianu stopnia $1$.
Gdyby $\nu$ miała regularną odwrotność, odpowiadający jej
homomorfizm $\nu^*:k[C]\to k[t]$ musiałby być izomorfizmem,
a nie jest nawet surjektywny.
\end{przyklad}

\begin{center}
\begin{tikzpicture}[>=Stealth,scale=.78]
\draw[->,gray!75] (-.25,0)--(4.8,0) node[right] {$x$};
\draw[->,gray!75] (0,-3.3)--(0,3.4) node[above] {$y$};
\draw[blue!75!black,very thick,domain=-1.63:1.63,samples=120,smooth]
 plot ({\x*\x},{\x*\x*\x});
\fill (0,0) circle (2.3pt);
\draw[densely dashed] (1,0)--(1,1);
\fill[red!75!black] (1,1) circle(2.3pt);
\node[right,red!75!black] at (1,1) {$(t^2,t^3)$ dla $t=1$};
\node[blue!75!black] at (3.1,2.8) {$y^2=x^3$};
% Objaśnienie przeniesiono do podpisu, aby nie przysłaniało dolnej gałęzi.
\end{tikzpicture}
\end{center}
\begin{center}\small Rysunek 1.6. Rzeczywisty obraz krzywej z ostrzem; nad $\C$ krzywa ma również punkty zespolone.\end{center}

\begin{uwaga}[Czym różni się ten przykład od paraboli?]
Dla paraboli parametr $t$ jest samą współrzędną $x$,
więc można go odzyskać wielomianowo. Na krzywej z ostrzem
parametr poza początkiem jest równy $y/x$.
Jest to funkcja wymierna, której zapis ma w początku
zerowy mianownik. Nie istnieje globalny opis tego parametru
przez element pierścienia $k[C]$.
W rozdziale o osobliwościach wrócimy do tego przykładu:
$\nu$ jest prototypem \emph{normalizacji} krzywej.
\end{uwaga}

\subsection{Obraz odwzorowania regularnego nie musi być domknięty}

\begin{przyklad}[Rzut hiperboli na oś]
Niech $k$ będzie algebraicznie domkniętym ciałem i
$H=\V_k(xy-1)\subseteq k^2$.
Rozważmy rzut $p:H\to k$, $p(x,y)=x$.
Jest to odwzorowanie regularne. Jego obraz jest równy
$k^\times=k\setminus\{0\}$: punkt hiperboli ma $x\ne0$,
a dla każdego $x\ne0$ punkt $(x,x^{-1})$ leży na $H$.

Zbiór $k\setminus\{0\}$ nie jest domknięty w topologii
Zariskiego na $k$, ponieważ $k$ jest nieskończone,
a jedyne właściwe zbiory domknięte na prostej są skończone.
Obraz odwzorowania regularnego może więc być otwarty, ale niedomknięty.
To ważne ostrzeżenie przed bezkrytycznym przenoszeniem własności
odwzorowań właściwych z topologii na wszystkie morfizmy algebraiczne.
\end{przyklad}

\begin{center}
\begin{tikzpicture}[>=Stealth,scale=.79]
\draw[->,gray!75] (-3.8,0)--(3.9,0) node[right] {$x$};
\draw[->,gray!75] (0,-3.25)--(0,3.4) node[above] {$y$};
\draw[blue!70!black,very thick,domain=.32:3.1,samples=100]
 plot (\x,{1/\x});
\draw[blue!70!black,very thick,domain=-3.1:-.32,samples=100]
 plot (\x,{1/\x});
\draw[->,red!70!black,thick] (2,0.5)--(2,0.08);
\draw[->,red!70!black,thick] (-2,-0.5)--(-2,-0.08);
\draw[fill=white,thick] (0,0) circle(4pt);
\node[blue!70!black] at (2.75,2.2) {$xy=1$};
% Objaśnienie rzutu jest w podpisie rysunku.
\end{tikzpicture}
\end{center}
\begin{center}\small Rysunek 1.7. Rzut punktów hiperboli na oś współrzędnych pomija zero.\end{center}

\section{Kolejne intuicje: wymiar i kierunek dalszej teorii}

\begin{definicja}[Wymiar przez łańcuchy nierozkładalnych podzbiorów]
Niech $k$ będzie ciałem algebraicznie domkniętym,
$n\ge1$, a $X\subseteq k^n$ zbiorem algebraicznym.
Dla $X\ne\varnothing$ definiujemy \emph{wymiar} $X$ jako kres górny liczb
całkowitych $r\ge0$, dla których istnieją
nierozkładalne, niepuste, domknięte podzbiory
$X_0\subsetneq X_1\subsetneq\cdots\subsetneq X_r\subseteq X$.
Zapisujemy $\dim X$. Gdy $X=\varnothing$, przyjmujemy
umownie $\dim X=-\infty$.
\end{definicja}

\begin{przyklad}[Punkt ma wymiar zero]
Jeśli $X=\{a\}$ jest pojedynczym punktem, to jedynym
niepustym podzbiorem domkniętym $X$ jest $X$.
Zatem istnieje łańcuch długości $0$ (samo $X$),
ale nie istnieje łańcuch długości $1$.
Stąd $\dim\{a\}=0$.
\end{przyklad}

\begin{przyklad}[Prosta afiniczna ma wymiar jeden]
Niech $k$ będzie algebraicznie domkniętym ciałem.
Ponieważ $k$ jest nieskończone, właściwe domknięte
podzbiory $\A^1(k)$ są skończone.
Cała prosta jest nierozkładalna: nie można zapisać
nieskończonego $k$ jako sumy dwóch skończonych zbiorów.
Wybierzmy dowolne $a\in k$. Łańcuch
$\{a\}\subsetneq\A^1(k)$ pokazuje, że wymiar jest co najmniej $1$.
Nie istnieje jednak dłuższy łańcuch nierozkładalnych
podzbiorów: każdy właściwy niepusty nierozkładalny zbiór
domknięty na prostej jest singletonem, ponieważ skończony
zbiór o co najmniej dwóch punktach można rozłożyć na dwa
właściwe domknięte podzbiory.
Zatem $\dim\A^1(k)=1$.
\end{przyklad}

Dla $\A^n(k)$ zachodzi $\dim\A^n(k)=n$, ale pełne
udowodnienie tego faktu wymaga dalszych narzędzi algebry przemiennej
(np. normalizacji Noethera lub teorii wysokości ideałów).
Nie będziemy zastępować brakującego dowodu intuicją z rysunku.
Szczegółową teorię wymiaru wprowadzimy przed badaniem
osobliwości i krzywych algebraicznych.

\section{Podsumowanie: sieć zależności między pojęciami}

Oto logika rozdziału. Dla dowolnego ciała $k$ i
$R=k[x_1,\dots,x_n]$ możemy przejść od zbioru równań
$S\subseteq R$ do zbioru ich zer $\V_k(S)$, a od zbioru
punktów $X\subseteq k^n$ do ideału znikania $\I_k(X)$.
Każde z tych przejść odwraca inkluzje. Ponadto
$\overline X^{\mathrm{Zar}}=\V_k(\I_k(X))$,
a $\I_k(X)$ jest zawsze radykalny.

Gdy \emph{dodatkowo} $k$ jest algebraicznie domknięte,
silne twierdzenie Hilberta o zerach gwarantuje
$\I_k(\V_k(I))=\sqrt I$ dla każdego ideału $I$.
Dlatego zbiory algebraiczne odpowiadają dokładnie
ideałom radykalnym, a nie wszystkim ideałom.
Z kolei zbiory nierozkładalne odpowiadają ideałom pierwszym.
Pierścienie współrzędnych przenoszą geometrię na grunt algebry,
a odwzorowaniom regularnym odpowiadają homomorfizmy
pierścieni biegnące w przeciwnym kierunku.
% ======================================================================
% TRZECI ROZDZIAŁ CAŁOŚCI: NUMER 2, PONIEWAŻ SKRYPT ZACZYNA SIĘ OD 0.
% ======================================================================

\chapter{Geometria rzutowa: punkty w nieskończoności, krzywe i przecięcia}
\label{ch:rzutowa}

\section{Po co domykać przestrzeń afiniczną?}

Na płaszczyźnie afinicznej $\A^2(k)$ proste $y=0$ i $y=1$ nie
mają wspólnego punktu. Z drugiej strony proste $y=0$ i $y=\varepsilon x$
przecinają się w początku układu współrzędnych dla każdego
$\varepsilon\ne0$. Wyniki dotyczące liczby punktów przecięcia
stają się przez to niewygodne: musimy osobno omawiać proste równoległe,
a dla krzywych wyższych stopni mogą znikać kolejne przecięcia.

W geometrii rzutowej dodajemy punkty reprezentujące \emph{kierunki}.
Nie należy przy tym wyobrażać sobie dodatkowej zwykłej prostej
w tej samej płaszczyźnie euklidesowej. Otrzymujemy nową przestrzeń,
w której punkt jest \emph{klasą niezerowego wektora}.

\begin{center}
\begin{tikzpicture}[>=Stealth,scale=.93]
  % Schematic of two parallel affine lines: equations agree with vertical coordinates.
  \draw[->,gray!65] (-3.2,-.55)--(-3.2,1.6) node[above] {$y$};
  \draw[->,gray!65] (-3.25,0)--(3.65,0) node[right] {$x$};
  \draw[blue!70!black,very thick] (-2.9,0)--(3.25,0);
  \draw[red!75!black,very thick] (-2.9,1)--(3.25,1);
  \node[blue!70!black,anchor=south west,fill=white,inner sep=1pt]
    at (-2.7,0.05) {$y=0$};
  \node[red!75!black,anchor=south west,fill=white,inner sep=1pt]
    at (-2.7,1.05) {$y=1$};
  \draw[->,dashed,gray!70] (3.3,0.02)--(4.35,.42);
  \draw[->,dashed,gray!70] (3.3,1)--(4.35,.63);
  \node[draw,circle,fill=white,inner sep=2pt] at (4.75,.54) {$p_\infty$};
\end{tikzpicture}

\small Rysunek 2.1. Szkic idei: dwie proste o tym samym kierunku mają
jeden wspólny punkt na prostej w nieskończoności. Nie jest to rysunek
w metryce przestrzeni rzutowej.
\end{center}

\section{Przestrzeń rzutowa jako zbiór kierunków}

\begin{definicja}[Przestrzeń rzutowa]
\label{def:projective-space}
Niech $k$ będzie ciałem, a $n\ge0$. Na zbiorze
$k^{n+1}\setminus\{0\}$ definiujemy relację
\[
 (a_0,\ldots,a_n)\sim(b_0,\ldots,b_n)
 \quad\Longleftrightarrow\quad
 \exists\lambda\in k^\times\ \forall i\in\{0,\ldots,n\}:\ b_i=\lambda a_i.
\]
Jej klasy oznaczamy $[a_0:\cdots:a_n]$. \emph{Przestrzenią
rzutową wymiaru $n$ nad $k$} na poziomie punktów $k$-wymiernych
nazywamy $\PP^n(k)=(k^{n+1}\setminus\{0\})/\!\sim$.
Współrzędne $[a_0:\cdots:a_n]$ nazywamy \emph{jednorodnymi}.
\end{definicja}

W definicji występuje kwantyfikator $\lambda\ne0$. Nie możemy
utożsamiać wszystkich niezerowych wektorów ze sobą: dwa wektory
reprezentują ten sam punkt tylko wtedy, gdy leżą na tej samej
jednowymiarowej podprzestrzeni liniowej. Wektor zerowy nie reprezentuje
żadnego punktu. Dlatego $[1:2]=[2:4]$ nad $\R$, ale $[1:0]\ne[0:1]$.

\begin{stwierdzenie}[Dobrze określona relacja]
Dla każdego ciała $k$ i $n\ge0$ relacja z definicji
\ref{def:projective-space} jest relacją równoważności.
\end{stwierdzenie}
\begin{proof}
Zwrotność otrzymujemy dla $\lambda=1$. Jeśli $b=\lambda a$
przy $\lambda\ne0$, to $a=\lambda^{-1}b$, więc relacja jest
symetryczna. Jeżeli $b=\lambda a$ oraz $c=\mu b$, to
$c=(\mu\lambda)a$, a $\mu\lambda\ne0$, co dowodzi
przechodniości.
\end{proof}

\begin{przyklad}[Prosta rzutowa i jeden punkt uzupełniający]
W $\PP^1(k)$ każdy punkt ma co najmniej jedną niezerową
współrzędną. Jeśli $a_0\ne0$, otrzymujemy
$[a_0:a_1]=[1:a_1/a_0]$, a więc punkt odpowiada dokładnie
jednemu $t\in k$. Jedyny punkt, którego nie można tak zapisać,
to $[0:1]$. Stąd bijekcja zbiorów
$\PP^1(k)\cong k\sqcup\{\infty\}$.

Przy $k=\R$ nie jest to jeszcze twierdzenie o homeomorfizmie:
trzeba osobno określić, czy używamy topologii Zariskiego,
czy topologii otrzymanej z rzeczywistej przestrzeni wektorowej.
\end{przyklad}

\begin{przyklad}[Projektowe spotkanie dwóch prostych]
W $\PP^2(k)$ ustalmy odtąd kolejność współrzędnych
$[X:Y:Z]$, gdzie $Z$ jest współrzędną homogenizującą.
Równanie afiniczne $y=0$ przepisuje się jako $Y=0$,
a $y=1$ jako $Y-Z=0$. Wspólne rozwiązania spełniają
$Y=Z=0$, co daje dokładnie punkt $[1:0:0]$.
Jest on niedostępny w części $Z\ne0$ i reprezentuje kierunek
osi $x$.
\end{przyklad}

\section{Mapy afiniczne przestrzeni rzutowej}

\begin{definicja}[Standardowe mapy afiniczne]
Niech $k$ będzie ciałem, a $n\ge1$ liczbą całkowitą.
Dla $i\in\{0,\ldots,n\}$ niech
$U_i=\{[a_0:\cdots:a_n]\in\PP^n(k):a_i\ne0\}$.
Definiujemy $\varphi_i:U_i\to k^n$ przez podzielenie
wszystkich współrzędnych oprócz $a_i$ przez $a_i$,
z zachowaniem ich kolejności.
\end{definicja}

\begin{stwierdzenie}[Mapy pokrywają przestrzeń]
Dla każdego ciała $k$ oraz $n\ge1$ zachodzi
$\PP^n(k)=\bigcup_{i=0}^n U_i$; każde $\varphi_i$
jest bijekcją. Na $U_i\cap U_j$ przejścia współrzędnych
są opisane ilorazami współrzędnych z nieznikającym mianownikiem.
\end{stwierdzenie}
\begin{proof}
Punkt rzutowy ma niezerową współrzędną, więc należy do
przynajmniej jednego $U_i$. Ilorazy $a_j/a_i$ nie zależą
od reprezentanta, ponieważ przemnożenie wszystkich współrzędnych
przez $\lambda\ne0$ skraca się w liczniku i mianowniku.
Odwrotność $\varphi_i$ otrzymujemy, wstawiając $1$ na
i-tej pozycji. Oba złożenia są identycznościami, gdyż w jedną
stronę dzielimy przez $a_i$, a w drugą wybieramy reprezentanta
z tą współrzędną równą $1$. Na przecięciu dwóch map trzeba
wykonać te operacje kolejno, dlatego powstają ilorazy.
\end{proof}

\begin{przyklad}[Zmiana map na $\PP^1$]
Dla $[s:t]\in\PP^1(k)$ mamy $U_s=\{s\ne0\}$ i
$U_t=\{t\ne0\}$. Współrzędna na pierwszej mapie to
$u=t/s$, a na drugiej $v=s/t$. Na przecięciu $u\ne0$
i zachodzi $v=1/u$. Nie jest to wielomian na całej prostej,
ale jest funkcją regularną na otwartym zbiorze, na którym
mianownik nie znika. To powód, dla którego w geometrii
rzutowej nie wystarczą globalne wielomiany afiniczne.
\end{przyklad}

Dla $[X:Y:Z]\in\PP^2(k)$ utożsamiamy
$U_Z=\{Z\ne0\}$ z $\A^2(k)$ przez
$[X:Y:Z]\mapsto(X/Z,Y/Z)$. Dopełnieniem tej części
jest \emph{prosta w nieskończoności} $H_\infty=\{Z=0\}$,
izomorficzna z $\PP^1(k)$ przez $[X:Y:0]\mapsto[X:Y]$.

\section{Wielomiany jednorodne i algebraiczne podzbiory rzutowe}

\begin{definicja}[Wielomian jednorodny]
Niech $k$ będzie ciałem, $n\ge0$ liczbą całkowitą
i $S=k[X_0,\ldots,X_n]$.
Wielomian $F\in S$ nazywamy \emph{jednorodnym stopnia $d$},
gdy jest sumą jednomianów o łącznym stopniu $d$; wielomian
zerowy traktujemy jako jednorodny dowolnego stopnia.
Dla niezerowego $F$ oznaczamy jego stopień przez $\deg F=d$.
\end{definicja}

\begin{lemat}[Dlaczego jednorodność jest konieczna]
Niech $k$ będzie ciałem, a $n,d\ge0$ liczbami całkowitymi.
Jeśli $F\in k[X_0,\dots,X_n]$ jest jednorodny stopnia $d$,
to dla każdego $a\in k^{n+1}$ i $\lambda\in k$ zachodzi
$F(\lambda a)=\lambda^dF(a)$. W szczególności warunek
$F(a)=0$ nie zależy od wybranego niezerowego reprezentanta
punktu $[a]\in\PP^n(k)$.
\end{lemat}
\begin{proof}
Dla jednomianu $cX_0^{e_0}\cdots X_n^{e_n}$
o $e_0+\cdots+e_n=d$ podstawienie $\lambda a_i$
mnoży wartość przez $\lambda^{\sum e_i}=\lambda^d$.
Po zsumowaniu jednomianów otrzymujemy ten sam wzór dla $F$.
Gdy $\lambda\ne0$, mamy $\lambda^d\ne0$,
więc $F(\lambda a)=0$ wtedy i tylko wtedy, gdy $F(a)=0$.
\end{proof}

\begin{uwaga}[Wielomian niejednorodny nie określa równania rzutowego]
Nad $k=\R$ warunek $X+1=0$ nie jest dobrze określony dla
$[Z:X]\in\PP^1(\R)$: reprezentant $(1,-1)$ go spełnia,
a przemnożony przez $2$ reprezentant $(2,-2)$ już nie.
Poprawne równanie jednorodne ma postać $X+Z=0$.
\end{uwaga}

\begin{definicja}[Zbiór algebraiczny rzutowy]
Niech $k$ będzie ciałem, $n\ge0$, $S=k[X_0,\dots,X_n]$ i niech $T\subseteq S$ będzie
rodziną wielomianów jednorodnych. Definiujemy
$\V_+(T)=\{[a]\in\PP^n(k):F(a)=0\text{ dla każdego }F\in T\}$.
Podzbiór $X\subseteq\PP^n(k)$ nazywamy \emph{rzutowym
zbiorem algebraicznym}, gdy $X=\V_+(T)$ dla pewnej takiej rodziny.
Ideał $I\subseteq S$ jest \emph{jednorodny}, gdy jest generowany
przez wielomiany jednorodne.
\end{definicja}

\begin{twierdzenie}[Topologia Zariskiego na przestrzeni rzutowej]
\label{tw:projective-zariski}
Dla każdego ciała $k$ i liczby całkowitej $n\ge0$ zbiory algebraiczne rzutowe
są zbiorami domkniętymi pewnej topologii na $\PP^n(k)$.
Dla ideałów jednorodnych $I,J\subseteq S=k[X_0,\dots,X_n]$ zachodzą
$\V_+(I+J)=\V_+(I)\cap\V_+(J)$ oraz
$\V_+(IJ)=\V_+(I)\cup\V_+(J)$.
\end{twierdzenie}
\begin{proof}
Wielomian zerowy określa całą przestrzeń, a stała $1$
zbiór pusty. Dowolne przecięcie zbiorów zerowych odpowiada
połączeniu rodzin wielomianów. Dla sumy dwóch zbiorów
można powtórzyć argument z twierdzenia~\ref{tw:zariski-top}:
$F(a)G(a)=0$ w ciele wymusza $F(a)=0$ lub $G(a)=0$.
Mamy więc wymaganą stabilność na przecięcia dowolne i sumy
skończone. Równości dotyczące ideałów wynikają z tych samych
własności, ponieważ ideały sumy i iloczynu ideałów jednorodnych
są również jednorodne.
\end{proof}

\begin{uwaga}[Punkty a ideał nieistotny]
Ideał $\mathfrak m_+=(X_0,\dots,X_n)$ nazywa się
\emph{ideałem nieistotnym} (ang. \emph{irrelevant ideal}).
Mamy $\V_+(\mathfrak m_+)=\varnothing$: wektor reprezentujący
punkt rzutowy nie może mieć wszystkich współrzędnych zerowych.
W przeciwieństwie do przypadku afinicznego nie każdy właściwy
ideał jednorodny ma zatem niepusty zbiór punktów rzutowych.
\end{uwaga}

\begin{stwierdzenie}[Topologia na mapach]
Niech $k$ będzie ciałem, $n\ge1$ i $i\in\{0,\dots,n\}$.
Wtedy bijekcja $\varphi_i:U_i\to k^n$
przenosi topologię Zariskiego indukowaną z $\PP^n(k)$
na topologię Zariskiego na $k^n$.
\end{stwierdzenie}
\begin{proof}
Wystarczy omówić $i=0$. Ograniczenie jednorodnego równania
$F(X_0,\dots,X_n)=0$ do $X_0\ne0$ jest równoważne
afinicznemu równaniu $f(x_1,\dots,x_n)=F(1,x_1,\dots,x_n)=0$.
Zatem przecięcie rzutowego zbioru algebraicznego z $U_0$
przechodzi w zbiór afinicznie domknięty. Odwrotnie,
dla wielomianu zerowego wybieramy jednorodny wielomian
zerowy. Dla niezerowego $f\in k[x_1,\dots,x_n]$ stopnia $d$
wielomian $X_0^df(X_1/X_0,\dots,X_n/X_0)$
(po rozwinięciu) jest jednorodny i wyznacza na $U_0$
taki sam zbiór zerowy jak $f$. Dowodzi to równości
topologii na podprzestrzeniach.
\end{proof}

\section{Homogenizacja i domknięcie rzutowe}

W tej części współrzędna homogenizująca $Z$ zajmuje
\emph{ostatnią} pozycję: punkty zapisujemy jako
$[X_1:\cdots:X_n:Z]$, a część afiniczna odpowiada $Z\ne0$.
W szczególności na płaszczyźnie mamy uzgodniony zapis $[X:Y:Z]$.

\begin{definicja}[Homogenizacja]
Niech $k$ będzie ciałem, a $n\ge1$ liczbą całkowitą.
Dla niezerowego wielomianu $f\in k[x_1,\ldots,x_n]$
stopnia całkowitego $d$ definiujemy
\[
 f^h(X_1,\ldots,X_n,Z)
 =Z^d f(X_1/Z,\ldots,X_n/Z),
\]
rozumiejąc prawą stronę po skróceniu jako wielomian;
$f^h$ jest jednorodny stopnia $d$. Dla $f=0$ przyjmujemy
$f^h=0$. Operację $F\mapsto F(x_1,\ldots,x_n,1)$
nazywamy \emph{dehomogenizacją} względem $Z$.
\end{definicja}

\begin{przyklad}[Dwa równoległe równania]
Na płaszczyźnie afinicznej $y=0$ i $y=1$ mają równania
$f=y$ i $g=y-1$. Ich homogenizacje względem $Z$ to
$F=Y$ oraz $G=Y-Z$.
W $\PP^2(k)$ przecinają się w punkcie $[1:0:0]$.
Zatem dodanie punktów w nieskończoności odzyskuje
przecięcie, którego brakowało na płaszczyźnie afinicznej.
\end{przyklad}

\begin{definicja}[Homogenizacja ideału]
Niech $k$ będzie ciałem, $n\ge1$, a
$I\subseteq k[x_1,\ldots,x_n]$ ideałem.
Przez $I^h$ oznaczamy ideał jednorodny w
$k[X_1,\ldots,X_n,Z]$ generowany przez wszystkie $f^h$
z $f\in I$ (a nie tylko przez homogenizacje dowolnie
wybranego układu generatorów $I$).
\end{definicja}

\begin{twierdzenie}[Domknięcie rzutowe zbioru afinicznego]
\label{tw:projective-closure}
Niech $k$ będzie algebraicznie domkniętym ciałem,
$n\ge1$, $X\subseteq\A^n(k)$ zbiorem algebraicznym,
a $\iota:\A^n(k)\to\PP^n(k)$ niech będzie odwzorowaniem
$\iota(a_1,\dots,a_n)=[a_1:\dots:a_n:1]$.
Wówczas domknięcie $\iota(X)$ w topologii Zariskiego jest równe
$\V_+(\I_k(X)^h)$.
\end{twierdzenie}
\begin{proof}
Niech $Y=\overline{\iota(X)}$.
Ponieważ $\V_+(\I_k(X)^h)$ jest domknięty i każdy
$f\in\I_k(X)$ spełnia $f^h(a,1)=f(a)=0$ dla $a\in X$,
otrzymujemy $Y\subseteq\V_+(\I_k(X)^h)$.

Dla inkluzji odwrotnej weźmy dowolny zbiór domknięty
$C=\V_+(T)$ zawierający $\iota(X)$.
Każdy jednorodny $F\in T$ stopnia $d$ ma dehomogenizację
$f(x)=F(x,1)$, która znika na $X$; zatem $f\in\I_k(X)$.
Jeśli $f\ne0$ ma stopień $m\le d$, to
$F=Z^{d-m}f^h$. Gdy $f=0$, jednorodność $F$ wymusza
$F=0$: w jednorodnym $F$ wszystkie jednomiany po podstawieniu
$Z=1$ są różnymi jednomianami w pozostałych zmiennych.
Stąd każdy $F\in T$ znika na $\V_+(\I_k(X)^h)$.
Mamy więc $\V_+(\I_k(X)^h)\subseteq C$.
Po przecięciu wszystkich domkniętych $C$ zawierających
$\iota(X)$ dostajemy $\V_+(\I_k(X)^h)\subseteq Y$.
\end{proof}

\begin{uwaga}[Czy naprawdę potrzeba tu algebraicznej domkniętości?]
Powyższy dowód samej tożsamości domknięcia działa dla
każdego ciała $k$, jeśli $\I_k(X)$ oznacza rzeczywisty
ideał znikania na punktach $k$-wymiernych. Założenie
algebraicznej domkniętości jest przydatne, gdy chcemy
zastąpić $\I_k(X)$ pierwiastnikiem pierwotnego ideału równań
na mocy Nullstellensatz. Nie należy tych dwóch stwierdzeń
utożsamiać.
\end{uwaga}

\begin{przyklad}[Pułapka: homogenizacja samych generatorów]
W $k[x,y]$ rozważmy ideał
$I=(x,xy-y)=(x,y)$. Równość wynika z
$y=xy-(xy-y)\in I$. Jego zbiór zerowy to tylko punkt $(0,0)$,
którego domknięcie rzutowe powinno być punktem $[0:0:1]$.

Jeśli jednak homogenizujemy \emph{tylko} podane dwa generatory,
otrzymamy w $k[X,Y,Z]$ ideał
$J=(X,XY-YZ)$. Z jego równań wynika $X=0$ i $YZ=0$,
a więc $\V_+(J)=\{[0:0:1],[0:1:0]\}$.
Drugi punkt jest sztucznym punktem w nieskończoności.
Homogenizacja także elementu $y\in I$ daje $Y$ i usuwa
ten niepożądany punkt. Właściwy ideał $I^h$ to $(X,Y)$.
\end{przyklad}

\begin{przyklad}[Parabola i jej punkt w nieskończoności]
Równanie $y=x^2$ homogenizuje się do $YZ=X^2$.
W części $Z=1$ odzyskujemy zwykłą parabolę, a dla $Z=0$
mamy $X=0$, więc jedyny punkt w nieskończoności to $[0:1:0]$.
Dla hiperboli $xy=1$ otrzymujemy $XY=Z^2$;
na prostej $Z=0$ mamy $XY=0$, czyli dwa punkty
$[0:1:0]$ i $[1:0:0]$.
\end{przyklad}

\section{Krzywe rzutowe, proste, stożkowe i parametryzacje}

\begin{definicja}[Krzywa płaska i jej stopień]
Niech $k$ będzie algebraicznie domkniętym ciałem.
Niezerowy jednorodny $F\in k[X,Y,Z]$ dodatniego stopnia $d$
wyznacza rzutową krzywą płaską $C=\V_+(F)\subseteq\PP^2(k)$.
Jeśli $F$ jest bezkwadratowy (żaden nierozkładalny czynnik
nie występuje z wykładnikiem większym od $1$), to $d$
nazywamy \emph{stopniem} takiej krzywej z jej naturalnym
zredukowanym równaniem. Jeśli $F$ ma powtórzone czynniki,
stopień $F$ uwzględnia również te krotności.
\end{definicja}

\begin{uwaga}[Dlaczego dodajemy warunek bezkwadratowości?]
Wielomiany $X$ i $X^3$ opisują ten sam \emph{zbiór} prostej,
ale ich stopnie są różne. Stopień samego zredukowanego
zbioru algebraicznego nie może więc zależeć od arbitralnie
wybranego równania z potęgami jego czynników.
\end{uwaga}

\begin{przyklad}[Stożkowa jest obrazem prostej rzutowej]
Niech $k$ będzie ciałem, a $C=\V_+(XZ-Y^2)\subseteq\PP^2(k)$.
Odwzorowanie
$\nu:\PP^1(k)\to C$, $[s:t]\mapsto[s^2:st:t^2]$
jest dobrze określone: po zamianie $(s,t)$ na
$(\lambda s,\lambda t)$ wszystkie trzy współrzędne
mnożą się przez $\lambda^2$. Równanie krzywej zachodzi,
bo $s^2t^2-(st)^2=0$.

Jeżeli $X\ne0$, to $s\ne0$, a parametr odtwarzamy przez
$[s:t]=[1:Y/X]$. Jeżeli $Z\ne0$, analogicznie
$[s:t]=[Y/Z:1]$. Na krzywej nie ma punktu z $X=Z=0$:
wtedy $Y^2=0$ i wszystkie współrzędne byłyby zerowe.
Formuły odwrotne mają regularne współrzędne na odpowiednich
mapach, zatem $C$ jest izomorficzna z $\PP^1$.
\end{przyklad}

\begin{center}
\begin{tikzpicture}[>=Stealth,scale=.83]
 \draw[->,gray!70] (-3.1,0)--(3.45,0) node[right] {$x$};
 \draw[->,gray!70] (0,-.5)--(0,4.5) node[above] {$y$};
 \draw[blue!70!black,thick,domain=-1.9:1.9,samples=100] plot (\x,{\x*\x});
 \fill[red!80!black] (0,0) circle(2.1pt);
 \draw[densely dashed,->] (2.15,4.05)--(3,4.4);
 \node[font=\small,align=left,anchor=west,fill=white,inner sep=2pt] at (3.5,3.3)
 {Jeden punkt\\w nieskończoności\\$[0:1:0]$};
 \node[blue!70!black] at (-2.05,3.9) {$YZ=X^2$};
\end{tikzpicture}

\small Rysunek 2.2. Afiniczny fragment stożkowej $YZ=X^2$ i kierunek
jej jedynego punktu w nieskończoności.
\end{center}

\section{Odwzorowania rzutowe i ich miejsca nieokreśloności}

\begin{definicja}[Odwzorowanie zadane formami]
Niech $k$ będzie ciałem, $n,m\ge0$, a
$F_0,\dots,F_m\in k[X_0,\dots,X_n]$ formami jednorodnymi tego samego dodatniego stopnia $d$.
Na zbiorze $U=\PP^n(k)\setminus\V_+(F_0,\dots,F_m)$
definiujemy odwzorowanie
$\Phi:U\to\PP^m(k)$ wzorem
$\Phi([a])=[F_0(a):\cdots:F_m(a)]$.
Zbiór $B=\V_+(F_0,\dots,F_m)$ nazywamy
\emph{zbiorem bazowym} tego zapisu.
\end{definicja}

\begin{stwierdzenie}[Niezależność od reprezentanta]
Niech $k$ będzie ciałem, $n,m\ge0$, a
$F_0,\ldots,F_m\in k[X_0,\ldots,X_n]$
formami jednorodnymi tego samego dodatniego stopnia $d$.
Jeśli $U=\PP^n(k)\setminus\V_+(F_0,\ldots,F_m)$, to wzór
$\Phi([a])=[F_0(a):\cdots:F_m(a)]$ wyznacza dobrze określone
odwzorowanie $\Phi:U\to\PP^m(k)$.
\end{stwierdzenie}
\begin{proof}
Na $U$ wartości $F_0(a),\dots,F_m(a)$ nie są jednocześnie
zerowe, więc definiują punkt rzutowy. Zamiana $a$ na
$\lambda a$ mnoży każdą z tych wartości przez $\lambda^d$,
co nie zmienia klasy jednorodnej w $\PP^m$.
\end{proof}

\begin{przyklad}[Przekształcenie Cremony]
Niech $k$ będzie ciałem. Na $\PP^2(k)$ rozważmy
$\Phi([X:Y:Z])=[YZ:XZ:XY]$.
Zbiór bazowy składa się z trzech punktów
$[1:0:0]$, $[0:1:0]$, $[0:0:1]$, bo wszystkie trzy iloczyny
są zerowe dokładnie wtedy, gdy co najwyżej jedna współrzędna
jest niezerowa. Na otwartym zbiorze $XYZ\ne0$ mamy
\[
 (\Phi\circ\Phi)([X:Y:Z])
 =[X^2YZ:XY^2Z:XYZ^2]=[X:Y:Z].
\]
Zatem $\Phi$ jest \emph{biwymiernym} przekształceniem
płaszczyzny rzutowej, ale nie jest zdefiniowana w trzech
punktach bazowych. Wymierne przekształcenie i wszędzie
określony morfizm nie są tym samym.
\end{przyklad}

\section{Punkty gładkie, styczne i osobliwości krzywych}

\begin{definicja}[Punkt gładki krzywej płaskiej]
Niech $k$ będzie ciałem algebraicznie domkniętym,
$F\in k[X,Y,Z]$ niezerowym jednorodnym wielomianem,
a $P=[a:b:c]\in\V_+(F)$. Punkt $P$ nazywamy
\emph{gładkim dla równania $F$}, jeśli przynajmniej jedna
z liczb $F_X(a,b,c)$, $F_Y(a,b,c)$, $F_Z(a,b,c)$ jest
niezerowa. W przeciwnym razie jest \emph{osobliwy dla $F$}.
Dla zredukowanej krzywej podane kryterium jest niezależne
od przemnożenia jej równania przez niezerową stałą.
\end{definicja}

\begin{uwaga}[Zredukowane równanie ma znaczenie]
Dla $F=X^2$ wszystkie pochodne zerują się na prostej $X=0$,
choć sama prosta zredukowana jest gładka.
Dlatego w dyskusji gładkości \emph{zbioru} krzywej wybieramy
jej zredukowane równanie, a nie dowolną potęgę.
\end{uwaga}

\begin{lemat}[Tożsamość Eulera]
Dla każdego ciała $k$, liczby całkowitej $n\ge0$
i jednorodnego wielomianu
$F\in k[X_0,\dots,X_n]$ stopnia $d$ zachodzi
$\sum_{i=0}^nX_i\frac{\partial F}{\partial X_i}=dF$.
\end{lemat}
\begin{proof}
Dla jednomianu $M=cX_0^{e_0}\cdots X_n^{e_n}$
wyrażenie po lewej jest równe
$(e_0+\cdots+e_n)M=dM$.
Równość zachowuje się przy sumowaniu, więc dotyczy $F$.
\end{proof}

\begin{twierdzenie}[Styczna rzutowa w punkcie gładkim]
\label{tw:tangent-projective}
Niech $k$ będzie ciałem algebraicznie domkniętym,
$F\in k[X,Y,Z]$ niezerową jednorodną formą,
a $P=[a:b:c]\in\V_+(F)$ punktem gładkim dla $F$.
Wtedy rzutowa prosta styczna do krzywej w $P$ ma równanie
\[
 F_X(a,b,c)X+F_Y(a,b,c)Y+F_Z(a,b,c)Z=0.
\]
Równanie to jest niezależne od reprezentanta punktu $P$.
\end{twierdzenie}
\begin{proof}
Co najmniej jeden współczynnik nie znika, więc wzór
opisuje prostą rzutową. Jeśli $F$ ma stopień $d$,
to zmiana reprezentanta $(a,b,c)$ na
$(\lambda a,\lambda b,\lambda c)$ mnoży pochodne
przez $\lambda^{d-1}$. Prosta się nie zmienia.
Tożsamość Eulera i $F(a,b,c)=0$ pokazują,
że punkt $P$ leży na podanej prostej.

Aby sprawdzić styczność, wybierzmy mapę, na której jedna
współrzędna $P$ jest niezerowa, i sprowadźmy ją do $1$.
W tej mapie równanie krzywej ma postać $f(u,v)=0$.
Wyraz liniowy rozwinięcia $f$ w punkcie $p=(u_0,v_0)$
wynosi $f_u(p)(u-u_0)+f_v(p)(v-v_0)$.
Jest to równanie afinicznej stycznej w gładkim punkcie.
Po przepisaniu do współrzędnych jednorodnych uzyskujemy
powyższy wzór; brakujący wyraz stały jest dokładnie
zapewniony przez tożsamość Eulera. W charakterystyce
dzielącej $d$ tożsamość Eulera wciąż obowiązuje,
choć $d=0$ jako element ciała, a rachunek lokalny
nadal daje poprawne równanie stycznej.
\end{proof}

\begin{przyklad}[Styczna do stożkowej]
Niech $k$ będzie ciałem. Dla $F=XZ-Y^2$ otrzymujemy $F_X=Z$, $F_Y=-2Y$,
$F_Z=X$. W punkcie $P=[1:t:t^2]$ prosta styczna ma
równanie $t^2X-2tY+Z=0$. Podstawienie tego punktu daje
$t^2-2t^2+t^2=0$, jak powinno. Gdy charakterystyka
wynosi $2$, współczynnik $-2t$ jest zerem i wzór
wciąż poprawnie opisuje styczną.
\end{przyklad}

\begin{przyklad}[Ostrze i węzeł]
Dla równania afinicznego $y^2=x^3$ obie pochodne
$-3x^2$ i $2y$ znikają w $(0,0)$; zatem początek
jest punktem osobliwym. Dla równania
$y^2=x^2(x+1)$ w początku również znikają obie pochodne.
Nad $\R$ pierwsza krzywa ma ostrze, a druga lokalnie
dwie gałęzie o różnych stycznych. To różne typy osobliwości,
choć w obu przypadkach kryterium Jacobiego wykrywa
niegładkość. Bardziej systematyczna teoria będzie
w następnym rozdziale o punktach osobliwych.
\end{przyklad}

\section{Przecięcia z prostą i pierwsza wersja twierdzenia Bézouta}

\begin{definicja}[Krotność pierwiastka formy binarnej]
Niech $k$ będzie algebraicznie domkniętym ciałem,
a $H(s,t)\in k[s,t]$ niezerową formą jednorodną stopnia $d$.
Każdy punkt $P\in\PP^1(k)$ odpowiada formie liniowej
$\ell_P(s,t)$, której zbiór zerowy w $\PP^1(k)$ jest
równy $\{P\}$. \emph{Krotnością zera $H$ w $P$} nazywamy
największe $m\ge0$ takie, że $\ell_P^m$ dzieli $H$.
\end{definicja}

\begin{lemat}[Rozkład formy binarnej]
Niech $k$ będzie ciałem algebraicznie domkniętym,
a $H\in k[s,t]$ niezerową formą jednorodną stopnia $d\ge1$.
Wtedy istnieją formy liniowe $\ell_1,\ldots,\ell_d$
i $c\in k^\times$ takie, że $H=c\ell_1\cdots\ell_d$.
W szczególności suma krotności wszystkich zer $H$ w $\PP^1(k)$
wynosi $d$.
\end{lemat}
\begin{proof}
Podstawmy $s=1$. Otrzymujemy wielomian
$h(t)=H(1,t)$ stopnia $m\le d$. Jest on niezerowy:
jednorodna forma $H$ jest jednoznacznie wyznaczona przez
współczynniki jednomianów $s^{d-j}t^j$, które po
podstawieniu $s=1$ stają się współczynnikami $t^j$.
Dzięki algebraicznej domkniętości mamy
$h(t)=c\prod_{j=1}^m(t-\alpha_j)$, z uwzględnieniem krotności.
Homogenizując z powrotem do stopnia $d$, mamy
$H(s,t)=c\,s^{d-m}\prod_{j=1}^m(t-\alpha_js)$.
Współczynnik przy $s^{d-m}$ opisuje zero $[0:1]$
o krotności $d-m$. Gdy $h$ jest stały, otrzymujemy
$H=cs^d$; przypadek $d=0$ nie jest tu potrzebny.
Suma wykładników czynników liniowych jest zawsze równa $d$.
\end{proof}

\begin{twierdzenie}[Bézout dla prostej i krzywej]
\label{tw:bezout-line}
Niech $k$ będzie algebraicznie domkniętym ciałem,
$F\in k[X,Y,Z]$ niezerową formą stopnia $d\ge1$,
a $L\subseteq\PP^2(k)$ prostą rzutową niebędącą
składową krzywej $\V_+(F)$. Wtedy $L\cap\V_+(F)$
jest niepusty i skończony, a suma krotności przecięć
w jego punktach wynosi $d$. Krotność w $P\in L$
definiujemy jako krotność zera otrzymanej formy binarnej
po parametryzacji liniowej $L\cong\PP^1(k)$.
\end{twierdzenie}
\begin{proof}
Prostą $L$ możemy sparametryzować jako
$[s:t]\mapsto[s\,u+t\,v]$, gdzie $u,v\in k^3$
są liniowo niezależne i rozpinają dwuwymiarową
podprzestrzeń odpowiadającą $L$. Po podstawieniu
otrzymujemy jednorodną formę $H(s,t)=F(su+tv)$ stopnia $d$.
Nie jest ona zerowa, ponieważ w przeciwnym razie $F$
znikałby na całej prostej $L$, a jej równanie liniowe
dzieliłoby $F$ (po liniowej zmianie współrzędnych
byłaby to podzielność przez jedną zmienną). To oznaczałoby,
że $L$ jest składową $\V_+(F)$, wbrew założeniu.
Na mocy poprzedniego lematu suma krotności zer $H$
wynosi $d$. Punkty zerowe $H$ odpowiadają dokładnie
punktom przecięcia $L$ i krzywej.
\end{proof}

\begin{przyklad}[Dlaczego styczność liczy się podwójnie?]
Dla stożkowej $XZ-Y^2=0$ i prostej $Z=X$ podstawiamy
$[X:Y:Z]=[s:t:s]$. Otrzymujemy $s^2-t^2$.
Jeśli $\operatorname{char}k\ne2$, są to dwa przecięcia
$[1:1:1]$ i $[1:-1:1]$, każde krotności $1$.
Natomiast prosta styczna w punkcie $[1:0:0]$ ma równanie
$Z=0$. Jej ograniczenie daje $-Y^2$, więc jedyny punkt
przecięcia $[1:0:0]$ ma krotność $2$.
\end{przyklad}

\section{Ogólne twierdzenie Bézouta: po co nam algebra ilorazowa?}

Dla dwóch krzywych stopni $d$ i $e$ oczekiwalibyśmy
$d\cdot e$ punktów przecięcia. Bez krotności jest to
fałszywe nawet dla prostej stycznej do stożkowej.
Potrzebujemy więc definicji uwzględniającej zarówno
zwykłe przecięcia, jak i styczność oraz punkty osobliwe.
Poniżej dowodzimy pełnej wersji płaskiego twierdzenia Bézouta
bez odwoływania się do nieudowodnionego wzoru na wymiar krzywej.
Przyda się jednak kilka lematów algebraicznych.

\begin{definicja}[Pierścień stopniowany i szereg Hilberta]
Niech $k$ będzie ciałem. Pierścień $S=k[X,Y,Z]$ ma \emph{stopniowanie}
$S=\bigoplus_{m\ge0}S_m$, gdzie $S_m$ składa się z form
jednorodnych stopnia $m$. Dla jednorodnego ideału $J$
iloraz $S/J$ dziedziczy stopniowanie.
Dla stopniowanej $k$-przestrzeni $M=\bigoplus_{m\ge0}M_m$
o skończenie wymiarowych składowych \emph{szeregiem Hilberta}
nazywamy formalny szereg
$H_M(t)=\sum_{m\ge0}(\dim_kM_m)t^m$.
\end{definicja}

\begin{stwierdzenie}[Szereg Hilberta przekroju dwóch form]
\label{st:hilbert-complete-intersection}
Niech $k$ będzie ciałem i niech $F,G\in S=k[X,Y,Z]$
będą niezerowymi formami odpowiednio stopni $d,e\ge1$,
bez wspólnego nierozkładalnego czynnika. Wtedy dla
$M=S/(F,G)$ zachodzi tożsamość szeregów formalnych
\[
 H_M(t)=\frac{(1-t^d)(1-t^e)}{(1-t)^3}.
\]
\end{stwierdzenie}
\begin{proof}
Ponieważ monomy $X^aY^bZ^c$, $a+b+c=m$, stanowią bazę
$S_m$, mamy $\dim_kS_m=\binom{m+2}{2}$, a więc
$H_S(t)=(1-t)^{-3}$ (to wynika z trzykrotnego
pomnożenia szeregu geometrycznego $\sum_{m\ge0}t^m$).
Mnożenie przez niezerowe $F$ jest różnowartościowe,
bo $S$ jest dziedziną. Z dokładnego ciągu przestrzeni
stopniowanych $0\to S(-d)\xrightarrow{\cdot F}S\to S/(F)\to0$
otrzymujemy $H_{S/(F)}(t)=(1-t^d)H_S(t)$.

Mnożenie przez $G$ na $S/(F)$ jest także różnowartościowe:
jeśli $Gh\in(F)$, to $F\mid Gh$. Wielomiany w $S$
mają jednoznaczny rozkład na czynniki nierozkładalne
(twierdzenie~\ref{tw:polynomial-UFD}),
a $F$ i $G$ nie mają wspólnego czynnika. Każdy czynnik
$F$ wraz ze swoją krotnością musi zatem dzielić $h$;
czyli $F\mid h$. Wobec tego $[h]=0$ w $S/(F)$.
Drugi dokładny ciąg,
$0\to(S/(F))(-e)\xrightarrow{\cdot G}S/(F)\to M\to0$,
daje $H_M(t)=(1-t^e)H_{S/(F)}(t)$, co kończy dowód.
\end{proof}

\begin{uwaga}[Dokładność i przesunięcie stopnia]
Przesunięcie $N(-d)$ oznacza ten sam moduł z nowym
stopniowaniem $(N(-d))_m=N_{m-d}$ (i $N_j=0$ dla $j<0$).
Ciąg $0\to A\to B\to C\to0$ jest \emph{dokładny},
gdy obraz każdego odwzorowania jest jądrem następnego.
Dla każdej składowej stopnia $m$ mamy wówczas
$\dim B_m=\dim A_m+\dim C_m$, stąd dodawanie
szeregów Hilberta. Nie używamy tu jeszcze teorii
homologii ani funktorów pochodnych.
\end{uwaga}

\begin{lemat}[Skończoność przecięcia i wybór prostej w nieskończoności]
\label{lem:finiteness-plane-intersection}
Niech $k$ będzie ciałem algebraicznie domkniętym, a niezerowe
formy $F,G\in k[X,Y,Z]$ nie mają wspólnego czynnika.
Wtedy $\V_+(F,G)\subseteq\PP^2(k)$ jest zbiorem skończonym.
Istnieje więc prosta rzutowa $L$ nieprzechodząca przez
żaden z tych punktów.
\end{lemat}
\begin{proof}
Wybierzmy punkt $Q\in\PP^2(k)$, który nie należy do
$\V_+(F)\cup\V_+(G)$. Taki punkt istnieje:
oba wielomiany są niezerowe, a nad nieskończonym ciałem
iloczyn $FG$ nie może znikać na całym $k^3$.
Po liniowej zmianie współrzędnych możemy przyjąć
$Q=[0:1:0]$, a wtedy współczynniki przy $Y^d$ w $F$
i przy $Y^e$ w $G$ są niezerowymi stałymi.
Postrzegamy $F,G$ jako wielomiany w $Y$ o współczynnikach
w $k[X,Z]$ i przechodzimy do ciała $K=k(X,Z)$.

Nie mają wspólnego czynnika w $K[Y]$. Rzeczywiście,
każdy z nich jest \emph{pierwotny} w $k[X,Z][Y]$,
ponieważ jego współczynnik wiodący jest niezerową stałą.
Gdyby miały wspólny niejednostkowy czynnik nad $K$,
po usunięciu mianowników i zastosowaniu lematu Gaussa
dla wielomianów nad pierścieniem z jednoznacznym rozkładem
otrzymalibyśmy wspólny czynnik w $k[X,Y,Z]$.

Dla dwóch wielomianów $A,B\in K[Y]$ stopni $d,e$
definiujemy \emph{rezultant} $\operatorname{Res}_Y(A,B)$
jako wyznacznik macierzy liniowego odwzorowania
\[
 K[Y]_{<e}\oplus K[Y]_{<d}\longrightarrow K[Y]_{<d+e},
 \qquad (U,V)\longmapsto UA+VB.
\]
Tutaj $K[Y]_{<r}$ oznacza wielomiany stopnia mniejszego niż $r$.
Jeśli wyznacznik jest zerowy, istnieje niezerowa para
$(U,V)$ z $UA+VB=0$. Dla wielomianów względnie pierwszych
wynikałoby $B\mid U$, co jest niemożliwe przy $\deg U<e$,
chyba że $U=0$, a wtedy także $V=0$; sprzeczność.
Odwrotnie, przy wspólnym czynniku $D$ można wybrać
$U=B/D$ i $V=-A/D$, więc wyznacznik jest zerowy.
Zatem rezultat $R(X,Z)=\operatorname{Res}_Y(F,G)$
jest \emph{niezerowym} wielomianem w $k[X,Z]$.

Jeśli $P=[x:y:z]$ jest wspólnym zerem $F,G$, to
$[x:z]\ne[0:0]$, gdyż inaczej $P=Q$, a $F(Q)G(Q)\ne0$.
Podstawienie $(X,Z)=(x,z)$ do macierzy rezultantu
daje dwa wielomiany w $Y$ o wspólnym pierwiastku $y$,
więc $R(x,z)=0$.
Wielomian $R$ jest jednorodny stopnia $de$.
Aby uzasadnić stopień, zapiszmy
$F=\sum_{j=0}^d f_j(X,Z)Y^j$ i
$G=\sum_{j=0}^e g_j(X,Z)Y^j$; współczynniki $f_j$
i $g_j$ są jednorodne stopni odpowiednio $d-j$ i $e-j$.
W macierzy rezultantu jest $e$ wierszy odpowiadających
$Y^iF$, $0\le i<e$, i $d$ wierszy odpowiadających
$Y^iG$, $0\le i<d$. W każdym składniku wyznacznika
wybiera się po jednym współczynniku z każdego wiersza
i z każdej kolumny potęg $Y^0,\dots,Y^{d+e-1}$.
Suma wykładników $Y$ po odjęciu przesunięć wierszy
wynosi zatem
$\sum_{j=0}^{d+e-1}j-\sum_{i=0}^{e-1}i-\sum_{i=0}^{d-1}i=de$.
Łączny stopień wybranych współczynników w $X,Z$
wynosi więc $ed+de-de=de$, niezależnie od składnika
wyznacznika.
Niezerowa forma binarna ma tylko skończenie wiele zer na
$\PP^1(k)$, jak wykazaliśmy wcześniej. Dla każdego z tych
skończenie wielu $[x:z]$ wielomian $F(x,Y,z)$ jest
niezerowy (ma niezerowy współczynnik przy $Y^d$),
więc ma skończenie wiele pierwiastków $y$.
W rezultacie $\V_+(F,G)$ jest skończony.

Ponieważ $k$ jest nieskończone, spośród prostych rzutowych
możemy wybrać prostą omijającą wszystkie te punkty:
każdy zadany punkt wyznacza tylko właściwy podzbiór
przestrzeni współczynników możliwych równań liniowych,
a skończona suma takich właściwych zbiorów nie wyczerpuje
$k^3$. Po zmianie współrzędnych prostą tę możemy przyjąć
jako $Z=0$.
\end{proof}

\begin{lemat}[Jednorodna algebra przecięcia ma rangę $de$]
\label{lem:bezout-rank}
Niech $k$ będzie ciałem algebraicznie domkniętym,
a $F,G\in k[X,Y,Z]$ będą formami stopni $d,e\ge1$
bez wspólnego czynnika. Załóżmy, że na prostej $Z=0$
nie ma ich wspólnego zera rzutowego.
Wtedy $A=k[x,y]/(f,g)$, gdzie
$f(x,y)=F(x,y,1)$, $g(x,y)=G(x,y,1)$,
jest skończenie wymiarową $k$-algebrą i
$\dim_kA=de$.
\end{lemat}
\begin{proof}
Niech $S=k[X,Y,Z]$, $M=S/(F,G)$ oraz
$\bar F=F(X,Y,0)$, $\bar G=G(X,Y,0)$.
Formy $\bar F,\bar G$ nie są zerowe: gdyby np.
$\bar F=0$, to $Z\mid F$. Ponieważ $F,G$ nie mają
wspólnego czynnika, koniecznie $Z\nmid G$,
czyli $\bar G\ne0$. Niezerowa forma $\bar G$
dodatniego stopnia ma nad ciałem algebraicznie
domkniętym zero na prostej $Z=0$, sprzecznie
z założeniem. Podobnie dla $\bar G=0$. Formy
$\bar F,\bar G$ są względnie pierwsze w $k[X,Y]$:
inaczej wspólny czynnik liniowy dawałby wspólny punkt
na prostej w nieskończoności.

Pokażemy, że mnożenie przez $Z$ jest różnowartościowe
w $M$. Jeśli $ZH=UF+VG$, to po podstawieniu $Z=0$
mamy $\bar U\bar F+\bar V\bar G=0$.
Ponieważ $\gcd(\bar F,\bar G)=1$, wynika stąd
$\bar G\mid\bar U$ i $\bar F\mid\bar V$.
Istnieje zatem $\bar W\in k[X,Y]$ z
$\bar U=\bar W\bar G$ i $\bar V=-\bar W\bar F$.
Wybierając dowolne podniesienie $W\in S$, możemy zapisać
$U=WG+ZU_1$, $V=-WF+ZV_1$. Otrzymujemy
$ZH=Z(U_1F+V_1G)$, a po skróceniu $Z$
w dziedzinie $S$ mamy $H\in(F,G)$.

Zatem mamy dokładny ciąg stopniowany
$0\to M(-1)\xrightarrow{\cdot Z}M\to M/ZM\to0$.
Szereg Hilberta z poprzedniego stwierdzenia daje
\[
 H_{M/ZM}(t)=(1-t)H_M(t)
 =\frac{(1-t^d)(1-t^e)}{(1-t)^2}
 =(1+t+\cdots+t^{d-1})(1+t+\cdots+t^{e-1}).
\]
Ponieważ $M/ZM\cong k[X,Y]/(\bar F,\bar G)$,
jest to skończenie wymiarowa $k$-przestrzeń, a jej
wymiar jest sumą współczynników podanego wielomianu,
czyli $de$.

Wyjaśnijmy, dlaczego stąd wynika $\dim_kM/(Z-1)M=de$.
Wybierzmy jednorodne elementy $b_1,\ldots,b_{de}\in M$
reprezentujące bazę stopniowanego ilorazu $M/ZM$.
Pokażemy, że są bazą $M$ jako modułu nad $k[Z]$.
Dla dowolnego elementu jednorodnego $h\in M$ jego klasa
modulo $ZM$ jest kombinacją klas tych $b_i$, których
stopień jest równy $\deg h$. Po odjęciu tej kombinacji
pozostaje $Zh_1$ z $\deg h_1=\deg h-1$.
Indukcja względem stopnia pokazuje, że $b_i$ generują $M$
nad $k[Z]$. Mamy więc surjekcję stopniowanych modułów
$\bigoplus_i k[Z](-\deg b_i)\twoheadrightarrow M$.
Szereg Hilberta źródła wynosi
$\sum_i t^{\deg b_i}/(1-t)=H_{M/ZM}(t)/(1-t)$,
czyli $H_M(t)$ dzięki dokładnemu ciągowi powyżej.
W każdej składowej stopnia mapa jest więc surjekcją
między przestrzeniami tego samego skończonego wymiaru;
jest izomorfizmem. Dlatego $M\cong\bigoplus_i k[Z]b_i$.

Iloraz przez $Z-1$ zamienia każdy składnik $k[Z]b_i$
w jednowymiarową przestrzeń nad $k$, a ponadto
$M/(Z-1)M\cong k[x,y]/(f,g)=A$, ponieważ podstawienie
$Z=1$ dehomogenizuje $F$ i $G$. Stąd $\dim_kA=de$.
\end{proof}

\begin{definicja}[Lokalna krotność przecięcia]
\label{def:intersection-multiplicity}
Niech $k$ będzie ciałem algebraicznie domkniętym,
a $f,g\in k[x,y]$ nie mają wspólnej składowej algebraicznej
i mają skończenie wiele wspólnych zer.
Załóżmy, że algebra $A=k[x,y]/(f,g)$ jest skończenie wymiarowa
nad $k$. Dla wspólnego zera $p=(a,b)$ oznaczmy przez
$\mathfrak m_p=(x-a,y-b)A$ ideał maksymalny $A$.
Przez $A_{\mathfrak m_p}$ oznaczamy lokalizację $A$
względem zbioru $A\setminus\mathfrak m_p$:
jej elementy to ułamki $u/s$ z $s\notin\mathfrak m_p$,
z relacją zwykłego mnożenia na krzyż po ewentualnym
pomnożeniu przez dodatkowy mianownik.
Definiujemy \emph{lokalną krotność przecięcia}
$I_p(f,g):=\dim_k A_{\mathfrak m_p}$.
\end{definicja}

\begin{lemat}[Algebra skończona rozkłada się na części lokalne]
\label{lem:artin-decomposition}
Niech $k$ będzie algebraicznie domkniętym ciałem,
a $A$ niezerową, skończenie wymiarową przemienną $k$-algebrą.
Wtedy $A$ ma skończenie wiele ideałów maksymalnych
$\mathfrak m_1,\ldots,\mathfrak m_r$ i
$\dim_kA=\sum_{i=1}^r\dim_k A_{\mathfrak m_i}$.
\end{lemat}
\begin{proof}
Ponieważ $A$ jest skończenie wymiarowa nad $k$, jest
pierścieniem noetherowskim. Jest też artinowski:
każdy malejący ciąg ideałów stabilizuje się, ponieważ
ideały są podprzestrzeniami $k$-liniowymi o nierosnących
wymiarach. Każdy jego iloraz będący dziedziną jest ciałem:
dla $0\ne a$ w skończenie wymiarowej dziedzinie
mnożenie przez $a$ jest różnowartościowym endomorfizmem
przestrzeni wektorowej, więc jest też surjektywne
i istnieje element odwrotny do $a$.
Zatem każdy ideał pierwszy $A$ jest maksymalny.

Każdy ideał maksymalny $\mathfrak m$ ma iloraz $A/\mathfrak m$
będący skończenie wymiarowym rozszerzeniem ciała $k$.
Algebraiczna domkniętość i lemat~\ref{lem:extension-trivial}
dają $A/\mathfrak m\cong k$.
Różne ideały maksymalne są parami współmaksymalne.
Dla dowolnych $r$ różnych ideałów maksymalnych chińskie
twierdzenie o resztach daje surjekcję $A\to k^r$,
a więc $r\le\dim_k A$. Liczba takich ideałów jest
zatem skończona.

Niech $N$ będzie nilradykałem $A$, czyli ideałem wszystkich
elementów nilpotentnych. Standardową równość
$N=\bigcap_i\mathfrak m_i$ uzasadnimy następująco:
każdy nilpotent leży w każdym ideale pierwszym;
odwrotnie, jeśli $a$ nie jest nilpotentny, zbiór
$\{1,a,a^2,\dots\}$ nie zawiera zera, a więc istnieje
ideał maksymalny w rodzinie ideałów rozłącznych z tym
zbiorem (lemat o maksymalności można otrzymać z
noetherowskości przez argument rosnącego łańcucha).
Taki maksymalny ideał rozłączny $J$ jest pierwszy.
Istotnie, jeśli $uv\in J$, lecz $u,v\notin J$, to oba
większe ideały $(J,u)$ i $(J,v)$ przecinają zbiór
$\{1,a,a^2,\dots\}$. Istnieją zatem liczby $r,s\ge0$
i elementy $j_1,j_2\in J$ oraz $b,c\in A$ takie, że
$a^r=j_1+bu$ i $a^s=j_2+cv$.
Po wymnożeniu otrzymujemy
$a^{r+s}=j_1j_2+j_1cv+j_2bu+bcuv\in J$,
co przeczy rozłączności $J$ ze zbiorem potęg $a$.
Ponieważ w $A$ wszystkie ideały pierwsze są maksymalne,
otrzymujemy wymaganą równość.

Ideał $N$ jest nilpotentny. Jeśli $N=(0)$, wystarcza
$N^1=(0)$. W przeciwnym przypadku jest skończenie
generowany przez niepusty układ $n_1,\dots,n_s$,
a każdy $n_i$ ma potęgę $n_i^{q_i}=0$.
W każdym iloczynie $\sum_i(q_i-1)+1$ generatorów
któryś $n_i$ wystąpi co najmniej $q_i$ razy,
więc iloczyn jest zerowy.
Stąd $N^M=0$ dla odpowiednio dużego $M$.
Ponieważ ideały $\mathfrak m_i$ są parami współmaksymalne,
$N=\bigcap_i\mathfrak m_i=\prod_i\mathfrak m_i$.
Wynika stąd $\prod_i\mathfrak m_i^M=0$.
Ideały $\mathfrak m_i^M$ też są parami współmaksymalne,
więc chińskie twierdzenie o resztach daje
$A\cong\prod_i A/\mathfrak m_i^M$.

W ilorazie $B_i=A/\mathfrak m_i^M$ jedynym ideałem maksymalnym
jest obraz $\mathfrak m_i$: każdy inny ideał maksymalny
$\mathfrak m_j$ nie zawiera $\mathfrak m_i^M$.
Ponadto każdy element spoza tego ideału jest odwracalny:
pozostaje poza wszystkimi ideałami maksymalnymi, a element
nieodwracalny zawsze należy do pewnego ideału maksymalnego.
Zatem $B_i$ jest pierścieniem lokalnym.
Lokalizacja iloczynu $\prod_j B_j$ względem dopełnienia
i-tego ideału maksymalnego usuwa wszystkie pozostałe czynniki
(idempotenty w nich stają się zerami), a zachowuje $B_i$.
Dlatego $A_{\mathfrak m_i}\cong B_i$. Biorąc wymiary
wektorowe rozkładu produktu, otrzymujemy tezę.
\end{proof}

\begin{twierdzenie}[Bézout dla dwóch krzywych płaskich]
\label{tw:bezout-general}
Niech $k$ będzie ciałem algebraicznie domkniętym,
a $F,G\in k[X,Y,Z]$ niezerowymi wielomianami jednorodnymi
stopni odpowiednio $d,e\ge1$, bez wspólnego
nierozkładalnego czynnika. Wtedy wspólnych punktów
krzywych $\V_+(F)$ i $\V_+(G)$ jest skończenie wiele,
a suma ich lokalnych krotności przecięcia wynosi $de$:
\[
 \sum_{P\in\V_+(F,G)} I_P(F,G)=de.
\]
Lokalną krotność definiujemy w dowolnej mapie afinicznej
zawierającej $P$ przez lokalizację ilorazu odpowiednich
dehomogenizacji, jak w definicji~\ref{def:intersection-multiplicity}.
\end{twierdzenie}
\begin{proof}
Z lematu~\ref{lem:finiteness-plane-intersection} zbiór
wspólnych punktów jest skończony i można wybrać prostą
w nieskończoności, która go omija. Po rzutowej liniowej
zmianie współrzędnych możemy przyjąć, że tą prostą
jest $Z=0$. Każdy wspólny punkt leży więc w mapie
$Z\ne0$ i odpowiada wspólnemu zeru wielomianów
$f(x,y)=F(x,y,1)$ i $g(x,y)=G(x,y,1)$.

Na mocy lematu~\ref{lem:bezout-rank} algebra
$A=k[x,y]/(f,g)$ ma wymiar $de$ nad $k$.
Jej ideały maksymalne odpowiadają dokładnie wspólnym
punktom $p=(a,b)$: wynika to z opisu maksymalnych ideałów
w $k[x,y]$ nad ciałem algebraicznie domkniętym
(wniosek~\ref{wn:maximal-point}) oraz z twierdzenia
o odpowiedniości ideałów w ilorazie. Wreszcie lemat
\ref{lem:artin-decomposition} daje
$de=\dim_k A=\sum_{p}\dim_k A_{\mathfrak m_p}
=\sum_p I_p(f,g)$. Przy zmianie mapy rzutowej
lokalne pierścienie ilorazowe są izomorficzne.
Istotnie, w przecięciu map $X_i\ne0$ i $X_j\ne0$
zmiana współrzędnych jest regularna i ma regularną odwrotność.
Jeśli $F$ ma stopień $d$, jego dehomogenizacje na tych
mapach spełniają $F/X_j^d=(X_i/X_j)^d(F/X_i^d)$;
analogicznie dla $G$ stopnia $e$. Czynniki $X_i/X_j$
i $X_j/X_i$ są na przecięciu odwracalne, więc
oba układy równań generują ten sam ideał po
przeniesieniu do wspólnego pierścienia funkcji
regularnych. Lokalizacja w punkcie $P$ daje zatem
izomorficzne lokalne pierścienie ilorazowe.
Ich wymiar nad $k$, a więc lokalna krotność przecięcia,
nie zależy od wybranej mapy. To dowodzi tezy.
\end{proof}

\begin{przyklad}[Dwa przecięcia, łącznie stopień cztery]
Niech $k$ będzie ciałem algebraicznie domkniętym
charakterystyki różnej od $2$. Dla $F=XZ-Y^2$
i $G=XZ+Y^2$ oba wielomiany mają
stopień $2$. Z równań $XZ-Y^2=XZ+Y^2=0$
otrzymujemy $Y^2=0$ oraz $XZ=0$, czyli punkty
$P=[1:0:0]$ i $Q=[0:0:1]$. W mapie $X=1$ wokół $P$
równania to $z-y^2=0$ i $z+y^2=0$.
W pierścieniu lokalnym iloraz jest izomorficzny
$k[y]_{(y)}/(y^2)$ (ponieważ $2$ jest odwracalne),
o wymiarze $2$ nad $k$. Analogicznie w mapie $Z=1$
wokół $Q$ krotność wynosi $2$. Suma $2+2=4$ zgadza się
z $\deg F\deg G=2\cdot2$.
\end{przyklad}

\section{Co należy zapamiętać przed teorią osobliwości i krzywych?}

Punkty przestrzeni rzutowej to kierunki w przestrzeni o jeden
wymiar większej. Równania geometryczne są jednorodne,
ponieważ tylko wtedy ich zbiory zerowe nie zależą od
reprezentanta punktu. Mapy $U_i\cong\A^n$ pozwalają
prowadzić zwykłe rachunki w lokalnych współrzędnych.
Przejście od zbioru afinicznego do jego domknięcia rzutowego
wymaga homogenizacji \emph{całego ideału znikania}.

Przecięcia rzutowe nie gubią punktów w nieskończoności,
ale ich liczba musi być liczona z krotnościami.
Dla prostej i krzywej stopnia $d$ dostajemy $d$,
a dla dwóch krzywych bez wspólnej składowej stopni $d,e$
dostajemy $de$. Ta sama algebra, która wcześniej służyła
do badania ideałów, pozwala teraz rachować przecięcia.
Kolejny naturalny temat to punkty osobliwe, stożki styczne,
normalizacja krzywych i ich rodzaj topologiczny.

% ======================================================================
% ROZDZIAŁ 3: PUNKTY OSOBLIWE I GEOMETRIA LOKALNA KRZYWYCH
% ======================================================================

\chapter{Punkty osobliwe i geometria lokalna krzywych algebraicznych}
\label{ch:osobliwosci}

\section{Dlaczego obraz krzywej może być mylący?}

Prosta, parabola i okrąg mają w zwykłych punktach jeden wyraźny
kierunek styczny. Natomiast krzywe $y^2=x^3$ oraz
$y^2=x^2(x+1)$ mają w początku układu współrzędnych
zachowanie zupełnie odmienne od prostej: pierwsza ma
\emph{ostrze}, a druga \emph{węzeł} z dwiema gałęziami.
Sama informacja, że jest to zbiór zerowy jednego wielomianu,
nie mówi jeszcze, jak rozróżnić takie punkty.

Celem tego rozdziału jest zrozumienie trzech poziomów opisu:
\emph{przestrzeni stycznej}, \emph{stożka stycznego} oraz
\emph{lokalnego pierścienia funkcji}. Zobaczymy, że przestrzeń
styczna może mieć zbyt duży wymiar, stożek styczny rozróżnia
kierunki dochodzenia, a pierścień lokalny zachowuje także
ich krotności i inne informacje algebraiczne.

W całym rozdziale $k$ oznacza ciało algebraicznie domknięte;
gdy odwołujemy się do lokalnej intuicji analitycznej,
przyjmujemy $k=\C$. W przykładzie węzła i w obliczeniach
z dwiema odmiennymi stycznymi dodatkowo zakładamy
$\operatorname{char}k\ne2$.

\begin{center}
\begin{tikzpicture}[>=Stealth,scale=.83]
  \begin{scope}[shift={(-4,0)}]
   \draw[->,gray!70] (-.4,0)--(2.7,0) node[right] {$x$};
   \draw[->,gray!70] (0,-1.65)--(0,1.85) node[above] {$y$};
   \draw[blue!75!black,very thick,domain=-1.25:1.25,samples=120]
    plot ({\x*\x},{\x*\x*\x});
   \fill (0,0) circle(2pt); \node at (1.5,-1.55) {ostrze $y^2=x^3$};
  \end{scope}
  \begin{scope}[shift={(2.3,0)}]
   \draw[->,gray!70] (-1.7,0)--(2.8,0) node[right] {$x$};
   \draw[->,gray!70] (0,-1.7)--(0,1.9) node[above] {$y$};
   \draw[red!75!black,very thick,domain=-1.7:1.7,samples=145]
    plot ({\x*\x-1},{\x*(\x*\x-1)});
   \fill (0,0) circle(2pt); \node at (.7,-2.05) {węzeł $y^2=x^2(x+1)$};
  \end{scope}
\end{tikzpicture}

\small Rysunek 3.1. Dwie krzywe o osobliwym początku, lecz o różnej
liczbie lokalnych gałęzi. Na prawym rysunku parametry $t=1$ i $t=-1$
dają ten sam punkt $(0,0)$.
\end{center}

\section{Pierścień lokalny: funkcje w pobliżu punktu}

\begin{definicja}[Pierścień lokalny w punkcie]
\label{def:local-ring-point}
Niech $X\subseteq k^n$ będzie niepustym zbiorem algebraicznym,
$A=k[X]$ jego pierścieniem współrzędnych, a $p\in X$.
Przez $\mathfrak m_p\subseteq A$ oznaczamy ideał klas
funkcji wielomianowych znikających w $p$.
\emph{Pierścieniem lokalnym} $X$ w punkcie $p$ nazywamy
\[
 \mathcal O_{X,p}=A_{\mathfrak m_p}
 =\left\{\frac fg:\ f,g\in A,\ g(p)\ne0\right\}/\!\sim,
\]
gdzie $f/g=f'/g'$ wtedy i tylko wtedy, gdy dla pewnego
$h\in A$ z $h(p)\ne0$ mamy $h(fg'-f'g)=0$ w $A$.
\end{definicja}

Warunek dodatkowego czynnika $h$ w relacji równości jest
istotny w pierścieniach z dzielnikami zera. Jeśli $A$
jest dziedziną, wystarczy zwykła równość $fg'=f'g$.
Funkcja $f/g$ ma sens na pewnym otoczeniu punktu $p$,
ponieważ $g(p)\ne0$, a zbiór $D(g)$ jest otwarty
w topologii Zariskiego.

\begin{twierdzenie}[Jedyny ideał maksymalny]
\label{tw:local-unique-max}
Dla każdego algebraicznego $X\subseteq k^n$ i $p\in X$
pierścień $\mathcal O_{X,p}$ ma dokładnie jeden ideał
maksymalny, równy
$\mathfrak m_p\mathcal O_{X,p}=\{f/g:f(p)=0,\ g(p)\ne0\}$.
Element $f/g$ jest odwracalny dokładnie wtedy, gdy $f(p)\ne0$.
\end{twierdzenie}
\begin{proof}
Podstawienie w $p$ jest dobrze określonym homomorfizmem
$\operatorname{ev}_p:\mathcal O_{X,p}\to k$,
$f/g\mapsto f(p)/g(p)$: jeśli ułamki są równe,
po podstawieniu w $p$ ich iloczyny na krzyż są równe.
Homomorfizm jest surjektywny, bo przyjmuje każdą stałą.
Jego jądro to dokładnie ułamki o liczniku znikającym w $p$,
a więc jest to ideał maksymalny.

Jeżeli $f(p)\ne0$, to $g/f$ jest poprawnym ułamkiem
lokalnym i jest odwrotnością $f/g$. Jeśli $f(p)=0$,
to $\operatorname{ev}_p(f/g)=0$, a więc element ten nie
może być odwracalny: homomorfizm zachowujący jedynkę
przenosi jednostki na jednostki.
W każdym pierścieniu ideał maksymalny nie zawiera jednostek.
Dowolny ideał właściwy $J$ może więc składać się tylko
z niejednostek i musi zawierać się w jądrze
$\operatorname{ev}_p$. Jedyny ideał maksymalny jest zatem
właśnie tym jądrem.
\end{proof}

\begin{przyklad}[Regularna funkcja na otoczeniu, ale nie na całej krzywej]
Na hiperboli $H=\V_k(xy-1)$ mamy
$k[H]\cong k[x,x^{-1}]$; element $1/x$ jest nawet
globalną funkcją regularną, bo w pierścieniu współrzędnych
$1/x=y$. Natomiast na prostej $\A^1(k)$ funkcja $1/x$
nie należy do $k[x]$, lecz należy do lokalnego pierścienia
w każdym punkcie $p\ne0$. Nie należy do
$\mathcal O_{\A^1,0}$, ponieważ mianownik $x$ znika w $0$.
\end{przyklad}

\begin{przyklad}[Lokalne równania węzła]
Dla $X=\V_k(y^2-x^2(x+1))$ i $p=(0,0)$
mamy $\mathcal O_{X,p}=(k[x,y]/(y^2-x^2(x+1)))_{(x,y)}$.
Element $1+x$ jest jednostką tego pierścienia, ponieważ
$(1+x)(p)=1$. Nie znaczy to, że $1+x$ jest jednostką
całego pierścienia $k[X]$.
\end{przyklad}

\section{Przestrzeń kostyczna i przestrzeń styczna Zariskiego}

\begin{definicja}[Przestrzeń kostyczna i styczna]
\label{def:tangent-zariski}
Niech $X\subseteq k^n$ będzie zbiorem algebraicznym,
$p\in X$, a $\mathfrak m$ jedynym ideałem maksymalnym
$\mathcal O_{X,p}$. Iloraz $\mathfrak m/\mathfrak m^2$
traktujemy jako przestrzeń wektorową nad
$\mathcal O_{X,p}/\mathfrak m\cong k$; nazywa się ją
\emph{przestrzenią kostyczną} w $p$.
\emph{Przestrzeń styczna Zariskiego} to
\[
 T_pX:=\operatorname{Hom}_k(\mathfrak m/\mathfrak m^2,k).
\]
\end{definicja}

Dlaczego modulo $\mathfrak m^2$? Iloczyny dwóch funkcji
znikających w $p$ są \emph{drugiego rzędu małości}.
W zwykłym rachunku różniczkowym pochodna iloczynu
takich funkcji w $p$ jest zerowa. Gdy więc zostawimy
jedynie klasy modulo $\mathfrak m^2$, zachowamy
informację o wyrazach pierwszego rzędu.

\begin{twierdzenie}[Jawny wzór na przestrzeń styczną]
\label{tw:zariski-tangent-formula}
Niech $k$ będzie ciałem, $X=\V_k(I)\subseteq k^n$
zbiorem algebraicznym, $p=(a_1,\dots,a_n)\in X$,
a $I(X)$ jego ideałem znikania. Wtedy istnieje
naturalna identyfikacja
\[
 T_pX\cong\left\{v\in k^n:
 \sum_{j=1}^n\frac{\partial f}{\partial x_j}(p)v_j=0
 \quad\text{dla każdego }f\in I(X)\right\}.
\]
Jeżeli $I(X)=(f_1,\dots,f_r)$, to wystarczy sprawdzać
te równania dla generatorów $f_i$.
\end{twierdzenie}
\begin{proof}
Przesuńmy współrzędne tak, aby $p$ stał się początkiem.
Niech $R=k[x_1,\dots,x_n]$ i
$\mathfrak n=(x_1-a_1,\dots,x_n-a_n)$.
Każdy wielomian z $\mathfrak n$ można zapisać jako
sumę wyrazów liniowych w $x_j-a_j$ i wyrazów należących
do $\mathfrak n^2$. W rezultacie
$\mathfrak n/\mathfrak n^2$ ma bazę klas
$x_1-a_1,\dots,x_n-a_n$ i jest izomorficzny z $k^n$.

W pierścieniu $A=k[X]=R/I(X)$ ideał punktu jest obrazem
$\mathfrak n$. Lokalizacja w $p$ nie zmienia ilorazu
$\mathfrak m/\mathfrak m^2$: każdy mianownik nieznikający
w $p$ ma w tym ilorazie tę samą wartość co jego stała
wartość w $p$, ponieważ różnica należy do $\mathfrak m$,
a jej iloczyn przez element $\mathfrak m$ należy do
$\mathfrak m^2$. Zatem przestrzeń kostyczna na $X$
jest ilorazem $\mathfrak n/\mathfrak n^2$ przez obrazy
wielomianów $f\in I(X)$, których wyrazy stałe są równe zero.

Wyraz liniowy rozwinięcia $f$ w $p$ ma postać
$\sum_j(\partial f/\partial x_j)(p)(x_j-a_j)$.
Przestrzeń dualna ilorazu jest więc anihilatorem
wszystkich tych form liniowych w $(k^n)^*$.
To dokładnie podany zbiór wektorów $v$.
Jeśli $f=\sum_i h_if_i$ oraz $f_i(p)=0$, to
reguła Leibniza daje
$df(p)=\sum_i h_i(p)\,df_i(p)$, więc wystarcza
sprawdzenie generatorów ideału.
\end{proof}

\begin{uwaga}[Dlaczego nie wystarczą generatory dowolnego ideału równań?]
Jeżeli $I$ nie jest radykalny, zbiór punktów $\V_k(I)$
może mieć większy ideał znikania $I(X)=\sqrt I$.
Dla zbioru $X=\{0\}\subseteq\A^1$ wyznaczonego
równaniem $x^2=0$ pochodna podanego równania w $0$
wynosi zero i błędnie sugeruje jednowymiarową przestrzeń
styczną. Jednak jako \emph{zredukowany zbiór punktów}
$X$ ma ideał znikania $(x)$ i $T_0X=0$.
Jeśli świadomie badamy niezredukowany obiekt
$\operatorname{Spec}k[x]/(x^2)$, jego przestrzeń
styczna ma wymiar $1$. Rozróżnienie to będzie kluczowe
w geometrii schematów, choć nie rozwijamy tu ich teorii.
\end{uwaga}

\begin{przyklad}[Przestrzeń styczna w punkcie gładkiej paraboli]
Niech $X=\V_k(y-x^2)$ i $p=(a,a^2)$.
Ponieważ $f_x=-2x$ oraz $f_y=1$, mamy
$T_pX=\{(u,v)\in k^2:-2au+v=0\}$,
czyli jest to prosta o kierunku $(1,2a)$.
Zgadza się to z pochodną parametryzacji $t\mapsto(t,t^2)$.
\end{przyklad}

\begin{przyklad}[Przestrzeń styczna w ostrzu i węźle]
W początku krzywej $y^2=x^3$ pochodne $-3x^2$ i $2y$
są zerowe. Podobnie dla $y^2=x^2(x+1)$
obie pochodne znikają w $(0,0)$. W obu przypadkach
wzór na przestrzeń styczną nie nakłada żadnego warunku
na $(u,v)$, a więc $T_{(0,0)}X=k^2$.
Krzywa ma jednak lokalnie wymiar $1$. Otrzymaliśmy
zatem \emph{za dużą} przestrzeń styczną, co sygnalizuje
osobliwość, ale nie rozróżnia jeszcze dwóch jej typów.
\end{przyklad}

\section{Kryterium Jacobiego dla krzywych płaskich}

\begin{definicja}[Gładki punkt krzywej afinicznej]
\label{def:smooth-affine-curve}
Niech $k$ będzie algebraicznie domkniętym ciałem,
a $X=\V_k(f)\subseteq k^2$ zredukowaną krzywą płaską,
czyli $f\in k[x,y]$ jest niezerowy, niestały i bezkwadratowy.
Punkt $p\in X$ nazywamy \emph{gładkim}, jeżeli
$(f_x(p),f_y(p))\ne(0,0)$; w przeciwnym razie
jest \emph{osobliwy}. Gładkość może być też opisana
warunkiem $\dim_kT_pX=1$.
\end{definicja}

\begin{stwierdzenie}[Równoważność kryteriów dla jednej krzywej]
Dla krzywej z definicji~\ref{def:smooth-affine-curve}
warunki $\dim_k T_pX=1$ i
$(f_x(p),f_y(p))\ne(0,0)$ są równoważne.
W punkcie osobliwym mamy $\dim_kT_pX=2$.
\end{stwierdzenie}
\begin{proof}
Ponieważ $f$ jest bezkwadratowy i $k$ jest algebraicznie
domknięte, silne Nullstellensatz daje $I(X)=(f)$:
ideał główny $(f)$ jest radykalny w UFD, gdy $f$ nie ma
powtórzonych czynników nierozkładalnych.
Z twierdzenia~\ref{tw:zariski-tangent-formula} dostajemy
$T_pX=\{(u,v):f_x(p)u+f_y(p)v=0\}$.
Jest to jądro jednego funkcjonału liniowego $k^2\to k$.
Jeżeli funkcjonał jest niezerowy, jego jądro ma wymiar $1$;
jeżeli jest zerowy, jądro jest całym $k^2$ i ma wymiar $2$.
\end{proof}

\begin{uwaga}[Jak działa twierdzenie o funkcji uwikłanej nad $\C$?]
Jeżeli $k=\C$ i $f_y(p)\ne0$, holomorficzne twierdzenie
o funkcji uwikłanej pokazuje, że w dostatecznie małym
otoczeniu $p=(a,b)$ zbiór $f(x,y)=0$ jest wykresem
funkcji holomorficznej $y=\psi(x)$ z $\psi(a)=b$.
Wówczas $f(x,\psi(x))=0$ i po zróżniczkowaniu
$\psi'(a)=-f_x(p)/f_y(p)$. Nie jest to tylko formalna
analogia z analizą: w gładkim punkcie zespolona krzywa
algebraiczna jest lokalnie jednowymiarową rozmaitością
zespoloną. W punkcie osobliwym twierdzenie o funkcji
uwikłanej nie ma zastosowania w takiej postaci.
\end{uwaga}

\section{Stożek styczny: rozróżniamy kierunki i krotności}

\begin{definicja}[Najniższa część jednorodna]
Niech $f\in k[x,y]$ będzie niezerowym wielomianem
spełniającym $f(0,0)=0$. Zapiszmy go jednoznacznie jako
$f=f_m+f_{m+1}+\cdots+f_d$, gdzie każde $f_j$
jest jednorodne stopnia $j$, $f_m\ne0$ i $m\ge1$.
Wielomian $f_m$ nazywamy \emph{formą początkową} $f$
w początku, a liczbę $m$ \emph{rzędem znikania} $f$
w tym punkcie i oznaczamy $\operatorname{ord}_0f=m$.
Dla punktu $p=(a,b)$ przesuwamy współrzędne o $p$.
\end{definicja}

\begin{definicja}[Stożek styczny krzywej]
Niech $X=\V_k(f)\subseteq k^2$ będzie krzywą zredukowaną,
a $p\in X$. Po przesunięciu $p$ do początku wybierzmy
formę początkową $f_m$ jej zredukowanego równania.
\emph{Stożkiem stycznym} krzywej w $p$ na poziomie zbioru
nazywamy $C_pX=\V_k(f_m)\subseteq k^2$.
Wielomian $f_m$, wraz z krotnościami czynników,
nazywamy \emph{równaniem stożka stycznego z krotnościami}.
\end{definicja}

\begin{lemat}[Forma początkowa iloczynu]
Niech $f,g\in k[x,y]$ będą niezerowe i niech
$\operatorname{ord}_0f=m$, $\operatorname{ord}_0g=r$.
Wtedy $\operatorname{ord}_0(fg)=m+r$ i
$(fg)_{m+r}=f_mg_r$.
\end{lemat}
\begin{proof}
W iloczynie $(f_m+f_{m+1}+\cdots)(g_r+g_{r+1}+\cdots)$
nie ma wyrazów stopnia mniejszego niż $m+r$,
a wyraz stopnia $m+r$ jest dokładnie $f_mg_r$.
Ponieważ pierścień wielomianów nad ciałem jest dziedziną,
$f_mg_r\ne0$. Nie następuje więc skasowanie całej
najniższej części.
\end{proof}

Z lematu wynika, że jeśli ideał krzywej jest główny
$I(X)=(f)$, formy początkowe jego elementów generują
ideał $(f_m)$. Dla ideału o wielu generatorach nie wystarcza
jednak branie najniższej części każdego \emph{podanego}
generatora: w kombinacjach mogą znikać wyrazy najniższych
stopni. W ogólnym przypadku definiuje się ideał początkowy
z \emph{wszystkich} elementów ideału.

\begin{przyklad}[Trzy rodzaje lokalnego zachowania]
Dla $f=y-x^2$ najniższa część ma stopień $1$ i wynosi $y$.
Stożek styczny jest prostą $y=0$, zgodną ze zwykłą styczną.
Dla $g=y^2-x^3$ najniższa część to $y^2$:
stożek styczny jako \emph{zbiór} jest prostą $y=0$,
ale z krotnością $2$. Dla
$h=y^2-x^2(x+1)=y^2-x^2-x^3$ najniższa część
wynosi $y^2-x^2=(y-x)(y+x)$, co daje
\emph{dwie różne} styczne $y=x$ i $y=-x$,
gdy $\operatorname{char}k\ne2$.
\end{przyklad}

\begin{center}
\begin{tikzpicture}[scale=.75,>=Stealth]
 \begin{scope}[shift={(-4.15,0)}]
 \draw[->,gray!65] (-1.7,0)--(1.7,0) node[right] {$x$};
 \draw[->,gray!65] (0,-1.45)--(0,1.6) node[above] {$y$};
 \draw[blue!75!black,very thick] (-1.4,0)--(1.4,0);
 \node[below,align=center,font=\small] at (0,-1.65) {ostrze\\podwójna styczna};
 \end{scope}
 \begin{scope}[shift={(2.3,0)}]
 \draw[->,gray!65] (-1.7,0)--(1.7,0) node[right] {$x$};
 \draw[->,gray!65] (0,-1.45)--(0,1.6) node[above] {$y$};
 \draw[red!75!black,very thick] (-1.4,-1.4)--(1.4,1.4);
 \draw[red!75!black,very thick] (-1.4,1.4)--(1.4,-1.4);
 \node[below,align=center,font=\small] at (0,-1.65) {węzeł\\dwie styczne};
 \end{scope}
\end{tikzpicture}

\small Rysunek 3.2. Stożek styczny ostrza i węzła. W obu przypadkach
przestrzeń styczna Zariskiego ma wymiar $2$, ale stożki styczne
rozróżniają lokalne kierunki gałęzi.
\end{center}

\section{Krotność punktu krzywej i styczność przecięć}

\begin{definicja}[Krotność punktu krzywej płaskiej]
Niech $X=\V_k(f)\subseteq k^2$ będzie zredukowaną krzywą,
a $p\in X$. \emph{Krotnością krzywej w punkcie $p$}
nazywamy $\operatorname{mult}_pX=\operatorname{ord}_pf$,
czyli stopień najniższej niezerowej jednorodnej części
przesuniętego wielomianu $f$.
\end{definicja}

\begin{stwierdzenie}[Krotność a gładkość]
Dla zredukowanej płaskiej krzywej $X=\V_k(f)$
i punktu $p\in X$ następujące warunki są równoważne:
$p$ jest gładki, $\operatorname{mult}_pX=1$ oraz
najniższa część jednorodna równania $f$ jest niezerową
formą liniową.
\end{stwierdzenie}
\begin{proof}
Po przesunięciu $p$ do początku część liniowa $f$
ma postać $f_x(p)x+f_y(p)y$. Jest niezerowa dokładnie wtedy,
gdy co najmniej jedna z pochodnych jest niezerowa.
Z definicji to właśnie warunek $\operatorname{ord}_pf=1$.
\end{proof}

\begin{przyklad}[Ostrze i węzeł mają tę samą krotność]
Dla ostrza $y^2=x^3$ i węzła $y^2=x^2(x+1)$
mamy $\operatorname{mult}_{(0,0)}=2$.
Krotność nie rozróżnia tych osobliwości; potrzebny jest
stożek styczny lub normalizacja.
\end{przyklad}

\begin{przyklad}[Styczna do ostrza ma krotność przecięcia trzy]
Niech $C=\V_k(y^2-x^3)$ oraz $L=\V_k(y)$.
W punkcie $(0,0)$ lokalny pierścień przecięcia ma postać
\[
 \bigl(k[x,y]/(y^2-x^3,y)\bigr)_{(x,y)}
 \cong k[x]_{(x)}/(x^3).
\]
Ma on bazę $1,x,x^2$ nad $k$, więc
$I_{(0,0)}(C,L)=3$. Dla prostej $x=0$ analogicznie
otrzymujemy $k[y]_{(y)}/(y^2)$, a więc krotność $2$.
Prosta styczna $y=0$ spotyka ostrze z wyższą krotnością
niż prosta poprzeczna $x=0$.
\end{przyklad}

\section{Normalizacja: rozdzielanie gałęzi}

\begin{definicja}[Element całkowity i domknięcie całkowite]
Niech $A$ będzie dziedziną całkowitości, a
$K=\Frac(A)$ jej ciałem ułamków. Element $u\in K$
jest \emph{całkowity nad $A$}, gdy spełnia moniczne
równanie wielomianowe o współczynnikach z $A$.
Zbiór wszystkich elementów $K$ całkowitych nad $A$
nazywamy \emph{domknięciem całkowitym} $A$ w $K$
i oznaczamy $\overline A$. Dziedzinę $A$ nazywamy
\emph{normalną}, gdy $\overline A=A$.
\end{definicja}

\begin{uwaga}[Dlaczego jest to nowa informacja?]
Pierścień współrzędnych krzywej z ostrzem jest dziedziną,
a krzywa jest nierozkładalna. Mimo to pierścień ten
nie jest normalny. Brak normalności wykrywa więc
zjawisko, którego nie wykrywa sam brak dzielników zera.
\end{uwaga}

\begin{twierdzenie}[Normalizacja krzywej z ostrzem]
\label{tw:normalize-cusp}
Niech $k$ będzie algebraicznie domknięte i
$A=k[x,y]/(y^2-x^3)$. Wtedy $A$ jest dziedziną,
$\Frac(A)\cong k(t)$ przez $x=t^2$, $y=t^3$,
a domknięcie całkowite $A$ w jego ciele ułamków jest
izomorficzne z $k[t]$. Odpowiadające odwzorowanie
krzywych to $\nu:\A^1(k)\to\V_k(y^2-x^3)$,
$t\mapsto(t^2,t^3)$, które jest bijekcją, lecz nie jest
izomorfizmem algebraicznym.
\end{twierdzenie}
\begin{proof}
W poprzednim rozdziale afinicznym pokazaliśmy szczegółowo,
że homomorfizm $k[x,y]\to k[t]$,
$x\mapsto t^2$, $y\mapsto t^3$, ma jądro $(y^2-x^3)$.
Dlatego $A\cong k[t^2,t^3]\subseteq k[t]$ i jest dziedziną.
W ciele ułamków $A$ mamy $t=y/x$, ponieważ
$t^3/t^2=t$ i $x\ne0$ jako element dziedziny.
Zatem $\Frac(A)=k(t)$.

Element $t$ spełnia moniczne równanie $T^2-x=0$
o współczynnikach z $A$ i jest całkowity nad $A$.
Zatem cały $k[t]=A[t]$ jest skończonym modułem nad $A$:
każda jego potęga redukuje się do kombinacji
$1,t$ ze współczynnikami w $A$.
W szczególności każdy element $k[t]$ jest całkowity nad $A$
na mocy lematu~\ref{lem:determinant-integral}.

Odwrotnie, niech $u\in k(t)$ będzie całkowity nad $A$.
Ponieważ $A\subseteq k[t]$, to samo moniczne równanie
pokazuje, że $u$ jest całkowity nad $k[t]$.
Pierścień $k[t]$ jest pierścieniem euklidesowym, a więc
UFD; pierścień UFD jest całkowicie domknięty w swoim
ciele ułamków. Krótko uzasadnijmy ostatni fakt:
jeśli $u=a/b$ jest ułamkiem nieskracalnym i spełnia
$u^m+c_{m-1}u^{m-1}+\cdots+c_0=0$, to po pomnożeniu
przez $b^m$ otrzymujemy
$a^m=-b(c_{m-1}a^{m-1}+\cdots+c_0b^{m-1})$.
Jeżeli $b$ byłby nieodwracalny, pewien jego czynnik
nierozkładalny dzieliłby $a^m$, a więc $a$, wbrew
nieskracalności. Zatem $b$ jest jednostką i $u\in k[t]$.
Otrzymaliśmy $\overline A=k[t]$.

Ostatnie zdanie o bijektywności i braku izomorfizmu
udowodniliśmy w przykładzie~\ref{pr:cusp}.
\end{proof}

\begin{przyklad}[Normalizacja węzła rozdziela dwa punkty]
\label{pr:node-normalization}
Przyjmijmy $\operatorname{char}k\ne2$ i rozważmy
$C=\V_k(y^2-x^2(x+1))$. Odwzorowanie
\[
 \nu:\A^1(k)\to C,\qquad
 t\longmapsto(t^2-1,\ t(t^2-1))
\]
ma wartości w $C$, ponieważ dla $x=t^2-1$ zachodzi
$y^2=t^2x^2=(x+1)x^2$.
Dla $x\ne0$ parametr jest określony jednoznacznie
przez $t=y/x$. W punkcie $x=0$, $y=0$
warunek $t^2-1=0$ daje \emph{dwa} parametry $t=1$ i $t=-1$.
Dla porównania ostrze ma nad osobliwym punktem tylko
jeden parametr $t=0$.

Homomorfizm $k[x,y]/(y^2-x^2(x+1))\to k[t]$
jest różnowartościowy: każdy wielomian po redukcji modulo
równanie krzywej ma postać $A(x)+yB(x)$, a jego obraz
ma postać $A(t^2-1)+t(t^2-1)B(t^2-1)$.
Pierwszy składnik jest parzysty w $t$, drugi nieparzysty;
przy $\operatorname{char}k\ne2$ ich suma może znikać
jako wielomian tylko wtedy, gdy oba są zerowe.
Podstawienie $x=t^2-1$ jest różnowartościowe, więc
$A=B=0$. Otrzymujemy w ten sposób pierścień współrzędnych
jako podpierścień $k[t]$.
Ponieważ $t^2=x+1$ i $t=y/x$ w ciele ułamków,
rozszerzenie do $k[t]$ jest skończone i całkowite,
a ten ostatni pierścień jest normalny. Tym samym
$\nu$ realizuje normalizację krzywej z węzłem.
\end{przyklad}

\section{Rozdmuchanie punktu: zastąpienie punktu kierunkami}

\begin{definicja}[Rozdmuchanie płaszczyzny w początku]
\label{def:blowup-plane}
W przestrzeni $\A^2(k)\times\PP^1(k)$
o współrzędnych $((x,y),[s:t])$ rozważmy zbiór
\[
 B=\{((x,y),[s:t]):xt=ys\}.
\]
Odwzorowanie $\pi:B\to\A^2(k)$,
$\pi((x,y),[s:t])=(x,y)$, nazywamy
\emph{rozdmuchaniem płaszczyzny w początku}.
Włókno $E=\pi^{-1}(0,0)\cong\PP^1(k)$ to
\emph{dzielnik wyjątkowy}.
\end{definicja}

\begin{stwierdzenie}[Co dzieje się z punktami?]
Dla każdego $(x,y)\ne(0,0)$ włókno
$\pi^{-1}(x,y)$ składa się z jednego punktu,
o współrzędnej rzutowej $[s:t]=[x:y]$.
Nad $(0,0)$ całe $\PP^1(k)$ jest włóknem.
Zbiór $B$ pokrywają dwie mapy afiniczne:
przy $s\ne0$ mamy współrzędne $(u,v)$ określone przez
$x=u$, $y=uv$; przy $t\ne0$ mamy
$x=uv$, $y=u$.
\end{stwierdzenie}
\begin{proof}
Jeśli $(x,y)\ne0$, równanie $xt=ys$ oznacza,
że wektory $(x,y)$ i $(s,t)$ są proporcjonalne,
a więc wyznaczają tę samą klasę $[x:y]$.
Dla $(x,y)=0$ równanie jest tożsamością $0=0$,
więc nie ogranicza $[s:t]$.
Gdy $s\ne0$, piszemy $[s:t]=[1:v]$;
z $xt=y$ otrzymujemy $y=xv$. Ustawienie $u=x$
daje zapowiedziane współrzędne. Gdy $t\ne0$, piszemy
$[s:t]=[v:1]$, skąd $x=yv$ i można przyjąć $u=y$.
W obu mapach dzielnik wyjątkowy ma równanie $u=0$.
\end{proof}

\noindent\begin{minipage}{\linewidth}\centering
\begin{tikzpicture}[>=Stealth,scale=.87]
  % Left: image in A^2; right: symbolic local chart of blow-up.
  \draw[->,gray!70] (-2.45,0)--(2.35,0) node[right] {$x$};
  \draw[->,gray!70] (0,-1.85)--(0,1.95) node[above] {$y$};
  \fill (0,0) circle (2.3pt);
  \draw[blue!70!black,thick] (-1.65,-1.2)--(1.65,1.2);
  \draw[red!70!black,thick] (-1.6,1.3)--(1.6,-1.3);
  \node[anchor=north,font=\small] at (0,-2.02) {płaszczyzna $\A^2$};
  % A genuine blow-up map points from the new space to the original plane.
  \draw[<-,thick] (2.65,0)--(4.35,0) node[midway,above] {$\pi$};
  \draw[rounded corners=5pt,thick,gray!70,fill=gray!6]
      (4.55,-1.7) rectangle (8.35,1.65);
  \draw[thick,densely dashed] (6.15,-1.25)--(6.15,1.25);
  \draw[blue!70!black,thick] (4.77,-1.12)--(6.15,-.23);
  \draw[red!70!black,thick] (4.77,.95)--(6.15,.23);
  \fill[blue!70!black] (6.15,-.23) circle (2pt);
  \fill[red!70!black] (6.15,.23) circle (2pt);
  \node[anchor=west,font=\small,fill=white,inner sep=1pt]
      at (6.35,.95) {$E\cong\PP^1$};
  \node[anchor=north,font=\small] at (6.45,-1.83) {lokalny szkic $\widetilde{\A^2}$};
\end{tikzpicture}

\small Rysunek 3.3. Szkic lokalny (nie obraz metryczny): morfizm $\pi$ wysyła cały dzielnik wyjątkowy na początek; rozdmuchanie zastępuje początek
całą prostą rzutową kierunków. Rzeczywiste rozdmuchanie nie jest
po prostu płaską figurą naszkicowaną po prawej.
\end{minipage}

\begin{definicja}[Transformata właściwa krzywej]
Niech $C\subseteq\A^2(k)$ będzie krzywą zawierającą
$(0,0)$. \emph{Transformatą właściwą} $\widetilde C\subseteq B$
nazywamy domknięcie zbioru $\pi^{-1}(C\setminus\{(0,0)\})$
w topologii Zariskiego na $B$.
\end{definicja}

\begin{przyklad}[Rozdmuchanie ostrza]
Dla $C=\V_k(y^2-x^3)$ w mapie $x=u$, $y=uv$
podstawiamy do równania krzywej:
$u^2v^2-u^3=u^2(v^2-u)=0$.
Czynnik $u^2$ opisuje składową wspartą na dzielniku
wyjątkowym i nie należy do równania \emph{transformaty
właściwej}. Aby otrzymać jej równanie, na zbiorze $u\ne0$
dzielimy przez $u^2$, a następnie bierzemy domknięcie.
Otrzymujemy $\widetilde C:\ v^2=u$.

Krzywa $v^2=u$ jest gładka: pochodna równania
$v^2-u$ względem $u$ jest stałą $-1$.
W szczególności osobliwość samej krzywej zniknęła po
jednym rozdmuchaniu. Jednak transformata właściwa
jest styczna do dzielnika wyjątkowego $E:\ u=0$
w ich wspólnym punkcie $(u,v)=(0,0)$, ponieważ
wzór $u=v^2$ ma w początku kierunek $u=0$.
Gdy celem jest uzyskanie przecięć normalnych, potrzebne
są dalsze operacje.
\end{przyklad}

\begin{przyklad}[Rozdmuchanie węzła]
Dla $C=\V_k(y^2-x^2(x+1))$ i $\operatorname{char}k\ne2$
w tej samej mapie podstawienie $x=u$, $y=uv$ daje
$u^2(v^2-u-1)=0$; transformata właściwa spełnia
$v^2=u+1$. Na dzielniku wyjątkowym $u=0$ pojawiają się
dwa punkty $(0,1)$ i $(0,-1)$. Są różne i gładkie,
a styczne transformaty nie pokrywają się tam ze styczną
do dzielnika $u=0$: różniczka
$2v\,dv-du$ ma w każdym z tych punktów niezerowy
współczynnik przy $dv$. Rozdmuchanie rozdzieliło
obie gałęzie węzła.
\end{przyklad}

\section{Wstęp do rodzaju krzywej: algebra spotyka topologię}

Dla gładkich zespolonych krzywych rzutowych geometria
algebraiczna łączy się z topologią powierzchni zwartych.
Krzywa zredukowana, nierozkładalna i gładka nad $\C$
jest zwartą powierzchnią Riemanna, a więc ma rodzaj
$g\in\Z_{\ge0}$, oznaczający liczbę jej uchwytów.
Dla krzywych osobliwych należy rozróżnić rodzaj
arytmetyczny i rodzaj gładkiej normalizacji.

\begin{definicja}[Rodzaj arytmetyczny płaskiej krzywej z ustaloną strukturą]
Niech $k$ będzie ciałem, $S=k[X,Y,Z]$, a $F\in S$
niezerową formą jednorodną stopnia $d\ge1$.
Rozpatrujemy \emph{krzywą określoną przez jednorodny
ideał $(F)$}, czyli nie tylko zbiór punktów
$\V_+(F)$, lecz także stopniowany pierścień $S/(F)$.
Dla $m$ dostatecznie dużego jego funkcja Hilberta ma
postać $\dim_k(S/(F))_m=dm+1-p_a(F)$.
Liczbę $p_a(F)$ nazywamy \emph{rodzajem arytmetycznym}
tej krzywej ze strukturą określoną przez $(F)$.
Dla zredukowanej krzywej płaskiej jej równanie jest
wyznaczone z dokładnością do niezerowego skalara,
więc zapis $p_a(C)$ nie powoduje niejednoznaczności.
\end{definicja}

\begin{przyklad}[Dlaczego sam zbiór punktów nie wystarcza?]
Równania $F=X$ oraz $G=X^3$ wyznaczają ten sam
zbiór punktów $\V_+(F)=\V_+(G)$ w $\PP^2(k)$.
Jednak $S/(X)$ i $S/(X^3)$ nie są izomorficznymi
pierścieniami stopniowanymi: drugi zawiera niezerowe
nilpotenty. Z poniższego twierdzenia wynika
$p_a(F)=0$, natomiast $p_a(G)=1$.
Dlatego rodzaju arytmetycznego struktury \emph{nie wolno}
przypisywać dowolnemu równaniu tylko na podstawie
jego zbioru zer, bez określenia krotności.
\end{przyklad}

\begin{twierdzenie}[Rodzaj arytmetyczny krzywej płaskiej]
\label{tw:arithmetic-genus-plane}
Niech $k$ będzie ciałem, $S=k[X,Y,Z]$, a $F\in S$
niezerową formą stopnia $d\ge1$. Dla każdego $m\ge d$
funkcja Hilberta ilorazu $S/(F)$ jest równa $dm+1-p_a$,
gdzie $p_a=(d-1)(d-2)/2$.
\end{twierdzenie}
\begin{proof}
Mnożenie przez $F$ jest różnowartościowe,
ponieważ pierścień $k[X,Y,Z]$ jest dziedziną.
W dokładnym ciągu stopniowanym
$0\to S(-d)\xrightarrow{\cdot F}S\to S/(F)\to0$
w składowej stopnia $m\ge d$ dostajemy
\[
 \dim_k(S/(F))_m=\binom{m+2}{2}-\binom{m-d+2}{2}
 =dm+\frac{3d-d^2}{2}.
\]
Porównując wyraz stały z zapisem $dm+1-p_a$,
otrzymujemy $p_a=1-(3d-d^2)/2=(d-1)(d-2)/2$.
\end{proof}

\begin{uwaga}[Rodzaj topologiczny może być mniejszy]
Dla gładkiej krzywej płaskiej nad $\C$ rodzaj
arytmetyczny zgadza się z topologicznym rodzajem
powierzchni Riemanna; dowód tej głębokiej zgodności
przeprowadzimy dopiero po wprowadzeniu dywizorów
i twierdzenia Riemanna--Roch. Dla osobliwej krzywej
kubicznej $p_a=1$, chociaż normalizacje zarówno
kubiki z ostrzem, jak i z węzłem są izomorficzne
z $\PP^1$, której rodzaj topologiczny jest równy $0$.
Ta rozbieżność jest informacją o osobliwości, nie błędem
w rachunku rodzaju.
\end{uwaga}

\section{Podsumowanie: trzy poziomy badania osobliwości}

Przestrzeń styczna $T_pX$ wykrywa brak zwykłego kierunku
stycznego, lecz dla dwóch różnych osobliwości może być
jednakowo duża. Stożek styczny odróżnia pojedynczą
wielokrotną styczną od kilku różnych kierunków.
Lokalny pierścień zachowuje krotności przecięć, a jego
normalizacja może rozdzielić gałęzie przechodzące przez
ten sam punkt. Rozdmuchanie zastępuje punkt przestrzenią
kierunków i pozwala geometrycznie badać powstawanie
osobliwości oraz ich rozwiązywanie.

W kolejnym rozdziale naturalnie przejdziemy do
całościowej teorii krzywych algebraicznych:
dywizorów, funkcji wymiernych, różniczek regularnych,
twierdzenia Riemanna--Roch i rodzaju krzywej.


% Rozdział 4. Algebra przemienna (wraz z going up i płaskością), snopy i krzywe.
% Numeracja zachowuje dotychczasowy układ rozdziałów 0--3.

\chapter{Algebra przemienna, snopy i krzywe algebraiczne}
\label{ch:krzywe-snopy}

\section{Plan: dlaczego potrzebujemy kolejnych narzędzi?}
\label{sec:plan-krzywe}
W rozdziale~\ref{ch:pierscienie} nauczyliśmy się przechodzić od równań
wielomianowych do pierścieni, a w rozdziałach o geometrii afinicznej,
rzutowej i osobliwościach---od tych pierścieni do punktów i ich
otoczeń. Pozostało kilka pytań, których nie da się dobrze postawić
w języku \emph{jednego} pierścienia współrzędnych:
\begin{enumerate}[label=\arabic*)]
\item Co dokładnie oznacza ``dzielenie przez funkcję, która nie znika
w pobliżu danego punktu''? Odpowiedzią jest \emph{lokalizacja}.
\item Jak w sposób algebraiczny opisać wielkości złożone z kilku
współrzędnych lub różniczki? Pojawiają się \emph{moduły}.
\item Jak skleić dane, które są określone w różnych otwartych częściach
krzywej? Potrzebujemy \emph{snopów}.
\item Jak mierzyć zera i bieguny funkcji wymiernej? Użyjemy
\emph{pierścieni lokalnych, wartościowań i dywizorów}.
\item Kiedy istnieje funkcja o zadanych biegunach? Odpowiada na to
\emph{twierdzenie Riemanna--Roch}.
\end{enumerate}
W tej kolejności przeprowadzimy dalszą teorię. Pierwsza część
rozdziału stanowi systematyczne uzupełnienie algebry przemiennej.
Następnie konstruujemy snopy, także od strony przykładów oraz
kontrprzykładów, i dopiero wówczas przechodzimy do krzywych.

\section{Lokalizacja pierścienia: ułamki zależne od otoczenia}
\label{sec:lokalizacja}

\subsection{Dlaczego zwykłe ułamki nie wystarczą?}
Dla dziedziny $A$ można skonstruować ciało ułamków $\Frac(A)$,
dzieląc przez wszystkie niezerowe elementy. W geometrii chcemy zwykle
zrobić mniej: dzielić tylko przez takie funkcje, które \emph{w
interesującym nas miejscu} nie znikają. Na paraboli $y=x^2$
wyrażenie $1/x$ nie jest regularne w początku, ale jest regularne
na części paraboli, gdzie $x\ne0$. Z kolei wyrażenie $1/(1+x^2)$ ma sens i nie ma biegunów
w żadnym punkcie rzeczywistym paraboli, ale ma bieguny w dwóch
punktach zespolonych. Uwaga: choć na każdym punkcie
rzeczywistym funkcję tę można lokalnie zapisać jako iloraz
z niezerowym mianownikiem, $1+x^2$ nie jest jednostką w
$\R[x]$; nie jest to więc funkcja globalna na schemacie
afinicznym $\operatorname{Spec}\R[x]$. Właśnie dlatego
trzeba rozróżniać lokalną regularność na zbiorze punktów
rzeczywistych od regularności schematowej.

\begin{definicja}[Zbiór multiplikatywny]
Niech $A$ będzie pierścieniem przemiennym z jedynką.
Podzbiór $S\subseteq A$ nazywamy \emph{multiplikatywnym},
jeśli $1\in S$ i $st\in S$ dla wszystkich $s,t\in S$.
Dopuszczamy $0\in S$; lokalizacja będzie wówczas pierścieniem zerowym.
\end{definicja}

\begin{definicja}[Lokalizacja]
\label{def:loc}
Niech $A$ będzie pierścieniem, a $S\subseteq A$ zbiorem
multiplikatywnym. Na $A\times S$ definiujemy relację
\[
(a,s)\sim(b,t) \quad\Longleftrightarrow\quad
\text{istnieje }u\in S\text{ takie, że }u(at-bs)=0.
\]
Klasę pary $(a,s)$ oznaczamy $a/s$, a zbiór klas $S^{-1}A$.
Dodawanie i mnożenie definiujemy wzorami
$a/s+b/t=(at+bs)/(st)$ i $(a/s)(b/t)=ab/(st)$.
\end{definicja}

\begin{uwaga}[Po co czynnik $u$?]
Gdy $A$ jest dziedziną i $0\notin S$, warunek równości
upraszcza się do $at=bs$. W pierścieniu z dzielnikami zera
trzeba dopuścić dodatkowe $u$: po odwróceniu elementu $u$
równość $u(at-bs)=0$ staje się rzeczywistą równością
ułamków. Bez tego poprawnego warunku konstrukcja nie
miałaby wszystkich oczekiwanych własności.
\end{uwaga}

\begin{stwierdzenie}[Poprawność konstrukcji]
\label{st:loc-popr}
Dla dowolnego pierścienia przemiennego $A$ i zbioru
multiplikatywnego $S\subseteq A$ relacja z definicji~\ref{def:loc}
jest relacją równoważności, a powyższe działania na klasach
są dobrze określone. Zbiór $S^{-1}A$ jest pierścieniem.
\end{stwierdzenie}
\begin{proof}
Zwrotność: dla $(a,s)$ wybieramy $u=1$ i mamy
$u(as-as)=0$. Symetria wynika ze zmiany znaku.
Jeśli $(a,s)\sim(b,t)$ oraz $(b,t)\sim(c,r)$,
to istnieją $u,v\in S$ z $u(at-bs)=0$ i $v(br-ct)=0$.
Po pomnożeniu pierwszej równości przez $vr$,
a drugiej przez $us$, i dodaniu dostajemy
$uvt(ar-cs)=0$, a ponieważ $uvt\in S$, relacja jest
przechodnia.

Sprawdzimy niezależność sumy od reprezentantów. Wystarczy
zmienić najpierw jeden składnik: niech $u(at-bs)=0$.
Dla dowolnego $c/r$ porównujemy
$(ar+cs)/(sr)$ i $(br+ct)/(tr)$.
Różnica iloczynów krzyżowych liczników i mianowników
jest równa
\[
 (ar+cs)tr-(br+ct)sr=r^2(at-bs).
\]
Po pomnożeniu przez $u$ otrzymujemy zero. Dla iloczynu
porównanie $ac/(sr)$ i $bc/(tr)$ daje różnicę
$cr(at-bs)$, którą także anuluje $u$. Drugiego
reprezentanta zmienia się identycznie.
Łączność, przemienność i rozdzielność działań
wynikają teraz z odpowiednich równości w $A$;
zerem jest $0/1$, a jedynką $1/1$.
\end{proof}

\begin{twierdzenie}[Własność uniwersalna lokalizacji]
\label{tw:loc-up}
Niech $A,B$ będą pierścieniami przemiennymi z jedynką,
$S\subseteq A$ zbiorem multiplikatywnym i
$\varphi:A\to B$ homomorfizmem takim, że każdy
$\varphi(s)$, $s\in S$, jest odwracalny w $B$.
Istnieje dokładnie jeden homomorfizm
$\widetilde\varphi:S^{-1}A\to B$ spełniający
$\widetilde\varphi(a/1)=\varphi(a)$ dla $a\in A$.
Jest nim $\widetilde\varphi(a/s)=\varphi(a)\varphi(s)^{-1}$.
\end{twierdzenie}
\begin{proof}
Skoro $a/s=(a/1)(s/1)^{-1}$ w $S^{-1}A$,
każdy szukany homomorfizm musi spełniać podany wzór;
stąd jednoznaczność. Jeśli $(a,s)\sim(b,t)$,
istnieje $u\in S$ takie, że $u(at-bs)=0$. Po
zastosowaniu $\varphi$ i pomnożeniu przez odwracalne
$\varphi(u)^{-1}\varphi(s)^{-1}\varphi(t)^{-1}$
dostajemy $\varphi(a)\varphi(s)^{-1}
=\varphi(b)\varphi(t)^{-1}$.
Wzór nie zależy więc od reprezentanta i bezpośredni
rachunek pokazuje, że zachowuje dodawanie, mnożenie
i jedynkę.
\end{proof}

\begin{przyklad}[Trzy lokalizacje tych samych liczb całkowitych]
\begin{enumerate}[label=\alph*)]
\item Dla $S=\{2^n:n\ge0\}$ otrzymujemy
$S^{-1}\Z=\Z[1/2]$, czyli ułamki o mianownikach
będących potęgami $2$.
\item Dla $S=\Z\setminus\{0\}$ mamy
$S^{-1}\Z=\Q$.
\item Dla $S=\Z\setminus(5)$ otrzymujemy
$\Z_{(5)}=\{a/b\in\Q:5\nmid b\}$.
W tym pierścieniu $2$ i $3$ są odwracalne, a $5$ nie.
\end{enumerate}
Zwróćmy uwagę, że $\Z[1/5]$ i $\Z_{(5)}$ są
\emph{różne}: w pierwszym odwracamy potęgi $5$,
a w drugim wszystkie liczby niepodzielne przez $5$.
\end{przyklad}

\subsection{Lokalizacja w elemencie i w ideale pierwszym}
\begin{definicja}
Dla $f\in A$ piszemy $A_f=S^{-1}A$, gdzie
$S=\{1,f,f^2,\ldots\}$. Dla ideału pierwszego
$\mathfrak p\subsetneq A$ piszemy $A_{\mathfrak p}$,
gdzie $S=A\setminus\mathfrak p$. To rzeczywiście
zbiór multiplikatywny: z $ab\in\mathfrak p$ wynika
$a\in\mathfrak p$ lub $b\in\mathfrak p$.
\end{definicja}

\begin{twierdzenie}[Lokalizacja przy ideale pierwszym jest lokalna]
\label{tw:Ap-local}
Dla dowolnego pierścienia $A$ i ideału pierwszego
$\mathfrak p\subset A$ pierścień $A_{\mathfrak p}$
ma dokładnie jeden ideał maksymalny:
$\mathfrak p A_{\mathfrak p}$.
\end{twierdzenie}
\begin{proof}
Każdy ułamek $a/s$ z $a\notin\mathfrak p$ jest
odwracalny: jego odwrotność wynosi $s/a$, ponieważ
$a\in A\setminus\mathfrak p$. Jeśli $a\in\mathfrak p$,
ułamek $a/s$ nie jest odwracalny, ponieważ obraz
$a/s$ w ciele $\Frac(A/\mathfrak p)$ przez
naturalny homomorfizm wynosi zero. Stąd
nieodwracalne elementy $A_{\mathfrak p}$ tworzą
dokładnie zbiór $\mathfrak p A_{\mathfrak p}$.
Jest to ideał właściwy, a każdy ideał właściwy
nie zawiera jednostek, więc w nim się zawiera.
Jest zatem jedynym ideałem maksymalnym.
\end{proof}

\begin{uwaga}[Pierścień lokalny]
Pierścień $R$ z dokładnie jednym ideałem maksymalnym
$\mathfrak m$ nazywamy \emph{lokalnym} i piszemy
$(R,\mathfrak m)$. Element $r\in R$ jest jednostką
dokładnie wtedy, gdy $r\notin\mathfrak m$.
W geometrii $A_{\mathfrak p}$ zawiera informacje
w bezpośrednim sąsiedztwie punktu opisanego przez
$\mathfrak p$: elementy spoza $\mathfrak p$ uznajemy
za odwracalne, bo ``nie znikają w tym punkcie''.
\end{uwaga}

\begin{przyklad}[Dwie funkcje w jednym punkcie]
Niech $A=k[x,y]/(y-x^2)\cong k[x]$ i niech
$\mathfrak m=(x,y)$ opisuje początek paraboli.
W pierścieniu $A_{\mathfrak m}\cong k[x]_{(x)}$
element $(1+x)^{-1}$ jest dozwolony, bo $1+x$
nie znika w początku. Natomiast $x^{-1}$
nie jest elementem tego pierścienia: $x$ należy
do jedynego ideału maksymalnego.
\end{przyklad}

\begin{twierdzenie}[Ideały i lokalizacja]
\label{tw:idealy-loc}
Niech $S$ będzie zbiorem multiplikatywnym pierścienia $A$.
Ideały pierwsze pierścienia $S^{-1}A$ odpowiadają
bijektywnie ideałom pierwszym $\mathfrak p\subset A$
takim, że $\mathfrak p\cap S=\varnothing$.
Odpowiedniość jest dana przez
$\mathfrak p\mapsto S^{-1}\mathfrak p$ i przeciwobraz
względem $A\to S^{-1}A$.
\end{twierdzenie}
\begin{proof}
Jeśli $\mathfrak p\cap S=\varnothing$, to
$S^{-1}\mathfrak p$ jest właściwy: inaczej $1=a/s$
dla pewnego $a\in\mathfrak p$, co oznaczałoby
$u(a-s)=0$ dla pewnego $u\in S$, skąd $us=ua\in\mathfrak p$,
sprzecznie z rozłącznością. Jego pierwszość
wynika z pierwszości $\mathfrak p$: jeżeli
$(a/s)(b/t)\in S^{-1}\mathfrak p$, to istnieje
$u\in S$ z $uab\in\mathfrak p$; ponieważ
$u\notin\mathfrak p$, otrzymujemy $a\in\mathfrak p$
lub $b\in\mathfrak p$.
Jeśli $\mathfrak q$ jest pierwszy w $S^{-1}A$,
jego przeciwobraz $\mathfrak p$ jest pierwszy
i nie przecina $S$, bo obrazy elementów $S$ są
jednostkami. Każdy $a/s\in\mathfrak q$ ma
$a/1=(s/1)(a/s)\in\mathfrak q$, stąd
$\mathfrak q=S^{-1}\mathfrak p$. Analogicznie
przeciwobraz $S^{-1}\mathfrak p$ jest dokładnie
$\mathfrak p$: jeśli $a/1\in S^{-1}\mathfrak p$,
to $ua\in\mathfrak p$ dla pewnego $u\in S$,
a z pierwszości wynika $a\in\mathfrak p$.
\end{proof}

\begin{center}
\begin{tikzpicture}[>=Stealth, every node/.style={font=\small}]
\node[draw,rounded corners,minimum width=3.2cm,minimum height=1cm] (a) at (0,0) {$A$};
\node[draw,rounded corners,minimum width=3.2cm,minimum height=1cm] (b) at (6,0) {$A_{\mathfrak p}$};
\draw[->,thick] (a) -- node[above]{odwracamy $A\setminus\mathfrak p$} (b);
\node[below=3mm of a,text width=3.4cm,align=center] {informacja globalna};
\node[below=3mm of b,text width=3.4cm,align=center] {informacja w pobliżu $\mathfrak p$};
\end{tikzpicture}

{\small Rysunek 4.1. Lokalizacja zachowuje informacje w sąsiedztwie ideału pierwszego, odwracając elementy spoza tego ideału.}\par
\end{center}

\section{Moduły: przestrzenie liniowe nad pierścieniem}
\label{sec:moduly}

\begin{definicja}[Moduł]
Niech $A$ będzie pierścieniem przemiennym. \emph{$A$-modułem}
nazywamy grupę abelową $(M,+)$ wyposażoną w działanie
$A\times M\to M$, $(a,m)\mapsto am$, dla którego
$(a+b)m=am+bm$, $a(m+n)=am+an$,
$(ab)m=a(bm)$ oraz $1m=m$ dla wszystkich
$a,b\in A$, $m,n\in M$. Homomorfizm $A$-modułów
jest homomorfizmem grup respektującym mnożenie przez $A$.
\end{definicja}

\begin{przyklad}[Moduł nie jest zwykle przestrzenią wektorową]
\begin{enumerate}[label=\alph*)]
\item Każda przestrzeń wektorowa nad ciałem $k$ jest $k$-modułem.
\item Grupa abelowa $M$ jest $\Z$-modułem przez
$nm=m+\cdots+m$ dla $n\ge0$ i przez elementy przeciwne
dla $n<0$. Odwrotnie, aksjomaty $\Z$-modułu wymuszają
takie mnożenie.
\item $A^r$ jest $A$-modułem wolnym z bazą $e_1,\ldots,e_r$.
\item Iloraz $A/I$ jest $A$-modułem. Na przykład
$\Z/6\Z$ nie jest wolnym $\Z$-modułem: $6\bar1=0$,
podczas gdy w wolnym module nad $\Z$ niezerowy element
nie może być anihilowany przez $6$.
\end{enumerate}
\end{przyklad}

\begin{definicja}[Generatory i relacje]
Moduł $M$ jest \emph{skończenie generowany}, jeżeli istnieją
$m_1,\ldots,m_r\in M$ takie, że każdy $m\in M$ ma postać
$m=\sum_{i=1}^r a_i m_i$ z $a_i\in A$.
Równoważnie istnieje surjektywny homomorfizm
$A^r\twoheadrightarrow M$. Jego jądro opisuje
\emph{relacje} pomiędzy generatorami.
\end{definicja}

\begin{przyklad}[Relacje mają znaczenie geometryczne]
Niech $A=k[x,y]$ i $M=A^2/A\,(y,-x)$.
W $M$ zachodzi relacja $y\bar e_1=x\bar e_2$.
Formalnie mamy dwa generatory, ale nie są one niezależne.
Moduły tego rodzaju pojawiają się przy badaniu pól
stycznych, różniczek i wiązek wektorowych: lokalnie
mogą zachowywać się jak $A^r$, choć globalnie
nie muszą mieć jednej bazy.
\end{przyklad}

\begin{definicja}[Lokalizacja modułu]
Niech $S\subseteq A$ będzie multiplikatywny, a $M$
$A$-modułem. Definiujemy $S^{-1}M$ jako klasy par
$(m,s)$ modulo relacji
$(m,s)\sim(n,t)$ wtedy i tylko wtedy, gdy istnieje
$u\in S$ takie, że $u(tm-sn)=0$ w $M$.
Działania to $(m/s)+(n/t)=(tm+sn)/(st)$ oraz
$(a/s)(m/t)=(am)/(st)$. Jest to $S^{-1}A$-moduł.
\end{definicja}

\begin{twierdzenie}[Lokalizacja zachowuje dokładność]
\label{tw:loc-exact}
Niech $A$ będzie pierścieniem, $S\subseteq A$ zbiorem
multiplikatywnym i
$M'\xrightarrow{f}M\xrightarrow{g}M''$ ciągiem
$A$-modułów. Jeśli $\im f=\ker g$, to
$\im(S^{-1}f)=\ker(S^{-1}g)$. W szczególności
lokalizacja zachowuje krótkie ciągi dokładne.
\end{twierdzenie}
\begin{proof}
Ponieważ $gf=0$, również $(S^{-1}g)(S^{-1}f)=0$,
więc obraz zawiera się w jądrze. Odwrotnie,
niech $m/s\in S^{-1}M$ należy do jądra
$S^{-1}g$. Z definicji równości ułamków istnieje
$u\in S$ takie, że $u\,g(m)=0$. Zatem
$g(um)=0$, czyli $um=f(m')$ dla pewnego $m'\in M'$.
W lokalizacji mamy
$m/s=(um)/(us)=f(m')/(us)$, więc $m/s$
należy do obrazu. Aby uzyskać pełną dokładność
krótkiego ciągu, stosujemy argument do obu kolejnych
par odwzorowań; na lewym końcu bierzemy $M'=0$.
\end{proof}

\begin{uwaga}[Dlaczego to ważne?]
Lokalizując dokładny ciąg
$0\to N\to M\to M/N\to0$,
widzimy, że $(M/N)_{\mathfrak p}\cong
M_{\mathfrak p}/N_{\mathfrak p}$. Możemy więc
sprawdzić wiele własności modułów \emph{punkt po punkcie}.
Z drugiej strony samo badanie modułów tylko globalnie
może ukrywać strukturę w konkretnych punktach.
\end{uwaga}

\subsection{Moduły noetherowskie i lemat Nakayamy}
\begin{definicja}
$A$-moduł $M$ nazywamy \emph{noetherowskim},
gdy każdy rosnący łańcuch podmodułów
$M_1\subseteq M_2\subseteq\cdots$ stabilizuje się.
Pierścień $A$ nazywamy \emph{noetherowskim},
gdy jest noetherowski jako moduł nad samym sobą,
czyli każdy jego ideał jest skończenie generowany.
\end{definicja}

\begin{twierdzenie}[Własności modułów noetherowskich]
\label{tw:noeth-mod}
Jeśli $A$ jest noetherowski, to każdy skończenie
generowany $A$-moduł $M$ jest noetherowski.
W szczególności każdy jego podmoduł jest skończenie
generowany.
\end{twierdzenie}
\begin{proof}
Najpierw dowodzimy, że dla każdego krótkiego
ciągu dokładnego $0\to N\to M\to Q\to0$
moduł $M$ jest noetherowski, gdy $N$ i $Q$ są
noetherowskie. Istotnie, dla rosnącego łańcucha
$M_i\subseteq M$ zarówno łańcuchy $M_i\cap N$,
jak i obrazów $M_i$ w $Q$ stabilizują się.
Jeśli od indeksu $i_0$ oba są stałe i $m\in M_j$
dla $j\ge i_0$, to obraz $m$ w $Q$ ma reprezentanta
$m_0\in M_{i_0}$, zatem $m-m_0\in M_j\cap N
=M_{i_0}\cap N$, czyli $m\in M_{i_0}$.
To dowodzi stabilizacji $M_i$.

Pierścień $A$ jest noetherowski jako moduł
na mocy założenia. Przez indukcję po $r$, stosując
$0\to A^{r-1}\to A^r\to A\to0$, otrzymujemy
noetherowskość $A^r$. Jeśli $M$ jest generowany
przez $r$ elementów, jest obrazem $A^r$;
każdy rosnący łańcuch podmodułów $M$ ma rosnący
łańcuch przeciwobrazów w $A^r$, więc stabilizuje się.
Na koniec: jeśli podmoduł $N\subseteq M$ nie byłby
skończenie generowany, moglibyśmy wybrać
$n_1\in N$, $n_2\in N\setminus An_1$, itd.,
tworząc nieskończony ściśle rosnący łańcuch
$An_1\subsetneq An_1+An_2\subsetneq\cdots$.
Sprzeczność kończy dowód.
\end{proof}

\begin{lemat}[Nakayama---wersja podstawowa]
\label{lem:nakayama}
Niech $(A,\mathfrak m)$ będzie pierścieniem lokalnym,
a $M$ skończenie generowanym $A$-modułem.
Jeżeli $\mathfrak m M=M$, to $M=0$.
\end{lemat}
\begin{proof}
Załóżmy przeciwnie, że $M\ne0$. Wybierzmy
minimalną liczbę generatorów $m_1,\ldots,m_r$,
$r\ge1$. Z równości $M=\mathfrak m M$ mamy
$m_r=\sum_{i=1}^r a_i m_i$ z $a_i\in\mathfrak m$.
Przenosząc ostatni wyraz, otrzymujemy
$(1-a_r)m_r=\sum_{i<r}a_i m_i$.
Element $1-a_r$ nie należy do $\mathfrak m$:
gdyby należał, to także $1=(1-a_r)+a_r$
należałby do $\mathfrak m$. Jest zatem jednostką,
więc $m_r$ jest kombinacją $m_1,\ldots,m_{r-1}$.
Sprzeczność z minimalnością liczby generatorów.
\end{proof}

\begin{wniosek}[Minimalna liczba generatorów]
\label{wn:generatory-local}
Jeśli $(A,\mathfrak m)$ jest noetherowskim pierścieniem
lokalnym i $M$ jest skończenie generowanym $A$-modułem,
to minimalna liczba generatorów $M$ jest równa
$\dim_{A/\mathfrak m}(M/\mathfrak mM)$.
\end{wniosek}
\begin{proof}
Jeśli $M$ ma $r$ generatorów, ich klasy generują
przestrzeń wektorową $M/\mathfrak mM$, zatem jej
wymiar nie przekracza $r$. Odwrotnie, niech
$\bar m_1,\ldots,\bar m_d$ będzie bazą tej przestrzeni.
Dla $N=Am_1+\cdots+Am_d$ mamy
$M=N+\mathfrak mM$, więc
$(M/N)=\mathfrak m(M/N)$. Moduł $M/N$ jest
skończenie generowany; lemat~\ref{lem:nakayama}
daje $M/N=0$. Zatem $m_i$ generują $M$.
\end{proof}

\begin{uwaga}[Ostrożnie z hipotezą skończonej generacji]
Dla $A=k[[t]]$ i $M=k((t))$ mamy $tM=M$,
ale $M\ne0$. Nie przeczy to Nakayamie: $M$
nie jest skończenie generowany jako $A$-moduł.
Tutaj $k[[t]]$ oznacza pierścień formalnych
szeregów potęgowych, a $k((t))$ ich ciało ułamków
(szeregi Laurenta z ograniczoną od dołu liczbą wykładników).
\end{uwaga}

\section{Iloczyn tensorowy i płaskość: kiedy równania pozostają dokładne?}
\label{sec:flatness}

\begin{definicja}[Iloczyn tensorowy modułów]
Niech $M,N$ będą modułami nad pierścieniem przemiennym $A$.
Iloczyn tensorowy $M\otimes_A N$ jest $A$-modułem,
wyposażonym w odwzorowanie $A$-dwuliniowe
$(m,n)\mapsto m\otimes n$, o następującej własności:
dla każdego $A$-modułu $P$ i każdego odwzorowania
$A$-dwuliniowego $b:M\times N\to P$ istnieje dokładnie
jeden homomorfizm $A$-modułów
$\widehat b:M\otimes_A N\to P$ spełniający
$\widehat b(m\otimes n)=b(m,n)$.
\end{definicja}

\begin{stwierdzenie}[Konstrukcja iloczynu tensorowego]
Niech $A$ będzie pierścieniem przemiennym z jedynką, a $M,N$
modułami nad $A$. Istnieje $A$-moduł $M\otimes_A N$
wraz z odwzorowaniem $A$-dwuliniowym
$M\times N\to M\otimes_A N$, spełniający własność
uniwersalną z poprzedniej definicji. Jest on jednoznaczny
z dokładnością do jednoznacznego izomorfizmu zachowującego
kanoniczne odwzorowanie dwuliniowe.
\end{stwierdzenie}
\begin{proof}
Niech $F$ będzie wolnym $A$-modułem z bazą
formalnych symboli $[m,n]$, po jednym dla każdej
pary $(m,n)\in M\times N$. Niech $R$ będzie
podmodułem generowanym przez
\[
\begin{aligned}
 &[m+m',n]-[m,n]-[m',n],\\
 &[m,n+n']-[m,n]-[m,n'],\\
 &[am,n]-a[m,n],\qquad [m,an]-a[m,n]
\end{aligned}
\]
dla wszystkich dopuszczalnych argumentów.
Definiujemy $M\otimes_A N=F/R$ i
$m\otimes n=[m,n]+R$. Przez konstrukcję
to odwzorowanie jest dwuliniowe.
Jeżeli $b:M\times N\to P$ jest dwuliniowe,
to funkcja na bazie $[m,n]\mapsto b(m,n)$
przedłuża się liniowo do $F\to P$, a każdą
wypisaną relację wysyła na zero.
Faktoryzuje się więc jednoznacznie przez $F/R$.
Ta własność uniwersalna wymusza też izomorfizm
dowolnych dwóch konstrukcji o takiej własności.
\end{proof}

\begin{stwierdzenie}[Iloczyn tensorowy z modułem wolnym]
Dla dowolnego $A$-modułu $M$ i liczby $r\ge0$
istnieje naturalny izomorfizm
$M\otimes_A A^r\cong M^r$.
\end{stwierdzenie}
\begin{proof}
Niech $e_1,\ldots,e_r$ będzie bazą $A^r$.
Odwzorowanie $m\otimes(a_1,\ldots,a_r)
\mapsto(a_1m,\ldots,a_rm)$ jest dobrze określone
z własności uniwersalnej. W przeciwną stronę
kładziemy $(m_1,\ldots,m_r)\mapsto
\sum_i m_i\otimes e_i$. Zależność
$m\otimes\sum_i a_i e_i=\sum_i a_i m\otimes e_i$
pokazuje, że odwzorowania są wzajemnie odwrotne.
\end{proof}

\begin{definicja}[Moduł płaski]
$A$-moduł $N$ nazywamy \emph{płaskim},
jeżeli dla każdego krótkiego ciągu dokładnego
$A$-modułów $0\to M'\to M\to M''\to0$
ciąg po zastosowaniu iloczynu tensorowego
$0\to M'\otimes_A N\to M\otimes_A N
\to M''\otimes_A N\to0$ pozostaje dokładny.
\end{definicja}

\begin{przyklad}[Dlaczego płaskość nie jest automatyczna?]
Weźmy $A=\Z$ i dokładny ciąg
$0\to\Z\xrightarrow{\cdot2}\Z\to\Z/2\Z\to0$.
Po zastosowaniu iloczynu tensorowego z $N=\Z/2\Z$ pierwszy
homomorfizm staje się mnożeniem przez $2$
na $\Z/2\Z$, czyli odwzorowaniem zerowym.
Nie jest różnowartościowy, zatem $\Z/2\Z$
\emph{nie jest płaskim $\Z$-modułem}.
Z kolei każdy moduł wolny $A^r$ jest płaski:
iloczyn tensorowy z $A^r$ jest naturalnie izomorficzny
z braniem $r$ kopii wyjściowego ciągu,
a bezpośrednia suma dokładnych ciągów pozostaje
dokładna.
\end{przyklad}

\begin{uwaga}[Związek z geometrią]
Moduły lokalnie wolne reprezentują na schemacie
wiązki wektorowe i są lokalnie płaskie.
Dla rodzin obiektów algebraicznych płaskość
odgrywa rolę warunku wykluczającego patologiczne
``nagłe zmiany'' struktury włókien. Nie utożsamiamy
jednak płaskości z identyczną geometrią wszystkich
włókien: płaska rodzina krzywych może zawierać
zarówno krzywe gładkie, jak i osobliwe.
\end{uwaga}


\begin{twierdzenie}[Lokalizacja jako płaski moduł]
\label{tw:lokalizacja-plaska}
Niech $A$ będzie pierścieniem przemiennym, a $S\subseteq A$
zbiorem multiplikatywnym. Wtedy $S^{-1}A$ jest płaskim
$A$-modułem. Dokładniej, dla każdego $A$-modułu $M$
istnieje naturalny izomorfizm
$M\otimes_A S^{-1}A\cong S^{-1}M$.
\end{twierdzenie}
\begin{proof}
Odwzorowanie $M\times S^{-1}A\to S^{-1}M$ zadane
$(m,a/s)\mapsto am/s$ jest addytywne względem obu
argumentów i spełnia warunek $A$-zrównoważenia:
$(bm)(a/s)=(m)(ba/s)$. Z własności uniwersalnej
tensora otrzymujemy homomorfizm
$\Phi:M\otimes_A S^{-1}A\to S^{-1}M$,
$\Phi(m\otimes a/s)=am/s$.

W drugą stronę niech $\Psi(m/s)=m\otimes 1/s$.
Sprawdzimy, że jest to określenie niezależne od
reprezentanta. Jeśli $m/s=n/t$, istnieje $u\in S$
z $u(tm-sn)=0$. W module tensorowym mamy
$(tm-sn)\otimes 1/(st)=0$, bo mnożenie przez
$u$ na drugim czynniku jest odwracalne:
\[
 (tm-sn)\otimes\frac1{st}
 =u(tm-sn)\otimes\frac1{ust}=0.
\]
Z dwuliniowości
$(tm)\otimes 1/(st)=m\otimes 1/s$ i
$(sn)\otimes 1/(st)=n\otimes 1/t$;
stąd $\Psi$ jest dobrze określone. Sprawdzenie
addytywności wynika ze wspólnego mianownika.
Dla elementarnych tensorów mamy
$\Psi\Phi(m\otimes a/s)=am\otimes1/s
=m\otimes a/s$, a dla ułamków
$\Phi\Psi(m/s)=m/s$. Są to więc izomorfizmy
wzajemnie odwrotne. Na mocy twierdzenia
\ref{tw:loc-exact} funktor $S^{-1}(-)$ zachowuje
krótkie ciągi dokładne, zatem zachowuje je też
zastosowanie iloczynu tensorowego z $S^{-1}A$. Jest to definicja
płaskości.
\end{proof}

\begin{przyklad}[Co oznacza płaskość geometrycznie?]
Niech $A=k[t]$, $S=\{1,t,t^2,\ldots\}$.
Moduł $A_t=k[t,t^{-1}]$ jest płaski nad $A$:
po usunięciu punktu $t=0$ wszystkie relacje,
które były dokładne przed lokalizacją, pozostają
dokładne. Natomiast przejście do włókna
$A/(t)\cong k$ nie jest płaskim $A$-modułem:
ciąg $0\to A\xrightarrow{\cdot t}A$ po
zastosowaniu iloczynu tensorowego z $k$ traci różnowartościowość.
\end{przyklad}

\section{Zależność całkowita, normalizacja i wymiar Krulla}
\label{sec:integral}
\begin{definicja}[Element całkowity]
Niech $A\subseteq B$ będą pierścieniami przemiennymi.
Element $b\in B$ jest \emph{całkowity nad $A$},
gdy istnieją $a_0,\ldots,a_{n-1}\in A$ takie, że
$b^n+a_{n-1}b^{n-1}+\cdots+a_0=0$.
Wielomian jest tu \emph{moniczny}, czyli współczynnik
przy $b^n$ jest równy $1$. Rozszerzenie $A\subseteq B$
nazywamy całkowitym, gdy każdy $b\in B$ jest całkowity.
\end{definicja}

\begin{przyklad}[Dlaczego moniczność jest ważna?]
Liczba $\sqrt2$ jest całkowita nad $\Z$, ponieważ
spełnia $T^2-2=0$. Element $1/2\in\Q$ spełnia
$2T-1=0$, ale nie jest całkowity nad $\Z$.
Gdyby spełniał moniczne równanie stopnia $n$,
po pomnożeniu przez $2^n$ otrzymalibyśmy
$1+2a_{n-1}+\cdots+2^na_0=0$, sprzeczność
modulo $2$. Sama możliwość znalezienia jakiegokolwiek
równania wielomianowego nic tu nie mówi.
\end{przyklad}

\begin{twierdzenie}[Kryterium skończonego modułu]
\label{tw:integrality-module}
Niech $A\subseteq B$ będą pierścieniami, a $b\in B$.
Wówczas $b$ jest całkowity nad $A$ wtedy i tylko wtedy,
gdy podpierścień $A[b]$ jest skończenie generowanym
$A$-modułem.
\end{twierdzenie}
\begin{proof}
Jeśli $b^n+a_{n-1}b^{n-1}+\cdots+a_0=0$,
to każdą potęgę $b^m$, $m\ge n$, można wyrazić
jako kombinację $1,b,\ldots,b^{n-1}$: najpierw
zapisujemy $b^n=-\sum_{i<n}a_ib^i$,
a potem mnożymy przez kolejne potęgi $b$ i
indukcyjnie redukujemy wykładnik. Zatem
$A[b]$ jest generowany przez $n$ elementów.

W przeciwną stronę niech $e_1,\ldots,e_r$ generują
$A[b]$ jako $A$-moduł. Mnożenie przez $b$
daje endomorfizm tego modułu; istnieje macierz
$C=(c_{ij})\in M_r(A)$ taka, że
$b e_j=\sum_i c_{ij}e_i$. Po przeniesieniu na jedną
stronę mamy $(bI-C)e=0$ dla wektora generatorów $e$.
Mnożąc przez macierz dołączoną, otrzymujemy
$\det(bI-C)e_j=0$ dla każdego $j$. Ponieważ
$1\in A[b]$ jest kombinacją $e_j$, wynika
$\det(bI-C)=0$ w $B$. Wielomian
$\det(TI-C)\in A[T]$ jest moniczny stopnia $r$,
a więc $b$ jest całkowity.
\end{proof}

\begin{uwaga}
Dowód w drugą stronę nie wymaga istnienia bazy modułu
$A[b]$: wystarcza skończony układ \emph{generatorów}.
Dzięki macierzy dołączonej nie użyliśmy też dzielenia
przez wyznacznik, co byłoby niedozwolone w ogólnym
pierścieniu.
\end{uwaga}

\begin{definicja}[Pierścień normalny]
Niech $A$ będzie dziedziną. Nazywamy ją
\emph{całkowicie domkniętą} lub \emph{normalną},
gdy każdy element $\Frac(A)$ całkowity nad $A$
już należy do $A$. Podpierścień $\Frac(A)$
utworzony przez wszystkie elementy całkowite nad $A$
nazywamy \emph{domknięciem całkowitym} lub
\emph{normalizacją algebraiczną} $A$.
\end{definicja}

\begin{przyklad}[Normalizacja ostrza---rachunek]
\label{prz:normalizacja-ostrza}
Niech $A=k[t^2,t^3]\subseteq k[t]$, gdzie $k$ jest ciałem.
Element $t$ należy do ciała ułamków $A$, bo
$t=t^3/t^2$, a spełnia równanie $T^2-t^2=0$,
w którym $t^2\in A$. Jest więc całkowity nad $A$.
Jednak $t\notin A$: każdy niezerowy jednomian $t^m$
w $A$ ma wykładnik $m=0$ albo $m\ge2$.
Zatem $A$ nie jest normalny, mimo że jest dziedziną.

Pierścień $k[t]$ jest całkowicie domknięty, ponieważ
jest pierścieniem euklidesowym, a zatem UFD: jeśli
$a/b\in\Frac(k[t])$ zapisano w postaci nieskracalnej
i $a/b$ spełnia moniczne równanie nad $k[t]$, to
$b\mid a^n$ i z pierwszości czynników nierozkładalnych
wynika, że $b$ jest jednostką. Każdy element $k[t]$
jest całkowity nad $A$, ponieważ
$k[t]=A+At$ jako $A$-moduł. A jeśli $z\in k(t)$
jest całkowity nad $A$, jest też całkowity nad $k[t]$
(każde moniczne równanie nad $A$ jest równaniem nad
$k[t]$), więc $z\in k[t]$. Wobec tego normalizacją
$A$ jest dokładnie $k[t]$.
Geometrycznie parametryzacja $(x,y)=(t^2,t^3)$
prostuje ostrze $y^2=x^3$.
\end{przyklad}

\subsection{Ideały prymarne: zachowanie krotności}
\begin{definicja}[Ideał prymarny]
Ideał $Q\subsetneq A$ nazywamy \emph{prymarnym},
gdy z $ab\in Q$ i $a\notin Q$ wynika
$b^n\in Q$ dla pewnego $n\ge1$.
Dla porównania: dla ideału \emph{pierwszego} można
w powyższym warunku wziąć zawsze $n=1$.
\end{definicja}

\begin{stwierdzenie}[Pierwiastek ideału prymarnego jest pierwszy]
Jeśli $Q\subsetneq A$ jest prymarny, to
$\sqrt Q=\{a\in A:a^n\in Q\text{ dla pewnego }n\ge1\}$
jest ideałem pierwszym.
\end{stwierdzenie}
\begin{proof}
Załóżmy, że $ab\in\sqrt Q$, czyli
$(ab)^r=a^rb^r\in Q$ dla pewnego $r\ge1$.
Jeśli $a\notin\sqrt Q$, to $a^r\notin Q$.
Z prymarności $Q$ wynika, że
$(b^r)^n=b^{rn}\in Q$ dla pewnego $n$.
Zatem $b\in\sqrt Q$. Ponieważ $Q$ jest właściwy,
również $\sqrt Q$ jest właściwy: inaczej
$1\in\sqrt Q$, skąd $1\in Q$. To dokładnie
warunek pierwszości.
\end{proof}

\begin{przyklad}[Punkt i punkt pogrubiony]
W $k[x]$ ideały $(x)$ i $(x^3)$ mają ten sam
pierwiastek $(x)$, więc ich miejsca zerowe
w $\A^1_k$ są identyczne: to punkt $0$.
Jednak w ilorazie $k[x]/(x^3)$ klasa $x$
jest niezerowa i nilpotentna rzędu trzy.
Ideał $(x^3)$ jest prymarny: jeśli $x^3\mid ab$
i $x^3\nmid a$, to rozkład na potęgi czynnika
$x$ pokazuje, że pewna potęga $b$ jest podzielna
przez $x^3$. Algebra zachowuje więc krotność,
której sam zbiór punktów nie pamięta.
\end{przyklad}

\begin{twierdzenie}[Rozkład prymarny ideałów wielomianowych]
\label{tw:primary-kx}
Niech $k$ będzie ciałem, a $0\ne f\in k[x]$
niech ma rozkład
$f=u\,p_1^{e_1}\cdots p_r^{e_r}$,
gdzie $u\in k^\times$, $p_i$ są parami
różnymi monicznymi wielomianami nierozkładalnymi,
a $e_i\ge1$. Wówczas
\[
 (f)=\bigcap_{i=1}^r(p_i^{e_i}),\qquad
 k[x]/(f)\cong\prod_{i=1}^r k[x]/(p_i^{e_i}).
\]
Każdy ideał $(p_i^{e_i})$ jest prymarny.
\end{twierdzenie}
\begin{proof}
Ponieważ różne wielomiany nierozkładalne $p_i,p_j$
są względnie pierwsze, również ich potęgi są
względnie pierwsze. Z algorytmu Euklidesa
istnieją wielomiany $a,b$ takie, że
$ap_i^{e_i}+bp_j^{e_j}=1$.
Ideały $(p_i^{e_i})$ są więc parami współmaksymalne.
Dla dwóch współmaksymalnych ideałów $I,J$
zachodzi $I\cap J=IJ$: jeśli $z\in I\cap J$
i $1=i+j$ z $i\in I$, $j\in J$, to
$z=zi+zj\in IJ$. Indukcja daje
$\bigcap_i(p_i^{e_i})=\prod_i(p_i^{e_i})=(f)$,
a chińskie twierdzenie o resztach daje izomorfizm
pierścieni ilorazowych.

Jeśli $p_i^{e_i}\mid ab$, to suma wykładników
$p_i$ w rozkładzie $a,b$ jest co najmniej $e_i$.
Jeśli $p_i^{e_i}\nmid a$, to $b$ musi zawierać
co najmniej jedną kopię czynnika $p_i$.
Wybierając odpowiednio wysoką potęgę $b$,
uzyskujemy $p_i^{e_i}\mid b^n$.
To dowodzi prymarności.
\end{proof}

\begin{uwaga}[Czego nie wolno wywnioskować z przykładu]
Rozkład ideału na przecięcie prymarnych w ogólnym
pierścieniu noetherowskim (twierdzenie Laskera--Noether)
jest znacznie trudniejszy niż przypadek $k[x]$:
nie wystarczy rozłożyć \emph{jeden wielomian}.
Wielowymiarowe odpowiedniki wymagają
algebry modułów, ideałów skojarzonych i argumentu
noetherowskiego.
\end{uwaga}

\subsection{Normalizacja Noether: dlaczego skończenie wiele równań ma skończony wymiar?}
\begin{twierdzenie}[Normalizacja Noether, wersja dla ciał nieskończonych]
\label{tw:noether-norm}
Niech $k$ będzie ciałem nieskończonym, a $A$
niezerową, skończenie generowaną $k$-algebrą,
która jest dziedziną. Istnieją elementy
$y_1,\ldots,y_d\in A$, algebraicznie niezależne
nad $k$, takie że $A$ jest skończenie generowanym
modułem nad podpierścieniem
$k[y_1,\ldots,y_d]\cong k[T_1,\ldots,T_d]$.
\end{twierdzenie}
\begin{proof}
Niech $A=k[a_1,\ldots,a_n]$. Prowadzimy indukcję
względem liczby generatorów $n$. Jeśli
$a_1,\ldots,a_n$ są algebraicznie niezależne,
przyjmujemy $d=n$ i te elementy jako $y_i$;
$A$ jest wtedy wolnym modułem rzędu jeden nad sobą.

W przeciwnym razie istnieje niezerowy wielomian
$F\in k[T_1,\ldots,T_n]$ spełniający
$F(a_1,\ldots,a_n)=0$. Niech $F_m$ oznacza
jego niezerową część jednorodną najwyższego stopnia.
Ponieważ $k$ jest nieskończone, możemy wybrać
$c_1,\ldots,c_{n-1}\in k$ tak, aby
$F_m(c_1,\ldots,c_{n-1},1)\ne0$.
Uzasadnijmy ten krok: wielomian
$F_m(T_1,\ldots,T_{n-1},1)$ nie jest zerowy,
bo w przeciwnym razie jednorodność $F_m$
wymuszałaby $F_m=0$; niezerowy wielomian
nad ciałem nieskończonym nie znika identycznie
na wszystkich jego punktach.

Wprowadzamy nowe generatory
$b_i=a_i-c_i a_n$ dla $i<n$ oraz $b_n=a_n$.
Są to współrzędne odwracalnej liniowej zmiany
zmiennych, więc nadal $A=k[b_1,\ldots,b_n]$.
Po podstawieniu $a_i=b_i+c_i b_n$ do równania
$F=0$ otrzymujemy równanie wielomianowe
na $b_n$, którego współczynnik najwyższej potęgi
$b_n^m$ jest niezerową \emph{stałą} z $k$.
Po podzieleniu przez tę stałą jest to równanie
moniczne o współczynnikach w
$B=k[b_1,\ldots,b_{n-1}]$. Zatem $b_n$
jest całkowity nad $B$, a $A=B[b_n]$ jest
skończenie generowanym $B$-modułem na mocy
twierdzenia~\ref{tw:integrality-module}.

Pierścień $B$ jest dziedziną i ma nie więcej
niż $n-1$ generatorów, więc z założenia indukcyjnego
jest skończony nad pewnym pierścieniem wielomianów
$k[y_1,\ldots,y_d]$ o algebraicznie niezależnych
generatorach. Jeśli $e_1,\ldots,e_r$ generują
$A$ nad $B$, a $b_1',\ldots,b_s'$ generują $B$
nad $k[y_1,\ldots,y_d]$, to iloczyny $e_i b_j'$
generują $A$ nad $k[y_1,\ldots,y_d]$.
To kończy indukcję.
\end{proof}

\begin{uwaga}[Co dokładnie wynika z normalizacji?]
Geometrycznie skończenie generowana dziedzina
$A$ jest skończona nad pierścieniem współrzędnych
przestrzeni afinicznej pewnego wymiaru $d$.
W szczególności $\operatorname{trdeg}_k\Frac(A)=d$:
każdy element $\Frac(A)$ jest algebraiczny nad
$k(y_1,\ldots,y_d)$, a $y_i$ są algebraicznie
niezależne. Równość $d=\dim A$ (wymiar Krulla)
wymaga ponadto twierdzeń o łańcuchach ideałów
pierwszych w rozszerzeniach całkowitych
(własności \emph{going up} i \emph{incomparability});
nie wynika ona formalnie wyłącznie z definicji
normalizacji Noether.
\end{uwaga}


\begin{definicja}[Wymiar Krulla]
Niech $A$ będzie pierścieniem przemiennym.
Jego \emph{wymiarem Krulla} jest supremum długości
$n$ łańcuchów ściśle rosnących ideałów pierwszych
$\mathfrak p_0\subsetneq\mathfrak p_1\subsetneq
\cdots\subsetneq\mathfrak p_n$. Długością jest
liczba zawierań, czyli $n$, nie liczba ideałów.
Dla pierścienia zerowego przyjmujemy $\dim A=-\infty$.
\end{definicja}

\begin{przyklad}[Punkty, krzywe i powierzchnie]
Dla ciała $k$ jedynym ideałem pierwszym jest $(0)$,
więc $\dim k=0$. W $k[x]$ mamy łańcuch
$(0)\subsetneq(x-a)$ dla każdego $a\in k$.
Ponieważ $k[x]$ jest PID, każdy niezerowy ideał
pierwszy jest maksymalny, więc $\dim k[x]=1$.
W $k[x,y]$ znajdujemy
$(0)\subsetneq(x)\subsetneq(x,y)$, zatem
$\dim k[x,y]\ge2$; równość $2$ jest szczególnym
przypadkiem twierdzenia $\dim k[x_1,\ldots,x_n]=n$.
Dowód pełnego twierdzenia wymaga dodatkowego
argumentu, np. normalizacji Noether, i nie wynika
wyłącznie z wypisania jednego łańcucha.
\end{przyklad}

\subsection{Rozszerzenia całkowite a łańcuchy ideałów pierwszych}
\label{subsec:going-up}
Dlaczego geometryczny wymiar mierzymy długością łańcuchów
ideałów pierwszych, a nie liczbą współrzędnych w
zapisanym równaniu? Kluczowa jest niezmienniczość
tego wymiaru przy przejściu do rozszerzenia
całkowitego. Zanim ją wykażemy, potrzebujemy
trzech faktów. Wszystkie rozpatrywane pierścienie
są przemienne z jedynką, a rozszerzenia $A\subseteq B$
są injektywnymi homomorfizmami zachowującymi jedynkę.

\begin{lemat}[Całkowitość nad ciałem]
\label{lem:integral-field}
Niech $R\subseteq L$ będzie rozszerzeniem dziedzin,
przy czym $L$ jest ciałem i każdy element $L$ jest
całkowity nad $R$. Wtedy $R$ jest ciałem.
\end{lemat}
\begin{proof}
Niech $0\ne a\in R$. Ponieważ $a^{-1}\in L$ jest
całkowity nad $R$, istnieje równanie
$(a^{-1})^n+c_{n-1}(a^{-1})^{n-1}+\cdots+c_0=0$
z $c_i\in R$. Pomnożenie przez $a^{n-1}$
daje $a^{-1}=-c_{n-1}-c_{n-2}a-\cdots-c_0a^{n-1}\in R$.
Zatem każdy niezerowy element $R$ ma odwrotność.
\end{proof}

\begin{twierdzenie}[Istnienie ideału leżącego nad danym: lying over]
\label{tw:lying-over}
Niech $A\subseteq B$ będzie rozszerzeniem całkowitym
dziedzin, a $\mathfrak p\subset A$ ideałem pierwszym.
Istnieje ideał pierwszy $\mathfrak q\subset B$ taki,
że $\mathfrak q\cap A=\mathfrak p$.
\end{twierdzenie}
\begin{proof}
Niech $S=A\setminus\mathfrak p$. Lokalizacja
$S^{-1}B$ jest niezerowa, gdyż $B$ jest dziedziną
i $S$ nie zawiera zera. Lokalizacja zachowuje
zależność całkowitą: jeśli $b\in B$ spełnia
równanie moniczne nad $A$, to $b/s$ spełnia
równanie moniczne nad $S^{-1}A$ po podzieleniu
kolejnych współczynników przez właściwe potęgi $s$.
Wobec tego $S^{-1}B$ jest całkowity nad
$A_{\mathfrak p}=S^{-1}A$.

Wybierzmy ideał maksymalny $\mathfrak n$ w
$S^{-1}B$ (jego istnienie wynika z lematu Zorna).
Kontrakcja $\mathfrak n\cap A_{\mathfrak p}$
jest maksymalna: iloraz
$(S^{-1}B)/\mathfrak n$ jest ciałem i jest
całkowity nad dziedziną
$A_{\mathfrak p}/(\mathfrak n\cap A_{\mathfrak p})$,
więc lemat~\ref{lem:integral-field} pokazuje,
że ta dziedzina jest ciałem. Ponieważ
$A_{\mathfrak p}$ jest lokalny z ideałem maksymalnym
$\mathfrak p A_{\mathfrak p}$, otrzymujemy
$\mathfrak n\cap A_{\mathfrak p}
=\mathfrak p A_{\mathfrak p}$.
Przeciwobraz $\mathfrak n$ przy odwzorowaniu
$B\to S^{-1}B$ jest ideałem pierwszym
$\mathfrak q$, a jego kontrakcja do $A$
wynosi $\mathfrak p$.
\end{proof}

\begin{twierdzenie}[Nieporównywalność: incomparability]
\label{tw:incomparability}
Niech $A\subseteq B$ będzie rozszerzeniem całkowitym
dziedzin. Jeżeli ideały pierwsze $\mathfrak q_1
\subseteq\mathfrak q_2$ w $B$ mają jednakową
kontrakcję $\mathfrak p=\mathfrak q_i\cap A$,
to $\mathfrak q_1=\mathfrak q_2$.
\end{twierdzenie}
\begin{proof}
Przechodzimy do rozszerzenia całkowitego
dziedzin $A/\mathfrak p\subseteq B/\mathfrak q_1$.
Obraz $\mathfrak q_2$ jest ideałem pierwszym,
którego kontrakcją w $A/\mathfrak p$ jest zero.
Gdyby ten obraz zawierał niezerowy element $b$,
to wybierzmy moniczne równanie całkowite
$b^n+a_{n-1}b^{n-1}+\cdots+a_0=0$ o
najmniejszym możliwym stopniu $n$. Jego wyraz
wolny $a_0$ jest niezerowy: gdyby był zerowy,
można by podzielić równanie przez $b\ne0$ w
dziedzinie, otrzymując krótsze równanie moniczne.
Ponieważ $b$ należy do rozważanego ideału,
z równania wynika, że $a_0$ też należy do
tego ideału; jest więc zerem w $A/\mathfrak p$,
co przeczy $a_0\ne0$. Zatem obraz
$\mathfrak q_2$ jest zerowy i oba ideały są równe.
\end{proof}

\begin{twierdzenie}[Przedłużanie łańcuchów: going up]
\label{tw:going-up}
Niech $A\subseteq B$ będzie rozszerzeniem całkowitym
dziedzin. Jeśli $\mathfrak p\subseteq\mathfrak p'$
są ideałami pierwszymi w $A$ oraz $\mathfrak q$
jest ideałem pierwszym w $B$ o kontrakcji
$\mathfrak p$, to istnieje ideał pierwszy
$\mathfrak q'\supseteq\mathfrak q$ w $B$ o
kontrakcji $\mathfrak p'$.
\end{twierdzenie}
\begin{proof}
Rozszerzenie $A/\mathfrak p\subseteq B/\mathfrak q$
jest całkowite. Ideał $\mathfrak p'/\mathfrak p$
jest pierwszy w $A/\mathfrak p$. Z twierdzenia
\ref{tw:lying-over} istnieje ideał pierwszy w
$B/\mathfrak q$ leżący nad nim. Jego przeciwobraz
w $B$ jest poszukiwanym ideałem $\mathfrak q'$.
\end{proof}

\begin{wniosek}[Zachowanie wymiaru Krulla]
\label{wn:dim-integral}
Niech $A\subseteq B$ będzie rozszerzeniem całkowitym
dziedzin. Wówczas $\dim A=\dim B$.
\end{wniosek}
\begin{proof}
Każdy skończony ściśle rosnący łańcuch
ideałów pierwszych w $A$ można unieść do
łańcucha w $B$: zaczynamy od ideału danego
przez twierdzenie~\ref{tw:lying-over}, a kolejne
wybieramy przez twierdzenie~\ref{tw:going-up}.
Ponieważ kontrakcje kolejnych ideałów są różne,
uniesiony łańcuch też jest ściśle rosnący.
Stąd $\dim B\ge\dim A$.

Z drugiej strony kontrakcje każdego ściśle
rosnącego łańcucha ideałów pierwszych w $B$
tworzą ściśle rosnący łańcuch w $A$: gdyby dwie
kolejne kontrakcje były równe, twierdzenie
\ref{tw:incomparability} wymusiłoby równość
samych ideałów. Zatem $\dim B\le\dim A$.
\end{proof}

\begin{twierdzenie}[Wymiar a liczba niezależnych parametrów]
\label{tw:dim-trdeg}
Niech $k$ będzie ciałem nieskończonym, a $A$
skończenie generowaną $k$-algebrą, która jest dziedziną.
Wówczas
$\dim A=\operatorname{trdeg}_k\operatorname{Frac}(A)$.
W szczególności dla każdej liczby całkowitej $n\ge0$
zachodzi $\dim k[x_1,\ldots,x_n]=n$.
\end{twierdzenie}
\begin{proof}
Pokażemy najpierw, że jeśli $R$ jest taką dziedziną,
$d=\operatorname{trdeg}_k\operatorname{Frac}(R)$,
a $\mathfrak q\subsetneq R$ jest \emph{niezerowym}
ideałem pierwszym, to
$\operatorname{trdeg}_k\operatorname{Frac}(R/\mathfrak q)<d$.
Przez normalizację Noether wybierzmy skończone
rozszerzenie całkowite
$B=k[y_1,\ldots,y_d]\subseteq R$.
Ponieważ $R$ jest dziedziną, ideał $(0)$ jest
pierwszy. Kontrakcja $\mathfrak q\cap B$ nie może
być zerowa: w przeciwnym razie porównywalne ideały
$(0)\subsetneq\mathfrak q$ leżałyby nad $(0)$ w $B$,
wbrew twierdzeniu o nieporównywalności
\ref{tw:incomparability}.
Zatem istnieje niezerowy wielomian
$h(y_1,\ldots,y_d)\in\mathfrak q\cap B$.
Obrazy $y_i$ w $B/(\mathfrak q\cap B)$ są więc
algebraicznie zależne, a ich ciało ułamków ma
stopień przestępny co najwyżej $d-1$ nad $k$.
Ponieważ $R/\mathfrak q$ jest całkowity nad
$B/(\mathfrak q\cap B)$, ich ciała ułamków tworzą
rozszerzenie algebraiczne, co nie zmienia stopnia
przestępnego. Otrzymujemy żądaną nierówność.

Rozważmy teraz dowolny ściśle rosnący łańcuch
$\mathfrak p_0\subsetneq\cdots\subsetneq
\mathfrak p_r$ ideałów pierwszych w $A$.
Dla każdego $i<r$ przechodzimy do skończenie
generowanej dziedziny $A/\mathfrak p_i$ i jej
niezerowego ideału pierwszego
$\mathfrak p_{i+1}/\mathfrak p_i$. Poprzedni
argument pokazuje, że stopień przestępny ciała
ułamków ściśle maleje wzdłuż łańcucha.
Zaczynamy od liczby nie większej niż $d$ i kończymy
na liczbie nieujemnej, zatem $r\le d$.
To daje $\dim A\le d$.

Z normalizacji Noether $A$ jest całkowity nad
$B=k[y_1,\ldots,y_d]$; z wniosku
\ref{wn:dim-integral} mamy $\dim A=\dim B$.
Dla pierścienia wielomianów $B$ istnieje łańcuch
$(0)\subsetneq(y_1)\subsetneq(y_1,y_2)
\subsetneq\cdots\subsetneq(y_1,\ldots,y_d)$
długości $d$ (dla $d=0$ jest to jednoelementowy
łańcuch). Zatem $\dim B\ge d$, a udowodniona
nierówność daje $\dim B\le d$. W konsekwencji
$\dim A=\dim B=d$.
\end{proof}

\begin{przyklad}[Łańcuchy ideałów a obiekty geometryczne]
W $k[x,y]$ łańcuch
$(0)\subsetneq(x)\subsetneq(x,y)$
ma długość $2$. Odpowiada mu \emph{odwrócony}
łańcuch nierozkładalnych zbiorów algebraicznych
$\{(0,0)\}\subsetneq\{x=0\}\subsetneq\A^2_k$.
Wzrost ideału oznacza dokładanie równań i
zmniejszanie geometrycznego zbioru rozwiązań.
\end{przyklad}

\begin{przyklad}[Normalizacja nie zmienia wymiaru]
Dla ostrza $A=k[t^2,t^3]\subseteq B=k[t]$ rozszerzenie
jest całkowite, gdyż $B=A[t]$, a $t$ spełnia
$T^2-t^2=0$. Dlatego $\dim A=\dim B=1$.
Geometrycznie normalizacja rozplątuje
osobliwość, ale nie zmienia wymiaru krzywej.
\end{przyklad}

\begin{definicja}[Pierścień regularny w wymiarze jeden]
Niech $(A,\mathfrak m)$ będzie noetherowską lokalną
dziedziną wymiaru Krulla $1$, a $\kappa=A/\mathfrak m$
jej ciałem resztowym. Pierścień nazywamy
\emph{regularnym}, jeśli
$\dim_\kappa(\mathfrak m/\mathfrak m^2)=1$.
Liczbę $\dim_\kappa(\mathfrak m/\mathfrak m^2)$
nazywamy wymiarem przestrzeni kostycznej w punkcie.
\end{definicja}

\begin{lemat}[Jednowymiarowy pierścień regularny ma parametr]
\label{lem:param-local}
Jeśli $(A,\mathfrak m)$ jest noetherowską lokalną
dziedziną i $\dim_{A/\mathfrak m}
(\mathfrak m/\mathfrak m^2)=1$, to istnieje
element $t\in\mathfrak m$ taki, że $\mathfrak m=(t)$.
\end{lemat}
\begin{proof}
Ponieważ $A$ jest noetherowski, ideał $\mathfrak m$
jest skończenie generowany jako $A$-moduł.
Wniosek~\ref{wn:generatory-local} zastosowany do
$M=\mathfrak m$ mówi, że minimalna liczba jego
generatorów wynosi
$\dim_{A/\mathfrak m}(\mathfrak m/\mathfrak m^2)=1$.
Wybrany generator oznaczamy $t$.
\end{proof}

\begin{twierdzenie}[Od parametru do wartościowania]
\label{tw:dvr}
Niech $(A,\mathfrak m)$ będzie noetherowską lokalną
dziedziną, $\mathfrak m=(t)\ne(0)$.
Wówczas każdy niezerowy element $a\in A$ ma
jednoznaczną postać $a=t^n u$, gdzie $n\in\Z_{\ge0}$
i $u\in A^\times$. Każdy niezerowy ideał $A$
jest potęgą $(t^n)$. Pierścień taki nazywamy
\emph{pierścieniem z wartościowaniem dyskretnym}
(DVR), a funkcję $\operatorname{ord}_A(a)=n$
rozszerzamy na $\Frac(A)^\times$ przez
$\operatorname{ord}_A(a/b)=\operatorname{ord}_A(a)-
\operatorname{ord}_A(b)$.
\end{twierdzenie}
\begin{proof}
Najpierw pokażemy, że $\bigcap_{n\ge0}(t^n)=0$.
Załóżmy przeciwnie, że $a\ne0$ należy do wszystkich
$(t^n)$. Dla każdego $n$ istnieje $a_n\in A$ z
$a=t^na_n$. Ponieważ $A$ jest dziedziną i $t\ne0$,
$a_n\ne0$ i $a_n=t a_{n+1}$. Mamy więc rosnący
łańcuch ideałów $(a_n)\subseteq(a_{n+1})$.
Noetherowskość wymusza, że od pewnego miejsca
$(a_n)=(a_{n+1})$. Istnieje wtedy $r\in A$
takie, że $a_{n+1}=r a_n=r t a_{n+1}$.
Przez skrócenie niezerowego $a_{n+1}$ otrzymujemy
$1=rt$, sprzecznie z $t\in\mathfrak m$.

Każdy niezerowy $a$ należy do $(t^0)=A$, ale nie
do wszystkich $(t^n)$; istnieje największe $n$
takie, że $a\in(t^n)$. Zatem $a=t^n u$,
przy czym $u\notin(t)=\mathfrak m$, bo w przeciwnym
razie $a\in(t^{n+1})$. Element spoza jedynego
ideału maksymalnego jest jednostką, więc $u$ jest
odwracalny. Jeśli $t^nu=t^mv$ i $n<m$, otrzymalibyśmy
$u=t^{m-n}v\in\mathfrak m$, sprzeczność;
analogicznie dla $m<n$. To daje jednoznaczność $n$.

Niech $I\ne0$ będzie ideałem. Zbiór wartościowań
niezerowych elementów $I$ jest niepustym podzbiorem
$\Z_{\ge0}$, ma więc minimum $n$. Wybieramy
$a=t^nu\in I$ z $u$ odwracalnym; stąd $t^n\in I$.
Z drugiej strony każdy $b\in I$ ma wartościowanie
co najmniej $n$ i należy do $(t^n)$. Zatem
$I=(t^n)$. Równość wartościowania na ułamkach
nie zależy od ich przedstawienia, ponieważ z
$ab'=a'b$ wynika addytywność rzędów po obu stronach.
\end{proof}

\begin{wniosek}[Lokalny pierścień gładkiej krzywej]
\label{wn:krzywa-dvr}
Niech $k$ będzie algebraicznie domknięte, a $C$
gładką nierozkładalną krzywą algebraiczną nad $k$.
Dla każdego punktu zamkniętego $P\in C$ pierścień
lokalny $\mathcal O_{C,P}$ jest DVR.
\end{wniosek}
\begin{proof}
Z lokalnej teorii afinicznych zbiorów algebraicznych
mamy, że $\mathcal O_{C,P}$ jest noetherowską
lokalną dziedziną. Ponieważ $C$ ma wymiar jeden,
jej pierścień lokalny w punkcie zamkniętym ma
wymiar Krulla jeden. Gładkość w $P$ oznacza,
że przestrzeń styczna, dualna do
$\mathfrak m_P/\mathfrak m_P^2$, ma wymiar jeden.
Zatem $\mathfrak m_P$ jest główny na mocy
lematu~\ref{lem:param-local}. Twierdzenie~\ref{tw:dvr}
kończy dowód. Przywołane utożsamienie wymiaru
przestrzeni stycznej z wymiarem kostycznym
jest kryterium Jacobiego z rozdziału~\ref{ch:osobliwosci}.
\end{proof}

\begin{uwaga}[Zakres wyniku]
Twierdzenie~\ref{tw:dvr} jest całkowicie algebraiczne:
nie wymaga charakterystyki zero. Dla krzywej z ostrzem
lokalny pierścień w osobliwości nie jest DVR;
normalizacja z przykładu~\ref{prz:normalizacja-ostrza}
przywraca lokalny parametr $t$.
\end{uwaga}

\section{Przedsnopy i snopy: dane lokalne, które można sklejać}
\label{sec:snopy}

\subsection{Punkt wyjścia: funkcja lokalna nie musi być globalna}
Rozważmy okrąg $S^1$ i funkcję ``kąt'' mierzoną w radianach.
Na małych łukach można wybrać gładką gałąź kąta, lecz
nie istnieje globalna ciągła funkcja $\theta:S^1\to\R$
przyjmująca kąty bez skoków i spełniająca
$e^{i\theta(z)}=z$ dla każdego $z\in S^1$.
Problem nie leży w złych wzorach lokalnych: polega na tym,
że wzory na nakładających się łukach nie zawsze
uzgadniają się jako funkcje rzeczywiste. Snop formalizuje
tę różnicę między \emph{istnieniem lokalnym} a
\emph{sklejalnością globalną}.

\begin{definicja}[Przedsnop zbiorów]
Niech $X$ będzie przestrzenią topologiczną.
\emph{Przedsnopem} $\mathcal F$ zbiorów na $X$ nazywamy
przyporządkowanie każdemu otwartemu $U\subseteq X$
zbioru $\mathcal F(U)$ (jego elementy to \emph{sekcje}),
a każdemu zawieraniu otwartych $V\subseteq U$
odwzorowania ograniczenia
$\rho^U_V:\mathcal F(U)\to\mathcal F(V)$, tak aby
$\rho^U_U=\id_{\mathcal F(U)}$ i
$\rho^V_W\circ\rho^U_V=\rho^U_W$ dla $W\subseteq V\subseteq U$.
Przedsnopem grup, pierścieni lub $A$-modułów nazywamy
analogiczną konstrukcję, w której sekcje i ograniczenia
mają odpowiednią strukturę algebraiczną.
\end{definicja}

\begin{definicja}[Snop]
\label{def:snop}
Przedsnop $\mathcal F$ na $X$ jest \emph{snopem},
jeśli dla każdego otwartego $U\subseteq X$ i każdego
pokrycia otwartego $U=\bigcup_{i\in I}U_i$ zachodzą:
\begin{enumerate}[label=(S\arabic*)]
\item \emph{Jednoznaczność:} jeśli $s,t\in\mathcal F(U)$
i $s|_{U_i}=t|_{U_i}$ dla każdego $i$, to $s=t$.
\item \emph{Sklejanie:} jeśli dane są $s_i\in\mathcal F(U_i)$
takie, że $s_i|_{U_i\cap U_j}=s_j|_{U_i\cap U_j}$
dla wszystkich $i,j$, to istnieje $s\in\mathcal F(U)$
spełniające $s|_{U_i}=s_i$ dla wszystkich $i$.
\end{enumerate}
Przez konwencję snop zbiorów ma jedną sekcję nad
zbiorem pustym, gdy pokrycie jest puste.
\end{definicja}

\begin{center}
\begin{tikzpicture}[scale=.91,>=Stealth]
% Zbior U jest suma ponizszych elips, nie wieksza elipsa.
\draw[fill=green!12,draw=green!55!black,thick] (-.9,.1) ellipse (1.75 and 1.03);
\draw[fill=orange!14,draw=orange!70!black,thick] (.9,.08) ellipse (1.75 and 1.03);
\node at (-1.95,.5) {$U_1$}; \node at (1.98,.5) {$U_2$};
\node at (0,-.12) {$U_1\cap U_2$};
\node[fill=white,inner sep=1pt] at (0,-1.24) {$U=U_1\cup U_2$};
\node[align=center,font=\small,anchor=north] at (0,-1.95)
{$s_1|_{U_1\cap U_2}=s_2|_{U_1\cap U_2}$\\[2pt]
 $\Longrightarrow$ istnieje dokładnie jedna sekcja $s\in\mathcal F(U)$};
\end{tikzpicture}

{\small Rysunek 4.2. Zgodne sekcje na pokryciu $U=U_1\cup U_2$ sklejamy jednoznacznie do sekcji na $U$.}\par
\end{center}

\begin{przyklad}[Podstawowe snopy]
Niech $X$ będzie przestrzenią topologiczną.
\begin{enumerate}[label=\alph*)]
\item $\mathcal C_X(U)$, zbiór funkcji ciągłych $U\to\R$,
tworzy snop pierścieni z ograniczeniami funkcji.
Jeśli ciągłe funkcje zgadzają się na przecięciach,
wówczas funkcja sklejona jest ciągła, bo ciągłość
jest własnością lokalną.
\item Na rozmaitości gładkiej $M$ przyporządkowanie
$U\mapsto C^\infty(U)$ jest snopem.
\item Na rozmaitości zespolonej $X$ funkcje holomorficzne
tworzą snop $\mathcal O_X$.
\item Na otwartej przestrzeni $U\subseteq\R^n$ przyporządkowanie
$V\mapsto\{f:V\to\R\ \text{lokalnie stała}\}$
tworzy \emph{snop lokalnie stałych funkcji}.
\end{enumerate}
\end{przyklad}

\begin{przyklad}[Przedsnop, który nie jest snopem]
\label{prz:presheaf-not-sheaf}
Niech $X=\{a,b\}$ z topologią dyskretną i ustalmy
zbiór $E=\{0,1\}$. Każdemu \emph{niepustemu}
otwartemu $U\subseteq X$ przypiszmy ten sam zbiór $E$,
a każdemu zawieraniu niepustych otwartych przypiszmy
identyczność; nad pustym zbiorem bierzemy singleton.
To przedsnop. Dla pokrycia $X=\{a\}\cup\{b\}$
sekcje $s_a=0$ i $s_b=1$ są zgodne na przecięciu,
bo jest ono puste. Nie można ich jednak skleić:
każda sekcja nad $X$ jest jednym elementem $E$,
więc ogranicza się do \emph{tej samej} wartości
na obu częściach. Snop skojarzony z tym przedsnopem
to snop lokalnie stałych funkcji $U\to E$.
\end{przyklad}

\subsection{Źdźbła i kiełki sekcji}
\begin{definicja}[Źdźbło]
\label{def:stalk}
Niech $\mathcal F$ będzie przedsnopem na $X$,
a $x\in X$. \emph{Źdźbłem} $\mathcal F_x$
nazywamy zbiór klas par $(U,s)$, gdzie $x\in U$
i $s\in\mathcal F(U)$, przy relacji
$(U,s)\sim(V,t)$ wtedy i tylko wtedy, gdy istnieje
otwarte $W$ z $x\in W\subseteq U\cap V$
takie, że $s|_W=t|_W$. Klasę oznaczamy $s_x$
i nazywamy \emph{kiełkiem} $s$ w $x$.
\end{definicja}

\begin{stwierdzenie}
Niech $X$ będzie przestrzenią topologiczną, $x\in X$,
a $\mathcal F$ przedsnopem zbiorów na $X$. Relacja na
parach $(U,s)$, gdzie $x\in U$ i $s\in\mathcal F(U)$,
określona warunkiem zgodności sekcji na pewnym otoczeniu
$x$ zawartym w $U\cap V$, jest relacją równoważności.
Jeśli $\mathcal F$ jest przedsnopem grup, pierścieni albo
modułów nad ustalonym pierścieniem, to zbiór klas
$\mathcal F_x$ dziedziczy odpowiednią strukturę algebraiczną.
\end{stwierdzenie}
\begin{proof}
Zwrotność i symetria są oczywiste. Jeśli
$s|_W=t|_W$ oraz $t|_{W'}=u|_{W'}$, to
na otwartym $W\cap W'$ (zawierającym $x$)
mamy $s=u$. Działania na kiełkach definiujemy
przez ograniczenie obu przedstawicieli do
wspólnego otoczenia $x$ i wykonanie działania tam.
Zmiana przedstawicieli nie wpływa na wynik,
bo zgodność lokalna zachowuje się przy działaniach.
\end{proof}

\begin{przyklad}[Dwa różne wzory, ten sam kiełek]
Niech $f:\R\to\R$ będzie dowolną funkcją
identycznie równą zero w przedziale $(-1,1)$,
ale niezerową np. dla $x>2$. W snopie funkcji
ciągłych (jeśli $f$ jest ciągła) jej kiełek
$f_0$ jest równy kiełkowi funkcji zerowej,
chociaż $f$ nie jest globalnie funkcją zerową.
Źdźbło zapomina wszystko poza dowolnie małym
otoczeniem wybranego punktu.
\end{przyklad}

\begin{twierdzenie}[Równość kiełków wykrywa równość sekcji]
\label{tw:germ-zero}
Niech $\mathcal F$ będzie snopem zbiorów na $X$,
a $U\subseteq X$ otwartym. Dla sekcji
$s,t\in\mathcal F(U)$ zachodzi $s=t$ wtedy
i tylko wtedy, gdy $s_x=t_x$ dla każdego $x\in U$.
\end{twierdzenie}
\begin{proof}
Kierunek ``tylko wtedy'' wynika z definicji.
Jeśli dla każdego $x$ kiełki są równe,
istnieje otwarte otoczenie $W_x\subseteq U$, na
którym $s|_{W_x}=t|_{W_x}$. Zbiory $W_x$ pokrywają $U$.
Warunek jednoznaczności snopa daje $s=t$.
\end{proof}

\section{Spektrum pierścienia i jego snop strukturalny}
\label{sec:spek-snop}

\begin{definicja}[Spektrum]
Niech $A$ będzie pierścieniem przemiennym z jedynką.
$\Spec A$ to zbiór wszystkich ideałów pierwszych
$\mathfrak p\subsetneq A$. Dla ideału $I\subseteq A$
niech $V(I)=\{\mathfrak p\in\Spec A:I\subseteq\mathfrak p\}$.
Zbiory $V(I)$ definiują zbiory domknięte
\emph{topologii Zariskiego} na $\Spec A$.
Zbiory $D(f)=\{\mathfrak p\in\Spec A:f\notin\mathfrak p\}$
stanowią bazę jej zbiorów otwartych.
\end{definicja}

\begin{stwierdzenie}[Rachunek zbiorów podstawowo otwartych]
Dla $f,g\in A$ zachodzi
$D(f)\cap D(g)=D(fg)$. Ponadto
$D(f)\subseteq\bigcup_{i\in I}D(g_i)$ wtedy
i tylko wtedy, gdy $f\in\sqrt{(g_i:i\in I)}$.
W szczególności każdy $D(f)$ jest quasi-zwarty:
każde jego pokrycie otwarte ma podpokrycie skończone.
\end{stwierdzenie}
\begin{proof}
Pierwsza równość wynika z pierwszości
$\mathfrak p$: $fg\notin\mathfrak p$ wtedy
i tylko wtedy, gdy $f,g\notin\mathfrak p$.
Inkluzja w tezie jest równoważna temu, że
każdy ideał pierwszy zawierający wszystkie $g_i$
zawiera również $f$. Przecięcie ideałów pierwszych
zawierających ideał $J$ jest równe $\sqrt J$;
aby to udowodnić, dla $f\notin\sqrt J$ lokalizujemy
$A/J$ w potęgach obrazu $f$ (lokalizacja jest
niezerowa), wybieramy w niej ideał maksymalny
i cofamy go do ideału pierwszego $A$,
który zawiera $J$, lecz nie $f$.

Jeżeli $f^N\in(g_i:i\in I)$, to z definicji
ideału generowanego wynika, że $f^N$ jest
skończoną sumą $\sum_{j=1}^r a_j g_{i_j}$.
Ta skończona podrodzina nadal pokrywa $D(f)$.
Ponieważ zbiory $D(g)$ tworzą bazę, ten argument
wystarcza również dla dowolnego otwartego pokrycia.
\end{proof}

\begin{uwaga}[Nowe punkty w porównaniu z $k^n$]
W $\Spec k[x]$ dla algebraicznie domkniętego $k$
są nie tylko punkty $(x-a)$, $a\in k$, lecz także
punkt $(0)$. Punkt $(0)$ jest \emph{generyczny}:
jego domknięciem jest całe $\Spec k[x]$.
To nie jest kolejna liczba z $k$; reprezentuje
całą nierozkładalną prostą. Z drugiej strony
$\Spec k[\varepsilon]/(\varepsilon^2)$ ma jeden
punkt, ale nie jest tym samym schematem co $\Spec k$:
w pierwszym pierścieniu istnieje niezerowy nilpotent.
Snop strukturalny zachowuje tę informację.
\end{uwaga}

\begin{definicja}[Snop strukturalny na bazie]
\label{def:O-Spec}
Dla $X=\Spec A$ każdemu otwartemu podstawowemu
$D(f)$ przypisujemy pierścień
$\mathcal O_X(D(f)):=A_f$.
Jeżeli $D(g)\subseteq D(f)$, odpowiadające
odwzorowanie ograniczenia jest jedynym homomorfizmem
z $A_f$ do $A_g$ przedłużającym $A\to A_g$:
istnieje on, bo $f$ staje się odwracalny w $A_g$
(istotnie $g\in\sqrt{(f)}$, czyli $g^N=af$
dla pewnego $a\in A$, a więc $f$ jest jednostką
w $A_g$). Dla dowolnego otwartego $U$ sekcją
$\mathcal O_X(U)$ nazywamy rodzinę ułamków
na pewnym pokryciu $U$ zbiorami $D(f_i)$,
zgodną na przecięciach, przy utożsamieniu
rodzin mających wspólne lokalne ograniczenia.
\end{definicja}

\begin{twierdzenie}[Dlaczego ta konstrukcja jest snopem]
\label{tw:structural-sheaf}
Niech $A$ będzie pierścieniem przemiennym z jedynką,
$X=\Spec A$, a $D(f)=\{\mathfrak p\in X:f\notin\mathfrak p\}$.
Przyporządkowanie $D(f)\mapsto A_f$ z naturalnymi
odwzorowaniami ograniczenia wyznacza snop pierścieni
$\mathcal O_X$ na $X$.
Dla każdego $f\in A$ naturalne odwzorowanie
$A_f\to\mathcal O_X(D(f))$ jest izomorfizmem,
a dla każdego ideału pierwszego $\mathfrak p\subset A$
zachodzi kanoniczny izomorfizm pierścieni lokalnych
$(\mathcal O_X)_{\mathfrak p}\cong A_{\mathfrak p}$.
\end{twierdzenie}
\begin{proof}
Warunek sklejania i jednoznaczności dla \emph{rodzin
lokalnych}, których użyliśmy do zdefiniowania
$\mathcal O_X(U)$, wynika bezpośrednio z możliwości
zastąpienia dwóch pokryć ich wspólnym rozdrobnieniem.
Pozostaje wykazać, że nie dodaliśmy nowych sekcji
nad samym $D(f)$.

Zastępujemy $A$ przez $A_f$; wtedy $D(f)$ staje się
całym jego spektrum. Pokrycie przez podstawowe zbiory
$D(g_1),\ldots,D(g_r)$ jest skończone po wybraniu
podpokrycia dzięki quasi-zwartości. Ponieważ zbiory
pokrywają $\Spec A$, ideał $(g_1,\ldots,g_r)$
jest równy $A$ (gdyby był właściwy, leżałby w ideale
maksymalnym, dając niepokryty punkt).

Niech zgodne sekcje na tych zbiorach mają postać
$x_i=m_i/g_i^{N_i}\in A_{g_i}$.
Wybieramy wspólny dostatecznie duży wykładnik $N$
i zastępujemy $m_i$ przez odpowiednie iloczyny
z $g_i$, tak by wszystkie sekcje były zapisane
$m_i/g_i^N$. Zgodność na $D(g_ig_j)$ mówi,
że po pomnożeniu przez dostatecznie wysoką
potęgę $g_ig_j$ zachodzi równość
$g_j^N m_i=g_i^N m_j$ w $A$; zwiększając wspólnie
$N$ i dostosowując liczniki można ją osiągnąć
jednocześnie dla wszystkich par $(i,j)$.
Ponieważ pierwiastnik ideału $(g_1^N,\ldots,g_r^N)$
jest równy $A$, ten ideał także jest całym $A$.
Istnieją zatem $c_1,\ldots,c_r\in A$ takie,
że $\sum_i c_i g_i^N=1$. Kładziemy
$a=\sum_i c_i m_i$. Wtedy dla każdego $j$
\[
 g_j^N a=\sum_i c_i g_j^N m_i
 =\sum_i c_i g_i^N m_j=m_j.
\]
A więc obraz $a$ w $A_{g_j}$ jest równy
$m_j/g_j^N$. To dowodzi istnienia globalnej sekcji
sklejającej. Jeśli $a\in A$ znika w każdym $A_{g_i}$,
po przemnożeniu przez wystarczająco wysokie potęgi
$g_i$ otrzymujemy $g_i^N a=0$ dla wszystkich $i$;
z równości $\sum_i c_i g_i^N=1$ wynika $a=0$.
To dowodzi jednoznaczności, a więc
$\mathcal O_X(D(f))=A_f$.

Na koniec otoczenia podstawowe $D(f)$ punktu
$\mathfrak p$ są dokładnie te, dla których
$f\notin\mathfrak p$. Kiełki sekcji, czyli
ułamki z tych $A_f$ modulo dalsze ograniczenia,
są więc dokładnie ułamkami w $A_{\mathfrak p}$.
\end{proof}

\begin{definicja}[Przestrzeń lokalnie pierścieniowana i schemat afiniczny]
\emph{Przestrzenią lokalnie pierścieniowaną} jest para
$(X,\mathcal O_X)$, w której $X$ jest przestrzenią
topologiczną, $\mathcal O_X$ snopem pierścieni,
a każde źdźbło $(\mathcal O_X)_x$ pierścieniem lokalnym.
Parę $(\Spec A,\mathcal O_{\Spec A})$ nazywamy
\emph{schematem afinicznym}. Schemat jest przestrzenią
lokalnie pierścieniowaną, którą można pokryć otwartymi
podprzestrzeniami izomorficznymi ze schematami afinicznymi.
\end{definicja}

\begin{uwaga}[Dlaczego w ogóle wspominamy o schematach?]
Nie rozwijamy tu całej teorii schematów. Konstrukcja
$\Spec A$ pomaga jednak zrozumieć, dlaczego snop
strukturalny jest czymś więcej niż ozdobnikiem:
$A=\mathcal O_{\Spec A}(\Spec A)$ jest pierścieniem
sekcji globalnych, zaś $A_{\mathfrak p}$ jest
pierścieniem \emph{kiełków} w konkretnym punkcie.
Dla klasycznej nierozkładalnej rozmaitości afinicznej
nad algebraicznie domkniętym $k$ można ograniczyć
tę strukturę do zwykłych punktów domkniętych.
\end{uwaga}

\section{Snopy modułów i wiązki liniowe}
\label{sec:wiazki}

\begin{definicja}[Snop modułów]
Na przestrzeni pierścieniowanej $(X,\mathcal O_X)$
\emph{snop $\mathcal O_X$-modułów} $\mathcal F$
to snop grup abelowych, w którym dla każdego
otwartego $U$ grupa $\mathcal F(U)$ jest modułem
nad $\mathcal O_X(U)$, a ograniczenia są zgodne
z mnożeniem przez funkcje. Snop jest \emph{lokalnie
wolny rzędu $r$}, gdy $X$ ma otwarte pokrycie
$U_i$ takie, że
$\mathcal F|_{U_i}\cong\mathcal O_{U_i}^{\,r}$.
Snop lokalnie wolny rzędu $1$ nazywamy
\emph{snopem odwracalnym} lub \emph{wiązką liniową}.
\end{definicja}

\begin{przyklad}[Snop modułu]
Niech $X=\Spec A$ i $M$ będzie $A$-modułem.
Na podstawowych otwartych definiujemy
$\widetilde M(D(f))=M_f$ i sklejamy, tak jak w
dowodzie twierdzenia~\ref{tw:structural-sheaf}.
Otrzymujemy snop $\mathcal O_X$-modułów, którego
źdźbło w $\mathfrak p$ jest równe $M_{\mathfrak p}$.
Dla $M=A$ odzyskujemy $\mathcal O_X$; dla
$M=A^r$ mamy $\widetilde M\cong\mathcal O_X^r$.
\end{przyklad}

\begin{przyklad}[Wiązka liniowa o niebanalnym sklejaniu]
Niech $X=\PP^1_k$ i oznaczmy
$U_0=\{[X_0:X_1]:X_0\ne0\}$,
$U_\infty=\{[X_0:X_1]:X_1\ne0\}$.
Współrzędną na $U_0$ jest $t=X_1/X_0$,
a na $U_\infty$ współrzędną $s=X_0/X_1=1/t$.
Przyjmijmy formalne bazy $e_0$ na $U_0$
i $e_\infty$ na $U_\infty$. Dla ustalonego
$n\in\Z$ sklejamy wolne moduły rzędu jeden
za pomocą wzoru $e_\infty=t^n e_0$
na $U_0\cap U_\infty$. Otrzymujemy wiązkę
liniową $\mathcal O_{\PP^1}(n)$.
Sekcja $h_0(t)e_0$ i sekcja $h_\infty(s)e_\infty$
zgadzają się wtedy i tylko wtedy, gdy
$h_0(t)=t^n h_\infty(t^{-1})$.
To, że lokalnie każda wiązka jest trywialna,
nie oznacza, że jej sekcje można skleić bez
uwzględnienia funkcji przejścia.
\end{przyklad}

\begin{center}
\begin{tikzpicture}[>=Stealth,every node/.style={font=\small}]
  \node[draw,rounded corners=3pt,fill=blue!6,minimum height=1cm,
     text width=2.6cm,align=center] (u) at (0,0)
     {$U_0\cong\A^1$\\baza $e_0$};
  \node[draw,rounded corners=3pt,fill=gray!5,minimum height=1cm,
     text width=3.2cm,align=center] (w) at (4.8,0)
     {$U_0\cap U_\infty\cong\mathbb G_m$\\$t\ne0$};
  \node[draw,rounded corners=3pt,fill=orange!8,minimum height=1cm,
     text width=2.7cm,align=center] (v) at (9.6,0)
     {$U_\infty\cong\A^1$\\baza $e_\infty$};
  \draw[->,thick] (u)--(w);
  \draw[<-,thick] (w)--(v);
  \node[below=10mm of w,font=\small,fill=white,inner sep=1pt]
     {$e_\infty=t^n e_0$ na $U_0\cap U_\infty$};
\end{tikzpicture}

{\small Rysunek 4.3. Wiązka $\mathcal O_{\PP^1}(n)$ powstaje ze sklejenia dwóch trywialnych wiązek funkcją przejścia na przecięciu map.}\par
\end{center}

\begin{definicja}[Grupa Picarda]
Dla przestrzeni pierścieniowanej $(X,\mathcal O_X)$
\emph{grupą Picarda} $\operatorname{Pic}(X)$
nazywamy zbiór klas izomorfizmu snopów odwracalnych
na $X$. Działanie to iloczyn tensorowy
$[\mathcal L]+[\mathcal M]=[\mathcal L\otimes_{\mathcal O_X}\mathcal M]$;
elementem neutralnym jest $[\mathcal O_X]$, a odwrotnym
$[\mathcal L^\vee]$, gdzie
$\mathcal L^\vee=\mathcal{H}om_{\mathcal O_X}
(\mathcal L,\mathcal O_X)$.
\end{definicja}

\begin{uwaga}[Iloczyn tensorowy: wystarczające przypomnienie]
Dla $A$-modułów $M,N$ iloczyn $M\otimes_A N$ powstaje
z wolnego modułu na symbolach $m\otimes n$ przez
narzucenie relacji addytywności w obu argumentach
i $am\otimes n=m\otimes an$. Jest uniwersalny
wśród odwzorowań $A$-dwuliniowych z $M\times N$.
Dla wolnych modułów rzędu jeden
$Ae\otimes_A Af\cong A(e\otimes f)$, dlatego
iloczyn tensorowy wiązek liniowych jest znów
lokalnie wolny rzędu jeden.
\end{uwaga}

\section{Krzywe algebraiczne: warunki i ciało funkcji}
\label{sec:curves-intro}

W dalszej części, \emph{o ile wyraźnie nie zaznaczono
inaczej}, $k$ jest ciałem algebraicznie domkniętym,
a $C$ \emph{gładką, nierozkładalną, rzutową krzywą}
nad $k$. Warunek rzutowości odpowiada geometrycznej
zwartości, a gładkość gwarantuje istnienie lokalnego
parametru w każdym punkcie. Są to istotne hipotezy,
nie dekoracja: dla krzywych afinicznych albo osobliwych
nie wszystkie późniejsze twierdzenia mają identyczną postać.

\begin{definicja}[Krzywa i ciało funkcji]
\emph{Krzywą algebraiczną} nad $k$ nazywamy
jednowymiarową nierozkładalną rozmaitość algebraiczną.
Jej \emph{ciało funkcji wymiernych} $k(C)$ jest
ciałem ułamków pierścienia współrzędnych dowolnego
niepustego afinicznego otwartego podzbioru $U\subseteq C$.
Definicja nie zależy od $U$: dla dwóch niepustych
afinicznych otwartych $U,V$ ich część wspólna jest
niepusta i zawiera niepusty afiniczny otwarty zbiór,
na którym obie klasy funkcji przyjmują te same wartości.
\end{definicja}

\begin{przyklad}[Ciało funkcji paraboli, prostej i kubiki]
Dla $C=\V(y-x^2)\subseteq\A^2_k$ mamy
$k[C]\cong k[t]$ przez $x\mapsto t$, $y\mapsto t^2$,
zatem $k(C)\cong k(t)$. Dla krzywej z ostrzem
$y^2=x^3$ pierścień współrzędnych jest
$k[t^2,t^3]$, ale ciało funkcji znowu wynosi $k(t)$.
Dwie krzywe mogą więc mieć to samo ciało funkcji
(\emph{być biwymierne}), chociaż jedna jest gładka,
a druga osobliwa.
\end{przyklad}

\begin{definicja}[Rząd zera lub bieguna]
\label{def:ord}
Niech $P\in C$ będzie punktem zamkniętym.
Ponieważ $\mathcal O_{C,P}$ jest DVR na mocy
wniosku~\ref{wn:krzywa-dvr}, istnieje lokalny parametr
$t_P$ generujący ideał maksymalny tego pierścienia.
Dla niezerowej funkcji wymiernej $f\in k(C)^\times$
liczbę całkowitą $\operatorname{ord}_P(f)$ definiujemy
tak, by $f=t_P^{\operatorname{ord}_P(f)}u$,
gdzie $u\in\mathcal O_{C,P}^\times$.
Gdy rząd jest dodatni, $f$ ma \emph{zero}; gdy
ujemny---\emph{biegun}; gdy wynosi zero, $f$
jest regularna i niezerowa w $P$.
\end{definicja}

\begin{stwierdzenie}[Arytmetyka rzędów]
Dla $f,g\in k(C)^\times$ i $P\in C$ zachodzą
$\operatorname{ord}_P(fg)=\operatorname{ord}_P(f)+
\operatorname{ord}_P(g)$ oraz
$\operatorname{ord}_P(f+g)\ge\min
\{\operatorname{ord}_P(f),\operatorname{ord}_P(g)\}$,
gdy $f+g\ne0$.
\end{stwierdzenie}
\begin{proof}
Zapisujemy $f=t^m u$ i $g=t^n v$ z jednostkami
$u,v\in\mathcal O_{C,P}^\times$. Wówczas
$fg=t^{m+n}uv$ i $uv$ jest jednostką, skąd
pierwsza równość. Jeśli $m\le n$, to
$f+g=t^m(u+t^{n-m}v)$, a czynnik w nawiasie
należy do $\mathcal O_{C,P}$; jego rząd jest
nieujemny (o ile nie jest zerem). Stąd nierówność.
\end{proof}

\begin{przyklad}[Punkt w nieskończoności na prostej rzutowej]
W $\PP^1_k$ przyjmijmy współrzędną afiniczną
$t=X_1/X_0$. Dla $a\in k$ funkcja $t-a$ ma
jedno zero rzędu $1$ w punkcie $P_a=[1:a]$.
W punkcie $\infty=[0:1]$ lokalnym parametrem jest
$s=1/t$, a $t-a=s^{-1}(1-as)$.
Ponieważ $1-as$ jest jednostką w $\mathcal O_{\PP^1,\infty}$,
$\operatorname{ord}_\infty(t-a)=-1$.
Nie można więc liczyć zer funkcji wymiernej,
pomijając punkty w nieskończoności.
\end{przyklad}

\section{Dywizory: uporządkowany zapis zer i biegunów}
\label{sec:divisors}

\begin{definicja}[Dywizor Weila na gładkiej krzywej]
\emph{Dywizorem} na $C$ jest formalna skończona suma
$D=\sum_{P\in C}n_P P$, gdzie $n_P\in\Z$
i tylko skończenie wiele $n_P$ jest różnych od zera.
Dywizory tworzą grupę abelową $\operatorname{Div}(C)$
z dodawaniem współczynników. Dywizor jest
\emph{efektywny}, co zapisujemy $D\ge0$,
gdy $n_P\ge0$ dla każdego $P$.
Jego \emph{stopniem} jest $\deg D=\sum_Pn_P$.
\end{definicja}

\begin{definicja}[Dywizor funkcji i równoważność liniowa]
Dla $0\ne f\in k(C)$ określamy
$\operatorname{div}(f)=\sum_P\operatorname{ord}_P(f)P$.
Jest to suma skończona: na krzywej rzutowej
$f$ jest regularna i niezerowa poza skończonym
zbiorem punktów. Dwa dywizory $D,E$ są
\emph{liniowo równoważne}, zapis $D\sim E$,
gdy $D-E=\operatorname{div}(f)$ dla pewnego
$f\in k(C)^\times$. Podgrupę dywizorów głównych
oznaczamy $\operatorname{Prin}(C)$, a iloraz
$\operatorname{Div}(C)/\operatorname{Prin}(C)$
przez $\operatorname{Cl}(C)$.
\end{definicja}

\begin{uwaga}[Dlaczego dywizor ma skończony nośnik?]
Dla niezerowej funkcji wymiernej $f$ istnieje niepuste
afiniczne otwarte $U\subseteq C$, na którym $f$
i $f^{-1}$ są regularne: po wyborze przedstawień
ułamkowych usuwamy miejsca zerowe mianowników
i liczników. Dopełnienie niepustego otwartego
podzbioru nierozkładalnej krzywej jest zbiorem
skończonym, gdyż każdy jego składnik nierozkładalny
ma wymiar zero. Poza dopełnieniem rząd wynosi zero.
\end{uwaga}

\begin{przyklad}[Dywizor funkcji na $\PP^1$]
Na $\PP^1_k$ dla $a\ne b$ mamy
$\operatorname{div}((t-a)/(t-b))=P_a-P_b$,
gdyż oba czynniki mają w nieskończoności
pojedynczy biegun i znoszą się w ilorazie.
W szczególności wszystkie punkty $\PP^1_k$
są liniowo równoważne.
\end{przyklad}

\begin{definicja}[Przestrzeń funkcji o ograniczonych biegunach]
Dla dywizora $D=\sum_Pn_PP$ definiujemy
\[
 L(D)=\{0\}\cup\{f\in k(C)^\times:
 \operatorname{div}(f)+D\ge0\},\qquad
 \ell(D)=\dim_k L(D).
\]
Warunek przynależności oznacza dokładnie
$\operatorname{ord}_P(f)\ge-n_P$ dla każdego $P$.
Jeśli $n_P>0$, dopuszczamy tam biegun rzędu
co najwyżej $n_P$; jeśli $n_P<0$, wymuszamy
zero rzędu co najmniej $-n_P$.
\end{definicja}

\begin{stwierdzenie}[Liniowo równoważne dywizory]
Jeśli $D'=D+\operatorname{div}(h)$ dla
$h\in k(C)^\times$, to
$L(D')=h^{-1}L(D)$. W szczególności
$\ell(D')=\ell(D)$.
\end{stwierdzenie}
\begin{proof}
Dla $g\ne0$ warunek $g\in L(D')$ jest równoważny
$\operatorname{div}(g)+D+\operatorname{div}(h)\ge0$,
czyli $\operatorname{div}(hg)+D\ge0$, a więc
$hg\in L(D)$. Mnożenie przez $h$ jest
izomorfizmem przestrzeni wektorowych.
\end{proof}

\begin{definicja}[Snop dywizora]
\label{def:O-D}
Dla dywizora $D=\sum_Pn_PP$ na gładkiej krzywej
$C$ definiujemy snop $\mathcal O_C(D)$ przez
\[
 \mathcal O_C(D)(U)=\{0\}\cup\{f\in k(C)^\times:
 \operatorname{ord}_P(f)+n_P\ge0
 \text{ dla wszystkich }P\in U\}.
\]
Ograniczenia są włączeniami funkcji wymiernych.
W szczególności $H^0(C,\mathcal O_C(D))=L(D)$.
Symbol $H^0$ oznacza przestrzeń sekcji globalnych.
\end{definicja}

\begin{twierdzenie}[Snop dywizora jest wiązką liniową]
\label{tw:O-D-invertible}
Dla dowolnego dywizora $D$ na gładkiej krzywej $C$
snop $\mathcal O_C(D)$ jest odwracalny.
Ponadto $\mathcal O_C(D+E)\cong
\mathcal O_C(D)\otimes_{\mathcal O_C}\mathcal O_C(E)$
oraz jeśli $D\sim E$, to
$\mathcal O_C(D)\cong\mathcal O_C(E)$.
\end{twierdzenie}
\begin{proof}
Wybierzmy punkt $P$ i lokalny parametr $t_P$.
Źdźbło w $P$ jest zbiorem ułamków o rzędzie
co najmniej $-n_P$, czyli wolnym
$\mathcal O_{C,P}$-modułem $t_P^{-n_P}\mathcal O_{C,P}$
rangi $1$. W otoczeniu $P$ możemy wybrać
funkcję regularną będącą lokalnym parametrem
i usunąć skończoną liczbę innych punktów nośnika
$D$; wtedy $\mathcal O_C(D)$ jest generowany
przez $t_P^{-n_P}$. To dowodzi lokalnej wolności.

W lokalnych generatorach iloczyn tensorowy
identyfikuje się z mnożeniem funkcji wymiernych:
$t^{-n_P}\otimes t^{-m_P}\mapsto t^{-(n_P+m_P)}$.
Daje to pierwszy izomorfizm. Jeśli
$E=D+\operatorname{div}(h)$, mnożenie przez
$h^{-1}$ identyfikuje lokalne sekcje
$\mathcal O_C(D)$ i $\mathcal O_C(E)$,
co daje drugi izomorfizm.
\end{proof}


\begin{twierdzenie}[Dywizory i grupa Picarda]
\label{tw:Pic-div}
Niech $C$ będzie gładką nierozkładalną krzywą
rzutową nad algebraicznie domkniętym ciałem $k$.
Odwzorowanie
$\operatorname{Div}(C)\to\operatorname{Pic}(C)$,
$D\mapsto[\mathcal O_C(D)]$ jest surjektywnym
homomorfizmem grup, którego jądrem jest
$\operatorname{Prin}(C)$. W szczególności
$\operatorname{Cl}(C)\cong\operatorname{Pic}(C)$.
\end{twierdzenie}
\begin{proof}
Z twierdzenia~\ref{tw:O-D-invertible} wiemy, że
$\mathcal O_C(D+E)\cong\mathcal O_C(D)\otimes
\mathcal O_C(E)$, więc przyporządkowanie jest
homomorfizmem.

Wykażemy surjektywność. Niech $\mathcal L$ będzie
snopem odwracalnym. W punkcie generycznym krzywej
jego źdźbło jest jednowymiarową przestrzenią
wektorową nad $k(C)$. Wybieramy jej niezerowy
element $s$ (\emph{sekcję wymierną}). W dowolnym
punkcie $P$ wybieramy lokalny generator $e_P$
snopa $\mathcal L$; wówczas $s=f_Pe_P$ z
$f_P\in k(C)^\times$. Określamy
$n_P=\operatorname{ord}_P(f_P)$. Ta liczba nie
zależy od zmiany generatora, bo dwie lokalne
bazy różnią się przez funkcję odwracalną,
mającą rząd zero. Poza skończoną liczbą punktów
$s$ jest lokalną bazą (wystarczy rozważyć
skończone afiniczne pokrycie i wyzerować na nim
skończone zbiory zer i biegunów współczynników).
Suma $D=\sum_P n_PP$ jest więc dywizorem.

W lokalnej bazie $e_P$ odwzorowanie
$\mathcal L\to\mathcal O_C(D)$ dane przez
$e_P\mapsto f_P^{-1}$ jest izomorfizmem:
$\operatorname{ord}_P(f_P^{-1})=-n_P$, a
$\mathcal O_C(D)_P=t_P^{-n_P}\mathcal O_{C,P}$.
Przy zmianie bazy współczynniki zmieniają się
tak, że te lokalne odwzorowania pozostają
zgodne. Otrzymujemy $\mathcal L\cong\mathcal O_C(D)$.

Pozostaje obliczyć jądro. Jeśli
$\mathcal O_C(D)\cong\mathcal O_C$, izomorfizm
w punkcie generycznym jest mnożeniem przez
pewien $h\in k(C)^\times$. Warunek, aby
mnożenie przez $h$ identyfikowało wszystkie
źdźbła $\mathcal O_{C,P}$ z
$t_P^{-n_P}\mathcal O_{C,P}$, wymusza
$\operatorname{ord}_P(h)=-n_P$ dla każdego $P$.
Zatem $D=\operatorname{div}(h^{-1})$ jest główny.
Odwrotnie, dla dywizora głównego izomorfizm
z $\mathcal O_C$ mamy już z
~\ref{tw:O-D-invertible}. To kończy dowód.
\end{proof}

\begin{uwaga}[Dlaczego gładkość jest ważna?]
Na krzywej osobliwej lokalne ideały punktów
nie muszą być główne. Wtedy każdy punkt nadal
można zapisać jako formalny dywizor Weila,
ale odpowiadający mu snop nie musi być odwracalny.
Na gładkiej krzywej DVR sprawia, że każdy
dywizor Weila jest również dywizorem Cartiera.
\end{uwaga}

\section{Różniczki Kählera i dywizor kanoniczny}
\label{sec:differentials}

\begin{definicja}[Różniczka uniwersalna]
Niech $A$ będzie $k$-algebrą. $k$-\emph{derywacją}
$A$ w $A$-moduł $M$ jest odwzorowanie
$d:A\to M$ spełniające $d(a+b)=d(a)+d(b)$,
$d(ab)=a\,d(b)+b\,d(a)$ i $d(c)=0$ dla $c\in k$.
Moduł \emph{różniczek Kählera} $\Omega^1_{A/k}$
wraz z $d:A\to\Omega^1_{A/k}$ charakteryzuje
następująca własność: każda $k$-derywacja
$\delta:A\to M$ rozkłada się w jedyny sposób
jako $\delta=\varphi\circ d$ przez homomorfizm
$A$-modułów $\varphi:\Omega^1_{A/k}\to M$.
\end{definicja}

\begin{stwierdzenie}[Różniczki w pierścieniu ilorazowym]
\label{st:kahler}
Niech $R=k[x_1,\ldots,x_n]$, $I=(f_1,\ldots,f_r)$
oraz $A=R/I$. Wówczas moduł $\Omega^1_{A/k}$ ma prezentację
\[
 A^r\xrightarrow{J^T}A^n\longrightarrow
 \Omega^1_{A/k}\longrightarrow0,
\]
gdzie $J=(\partial f_i/\partial x_j)$ jest
macierzą Jacobiego. W szczególności dla
$A=k[x,y]/(f)$ mamy relację $f_x\,dx+f_y\,dy=0$.
\end{stwierdzenie}
\begin{proof}
W pierścieniu wielomianów pochodne formalne
spełniają regułę Leibniza i każda derywacja
jest wyznaczona przez obrazy $x_i$;
wolny moduł $\bigoplus_i A\,dx_i$ realizuje więc
uniwersalną derywację przed przejściem do ilorazu.
Derywacja na ilorazie musi wysyłać każde $f_i$
na zero, zatem $df_i=\sum_j(\partial f_i/\partial x_j)dx_j=0$.
Jeśli te relacje nałożymy na moduł wolny,
to reguła Leibniza zapewnia $d(hf_i)=h\,df_i+f_i\,dh=0$
w ilorazie dla każdego $h$, czyli wyzerujemy
cały ideał $(f_1,\ldots,f_r)$. Powstały moduł
ma wymaganą własność uniwersalną; prezentacja
jest dokładnie podanym ilorazem.
\end{proof}

\begin{przyklad}[Dlaczego kryterium Jacobiego wykrywa osobliwość?]
Dla ostrza $f=y^2-x^3$ mamy
$df=2y\,dy-3x^2\,dx$. W punkcie $(0,0)$
relacja ta po redukcji modulo ideał maksymalny
staje się zerowa, więc przestrzeń kostyczna
ma dwa generatory $dx,dy$. Na gładkiej krzywej
płaskiej przynajmniej jedna z pochodnych $f_x,f_y$
jest w danym punkcie niezerowa, więc jedna
różniczka daje się wyrazić przez drugą;
wymiar przestrzeni kostycznej wynosi jeden.
\end{przyklad}

\begin{definicja}[Snop różniczek i dywizor kanoniczny]
Na gładkiej krzywej $C$ lokalne moduły
$\Omega^1_{A/k}$ sklejają się w snop
$\Omega^1_{C/k}$. Jest on lokalnie wolny rzędu
jeden. \emph{Dywizorem kanonicznym} nazywamy
dywizor $K_C=\operatorname{div}(\omega)$
dowolnej niezerowej \emph{wymiernej} różniczki
$\omega$, tzn. niezerowego elementu
jednowymiarowej $k(C)$-przestrzeni
$\Omega^1_{k(C)/k}$. Jego klasa liniowej
równoważności nie zależy od wyboru $\omega$,
a $\Omega^1_{C/k}\cong\mathcal O_C(K_C)$.
\end{definicja}

\begin{proof}[Uzasadnienie niezależności wyboru]
Jeśli $\eta$ jest drugą niezerową różniczką wymierną,
to $\eta=f\omega$ dla jednoznacznego $f\in k(C)^\times$,
bo przestrzeń różniczek nad ciałem funkcji jednej
zmiennej jest jednowymiarowa. W lokalnej bazie $dt_P$
piszemy $\omega=u\,dt_P$, więc
$\eta=fu\,dt_P$ i
$\operatorname{ord}_P(\eta)=\operatorname{ord}_P(f)+
\operatorname{ord}_P(\omega)$. Stąd
$\operatorname{div}(\eta)=\operatorname{div}(\omega)+
\operatorname{div}(f)$. Warunek regularności $g\omega$
w punkcie $P$ brzmi
$\operatorname{ord}_P(g)+\operatorname{ord}_P(\omega)\ge0$,
co identyfikuje snop różniczek ze snopem dywizora.
\end{proof}

\begin{przyklad}[Różniczka $dt$ na prostej rzutowej]
Dla $C=\PP^1_k$ i $t=X_1/X_0$ mamy na
części afinicznej regularną różniczkę $dt$.
W nieskończoności parametr $s=1/t$ daje
$dt=-s^{-2}ds$. Zatem
$\operatorname{div}(dt)=-2\infty$ i
$\Omega^1_{\PP^1/k}\cong\mathcal O_{\PP^1}(-2)$.
Stopień dywizora kanonicznego wynosi tu $-2$.
\end{przyklad}

\section{Kohomologia snopów: dlaczego sekcje globalne nie wystarczą?}
\label{sec:cech-curve}

\begin{definicja}[Kohomologia Čecha dla pokrycia dwoma zbiorami]
Niech $X=U\cup V$ i $\mathcal F$ będzie snopem
grup abelowych na $X$. Rozważmy homomorfizm
\[
 \delta:\mathcal F(U)\oplus\mathcal F(V)
 \longrightarrow\mathcal F(U\cap V),\quad
 \delta(s,t)=s|_{U\cap V}-t|_{U\cap V}.
\]
Definiujemy $\check H^0(\{U,V\},\mathcal F)=\ker\delta$
i $\check H^1(\{U,V\},\mathcal F)=\operatorname{coker}\delta$.
Pierwsza przestrzeń identyfikuje się z $\mathcal F(X)$
dzięki aksjomatowi snopa. Druga mierzy \emph{przeszkodę
w rozłożeniu} zadanej sekcji na przecięciu na
różnicę dwóch sekcji określonych osobno na $U$ i $V$.
\end{definicja}

\begin{uwaga}[Czym są ogólne grupy $H^i$?]
Dla dowolnego snopa grup abelowych $\mathcal F$
przestrzenie (grupy) $H^i(X,\mathcal F)$
można zdefiniować jako funktory pochodne funktora
sekcji globalnych $\Gamma(X,-)$. Dla snopów
quasi-koherentnych na oddzielonym schemacie
rozkrytym przez afiniczne zbiory otwarte kohomologia
Čecha takiego pokrycia oblicza kohomologię snopa.
To twierdzenie wymaga \emph{znikania wyższej kohomologii
snopów quasi-koherentnych na schematach afinicznych}.
Nie dowodzimy tutaj całej teorii funktorów pochodnych;
w obliczeniach na $\PP^1$ wykonamy jednak każdą
operację bezpośrednio z definicji kompleksu Čecha.
\end{uwaga}

\begin{stwierdzenie}[Dwuelementowe pokrycie afiniczne krzywej]
Niech $k$ będzie ciałem nieskończonym, a
$C\subseteq\PP^N_k$ nierozkładalną krzywą rzutową.
Wówczas $C$ ma pokrycie dwoma otwartymi zbiorami
afinicznymi $U,V$ o afinicznym przecięciu.
\end{stwierdzenie}
\begin{proof}
Wybieramy hiperpłaszczyznę $H_0$, która nie zawiera
$C$. Jej przecięcie z krzywą jest skończone,
a $U=C\setminus H_0$ jest afiniczny, gdyż
jest zamkniętym podzbiorem mapy afinicznej
$\PP^N\setminus H_0\cong\A^N$.
Istnieje druga hiperpłaszczyzna $H_1$, która
omija wszystkie punkty skończonego zbioru
$C\cap H_0$: dla każdego z tych punktów
warunek zawierania go przez hiperpłaszczyznę
jest jednym właściwym liniowym warunkiem na
współczynniki równania hiperpłaszczyzny,
a skończona suma właściwych podprzestrzeni
nie pokrywa całej przestrzeni liniowej nad
nieskończonym ciałem $k$.
Zatem $V=C\setminus H_1$ jest afiniczny i
$U\cup V=C$. Zbiór $U\cap V$ to podstawowy
otwarty podzbiór afinicznego $U$, otrzymany
przez nieznikanie równania $H_1$ podzielonego
przez odpowiednią potęgę równania $H_0$.
Jest zatem afiniczny, gdyż $D(f)$ w
$\Spec A$ jest izomorficzny ze $\Spec A_f$.
\end{proof}

\begin{twierdzenie}[Skończoność i znikanie kohomologii krzywej]
\label{tw:curve-coh}
Niech $k$ będzie ciałem algebraicznie domkniętym,
$C$ gładką, nierozkładalną krzywą rzutową nad $k$,
a $\mathcal L$ snopem odwracalnym na $C$.
Wówczas $H^0(C,\mathcal L)$ i $H^1(C,\mathcal L)$ są
skończeniewymiarowymi przestrzeniami nad $k$,
a $H^i(C,\mathcal L)=0$ dla $i\ge2$.
\end{twierdzenie}
\begin{proof}[Redukcja do standardowych twierdzeń kohomologicznych]
Pokryjmy $C$ dwoma afinicznymi otwartymi $U,V$
jak wyżej. Znikanie kohomologii quasi-koherentnej
na schematach afinicznych oraz twierdzenie
porównawcze Čecha dają dokładny kompleks
\[
 0\longrightarrow H^0(C,\mathcal L)
 \longrightarrow\mathcal L(U)\oplus\mathcal L(V)
 \xrightarrow{\delta}\mathcal L(U\cap V)
 \longrightarrow H^1(C,\mathcal L)\longrightarrow0.
\]
Ponieważ kompleks ma długość jeden,
$H^i=0$ dla $i\ge2$.
Skończoność wymiarów jest standardowym
twierdzeniem o kohomologii snopów koherentnych
na przestrzeniach rzutowych. Jego pełny dowód
wymaga twierdzenia Serre'a o skończoności
kohomologii i nie wynika z samej długości
kompleksu (jego wyrazy mogą być nieskończeniewymiarowe).
W dalszym ciągu będziemy używać tego rezultatu
jako jawnie wskazanego twierdzenia pomocniczego.
\end{proof}

\begin{definicja}[Charakterystyka Eulera i rodzaj]
Dla snopa odwracalnego $\mathcal L$ na gładkiej
rzutowej krzywej $C$ definiujemy
$\chi(C,\mathcal L)=h^0(C,\mathcal L)-h^1(C,\mathcal L)$,
gdzie $h^i=\dim_kH^i$. \emph{Rodzaj algebraiczny}
krzywej $C$ określamy wzorem
$g(C)=h^1(C,\mathcal O_C)$.
\end{definicja}

\section{Riemann--Roch: ile funkcji ma dozwolone bieguny?}
\label{sec:rr}

\begin{lemat}[Zmiana dywizora o punkt]
\label{lem:exact-point}
Niech $k$ będzie ciałem algebraicznie domkniętym,
$C$ gładką, nierozkładalną, rzutową krzywą nad $k$,
$D$ dywizorem na $C$, a $P\in C(k)$ punktem.
Istnieje krótki ciąg dokładny snopów
\[
 0\longrightarrow\mathcal O_C(D)
 \longrightarrow\mathcal O_C(D+P)
 \longrightarrow k_P\longrightarrow0,
\]
gdzie $k_P$ jest snopem mającym źdźbło $k$
w punkcie $P$ i zero poza $P$.
\end{lemat}
\begin{proof}
W punkcie $Q\ne P$ oba snopy mają takie same
źdźbła, ponieważ współczynniki $D$ i $D+P$
są tam równe. W punkcie $P$ wybieramy parametr
$t$ i oznaczamy $n=n_P$. Źdźbła mają postać
$t^{-n}R\subseteq t^{-(n+1)}R$, gdzie
$R=\mathcal O_{C,P}$ i jego ideał maksymalny
wynosi $(t)$. Iloraz jest izomorficzny z
$R/(t)=k$ przez wybranie współczynnika przy
$t^{-(n+1)}$. Ponieważ dokładność snopów
sprawdza się na źdźbłach, otrzymujemy tezę.
\end{proof}

\begin{twierdzenie}[Riemann--Roch w postaci Eulera]
\label{tw:rr-euler}
Niech $C$ będzie gładką, nierozkładalną,
rzutową krzywą nad algebraicznie domkniętym
ciałem $k$. Dla każdego dywizora $D$ na $C$
zachodzi
\[
 \chi\bigl(C,\mathcal O_C(D)\bigr)
   =\deg D+1-g(C).
\]
\end{twierdzenie}
\begin{proof}
Z krótkiego ciągu dokładnego
lematu~\ref{lem:exact-point} i długiego ciągu
kohomologii otrzymujemy
\[
\begin{aligned}
0&\longrightarrow H^0(C,\mathcal O_C(D))
\longrightarrow H^0(C,\mathcal O_C(D+P))
\longrightarrow k\\
&\longrightarrow H^1(C,\mathcal O_C(D))
\longrightarrow H^1(C,\mathcal O_C(D+P))
\longrightarrow0.
\end{aligned}
\]
Z addytywności sumy naprzemiennej wymiarów
w skończonym ciągu dokładnym wynika
$\chi(C,\mathcal O_C(D+P))=
\chi(C,\mathcal O_C(D))+1$.
Tę samą równość wolno zastosować w odwrotną
stronę, zastępując $D$ przez $D-P$.
Każdy dywizor jest skończoną sumą punktów
z całkowitymi współczynnikami. Indukcyjnie
$\chi(C,\mathcal O_C(D))=
\chi(C,\mathcal O_C)+\deg D$.

Pozostaje obliczyć $\chi(C,\mathcal O_C)$.
Każda globalna funkcja regularna na nierozkładalnej
rzutowej krzywej jest stała: jeśli $f$ byłaby
niestała, określałaby morfizm $C\to\A^1$,
a także morfizm $C\to\PP^1$ omijający
$\infty$. W tym miejscu używamy standardowego
wyniku o \emph{właściwości morfizmów rzutowych}:
obraz zbioru domkniętego przez morfizm właściwy
jest domknięty. Dowód tego wyniku (przez eliminację
rzutową lub kryterium wartościujące) wykracza poza
dotychczasowy aparat rozdziałów 0--4 i jest tutaj
jawnie przyjęty. W konsekwencji obraz krzywej
rzutowej przez morfizm do $\PP^1$ jest domknięty,
a obraz morfizmu niestałego jest nieskończony;
na nierozkładalnej prostej rzutowej jedynym
nieskończonym zbiorem domkniętym jest cała
$\PP^1$, sprzecznie z omijaniem $\infty$.
Zatem $H^0(C,\mathcal O_C)=k$, skąd
$\chi(C,\mathcal O_C)=1-h^1(C,\mathcal O_C)=1-g(C)$.
Podstawienie kończy dowód.
\end{proof}

\begin{wniosek}[Dywizory główne mają stopień zero]
\label{wn:div-zero}
Dla każdego $0\ne f\in k(C)$ zachodzi
$\deg\operatorname{div}(f)=0$.
\end{wniosek}
\begin{proof}
Snopy $\mathcal O_C$ i
$\mathcal O_C(\operatorname{div}(f))$ są izomorficzne
przez mnożenie przez $f^{-1}$
(twierdzenie~\ref{tw:O-D-invertible}),
więc mają tę samą charakterystykę Eulera.
Z twierdzenia~\ref{tw:rr-euler} różnica
ich charakterystyk wynosi
$\deg\operatorname{div}(f)$, musi więc być zerem.
\end{proof}

\begin{twierdzenie}[Dualność Serre'a dla krzywych]
\label{tw:serre-curve}
Niech $C$ będzie gładką nierozkładalną rzutową
krzywą nad algebraicznie domkniętym ciałem $k$,
a $D$ dywizorem na $C$. Istnieje kanoniczne
niezdegenerowane parowanie
\[
 H^1\bigl(C,\mathcal O_C(D)\bigr)
 \times H^0\bigl(C,\Omega^1_{C/k}\otimes
 \mathcal O_C(-D)\bigr)\longrightarrow k.
\]
W szczególności
$h^1(C,\mathcal O_C(D))=\ell(K_C-D)$.
\end{twierdzenie}

\begin{uwaga}[Zakres dowodu dualności]
Dualność Serre'a jest istotnym, głębszym rezultatem,
a nie tożsamością wynikającą z samej definicji
$H^1$. Można ją udowodnić przez residua:
klasa kohomologii ma przedstawicieli w postaci
lokalnych części głównych funkcji wymiernych,
a różniczka $\omega$ działa na taką klasę przez
$\sum_P\operatorname{Res}_P(f_P\omega)$.
Globalne twierdzenie o residuach zapewnia
niezależność przedstawicieli, zaś lokalna
niezdegenerowaność parowania części głównych
i różniczek, wraz z globalnym argumentem
skończoności, prowadzi do izomorfizmu dualnego.
Ten ciąg argumentów wymaga osobnego rozwinięcia
teorii residuów i kohomologii; \emph{nie traktujemy
powyższego opisu jako pełnego dowodu}.
Poniżej podajemy pełny dowód redukcji
Riemanna--Rocha do dualności Serre'a.
\end{uwaga}

\begin{twierdzenie}[Riemann--Roch w klasycznej postaci]
\label{tw:rr}
Niech $k$ będzie ciałem algebraicznie domkniętym,
$C$ gładką nierozkładalną rzutową krzywą nad $k$,
$g=g(C)$ jej rodzajem, $K_C$ dywizorem kanonicznym,
a $D$ dowolnym dywizorem na $C$. Wówczas
\[
 \ell(D)-\ell(K_C-D)=\deg D+1-g.
\]
\end{twierdzenie}
\begin{proof}
Z definicji~\ref{def:O-D} mamy
$\ell(D)=h^0(C,\mathcal O_C(D))$.
Dualność Serre'a (twierdzenie~\ref{tw:serre-curve})
identyfikuje drugi wymiar
$h^1(C,\mathcal O_C(D))$ z
$h^0(C,\Omega^1_{C/k}\otimes\mathcal O_C(-D))$.
Ponieważ $\Omega^1_{C/k}\cong\mathcal O_C(K_C)$,
ten ostatni wymiar to $\ell(K_C-D)$.
Zatem lewa strona w tezie jest dokładnie
$\chi(C,\mathcal O_C(D))$, a twierdzenie
~\ref{tw:rr-euler} oblicza ją jako
$\deg D+1-g$.
\end{proof}

\begin{wniosek}[Stopień dywizora kanonicznego]
\label{wn:degK}
Dla gładkiej nierozkładalnej rzutowej krzywej
$C$ rodzaju $g$ mamy $\ell(K_C)=g$
oraz $\deg K_C=2g-2$.
\end{wniosek}
\begin{proof}
Podstawiamy $D=0$ do Riemanna--Rocha:
$\ell(0)-\ell(K_C)=1-g$. Ponieważ
$\ell(0)=h^0(C,\mathcal O_C)=1$,
wynika $\ell(K_C)=g$. Następnie wstawiamy
$D=K_C$: $\ell(K_C)-\ell(0)=\deg K_C+1-g$.
Stąd $g-1=\deg K_C+1-g$, czyli
$\deg K_C=2g-2$.
\end{proof}

\begin{wniosek}[Dywizory dużego stopnia]
Jeśli $D$ jest dywizorem na $C$ i $\deg D>2g-2$,
to $\ell(D)=\deg D+1-g$.
Jeśli $\deg D<0$, to $\ell(D)=0$.
\end{wniosek}
\begin{proof}
Dla dywizora $E$ ujemnego stopnia $L(E)=0$:
gdyby istniało $0\ne f\in L(E)$, to
$\operatorname{div}(f)+E$ byłby efektywny,
a z wniosku~\ref{wn:div-zero} jego stopień byłby
równy $\deg E<0$, co jest niemożliwe.
Jeżeli $\deg D>2g-2$, to
$\deg(K_C-D)=2g-2-\deg D<0$,
zatem $\ell(K_C-D)=0$; teza wynika z Riemanna--Rocha.
\end{proof}

\section{Rachunek kompletny: \texorpdfstring{$\PP^1$}{P1} i wiązki \texorpdfstring{$\mathcal O(n)$}{O(n)}}
\label{sec:pp1-rr}

\begin{twierdzenie}[Kohomologia $\mathcal O_{\PP^1}(n)$]
\label{tw:p1-coh}
Dla dowolnego ciała $k$ i liczby całkowitej $n$
mamy
\[
 h^0\bigl(\PP^1_k,\mathcal O(n)\bigr)=\max\{n+1,0\},\qquad
 h^1\bigl(\PP^1_k,\mathcal O(n)\bigr)=\max\{-n-1,0\}.
\]
W szczególności $g(\PP^1)=0$ i
$\chi(\PP^1,\mathcal O(n))=n+1$.
\end{twierdzenie}
\begin{proof}
Korzystamy ze standardowego pokrycia
$\PP^1=U_0\cup U_\infty$, gdzie
$U_0=\Spec k[t]$, $U_\infty=\Spec k[s]$,
$s=1/t$, a przecięcie ma pierścień
$k[t,t^{-1}]$. Dla $\mathcal O(n)$
funkcja przejścia wynosi $e_\infty=t^ne_0$
(zob. sekcję~\ref{sec:wiazki}). Zatem kompleks
Čecha ma postać
\[
 k[t]\oplus k[t^{-1}]
 \xrightarrow{\delta}
 k[t,t^{-1}],\qquad
 \delta(a,b)=a(t)-t^n b(t^{-1}).
\]
Element jądra jest parą $(a,b)$, dla której
$a(t)=t^n b(t^{-1})$. Ponieważ $a$ jest
wielomianem o nieujemnych wykładnikach,
a $b$ wielomianem w $t^{-1}$,
taki $a$ istnieje wtedy i tylko wtedy, gdy
$n\ge0$, a wtedy można wybrać dowolny
$a(t)=a_0+\cdots+a_nt^n$ i odzyskać
$b(t^{-1})=t^{-n}a(t)$. Wymiar jądra wynosi
$n+1$ dla $n\ge0$, a zero dla $n<0$.

Kokernel jest ilorazem
$k[t,t^{-1}]/\bigl(k[t]+t^n k[t^{-1}]\bigr)$.
Wielomiany Laurenta mają bazę $\{t^m:m\in\Z\}$.
Pierwszy składnik zabija wszystkie potęgi
$t^m$ dla $m\ge0$, a drugi wszystkie dla
$m\le n$. Niezredukowane pozostają dokładnie
potęgi o wykładnikach $n<m<0$.
Jeśli $n\ge-1$, nie pozostaje żadna;
jeśli $n\le-2$, jest ich $-n-1$.
To daje drugi wzór. Dla $n=0$ mamy
$h^1(\PP^1,\mathcal O)=0$, zatem rodzaj jest zerowy,
a różnica podanych wymiarów wynosi $n+1$
dla wszystkich $n\in\Z$.
\end{proof}

\begin{przyklad}[Wzór Riemanna--Rocha dla funkcji racjonalnych]
Na $\PP^1$ mamy $\mathcal O(n\infty)\cong\mathcal O(n)$.
Dla $n\ge0$ przestrzeń $L(n\infty)$ składa się
dokładnie z wielomianów stopnia nie większego niż $n$:
$f(t)=a_0+\cdots+a_nt^n$. Stąd $\ell(n\infty)=n+1$.
Zjawisko jest dokładnie takie, jak przewiduje Riemann--Roch
przy $g=0$ i $K\sim-2\infty$.
Dla $n=-1$ nie ma niezerowej funkcji, bo wymagałoby
to jednego zera bez żadnych biegunów.
\end{przyklad}

\section{Rodzaj gładkiej krzywej płaskiej i wzór adjunkcji}
\label{sec:genus-plane}

\begin{twierdzenie}[Adjunkcja dla gładkiej krzywej płaskiej]
\label{tw:adjunction-plane}
Niech $k$ będzie ciałem algebraicznie domkniętym,
a $C=V_+(F)\subseteq\PP^2_k$ gładką nierozkładalną
krzywą zadaną jednorodnym wielomianem $F$
stopnia $d\ge1$. Wówczas
$\Omega^1_{C/k}\cong\mathcal O_C(d-3)$.
Jeżeli $K_C$ jest dywizorem kanonicznym,
to $\deg K_C=d(d-3)$.
\end{twierdzenie}
\begin{proof}
Wyprowadzimy lokalny wzór, z którego widać
zarówno skąd bierze się liczba $3$, jak i dlaczego
wymagamy gładkości. Na mapie $Z\ne0$ wprowadzamy
$x=X/Z$, $y=Y/Z$ i $f(x,y)=F(x,y,1)$.
Relacja dla różniczek jest równa
$f_x\,dx+f_y\,dy=0$ (stwierdzenie~\ref{st:kahler}).
Na zbiorze $f_y\ne0$ wybieramy bazę
$\omega_Z=dx/f_y$, a na zbiorze $f_x\ne0$
możemy wybrać $\omega_Z=-dy/f_x$.
Na ich części wspólnej te formy są równe
dzięki relacji $df=0$. Ponieważ $C$ jest gładka,
zbiory te pokrywają całą część $C\cap\{Z\ne0\}$.
Otrzymujemy lokalny nieznikający generator
snopa $\Omega^1_{C/k}$.

Teraz przechodzimy do mapy $X\ne0$ z
$u=Y/X$, $v=Z/X$ i $G(u,v)=F(1,u,v)$.
W części $X\ne0$, gdzie $G_u\ne0$,
analogicznym generatorem jest $\omega_X=-dv/G_u$.
Na przecięciu $X\ne0$, $Z\ne0$ mamy
$x=1/v$, $y=u/v$, a z jednorodności $F$ wynika
$G(u,v)=v^d f(1/v,u/v)$ i
$G_u=v^{d-1}f_y$.
Ponieważ $dx=-v^{-2}dv$, otrzymujemy
\[
 \omega_Z=\frac{dx}{f_y}
   =-\,v^{d-3}\frac{dv}{G_u}
   =v^{d-3}\omega_X.
\]
Na pozostałych częściach map otrzymujemy tę samą
regułę przejścia przez użycie alternatywnych
postaci $-dy/f_x$ i $du/G_v$; wzory są zgodne
na częściach wspólnych, a więc definiują
izomorfizm snopów. Tymczasem standardowe lokalne
trywializacje $\mathcal O_C(d-3)$ w mapach
$Z\ne0$ i $X\ne0$ mają dokładnie przejście
$e_Z=v^{d-3}e_X$. Stąd
$\Omega^1_{C/k}\cong\mathcal O_C(d-3)$.

Na płaskiej krzywej stopnia $d$ przecięcie
z ogólną prostą ma $d$ punktów z uwzględnieniem
krotności (twierdzenie Bézouta z
rozdziału~\ref{ch:rzutowa}); to znaczy
$\deg\mathcal O_C(1)=d$. Z własności tensorowych
wynika $\deg\mathcal O_C(d-3)=d(d-3)$.
Ponieważ snop kanoniczny jest izomorficzny
z $\mathcal O_C(K_C)$, otrzymujemy tezę.
\end{proof}

\begin{twierdzenie}[Rodzaj gładkiej krzywej płaskiej]
\label{tw:plane-genus}
Niech $C\subseteq\PP^2_k$ będzie gładką,
nierozkładalną krzywą stopnia $d\ge1$
nad algebraicznie domkniętym ciałem $k$.
Wówczas jej rodzaj jest równy
\[
 g(C)=\frac{(d-1)(d-2)}2.
\]
\end{twierdzenie}
\begin{proof}
Wniosek~\ref{wn:degK} daje $2g-2=\deg K_C$,
a twierdzenie~\ref{tw:adjunction-plane} mówi,
że $\deg K_C=d(d-3)$. Zatem
$2g-2=d^2-3d$ i po przekształceniu
$2g=d^2-3d+2=(d-1)(d-2)$.
\end{proof}

\begin{przyklad}[Stopień wielomianu określa topologię]
Nad $\C$ gładka krzywa rzutowa jest zwartą
powierzchnią Riemanna. Jej rodzaj algebraiczny
jest równy liczbie uchwytów powierzchni
(identyfikacja wynika z porównania algebraicznej
kohomologii snopów z analityczną i z klasycznej
teorii powierzchni Riemanna). Otrzymujemy:
\begin{center}
\begin{tabular}{c|c|l}
stopień $d$&rodzaj $g$&topologia nad $\C$\\\hline
$1$&$0$&sfera\\
$2$&$0$&sfera\\
$3$&$1$&torus\\
$4$&$3$&powierzchnia z trzema uchwytami\\
$5$&$6$&powierzchnia z sześcioma uchwytami
\end{tabular}
\end{center}
Ważne zastrzeżenie: dla krzywej \emph{osobliwej}
stopnia $d$ wzór nadal oblicza rodzaj \emph{arytmetyczny},
ale rodzaj gładkiej normalizacji może być mniejszy.
\end{przyklad}

\begin{uwaga}[Jaka jest różnica między dwoma wzorami na rodzaj?]
W rozdziale~\ref{ch:osobliwosci} wykazaliśmy z
wielomianu Hilberta, że każda płaska krzywa
stopnia $d$ ma rodzaj arytmetyczny
$(d-1)(d-2)/2$. W niniejszej sekcji założyliśmy
gładkość i użyliśmy dywizora kanonicznego oraz
Riemanna--Rocha, aby wykazać, że \emph{rodzaj
geometryczny} gładkiej krzywej jest taki sam.
Na krzywej z ostrzem lub węzłem oba pojęcia
nie są automatycznie jednakowe.
\end{uwaga}

\section{Krzywe eliptyczne: dodawanie punktów z równania kubicznego}
\label{sec:elliptic}

\begin{definicja}[Krzywa eliptyczna]
Krzywą eliptyczną nad ciałem $k$ nazywamy
parę $(E,O)$, gdzie $E$ jest gładką
rzutową geometrycznie nierozkładalną krzywą
rodzaju jeden, a $O\in E(k)$ wyróżnionym
punktem. Gdy $\operatorname{char}k\ne2,3$,
klasyczny model to domknięcie rzutowe równania
$y^2=x^3+ax+b$, przy warunku
$4a^3+27b^2\ne0$. Punkt w nieskończoności
$O=[0:1:0]$ jest elementem wyróżnionym.
\end{definicja}

\begin{przyklad}[Sprawdzenie gładkości]
Dla $E:y^2=x^3-x$ nad $\C$ mamy
$f(x,y)=y^2-x^3+x$. Równania osobliwości
w części afinicznej wymagają jednocześnie
$f=0$, $f_x=-3x^2+1=0$ i $f_y=2y=0$.
Z $y=0$ i $f=0$ otrzymujemy
$x\in\{0,1,-1\}$, a dla żadnego z tych $x$
nie zachodzi $-3x^2+1=0$. Punkt w nieskończoności
jest również gładki (pochodna rzutowego równania
po $Z$ jest tam niezerowa). Krzywa ma stopień trzy,
a więc rodzaj jeden na mocy
~\ref{tw:plane-genus}.
\end{przyklad}

\begin{twierdzenie}[Punkty krzywej eliptycznej tworzą grupę]
\label{tw:elliptic-group}
Niech $(E,O)$ będzie krzywą eliptyczną nad
algebraicznie domkniętym ciałem $k$.
Odwzorowanie $P\mapsto[P-O]$ jest bijekcją
$E(k)\to\operatorname{Pic}^0(E)$, gdzie
$\operatorname{Pic}^0(E)$ oznacza podgrupę
$\operatorname{Pic}(E)$ odpowiadającą klasom
dywizorów stopnia zero przez izomorfizm
z twierdzenia~\ref{tw:Pic-div}.
Transport działania z $\operatorname{Pic}^0(E)$
nadaje $E(k)$ strukturę grupy abelowej
z elementem neutralnym $O$.
\end{twierdzenie}
\begin{proof}
Dla dowolnego dywizora $D$ stopnia zero
$D+O$ ma stopień jeden. Ponieważ $g(E)=1$
i $\deg K_E=0$,
$\deg(K_E-(D+O))=-1$. Wniosek po Riemannie--Rochu
daje $\ell(D+O)=1$. Istnieje więc niezerowa funkcja
$f\in L(D+O)$, a dywizor
$\operatorname{div}(f)+D+O$ jest efektywny
stopnia jeden. Każdy taki dywizor ma postać
jednego punktu $P$, skąd $D\sim P-O$.
To dowodzi surjektywności.

Jeśli $P-O\sim Q-O$, to $P\sim Q$.
Gdyby $P\ne Q$, istniałaby niezerowa funkcja
$f$ z $\operatorname{div}(f)=Q-P$.
Wówczas $1$ oraz $f$ byłyby liniowo niezależnymi
elementami $L(P)$, gdyż $f$ ma biegun w $P$.
Przeczyłoby to równości $\ell(P)=1$,
udowodnionej wyżej dla dywizora stopnia jeden.
Stąd $P=Q$, co daje iniektywność.

Zbiór klas stopnia zero jest podgrupą grupy
$\operatorname{Div}(E)/\operatorname{Prin}(E)$,
ponieważ dywizory główne mają stopień zero
(wniosek~\ref{wn:div-zero}). Jest więc grupą
abelową. Transport działania przez bijekcję
daje grupę na $E(k)$, a klasa $[O-O]$ jest
neutralna, zatem odpowiadający jej punkt to $O$.
\end{proof}

\begin{przyklad}[Skąd bierze się rysunkowa reguła siecznych?]
Na gładkiej kubice w postaci Weierstrassa
prosta przecina krzywą w trzech punktach
$P,Q,R$ z uwzględnieniem krotności
(twierdzenie Bézouta). Iloraz liniowego
równania tej prostej przez równanie prostej
w nieskończoności wyznacza funkcję wymierną
z dywizorem $P+Q+R-3O$.
Zatem $[P-O]+[Q-O]=-[R-O]$ w
$\operatorname{Pic}^0(E)$. Współrzędne punktu
przeciwnego do $R=(x_R,y_R)$ wynoszą
$-R=(x_R,-y_R)$, ponieważ linia pionowa ma
dywizor przecięcia $R+(-R)+O$.
Zatem $P+Q=-R$, gdzie $R$ jest trzecim
punktem przecięcia prostej z krzywą.
\end{przyklad}

\begin{center}
\begin{tikzpicture}[>=Stealth,scale=.90]
\draw[->,gray!70] (-2.25,0)--(3.0,0) node[right] {$x$};
\draw[->,gray!70] (0,-2.45)--(0,2.5) node[above] {$y$};
% An actual smooth real cubic: y^2=(x+1)(x+1/2)(x-1).
\draw[blue!65!black,very thick,domain=-1:-.5,samples=100]
 plot (\x,{sqrt(max(0,(\x+1)*(\x+.5)*(\x-1)))});
\draw[blue!65!black,very thick,domain=-1:-.5,samples=100]
 plot (\x,{-sqrt(max(0,(\x+1)*(\x+.5)*(\x-1)))});
\draw[blue!65!black,very thick,domain=1:1.88,samples=100]
 plot (\x,{sqrt(max(0,(\x+1)*(\x+.5)*(\x-1)))});
\draw[blue!65!black,very thick,domain=1:1.88,samples=100]
 plot (\x,{-sqrt(max(0,(\x+1)*(\x+.5)*(\x-1)))});
% Etykieta przeniesiona do podpisu.
\end{tikzpicture}

{\small Rysunek 4.4. Dwie składowe rzeczywiste kubiki $y^2=(x+1)(x+\tfrac12)(x-1)$; jest to wykres konkretnego równania, a nie schemat dodawania.}\par
\end{center}

\begin{stwierdzenie}[Wzory na sumę punktów]
Niech $k$ będzie ciałem charakterystyki różnej od $2,3$,
a $E:y^2=x^3+ax+b$ krzywą eliptyczną.
Jeśli $P=(x_1,y_1)$ i $Q=(x_2,y_2)$ są punktami
afinicznymi, $P\ne-Q$, to współrzędne
$P+Q=(x_3,y_3)$ wyrażają się wzorami
\[
 m=\begin{cases}
 \dfrac{y_2-y_1}{x_2-x_1},&x_1\ne x_2,\\[4pt]
 \dfrac{3x_1^2+a}{2y_1},&P=Q,\ y_1\ne0,
 \end{cases}
 \qquad x_3=m^2-x_1-x_2,\qquad
 y_3=m(x_1-x_3)-y_1.
\]
Jeśli $Q=-P$, to $P+Q=O$.
\end{stwierdzenie}
\begin{proof}
Prosta przez $P,Q$ ma równanie
$y=m(x-x_1)+y_1$, a dla $P=Q$ jest to
styczna o nachyleniu uzyskanym przez
różniczkowanie $2y\,dy=(3x^2+a)dx$.
Po podstawieniu równania prostej do równania
kubiki otrzymujemy moniczny wielomian trzeciego
stopnia względem $x$ z przeciwnym współczynnikiem
przy $x^2$ równym $m^2$. Suma jego trzech
pierwiastków, z krotnościami, wynosi $m^2$.
Dwa znane pierwiastki to $x_1,x_2$, a trzeci
$x_R=m^2-x_1-x_2$. Współrzędna $y_R$ jest
wartością równania prostej w $x_R$, a odbicie
$R\mapsto-R$ zmienia znak $y$. Stąd
$y_3=-[m(x_R-x_1)+y_1]=m(x_1-x_R)-y_1$.
Jeśli $P=-Q$, prosta jest pionowa i trzecim
punktem przecięcia w modelu rzutowym jest $O$;
stąd suma jest neutralna.
\end{proof}

\section{Rozgałęzienia, formuła Riemanna--Hurwitza i krzywe hipereliptyczne}
\label{sec:hurwitz}

\begin{definicja}[Indeks rozgałęzienia]
Niech $k$ będzie ciałem charakterystyki zero, a
$f:C\to B$ niestałym morfizmem gładkich
nierozkładalnych krzywych rzutowych nad $k$.
Dla punktu zamkniętego $P\in C$ i $Q=f(P)$
wybieramy lokalne parametry $t$ w $P$ i $u$ w $Q$. Istnieje
jednostka $v\in\mathcal O_{C,P}^\times$ i
liczba całkowita $e_P\ge1$ taka, że
$f^*u=t^{e_P}v$. Liczbę $e_P$ nazywamy
\emph{indeksem rozgałęzienia} morfizmu $f$ w $P$.
Jeśli $e_P>1$, punkt $P$ nazywamy rozgałęzionym.
\end{definicja}

\begin{twierdzenie}[Riemann--Hurwitz, charakterystyka zero]
\label{tw:RH}
Niech $k$ będzie algebraicznie domknięte
charakterystyki zero, a $f:C\to B$ niestałym
morfizmem gładkich nierozkładalnych krzywych
rzutowych, stopnia $n=[k(C):k(B)]$.
Wówczas
\[
 2g(C)-2=n\bigl(2g(B)-2\bigr)
  +\sum_{P\in C}(e_P-1).
\]
Suma jest skończona.
\end{twierdzenie}
\begin{proof}
Niech $\omega$ będzie niezerową wymierną
różniczką na $B$. W punkcie $Q=f(P)$ możemy
zapisać $\omega=u^m w\,du$, gdzie $w$
jest jednostką w pierścieniu lokalnym $\mathcal O_{B,Q}$.
Po podstawieniu $u=t^{e_P}v$ dostajemy
\[
 f^*\omega=t^{e_Pm}v^m(f^*w)
 (e_Pt^{e_P-1}v+t^{e_P}v')\,dt.
\]
Ponieważ $\operatorname{char}k=0$, liczba
$e_P$ jest niezerowa w $k$, a wyrażenie
$e_Pv+tv'$ jest jednostką. Zatem
$\operatorname{ord}_P(f^*\omega)=
 e_P\operatorname{ord}_Q(\omega)+(e_P-1)$.
Sumując po punktach $P$, otrzymujemy tożsamość
dywizorów
$K_C\sim f^*K_B+\sum_P(e_P-1)P$.

Wyjaśnijmy teraz dokładniej, skąd bierze się
$\deg(f^*D)=n\deg D$. Niestały morfizm
rzutowych nierozkładalnych krzywych jest skończony:
jego włókna są skończone (w przeciwnym razie
cała krzywa byłaby włóknem), a morfizm właściwy
o włóknach skończonych jest skończony.
Dla punktu $Q\in B$ jego pierścień lokalny
$R=\mathcal O_{B,Q}$ jest DVR z parametrem $u$.
Algebra funkcji na przeciwobrazie
$\Spec R$ jest skończonym, beztorsyjnym
$R$-modułem $S$. Jako moduł nad DVR,
a więc nad pierścieniem ideałów głównych, jest
wolna. Jej ranga wynosi
$[k(C):k(B)]=n$, stąd
$\dim_k(S/uS)=n$. Z drugiej strony skończona algebra $S/uS$ jest
artinowska. Twierdzenie o rozkładzie skończonej algebry
artinowskiej na czynniki lokalne daje jej rozkład
na produkt pierścieni lokalnych odpowiadających
punktom włókna nad $Q$. W każdym takim
punkcie $P$ lokalny parametr $t_P$ spełnia
$u=t_P^{e_P}v$ dla jednostki $v$, więc
$\dim_k(\mathcal O_{C,P}/u\mathcal O_{C,P})=e_P$.
Sumując wymiary, otrzymujemy
$\sum_{P\in f^{-1}(Q)}e_P=n$.
Jest to dokładnie wzór na stopień przeciwobrazu
punktu. Przez liniowość
$\deg(f^*D)=n\deg D$ dla dowolnego dywizora.
Zastosowanie stopnia do tożsamości dywizorów
i wzoru $\deg K=2g-2$ daje tezę.
\end{proof}

\begin{przyklad}[Dlaczego krzywe hipereliptyczne mają zadany rodzaj?]
Niech $k=\C$ i $p(x)\in\C[x]$ będzie wielomianem
bez pierwiastków wielokrotnych stopnia $m\ge5$.
Niech $C$ będzie gładkim modelem rzutowym
krzywej afinicznej $y^2=p(x)$.
Odwzorowanie $(x,y)\mapsto x$ przedłuża się
do morfizmu $f:C\to\PP^1$ stopnia $2$.
Nad każdym pierwiastkiem $p$ występuje jedno
proste rozgałęzienie (lokalnie $x-a$ jest
kwadratem parametru). Jeśli $m$ jest nieparzyste,
dodatkowy punkt rozgałęzienia występuje
w nieskończoności; jeśli $m$ jest parzyste,
nie ma tam rozgałęzienia. Łączna liczba punktów
rozgałęzienia jest więc równa $m+1$ dla
nieparzystego $m$ i $m$ dla parzystego $m$.

Z twierdzenia~\ref{tw:RH} dostajemy
$2g(C)-2=2(-2)+r=-4+r$, gdzie $r$ oznacza
liczbę punktów rozgałęzienia. Stąd
$g(C)=(r-2)/2$.
Na przykład dla $\deg p=5$ albo $6$ mamy
$r=6$ i $g(C)=2$.
Tak otrzymujemy jawne równania krzywych o rodzaju
większym niż jeden.
\end{przyklad}


% ==============================================================
% Iteracja następująca po audycie: przykłady i residua.
% ==============================================================
\section{Residua na prostej rzutowej: pełny rachunek lokalny i globalny}
\label{sec:residua-p1}

Dotychczas wprowadziliśmy dualność Serre'a dla dowolnej gładkiej
krzywej jako osobne twierdzenie. Zanim spróbujemy posługiwać się nią
w pełnej ogólności, przeprowadźmy kompletny rachunek w najprostszym
przypadku: na $\PP^1$. Przekonamy się, że parowanie dualności można
\emph{zobaczyć} bezpośrednio jako wybranie jednego współczynnika
szeregu Laurenta. Wyjaśni to jednocześnie, dlaczego różniczki mają
na prostej rzutowej stopień $-2$.

\begin{definicja}[Residuum różniczki wymiernej w punkcie skończonym]
Niech $k$ będzie ciałem charakterystyki zero,
$t$ współrzędną na $\A^1_k\subset\PP^1_k$, $a\in k$,
a $\omega=R(t)\,dt$, gdzie $R\in k(t)$.
W ciele szeregów Laurenta $k((u))$, $u=t-a$, zapisujemy
$R(a+u)=\sum_{j\ge-N}c_j u^j$.
Definiujemy
$\operatorname{Res}_{a}(\omega):=c_{-1}$.
Jest to dobrze określone, bo mianownik $R(a+u)$ można
przedstawić jako $u^m q(u)$ z $q(0)\ne0$;
$1/q(u)$ istnieje jako formalny szereg potęgowy.
\end{definicja}

\begin{przyklad}[Residua obliczane bez całkowania]
Niech $\omega=\dfrac{2t+1}{t(t-1)}\,dt$ nad $\C$.
Rozkład na ułamki proste daje
$\dfrac{2t+1}{t(t-1)}=-\dfrac1t+\dfrac3{t-1}$.
Zatem $\operatorname{Res}_0(\omega)=-1$,
$\operatorname{Res}_1(\omega)=3$.
Sam fakt, że funkcja ma biegun, nie przesądza o niezerowym residuum:
$\operatorname{Res}_0(dt/t^2)=0$, bo brakuje składnika $t^{-1}dt$.
\end{przyklad}

\begin{definicja}[Residuum w nieskończoności]
Niech $\omega=R(t)\,dt$ będzie wymierną różniczką
na $\PP^1_k$. W mapie $\PP^1\setminus\{0\}$ wybieramy
współrzędną $s=1/t$. Ponieważ $dt=-s^{-2}ds$,
określamy
$\operatorname{Res}_{\infty}(\omega)$ jako współczynnik
przy $s^{-1}ds$ w szeregu Laurenta różniczki
$-R(s^{-1})s^{-2}ds$ w punkcie $s=0$.
\end{definicja}

\begin{przyklad}[Znaki w nieskończoności]
Dla $\omega=dt/t$ mamy $\omega=-ds/s$, więc
$\operatorname{Res}_0(\omega)=1$ i
$\operatorname{Res}_{\infty}(\omega)=-1$.
Dla różniczki z poprzedniego przykładu
$\operatorname{Res}_{\infty}(\omega)=-2$.
To drugie można także odczytać z rozwinięcia
$R(t)=2t^{-1}+O(t^{-2})$ dla $t\to\infty$:
wyrażenie $R(1/s)(-s^{-2})ds$ zaczyna się od
$-2s^{-1}ds$.
\end{przyklad}

\begin{twierdzenie}[Suma residuów na $\PP^1$]
\label{tw:sum-res-p1}
Niech $k$ będzie ciałem algebraicznie domkniętym charakterystyki
zero. Dla każdej wymiernej różniczki $\omega$ na
$\PP^1_k$ zachodzi
\[
 \sum_{P\in\PP^1(k)}\operatorname{Res}_{P}(\omega)=0.
\]
W powyższej sumie tylko skończenie wiele składników jest niezerowych.
\end{twierdzenie}
\begin{proof}
Każda wymierna różniczka ma postać $R(t)\,dt$, gdzie
$R(t)\in k(t)$. Rozłóżmy $R$ na ułamki proste:
\[
 R(t)=Q(t)+\sum_{i=1}^{r}\sum_{m=1}^{n_i}
 \frac{c_{im}}{(t-a_i)^m},
\]
gdzie $Q\in k[t]$, punkty $a_i$ są parami różne,
a wszystkie sumy są skończone. Residuum w skończonym
punkcie $a_i$ wynosi $c_{i1}$; inne punkty skończone
nie wnoszą nic. Wystarczy więc obliczyć residuum w nieskończoności
trzech typów składników.

Dla $Q(t)\,dt$ po podstawieniu $t=1/s$ dostajemy
$-Q(1/s)s^{-2}ds$, w którym występują tylko potęgi
$s^j$ dla $j\le-2$; residuum jest więc zerowe.
Dla $m\ge2$ rozwinięcie
\[
 \frac{dt}{(t-a)^m}
 =-\frac{s^{m-2}}{(1-as)^m}\,ds
\]
ma jedynie potęgi nieujemne, więc też nie daje residuum.
Natomiast dla $m=1$ mamy
$dt/(t-a)=-\bigl(s^{-1}+a+a^2s+\cdots\bigr)ds$;
residuum wynosi $-1$. Stąd
$\operatorname{Res}_\infty(\omega)=-\sum_{i=1}^r c_{i1}$,
co dokładnie równoważy residua w skończonych punktach.
\end{proof}

\begin{uwaga}[Dlaczego ten rachunek nie jest jeszcze dowodem
ogólnego twierdzenia o residuach?]
Użyliśmy współrzędnej wymiernej $t$, która parametryzuje
całą prostą rzutową. Na gładkiej rzutowej krzywej dodatniego rodzaju taka
współrzędna nie istnieje: gdyby istniała, krzywa
byłaby biwymierna z $\PP^1$, a więc miałaby rodzaj zero. Ogólne twierdzenie o residuach i
pełna dualność Serre'a wymagają dodatkowych argumentów globalnych.
Poniższy dowód dualności dotyczy zatem \emph{dokładnie $\PP^1$},
nie dowolnej krzywej.
\end{uwaga}

\section[Dualność Serre'a na prostej rzutowej]{Dualność Serre'a na $\PP^1$: parowanie policzone w bazach}
\label{sec:serre-p1}

Z sekcji~\ref{sec:pp1-rr} wiemy, że kohomologię
$\mathcal O(n)$ można obliczyć za pomocą wielomianów Laurenta.
Pokażemy teraz, jak z tych samych danych zbudować
\emph{niezdegenerowane parowanie}, stanowiące szczególny
przypadek dualności Serre'a.

\begin{twierdzenie}[Jawna dualność na prostej rzutowej]
\label{tw:serre-p1-explicit}
Niech $k$ będzie ciałem charakterystyki zero,
$n\in\Z$, a $\omega_{\PP^1}:=\Omega^1_{\PP^1/k}
\cong\mathcal O(-2)$. Istnieje naturalne parowanie
\[
 H^1(\PP^1,\mathcal O(n))\times
 H^0(\PP^1,\mathcal O(-n-2))\longrightarrow k
\]
wyrażone na przedstawicielach Čecha wzorem
\[
 ([\alpha],\beta)\longmapsto
 \operatorname{Res}_{t=0}\bigl(\alpha(t)\beta(t)\,dt\bigr).
\]
Jest ono dobrze określone i niezdegenerowane.
\end{twierdzenie}
\begin{proof}
Przypomnijmy dokładnie używane trywializacje.
Na $U_0=\Spec k[t]$ wybieramy bazę $e_0$ wiązki
$\mathcal O(n)$; na $U_\infty=\Spec k[s]$, $s=t^{-1}$,
wybieramy $e_\infty=t^n e_0$. Klasa z $H^1$ jest
przedstawiana przez wielomian Laurenta $\alpha(t)$ modulo
podprzestrzeń $k[t]+t^n k[t^{-1}]$.
Z kolei sekcja globalna $\mathcal O(-n-2)$,
gdy istnieje, jest przedstawiana na $U_0$ przez
wielomian $\beta(t)$ stopnia co najwyżej $-n-2$.

Sprawdźmy najpierw niezależność od wyboru przedstawiciela.
Dodanie do $\alpha(t)$ wielomianu $a(t)\in k[t]$
nie zmienia residuum, ponieważ $a(t)\beta(t)$
nie zawiera ujemnych potęg $t$. Jeśli natomiast do
$\alpha$ dodamy $t^n b(t^{-1})$, gdzie
$b\in k[t^{-1}]$, to iloczyn
$t^n b(t^{-1})\beta(t)$ ma wykładniki co najwyżej
$n+(-n-2)=-2$. Również nie ma składnika $t^{-1}$.
Parowanie jest zatem dobrze określone na ilorazie.

Jeśli $n\ge-1$, obie przestrzenie występujące
w tezie są zerowe, więc niezdegenerowaność zachodzi
w sposób pusty. Jeśli $n\le-2$, oznaczmy $r=-n-1$.
Z dowodu twierdzenia~\ref{tw:p1-coh} wynika, że
pierwsza przestrzeń ma bazę klas
\[
 [t^{n+1}],[t^{n+2}],\ldots,[t^{-1}],
\]
zaś druga ma bazę $1,t,\ldots,t^{r-1}$.
Dla $j\in\{n+1,\ldots,-1\}$ oraz
$i\in\{0,\ldots,r-1\}$ parowanie bazowych
wektorów wynosi $1$ dokładnie wtedy, gdy
$j+i=-1$, a w przeciwnym razie $0$.
Macierz parowania jest zatem antydiagonalną macierzą
permutacji: w każdym wierszu i każdej kolumnie występuje
jedna jedynka. Jej wyznacznik wynosi $\pm1\ne0$,
co dowodzi niezdegenerowaności.
\end{proof}

\begin{przyklad}[Cały rachunek dla $n=-4$]
Mamy
$H^1(\PP^1,\mathcal O(-4))
=\operatorname{span}\{[t^{-3}],[t^{-2}],[t^{-1}]\}$
oraz $H^0(\PP^1,\mathcal O(2))
=\operatorname{span}\{1,t,t^2\}$.
W tych bazach macierz parowania to
\[
 \begin{pmatrix}0&0&1\\0&1&0\\1&0&0\end{pmatrix}.
\]
Na przykład klasa $[t^{-3}]$ jest wykrywana
przez sekcję $t^2$, a $[t^{-1}]$ przez sekcję $1$.
Takiej klasy nie można ``zgubić'' przy sklejaniu:
każda niezerowa klasa ma niezerowe parowanie z którąś
sekcją dualnej wiązki.
\end{przyklad}

\begin{uwaga}[Dlaczego pojawia się $-2$?]
W przejściu $s=t^{-1}$ mamy $ds=-t^{-2}dt$.
Różniczki przekształcają się więc z czynnikiem
$t^{-2}$, co odpowiada wiązce kanonicznej
$\mathcal O(-2)$. Dwa wykładniki $n$ oraz $-n-2$
są komplementarne dokładnie ze względu na ten dodatkowy
czynnik. Nie jest on arbitralną normalizacją.
\end{uwaga}

\section{Krzywa eliptyczna: funkcje o zadanych biegunach}
\label{sec:elliptic-poles}

Prawo grupowe krzywej eliptycznej można badać
za pomocą przecinania kubiki prostymi. Chcemy jednak
zrozumieć jeszcze inną, zaskakującą własność:
\emph{na krzywej eliptycznej nie istnieje niezerowa
funkcja wymierna z dokładnie jednym prostym biegunem}.
Z kolei funkcje z biegunami rzędów dwa i trzy istnieją.
To zjawisko wynika z rodzaju krzywej.

\begin{twierdzenie}[Wymiary $L(mO)$ na krzywej rodzaju jeden]
\label{tw:ell-poles}
Niech $k$ będzie ciałem algebraicznie domkniętym,
$E$ gładką nierozkładalną krzywą rzutową rodzaju jeden,
a $O\in E(k)$. Wówczas
$\ell(0)=1$, $\ell(mO)=m$ dla każdego $m\ge1$,
oraz $\ell(mO)=0$ dla $m<0$.
\end{twierdzenie}
\begin{proof}
Dla $m<0$ stopień dywizora $mO$ jest ujemny,
więc $L(mO)=0$ na mocy wniosku po Riemannie--Rochu.
Dla $m=0$ stałe są wszystkimi globalnymi funkcjami
regularnymi na krzywej rzutowej, stąd $\ell(0)=1$.
Jeśli $m\ge1$, to $\deg(K_E-mO)=-m<0$,
ponieważ $\deg K_E=2g(E)-2=0$.
Wobec tego $\ell(K_E-mO)=0$ i Riemann--Roch daje
$\ell(mO)=m+1-g(E)=m$.
\end{proof}

\begin{wniosek}[Luka w możliwych rzędach biegunów]
Dla krzywej $(E,O)$ z twierdzenia~\ref{tw:ell-poles}
nie istnieje niezerowa funkcja $f\in k(E)$, której
\emph{jedyny} biegun miałby rząd $1$ w punkcie $O$.
Natomiast istnieją funkcje z jedynym biegunem
w $O$ dowolnego rzędu $m\ge2$.
\end{wniosek}
\begin{proof}
Gdyby $f$ miała jedyny prosty biegun w $O$,
$1,f$ byłyby liniowo niezależnymi elementami $L(O)$,
co przeczy $\ell(O)=1$.
Dla $m\ge2$ mamy
$L((m-1)O)\subsetneq L(mO)$, ponieważ
$\ell(mO)=m$ oraz $\ell((m-1)O)=m-1$
(także dla $m=2$). Wybieramy
$f\in L(mO)\setminus L((m-1)O)$.
Jej wszystkie bieguny mogą leżeć wyłącznie w $O$,
a rząd bieguna w $O$ musi być dokładnie równy $m$.
\end{proof}

\begin{center}
\begin{tikzpicture}[>=Stealth,x=1.08cm,y=.9cm]
\draw[->,gray!80] (-.55,0)--(7.8,0) node[right] {rząd bieguna};
\foreach \j in {0,...,7} {
  \draw (\j,.08)--(\j,-.08);
  \node[below,font=\small] at (\j,-.1) {$\j$};
}
\foreach \j in {2,...,7} {\fill[blue!70!black] (\j,0) circle (2.7pt);}
\fill[red!70!black] (1,0) circle (3pt);
\node[red!70!black,anchor=south,font=\small] at (1,.27) {niemożliwe};
\node[blue!70!black,anchor=south,font=\small] at (4.5,.27)
 {rzędy $2,3,4,\ldots$};
\end{tikzpicture}

{\small Rysunek 4.5. W punkcie $O$ na krzywej rodzaju jeden
rząd $1$ jest luką w zbiorze możliwych rzędów jedynego bieguna.}
\end{center}

\begin{przyklad}[Ciało funkcji krzywej Weierstrassa: konkretna baza]
\label{prz:weierstrass-LmO}
Załóżmy $\operatorname{char}k\ne2,3$ i niech $E$ ma równanie
$y^2=x^3+ax+b$, przy $4a^3+27b^2\ne0$.
Jej domknięcie rzutowe ma równanie
$Y^2Z=X^3+aXZ^2+bZ^3$ i jeden punkt
w nieskończoności $O=[0:1:0]$.

W mapie $Y\ne0$ zapisujemy $u=X/Y$, $v=Z/Y$.
Równanie krzywej ma w tych współrzędnych postać
\[
 v=u^3+auv^2+bv^3.
\]
W punkcie $O$, gdzie $(u,v)=(0,0)$,
pochodna lewej strony minus prawej po $v$
wynosi $1$. Z kryterium Jacobiego i rachunku
w pierścieniu lokalnym wynika, że $u$ jest
parametrem lokalnym. Ponieważ
$v(1-auv-bv^2)=u^3$, a drugi czynnik ma w $O$
wartość $1$, jest jednostką; zatem
$\operatorname{ord}_O(v)=3$.
Wracając do współrzędnych afinicznych,
$x=X/Z=u/v$ i $y=Y/Z=1/v$, skąd
$\operatorname{ord}_O(x)=-2$ oraz
$\operatorname{ord}_O(y)=-3$.
Ponieważ $x,y$ są regularne w części afinicznej $E\setminus\{O\}$,
to ich jedyne bieguny znajdują się w $O$.

Dla dowolnego $m\ge1$ rozważmy zbiór jednomianów
\[
 \mathcal B_m=
 \{x^i: i\ge0,\ 2i\le m\}\ \cup\
 \{x^j y:j\ge0,\ 2j+3\le m\}.
\]
Każdy element należy do $L(mO)$. Elementy te mają
parami różne rzędy bieguna w $O$ (parzyste w pierwszej
rodzinie, nieparzyste w drugiej), dlatego są liniowo
niezależne: w niezerowej kombinacji liniowej składnik
o największym rzędzie bieguna nie może się skrócić z innymi.
Dla $m=1$ zbiór zawiera tylko $1$, a dla $m\ge2$
jego liczność wynosi $m$. Poprzednie twierdzenie mówi,
że $\dim_kL(mO)=m$, a więc $\mathcal B_m$ jest bazą.

Konkretnie:
\[
\begin{array}{c|l|c}
 m & \text{baza }L(mO) & \ell(mO)\\ \hline
 1 & 1 & 1\\
 2 & 1,x & 2\\
 3 & 1,x,y & 3\\
 4 & 1,x,y,x^2 & 4\\
 5 & 1,x,y,x^2,xy & 5\\
 6 & 1,x,y,x^2,xy,x^3 & 6
\end{array}
\]
Tutaj wyrażenie $y^2$ nie jest potrzebnym nowym generatorem,
gdyż równanie Weierstrassa zastępuje je przez $x^3+ax+b$.
To pokazuje, jak geometria biegunów i algebra pierścienia
współrzędnych wzajemnie się uzupełniają.
\end{przyklad}

\section{Rodzaj, rozgałęzienia i nowe przykłady z wyliczeniem}
\label{sec:ramification-examples}


\begin{center}
\begin{minipage}{.96\linewidth}
\centering
\begin{tikzpicture}[>=Stealth,every node/.style={font=\small}]
 \draw[gray!70,thick] (0,0) circle (1.1);
 \draw[gray!70,thick] (5.5,0) circle (1.1);
 \fill[blue!70!black] (.62,.36) circle (2.6pt);
 \fill[blue!70!black] (-.62,-.36) circle (2.6pt);
 \fill[red!75!black] (5.755,.446) circle (2.6pt);
 \fill[black] (0,0) circle (2pt);
 \fill[black] (5.5,0) circle (2pt);
 \draw[->,thick] (1.4,0)--(4.05,0)
     node[midway,above] {$w=z^2$};
 \node[below=1.15cm,align=center] at (0,-.2) {dziedzina: $z$ i $-z$};
 \node[below=1.15cm,align=center] at (5.5,-.2) {obraz: $w$};
 \node[above=1.18cm] at (0,.2) {$z$-płaszczyzna};
 \node[above=1.18cm] at (5.5,.2) {$w$-płaszczyzna};
\end{tikzpicture}

{\small Rysunek 4.6. Lokalny model rozgałęzienia indeksu $2$:
poza początkiem dwa punkty $z$ i $-z$ mają ten sam obraz,
natomiast nad $w=0$ pozostaje jeden punkt $z=0$.}
\end{minipage}
\end{center}

\begin{przyklad}[Podwójne nakrycie prostej z czterema punktami rozgałęzienia]
Niech $a,b,c,d\in\C$ będą parami różne, a $C$ gładkim
modelem rzutowym krzywej
$y^2=(x-a)(x-b)(x-c)(x-d)$.
Funkcja $x$ wyznacza morfizm $f:C\to\PP^1$ stopnia $2$.
Nad każdym z czterech pierwiastków po prawej stronie
jest jeden punkt o indeksie rozgałęzienia $2$.
W nieskończoności są dwa punkty nierozgałęzione,
gdyż stopień wielomianu jest parzysty.
Formuła Riemanna--Hurwitza daje
$2g(C)-2=2(-2)+4=0$, stąd $g(C)=1$.
W ten sposób otrzymujemy krzywą rodzaju jeden także
z równania o prawej stronie stopnia czwartego.
Po wybraniu jednego z jej punktów jako $O$
staje się ona krzywą eliptyczną.
\end{przyklad}

\begin{przyklad}[Dlaczego do równań osobliwych używamy normalizacji?]
Rozważmy $y^2=x^5$ nad $\C$.
Parametryzacja $x=t^2$, $y=t^5$ jest bijektywna
na punktach zespolonych, ale nie jest izomorfizmem
krzywych afinicznych, ponieważ $t$ nie należy do
$\C[t^2,t^5]$ (w tym pierścieniu występują tylko potęgi
$t$ z wykładnikami w półgrupie generowanej przez $2$ i $5$).
Jednocześnie $t=t^5/(t^2)^2$ należy do ciała ułamków i jest
całkowity nad tym pierścieniem, ponieważ spełnia
$T^2-t^2=0$. Domknięcie całkowite wynosi $\C[t]$
(uzasadnienie jak dla normalizacji ostrza w~\ref{prz:normalizacja-ostrza}).
Domknięcie rzutowe danej krzywej jest osobliwe;
formuły Riemanna--Hurwitza i Riemanna--Rocha
w postaci używanej w rozdziale dotyczą \emph{gładkiego
modelu rzutowego}. Przez ciało funkcji $\C(t)$ widzimy,
że jego gładkim modelem rzutowym jest $\PP^1$,
zatem rodzaj geometryczny wynosi $0$. Wykorzystujemy tu
standardowy fakt, że gładkie uzupełnienie rzutowe
normalizacji $\A^1$ jest $\PP^1$; pełnej teorii
gładkich modeli rzutowych nie rozwijamy w tym rozdziale.
Natomiast rodzaj arytmetyczny płaskiej krzywej
stopnia $5$ wynosi $(5-1)(5-2)/2=6$.
Nie ma sprzeczności: różnicę między tymi
rodzajami powodują osobliwości modelu płaskiego.
\end{przyklad}

\begin{uwaga}[Co zostało udowodnione, a co nadal jest narzędziem zewnętrznym?]
Pełne, bezpośrednie dowody z tej iteracji obejmują
sumę residuów na $\PP^1$, jawne parowanie dualności
na $\PP^1$ oraz obliczenia przestrzeni funkcji
$L(mO)$ przy przyjęciu ogólnego twierdzenia
Riemanna--Rocha. Nadal \emph{nie twierdzimy}, że
udowodniliśmy od podstaw ogólną dualność Serre'a
ani twierdzenie o skończoności kohomologii koherentnej.
Uczciwe domknięcie tych fundamentów wymagałoby
osobnego działu o kohomologii snopów na schematach,
który wykracza poza założone minimum teorii krzywych.
\end{uwaga}

% ==============================================================
% Iteracja 6: od konstrukcji Cecha do twierdzenia o residuach
% i jawnego parowania na gladkich krzywych zespolonych.
% ==============================================================
\section{Kohomologia Čecha od podstaw: pełny kompleks i jego sens}
\label{sec:cech-full}

W sekcji~\ref{sec:cech-curve} użyliśmy pokrycia dwiema częściami
otwartymi. Zanim przejdziemy do ogólnych krzywych, wyjaśnijmy,
co naprawdę kryje się za symbolem $H^1$, co trzeba udowodnić,
żeby wolno było liczyć go metodą Čecha, oraz dlaczego lokalne
rozwiązania równań nie muszą się sklejać globalnie.

\begin{definicja}[Kompleks Čecha skończonego pokrycia]
Niech $X$ będzie przestrzenią topologiczną, $\mathcal F$
snopem grup abelowych na $X$, a
$\mathfrak U=(U_0,\ldots,U_r)$ skończonym pokryciem
otwartym. Dla $i_0<\cdots<i_q$ piszemy
$U_{i_0\ldots i_q}=U_{i_0}\cap\cdots\cap U_{i_q}$.
Niech
\[
 C^q(\mathfrak U,\mathcal F)
 =\prod_{0\le i_0<\cdots<i_q\le r}
       \mathcal F(U_{i_0\ldots i_q})\qquad(q\ge0).
\]
Element $c\in C^q$ jest rodziną sekcji na przecięciach
$(q+1)$ części pokrycia. Różniczkę
$d^q:C^q\to C^{q+1}$ definiujemy wzorem
\[
 (d^qc)_{i_0\ldots i_{q+1}}=
 \sum_{j=0}^{q+1}(-1)^j
 c_{i_0\ldots\widehat{i_j}\ldots i_{q+1}}
 \big|_{U_{i_0\ldots i_{q+1}}},
\]
gdzie daszek oznacza pominięcie indeksu. Ustalamy
$\check H^0(\mathfrak U,\mathcal F)=\ker d^0$ oraz
$\check H^q(\mathfrak U,\mathcal F)
=\ker d^q/\operatorname{im}d^{q-1}$ dla $q\ge1$.
\end{definicja}

\begin{lemat}[Dlaczego $d^{q+1}d^q=0$?]
Dla każdego $q\ge0$ złożenie kolejnych różniczek
kompleksu Čecha jest odwzorowaniem zerowym.
\end{lemat}
\begin{proof}
Ustalmy indeksy $i_0<\cdots<i_{q+2}$. W rozwinięciu
$(d^{q+1}d^qc)_{i_0\ldots i_{q+2}}$ każda sekcja
otrzymana przez usunięcie dwóch różnych indeksów
$i_a,i_b$ występuje dwukrotnie. Jeśli $a<b$, usunięcie
najpierw $i_a$, a potem $i_b$ daje znak
$(-1)^a(-1)^{b-1}$; w odwrotnej kolejności otrzymujemy
$(-1)^b(-1)^a$. Są to liczby przeciwne. Ograniczenia
sekcji na wspólne przecięcie są takie same, bo
ograniczanie jest zgodne ze składaniem. Wszystkie składniki
znoszą się parami. Właśnie to umożliwia tworzenie ilorazu
$\ker d^q/\operatorname{im}d^{q-1}$.
\end{proof}

\begin{stwierdzenie}[Sekcje globalne jako zerowa kohomologia]
Dla każdego snopa grup abelowych $\mathcal F$ i dowolnego
otwartego pokrycia $\mathfrak U$ przestrzeni $X$ zachodzi
kanoniczny izomorfizm
$\check H^0(\mathfrak U,\mathcal F)\cong\mathcal F(X)$.
\end{stwierdzenie}
\begin{proof}
Równanie $d^0c=0$ oznacza dokładnie
$c_j|_{U_i\cap U_j}=c_i|_{U_i\cap U_j}$ dla dowolnych
$i,j$. Aksjomat sklejania snopa mówi, że taka rodzina
pochodzi od dokładnie jednej sekcji globalnej.
\end{proof}

\begin{przyklad}[Trzy części pokrycia i warunek na trójkrotne przecięcie]
Niech $X=U_0\cup U_1\cup U_2$. 1-kołańcuchem jest
trójka $(g_{01},g_{02},g_{12})$ sekcji określonych
na odpowiednich parach. Warunek $dg=0$ brzmi
$g_{12}-g_{02}+g_{01}=0$ na $U_0\cap U_1\cap U_2$.
Nie jest to warunek dowolny: jeśli $g_{ij}=f_j-f_i$
pochodzi z lokalnych sekcji $f_i$, to
$(f_2-f_1)-(f_2-f_0)+(f_1-f_0)=0$.
Kohomologia $H^1$ rozróżnia więc układy lokalnych
\emph{różnic}, których nie można w ten sposób wyrazić,
od tych, które można.
\end{przyklad}

\begin{center}
\begin{tikzpicture}[scale=.9,every node/.style={font=\small}]
  \fill[blue!8] (-1.5,0) ellipse (1.65 and 1.05);
  \fill[red!8] (1.5,0) ellipse (1.65 and 1.05);
  \draw[blue!70!black,thick] (-1.5,0) ellipse (1.65 and 1.05);
  \draw[red!70!black,thick] (1.5,0) ellipse (1.65 and 1.05);
  \node[blue!70!black] at (-1.8,.3) {$U$};
  \node[red!70!black] at (1.8,.3) {$V$};
  \draw[->,thick] (0,-1.42)--(0,-.18);
  \node[below,align=center] at (0,-1.52)
     {$U\cap V$: sekcja $g$\\[1pt]czy $g=f_U-f_V$?};
\end{tikzpicture}

{\small Rysunek 4.7. Problem kohomologiczny: lokalnie dane są
sekcje, lecz globalne sklejenie wymaga zgodności na przecięciach.}
\end{center}

\begin{uwaga}[Dlaczego nie wystarczy policzyć dowolnego pokrycia?]
Kompleks Čecha zależy od pokrycia. Dopiero po udowodnieniu
odpowiednich twierdzeń porównawczych można utożsamić jego
kohomologię z kohomologią snopa $H^q(X,\mathcal F)$.
Na schematach rozdzielonych pokrycia afiniczne są do tego
przydatne, gdy $\mathcal F$ jest quasi-koherentny:
wszystkie skończone przecięcia otwartych afinicznych są
afiniczne i wyższe kohomologie na nich znikają.
Nie należy natomiast automatycznie utożsamiać grup
$\check H^q$ z $H^q$ dla dowolnego snopa i dowolnego pokrycia.
\end{uwaga}

\section{Kluczowy rachunek afiniczny: dlaczego lokalizacje sklejają się dokładnie?}
\label{sec:affine-cech-proof}

W geometrii afinicznej sekcje $\widetilde M$ odpowiadają
ułamkom modułu $M$. Udowodnimy dokładność kompleksu
występującego przy pokryciu przez podstawowe otwarte.
Jest to rachunkowy fundament twierdzenia o znikaniu
kohomologii snopów quasi-koherentnych na schematach afinicznych.

\begin{twierdzenie}[Dokładny kompleks lokalizacji]
\label{tw:affine-cech-exact}
Niech $A$ będzie pierścieniem przemiennym z jedynką,
$M$ dowolnym $A$-modułem i $f_0,\ldots,f_r\in A$
takimi, że $(f_0,\ldots,f_r)=A$. Wtedy kompleks
\[
0\longrightarrow M\xrightarrow{\epsilon}
\prod_iM_{f_i}\xrightarrow{d^0}
\prod_{i<j}M_{f_if_j}\xrightarrow{d^1}\cdots
\xrightarrow{d^{r-1}}M_{f_0\cdots f_r}\longrightarrow0
\]
jest dokładny; $\epsilon(m)$ jest rodziną obrazów
$m$ we wszystkich lokalizacjach, a różniczki są
określone wzorem z poprzedniej sekcji.
\end{twierdzenie}
\begin{proof}
Pokażemy najpierw, że dokładność ciągu modułów nad
$A$ można sprawdzić po lokalizacji we wszystkich
ideałach pierwszych. Jeśli $N$ jest niezerowym
$A$-modułem, wybieramy $0\ne n\in N$ i ideał
maksymalny $\mathfrak m$ zawierający właściwy ideał
$\operatorname{Ann}(n)=\{a\in A:an=0\}$.
Wówczas $n/1\ne0$ w $N_{\mathfrak m}$:
gdyby $n/1=0$, pewien $s\notin\mathfrak m$
spełniałby $sn=0$, co przeczy zawieraniu
$\operatorname{Ann}(n)\subseteq\mathfrak m$.
W zastosowaniu do jąder i ilorazów wynika stąd,
że kompleks modułów jest dokładny wtedy i tylko wtedy,
gdy pozostaje dokładny po lokalizacji w każdym
ideale pierwszym (lokalizacja jest funktorem dokładnym;
jak wykazano wcześniej przy konstrukcji lokalizacji).

Ustalmy więc $\mathfrak p\in\Spec A$. Ponieważ
$(f_0,\ldots,f_r)=A$, istnieje indeks $j$ z
$f_j\notin\mathfrak p$. W pierścieniu $A_{\mathfrak p}$
element $f_j$ jest jednostką. Po lokalizacji całego
kompleksu w $\mathfrak p$ wprowadźmy dodatkowe działanie
$h:C^q\to C^{q-1}$, które do składowej indeksowanej
przez $(i_0,\ldots,i_{q-1})$ przypisuje składową
indeksowaną przez $(j,i_0,\ldots,i_{q-1})$ z odpowiednim
znakiem permutacji; jeśli indeks $j$ już występuje,
kładziemy tę składową równą zeru. Odwzorowanie jest
poprawnie określone, ponieważ dołączenie $f_j$
nie zmienia lokalizacji po przejściu do $A_{\mathfrak p}$.
Dla $q=0$ przyjmujemy $h:C^0\to M_{\mathfrak p}$,
$h((m_i))=m_j$.

Bezpośrednio z naprzemiennego wzoru na różniczkę
wynikają tożsamości
$h\epsilon=\operatorname{id}$ na $M_{\mathfrak p}$
oraz $dh+hd=\operatorname{id}$ na każdym $C^q$.
Można je sprawdzić, rozpisując składową o ustalonych
indeksach: składnik otrzymany przez usunięcie
wprowadzonego indeksu $j$ jest wyjściową składową,
a wszystkie inne występują parami z przeciwnymi znakami.
Zatem zlokalizowany kompleks jest ściągalny, a więc
dokładny. Stwierdzona na początku lokalna zasada
dokładności kończy dowód.
\end{proof}

\begin{wniosek}[Znikanie wyższej kohomologii Čecha na pokryciu afinicznym]
Niech $X=\Spec A$, $\mathcal F=\widetilde M$ i
$X=\bigcup_{i=0}^rD(f_i)$. Wtedy
$\check H^q(\{D(f_i)\},\mathcal F)=0$ dla
$q>0$, a $\check H^0\cong M$.
\end{wniosek}
\begin{proof}
Dla każdego wielowskaźnika $I$ zachodzi
$\mathcal F(D(\prod_{i\in I}f_i))=M_{\prod_{i\in I}f_i}$.
Zbiory $D(f_i)$ pokrywają $\Spec A$ dokładnie wtedy,
gdy $(f_0,\ldots,f_r)=A$. Teza jest więc
przeformułowaniem twierdzenia~\ref{tw:affine-cech-exact}.
\end{proof}

\begin{uwaga}[Przejście do rzeczywistej kohomologii snopa]
Wykazaliśmy dokładność \emph{konkretnego} kompleksu
Čecha. Pełne stwierdzenie
$H^q(\Spec A,\widetilde M)=0$ dla $q>0$
uzyskuje się, łącząc rachunek lokalizacji z
porównaniem kohomologii Čecha z kohomologią
snopową na pokryciach acyklicznych. Jest to
\emph{twierdzenie o acykliczności afinicznej};
w dalszych obliczeniach używamy go jawnie,
a nie wyprowadzamy z samego faktu, że pewne sekcje
potrafimy skleić. Jego typowe dowody korzystają
z rezolwent iniektywnych lub z teorii kohomologii
Čecha i rozdrobnienia pokryć.
\end{uwaga}

\begin{przyklad}[Dlaczego funktor sekcji globalnych nie jest prawostronnie dokładny?]
Na $\PP^1_k$ rozważmy ciąg snopów
\[
0\longrightarrow\mathcal O(-2)
\xrightarrow{\,(-X_1,X_0)\,}\mathcal O(-1)^{\oplus2}
\xrightarrow{\,(X_0,X_1)\,}\mathcal O\longrightarrow0.
\]
Mapy są dobrze określone, bo mnożenie przez współrzędną
jednorodną stopnia jeden przenosi $\mathcal O(m)$ w
$\mathcal O(m+1)$. Złożenie wynosi zero.
W otoczeniu $X_0\ne0$ pierwsza współrzędna drugiej
mapy jest lokalnie odwracalna, a jądro jest generowane
przez $(-X_1,X_0)$. Analogicznie na $X_1\ne0$.
Ciąg jest zatem dokładny na źdźbłach, czyli jako
ciąg snopów. Jednak
$H^0(\PP^1,\mathcal O(-1))=0$, podczas gdy
$H^0(\PP^1,\mathcal O)=k$.
Odwzorowanie $H^0(\mathcal O(-1)^{\oplus2})\to
H^0(\mathcal O)$ jest więc zerowe, a nie surjektywne.
Brakujące sekcje są wykrywane przez
$H^1(\PP^1,\mathcal O(-2))\cong k$.
Jest to konkretny powód wprowadzania kohomologii.
\end{przyklad}

\section{Residua na zwartej powierzchni Riemanna}
\label{sec:residue-compact}

W tej sekcji pracujemy nad $\C$ i korzystamy z
analityfikacji gładkiej rzutowej krzywej. Jej
przestrzeń punktów zespolonych jest zwartą
spójną powierzchnią Riemanna $S$. Zastosowanie
zespolonej analizy jednej zmiennej pozwoli
udowodnić globalne twierdzenie o residuach bez
zakładania istnienia globalnej współrzędnej wymiernej.

\begin{definicja}[Residuum meromorficznej formy]
Niech $S$ będzie powierzchnią Riemanna, $P\in S$,
$t$ holomorficzną współrzędną lokalną z $t(P)=0$,
a $\omega$ meromorficzną $1$-formą w otoczeniu $P$.
W zapisie
$\omega=(\sum_{j\ge-N}a_jt^j)\,dt$
definiujemy $\operatorname{Res}_P\omega=a_{-1}$.
Współczynnik ten nie zależy od wyboru współrzędnej.
\end{definicja}
\begin{proof}[Uzasadnienie niezależności współrzędnej]
Dla małego dodatnio zorientowanego okręgu
$\gamma:|t|=\varepsilon$ szereg Laurenta można
całkować wyraz po wyrazie. Wobec tego
$\int_\gamma t^j\,dt=0$ dla $j\ne-1$,
a $\int_\gamma dt/t=2\pi i$. Otrzymujemy wzór
\[
\operatorname{Res}_P\omega=
 \frac1{2\pi i}\int_\gamma\omega.
\]
Dwie lokalne współrzędne prowadzą do małych
okręgów otaczających $P$ i homologicznych w
przekłutym otoczeniu, a $\omega$ jest tam
holomorficzna. Z twierdzenia Cauchy'ego całki
są równe. Tak więc definicja jest wewnętrzna,
a nie zależna od parametryzacji.
\end{proof}

\begin{twierdzenie}[Globalne twierdzenie o sumie residuów]
\label{tw:residue-general}
Niech $S$ będzie zwartą spójną powierzchnią
Riemanna, a $\omega$ meromorficzną $1$-formą na $S$.
Wtedy ma ona tylko skończenie wiele biegunów i
\[
\sum_{P\in S}\operatorname{Res}_P\omega=0.
\]
\end{twierdzenie}
\begin{proof}
Bieguny formy meromorficznej są dyskretne.
Ich zbiór jest domknięty w zwartej przestrzeni $S$,
zatem jest skończony; oznaczmy go $P_1,\ldots,P_s$.
Wybierzmy wokół tych punktów parami rozłączne,
małe domknięte dyski współrzędnych $D_j$,
takie że na każdej krzywej brzegowej $\partial D_j$
forma $\omega$ jest holomorficzna.
Niech $S^\circ=S\setminus\bigcup_j\operatorname{int}D_j$.
Na zwartej rozmaitości z brzegiem $S^\circ$
forma $\omega$ jest holomorficzna, toteż zamknięta:
$\mathrm d\omega=0$. Z twierdzenia Stokesa
otrzymujemy
\[
0=\int_{S^\circ}\mathrm d\omega
 =\int_{\partial S^\circ}\omega
 =-\sum_{j=1}^s\int_{\partial D_j}\omega
 =-2\pi i\sum_{j=1}^s\operatorname{Res}_{P_j}\omega.
\]
Znak minus wynika stąd, że dodatnia orientacja
brzegu powierzchni z \emph{wyciętymi} dyskami
jest przeciwna do dodatniej orientacji okręgów
otaczających bieguny. Podzielenie przez $-2\pi i$
daje tezę.
\end{proof}

\begin{center}
\begin{tikzpicture}[scale=.9,every node/.style={font=\small},>=Stealth]
  \fill[blue!6] (0,0) ellipse (3.05 and 1.65);
  \draw[blue!70!black,thick] (0,0) ellipse (3.05 and 1.65);
  \fill[white] (-1.25,.1) circle (.38);
  \fill[white] (1.25,-.15) circle (.38);
  \draw[red!75!black,thick] (-1.25,.1) circle (.38);
  \draw[red!75!black,thick] (1.25,-.15) circle (.38);
  \fill[red!75!black] (-1.25,.1) circle (1.4pt);
  \fill[red!75!black] (1.25,-.15) circle (1.4pt);
  \node[above] at (-1.25,.53) {$P_1$};
  \node[above] at (1.25,.32) {$P_2$};
  \node at (0,1.08) {$S^\circ$};
  \node[below,align=center] at (0,-1.85)
     {Wycięte dyski mają brzeg przeciwnie zorientowany\\
      do małych dodatnich okręgów wokół biegunów.};
\end{tikzpicture}

{\small Rysunek 4.8. Schemat idei dowodu twierdzenia o sumie residuów.
Rysunek przedstawia powierzchnię \emph{schematycznie}, a nie jako
metrycznie wierne zanurzenie powierzchni Riemanna.}
\end{center}

\begin{wniosek}[Stopień dywizora funkcji wymiernej raz jeszcze]
Jeżeli $S$ jest zwartą spójną powierzchnią
Riemanna, a $f$ jest niezerową funkcją meromorficzną,
to suma rzędów jej zer i biegunów wynosi zero:
$\sum_{P\in S}\operatorname{ord}_P(f)=0$.
\end{wniosek}
\begin{proof}
W pobliżu dowolnego punktu $P$ piszemy
$f=t^mu(t)$, gdzie $m\in\Z$ oraz $u$ jest
holomorficzna i nigdzie nie znika. Zatem
$\mathrm df/f=m\,\mathrm dt/t+\mathrm du/u$.
Drugi składnik jest holomorficzny, a pierwszy
ma residuum $m$. Z twierdzenia~\ref{tw:residue-general}
dla formy $\mathrm df/f$ wynika
$0=\sum_P\operatorname{Res}_P(\mathrm df/f)
=\sum_P\operatorname{ord}_P(f)$.
Ten dowód jest niezależny od wcześniejszego
rachunku charakterystyki Eulera.
\end{proof}

\begin{przyklad}[Krzywa eliptyczna i regularność różniczki]
Niech $E$ będzie gładką zespoloną kubiką
$y^2=x^3+ax+b$ z $4a^3+27b^2\ne0$.
Forma $\omega=\mathrm dx/(2y)$ jest holomorficzna
na całym $E$, mimo że jej zapis wygląda osobliwie
w punktach, gdzie $y=0$. W takim punkcie
$3x^2+a\ne0$ (bo równanie jest bez pierwiastków
wielokrotnych), a różniczkowanie równania daje
$2y\,\mathrm dy=(3x^2+a)\,\mathrm dx$, czyli
$\omega=\mathrm dy/(3x^2+a)$; to wyrażenie jest
regularne. W punkcie w nieskończoności użyjmy
parametrów $u=X/Y$, $v=Z/Y$ z przykładu
\ref{prz:weierstrass-LmO}. Tam $v$ ma rząd $3$
i $u$ jest parametrem lokalnym, więc
$x=u/v$ ma rząd $-2$, a $y=1/v$ rząd $-3$.
Dokładniej $v=u^3(1+O(u^4))$, skąd
$x=u^{-2}(1+O(u^4))$ oraz
$y=u^{-3}(1+O(u^4))$ i
$\omega=(-1+O(u^4))\,\mathrm du$.
Forma nie ma bieguna również w nieskończoności.
W szczególności wszystkie jej residua są zerowe.
\end{przyklad}

\section{Części główne, residua i parowanie dualności}
\label{sec:principal-residue-pairing}

\begin{definicja}[Część główna w punkcie]
Niech $C$ będzie gładką rzutową krzywą nad $\C$,
$P\in C$, a $t$ parametrem lokalnym. Częścią główną
funkcji meromorficznej $f=\sum_{j\ge-N}a_jt^j$
w $P$ nazywamy klasę
$\operatorname{pp}_P(f)=\sum_{j=-N}^{-1}a_jt^j$
modulo funkcje holomorficzne w $P$.
Innymi słowy, jest to element ilorazu
$\mathcal M_P/\mathcal O_{C,P}$, gdzie
$\mathcal M_P$ oznacza kiełki funkcji meromorficznych.
\end{definicja}

\begin{przyklad}[Przeszkoda w wyznaczeniu funkcji o zadanym biegunie]
Na $\PP^1$ zadana część główna $t^{-3}+2t^{-1}$
w punkcie $t=0$ jest realizowana przez globalną
funkcję wymierną $t^{-3}+2t^{-1}$. Na krzywej
eliptycznej $E$ nie da się natomiast zrealizować
części głównej $t^{-1}$ w jednym punkcie $P$ przez
funkcję wymierną bez innych biegunów:
taka funkcja należałaby do $L(P)\setminus L(0)$,
a \ref{tw:ell-poles} daje $L(P)=L(0)=\C$.
Jest to przykład globalnej przeszkody przy
pozornie poprawnych danych lokalnych.
\end{przyklad}

\begin{twierdzenie}[Residua wykrywają konieczny warunek realizacji części głównych]
Niech $C$ będzie gładką rzutową krzywą nad $\C$,
$P_1,\ldots,P_s$ parami różnymi punktami, a $p_j$
częściami głównymi w $P_j$. Jeśli istnieje globalna
funkcja meromorficzna $f$ na $C$ mająca dokładnie
te części główne i bez innych biegunów, to
\[
\sum_{j=1}^{s}\operatorname{Res}_{P_j}
  (p_j\,\omega)=0
\quad\text{dla każdej globalnej holomorficznej
różniczki }\omega\in H^0(C,\Omega^1_C).
\]
\end{twierdzenie}
\begin{proof}
Lokalna różnica $f-p_j$ jest holomorficzna
w $P_j$. Ponieważ $\omega$ jest tam holomorficzna,
formy $f\omega$ i $p_j\omega$ mają takie samo
residuum. W każdym pozostałym punkcie $f\omega$
jest holomorficzna. Sumując residua formy
meromorficznej $f\omega$ i stosując
twierdzenie~\ref{tw:residue-general}, otrzymujemy tezę.
\end{proof}

\begin{przyklad}[Dlaczego prosty biegun na torusie nie istnieje — dowód z residuów]
Na krzywej eliptycznej $E$ forma $\omega=\mathrm dx/(2y)$
nie znika nigdzie: na części afinicznej w punktach
$y\ne0$ jest niezerowa, bo $x$ jest parametrem;
w punktach $y=0$ forma $\mathrm dy/(3x^2+a)$ jest
niezerowa, a w nieskończoności ma współczynnik $-1$
przy $\mathrm du$. Jeśli $f$ miałaby jedyny prosty
biegun w $P$, to w parametrze $t$ mielibyśmy
$f=c_{-1}t^{-1}+O(1)$, $c_{-1}\ne0$, oraz
$\omega=(a_0+O(t))\,\mathrm dt$, $a_0\ne0$.
Stąd $\operatorname{Res}_P(f\omega)=c_{-1}a_0\ne0$.
Wszystkie pozostałe residua $f\omega$ są zerowe,
co przeczy twierdzeniu o ich sumie.
To drugi, niezależny rachunek wyjaśniający lukę
w rzędach biegunów funkcji na krzywej eliptycznej.
\end{przyklad}

\begin{uwaga}[Co dodaje pełna dualność Serre'a?]
Poprzednie twierdzenie daje \emph{warunek konieczny}.
Twierdzenie Mittag--Lefflera dla zwartych powierzchni
Riemanna mówi, że warunek znikania residuów względem
wszystkich globalnych holomorficznych różniczek
jest również wystarczający. W sformułowaniu
kohomologicznym odpowiada to niezdegenerowanemu
parowaniu
$H^1(C,\mathcal O_C)\times H^0(C,\Omega^1_C)\to\C$.
W przeciwieństwie do udowodnionej sumy residuów
część o wystarczalności wymaga dodatkowej teorii:
można ją uzyskać z twierdzenia Dolbeaulta i teorii
harmonicznych form różniczkowych. Poniżej podajemy
konstrukcję parowania oraz precyzyjną rolę tych
narzędzi; nie zastępujemy ich pozornie elementarnym
argumentem.
\end{uwaga}

\begin{stwierdzenie}[Parowanie Čecha i residuów na krzywej zespolonej]
\label{st:cech-res-pairing}
Niech $C$ będzie gładką spójną krzywą rzutową nad $\C$,
$\mathcal L$ holomorficzną wiązką liniową na $C^{\mathrm{an}}$,
a $\omega\in H^0(C^{\mathrm{an}},K_C\otimes\mathcal L^{-1})$.
Wybierzmy punkt $P\in C$, dysk współrzędnych $V$
wokół $P$ oraz $U=C\setminus\{P\}$.
Jeśli $g_{UV}$ jest holomorficzną sekcją $\mathcal L$
na pierścieniu $U\cap V$, to liczba
\[
\langle[g_{UV}],\omega\rangle
 =\frac1{2\pi i}\int_{\gamma}g_{UV}\,\omega,
\]
gdzie $\gamma$ jest małym dodatnim okręgiem wokół $P$,
nie zmienia się po zastąpieniu $g_{UV}$ przez
$g_{UV}+s_U|_{U\cap V}-s_V|_{U\cap V}$,
$s_U\in\mathcal L(U)$, $s_V\in\mathcal L(V)$.
Przy mnożeniu używamy naturalnego parowania
$\mathcal L\otimes (K_C\otimes\mathcal L^{-1})\to K_C$.
\end{stwierdzenie}
\begin{proof}
Po dodaniu $s_V$ całka nie zmieni się, bo
$s_V\omega$ jest holomorficzną formą na całym dysku
$V$; jej całka po $\gamma$ jest zerowa z twierdzenia
Cauchy'ego. Z kolei $s_U\omega$ jest holomorficzną
formą na $U$, czyli na powierzchni zwartej pozbawionej
punktu $P$. Wybierzmy mały dysk $D$ wokół $P$.
Na zwartym $C\setminus\operatorname{int}D$
forma $s_U\omega$ jest zamknięta. Stokes daje
$-\int_{\partial D}s_U\omega=0$, więc także
jej całka po $\gamma$ wynosi zero. Są to dokładnie
dwa rodzaje zmian reprezentanta dopuszczalne
w grupie Čecha. Dlatego parowanie zależy tylko
od klasy $[g_{UV}]$.
\end{proof}

\begin{uwaga}[Dlaczego to jest model ogólnej dualności?]
Zbiór $U=C\setminus\{P\}$ jest niezwartą powierzchnią
Riemanna, a zatem przestrzenią Steina na mocy twierdzenia
Behnkego--Steina. Dysk $V$ i pierścień $U\cap V$
są dziedzinami Steina;
także ich przecięcie jest przestrzenią Steina.
Twierdzenie Cartana B daje znikanie wyższej
kohomologii snopów koherentnych na tych częściach,
a twierdzenie porównawcze Čecha pozwala obliczać
$H^1(C^{\mathrm{an}},\mathcal L)$ z takiego pokrycia.
Twierdzenie Dolbeaulta identyfikuje tę grupę
z klasami $\bar\partial$-zamkniętych $(0,1)$-form
o wartościach w $\mathcal L$ modulo obrazy $\bar\partial$.
Poprzez rozkład jedności reprezentant Čecha
przechodzi w taką formę, a powyższa całka konturowa
zgadza się (z uwzględnieniem normalizacji $2\pi i$)
z całką formy $(1,1)$ po $C$.

W konsekwencji \emph{analityczna dualność Serre'a}
może być dowodzona za pomocą teorii Hodge'a:
harmoniczne $(0,1)$-formy z wartościami w
$\mathcal L$ są antyliniowo dualne do
holomorficznych sekcji $K_C\otimes\mathcal L^{-1}$.
Niezerowa forma harmoniczna ma nieznikającą normę
$L^2$, co zapewnia niezdegenerowaność parowania.
\emph{Twierdzenia Steina, Cartana B, Dolbeaulta i Hodge'a
są tu używane jako jawnie nazwane rezultaty.}
Własność niezależności parowania od reprezentantów
oraz samą globalną sumę residuów udowodniliśmy
natomiast bez ich przywoływania.
\end{uwaga}


\begin{uwaga}[Czego wymaga przejście od analizy do geometrii algebraicznej?]
Niech $C$ będzie gładką rzutową krzywą nad $\C$, a
$\mathcal L$ algebraiczną wiązką liniową. Jej
\emph{analityfikacja} $\mathcal L^{\mathrm{an}}$ jest
wiązką holomorficzną na zwartej powierzchni Riemanna
$C^{\mathrm{an}}$ (podstawowa topologia to zwykła
topologia zespolona, a nie topologia Zariskiego).
\emph{Twierdzenie porównawcze GAGA Serre'a} daje
naturalne izomorfizmy
$H^i(C,\mathcal L)\cong
H^i(C^{\mathrm{an}},\mathcal L^{\mathrm{an}})$
dla wszystkich $i\ge0$. Nie dowodzimy tu twierdzenia
GAGA. Jest ono koniecznym, \emph{zewnętrznym}
mostem między poniższym analitycznym dowodem dualności
a jej użyciem dla algebraicznych snopów na krzywych
zespolonych. Nad dowolnym algebraicznie domkniętym
ciałem (również dodatniej charakterystyki)
algebraiczna dualność Serre'a jest osobnym twierdzeniem;
nie wynika z przedstawionych metod Hodge'a.
\end{uwaga}

\begin{twierdzenie}[Dlaczego teoria Hodge'a daje niezdegenerowanie]
\label{tw:serre-hodge-dependent}
Niech $C$ będzie zwartą powierzchnią Riemanna,
a $\mathcal L$ holomorficzną wiązką liniową.
Przyjmijmy \emph{twierdzenie Dolbeaulta} oraz
\emph{twierdzenie Hodge'a dla operatora $\bar\partial$}:
każda klasa $H^1(C,\mathcal L)$ ma dokładnie jednego
reprezentanta harmonicznego w wybranych metrykach
hermitowskich. Wówczas parowanie
\[
H^1(C,\mathcal L)\times
H^0(C,K_C\otimes\mathcal L^{-1})\longrightarrow\C,
\qquad([\alpha],\omega)\longmapsto
\int_C\alpha\wedge\omega
\]
jest niezdegenerowane.
\end{twierdzenie}
\begin{proof}[Dowód oparty na jawnie przyjętych twierdzeniach analitycznych]
Wybierzmy hermitowską metrykę na $\mathcal L$ i metrykę
na $C$. Niech $0\ne[\alpha]\in H^1(C,\mathcal L)$;
z twierdzenia Hodge'a ma on niezerowego harmonicznego
reprezentanta $\alpha$. W lokalnej współrzędnej $z$
i holomorficznej ramie $e$ wiązki piszemy
$\alpha=a(z)\,\mathrm d\bar z\otimes e$ i
$\lambda(z)=\lVert e\rVert^2>0$.
Harmoniczność oznacza $\bar\partial\alpha=0$ oraz
$\bar\partial^*\alpha=0$. Pierwsza równość jest
automatyczna w zespolonym wymiarze jeden.
Druga, po całkowaniu przez części w lokalnym iloczynie
skalarnym, jest równoważna równaniu
$\partial_z(\lambda a)=0$ (ewentualne dodatnie czynniki
pochodzące od metryki na $C$ nie wpływają na jądro).
Zatem funkcja $\lambda\overline a$ jest holomorficzna.

Formuły $\omega=\lambda\overline a\,\mathrm dz\otimes e^{-1}$
na lokalnych mapach tworzą globalną holomorficzną
sekcję $K_C\otimes\mathcal L^{-1}$. Sprawdźmy zgodność
przy zmianie ramy: jeśli $e'=g e$, to
$a'=a/g$ i $\lambda'=|g|^2\lambda$,
więc $\lambda'\overline{a'}(e')^{-1}
=\lambda\overline a\,e^{-1}$. Zgodność przy zmianie
współrzędnej wynika z przekształcania
$\mathrm d\bar z$ i $\mathrm dz$.
Wybierając tę sekcję w parowaniu, otrzymujemy
w lokalnych współrzędnych formę
$\alpha\wedge\omega=
\lambda|a|^2\,\mathrm d\bar z\wedge\mathrm dz$.
Jej całka różni się od dodatniej normy $L^2$ formy
$\alpha$ tylko stałym niezerowym czynnikiem,
więc nie jest zerowa. Pokazaliśmy niezdegenerowanie
względem pierwszego argumentu.

W odwrotną stronę, dla holomorficznej niezerowej
sekcji $\omega=b(z)\,\mathrm dz\otimes e^{-1}$
konstruujemy harmoniczną formę
$\alpha=\lambda^{-1}\overline b\,\mathrm d\bar z\otimes e$.
Równanie $\partial_z(\lambda a)=0$ zachodzi, ponieważ
$\overline b$ jest antyholomorficzna. Ta forma nie
jest zerowa; dzięki twierdzeniu Hodge'a reprezentuje
niezerową klasę. Jej parowanie z $\omega$ jest
niezerową wielokrotnością całki z $|b|^2/\lambda$.
To kończy dowód względnej niezdegenerowaności.

Dowód jest \emph{kompletny po przyjęciu dwóch
wyraźnie sformułowanych twierdzeń analitycznych}.
Nie jest niezależnym dowodem twierdzeń Dolbeaulta
ani Hodge'a; te wymagają teorii eliptycznych
operatorów różniczkowych na zwartych rozmaitościach.
\end{proof}

\begin{przyklad}[Trzy różne typy wiedzy o dualności]
Na $\PP^1$ dysponujemy \emph{pełnym algebraicznym}
dowodem przez bazy Laurenta
(twierdzenie~\ref{tw:serre-p1-explicit}).
Na dowolnej zwartej powierzchni Riemanna
mamy \emph{pełny dowód} globalnego twierdzenia
o sumie residuów (twierdzenie~\ref{tw:residue-general})
oraz bezpośrednią konstrukcję dobrze określonego
parowania konturowego. Niezdegenerowaność
tego parowania dla dowolnej krzywej jest
natomiast \emph{konsekwencją ogólnej dualności},
której fundamenty analityczne zostały powyżej
nazwane, lecz nie udowodnione od podstaw.
Rozróżnienie tych trzech poziomów zapobiega
błędnemu przekonaniu, że lokalny rachunek
residuum automatycznie stanowi dowód całej
dualności Serre'a.
\end{przyklad}


\section{Twierdzenie Mittag--Lefflera: od lokalnych biegunów do funkcji globalnej}
\label{sec:mittag-residue}

Możemy teraz domknąć ideę, od której zaczęła się
kohomologia krzywych. Dane lokalne mogą być zgodne
z definicją funkcji meromorficznej, a jednak nie
pochodzić od \emph{jednej} funkcji na całej krzywej.
Warunki residuowe opisują dokładnie tę przeszkodę.

\begin{definicja}[Snop części głównych ograniczonych dywizorem]
Niech $C$ będzie gładką rzutową krzywą nad $\C$,
$D=\sum_{j=1}^sm_jP_j$ efektywnym dywizorem,
gdzie $m_j\ge1$, a punkty $P_j$ są parami różne.
Oznaczmy przez $\mathcal P_D$ ilorazowy snop
$\mathcal O_C(D)/\mathcal O_C$.
W punkcie $P_j$ wybierzmy parametr $t_j$.
Wtedy $(\mathcal P_D)_{P_j}$ jest przestrzenią
wektorową nad $\C$ o bazie klas
$t_j^{-m_j},\ldots,t_j^{-1}$; poza nośnikiem $D$
źdźbła tego snopa są zerowe. Zatem
\[
 H^0(C,\mathcal P_D)\cong
 \bigoplus_{j=1}^s
 \operatorname{span}_{\C}
 \{t_j^{-m_j},\ldots,t_j^{-1}\}.
\]
To, że podana rodzina istotnie stanowi bazę,
wynika z faktu, że w dyskretnym pierścieniu
wartościowania $\mathcal O_{C,P_j}$ iloraz
$t_j^{-m_j}\mathcal O_{C,P_j}/\mathcal O_{C,P_j}$
ma dokładnie te $m_j$ klasy jednomianów.
\end{definicja}

\begin{twierdzenie}[Mittag--Leffler na gładkiej krzywej zespolonej]
\label{tw:mittag-curve}
Niech $C$ będzie gładką, spójną, rzutową krzywą
nad $\C$, $P_1,\ldots,P_s$ parami różnymi punktami,
a $p_j=\sum_{r=1}^{m_j}a_{jr}t_j^{-r}$
zadanymi częściami głównymi w wybranych parametrach
lokalnych $t_j$ przy $P_j$. Istnieje globalna
funkcja meromorficzna $f$ na $C$, której części
główne są równe $p_j$, i która nie ma innych
biegunów, wtedy i tylko wtedy, gdy
\[
  \sum_{j=1}^{s}\operatorname{Res}_{P_j}
       (p_j\,\omega)=0
  \qquad\text{dla każdej }\omega\in H^0(C,K_C).
\]
Dowód wykorzystuje analityczną dualność Serre'a
uzyskaną z twierdzeń Dolbeaulta i Hodge'a.
\end{twierdzenie}
\begin{proof}
Konieczność została udowodniona w poprzedniej
sekcji za pomocą globalnego twierdzenia o residuach.
Przeprowadźmy dokładnie dowód wystarczalności.

Niech $D=\sum_jm_jP_j$. Włączenie stałych warunków
na brak biegunów w warunki dopuszczające bieguny
rzędów nie większych niż $m_j$ prowadzi do krótkiego
ciągu dokładnego snopów
\[
 0\longrightarrow\mathcal O_C
 \longrightarrow\mathcal O_C(D)
 \longrightarrow\mathcal P_D\longrightarrow0.
\]
Jego dokładność wynika z lokalnych wzorów
$\mathcal O_{C,P_j}\subseteq
 t_j^{-m_j}\mathcal O_{C,P_j}$ oraz identyczności
tych źdźbeł poza punktami $P_j$.
Z długiego ciągu dokładnego kohomologii otrzymujemy
fragment
\[
 H^0(C,\mathcal O_C(D))\xrightarrow{\ q\ }
 H^0(C,\mathcal P_D)\xrightarrow{\ \partial\ }
 H^1(C,\mathcal O_C).
\]
Odwzorowanie $q$ przypisuje funkcji jej części główne.
Z dokładności wynika, że dana rodzina
$p=(p_1,\ldots,p_s)$ pochodzi od globalnej funkcji
meromorficznej z dopuszczonymi biegunami dokładnie wtedy,
gdy $\partial p=0$.

Pozostaje wyjaśnić, jak rozpoznać zerowość tej klasy.
Wybierzmy dla każdego $P_j$ mały dysk $V_j$, przy czym dyski są parami rozłączne, i weźmy
$U=C\setminus\{P_1,\ldots,P_s\}$. Na każdym
$V_j$ wybieramy meromorficzną funkcję $p_j$ realizującą
zadane części główne; na $U$ wybieramy lokalne podniesienie
równe zeru. Różnice podniesień na pierścieniach
$U\cap V_j$ tworzą $1$-kocykl reprezentujący
$\partial p$. Przy konwencji
$\partial p$ reprezentowanej przez $p_j|_{U\cap V_j}$,
parowanie konturowe z globalną holomorficzną różniczką
$\omega$ jest równe
\[
  \langle\partial p,\omega\rangle
  =\sum_{j=1}^s\frac1{2\pi i}
        \int_{\partial D_j}p_j\,\omega
  =\sum_{j=1}^s\operatorname{Res}_{P_j}(p_j\omega).
\]
Jeśli w definicji różniczki Čecha wybiera się
przeciwną konwencję znaku, obie strony zmieniają
znak, co nie wpływa na kryterium znikania.
Równość jest przypadkiem zgodności parowania
Čecha i residuów wykazanej przy
stwierdzeniu~\ref{st:cech-res-pairing}.

Z założenia wszystkie powyższe liczby są zerowe.
Niezdegenerowanie dualności Serre'a
(twierdzenie~\ref{tw:serre-hodge-dependent}) daje
$\partial p=0$. Z dokładności ciągu istnieje
$f\in H^0(C,\mathcal O_C(D))$ takie, że
$q(f)=p$. Sekcja ta jest szukaną funkcją
meromorficzną. To kończy dowód.
\end{proof}

\begin{przyklad}[Sfera i torus: dwa odmienne wyniki dla prostego bieguna]
Na $\PP^1(\C)$ nie ma niezerowych holomorficznych
różniczek, bo $K_{\PP^1}\cong\mathcal O(-2)$.
Warunek z twierdzenia~\ref{tw:mittag-curve}
jest wobec tego pusty: każdą skończoną kolekcję
części głównych funkcji można zrealizować
przez globalną funkcję wymierną.

Na krzywej eliptycznej $E$ istnieje natomiast
niezerowa holomorficzna forma $\omega$ bez zer.
Dla jednej zadanej części głównej $t^{-1}$ w punkcie
$P$ residuum $(t^{-1}\omega)$ jest niezerowe,
więc takiej funkcji nie ma. Dla dwóch punktów
$P,Q$ i części głównych
$a\,t_P^{-1}$ oraz $b\,t_Q^{-1}$
warunek przyjmuje postać
$a\,\omega_P(\partial/\partial t_P)
+b\,\omega_Q(\partial/\partial t_Q)=0$.
Jest to \emph{jedno} równanie liniowe, bo
$\dim H^0(E,K_E)=1$. Kiedy jest spełnione,
możemy wybrać globalną funkcję meromorficzną
z tymi dwiema zadanymi częściami głównymi.
\end{przyklad}

\section{Podsumowanie zależności i wskazówki do dalszej nauki}
\label{sec:summary-four}

Prześledźmy najważniejszy łańcuch pojęć. Lokalizacja
przekształca pierścień globalnych wielomianów
w pierścień funkcji regularnych przy wybranym
punkcie. W przypadku gładkiej krzywej ten lokalny
pierścień ma \emph{jeden parametr} i jest DVR;
dzięki temu możemy bez niejednoznaczności mierzyć
rzędy zer i biegunów. Dywizor łączy te lokalne
liczby w jeden obiekt globalny. Snop
$\mathcal O_C(D)$ tłumaczy warunki na zera i bieguny
na warunki o sekcjach, a kohomologia mierzy
deficyt możliwości sklejania tych sekcji.
Riemann--Roch zamienia tę informację na wzór
numeryczny: stopień dywizora i rodzaj krzywej
określają różnicę dwóch wymiarów.

\begin{center}
\begin{tikzpicture}[node distance=8mm,>=Stealth,
  every node/.style={font=\small},
  summaryblock/.style={draw,rounded corners,align=center,text width=7.5cm,
  inner sep=5pt}]
\node[summaryblock] (a) {Pierścień globalny $A$ i lokalizacja $A_{\mathfrak p}$};
\node[summaryblock,below=of a] (b) {Pierścień lokalny DVR $\to$ rzędy $\operatorname{ord}_P$};
\node[summaryblock,below=of b] (c) {Dywizory $D$ i snopy odwracalne $\mathcal O_C(D)$};
\node[summaryblock,below=of c] (d) {Kohomologia $H^0,H^1$ i dualność Serre'a};
\node[summaryblock,below=of d] (e) {Riemann--Roch, rodzaj, krzywe eliptyczne};
\draw[->,thick] (a)--(b); \draw[->,thick] (b)--(c);
\draw[->,thick] (c)--(d); \draw[->,thick] (d)--(e);
\end{tikzpicture}

{\small Rysunek 4.9. Zależności między narzędziami algebry przemiennej, snopami i geometrią krzywych.}\par
\end{center}

Po rozszerzeniu udowodniliśmy już bezpośrednio dokładność
kompleksu Čecha dla podstawowych pokryć afinicznych,
globalną sumę residuów na zwartej powierzchni Riemanna
oraz dobrze określoność parowania residuowego.
Jawnie oddzielamy od nich użyte głębsze rezultaty:
porównanie kohomologii Čecha z kohomologią snopów,
acykliczność afiniczną w pełnej ogólności, skończoność
kohomologii koherentnej oraz twierdzenia Dolbeaulta
i Hodge'a. Teoria dualności Serre'a została rozwinięta
do konstrukcji parowania i do dowodu jego
niezdegenerowania \emph{przy przyjęciu} dwóch
ostatnich twierdzeń, bez sugerowania, że są one elementarne.

% Źródła referencyjne (nie zastępują dowodów w tekście):
% The Stacks Project, Tags 0C4S, 0A22, 0B5B.
% R. Vakil, The Rising Sea: Foundations of Algebraic Geometry.

% =============================================================
% DODATEK DO ROZDZIALU 4: RESIDUA WIELU ZMIENNYCH
% Wprowadzany po podsumowaniu dotychczasowego rozdzialu 4.
% =============================================================
\section{Residua wielowymiarowe: od całek Cauchy'ego do geometrii przecięć}
\label{sec:ga4-residua-wielo}

\subsection{Dlaczego nie wystarcza residuum jednej zmiennej?}

W jednej zmiennej meromorficzna forma $h(z)\,\mathrm dz/f(z)$
ma izolowane bieguny i możemy całkować po małej okrężnej krzywej.
W $\C^n$ dla $n\ge2$ pojedyncze równanie $f=0$ opisuje na ogół
hiperpowierzchnię, a nie zbiór punktów. W szczególności funkcja
$1/f$ nie ma izolowanego bieguna w sensie jednej zmiennej.
Aby otrzymać \emph{liczbę} z lokalnego całkowania $n$-formy,
rozpatrzymy $n$ funkcji $f_1,\ldots,f_n$ o izolowanym wspólnym zerze.
Inna konstrukcja, \emph{residuum Poincarégo}, będzie przypisywać
formę różniczkową gładkiej hiperpowierzchni $f=0$.
Te dwie konstrukcje pełnią różne role.

\begin{uwaga}[Dwie różne dziedziny całkowania]
W $\C^2$ okrąg $S^1$ ma wymiar rzeczywisty $1$, sfera
$S^3$ ma wymiar $3$, natomiast powierzchnia całkowania
residuum Grothendiecka ma wymiar rzeczywisty $2$:
typowo jest torusem $S^1\times S^1$.
Nie należy bez wyjaśnienia zastępować tego torusa małą sferą $S^3$.
\end{uwaga}

\subsection{Wielowymiarowy wzór Cauchy'ego}
\label{subsec:ga4-cauchy-product}

\begin{twierdzenie}[Wzór Cauchy'ego na polidysku]
\label{tw:ga4-polydisc-cauchy}
Niech $a=(a_1,\ldots,a_n)\in\C^n$, $r_j>0$, a $h$ będzie
holomorficzna w otoczeniu domkniętego polidysku
$\overline P(a,r)=\{z\in\C^n:|z_j-a_j|\le r_j,\ 1\le j\le n\}$.
Wówczas dla $z\in P(a,r)$ zachodzi
\[
 h(z)=\frac{1}{(2\pi i)^n}
 \int_{|\zeta_1-a_1|=r_1}\!\cdots
 \int_{|\zeta_n-a_n|=r_n}
 \frac{h(\zeta_1,\ldots,\zeta_n)}
 {\prod_{j=1}^n(\zeta_j-z_j)}
 \,\mathrm d\zeta_n\cdots\mathrm d\zeta_1.
\]
Każdą okręgową krzywą całkujemy przeciwnie do ruchu wskazówek
zegara; porządek całkowań został tu jawnie ustalony.
\end{twierdzenie}
\begin{proof}
Traktujemy $\zeta_2,\ldots,\zeta_n$ jako ustalone parametry.
Jednowymiarowy wzór Cauchy'ego względem $\zeta_1$ zastępuje
pierwszą całkę wartością funkcji przy $\zeta_1=z_1$.
Pozostała funkcja jest holomorficzna w następnej zmiennej,
więc ponownie stosujemy wzór Cauchy'ego. Po $n$ krokach
otrzymujemy $h(z_1,\ldots,z_n)$.
Wszystkie całki są całkami funkcji ciągłych po zwartych torusach,
co uzasadnia zamianę kolejności ich obliczania w tej właśnie
sytuacji produktu niezależnych okręgów.
\end{proof}

\begin{wniosek}[Wzór na współczynnik Taylora]
\label{wn:ga4-monomial-res}
Niech $m_1,\ldots,m_n\ge1$ będą liczbami całkowitymi,
a $h$ funkcją holomorficzną w otoczeniu domkniętego
polidysku $\{z\in\C^n:|z_j|\le r_j,\ 1\le j\le n\}$
dla ustalonych $r_1,\ldots,r_n>0$. Wówczas
\[
\frac{1}{(2\pi i)^n}\int_{|z_1|=r_1}\!\cdots
\int_{|z_n|=r_n}
\frac{h(z)\,\mathrm dz_n\cdots\mathrm dz_1}
{z_1^{m_1}\cdots z_n^{m_n}}
=\frac{\partial_1^{m_1-1}\cdots\partial_n^{m_n-1}h(0)}
{(m_1-1)!\cdots(m_n-1)!}.
\]
W szczególności całka odczytuje współczynnik przy
$z_1^{m_1-1}\cdots z_n^{m_n-1}$ w szeregu Taylora funkcji $h$.
\end{wniosek}
\begin{proof}
Z jednowymiarowego wzoru Cauchy'ego dla pochodnych
obliczamy najpierw całkę po okręgu w pierwszej zmiennej.
Następnie powtarzamy ten rachunek w pozostałych zmiennych.
Wynik jest podaną pochodną mieszaną podzieloną przez
iloczyn silni. Kolejność różniczkowania funkcji
holomorficznej nie zmienia wyniku.
\end{proof}

\begin{przyklad}[Residuowe odczytywanie współczynników]
W $\C^2$ rozważmy $h(z,w)=e^{z+w}$.
Wielomian Taylora pokazuje, że współczynnik $zw^2$ wynosi $1/2$.
Dlatego
\[
\frac{1}{(2\pi i)^2}\int_{|z|=r}\int_{|w|=s}
\frac{e^{z+w}}{z^2w^3}\,\mathrm dw\,\mathrm dz=\frac12.
\]
Natomiast całka z $1/(z^2w^3)$ wynosi zero. Sam fakt,
że mianownik ma biegun w każdej zmiennej, nie wystarcza,
aby residuum było niezerowe.
\end{przyklad}

\begin{center}
\begin{tikzpicture}[>=Stealth,scale=0.95]
 \draw[->,gray!70] (-4.8,0)--(-2.25,0) node[right,font=\scriptsize] {$\Re z$};
 \draw[->,gray!70] (-3.5,-1.15)--(-3.5,1.25) node[above,font=\scriptsize] {$\Im z$};
 \draw[thick,blue!70!black] (-3.5,0) circle (0.72);
 \node[below,font=\footnotesize] at (-3.5,-1.35) {$|z|=r$};
 \draw[->,gray!70] (0.0,0)--(2.5,0) node[right,font=\scriptsize] {$\Re w$};
 \draw[->,gray!70] (1.2,-1.15)--(1.2,1.25) node[above,font=\scriptsize] {$\Im w$};
 \draw[thick,teal!70!black] (1.2,0) circle (0.72);
 \node[below,font=\footnotesize] at (1.2,-1.35) {$|w|=s$};
 \node[draw,rounded corners,align=center,font=\footnotesize,
       text width=9.0cm,minimum height=0.72cm] at (-1.15,-2.47)
 {$\Gamma=\{(z,w):|z|=r,\ |w|=s\}\cong S^1\times S^1$
 \quad (torus o wymiarze rzeczywistym $2$)};
\end{tikzpicture}

{\small Rysunek 4.10. Torus całkowania jako produkt dwóch dodatnio
zorientowanych okręgów. Nie jest on sferą trójwymiarową.}\par
\end{center}

\subsection{Residuum Grothendiecka w izolowanym wspólnym zerze}
\label{subsec:ga4-groth-res}

\begin{definicja}[Residuum lokalne Grothendiecka]
\label{def:ga4-groth-res}
Niech $U\subseteq\C^n$ będzie otwarte, $a\in U$, a funkcje
$f_1,\ldots,f_n,h\in\mathcal O(U)$ takie, że $a$ jest
\emph{izolowanym} punktem zbioru
$\{z\in U:f_1(z)=\cdots=f_n(z)=0\}$.
Wybierzmy dostatecznie małe otoczenie $V$ punktu $a$,
nie zawierające innego wspólnego zera, oraz dostatecznie małe regularne
wartości dodatnie $\varepsilon_j$ funkcji $|f_j|$,
tak aby rzeczywisty cykl
\[
\Gamma_\varepsilon=\{z\in V:|f_1(z)|=\varepsilon_1,\ldots,
|f_n(z)|=\varepsilon_n\}
\]
był zawarty w tym otoczeniu.
Orientację wyznacza dodatniość formy
$\mathrm d\arg f_1\wedge\cdots\wedge\mathrm d\arg f_n$.
Definiujemy
\[
\operatorname{Res}_a\!\left[\frac{h\,\mathrm dz_1\wedge\cdots
\wedge\mathrm dz_n}{f_1\cdots f_n}\right]
=\frac1{(2\pi i)^n}\int_{\Gamma_\varepsilon}
\frac{h(z)\,\mathrm dz_1\wedge\cdots\wedge\mathrm dz_n}
{f_1(z)\cdots f_n(z)}.
\]
Wybór otoczenia i dostatecznie małych $\varepsilon_j$
nie zmienia wartości całki; to własność dobrze określonego
lokalnego cyklu residuowego.
\end{definicja}

\begin{uwaga}[Co jest twierdzeniem, a co definicją?]
Istnienie odpowiedniego małego cyklu i niezależność jego
klasy homologii od dostatecznie małych parametrów wynika
z lokalnej teorii skończonych odwzorowań holomorficznych
lub z konstrukcji cyklu Leraya. To jest odrębny wynik
analityczny, przyjmowany tutaj jawnie; w przykładach
regularnych sprowadza się on do zmiany współrzędnych,
a w przypadku jednomianowym do torusa produktowego.
Zapis $\operatorname{Res}_a$ zawsze odnosi się do \emph{uporządkowanej}
listy $(f_1,\ldots,f_n)$, nie tylko do ideału, który generują.
\end{uwaga}

\begin{twierdzenie}[Proste wspólne zero: wzór jakobianowy]
\label{tw:ga4-jac-res}
Niech $f_1,\ldots,f_n,h$ będą holomorficzne w otoczeniu
$a\in\C^n$, spełniają $f_j(a)=0$ dla każdego $j$, a
\[
J_f(a)=\det\left(\frac{\partial f_i}{\partial z_j}(a)
\right)_{1\le i,j\le n}\ne0.
\]
Wówczas $a$ jest izolowanym wspólnym zerem i
\[
\operatorname{Res}_a\!\left[
\frac{h\,\mathrm dz_1\wedge\cdots\wedge\mathrm dz_n}
{f_1\cdots f_n}\right]=\frac{h(a)}{J_f(a)}.
\]
\end{twierdzenie}
\begin{proof}
Holomorficzne twierdzenie o odwzorowaniu odwrotnym
mówi, że $u=f(z)$ są współrzędnymi holomorficznymi
w otoczeniu $a$. Ponieważ
$\mathrm du_1\wedge\cdots\wedge\mathrm du_n
=J_f(z)\,\mathrm dz_1\wedge\cdots\wedge\mathrm dz_n$,
podcałkowa forma przyjmuje postać
$\bigl(h(z(u))/J_f(z(u))\bigr)\mathrm du_1\wedge\cdots
\wedge\mathrm du_n/(u_1\cdots u_n)$.
Cykl $\Gamma_\varepsilon$ przechodzi na standardowy torus
$|u_j|=\varepsilon_j$, z jego dodatnią orientacją.
Wielowymiarowy wzór Cauchy'ego daje wartość
$h(a)/J_f(a)$.
\end{proof}

\begin{przyklad}[Jakobian i znak residuum]
Weźmy $f_1(z,w)=z+w$, $f_2(z,w)=z-w$ i $h\equiv1$.
Wspólne zero to $(0,0)$, a jakobian ma wyznacznik
$\det\begin{pmatrix}1&1\\1&-1\end{pmatrix}=-2$.
Zatem
$\operatorname{Res}_{(0,0)}[\mathrm dz\wedge\mathrm dw/
((z+w)(z-w))]=-1/2$.
Po zamianie kolejności mianowników otrzymamy $+1/2$.
Nie jest to błąd rachunkowy: orientacja cyklu zależy od
kolejności funkcji definiujących przecięcie.
\end{przyklad}

\subsection{Zera wielokrotne i lokalna algebra}
\label{subsec:ga4-res-multiple}

Nie wolno stosować wzoru $h(a)/J_f(a)$, gdy $J_f(a)=0$.
Residuum może jednak nadal istnieć i być łatwe do obliczenia.

\begin{twierdzenie}[Wzór jednomianowy dla residuum]
\label{tw:ga4-res-monomial}
Niech $m_1,\ldots,m_n\ge1$, a $h$ będzie holomorficzna
w otoczeniu zera w $\C^n$. Wtedy
\[
\operatorname{Res}_0\!\left[
\frac{h(z)\,\mathrm dz_1\wedge\cdots\wedge\mathrm dz_n}
{z_1^{m_1}\cdots z_n^{m_n}}\right]
=\frac{\partial_1^{m_1-1}\cdots\partial_n^{m_n-1}h(0)}
{\prod_{j=1}^n(m_j-1)!}.
\]
\end{twierdzenie}
\begin{proof}
Cykl $|z_j^{m_j}|=\varepsilon_j$ jest torusem produktowym
$|z_j|=\varepsilon_j^{1/m_j}$; jego orientacja jest
zgodna ze standardową, ponieważ każda funkcja
$z_j\mapsto z_j^{m_j}$ zachowuje orientację okręgu.
Stosujemy wniosek~\ref{wn:ga4-monomial-res}.
\end{proof}

\begin{przyklad}[Jakobian znika, lecz residuum nie]
Dla $f_1=z^2$ i $f_2=w^2$ mamy $J_f(0,0)=0$.
Mimo to $\operatorname{Res}_0[zw\,\mathrm dz\wedge\mathrm dw/
(z^2w^2)]=1$ oraz
$\operatorname{Res}_0[\mathrm dz\wedge\mathrm dw/(z^2w^2)]=0$.
Lokalny iloraz $\C\{z,w\}/(z^2,w^2)$ ma bazę
$1,z,w,zw$, więc ma wymiar $4$ nad $\C$.
Wartość residuum nie jest po prostu tą długością:
jest funkcjonałem liniowym zależnym od licznika $h$.
\end{przyklad}

\begin{twierdzenie}[Twierdzenie o transformacji residuum — wynik zewnętrzny]
\label{tw:ga4-res-transform}
Niech $a$ będzie izolowanym wspólnym zerem
$f=(f_1,\ldots,f_n)$, gdzie $f_j$ są holomorficzne
w otoczeniu $a$. Niech $A=(a_{ij})$ będzie macierzą
funkcji holomorficznych, a $g_i=\sum_{j=1}^n a_{ij}f_j$.
Załóżmy, że $a$ jest również izolowanym wspólnym zerem
$g=(g_1,\ldots,g_n)$. Wówczas dla każdej holomorficznej
funkcji $h$ zachodzi
\[
\operatorname{Res}_a\!\left[
\frac{h\,\mathrm dz_1\wedge\cdots\wedge\mathrm dz_n}
{f_1\cdots f_n}\right]
=\operatorname{Res}_a\!\left[
\frac{h\det A\,\mathrm dz_1\wedge\cdots\wedge\mathrm dz_n}
{g_1\cdots g_n}\right].
\]
\end{twierdzenie}

\begin{uwaga}[Dlaczego trzeba uważać na kierunek wzoru?]
Jeżeli $g=Af$, to po \emph{prawej} stronie pojawia się
$\det A$ i mianowniki $g_i$. Nie należy bez sprawdzenia
przenosić wyznacznika na lewą stronę.
Dla $n=1$, $f=z$, $g=z^2=zf$ wzór daje
$\operatorname{Res}_0[h\,\mathrm dz/z]
=\operatorname{Res}_0[hz\,\mathrm dz/z^2]=h(0)$.
Ogólny dowód prawa transformacji korzysta z lokalnych
własności cyklu Grothendiecka (równoważnie z lokalnej
dualności); nie zastępujemy go nieuprawnioną zmianą
zmiennych, gdy $\det A(a)=0$.
\end{uwaga}

\begin{przyklad}[Nietrywialne przecięcie zdegenerowane]
\label{prz:ga4-res-degenerate}
Weźmy $f_1=z^2$ oraz $f_2=w^2+z^3$ w otoczeniu zera.
Wspólne zero jest izolowane, a jakobian $J_f(0)=0$.
Niech $g_1=z^2$ i $g_2=w^2=f_2-zf_1$. Wtedy
\[
\begin{pmatrix}g_1\\g_2\end{pmatrix}
=\begin{pmatrix}1&0\\-z&1\end{pmatrix}
\begin{pmatrix}f_1\\f_2\end{pmatrix},\qquad \det A=1.
\]
Wzór transformacji sprowadza zadanie do współczynnika
przy $zw$ w rozwinięciu licznika $h$:
\[
\operatorname{Res}_0\!\left[
\frac{h(z,w)\,\mathrm dz\wedge\mathrm dw}
{z^2(w^2+z^3)}\right]
=\partial_z\partial_w h(0,0).
\]
Dla $h=1$, $h=zw$ i $h=(z+w)^2$ otrzymujemy
odpowiednio $0$, $1$ i $2$.
Przy tym lokalny pierścień
$\C\{z,w\}/(z^2,w^2+z^3)$ jest izomorficzny z
$\C\{z,w\}/(z^2,w^2)$, z bazą $1,z,w,zw$.
Tak widzimy jednocześnie krotność przecięcia równą $4$
i subtelniejsze informacje zachowywane przez residuum.
\end{przyklad}

\begin{stwierdzenie}[Anulowanie ideału w liczniku]
\label{st:ga4-res-annihilates}
Jeżeli $a$ jest izolowanym wspólnym zerem funkcji
$f_1,\ldots,f_n\in\mathcal O(U)$, a
$h\in(f_1,\ldots,f_n)\mathcal O_{U,a}$, to
\[
\operatorname{Res}_a\!\left[
\frac{h\,\mathrm dz_1\wedge\cdots\wedge\mathrm dz_n}
{f_1\cdots f_n}\right]=0.
\]
\end{stwierdzenie}
\begin{proof}
Możemy przesunąć $a$ do zera. Ponieważ wspólne zero
jest izolowane, lokalny analityczny Nullstellensatz daje
$\sqrt{(f_1,\ldots,f_n)}=(z_1,\ldots,z_n)$,
a zatem dla pewnych $N_i\ge1$ funkcje $g_i=z_i^{N_i}$
należą do ideału $(f_1,\ldots,f_n)$.
Wybierzmy macierz holomorficzną $A$ taką, że $g=Af$.
Z prawa transformacji~\ref{tw:ga4-res-transform} mamy
$\operatorname{Res}_0(h;f)=\operatorname{Res}_0(h\det A;g)$.
Jeśli $h=\sum_j b_jf_j$, to tożsamość macierzowa
$\operatorname{adj}(A)A=(\det A)I$ pokazuje, że
$(\det A)f_j$ jest kombinacją funkcji $g_i$.
Stąd $h\det A\in(g_1,\ldots,g_n)$.
W zapisie $g_i=z_i^{N_i}$ każdy składnik licznika
jest podzielny przez odpowiednie $z_i^{N_i}$,
więc wzór Cauchy'ego na współczynniki daje residuum zero.
Jawnymi przesłankami tego dowodu są lokalny analityczny
Nullstellensatz i twierdzenie o transformacji.
\end{proof}

\begin{uwaga}[Dalsza lokalna dualność]
Dla izolowanego przecięcia zupełnego funkcjonał
$[h]\mapsto\operatorname{Res}_a(h;f)$ jest
szczególnym funkcjonałem na skończenie wymiarowej
algebrze lokalnej $\mathcal O_{\C^n,a}/(f_1,\ldots,f_n)$.
Silniejsze twierdzenie głosi, że parowanie
$([h],[q])\mapsto\operatorname{Res}_a(hq;f)$ jest
niezdegenerowane. To przykład \emph{lokalnej dualności}
Grothendiecka; jego ogólny dowód pozostawiamy poza
zakresem skryptu, a nie utożsamiamy go ze zwykłym
wzorem jakobianowym.
\end{uwaga}

\subsection{Residuum Poincarégo: forma na hiperpowierzchni}
\label{subsec:ga4-poincare-res}

\begin{definicja}[Lokalne residuum Poincarégo]
\label{def:ga4-poincare-res}
Niech $U\subseteq\C^n$ będzie otwarte, $n\ge2$, a
$f,g\in\mathcal O(U)$ takie, że $H=\{f=0\}$ jest
gładką hiperpowierzchnią w otoczeniu punktu $a\in H$
(tzn. $\mathrm df(a)\ne0$). Dla meromorficznej formy
$\omega=g\,\mathrm dz_1\wedge\cdots\wedge\mathrm dz_n/f$
wybierzmy indeks $j$ z $\partial f/\partial z_j\ne0$.
\emph{Residuum Poincarégo} na $H$ definiujemy lokalnie jako
\[
\operatorname{Res}_H^{\mathrm P}(\omega)
=\left.(-1)^{j-1}\frac{g\,\mathrm dz_1\wedge\cdots
\wedge\widehat{\mathrm dz_j}\wedge\cdots\wedge\mathrm dz_n}
{\partial f/\partial z_j}\right|_H.
\]
Symbol daszka oznacza pominięcie czynnika.
\end{definicja}
\begin{proof}[Uzasadnienie niezależności od współrzędnych]
Ponieważ $\mathrm df\ne0$, funkcja $f$ może służyć
jako jedna z holomorficznych współrzędnych lokalnych.
W takich współrzędnych istnieje holomorficzna $(n-1)$-forma
$\eta$ taka, że $\omega=(\mathrm df/f)\wedge\eta$.
Po ograniczeniu do $H$ składnik $\eta|_H$ jest
niezależny od wyboru uzupełnienia $f$ do układu
współrzędnych: dwie możliwości różnią się o składnik
zawierający $\mathrm df$, którego ograniczenie do $TH$
znika. Przestawienie $\mathrm df$ na pierwsze miejsce
w iloczynie zewnętrznym daje znak $(-1)^{j-1}$
i zapisany wzór. Konstrukcja lokalna skleja się na
przecięciach map współrzędnych na $H$.
\end{proof}

\begin{przyklad}[Residuum wzdłuż paraboli]
Dla $f(z,w)=w-z^2$ rozważmy
$\omega=\mathrm dz\wedge\mathrm dw/(w-z^2)$.
Ponieważ $\partial_w f=1$, z definicji przy $j=2$ mamy
$\operatorname{Res}^{\mathrm P}_{\{w=z^2\}}(\omega)=-\mathrm dz$.
Znak sprawdzamy bez pamięciowego wzoru:
$\mathrm df\wedge(-\mathrm dz)
=(\mathrm dw-2z\,\mathrm dz)\wedge(-\mathrm dz)
=\mathrm dz\wedge\mathrm dw$.
Residuum jest więc holomorficzną $1$-formą na paraboli,
nie liczbą przypisaną jej wybranemu punktowi.
\end{przyklad}

\begin{uwaga}[Związek z adjunkcją]
Dla gładkiego dzielnika $H$ w rozmaitości zespolonej $X$
residuum Poincarégo opisuje lokalnie izomorfizm
$(K_X\otimes\mathcal O_X(H))|_H\cong K_H$.
Jest to analityczna realizacja wzoru adjunkcji,
znanego już z geometrii krzywych i powierzchni.
Nie stosujemy tego izomorfizmu bezpośrednio do punktu
osobliwego hiperpowierzchni, gdzie warunek $\mathrm df\ne0$
nie zachodzi.
\end{uwaga}

\subsection{Globalne sumy residuów: prosty przykład zamiast pochopnego uogólnienia}
\label{subsec:ga4-global-multi-res}

\begin{przyklad}[Cztery wspólne zera i suma jakobianowa]
\label{prz:ga4-four-residues}
Niech $f_1(z,w)=z^2-1$ i $f_2(z,w)=w^2-1$.
Wspólne zera to $(\pm1,\pm1)$, a $J_f(z,w)=4zw$.
Dla licznika $h\in\C[z,w]$ suma lokalnych residuów
we wszystkich czterech punktach wynosi
\[
\sum_{\substack{z=\pm1\\w=\pm1}}
\operatorname{Res}_{(z,w)}\!\left[
\frac{h\,\mathrm dz\wedge\mathrm dw}{(z^2-1)(w^2-1)}\right]
=\sum_{z,w=\pm1}\frac{h(z,w)}{4zw}.
\]
Jeżeli $h=1$ albo $h=z+w$, suma znika.
Jeżeli $h=zw$, każdy składnik jest równy $1/4$,
zatem suma wynosi $1$.
Widać, że stwierdzenie \emph{„suma residuów wielowymiarowych
zawsze wynosi zero”} byłoby fałszywe bez warunków
na zachowanie formy w nieskończoności.
\end{przyklad}

\begin{twierdzenie}[Wersja Eulera--Jacobiego — wynik zewnętrzny]
\label{tw:ga4-euler-jacobi}
Niech $f_1,\ldots,f_n\in\C[z_1,\ldots,z_n]$ mają
stopnie $d_1,\ldots,d_n\ge1$. Załóżmy, że ich
najwyższe części jednorodne nie mają wspólnego zera
w $\PP^{n-1}(\C)$, a wszystkie wspólne zera afiniczne
są proste. Wówczas zbiór tych zer jest skończony,
a dla każdego wielomianu $h$ stopnia
$\deg h<\sum_{j=1}^n(d_j-1)$ zachodzi
\[
\sum_{a\in\{f_1=\cdots=f_n=0\}}\frac{h(a)}{J_f(a)}=0.
\]
\end{twierdzenie}

\begin{uwaga}[Dlaczego pojawia się warunek w nieskończoności?]
Warunek o najwyższych częściach jednorodnych wyklucza
wspólne punkty zerowe po homogenizacji na hiperpłaszczyźnie
w nieskończoności. Bez niego część informacji residuowej
może „uciekać do nieskończoności”. Pełny dowód powyższego
wzoru wykorzystuje globalną teorię residuów lub algebraiczną
postać teorii eliminacji i nie jest tu ukrywany pod
jednowymiarowym twierdzeniem o sumie residuów.
\end{uwaga}

\begin{przyklad}[Sprawdzenie progu stopnia]
W przykładzie~\ref{prz:ga4-four-residues} mamy $d_1=d_2=2$,
więc próg to $2$. Dla $h=1,z,w$ (stopień mniejszy niż $2$)
suma znika, lecz dla $h=zw$ (stopień równy $2$)
suma jest równa $1$. Ograniczenie stopnia w twierdzeniu
nie jest zatem czysto techniczne.
\end{przyklad}

\subsection{Które residuum stosować i kiedy?}
\begin{enumerate}[label=\arabic*)]
\item Jeżeli mamy $n$ równań holomorficznych w $\C^n$ z izolowanym
wspólnym zerem i chcemy otrzymać \emph{liczbę}, stosujemy residuum
Grothendiecka; przy niezerowym jakobianie korzystamy z wzoru
$h(a)/J_f(a)$, a przy zerowym z jednomianów i prawa transformacji.
\item Jeżeli mamy \emph{jedną} gładką hiperpowierzchnię $H=\{f=0\}$
i meromorficzną $n$-formę z prostym biegunem wzdłuż $H$,
stosujemy residuum Poincarégo, otrzymując $(n-1)$-formę na $H$.
\item Jeżeli chcemy dodawać residua ze wszystkich punktów wspólnych
zer, musimy kontrolować zachowanie przy nieskończoności;
w wyraźnie podanych hipotezach służy temu wzór Eulera--Jacobiego.
\item Dowolne \emph{iterowane residua} w wybranych zmiennych nie są
automatycznie tym samym co residuum Grothendiecka.
Wersja iterowana wymaga określenia kolejności oraz
obszarów rozwinięć (a więc i odpowiednich cykli całkowania).
\end{enumerate}

\begin{uwaga}[Świadome granice zakresu]
Prądy residuowe Coleffa--Herrery, residua toriczne i
pełna teoria lokalnej dualności Grothendiecka wykraczają
poza przyjęty tu zakres. Przedstawione definicje, wzory
wykonawcze i przykłady wystarczają jednak do pierwszych
zastosowań w analizie wielu zmiennych zespolonych
i lokalnej geometrii przecięć.
\end{uwaga}

% Referencje kontrolne: A. Tsikh, Multidimensional Residues and Their Applications;
% E. Cattani--A. Dickenstein, Introduction to Residues and Resultants;
% A. Yger, Residue Calculus and Applications;
% D. Arapura, Algebraic Geometry over the Complex Numbers.

% =========================================================
% ROZDZIAL 5. ZALEZNOSCI: ROZDZIALY 0--4
% Wszystkie etykiety nowego rozdzialu maja prefiks ga5:.
% =========================================================

\chapter{Rozdmuchania i rozwiązywanie osobliwości}
\label{ch:ga5-rozdmuchania}

\section{Dlaczego wygładzanie osobliwości wymaga zmiany przestrzeni?}
\label{sec:ga5-motywacja}

Rozważmy krzywą $C=\V(y^2-x^3)\subseteq\A^2_\C$. W początku układ
$y^2=x^3$ nie wyznacza gładkiej krzywej: obie pochodne cząstkowe
znikają. Nie jest to jednak argument, że nie można jej opisać
\emph{gładką} krzywą. Parametryzacja $t\mapsto(t^2,t^3)$ pokazuje, że
trudność polega na sposobie umieszczenia krzywej w płaszczyźnie.

Samo usunięcie punktu osobliwego nie jest rozwiązaniem: gubimy punkt,
a zależy nam na odwzorowaniu do wyjściowej krzywej, które nie zmienia
jej poza osobliwościami. Chcemy zastąpić punkt dodatkowymi danymi,
na przykład kierunkiem, z którego do niego dochodzimy. Prowadzi to
do \emph{rozdmuchania} (ang. \emph{blow-up}).

W rozdziale odróżniamy trzy cele:
\begin{enumerate}[label=\arabic*)]
\item \emph{rozdmuchanie otoczenia} punktu w gładkiej powierzchni;
      nowa powierzchnia może nadal być gładka, mimo że włókno nad
      punktem zawiera całą prostą rzutową;
\item \emph{rozwiązanie osobliwości krzywej}: uzyskanie gładkiej
      krzywej z odwzorowaniem właściwym i biwymiernym na $C$;
\item \emph{rozwiązanie zanurzone}: dodatkowo wymaga, aby wygładzona
      krzywa i dzielniki wyjątkowe przecinały się z normalnymi
      skrzyżowaniami.
\end{enumerate}
Trzecie zadanie jest silniejsze od drugiego. Pokażemy to rachunkiem
na ostrzu $y^2=x^3$.

W całym rozdziale $k$ oznacza ciało algebraicznie domknięte
charakterystyki $0$, chyba że przy danym twierdzeniu zaznaczono
słabsze założenia. Posługujemy się wcześniejszymi wiadomościami
o pierścieniach, lokalizacji i schematach afinicznych; potrzebne
uzupełnienia podajemy poniżej. Nie zakładamy znajomości ogólnej
teorii schematów projektowych.

\section{Schematy i morfizmy: minimalny, ale dokładny fundament}
\label{sec:ga5-schematy}

\subsection{Morfizmy i włókna nie są tylko odwzorowaniami punktów}

\begin{definicja}[Morfizm schematów]
\label{def:ga5-morfizm}
Niech $(X,\mathcal O_X)$ i $(Y,\mathcal O_Y)$ będą schematami.
Morfizmem $f:X\to Y$ nazywamy parę złożoną z ciągłego odwzorowania
$f:|X|\to|Y|$ oraz homomorfizmu snopów pierścieni
$f^\sharp:\mathcal O_Y\to f_*\mathcal O_X$ (równoważnie,
$f^{-1}\mathcal O_Y\to\mathcal O_X$), taką że dla każdego $x\in X$
indukowany homomorfizm lokalnych pierścieni
$\mathcal O_{Y,f(x)}\to\mathcal O_{X,x}$ jest \emph{lokalny}, tzn.
przeciwobraz ideału maksymalnego pierścienia $\mathcal O_{X,x}$
jest dokładnie ideałem maksymalnym $\mathcal O_{Y,f(x)}$.
\end{definicja}

Sformułowanie „homomorfizm lokalny” oznacza dokładnie
$\varphi^{-1}(\mathfrak m_B)=\mathfrak m_A$ dla homomorfizmu
$\varphi:A\to B$ pierścieni lokalnych $(A,\mathfrak m_A)$ i
$(B,\mathfrak m_B)$. Dzięki temu odwzorowanie respektuje funkcje
regularne i ich wartości, a nie tylko zbiór punktów.

\begin{stwierdzenie}[Morfizmy schematów afinicznych]
\label{st:ga5-affine-morphism}
Dla dowolnych pierścieni przemiennych $A,B$ z jedynką istnieje
naturalna bijekcja
\[
 \operatorname{Hom}_{\mathrm{Sch}}(\Spec B,\Spec A)
 \cong\operatorname{Hom}_{\mathrm{Ring}}(A,B).
\]
Homomorfizm $\varphi:A\to B$ odpowiada odwzorowaniu
$\mathfrak q\mapsto\varphi^{-1}(\mathfrak q)$ na ideałach pierwszych.
\end{stwierdzenie}
\begin{proof}
Przeciwobraz ideału pierwszego jest pierwszy: jeśli
$\varphi(a)\varphi(b)\in\mathfrak q$, to przynajmniej jeden
z obrazów należy do $\mathfrak q$. Ponadto
$\varphi^{-1}(D(a))=D(\varphi(a))$, więc odwzorowanie spektrów
jest ciągłe. Na podstawowych otwartych zbiorach $D(a)$
homomorfizm $A\to B$ indukuje
$A_a\to B_{\varphi(a)}$ przez $u/a^m\mapsto\varphi(u)/\varphi(a)^m$.
Te homomorfizmy są zgodne z dalszymi lokalizacjami i sklejają się
do homomorfizmu snopów. Na źdźbłach dostajemy
$A_{\mathfrak p}\to B_{\mathfrak q}$ dla
$\mathfrak p=\varphi^{-1}(\mathfrak q)$; jest to homomorfizm lokalny.

W drugą stronę morfizm schematów określa homomorfizm sekcji
globalnych $\Gamma(\Spec A,\mathcal O)=A\to
\Gamma(\Spec B,\mathcal O)=B$. Dla $\mathfrak q\in\Spec B$
lokalność homomorfizmu na źdźbłach wymusza obraz
$\varphi^{-1}(\mathfrak q)$. Zgodność z ograniczeniami do
$D(a)$ wymusza również opisany homomorfizm snopów.
Zatem obie konstrukcje są wzajemnie odwrotne.
\end{proof}

\begin{przyklad}[Dlaczego kierunek strzałki się odwraca?]
Homomorfizm $k[t]\to k[x]$, $t\mapsto x^2$, daje morfizm
$\A^1_k\to\A^1_k$ wysyłający $a\mapsto a^2$ na punktach
zamkniętych. Funkcje na przestrzeni docelowej cofamy na przestrzeń
źródłową: $t$ staje się funkcją $x^2$.
Ta kontrawariantność nie jest kwestią zapisu, lecz zasadą
geometryczną „funkcję można złożyć z odwzorowaniem”.
\end{przyklad}

\begin{definicja}[Iloczyn włóknisty i włókno]
\label{def:ga5-fibreproduct}
Dla morfizmów $f:X\to S$ i $g:Y\to S$ ich \emph{iloczynem
włóknistym} jest schemat $X\times_S Y$ wraz z rzutowaniami do $X$
i $Y$, komutującymi z odwzorowaniami do $S$, o następującej
własności uniwersalnej: każdy schemat $T$ z morfizmami $T\to X$
i $T\to Y$ zgodnymi po przejściu do $S$ ma dokładnie jeden
morfizm $T\to X\times_S Y$ zgodny z rzutowaniami.

Jeżeli $s\in S$ i $\kappa(s)$ oznacza ciało reszt
$\mathcal O_{S,s}/\mathfrak m_s$, to \emph{włóknem} $f$ nad $s$
jest schemat $X_s=X\times_S\Spec\kappa(s)$.
\end{definicja}

\begin{stwierdzenie}[Iloczyn włóknisty w przypadku afinicznym]
\label{st:ga5-tensor-fibre}
Dla homomorfizmów $R\to A$, $R\to B$ zachodzi
\[
 \Spec A\times_{\Spec R}\Spec B\cong\Spec(A\otimes_R B).
\]
\end{stwierdzenie}
\begin{proof}
Dla dowolnego pierścienia $C$ homomorfizm
$A\otimes_R B\to C$ odpowiada parze homomorfizmów
$A\to C$ i $B\to C$, których ograniczenia do $R$ są równe:
para wyznacza odwzorowanie $a\otimes b\mapsto\alpha(a)\beta(b)$,
a jego istnienie i jednoznaczność wynika z własności uniwersalnej
iloczynu tensorowego. Za pomocą stwierdzenia~\ref{st:ga5-affine-morphism}
jest to dokładnie własność uniwersalna iloczynu włóknistego
schematów afinicznych. Dla dowolnego schematu $T$ parę
morfizmów $T\to\Spec A$ i $T\to\Spec B$ zgodnych nad
$\Spec R$ ograniczamy do jego afinicznego pokrycia.
Na każdym otwartym afinicznym otrzymujemy jedyny morfizm
w $\Spec(A\otimes_R B)$; jednoznaczność zapewnia ich
zgodność na przecięciach, toteż sklejają się na całym $T$.
Dowodzi to wymaganej własności uniwersalnej bez
ograniczania schematu testowego do przypadku afinicznego.
\end{proof}

\begin{przyklad}[Włókno może być punktem „podwójnym”]
\label{prz:ga5-double-fibre}
W tym przykładzie $k$ jest dowolnym ciałem o charakterystyce różnej od $2$ (nie zakładamy jego algebraicznej domkniętości).
Niech $f:\A^1_k\to\A^1_k$ będzie dany przez $t=x^2$.
Nad punktem $t=0$ otrzymujemy
$\Spec(k[x]\otimes_{k[t]}k[t]/(t))\cong\Spec k[x]/(x^2)$.
To włókno ma jeden punkt topologiczny, ale zawiera niezerowy
nilpotent $\overline x$. Nad $t=a\ne0$ (przy $\operatorname{char}k\ne2$)
włókno jest zredukowanym schematem $\Spec k[x]/(x^2-a)$.
Ma dwa różne punkty $k$-wymierne, jeśli $a$ jest kwadratem
w $k$; jeśli $a$ nie jest kwadratem, wielomian $x^2-a$
jest nierozkładalny i włókno ma jeden punkt o ciele reszt
$k(\sqrt a)$. Po przejściu do domknięcia algebraicznego
włókno ma zawsze dwa różne punkty geometryczne.
Dlatego obraz geometryczny włókna nie zawiera całej informacji
o morfizmie.
\end{przyklad}

\subsection{Zanurzenia zamknięte, ideały i obrazy odwrotne}

\begin{definicja}[Zanurzenie zamknięte i centrum]
Morfizm $i:Z\to X$ jest \emph{zanurzeniem zamkniętym}, gdy
lokalnie na $X=\Spec A$ pochodzi od surjekcji $A\to A/I$.
Wtedy $Z$ jest \emph{podschematem zamkniętym} wyznaczonym
przez snop ideałów $\mathcal I\subseteq\mathcal O_X$.
Przeciwobrazem schematowym $f^{-1}(Z)$ względem $f:Y\to X$
nazywamy $Y\times_X Z$; jego snop ideałów jest
obrazem $f^*\mathcal I\to\mathcal O_Y$, oznaczanym
$\mathcal I\mathcal O_Y$.
\end{definicja}

\begin{uwaga}[Przeciwobraz zbiorowy i schematowy]
Przeciwobraz zbioru $|Z|$ ma te same punkty co
$|Y\times_X Z|$, lecz nie pamięta krotności. Na przykład
przeciwobraz schematowy zera przez $x\mapsto x^2$
jest właśnie $\Spec k[x]/(x^2)$, nie $\Spec k$.
Ta różnica powróci przy \emph{transformacie całkowitej} krzywej.
\end{uwaga}

\subsection{Gładkość, właściwość i biwymierność}

\begin{definicja}[Gładkość morfizmu]
\label{def:ga5-smooth}
Morfizm schematów $f:X\to S$ jest \emph{gładki w punkcie}
$x\in X$ o wymiarze względnym $d$, jeżeli w pewnych otoczeniach
afinicznych $x$ i $f(x)$ można przedstawić go jako morfizm
$\Spec B\to\Spec R$, gdzie
$B=(R[z_1,\ldots,z_n]/(F_1,\ldots,F_r))_h$,
$n-r=d$ i pewien minor rozmiaru $r$ macierzy
$(\partial F_i/\partial z_j)$ jest odwracalny w $B$.
Morfizm jest \emph{gładki}, jeśli jest gładki w każdym punkcie.
\end{definicja}

Jest to kryterium oparte na \emph{lokalnych prezentacjach
standardowo gładkich}; równoważność z innymi definicjami gładkości
morfizmów wymaga twierdzenia o lokalnym modelu gładkim.
Wynik ten przyjmujemy w ogólnym przypadku. W przykładach poniżej
wystarczą wprost podane prezentacje i rachunek Jacobiego.

\begin{przyklad}[Rodzina gładka poza jednym parametrem]
Dla $X=\V(xy-t)\subseteq\A^2_{x,y}\times\A^1_t$
morfizm $X\to\A^1_t$ ma nad $t=a\ne0$ gładką hiperbolę
$xy=a$, a nad $t=0$ sumę osi $xy=0$. Cała przestrzeń $X$
jest izomorficzna z $\A^2_{x,y}$, lecz morfizm do osi parametru
nie jest gładki w punkcie $(x,y,t)=(0,0,0)$: różniczki
względem zmiennych włókna to $(y,x)$, które tam znikają.
Gładkość całej przestrzeni nie oznacza gładkości każdego morfizmu.
\end{przyklad}

\begin{definicja}[Morfizm właściwy i rzutowy]
\label{def:ga5-proper}
Morfizm schematów $f:X\to S$ jest \emph{właściwy}, gdy jest
skończonego typu, rozdzielony i uniwersalnie domknięty:
dla każdego $S'\to S$ rzutowanie
$X\times_S S'\to S'$ wysyła zbiory domknięte na domknięte.
Morfizm jest \emph{rzutowy}, gdy nad $S$ jest izomorficzny
z zanurzeniem zamkniętym $X\hookrightarrow\PP^n_S$
(złożonym z projekcją $\PP^n_S\to S$) dla pewnego $n$.
\end{definicja}

\begin{uwaga}[Skąd bierze się właściwość?]
Każdy morfizm rzutowy jest właściwy; w ogólności dowód wymaga
kryterium wartościowaniowego i faktu, że $\PP^n_S\to S$ jest
właściwy. Dla czytelnika ważny jest konkretny sens:
morfizm $\A^1_k\to\Spec k$ nie jest właściwy, zaś
$\PP^1_k\to\Spec k$ jest właściwy, choć $\PP^1$ nie jest
afiniczna. Przykładowy dowód wartościowaniowy dla
$\PP^n$ podajemy w stwierdzeniu~\ref{st:ga5-valuative-projective}.
\end{uwaga}

\begin{stwierdzenie}[Rozszerzanie punktów rzutowych nad DVR]
\label{st:ga5-valuative-projective}
Niech $R$ będzie dyskretnym pierścieniem wartościowania,
$K=\Frac(R)$ i $n\ge0$. Każdy punkt $p\in\PP^n(K)$
ma dokładnie jedno przedłużenie do morfizmu
$\Spec R\to\PP^n_R$ zgodnego z $p$ nad $\Spec K$.
\end{stwierdzenie}
\begin{proof}
Weźmy jednorodne współrzędne $p=[a_0:\cdots:a_n]$
z $a_i\in K$, nie wszystkie zerowe. Niech $\pi$ oznacza
parametr jednostajny $R$. Po pomnożeniu wszystkich
$a_i$ przez odpowiednią potęgę $\pi$ możemy założyć,
że wszystkie $a_i\in R$ i co najmniej jeden z nich jest
jednostką. Dają zatem morfizm $\Spec R\to\PP^n_R$:
na otoczeniu $D_+(X_i)$ wybranego odwracalnego $a_i$
współrzędne afiniczne mają wartości $a_j/a_i\in R$.

Jeśli $[b_0:\cdots:b_n]$ określa inne przedłużenie,
możemy analogicznie wybrać współrzędne całkowite z
co najmniej jedną jednostką. Na punkcie ogólnym mamy
$b_i=\lambda a_i$ dla wspólnego $\lambda\in K^\times$.
Minimalne wartościowania współrzędnych obu reprezentacji
są równe zeru, więc $v(\lambda)=0$: $\lambda$ jest
jednostką $R$. Wobec tego obie reprezentacje wyznaczają
ten sam morfizm. Jest to część kryterium wartościowaniowego,
która wyjaśnia geometrycznie, dlaczego granice kierunków
pozostają w przestrzeni rzutowej.
\end{proof}

\begin{definicja}[Morfizm biwymierny i rozwiązanie osobliwości]
\label{def:ga5-resolution}
Niech $X,Y$ będą całkowitymi (nierozkładalnymi i zredukowanymi)
schematami skończonego typu nad $k$. Morfizm $f:Y\to X$
nazywamy \emph{biwymiernym}, jeśli indukuje izomorfizm ciał
funkcji $k(X)\cong k(Y)$; równoważnie istnieją niepuste
otwarte podzbiory, na których $f$ jest izomorfizmem.
\emph{Rozwiązaniem osobliwości} całkowitego schematu $X$
nazywamy właściwy morfizm biwymierny $\pi:\widetilde X\to X$,
gdzie $\widetilde X$ jest gładka nad $k$, a $\pi$ jest
izomorfizmem nad gładkim miejscem $X_{\mathrm{reg}}$.
\end{definicja}

Warunek „właściwy” zapobiega uznaniu samego otwartego
podzbioru gładkiego za rozwiązanie. Wymaganie izomorfizmu
nad $X_{\mathrm{reg}}$ zapobiega niepotrzebnemu zmienianiu
gładkich punktów. W dalszym ciągu będziemy rozpatrywać
także podschematy zredukowane o wielu składowych;
wtedy rozwiązuje się oddzielnie składowe lub używa
odpowiedniej wersji dla schematów zredukowanych.

\section{Konstrukcja \texorpdfstring{$\operatorname{Proj}$}{Proj}: przestrzeń z algebry stopniowanej}
\label{sec:ga5-proj}

\begin{definicja}[Spektrum jednorodne]
Niech $S=\bigoplus_{n\ge0}S_n$ będzie pierścieniem
stopniowanym, $S_mS_n\subseteq S_{m+n}$, i niech
$S_+=\bigoplus_{n>0}S_n$ będzie ideałem nieistotnym
(ang. \emph{irrelevant ideal}). Definiujemy
\[
 \operatorname{Proj}S=
 \{\mathfrak p\subset S:\mathfrak p\text{ jest jednorodnym ideałem
 pierwszym},\ S_+\not\subseteq\mathfrak p\}.
\]
Dla jednorodnego $f\in S_d$, $d>0$, oznaczamy
$D_+(f)=\{\mathfrak p\in\operatorname{Proj}S:f\notin\mathfrak p\}$.
\end{definicja}

Ideał jednorodny oznacza ideał generowany przez elementy
jednorodne. Pominięcie ideałów zawierających $S_+$ jest
istotne: w $k[T_0,T_1]$ ideał $(T_0,T_1)$ odpowiadałby
początkowi stożka, a nie kierunkowi na prostej rzutowej.

\begin{definicja}[Część stopnia zerowego lokalizacji]
Jeżeli $f\in S_d$ jest jednorodny, $d>0$, to
$S_{(f)}=(S_f)_0$ oznacza podpierścień elementów stopnia
$0$ w lokalizacji $S_f$. Każdy jego element można zapisać
jako $a/f^m$, gdzie $a\in S_{md}$ jest jednorodny.
\end{definicja}

\begin{stwierdzenie}[Podstawowe otwarte w $\operatorname{Proj}$]
\label{st:ga5-proj-charts}
Dla każdego jednorodnego $f\in S_+$ przestrzeń $D_+(f)$
ma kanoniczną strukturę schematu afinicznego
$D_+(f)\cong\Spec S_{(f)}$. Jeśli $S$ jest generowany
jako $S_0$-algebra przez skończenie wiele elementów
$f_1,\ldots,f_r$ dodatnich stopni, to zbiory $D_+(f_i)$
pokrywają $\operatorname{Proj}S$.
\end{stwierdzenie}
\begin{proof}
Wzajemne przyporządkowanie ideałów pierwszych uzyskuje się
przez lokalizację jednorodnych ideałów $\mathfrak p$
niezawierających $f$, a następnie wzięcie ich części
stopnia zero. W drugą stronę dla
$\mathfrak q\in\Spec S_{(f)}$ rozważa się jednorodne
$a\in S_e$, których odpowiednia potęga $a^d/f^e$
należy do $\mathfrak q$. Użycie potęg jest potrzebne,
gdy $\deg f>1$. Z definicji wynika, że otrzymany zbiór
jest jednorodnym ideałem pierwszym, niezawierającym $f$;
obie operacje są odwrotne. Lokalizacje funkcji regularnych
na dalszych otwartych podzbiorach są zgodne, co daje
strukturę snopa i identyfikację schematów afinicznych.

Jeśli $\mathfrak p\in\operatorname{Proj}S$ zawierałby
wszystkie $f_i$, to zawierałby również generowany
przez nie ideał $S_+$, wbrew definicji. Zatem
przynajmniej jedno $f_i$ nie należy do $\mathfrak p$
i otrzymujemy pokrycie.
\end{proof}

\begin{przyklad}[Powrót do zwykłej przestrzeni rzutowej]
Dla $S=k[U,V]$ ze zwykłym stopniowaniem
$D_+(U)=\Spec k[V/U]$, a $D_+(V)=\Spec k[U/V]$.
Na przecięciu współrzędne $v=V/U$ i $u=U/V$
spełniają $uv=1$. Sklejenie dwóch prostych afinicznych
według $u=v^{-1}$ daje $\PP^1_k$.
Podobnie $\operatorname{Proj}k[T_0,\ldots,T_n]=\PP^n_k$.
\end{przyklad}

\begin{uwaga}[Co jest tu twierdzeniem, a co konstrukcją?]
Pełne sprawdzenie zgodności snopów na przecięciach jest
standardowym lematem konstrukcji $\operatorname{Proj}$.
Dla rozdmuchań będziemy korzystać tylko z jawnych map
opisanych pierścieniami i sprawdzimy ich przejścia bezpośrednio.
Dlatego dalsze rachunki nie zależą od nieudowodnionych
własności kohomologicznych $\operatorname{Proj}$.
\end{uwaga}

\section{Algebra Reesa i właściwa definicja rozdmuchania}
\label{sec:ga5-rees}

\begin{definicja}[Algebra Reesa]
\label{def:ga5-rees}
Niech $A$ będzie pierścieniem przemiennym z jedynką,
a $I\subseteq A$ ideałem. Jego \emph{algebrą Reesa}
nazywamy pierścień stopniowany
\[
 \mathcal R_A(I)=\bigoplus_{n\ge0}I^nT^n\subseteq A[T],
 \qquad I^0=A,
\]
przy czym składnik $I^nT^n$ ma stopień $n$.
\emph{Rozdmuchaniem} $\Spec A$ wzdłuż centrum
$Z=\Spec(A/I)$ jest schemat
$\operatorname{Bl}_Z(\Spec A)=\operatorname{Proj}\mathcal R_A(I)$
z naturalnym morfizmem $\pi:\operatorname{Bl}_Z(\Spec A)\to\Spec A$.
\end{definicja}

\begin{uwaga}[Dlaczego nie wystarczy $A/I$?]
Iloraz $A/I$ pamięta samo centrum. Ciąg
$I,I^2,I^3,\ldots$ pamięta kolejne rzędy znikania
wzdłuż centrum. Algebra Reesa zbiera je w jednym
obiekcie stopniowanym. Po wzięciu $\operatorname{Proj}$
uzyskujemy kierunki, w jakich znikają elementy $I$.
To właśnie informacja tracona przy zastąpieniu
centrum pojedynczym punktem.
\end{uwaga}

\begin{stwierdzenie}[Mapy współrzędnych rozdmuchania dziedziny]
\label{st:ga5-blowup-charts}
Niech $A$ będzie dziedziną, a
$I=(f_1,\ldots,f_r)\ne(0)$ ideałem skończenie generowanym,
przy czym $f_i\ne0$. Rozdmuchanie $\operatorname{Bl}_I\Spec A$
jest pokryte przez mapy
\[
 U_i=\Spec A[I/f_i],\qquad
 A[I/f_i]=A\bigl[f_1/f_i,\ldots,f_r/f_i\bigr]\subseteq\Frac(A).
\]
Na $U_i$ ideał centrum po cofnięciu staje się główny:
$I\mathcal O_{U_i}=f_i\mathcal O_{U_i}$.
\end{stwierdzenie}
\begin{proof}
Algebra Reesa jest generowana nad $A$ przez jednorodne
elementy $f_1T,\ldots,f_rT$. Zatem z poprzedniego
stwierdzenia wystarczą $D_+(f_iT)$. Element stopnia
zero w $\mathcal R_A(I)_{f_iT}$ ma postać
$(aT^m)/(f_iT)^m=a/f_i^m$ dla $a\in I^m$.
Ponieważ $A$ jest dziedziną, traktujemy te ilorazy
jako elementy $\Frac(A)$. Są generowane przez $A$
oraz $f_j/f_i$, więc pierścień mapy to $A[I/f_i]$.
Dla każdego $j$ równość $f_j=f_i(f_j/f_i)$ pokazuje,
że $I\mathcal O_{U_i}\subseteq f_i\mathcal O_{U_i}$;
odwrotna inkluzja jest oczywista, bo $f_i\in I$.
\end{proof}

\begin{uwaga}[Ograniczenie ważne przy dzielnikach zera]
Jeśli $A$ ma dzielniki zera, zapis $A[I/f_i]$ jako
podpierścienia $A_{f_i}$ może zgubić torsję. Ogólne
mapy definiujemy wtedy wyłącznie przez
$(\mathcal R_A(I)_{f_iT})_0$; nie wolno bez dowodu
zastępować ich ilorazami w ciele ułamków.
W szczególności definicja przez algebrę Reesa działa
także dla schematów niezredukowanych.
\end{uwaga}

\begin{stwierdzenie}[Poza centrum nic się nie zmienia]
\label{st:ga5-off-centre}
Dla dowolnego schematu $X$ i quasi-koherentnego snopa
ideałów $\mathcal I$ skończonego typu, definiującego $Z\subseteq X$
morfizm rozdmuchania indukuje izomorfizm
$\operatorname{Bl}_Z X\setminus\pi^{-1}(Z)\cong X\setminus Z$.
\end{stwierdzenie}
\begin{proof}
Wystarczy sprawdzić to po ograniczeniu do otwartych
afinicznych. Jeżeli $f\in I$ jest odwracalne na $D(f)$,
to $I A_f=A_f$. Jego algebra Reesa jest wówczas
$A_f[T]$ ze zwykłym stopniowaniem, więc
$\operatorname{Proj}A_f[T]\cong\Spec A_f$: jedyna mapa
$D_+(T)$ jest całym $\operatorname{Proj}$, a jego
pierścień stopnia zero jest $A_f$. Zbiory $D(f)$
dla $f\in I$ pokrywają dopełnienie centrum; otrzymane
izomorfizmy są zgodne na przecięciach.
\end{proof}

\begin{stwierdzenie}[Rozdmuchanie jest lokalnie rzutowe i właściwe]
\label{st:ga5-blowup-proper}
Niech $X$ będzie schematem, a $\mathcal I\subseteq\mathcal O_X$
quasi-koherentnym snopem ideałów skończonego typu.
W otoczeniach afinicznych $X$ morfizm rozdmuchania jest
złożeniem zanurzenia zamkniętego w pewną przestrzeń
$\PP^{r-1}_X$ i projekcji na $X$. Zatem jest lokalnie
na bazie rzutowy, a w szczególności właściwy. Istnienie
jednego globalnego zanurzenia w $\PP^N_X$ wymaga
silniejszych założeń i nie wynika wyłącznie z tej lokalnej
konstrukcji.
\end{stwierdzenie}
\begin{proof}
Lokalnie niech $I=(f_1,\ldots,f_r)\subseteq A$.
Przyporządkowanie $T_i\mapsto f_iT$ daje surjektywny
homomorfizm stopniowanych $A$-algebr
$A[T_1,\ldots,T_r]\twoheadrightarrow\mathcal R_A(I)$.
Każdy jednorodny ideał relacji wycina domknięty
podschemat w $\operatorname{Proj}A[T_1,\ldots,T_r]
=\PP^{r-1}_A$; otrzymujemy więc zanurzenie zamknięte
rozdmuchania w przestrzeni rzutowej nad $\Spec A$.
Złożenie zanurzenia zamkniętego z projekcją rzutową
jest właściwe. Ostatnie przejście wykorzystuje standardowe
twierdzenie o właściwości morfizmów rzutowych (jego
mechanizm wartościowaniowy zilustrowaliśmy wcześniej).
\end{proof}

\begin{uwaga}[Drobna, ale istotna subtelność globalna]
Lokalny dowód daje \emph{morfizm lokalnie rzutowy}.
Jeśli ideał $\mathcal I$ jest generowany przez skończoną
rodzinę sekcji globalnych, otrzymujemy globalne zanurzenie
w $\PP^{r-1}_X$. Z właściwości jako własności lokalnej
na bazie wynika jednak, że lokalna rzutowość wystarcza
do właściwości również bez założenia noetherowskości.
Nie należy utożsamiać obu rodzajów rzutowości.
\end{uwaga}

\subsection{Własność uniwersalna i dzielniki Cartiera}

\begin{definicja}[Efektywny dzielnik Cartiera]
\label{def:ga5-cartier}
Podschemat zamknięty $D\subseteq Y$ jest \emph{efektywnym
dzielnikiem Cartiera}, jeżeli jego snop ideałów
$\mathcal I_D$ jest lokalnie generowany przez jeden
\emph{dzielnik regularny}, czyli przez element niebędący
dzielnikiem zera w odpowiednim pierścieniu lokalnym.
Wówczas $\mathcal I_D$ jest snopem odwracalnym.
\end{definicja}

\begin{przyklad}[Dzielnik Cartiera to nie dowolne równanie]
W $\A^2$ prosta $x=0$ jest efektywnym dzielnikiem Cartiera,
bo $x$ nie jest dzielnikiem zera w $k[x,y]$.
Natomiast w $\Spec k[x,y]/(xy)$ element $x$ jest dzielnikiem
zera; podschemat zadany przez $(x)$ nie jest tam efektywnym
dzielnikiem Cartiera. Znaczenie ma pierścień otoczenia,
a nie sam widoczny zbiór punktów.
\end{przyklad}

\begin{twierdzenie}[Własność uniwersalna rozdmuchania]
\label{tw:ga5-universal}
Niech $X$ będzie schematem, $\mathcal I\subseteq\mathcal O_X$
quasi-koherentnym snopem ideałów skończonego typu i
$\pi:\widetilde X=\operatorname{Bl}_{\mathcal I}X\to X$
rozdmuchaniem. Wtedy $\mathcal I\mathcal O_{\widetilde X}$
jest ideałem efektywnego dzielnika Cartiera. Ponadto
każdy morfizm $f:Y\to X$, dla którego
$\mathcal I\mathcal O_Y$ definiuje efektywny dzielnik
Cartiera, faktoryzuje się \emph{jednoznacznie} przez $\pi$:
$ f=\pi\circ\widetilde f$ dla dokładnie jednego
$\widetilde f:Y\to\widetilde X$.
\end{twierdzenie}
\begin{proof}
Własność jest lokalna na $X$ i $Y$. Na otoczeniu
afinicznym $X=\Spec A$ niech $I=(f_1,\ldots,f_r)$.
Na mapie $D_+(f_iT)$ ideał $I\mathcal O_{\widetilde X}$
jest generowany przez obraz $f_i$. Musimy uzasadnić,
że jest on niedzielnikiem zera również wtedy, gdy $A$
nie jest dziedziną. Oznaczmy przez $B_i$ część stopnia
zero $\mathcal R_A(I)_{f_iT}$. Jeśli $f_i b=0$ w $B_i$,
wybierzmy zapis $b=aT^m/(f_iT)^m$ z $a\in I^m$.
Zerowość $f_i b$ w lokalizacji znaczy, że istnieje
$N\ge0$ takie, że $(f_iT)^N(f_i aT^m)=0$
w $\mathcal R_A(I)$. Pierścień ten jest podpierścieniem
$A[T]$, a zatem $f_i^{N+1}a=0$. Mnożąc przez $f_i$
jeszcze raz, widzimy, że $(f_iT)^{N+1}(aT^m)=0$;
czyli już $b=0$. Wynika stąd, że $f_i$ jest
niedzielnikiem zera w $B_i$. Zatem cofnięty ideał
jest lokalnie główny i odwracalny.

Załóżmy teraz, że $f:Y\to X$ ma właściwość z tezy.
W otoczeniu danego $y\in Y$ obrazy $f_1,\ldots,f_r$
generują ideał $(a)$ z $a$ niedzielnikiem zera.
Po zapisaniu $f_j=ab_j$ przynajmniej jeden z $b_j$
ma wartość odwracalną w $\mathcal O_{Y,y}$: inaczej
wszystkie leżałyby w ideale maksymalnym i nie mogłyby
generować $(a)$ (stosujemy lemat Nakayamy do lokalnie
wolnego modułu rzędu jeden lub po prostu fakt, że
jedynka należy do ideału generowanego przez $b_j$).
Po zmniejszeniu otoczenia można więc przyjąć, że
$f_i$ generuje cofnięty ideał. Relacje pomiędzy
$f_j/f_i$ na mapie rozdmuchania mają wtedy jedyny
sensowny obraz $b_j/b_i=f_j/f_i$ w pierścieniu
funkcji na $Y$ (iloraz po stronie $Y$ jest dobrze
określony, bo $f_i$ jest niedzielnikiem zera).

Ściśle: homomorfizm
$(\mathcal R_A(I)_{f_iT})_0\to\mathcal O_Y$
wysyła klasę $aT^m/(f_iT)^m$ na jedyny element
$c$ spełniający $f_i^m c=f^\sharp(a)$.
Taki element istnieje, ponieważ $a\in I^m$ i
$I\mathcal O_Y=(f_i)$, i jest jednoznaczny,
ponieważ $f_i$ jest niedzielnikiem zera.
Jeśli różne reprezentacje są równe w lokalizacji,
po pomnożeniu przez potęgę $f_iT$ są równe już w
$\mathcal R_A(I)$; ich obrazy są więc równe po
pomnożeniu przez niedzielnik zera, a zatem równe.
Tak skonstruowany homomorfizm daje lokalny morfizm
$Y\to D_+(f_iT)$. Jednoznaczność ilorazów zapewnia
zgodność tych morfizmów na przecięciach, więc
sklejają się one do $\widetilde f:Y\to\widetilde X$.
Każda faktoryzacja musi mieć na mapach identyczne
obrazy elementów $f_j/f_i$, stąd jest jednoznaczna.
\end{proof}

\begin{wniosek}[Rozdmuchanie dzielnika Cartiera niczego nie zmienia]
\label{wn:ga5-cartier-trivial}
Jeżeli $Z\subseteq X$ jest efektywnym dzielnikiem
Cartiera, to $\operatorname{Bl}_Z X\to X$ jest izomorfizmem.
\end{wniosek}
\begin{proof}
Identyczność $\id_X:X\to X$ spełnia założenie własności
uniwersalnej, zatem podnosi się jednoznacznie do
$g:X\to\operatorname{Bl}_Z X$. Z konstrukcji
$\pi\circ g=\id_X$. Zastosujmy teraz tę samą własność
do $Y=\operatorname{Bl}_Z X$: zarówno $g\circ\pi$,
jak i identyczność są faktoryzacjami $\pi$ przez siebie.
Jednoznaczność daje $g\circ\pi=\id$. Są to więc
morfizmy wzajemnie odwrotne.
\end{proof}

Wniosek pokazuje, po co wymuszamy lokalną odwracalność
ideału centrum: rozdmuchanie jest \emph{minimalną}
modyfikacją, po której centrum staje się dzielnikiem Cartiera.

\subsection{Dzielnik wyjątkowy i stożek normalny}

\begin{definicja}[Dzielnik wyjątkowy i algebra stowarzyszona]
Dla rozdmuchania $\pi:\widetilde X\to X$ wzdłuż
$Z=\V(\mathcal I)$ \emph{dzielnikiem wyjątkowym}
(nie zawsze nierozkładalnym) nazywamy podschemat
$E=\widetilde X\times_X Z$.
Jeśli $X=\Spec A$ i $Z=\Spec(A/I)$, definiujemy
algebrę stopniowaną stowarzyszoną z filtracją $I$:
\[
 \operatorname{gr}_I(A)=\bigoplus_{n\ge0}I^n/I^{n+1}.
\]
Jej spektrum afiniczne
$C_{Z/X}=\Spec_Z\operatorname{gr}_I(A)$
nazywamy \emph{stożkiem normalnym} centrum w $X$.
\end{definicja}

\begin{twierdzenie}[Włókno wyjątkowe jest projektywizacją stożka]
\label{tw:ga5-normalcone}
Dla dowolnego ideału skończenie generowanego $I\subseteq A$
zachodzi kanoniczny izomorfizm podschematów nad $Z$:
\[
 E\cong\operatorname{Proj}_{A/I}\operatorname{gr}_I(A).
\]
\end{twierdzenie}
\begin{proof}
Podschemat $E$ jest przeciwobrazem $Z$, więc na poziomie
algebry stopniowanej odpowiada zmianie bazy z $A$ na
$A/I$. Składowa stopnia $n$ algebry Reesa po tej
zmianie to
$I^n\otimes_A(A/I)\cong I^n/(I\cdot I^n)=I^n/I^{n+1}$.
Wszystkie identyfikacje respektują mnożenie, dlatego
$\mathcal R_A(I)\otimes_A A/I\cong\operatorname{gr}_I(A)$.
Konstrukcja $\operatorname{Proj}$ jest zgodna z taką
zmianą bazy, co widać na mapach
$D_+(f_iT)$ przez lokalizację jednorodną i wzięcie
składowej stopnia zero. To daje żądany izomorfizm.
\end{proof}

\begin{przyklad}[Punkt w $\A^n$]
Dla $A=k[x_1,\ldots,x_n]$ i
$I=(x_1,\ldots,x_n)$ klasy jednomianów stopnia $m$
tworzą bazę $I^m/I^{m+1}$. Stąd
$\operatorname{gr}_I(A)\cong k[X_1,\ldots,X_n]$
i $E\cong\PP^{n-1}_k$.
Jeden punkt w $n$-wymiarowym otoczeniu zostaje
zastąpiony przestrzenią jednowymiarowych kierunków.
\end{przyklad}

\begin{uwaga}[Kiedy dzielnik wyjątkowy jest projektywizacją wiązki normalnej?]
Jeżeli $Z\hookrightarrow X$ jest regularnym zanurzeniem
(np. gładkim podschematem gładkiej rozmaitości nad $k$),
to standardowy izomorfizm algebr stopniowanych
$\operatorname{Sym}_{\mathcal O_Z}(\mathcal I/\mathcal I^2)
\cong\bigoplus_{n\ge0}\mathcal I^n/\mathcal I^{n+1}$
daje
$E\cong\operatorname{Proj}_Z
\operatorname{Sym}(\mathcal I/\mathcal I^2)$.
Jeżeli kowymiar $Z$ w $X$ wynosi $c$, jego włókna są
izomorficzne z $\PP^{c-1}$. Równość algebr
symetrycznej i stowarzyszonej dla regularnego zanurzenia
jest osobnym twierdzeniem algebry przemiennej;
poniżej sprawdzimy ją bezpośrednio dla centrum
zadawanego współrzędnymi.
\end{uwaga}

\section{Rozdmuchanie płaszczyzny w punkcie: pełny model}
\label{sec:ga5-plane}

Niech $A=k[x,y]$, $I=(x,y)$. Z generatorów
$U=xT$ oraz $V=yT$ algebry Reesa uzyskujemy
równanie jednorodne $xV-yU=0$. Opis geometryczny
rozdmuchania ma więc postać
\[
 \widetilde{\A^2}
 =\{((x,y),[U:V])\in\A^2\times\PP^1:\;xV=yU\},
 \qquad \pi((x,y),[U:V])=(x,y).
\]
Że to dokładnie algebra Reesa, a nie tylko zbiór
punktów spełniających potrzebne równanie, sprawdzimy
przez porównanie obu map z formułą
stwierdzenia~\ref{st:ga5-blowup-charts}.

\begin{twierdzenie}[Dwie mapy współrzędnych płaszczyzny rozdmuchanej]
\label{tw:ga5-plane-charts}
Niech $k$ będzie ciałem, $X=\A^2_k=\Spec k[x,y]$,
a $\pi:\widetilde X=\operatorname{Bl}_{(x,y)}X\to X$
rozdmuchaniem początku. Schemat $\widetilde X$ pokrywają
dwa otwarte podschematy afiniczne, każdy izomorficzny z
$\A^2_k$:
\begin{enumerate}[label=(\roman*)]
\item na $U\ne0$ mamy współrzędne $(x,v)$, gdzie
$v=V/U$ (na $x\ne0$ jest to $y/x$), równanie $y=xv$ i
$\pi(x,v)=(x,xv)$;
\item na $V\ne0$ mamy współrzędne $(u,y)$, gdzie
$u=U/V$ (na $y\ne0$ jest to $x/y$), równanie $x=uy$ i
$\pi(u,y)=(uy,y)$.
\end{enumerate}
Na przecięciu $u=v^{-1}$ oraz $y=xv$.
Dzielnik wyjątkowy jest równy $x=0$ na pierwszej
mapie, $y=0$ na drugiej i globalnie jest izomorficzny
z $\PP^1_k$. Całe rozdmuchanie jest gładką powierzchnią.
\end{twierdzenie}
\begin{proof}
W równaniu $xV=yU$ na $U\ne0$ możemy podzielić przez
$U$ i otrzymujemy $y=x(V/U)=xv$. Eliminacja $y$
daje pierścień współrzędnych $k[x,v]$, dokładnie
$k[x,y,y/x]$ z definicji algebraicznej. Analogicznie
na $V\ne0$ otrzymujemy $k[u,y]$. Na przecięciu
$u$ i $v$ są odwracalne oraz $uv=1$.

Przeciwobraz schematowy początku wyznacza ideał
$(x,y)$, który na pierwszej mapie jest równy $(x)$,
a na drugiej $(y)$. Oba jego równania opisują
$\A^1$; sklejenie przez $u=v^{-1}$ daje $\PP^1$.
Każdy z dwóch otwartych podschematów jest płaszczyzną afiniczną, więc gładką;
gładkość jest własnością lokalną i obejmuje całe
rozdmuchanie.
\end{proof}

\begin{center}
\begin{tikzpicture}[>=Stealth,scale=0.90,
 every node/.style={font=\small}]
 \draw[->] (-5.0,0)--(-1.5,0) node[below right] {$x$};
 \draw[->] (-3.25,-1.35)--(-3.25,1.35) node[above] {$y$};
 \fill (-3.25,0) circle (2.1pt);
 \draw[blue!65!black,thick] (-4.60,-1.05)--(-1.9,1.05);
 \draw[blue!65!black,thick] (-4.60,0.85)--(-1.9,-0.85);
 \node[align=center,text width=4.7cm] at (-3.25,-2.15)
 {Płaszczyzna przed rozdmuchaniem\\dwie proste, wspólny punkt};
 \draw[->,thick] (-0.90,0)--(0.90,0);
 \node[above] at (0,0.1) {$\pi$};
 \draw[very thick,red!70!black] (3.25,-1.25)--(3.25,1.25);
 \draw[blue!65!black,thick] (1.7,0.58)--(4.80,0.58);
 \draw[blue!65!black,thick] (1.7,-0.58)--(4.80,-0.58);
 \fill[blue!65!black] (3.25,0.58) circle (2.1pt);
 \fill[blue!65!black] (3.25,-0.58) circle (2.1pt);
 \node[above] at (3.25,1.39) {$E\cong\PP^1$};
 \node[align=center,text width=4.7cm] at (3.25,-2.15)
 {Po rozdmuchaniu\\dwa różne punkty na $E$};
\end{tikzpicture}

{\small Rysunek 5.1. Schemat dwóch prostych i dwóch różnych punktów na dzielniku wyjątkowym. W lokalnej mapie rozdmuchania ich transformaty są rozłączne i mają równania $v=m_1$ oraz $v=m_2$ dla $m_1\ne m_2$. Czerwony odcinek symbolizuje fragment $E\cong\PP^1$, a nie całą powierzchnię.}\par
\end{center}

\begin{przyklad}[Granica prostej o ustalonym nachyleniu]
Dla $m\in k$ i punktów $(x,mx)$ z $x\ne0$ na
pierwszej mapie mamy $v=m$. Przy $x\to0$ trafiamy
w punkt $(0,m)\in E$, czyli w kierunek $[1:m]$.
Dla pionowej prostej $x=0$ używamy drugiej mapy;
odpowiada kierunkowi $[0:1]$. W ten sposób wszystkie
kierunki tworzą $\PP^1$, a nie jedynie $\A^1$.
\end{przyklad}

\begin{przyklad}[Rozdmuchanie gładkiej powierzchni nie tworzy osobliwości]
Zbiór punktów nad początkiem jest nieskończony
($E\cong\PP^1$), dlatego $\pi$ nie jest
izomorfizmem. Mimo to obie mapy współrzędnych są
gładkimi płaszczyznami. Rozdmuchanie nie jest więc
operacją usuwającą \emph{zawsze} pewną osobliwość:
można je wykonać również na przestrzeni gładkiej,
aby rozdzielić kierunki lub uzyskać miejsce na
bardziej przejrzyste przecięcia krzywych.
\end{przyklad}

\section{Transformata całkowita, ścisła i wielomian styczny}
\label{sec:ga5-transform}

\begin{definicja}[Dwie transformaty podschematu]
\label{def:ga5-transform}
Niech $\pi:\widetilde X=\operatorname{Bl}_Z X\to X$
będzie rozdmuchaniem, a $C\subseteq X$ podschematem
zamkniętym, którego żadna składowa nierozkładalna
nie jest zawarta w $Z$. \emph{Transformatą całkowitą}
(ang. \emph{total transform}) $C$ jest przeciwobraz
schematowy $\pi^{-1}(C)=C\times_X\widetilde X$.
\emph{Transformatą ścisłą} $\widetilde C$ jest
schematowe domknięcie podschematu
$\pi^{-1}(C\setminus Z)$ w $\widetilde X$.
Dla krzywej zredukowanej można ją również opisać
jako domknięcie zredukowanego przeciwobrazu
$C\setminus Z$, z naturalną strukturą zredukowaną.
\end{definicja}

\begin{uwaga}[Dlaczego potrzebne są dwa pojęcia?]
Transformata całkowita zachowuje krotności równania:
jeśli $f=0$, to cofamy $f$ przez podstawienie.
Transformata ścisła usuwa składowe leżące w dzielniku
wyjątkowym. Nie wolno tego robić przez
„podzielenie przez $x$” bez sprawdzenia, jaka jest
maksymalna wspólna potęga i gdzie dane dzielenie
jest dozwolone.
\end{uwaga}

\begin{twierdzenie}[Lokalne równanie transformaty ścisłej]
\label{tw:ga5-strict-formula}
Niech $f\in k[x,y]$ będzie niezerowym wielomianem bez
wielokrotnych czynników, spełniającym $f(0,0)=0$,
i niech $m\ge1$ oznacza jego
krotność w początku, czyli najmniejszy stopień
niezerowej części jednorodnej. Zapiszmy
$f=f_m+f_{m+1}+\cdots$, gdzie każdy $f_j$
jest jednorodny stopnia $j$, a $f_m\ne0$.
Na mapie $y=xv$ istnieje wielomian $F(x,v)$ taki, że
\[
 f(x,xv)=x^mF(x,v),\qquad F(0,v)=f_m(1,v).
\]
Transformatę całkowitą wyznacza $x^mF=0$, a transformatę
ścisłą równanie $F=0$. Analogicznie w drugiej mapie
$x=uy$ transformatę ścisłą wyznacza wielomian
$f(uy,y)/y^m$, którego ograniczeniem do $E$ jest
$f_m(u,1)$.
\end{twierdzenie}
\begin{proof}
Dla jednomianu $x^ay^b$ stopnia $j=a+b$ podstawienie
$y=xv$ daje $x^jv^b$. Każdy składnik ma więc czynnik
$x^m$. Po jego wyłączeniu część stała względem $x$
pochodzi dokładnie ze składników stopnia $m$ i wynosi
$f_m(1,v)$. Ten ostatni jest niezerowym wielomianem,
więc $x$ nie dzieli $F$ w $k[x,v]$.

Na otwartym podzbiorze $x\ne0$ transformatę krzywej
wyznacza $F=0$, bo $x^m$ jest odwracalne.
Jej schematowe domknięcie jest wyznaczone przez ideał
$((F):x^\infty)$. Ponieważ $k[x,v]$ jest UFD,
$x$ jest elementem pierwszym i nie dzieli $F$.
Z $x^r h\in(F)$ wynika zatem $F\mid h$:
każdy czynnik nierozkładalny $F$ z odpowiednią
krotnością musi dzielić $h$, bo żaden nie jest
stowarzyszony z $x$. Stąd $((F):x^\infty)=(F)$.
To dowodzi twierdzenia na pierwszej mapie.
Drugi otrzymujemy przez zamianę ról $x$ i $y$.
\end{proof}

\begin{wniosek}[Kierunki styczne i wielomian najniższego stopnia]
\label{wn:ga5-tangent-directions}
Przy założeniach poprzedniego twierdzenia
przecięcie schematowe transformaty ścisłej z $E$
jest na dwóch mapach określone przez
$f_m(1,v)=0$ i $f_m(u,1)=0$. Jest to schemat zer
formy jednorodnej $f_m$ na $\PP^1$.
Zliczone z krotnościami przecięcia odpowiadają
$m$ kierunkom stycznym nad ciałem algebraicznie
domkniętym.
\end{wniosek}
\begin{proof}
Na pierwszej mapie dzielnik wyjątkowy ma równanie
$x=0$, a transformata ścisła równanie $F=0$.
Ich schematowe przecięcie ma pierścień współrzędnych
$k[x,v]/(x,F)\cong k[v]/(f_m(1,v))$.
Analogicznie na drugiej mapie otrzymujemy
$k[u]/(f_m(u,1))$. Są to mapy podschematu
$\V_+(f_m)\subseteq\PP^1$. Nad ciałem algebraicznie
domkniętym forma stopnia $m$ rozkłada się na
$m$ liniowych czynników z krotnościami, więc
schemat zer ma długość $m$.
\end{proof}

\begin{przyklad}[Gdzie widzimy informację o styczności?]
Dla $f=y^2-x^3$ mamy $f_2=y^2$, więc na $E$
pozostaje punkt $[1:0]$ \emph{z krotnością dwa}.
Dla $g=y^2-x^2-x^3$ otrzymujemy
$g_2=y^2-x^2$, a więc dwa różne kierunki
$[1:1]$ i $[1:-1]$. Dwie krzywe mają tę samą
krotność w początku, ale różne rozkłady kierunków
stycznych. To uprzedza rachunek o ostrzu i węźle.
\end{przyklad}

\section{Ostrze, węzeł i kolejne rozdmuchania: rachunek bez skrótów}
\label{sec:ga5-cusp-node}

\subsection[Ostrze: wygładzenie po jednym rozdmuchaniu]{Ostrze $y^2=x^3$: krzywa staje się gładka po jednym rozdmuchaniu}

\begin{przyklad}[Pełny rachunek transformaty ostrza]
\label{prz:ga5-cusp}
Niech $C=\V(y^2-x^3)\subseteq\A^2_k$. Na pierwszej
mapie rozdmuchania $y=xv$ mamy
\[
 y^2-x^3=x^2(v^2-x),\qquad
 \widetilde C:\ v^2-x=0,\qquad E:\ x=0.
\]
Transformata ścisła ma parametryzację
$(x,v)=(t^2,t)$, a więc jest izomorficzna z $\A^1_t$.
Jest gładka, ponieważ $\partial(v^2-x)/\partial x=-1$.
Jej odwzorowanie do pierwotnej krzywej ma postać
$t\mapsto(x,y)=(t^2,t^3)$.

Sprawdźmy, że druga mapa niczego nie ukrywa.
Po podstawieniu $x=uy$ mamy
$y^2-x^3=y^2(1-u^3y)$, więc transformata ścisła
spełnia $1-u^3y=0$. W szczególności $u\ne0$ i
$y=u^{-3}$, a zatem $x=u^{-2}$. Ta mapa odpowiada
otwartemu podzbiorowi $t\ne0$ poprzedniej mapy,
przy $t=1/u$. Tak więc cała transformata jest
rzeczywiście $\A^1_t$, a nie tylko jej część.

Mimo gładkości transformata ścisła \emph{nie przecina
transwersalnie} dzielnika wyjątkowego: w punkcie
$(x,v)=(0,0)$ jej wektor styczny spełnia $dx=0$,
a dzielnik $E$ jest dany przez $x=0$. Obie styczne
są więc jednakowe. Również rachunek przecięcia
daje $k[v]/(v^2)$, czyli długość $2$ zamiast $1$.
To jest powód wykonania dalszych rozdmuchań
w rozwiązaniu \emph{zanurzonym}.
\end{przyklad}

\begin{przyklad}[Węzeł: dwa różne punkty nad osobliwością]
\label{prz:ga5-node}
Niech $k$ będzie ciałem o charakterystyce różnej od $2$
i niech $C=\V(y^2-x^2-x^3)$. Na mapie $y=xv$
otrzymujemy
$y^2-x^2-x^3=x^2(v^2-1-x)$, więc
$\widetilde C:\ x=v^2-1$. Jest to gładka krzywa,
parametryzowana przez $v\in\A^1$.
Nad początkiem mamy $x=0$, a więc dwa punkty
$v=1$ i $v=-1$. W każdym z nich $E$ ma równanie
$x=0$, a $\widetilde C$ spełnia $dx=2v\,dv$.
Ponieważ $2v\ne0$, oba przecięcia z $E$ są
transwersalne. Druga mapa nie dodaje nowych punktów,
gdyż tam równanie transformaty
$1-u^2-yu^3=0$ wymusza $u\ne0$.

Pierwotny punkt węzłowy miał zatem \emph{dwie gałęzie},
które rozdmuchanie rozdzieliło. W przypadku ostrza
pozostawała tylko jedna gałąź, ale styczna do $E$.
\end{przyklad}

\begin{center}
\begin{tikzpicture}[>=Stealth,scale=0.93,every node/.style={font=\small}]
  % left panel
  \begin{scope}[shift={(-3.3,0)}]
    \draw[gray!65,->] (-1.9,0)--(1.9,0) node[below] {$v$};
    \draw[gray!65,->] (0,-1.5)--(0,1.45) node[left] {$x$};
    \draw[very thick,red!70!black] (-1.8,0)--(1.8,0);
    \draw[domain=-1.1:1.1,smooth,variable=\t,thick,blue!65!black]
      plot ({\t},{\t*\t});
    \fill (0,0) circle (2.5pt);
    \node[red!70!black,below right] at (1.22,0) {$E$};
    \node[above] at (0,1.55) {Ostrze};
    \node[align=center] at (0,-1.94) {$x=v^2$; jedno styczne\\przecięcie z $E$};
  \end{scope}
  % right panel
  \begin{scope}[shift={(2.45,0)}]
    \draw[gray!65,->] (-1.9,0)--(1.9,0) node[below] {$v$};
    \draw[gray!65,->] (0,-1.5)--(0,1.45) node[left] {$x$};
    \draw[very thick,red!70!black] (-1.8,0)--(1.8,0);
    \draw[domain=-1.37:1.37,smooth,variable=\t,thick,blue!65!black]
      plot ({\t},{\t*\t-1});
    \fill (-1,0) circle (2.5pt);
    \fill (1,0) circle (2.5pt);
    \node[red!70!black,below right] at (1.30,0) {$E$};
    \node[above] at (0,1.55) {Węzeł};
    \node[align=center] at (0,-1.94) {$x=v^2-1$; dwa różne\\przecięcia z $E$};
  \end{scope}
\end{tikzpicture}

{\small Rysunek 5.2. Dwa przykłady w tej samej mapie $U\ne0$ rozdmuchania płaszczyzny. Kolor czerwony oznacza dzielnik wyjątkowy, niebieski transformatę ścisłą. Węzeł ma dwa różne kierunki styczne; ostrze tylko jeden, liczony podwójnie. Skala pozioma w obu panelach jest umowna.}\par
\end{center}

\subsection{Dlaczego dla ostrza potrzeba jeszcze dwóch rozdmuchań?}

\begin{definicja}[Proste skrzyżowania normalne i skrzyżowania normalne]
\label{def:ga5-snc}
Niech $X$ będzie gładkim schematem nad ciałem $k$, a
$D\subseteq X$ zredukowanym efektywnym dzielnikiem Cartiera.
Mówimy, że $D$ ma \emph{proste skrzyżowania normalne}
(ang. \emph{simple normal crossings}, SNC), jeżeli dla każdego
$p\in |D|$ istnieje układ parametrów regularnych
$z_1,\ldots,z_d$ maksymalnego ideału pierścienia
regularnego $\mathcal O_{X,p}$, gdzie
$d=\dim\mathcal O_{X,p}$, taki że ideał $D$ w tym
pierścieniu jest generowany przez $z_1\cdots z_r$
dla pewnego $1\le r\le d$.

Dzielnik ma \emph{skrzyżowania normalne} (NC), jeżeli
warunek SNC zachodzi po odpowiednim przejściu étale:
dla każdego $p$ istnieje morfizm étale $U\to X$,
którego obraz zawiera $p$, i taki że $D\times_XU$ jest
na $U$ dzielnikiem SNC.
\end{definicja}

W przypadku gładkiej powierzchni nad $\C$ warunek SNC
oznacza, że lokalnie w pierścieniu algebraicznym występuje
jedna gładka gałąź $z_1=0$ albo dwie różne, gładkie,
transwersalne gałęzie $z_1z_2=0$. Dopuszczenie jedynie
współrzędnych étale jest słabsze: prowadzi do NC, a nie
koniecznie do SNC. W szczególności globalnie nierozkładalna
krzywa węzłowa, np. $y^2=x^2(1+x)$ w $\A^2_{\C}$,
ma lokalnie po odpowiednim pokryciu étale dwie gałęzie,
lecz nie jest dzielnikiem SNC w swojej pierwotnej
płaszczyźnie. Powodem jest nierozkładalność jej lokalnego
równania w pierścieniu algebraicznym
$\mathcal O_{\A^2,(0,0)}$; formalny lub étale rozkład
nie zastępuje rozkładu w tym pierścieniu.

W konfiguracji SNC na gładkiej powierzchni dopuszczalne jest
transwersalne przecięcie dwóch gładkich krzywych, lecz nie
przecięcie trzech różnych składowych w jednym punkcie.
Takie przecięcia są oczywiście możliwe na samej powierzchni,
ale ich suma nie ma SNC. Styczne przecięcie dwóch krzywych
również nie jest SNC.

\begin{center}
\begin{tikzpicture}[scale=0.82,every node/.style={font=\small}]
 \begin{scope}[shift={(-4.7,0)}]
   \draw[thick,blue!65!black] (-1.15,-1.0)--(1.15,1.0);
  \draw[thick,red!65!black] (-1.15,1.0)--(1.15,-1.0);
  \node[align=center,text width=3.7cm] at (0,-1.8)
  {Dwa kierunki niezależne\\SNC};
 \end{scope}
 \begin{scope}
  \draw[thick,red!65!black] (-1.3,0)--(1.3,0);
  \draw[domain=-1.1:1.1,smooth,variable=\t,thick,blue!65!black]
   plot ({\t},{0.9*\t*\t});
  \node[align=center,text width=3.7cm] at (0,-1.8)
  {Dwie krzywe styczne\\brak SNC};
 \end{scope}
 \begin{scope}[shift={(4.7,0)}]
  \draw[thick,blue!65!black] (-1.15,-1.0)--(1.15,1.0);
  \draw[thick,red!65!black] (-1.15,1.0)--(1.15,-1.0);
  \draw[thick,gray!75!black] (-1.3,0)--(1.3,0);
  \node[align=center,text width=3.7cm] at (0,-1.8)
  {Trzy proste przez punkt\\brak SNC na powierzchni};
 \end{scope}
\end{tikzpicture}

{\small Rysunek 5.3. Trzy lokalne konfiguracje dzielników zredukowanych na gładkiej powierzchni. Dopuszczalne są tylko proste skrzyżowania normalne, jak w lewym panelu.}\par
\end{center}

\begin{twierdzenie}[Trzy rozdmuchania dają rozwiązanie zanurzone ostrza]
\label{tw:ga5-embedded-cusp}
Dla krzywej $C=\V(y^2-x^3)$ w gładkiej powierzchni
$\A^2_k$ istnieje ciąg trzech rozdmuchań w punktach,
po którym transformata ścisła $C^{(3)}$ jest gładka,
a zredukowana suma $C^{(3)}$ i dzielników wyjątkowych
ma proste skrzyżowania normalne. Pierwsze rozdmuchanie
wygładza samą krzywą, ale nie wystarcza do SNC.
\end{twierdzenie}
\begin{proof}
\emph{Krok 1.} Współrzędne pierwszej mapy to
$(x,v)$, $y=xv$. Transformata ścisła ma równanie
$C^{(1)}:\ x=v^2$, a pierwszy dzielnik wyjątkowy
$E_1:\ x=0$. Obie krzywe są gładkie i styczne
w $(0,0)$. Dlatego rozdmuchujemy ten punkt.

\emph{Krok 2.} Na mapie drugiego rozdmuchania,
w którym $x=vw$, mamy
$C^{(1)}:\ vw=v^2$. Po usunięciu czynnika
wyjątkowego $v$ otrzymujemy
$C^{(2)}:\ w=v$. Nowy dzielnik wyjątkowy ma
równanie $E_2:\ v=0$, a transformata ścisła
poprzedniego dzielnika $E_1:\ x=0$ ma równanie
$E_1':\ w=0$. Zatem przez początek mapy
przechodzą \emph{trzy} różne gładkie krzywe
$w=v$, $v=0$ i $w=0$. Każde dwie przecinają się
transwersalnie, ale ich suma nie jest SNC na powierzchni.
Na drugiej mapie tego rozdmuchania, gdzie $v=xu$,
transformata ścisła $C^{(2)}$ spełnia $xu^2=1$,
więc nie ma tam punktu nad rozdmuchiwanym centrum.
Wszystkie pozostałe przecięcia rozważanych krzywych
były już SNC lub nie występowały.

\emph{Krok 3.} Rozdmuchujemy jedyny punkt potrójny
$(v,w)=(0,0)$. Na mapie $v=a$, $w=ab$
transformata $C^{(2)}$ spełnia
$a(b-1)=0$, zatem jej transformata ścisła ma
równanie $b=1$. Nowy dzielnik wyjątkowy ma
równanie $E_3:\ a=0$, natomiast transformata $E_1'$
ma równanie $b=0$. Krzywa $C^{(3)}$ przecina
$E_3$ w $(a,b)=(0,1)$, gdzie $E_1''$ nie przechodzi.
Przecięcie jest transwersalne, ponieważ lokalne
równania $a=0$ oraz $b-1=0$ mają niezależne
różniczki $da$ i $db$.

Na drugiej mapie $v=wc$, $w=d$ transformata
$E_2$ ma równanie $c=0$, a $E_3$ ma równanie
$d=0$; przecinają się transwersalnie. W tej mapie
$C^{(3)}$ spełnia $1=c$, więc przecina $E_3$
w punkcie $(d,c)=(0,1)$, który jest tym samym
punktem co $(a,b)=(0,1)$ w poprzedniej mapie:
na przecięciu $c=1/b$, a $b=1$.
Transformata $E_1''$ nie przechodzi przez ten punkt,
bo w pierwszej mapie ma $b=0$. Spotkania
$E_1''$ z $E_3$ i $E_2''$ z $E_3$ zachodzą
w dwóch innych punktach, odpowiadających różnym
kierunkom $[v:w]$, i każde jest transwersalne.
Nie ma już przecięć potrójnych ani stycznych.
Zatem zredukowany całkowity dzielnik ma SNC.
\end{proof}

\begin{uwaga}[Uwaga o krotnościach dzielnika całkowitego]
Po kolejnych podstawieniach transformata \emph{całkowita}
równania $y^2-x^3$ ma wielokrotne składowe wyjątkowe.
Po drugim rozdmuchaniu jej lokalne równanie to
$v^3w^2(v-w)=0$, a po trzecim (w mapie $v=a$,
$w=ab$) ma postać $a^6b^2(1-b)=0$.
Warunek SNC w poprzednim twierdzeniu dotyczy
\emph{nośnika zredukowanego} dzielnika całkowitego,
nie twierdzi, że jego składowe występują z krotnością jeden.
W teorii rozwiązywania osobliwości ta różnica jest
kluczowa dla obliczeń rozbieżności, całek lokalnych
i niezmienników singularności.
\end{uwaga}

\begin{przyklad}[Taknod: po pierwszym rozdmuchaniu problem nie znika]
Niech $C=\V(y^2-x^4)$. Nad $k$ o charakterystyce
różnej od $2$ ma dwie gałęzie $y=x^2$ i $y=-x^2$,
styczne do siebie w początku. Po podstawieniu
$y=xv$ otrzymujemy transformatę ścisłą
$v^2-x^2=(v-x)(v+x)=0$. Są to dwie gładkie,
przecinające się gałęzie, ale obie przechodzą
przez ten sam punkt pierwszego dzielnika wyjątkowego.
Drugie rozdmuchanie rozdziela ich kierunki
$v=x$ i $v=-x$, dzięki czemu obie transformaty
spotykają nowy dzielnik wyjątkowy w różnych punktach.
Jest to przykład, w którym zmniejszanie krotności
równania nie jest tym samym co zakończenie
rozwiązywania zanurzonego.
\end{przyklad}

\subsection{Normalizacja a rozdmuchanie krzywej}

\begin{stwierdzenie}[Ostrze i węzeł: normalizacje policzone jawnie]
\label{st:ga5-normalization}
Niech $k$ będzie ciałem algebraicznie domkniętym
charakterystyki różnej od $2$. Wówczas:
\begin{enumerate}[label=(\roman*)]
\item normalizacją $k[x,y]/(y^2-x^3)$ jest $k[t]$
przy $x=t^2$, $y=t^3$;
\item normalizacją $k[x,y]/(y^2-x^2-x^3)$ jest
$k[t]$ przy $x=t^2-1$, $y=t(t^2-1)$.
\end{enumerate}
W obu przypadkach transformata ścisła po pierwszym
rozdmuchaniu płaszczyzny jest izomorficzna z $\A^1_t$
i realizuje tę normalizację. Preobrazy punktu osobliwego
mają odpowiednio jeden punkt $t=0$ oraz dwa punkty
$t=1$ i $t=-1$.
\end{stwierdzenie}
\begin{proof}
W (i) z równania $y^2=x^3$ wynika w ciele ułamków
$t=y/x$, $t^2=x$ oraz $t^3=y$. Zatem
$A\subseteq k[t]\subseteq\Frac(A)$ i $t$ jest całkowity
nad $A$, bo spełnia równanie moniczne $T^2-x=0$.
W szczególności $k[t]$ jest skończonym modułem nad
$A$, generowanym przez $1,t$. Jest całkowicie
domknięty, bo $k[t]$ to pierścień euklidesowy.
Jeżeli $s\in\Frac(A)$ jest całkowity nad $A$, to jest
też całkowity nad $k[t]$ (to samo równanie moniczne
ma współczynniki w $k[t]$), więc $s\in k[t]$.
Stąd $k[t]$ jest domknięciem całkowitym $A$.
Porównanie ze wzorami z przykładu~\ref{prz:ga5-cusp}
pokazuje, że jest też pierścieniem transformaty ścisłej.

W (ii) podobnie $t=y/x$ spełnia $t^2=1+x$,
skąd $x=t^2-1$, $y=t(t^2-1)$; zatem
$A\subseteq k[t]\subseteq\Frac(A)$, a $t$ jest
całkowity nad $A$. Ten sam argument daje normalizację.
W pierwszej mapie rozdmuchania relacja
$x=v^2-1$ pokazuje, że transformata ścisła ma
pierścień $k[v]$ i odpowiada $t=v$.
Ostatnie twierdzenie o włóknach to bezpośrednie
rozwiązanie równań $t^2=0$ oraz $t^2-1=0$.
\end{proof}

\begin{przyklad}[Schematowe przecięcie a liczba gałęzi]
Dla węzła $y^2=x^2+x^3$ transformata ścisła po
rozdmuchaniu ma w mapie $y=xv$ równanie
$x=v^2-1$. Jej przecięcie z dzielnikiem wyjątkowym
$x=0$ ma pierścień $k[v]/(v^2-1)\cong k\times k$
(przy $\operatorname{char}k\ne2$): dwa różne punkty
zredukowane. Dla ostrza $y^2=x^3$ analogiczny
pierścień wynosi $k[v]/(v^2)$: jeden punkt
niezredukowany długości dwa. Równość sumy krotności
przecięć nie oznacza więc równości liczby gałęzi.
Ten rachunek pokazuje, dlaczego nie wolno zastępować
przecięcia schematowego samym zbiorem punktów.
\end{przyklad}

\begin{uwaga}[Dlaczego normalizacja nie rozwiązuje wszystkiego?]
Normalizacja jest skończonym morfizmem biwymiernym
z przestrzeni normalnej; dla krzywych skończonego typu
nad ciałem doskonałym normalność oznacza gładkość.
W wyższych wymiarach \emph{normalność nie oznacza
gładkości}. Na przykład stożek kwadratowy
$z^2=x^2+y^2$ w $\A^3_\C$ ma izolowaną osobliwość
w początku, mimo że jego pierścień współrzędnych
jest normalny (można to sprawdzić kryterium Serre'a
$R_1+S_2$, którego tu nie dowodzimy).
Dlatego rozwiązywanie osobliwości powierzchni
wymaga więcej niż samej normalizacji.
\end{uwaga}

\section{Przypadek wyższych wymiarów: gładkie centra i stożek}
\label{sec:ga5-higherdim}

\begin{twierdzenie}[Rozdmuchanie przestrzeni afinicznej wzdłuż podprzestrzeni]
\label{tw:ga5-smooth-center-model}
Niech $1\le c\le n$ i niech
$X=\A^n_k=\Spec k[x_1,\ldots,x_n]$.
Dla centrum $Z=\V(x_1,\ldots,x_c)$ rozdmuchanie
$\widetilde X=\operatorname{Bl}_Z X$ jest gładkie.
Na mapie odpowiadającej $x_i$, $1\le i\le c$,
ma współrzędne
\[
 x_i,\quad u_j=x_j/x_i\ (1\le j\le c,\ j\ne i),
 \quad x_{c+1},\ldots,x_n.
\]
Dzielnik wyjątkowy w tej mapie jest zadany
przez $x_i=0$. Dla każdego punktu $z\in Z$ włókno
$\widetilde X\times_X\Spec\kappa(z)$ jest izomorficzne
z $\PP^{c-1}_{\kappa(z)}$. W szczególności dla punktów
$k$-wymiernych jest to $\PP^{c-1}_k$.
\end{twierdzenie}
\begin{proof}
Z definicji algebry Reesa mapy mają postać
$k[x_1,\ldots,x_n][x_j/x_i:1\le j\le c]$.
Po zastąpieniu $x_j$ przez $x_i u_j$ dla $j\ne i$
otrzymujemy pierścień wielomianów w dokładnie
$n$ algebraicznie niezależnych zmiennych podanych
w tezie. Otrzymujemy otwarty podschemat izomorficzny
z $\A^n_k$, a więc gładki.
Ideał centrum staje się $(x_i)$ i wyznacza dzielnik
$x_i=0$. Po nałożeniu relacji $x_1=\cdots=x_c=0$
na wszystkie mapy pozostaje kolekcja współrzędnych
$u_j$ z przejściami jak dla $\PP^{c-1}$,
niezależna od pozostałych współrzędnych centrum.
Zatem $E\cong Z\times\PP^{c-1}$ w tym liniowym modelu.
\end{proof}

\begin{wniosek}[Gładkie centrum w gładkiej przestrzeni]
\label{wn:ga5-smooth-center}
Niech $X$ i $Z\subseteq X$ będą gładkimi rozmaitościami
algebraicznymi nad ciałem $k$ charakterystyki $0$,
a zanurzenie $Z\hookrightarrow X$ ma stały kowymiar $c$.
Wtedy $\operatorname{Bl}_Z X$ jest gładkie, a jego
dzielnik wyjątkowy jest lokalnie nad $Z$ wiązką
przestrzeni rzutowych, której włókno nad $z\in Z$
jest izomorficzne z $\PP^{c-1}_{\kappa(z)}$.
\end{wniosek}
\begin{proof}
Potrzebny jest standardowy lokalny model zanurzenia
gładkiego: w otoczeniu étale każdego punktu $Z$
para $(X,Z)$ jest równoważna parze
$(\A^n,\A^{n-c})$, z centrum wyznaczonym przez
pierwszych $c$ współrzędnych. Jest to geometryczna
wersja algebraicznego twierdzenia o funkcji uwikłanej.
Rozdmuchanie jest zgodne z płaską zmianą bazy
(twierdzenie~\ref{tw:ga5-flat-basechange} poniżej),
a morfizmy étale są płaskie. Po tej zmianie bazy
mapy mają postać z poprzedniego twierdzenia.
Gładkość jest lokalna w topologii étale, zatem
$\operatorname{Bl}_Z X$ jest gładkie. Mapy współrzędnych dzielnika
wyjątkowego sklejają się do projektywizacji wiązki
normalnej. Przyjęty zewnętrzny składnik dowodu to
wskazane twierdzenie o lokalnym modelu étale;
nie jest ono ukryte w rachunku współrzędnych.
\end{proof}

\begin{przyklad}[Rozdmuchanie prostej w $\A^3$]
Dla $X=\A^3_{x,y,z}$ i $Z=\{x=y=0\}$
mamy $I=(x,y)$. Na mapie $y=xv$ współrzędnymi
są $(x,v,z)$, a dzielnik wyjątkowy ma równanie
$x=0$. Nad każdym punktem $(0,0,z_0)$ w centrum
pojawia się $\PP^1$. Różnica wobec rozdmuchania
punktu jest taka, że dzielnik wyjątkowy jest teraz
powierzchnią, nie krzywą.
\end{przyklad}

\begin{przyklad}[Rozwiązanie izolowanej osobliwości stożka]
\label{prz:ga5-cone}
Niech $X\subseteq\A^3_k$ będzie dane przez
$z^2=x^2+y^2$. W charakterystyce różnej od $2$
pochodne równania to $-2x,-2y,2z$, więc jedynym
punktem osobliwym jest początek. Rozdmuchujemy
\emph{przestrzeń otaczającą} $\A^3$ w początku
i bierzemy transformatę ścisłą $X$.

Na mapie $x\ne0$ kierunku (nie należy mylić tego
z warunkiem $x\ne0$ na samej przestrzeni rozdmuchanej)
piszemy $y=xu$, $z=xv$; transformata ścisła ma
równanie $v^2=1+u^2$. Jej gradient względem
$(u,v)$ to $(-2u,2v)$, który nie może się zerować
na krzywej $v^2-u^2=1$.

Na mapie kierunku $z\ne0$ piszemy
$x=za$, $y=zb$ i otrzymujemy $1=a^2+b^2$,
również gładkie równanie; mapa kierunku $y\ne0$
jest analogiczna do pierwszej mapy. Cała transformata
ścisła jest więc gładką powierzchnią.
Jej przecięcie z wyjątkową $\PP^2$ jest stożkową
$[X:Y:Z]\in\PP^2$ o równaniu $Z^2=X^2+Y^2$,
gładką, a zatem izomorficzną z $\PP^1$.
To pokazuje, jak jeden punkt osobliwy powierzchni
zostaje zastąpiony całą gładką krzywą.
\end{przyklad}

\section{Zmiana bazy, geometria biwymierna i granice metody}
\label{sec:ga5-basechange}

\begin{twierdzenie}[Rozdmuchanie i płaska zmiana bazy]
\label{tw:ga5-flat-basechange}
Niech $A\to B$ będzie homomorfizmem płaskim,
a $I\subseteq A$ ideałem skończenie generowanym.
Wówczas istnieje naturalny izomorfizm schematów nad
$\Spec B$:
\[
 \operatorname{Bl}_{IB}\Spec B\cong
 \operatorname{Bl}_{I}\Spec A\times_{\Spec A}\Spec B.
\]
Analogiczne twierdzenie zachodzi dla płaskiej zmiany
bazy dowolnych schematów i quasi-koherentnych
snopów ideałów skończonego typu.
\end{twierdzenie}
\begin{proof}
Z płaskości wynika, że po zastosowaniu iloczynu tensorowego z $B$
każdy dokładny ciąg
$0\to I^n\to A\to A/I^n\to0$ pozostaje dokładny.
Stąd $I^n\otimes_A B$ identyfikuje się z ideałem
$I^nB=(IB)^n\subseteq B$. Mamy zatem izomorfizm
stopniowanych $B$-algebr
\[
 \mathcal R_A(I)\otimes_A B\cong
 \bigoplus_{n\ge0}(IB)^nT^n=\mathcal R_B(IB).
\]
Po zastosowaniu $\operatorname{Proj}$, który zachowuje
zmianę bazy, otrzymujemy izomorfizm z tezy.
Dla schematów argument stosujemy na afinicznych
otwartych i sklejamy otrzymane izomorfizmy.
\end{proof}

\begin{uwaga}[Dlaczego płaskość jest ważna?]
Bez płaskości odwzorowanie $I^n\otimes_A B\to B$
może mieć niezerowe jądro. Nie ma wówczas podstaw,
by utożsamiać wyrazy algebry Reesa po zmianie bazy
z potęgami ideału $IB$. Dlatego twierdzenie nie mówi,
że rozdmuchanie komutuje z \emph{dowolną} zmianą bazy.
\end{uwaga}

\begin{przyklad}[Konkretny brak zgodności przy zmianie bazy]
Niech $A=k[x,y]$, $I=(x,y)$ oraz
$B=A/(x)\cong k[y]$. Gdyby zmiana bazy była zawsze
zgodna z rozdmuchaniem, dostalibyśmy rozdmuchanie
prostej $\Spec k[y]$ w centrum $(y)$, które jest
izomorfizmem (wniosek~\ref{wn:ga5-cartier-trivial}).
Tymczasem
$\operatorname{Bl}_{(x,y)}\A^2\times_{\A^2}\V(x)$
na pierwszej mapie $y=xv$ spełnia po podstawieniu
$x=0$ także $y=0$, więc pozostaje cały
otwarty podschemat $\Spec k[v]$, odwzorowywany w początek.
Na drugiej mapie $x=uy$ po podstawieniu $x=0$
dostajemy równanie $uy=0$: oprócz transformaty
prostej $u=0$ znajduje się składowa wyjątkowa
$y=0$. Jest to większy schemat niż sama prosta.
\end{przyklad}

\begin{uwaga}[Wyjątkowość, samoprzecięcie i znak minus]
Dla rozdmuchania gładkiej powierzchni w punkcie
$E\cong\PP^1$ oraz
$\mathcal O_{\widetilde X}(E)|_E\cong\mathcal O_{\PP^1}(-1)$.
Gdy powierzchnia jest właściwa i ma zdefiniowane
liczby przecięcia, wynika stąd $E\cdot E=-1$.
Znak minus jest ważny w klasyfikacji powierzchni:
krzywe wyjątkowe pierwszego rodzaju można następnie
\emph{zdmuchiwać} (ang. \emph{blow-down}).
Nie traktujemy tu tego faktu jako udowodnionego:
wymaga teorii wiązek normalnych i liczb przecięcia,
którą rozwiniemy w rozdziałach o powierzchniach
i teorii przecięć.
\end{uwaga}

\subsection{Co mówi twierdzenie Hironaki?}

\begin{twierdzenie}[Hironaka: istnienie rozwiązania osobliwości,
wersja w charakterystyce zero]
\label{tw:ga5-hironaka}
Niech $k$ będzie ciałem charakterystyki $0$, a $X$
zredukowanym, rozdzielonym schematem skończonego typu nad $k$.
Wówczas istnieje właściwy morfizm
$\pi:\widetilde X\to X$ z gładkiego schematu
$\widetilde X$, który jest izomorfizmem nad gładkim
miejscem $X$ i jest biwymierny na każdej składowej
nierozkładalnej. W wersji zanurzonej dla zredukowanej
hiperpowierzchni $D$ w gładkiej rozmaitości $M$
można uzyskać ten efekt poprzez skończony ciąg
rozdmuchań odpowiednio dobranych gładkich centrów,
tak aby transformata ścisła $D$ była gładka, a
zredukowana suma transformaty ścisłej i dzielników
wyjątkowych miała proste skrzyżowania normalne.
\end{twierdzenie}

\begin{uwaga}[Granica samodzielności: status dowodu]
Nie podajemy tu dowodu pełnego twierdzenia Hironaki.
W przeciwieństwie do obliczeń krzywych płaskich nie jest
to konsekwencja kilku podstawień współrzędnych.
Konstrukcja wymaga wyboru centrów na podstawie
zmniejszających się niezmienników osobliwości,
kontroli ich zachowania przy rozdmuchaniach
i indukcji po wymiarze oraz po tych niezmiennikach.
Powyższy wynik cytujemy jawnie jako twierdzenie
zewnętrzne. Udowodniliśmy natomiast bezpośrednio
wszystkie szczególne obliczenia ostrza, węzła
i stożka podane w tym rozdziale.
\end{uwaga}

\begin{uwaga}[Dlaczego charakterystyka zero nie jest ozdobnikiem?]
W dodatniej charakterystyce istnieją rozległe
teorie rozwiązywania osobliwości w niskich wymiarach,
jednak pełna wersja twierdzenia w dowolnym wymiarze
nie jest znana w tej ogólności. W rachunku
współrzędnych dodatnia charakterystyka może powodować
znikanie pochodnych potęg $p$; w metodzie globalnej
trudniejszym problemem jest zachowanie wielkości
mierzących złożoność osobliwości po rozdmuchaniu.
\end{uwaga}

\begin{uwaga}[Rozdmuchania a biracjonalność]
Gdy centrum nie zawiera żadnej składowej $X$, rozdmuchanie
jest izomorfizmem nad gęstym otwartym podzbiorem
$X\setminus Z$, a zatem jest morfizmem biwymiernym,
gdy $X$ jest całkowity. Nie każdy morfizm biwymierny
jest pojedynczym rozdmuchaniem, a nie każda
biracjonalna modyfikacja ma postać rozdmuchania
gładkiego centrum. Rozkładanie ich na elementarne
operacje to osobny problem geometrii biwymiernej.
\end{uwaga}

\section{Mapa pojęć i perspektywa następnego rozdziału}
\label{sec:ga5-outlook}

W rozdziale przeprowadziliśmy następujący łańcuch idei:
\emph{ideał centrum $I$ $\longrightarrow$ potęgi $I^n$
$\longrightarrow$ algebra Reesa $\longrightarrow$
$\operatorname{Proj}$ $\longrightarrow$ rozdmuchanie
$\longrightarrow$ transformaty i dzielniki wyjątkowe}.
Właściwość uniwersalna mówi, że rozdmuchanie jest
kanoniczną operacją czyniącą ideał centrum odwracalnym.
W praktyce na gładkich powierzchniach rachunek ten
sprowadza się do dwóch prostych podstawień
$y=xv$ i $x=yu$. Dopiero śledzenie przecięć
z dzielnikami wyjątkowymi wyjaśnia, ile takich operacji
trzeba wykonać, aby uzyskać konfigurację SNC.

\begin{center}
\begin{minipage}{\linewidth}
\centering
\begin{tikzpicture}[>=Stealth,node distance=9mm,
 every node/.style={font=\small},
 b/.style={draw,rounded corners,align=center,text width=7.8cm,inner sep=5pt}]
 \node[b] (a) {Centrum $Z=\V(\mathcal I)$ i filtracja $\mathcal I^0,\mathcal I^1,\ldots$};
 \node[b,below=of a] (b1) {Algebra Reesa $\mathcal R(\mathcal I)=\bigoplus_{n\ge0}\mathcal I^nT^n$};
 \node[b,below=of b1] (c) {$\widetilde X=\operatorname{Proj}_X\mathcal R(\mathcal I)$,\quad $E=\operatorname{Proj}_Z\operatorname{gr}_{\mathcal I}(\mathcal O_X)$};
 \node[b,below=of c] (d) {Mapy współrzędnych $y=xv$, $x=yu$: transformata całkowita i ścisła};
 \node[b,below=of d] (e) {Rozwiązanie osobliwości\\oraz proste skrzyżowania normalne (SNC)};
 \draw[->,thick] (a)--(b1);
 \draw[->,thick] (b1)--(c);
 \draw[->,thick] (c)--(d);
 \draw[->,thick] (d)--(e);
\end{tikzpicture}

{\small Rysunek 5.4. Logiczna zależność pojęć, a nie algorytm automatycznego rozwiązywania wszystkich osobliwości.}\par
\end{minipage}
\end{center}

\begin{przyklad}[Dwa kolejne rozdmuchania ostrza $y^2=x^5$]
Rozważmy krzywą $C=\V(y^2-x^5)\subset\A^2_{\C}$.
W mapie $y=xv$ pierwszego rozdmuchania mamy
$y^2-x^5=x^2(v^2-x^3)$. Dzielnik wyjątkowy jest
określony równaniem $x=0$, a transformata ścisła
$C_1$ ma równanie $v^2=x^3$. W szczególności
$C_1$ nadal ma ostrze w $(0,0)$: obie pochodne
cząstkowe równania $v^2-x^3$ tam znikają.

Rozdmuchajmy teraz ten punkt. W mapie $v=xw$
otrzymujemy $v^2-x^3=x^2(w^2-x)$, a zatem
transformata ścisła $C_2$ jest dana równaniem
$x=w^2$. Jest ona gładka, ponieważ pochodna
$x-w^2$ względem $x$ jest równa $1$.
Nie wynika stąd jednak, że konfiguracja całkowitego
przeciwobrazu $C$ ma proste skrzyżowania normalne.
Rzeczywiście $C_2$ styka się z nowym dzielnikiem
wyjątkowym $x=0$ w punkcie $(0,0)$, gdyż ograniczenie
$x=w^2$ daje tam przecięcie o krotności $2$.
Jeśli celem jest rozwiązanie zanurzone, należy więc
kontynuować rozdmuchania i kontrolować wszystkie
składowe dzielnika wyjątkowego, nie tylko gładkość $C_2$.
\end{przyklad}

\begin{uwaga}[Przygotowanie do geometrii powierzchni]
W następnym rozdziale rozdmuchania staną się
narzędziem, a nie głównym obiektem. Zajmiemy się
powierzchniami algebraicznymi, przekształceniami
biwymiernymi, ich dzielnikami i krzywymi wyjątkowymi.
W rozdziale o teorii przecięć wyprowadzimy
liczbowy sens $E^2=-1$ i podstawowe formuły
dla przecięć po rozdmuchaniu. Tam również rozróżnimy
równość cykli od równości samych nośników.
\end{uwaga}

% Literatura kontrolna (niewymagana do zrozumienia rachunków):
% The Stacks Project: 01OF (Blowing up), 01V4 (Smooth morphisms),
% 01W7 (Projective morphisms), 0FUT (Blow-up of smooth center).
% Wykorzystane jako odniesienia do twierdzen zewnetrznych:
% Hironaka (resolution in characteristic zero), etale-local smooth centers.


% =========================================================
% ROZDZIAL 6. POWIERZCHNIE ALGEBRAICZNE
% =========================================================

\chapter{Powierzchnie algebraiczne: przecięcia, rozdmuchania i szczególne klasy}
\label{ch:ga6-powierzchnie}

\section{Dlaczego geometria powierzchni jest jakościowo nowa?}
\label{sec:ga6-motywacja}

W przypadku gładkiej krzywej algebraicznej punkty są dywizorami,
a odwzorowanie biwymierne między gładkimi krzywymi właściwymi jest
izomorfizmem. W wymiarze dwa sytuacja zmienia się zasadniczo:
można rozdmuchać punkt gładkiej powierzchni i otrzymać inną gładką
powierzchnię. Nowa powierzchnia ma nawet nową krzywą $E\cong\PP^1$,
ale jej ciało funkcji jest takie samo. Żeby odróżnić powierzchnie,
nie wystarcza zatem samo ciało funkcji ani zbiór punktów.

W tym rozdziale głównym obiektem będzie \emph{gładka, spójna,
rzutowa powierzchnia} $S$ nad ciałem algebraicznie domkniętym
$k$ charakterystyki $0$. Oznacza to całkowity schemat gładki i
rzutowy nad $k$, o wymiarze dwa. Gdy nie korzystamy z rzutowości,
powiemy to jawnie. Przypomnijmy, że \emph{całkowity} znaczy
nierozkładalny i zredukowany. W celu porównania z topologią
przyjmujemy $k=\C$.

\begin{przyklad}[Pierwsze cztery powierzchnie]
\label{prz:ga6-examples}
\begin{enumerate}[label=(\roman*)]
\item $\PP^2_k$ jest gładką powierzchnią rzutową; proste w niej
      mają wspólny stopień jeden.
\item $\PP^1_k\times\PP^1_k$ jest gładką powierzchnią z dwiema
      różnymi rodzinami krzywych będących włóknami rzutowań.
\item Gładka powierzchnia kubiczna $S=\V_+(F)\subseteq\PP^3_k$,
      gdzie $F$ jest jednorodny stopnia trzy, zawiera dokładnie
      $27$ prostych (udowodnimy część rachunkową tego wyniku).
\item Gładka powierzchnia stopnia czwartego w $\PP^3_\C$
      jest powierzchnią K3; zobaczymy, dlaczego jej wiązka
      kanoniczna jest trywialna i dlaczego nie jest wymierna.
\end{enumerate}
\end{przyklad}

\begin{definicja}[Wymierność i powierzchnia wymierna]
Niech $X,Y$ będą całkowitymi rozmaitościami algebraicznymi nad $k$.
Są \emph{biwymiernie równoważne}, gdy zawierają izomorficzne
niepuste otwarte podzbiory; równoważnie $k(X)\cong k(Y)$ jako
rozszerzenia ciała $k$. Powierzchnia $S$ jest \emph{wymierna},
jeżeli jest biwymiernie równoważna z $\PP^2_k$.
\end{definicja}

\begin{uwaga}[Dwa rodzaje klasyfikacji]
Klasyfikacja \emph{z dokładnością do izomorfizmu} odróżnia
$\PP^2$ od jego rozdmuchania w punkcie. Klasyfikacja
\emph{biwymierna} je utożsamia. Jeżeli chcemy zachować
liczby przecięć i układ krzywych, musimy kontrolować
konkretne modele, nie tylko ich ciało funkcji.
\end{uwaga}

\section{Przecięcia krzywych: od prostych rachunków do formy przecięć}
\label{sec:ga6-intersections}

\subsection{Krotność lokalna: dlaczego samo liczenie punktów nie działa?}

Niech $S$ będzie gładką powierzchnią nad $k$ oraz $C,D\subseteq S$
dwiema \emph{różnymi} całkowitymi krzywymi. W otoczeniu danego
punktu $p\in C\cap D$ krzywe są zadane przez pojedyncze równania
lokalne $f=0$ i $g=0$. Istnienie takich równań wynika z faktu,
że regularny pierścień lokalny powierzchni jest pierścieniem
z jednoznacznym rozkładem na czynniki. To ostatnie twierdzenie
o regularnych pierścieniach lokalnych przyjmujemy na tym etapie
jako wynik algebry przemiennej.

\begin{definicja}[Lokalna krotność przecięcia]
\label{def:ga6-intersection-local}
Jeżeli $p\in C\cap D$ jest izolowanym punktem ich przecięcia,
a $f,g\in\mathcal O_{S,p}$ są lokalnymi równaniami $C,D$,
to definiujemy
\[
 i_p(C,D;S):=\operatorname{length}_{\mathcal O_{S,p}}
   \bigl(\mathcal O_{S,p}/(f,g)\bigr).
\]
Długość modułu oznacza liczbę czynników w jego szeregu
kompozycyjnym; dla $k$ algebraicznie domkniętego, przy
skończonej długości, jest równa wymiarowi tej przestrzeni
jako przestrzeni wektorowej nad $k$. Krzywe przecinają się
\emph{transwersalnie} w $p$, gdy ich przestrzenie styczne w $p$
są różnymi prostymi w $T_pS$.
\end{definicja}

\begin{stwierdzenie}[Niezależność od wyboru równań]
\label{st:ga6-independent-local}
Niech $S$ będzie gładką powierzchnią nad ciałem $k$,
a $C,D\subseteq S$ efektywnymi dzielnikami Cartiera
bez wspólnej składowej w otoczeniu punktu zamkniętego
$p\in C\cap D$. Długość
$\operatorname{length}_{\mathcal O_{S,p}}
\bigl(\mathcal O_{S,p}/(f,g)\bigr)$, gdzie $f,g$
są lokalnymi równaniami $C,D$, nie zależy od ich wyboru.
\end{stwierdzenie}
\begin{proof}
Dwa równania definiujące ten sam efektywny dzielnik
Cartiera w otoczeniu $p$ różnią się o jednostkę:
$f'=uf$ i $g'=vg$, gdzie $u,v\in\mathcal O_{S,p}^{\times}$.
Zatem $(f',g')=(f,g)$ i pierścienie ilorazowe są takie same.
Wobec tego długość nie zależy od lokalnych generatorów.
\end{proof}

\begin{stwierdzenie}[Przecięcie transwersalne ma krotność jeden]
\label{st:ga6-transversal-one}
Niech $p$ będzie punktem gładkim obu krzywych $C,D$ na
gładkiej powierzchni $S$. Wówczas $i_p(C,D;S)=1$ wtedy
i tylko wtedy, gdy krzywe przecinają się transwersalnie w $p$.
\end{stwierdzenie}
\begin{proof}
Oznaczmy $R=\mathcal O_{S,p}$ i $\mathfrak m=\mathfrak m_p$.
Ponieważ $R$ jest regularnym pierścieniem lokalnym wymiaru dwa,
przestrzeń $\mathfrak m/\mathfrak m^2$ ma wymiar dwa nad $k$.
Różniczki $df(p),dg(p)$ są niezależne wtedy i tylko wtedy,
gdy klasy $f,g$ tworzą bazę $\mathfrak m/\mathfrak m^2$.
Z lematu Nakayamy otrzymujemy wówczas $(f,g)=\mathfrak m$,
więc $R/(f,g)=R/\mathfrak m\cong k$ ma długość jeden.
Odwrotnie, jeśli długość $R/(f,g)$ jest równa jeden,
to iloraz jest prostym $R$-modułem; ponieważ $(f,g)$ jest
$\mathfrak m$-pierwotny, musi zachodzić $(f,g)=\mathfrak m$.
Klasy $f,g$ generują wtedy dwuwymiarową przestrzeń
$\mathfrak m/\mathfrak m^2$ i są liniowo niezależne.
Jest to właśnie warunek transwersalności.
\end{proof}

\begin{przyklad}[Dwie krzywe styczne]
\label{prz:ga6-tangent}
Na $\A^2_{x,y}$ rozważmy $C=\V(y)$ i $D=\V(y-x^m)$,
gdzie $m\ge1$. Obie krzywe są gładkie w początku.
Ponieważ w pierścieniu lokalnym można najpierw podstawić
$y=0$, mamy
$\mathcal O_{\A^2,0}/(y,y-x^m)\cong k[x]_{(x)}/(x^m)$.
Jego bazę nad $k$ stanowią klasy $1,x,\ldots,x^{m-1}$;
więc $i_0(C,D;\A^2)=m$. Dla $m=1$ przecięcie jest
transwersalne, dla $m=2$ jest styczne pierwszego rzędu,
a dowolnie wysokie $m$ daje dowolnie dużą krotność mimo
obecności tylko jednego punktu przecięcia.
\end{przyklad}

\begin{center}
\begin{minipage}{\linewidth}
\centering
\begin{tikzpicture}[scale=0.88,>=Stealth,every node/.style={font=\small}]
  \begin{scope}[shift={(-3.2,0)}]
   \draw[gray!65,->] (-1.45,0)--(1.45,0) node[below right] {$x$};
   \draw[gray!65,->] (0,-0.9)--(0,1.45) node[left] {$y$};
   \draw[blue!70!black,thick] (-1.05,-0.7)--(1.05,0.7);
   \draw[red!70!black,thick] (-1.05,0.7)--(1.05,-0.7);
   \node[align=center,text width=4.3cm] at (0,-1.65)
     {Przecięcie transwersalne\\$i_0=1$};
  \end{scope}
  \begin{scope}[shift={(3.2,0)}]
   \draw[gray!65,->] (-1.45,0)--(1.45,0) node[below right] {$x$};
   \draw[gray!65,->] (0,-0.9)--(0,1.45) node[left] {$y$};
   \draw[red!70!black,thick] (-1.3,0)--(1.3,0);
   \draw[blue!70!black,thick,domain=-1.1:1.1,variable=\t]
    plot ({\t},{0.8*\t*\t});
   \node[align=center,text width=4.3cm] at (0,-1.65)
      {Przecięcie styczne\\$i_0=2$};
  \end{scope}
\end{tikzpicture}

{\small Rysunek 6.1. Jeden punkt przecięcia może mieć różną krotność algebraiczną. W lewym panelu różniczki równań są niezależne; w prawym krzywe mają wspólną styczną.}\par
\end{minipage}
\end{center}

\subsection{Globalne liczby przecięć i samoprzecięcie}

\begin{definicja}[Dywizory i ich równoważność liniowa]
\label{def:ga6-divisor-group}
Niech $S$ będzie gładką całkowitą powierzchnią nad ciałem $k$.
Grupa $\operatorname{Div}(S)$ jest wolną grupą abelową na
całkowitych domkniętych krzywych $C\subseteq S$. Jej elementy
$D=\sum_C n_C C$ są \emph{dywizorami Weila}; suma ma skończony
nośnik. Dla niezerowej funkcji wymiernej $f\in k(S)^\times$
definiujemy $\operatorname{div}(f)=\sum_C\operatorname{ord}_C(f)C$.
Rząd $\operatorname{ord}_C$ jest wartościowaniem w pierścieniu lokalnym
$\mathcal O_{S,\eta_C}$ punktu generycznego krzywej $C$;
jest to dyskretny pierścień wartościowania, ponieważ $S$ jest
regularna. Dwa dywizory $D,D'$ są \emph{liniowo równoważne},
co oznaczamy $D\sim D'$, gdy $D-D'=\operatorname{div}(f)$
dla pewnego $f\in k(S)^\times$.
\end{definicja}

\begin{twierdzenie}[Forma przecięć na powierzchni rzutowej]
\label{tw:ga6-intersection-global}
Niech $S$ będzie gładką rzutową powierzchnią nad $k$.
Istnieje jednoznaczna symetryczna forma dwuliniowa
$\operatorname{Div}(S)\times\operatorname{Div}(S)\to\Z$,
$(D_1,D_2)\mapsto D_1\cdot D_2$, zależna tylko od klas
równoważności liniowej, która spełnia:
\begin{enumerate}[label=(\roman*)]
\item dla różnych krzywych całkowitych $C,D$ zachodzi
$C\cdot D=\sum_{p\in C\cap D}i_p(C,D;S)$;
\item dla gładkiej krzywej całkowitej $C\subseteq S$
zachodzi $C^2:=C\cdot C=\deg\bigl(\mathcal O_S(C)|_C\bigr)$;
\item dla dowolnej krzywej całkowitej $C$ i dywizora $D$
zachodzi $D\cdot C=\deg(\mathcal O_S(D)|_{\widetilde C})$,
gdzie $\widetilde C\to C$ jest normalizacją.
\end{enumerate}
W (iii) stopień wiązki na normalizacji liczymy z jej
pullbacku. Na gładkiej krzywej $C$ wzór upraszcza się
do stopnia na samej $C$.
\end{twierdzenie}

\begin{uwaga}[Granica pełnego dowodu]
Nie będziemy udawać, że globalna postać formy przecięć
wynika wyłącznie z dwóch przykładów lokalnych. Jej konstrukcja
wymaga m.in. zgodności lokalnej długości z teorią klas dywizorów
i stopni wiązek oraz lematu o przesuwaniu cykli.
\emph{Twierdzenie~\ref{tw:ga6-intersection-global} przyjmujemy tutaj
jako podstawowy wynik teorii przecięć}. Jego lokalne części
są wyliczone bezpośrednio, natomiast kluczowe zastosowania
do rozdmuchań poniżej zostaną uzasadnione rachunkiem wiązek.
W rozdziale 7 wyprowadzimy formę przecięć przez
pierwszą klasę Cherna i stopnie na krzywych;
sam ogólny lemat o przesuwaniu cykli zostanie tam jawnie
odnotowany jako osobne twierdzenie.
\end{uwaga}

\begin{przyklad}[Przecięcia w $\PP^2$]
\label{prz:ga6-p2}
Niech $H\subseteq\PP^2$ będzie prostą. Dwie różne proste
przecinają się w jednym punkcie transwersalnie,
więc $H^2=1$ (zależność od klasy liniowej pozwala zastąpić
prostą inną prostą). Krzywa płaska zadana równaniem jednorodnym
stopnia $d$ ma klasę liniową $dH$: jest to dywizor zer
sekcji $\mathcal O_{\PP^2}(d)$. Z dwuliniowości wynika
$(dH)\cdot(eH)=de$, czyli wersja twierdzenia Bézouta
z krotnościami, jeśli dwie krzywe nie mają wspólnej składowej.
Dla krzywej stopnia $d$ jej \emph{samoprzecięcie} wynosi $d^2$,
choć geometrycznie nie przecinamy krzywej z samą sobą
w sensie zbiorowym.
\end{przyklad}

\begin{definicja}[Wiązka normalna krzywej]
Jeżeli $C\subseteq S$ jest gładką całkowitą krzywą
na gładkiej powierzchni, to jej \emph{wiązka normalna}
jest wiązką liniową
$N_{C/S}:=(\mathcal I_C/\mathcal I_C^2)^\vee$ na $C$.
Z identyfikacji snopa ideałów $\mathcal I_C\cong\mathcal O_S(-C)$
otrzymujemy $N_{C/S}\cong\mathcal O_S(C)|_C$, a zatem
$C^2=\deg N_{C/S}$. To nadaje samoprzecięciu sens
geometryczny: mierzy stopień normalnego kierunku do krzywej.
\end{definicja}

\section{Rozdmuchanie punktu: postać przecięć i klasa kanoniczna}
\label{sec:ga6-blowup}

Niech teraz $\pi:\widetilde S\to S$ będzie rozdmuchaniem
gładkiej rzutowej powierzchni w punkcie $p\in S$,
a $E=\pi^{-1}(p)\cong\PP^1$ jego dzielnikiem wyjątkowym.

\input{../wspolne/rozdmuchanie_punktu_normalna.tex}

\begin{stwierdzenie}[Ortogonalność i formuła przecięć po rozdmuchaniu]
\label{st:ga6-blowup-form}
Niech $D,D'$ będą dywizorami na $S$. Wtedy
\[
 (\pi^*D)\cdot E=0,\qquad
 (\pi^*D)\cdot(\pi^*D')=D\cdot D',\qquad E^2=-1.
\]
Jeżeli całkowita krzywa $C\subseteq S$ przechodzi
przez $p$ z krotnością $m$, to jej transformata ścisła
$\widetilde C$ spełnia
$\pi^*C=\widetilde C+mE$ jako dywizory, a zatem
\[
 \widetilde C^2=C^2-m^2,\qquad
 \widetilde C\cdot E=m.
\]
\end{stwierdzenie}
\begin{proof}
Ponieważ $\pi|_E:E\to\{p\}$ jest odwzorowaniem stałym,
$\mathcal O_{\widetilde S}(\pi^*D)|_E$
jest trywialną wiązką $\mathcal O_E$;
jej stopień wynosi zero. Z definicji formy przecięć
otrzymujemy $(\pi^*D)\cdot E=0$.
Druga równość jest szczególną postacią \emph{wzoru projekcji}
dla przecięć: $\pi_*(\pi^*D\cdot\pi^*D')=D\cdot D'$.
W przypadku $\PP^2$ można ją też zobaczyć bez ogólnego
wzoru projekcji: przesuwamy przedstawicieli klas
prostych poza $p$, gdzie $\pi$ jest izomorfizmem.
Ogólną wersję wzoru projekcji traktujemy jako część
przyjętej teorii przecięć.

Dla ostatnich równości weźmy w $p$ lokalne równanie
$f$ krzywej $C$. Jeśli jego najmniejsza część
jednorodna ma stopień $m$, to na charcie $y=xv$
znany z rozdziału 5 wzór daje
$f(x,xv)=x^mF(x,v)$. Wyrażenie $x^m$ oznacza $mE$,
a $F=0$ transformatę ścisłą. Analogiczny wzór na drugim
charcie skleja się do równości dywizorów
$\pi^*C=\widetilde C+mE$. Stąd, korzystając z
$\pi^*C\cdot E=0$ oraz $E^2=-1$, mamy
$\widetilde C\cdot E=m$. Wreszcie
$\widetilde C^2=(\pi^*C-mE)^2=C^2-m^2$.
\end{proof}

\begin{przyklad}[Zerowanie i zmiana samoprzecięcia]
\label{prz:ga6-blowup-curves}
W $S=\PP^2$ prosta ma klasę $H$ i $H^2=1$.
Po rozdmuchaniu punktu $p$ należącego do prostej
jej transformata ścisła ma klasę $H-E$ i
$(H-E)^2=1-1=0$. Po rozdmuchaniu \emph{dwóch różnych}
punktów tej samej prostej otrzymujemy klasę
$H-E_1-E_2$ o samoprzecięciu $-1$.
Jeżeli rozdmuchiwany punkt nie leży na prostej,
transformata ma klasę $H$ i samoprzecięcie nadal $1$.
Ta sama krzywa jako abstrakcyjna $\PP^1$ może więc
mieć w różnych powierzchniach wiązki normalne
różnych stopni. Samoprzecięcie nie jest niezmiennikiem
samej krzywej, lecz jej zanurzenia.
\end{przyklad}

\begin{stwierdzenie}[Klasa kanoniczna po rozdmuchaniu]
\label{st:ga6-Kblowup}
Niech $K_S$ oznacza dowolny dywizor kanoniczny
gładkiej powierzchni $S$. Wtedy
\[
 K_{\widetilde S}\sim\pi^*K_S+E,
 \qquad K_{\widetilde S}^2=K_S^2-1.
\]
Pierwsza równość jest równością klas liniowej równoważności.
\end{stwierdzenie}
\begin{proof}
W lokalnych współrzędnych $x,y$ na $S$ generator
wiązki kanonicznej jest reprezentowany przez
$dx\wedge dy$. Na charcie rozdmuchania $y=xv$
mamy $dy=v\,dx+x\,dv$, zatem
$\pi^*(dx\wedge dy)=x\,dx\wedge dv$.
Przy przeciąganiu lokalnej różniczki dwuwymiarowej
pojawia się dokładnie jeden czynnik $x$ — równanie
$E$ na tym charcie. Na drugim charcie analogicznie
$dx\wedge dy=-(y)\,du\wedge dy$ po cofnięciu,
więc również współczynnik zanika na $E$ rzędu jeden.
W ten sposób dywizor kanoniczny na $\widetilde S$
równa się klasie $\pi^*K_S+E$.
Równość kwadratów wynika z poprzedniego stwierdzenia:
$(\pi^*K_S+E)^2=K_S^2+E^2=K_S^2-1$.
\end{proof}

\begin{stwierdzenie}[Grupa Picarda po rozdmuchaniu]
\label{st:ga6-pic-blowup}
Dla $\pi:\widetilde S=\operatorname{Bl}_pS\to S$,
gdzie $S$ jest gładką rzutową powierzchnią, zachodzi
\[
 \operatorname{Pic}(\widetilde S)\cong
 \pi^*\operatorname{Pic}(S)\oplus\Z[E].
\]
W szczególności dla rozdmuchania $\PP^2$ w $r$
różnych punktach otrzymujemy wolną grupę abelową
z bazą $H,E_1,\ldots,E_r$; w tej bazie formę
przecięć opisuje diagonalna macierz
$\operatorname{diag}(1,-1,\ldots,-1)$.
\end{stwierdzenie}
\begin{proof}
Na gładkiej powierzchni każdy dzielnik Weila jest
Cartiera: wykorzystujemy standardowe twierdzenie,
że regularne pierścienie lokalne są UFD. Niech
$\widetilde D$ będzie dzielnikiem na $\widetilde S$.
Każdą jego nierozkładalną składową różną od $E$
można zrzutować na krzywą $D\subseteq S$; jej
transformata ścisła różni się od $\pi^*D$ o
całkowitą wielokrotność $E$. Po zsumowaniu składowych
otrzymujemy $\widetilde D\sim\pi^*D+aE$.
To dowodzi surjektywności odpowiedniego homomorfizmu.

Jeżeli $\pi^*D+aE$ jest liniowo równoważny zeru,
to po ograniczeniu wiązki $\mathcal O_{\widetilde S}(\pi^*D+aE)$
do $E$ otrzymujemy $\mathcal O_{\PP^1}(-a)$.
Musi ona być trywialna, więc $a=0$. Wtedy
$\pi^*D\sim0$. Zmniejszając powierzchnie do
dopełnień punktu $p$ i krzywej $E$, na których
$\pi$ jest izomorfizmem, otrzymujemy $D|_{S\setminus\{p\}}\sim0$.
Na gładkiej powierzchni usunięcie punktu kowymiaru
dwa nie zmienia klas dywizorów: przedstawiciele
Weila mają nośniki kowymiaru jeden, a funkcje
wymierne są te same. Wobec tego $D\sim0$.
Uzyskaliśmy również injektywność i rozkład bezpośredni.
Dla $S=\PP^2$ jej grupa Picarda jest generowana
przez $H=\mathcal O_{\PP^2}(1)$; kolejne rozdmuchania
powtarzają dowód i dają podaną bazę.
\end{proof}

\begin{uwaga}[Kiedy klasa $E_1$ nie oznacza już tej samej krzywej?]
\label{uw:ga6-infinitely-near}
Jeżeli po rozdmuchaniu $p$ rozdmuchujemy następnie punkt
$q$ \emph{na pierwszym dzielniku wyjątkowym}, to
$E_1$ w bazie grupy Picarda oznacza \emph{całkowity
przeciwobraz} poprzedniego dzielnika, nie jego aktualną
transformatę ścisłą. Ta ostatnia ma klasę $E_1-E_2$,
a więc samoprzecięcie $(E_1-E_2)^2=-2$.
Ponadto $(E_1-E_2)\cdot E_2=1$. Otrzymujemy dwie
gładkie krzywe przecinające się transwersalnie,
z samoprzecięciami $-2$ i $-1$. To właśnie typ
konfiguracji, który pojawia się przy rozwiązywaniu
osobliwości.
\end{uwaga}

\begin{center}
\begin{minipage}{\linewidth}
\centering
\begin{tikzpicture}[>=Stealth,every node/.style={font=\small}]
 \node[align=center,draw,rounded corners,inner sep=8pt,text width=4.1cm] (a) at (0,0)
 {Rozdmuchanie pierwszego punktu\\$E_1^2=-1$};
 \node[align=center,draw,rounded corners,inner sep=8pt,text width=4.6cm] (b) at (7.5,0)
 {Drugie rozdmuchanie na $E_1$\\$(E_1-E_2)^2=-2$, $E_2^2=-1$};
 \draw[->,thick] (a.east)--node[above] {punkt na $E_1$}(b.west);
 \draw[red!70!black,very thick] (-0.65,-1.55)--(0.65,-1.55);
 \node[below] at (0,-1.65) {$E_1$};
 \draw[red!70!black,very thick] (6.7,-1.55)--(7.85,-1.55);
 \draw[blue!65!black,very thick] (7.85,-2.05)--(7.85,-1.05);
 \fill (7.85,-1.55) circle (2pt);
 \node[below] at (6.95,-1.65) {$E_1'$};
 \node[right] at (7.95,-1.2) {$E_2$};
\end{tikzpicture}

{\small Rysunek 6.2. Po drugim rozdmuchaniu zmienia się samoprzecięcie transformaty ścisłej starego dzielnika. Nie należy utożsamiać jej klasy z klasą całkowitego przeciwobrazu.}\par
\end{minipage}
\end{center}

\section{Wzór adjunkcji i krzywe wyjątkowe}
\label{sec:ga6-adjunction}

\begin{twierdzenie}[Wzór adjunkcji dla gładkiej krzywej na powierzchni]
\label{tw:ga6-adjunction}
Niech $S$ będzie gładką rzutową powierzchnią, a
$C\subseteq S$ gładką, spójną krzywą. Wówczas
$K_C\sim (K_S+C)|_C$, a zatem
\[
 2g(C)-2=(K_S+C)\cdot C=K_S\cdot C+C^2.
\]
\end{twierdzenie}
\begin{proof}
Ponieważ $C$ jest gładkim efektywnym dzielnikiem
Cartiera, otrzymujemy dokładny ciąg snopów na $C$
\[
0\longrightarrow\mathcal I_C/\mathcal I_C^2
 \longrightarrow\Omega^1_S|_C\longrightarrow\Omega^1_C
 \longrightarrow0.
\]
Pierwszy wyraz jest wiązką liniową
$\mathcal O_C(-C)$; drugi ma rząd dwa, a trzeci
rząd jeden. Biorąc wyznaczniki tego dokładnego
ciągu wiązek lokalnie wolnych, otrzymujemy
$\det(\Omega^1_S|_C)\cong
 \mathcal O_C(-C)\otimes\Omega^1_C$.
Lewa strona jest $\mathcal O_C(K_S)$,
a więc po przeniesieniu czynnika
$\mathcal O_C(-C)$ otrzymujemy
$\Omega_C^1\cong\mathcal O_C(K_S+C)$.
To daje pierwszą równość. Druga wynika przez
wzięcie stopni oraz ze znanego z twierdzenia
Riemanna--Roch faktu $\deg K_C=2g(C)-2$.
\end{proof}

\begin{wniosek}[Samoprzecięcie wymiernej krzywej]
\label{wn:ga6-rational-adjunction}
Jeżeli $C\cong\PP^1$ jest gładką krzywą na gładkiej
rzutowej powierzchni $S$, to
$C^2=-2-K_S\cdot C$.
W szczególności jeżeli $K_S\cdot C=-1$, to $C^2=-1$.
\end{wniosek}
\begin{proof}
Wystarczy wstawić $g(C)=0$ do wzoru adjunkcji.
\end{proof}

\begin{definicja}[Krzywa wyjątkowa pierwszego rodzaju]
Gładką krzywą $E\subseteq S$, izomorficzną z
$\PP^1$ i spełniającą $E^2=-1$, nazywamy
\emph{krzywą $(-1)$}, czyli krzywą wyjątkową
pierwszego rodzaju. Rozdmuchanie punktu na gładkiej
powierzchni wytwarza taką krzywą.
\end{definicja}

\begin{twierdzenie}[Castelnuovo: kryterium zdmuchnięcia]
\label{tw:ga6-castelnuovo}
Niech $S$ będzie gładką rzutową powierzchnią nad
algebraicznie domkniętym ciałem charakterystyki zero.
Jeśli $E\subseteq S$ jest krzywą $(-1)$, to istnieje
gładka powierzchnia $S'$ i właściwy morfizm
$\rho:S\to S'$ kontraktujący dokładnie $E$ do jednego
gładkiego punktu i będący izomorfizmem poza $E$.
Morfizm $\rho$ jest odwrotnością rozdmuchania tego punktu.
\end{twierdzenie}

\begin{uwaga}[Status dowodu Castelnuova]
Nie wyprowadzamy tego twierdzenia z samego warunku
$E^2=-1$: wymaga ono argumentu o deformacjach,
wiązce normalnej i budowie lokalnego otoczenia krzywej.
Widzieliśmy natomiast \emph{model odwrotny}: wyjątkowy
dzielnik rozdmuchania gładkiego punktu ma właśnie
wiązkę normalną $\mathcal O(-1)$. Twierdzenie
Castelnuova mówi, że w kategorii gładkich powierzchni
rzutowych jest to jedyny możliwy lokalny model
gładkiego zdmuchnięcia. Pozostaje jawnie przyjętym
wynikiem zewnętrznym.
\end{uwaga}

\begin{przyklad}[Krzywa $(-2)$ nie może zostać zdmuchnięta do gładkiego punktu]
Rozważmy dwa kolejne rozdmuchania z
uwagi~\ref{uw:ga6-infinitely-near}. Transformata
ścisła pierwszego dzielnika ma klasę $E_1-E_2$
i kwadrat $-2$. Nie spełnia więc warunku
Castelnuova, choć nadal jest krzywą $\PP^1$.
Może być kontraktowana w innym kontekście do
punktu osobliwego, lecz nie jest obrazem wyjątkowej
krzywej jednego rozdmuchania \emph{gładkiego} punktu.
\end{przyklad}

\section{Powierzchnie wymierne i powierzchnie Hirzebrucha}
\label{sec:ga6-rational}

\begin{przyklad}[Dwie rodziny prostych na $\PP^1\times\PP^1$]
Niech $S=\PP^1\times\PP^1$ i niech
$F_1=\{p\}\times\PP^1$, $F_2=\PP^1\times\{q\}$.
Dwa różne włókna tego samego rzutowania są rozłączne,
więc $F_1^2=F_2^2=0$ przez równoważność liniową.
Natomiast $F_1$ i $F_2$ przecinają się dokładnie
w jednym punkcie transwersalnie, zatem
$F_1\cdot F_2=1$. Dla dywizora $aF_1+bF_2$
mamy więc $(aF_1+bF_2)^2=2ab$.
Kanoniczny dywizor jest równy
$K_S\sim-2F_1-2F_2$, ponieważ
$K_{\PP^1\times\PP^1}$ jest tensorowym iloczynem
pullbacków wiązek $K_{\PP^1}=\mathcal O(-2)$.
Stąd $K_S^2=8$.
\end{przyklad}

\begin{definicja}[Powierzchnia Hirzebrucha]
Dla liczby całkowitej $n\ge0$ \emph{powierzchnią
Hirzebrucha} $\mathbb F_n$ nazywamy
$\PP_{\PP^1}\bigl(\mathcal O_{\PP^1}\oplus
\mathcal O_{\PP^1}(n)\bigr)$, z konwencją, w której
$\PP$ parametryzuje jednowymiarowe \emph{ilorazy}
włókien wiązki wektorowej. Ma ona naturalne rzutowanie
$f:\mathbb F_n\to\PP^1$, którego włókna są prostymi
rzutowymi. Oznaczymy przez $F$ klasę włókna,
a przez $B$ sekcję odpowiadającą ilorazowi
$\mathcal O\oplus\mathcal O(n)\twoheadrightarrow\mathcal O$.
\end{definicja}

\begin{stwierdzenie}[Przecięcia na powierzchni Hirzebrucha]
\label{st:ga6-hirzebruch}
W $\mathbb F_n$ zachodzą równości
\[
 B^2=-n,\qquad B\cdot F=1,\qquad F^2=0,
 \qquad K_{\mathbb F_n}\sim-2B-(n+2)F.
\]
W szczególności $K_{\mathbb F_n}^2=8$ dla każdego $n\ge0$.
\end{stwierdzenie}
\begin{proof}
Rzutowanie $f$ ma rozłączne włókna liniowo równoważne,
więc $F^2=0$. Sekcja przecina każde włókno w jednym
punkcie transwersalnie, stąd $B\cdot F=1$.
Wiązka normalna sekcji określonej przez iloraz
$\mathcal O\oplus\mathcal O(n)\twoheadrightarrow\mathcal O$
jest homomorfizmami z jego jądra $\mathcal O(n)$
do ilorazu $\mathcal O$, a więc
$N_{B/\mathbb F_n}\cong\operatorname{Hom}
(\mathcal O(n),\mathcal O)\cong\mathcal O(-n)$.
Zatem $B^2=\deg N_{B/\mathbb F_n}=-n$.

Grupa Picarda $\mathbb F_n$ jest generowana przez
$B,F$ (to standardowa własność projektywizacji
wiązki; można ją zobaczyć z lokalnych trywializacji
i funkcji przejścia). Zapiszmy
$K_{\mathbb F_n}\sim aB+bF$. Adjunkcja na włóknie
$F\cong\PP^1$ daje $K\cdot F=-2-F^2=-2$,
czyli $a=-2$. Adjunkcja na sekcji $B\cong\PP^1$
daje $K\cdot B=-2-B^2=n-2$.
Zatem $-2(-n)+b=n-2$, skąd $b=-n-2$.
Podstawienie do formy przecięć daje
$K^2=(-2B-(n+2)F)^2=4(-n)+4(n+2)=8$.
\end{proof}

\begin{przyklad}[Dlaczego $\mathbb F_1$ jest rozdmuchaniem płaszczyzny?]
\label{prz:ga6-f1}
Ustalmy $p=[0:0:1]\in\PP^2$ i współrzędne $[x:y:z]$.
Rozdmuchanie można zapisać jako powierzchnię incydencji
\[
B=\{([x:y:z],[u:v])\in\PP^2\times\PP^1:\ xv=yu\}.
\]
Projekcja na $\PP^2$ jest izomorfizmem poza $p$, natomiast
jej włókno nad $p$ jest całym $\PP^1$.
Druga projekcja $q:B\to\PP^1$ ma włókna $\PP^1$.
Rzeczywiście, nad chartem $u\ne0$ piszemy $t=v/u$ i
\[
([\lambda:\lambda t:\mu],[1:t]),\qquad[\lambda:\mu]\in\PP^1.
\]
Nad chartem $v\ne0$ piszemy $s=u/v$ i otrzymujemy
$([s\lambda':\lambda':\mu'],[s:1])$.
Na przecięciu $s=1/t$ współrzędne włókien spełniają
$[\lambda':\mu']=[t\lambda:\mu]$. Ten wzór jest
projektywizacją przejścia wiązki
$\mathcal O_{\PP^1}(-1)\oplus\mathcal O_{\PP^1}$
w konwencji \emph{prostych} we włóknach: w lokalnych
ramkach $e_u=(1,t)$ i $e_v=(s,1)$ wiązki tautologicznej
zachodzi $e_u=t e_v$.
Dualizacja przekształca konwencję prostych w przyjętą
w definicji $\mathbb F_n$ konwencję jednowymiarowych ilorazów.
Dlatego
$B\cong\PP_{\PP^1}(\mathcal O\oplus\mathcal O(1))=\mathbb F_1$.
Wyjątkowa krzywa $E$ jest sekcją $q$, zadaną warunkiem
$\lambda=0$ ma $E^2=-1$; to jednocześnie
bezpośrednia identyfikacja powierzchni i jej sekcji,
bez odwoływania się do niewyjaśnionej jednoznaczności.
\end{przyklad}

\begin{uwaga}[Powierzchnie minimalne]
Gładką rzutową powierzchnię nazywamy \emph{minimalną},
jeżeli nie zawiera krzywych $(-1)$. Z twierdzenia
Castelnuova wynika, że każdą gładką rzutową powierzchnię
można, dopóki ma krzywe $(-1)$, próbować upraszczać
przez ich zdmuchiwanie; zakończenie tego procesu
uzasadnia np. zmiana liczby Picarda. Powierzchnia
$\mathbb F_1$ nie jest minimalna, natomiast
$\PP^2$ oraz $\mathbb F_n$ dla $n=0$ i $n\ge2$
są minimalne. Dla wymiernych powierzchni model minimalny
nie zawsze jest jednoznaczny, co odróżnia je od
odpowiedniego zagadnienia dla powierzchni o nieujemnym
wymiarze Kodairy.
\end{uwaga}

\section{Powierzchnie del Pezzo i antykanoniczne odwzorowania}
\label{sec:ga6-delpezzo}

\begin{definicja}[Powierzchnia del Pezzo i jej stopień]
Gładka rzutowa powierzchnia $S$ jest \emph{powierzchnią
del Pezzo}, jeśli $-K_S$ jest dywizorem obszernym
(ang. \emph{ample}). Przypomnijmy: wiązka liniowa $L$
jest obszerna, gdy dla pewnego $m>0$ liniowy układ
sekcji $L^{\otimes m}$ wyznacza zanurzenie zamknięte
$S\hookrightarrow\PP^N$. Jej \emph{stopień del Pezzo}
to liczba $d=K_S^2$.
\end{definicja}

\begin{przyklad}[Płaszczyzna i powierzchnia kwadrykowa]
Dla $\PP^2$ mamy $K_{\PP^2}\sim-3H$,
więc $(-K_{\PP^2})^2=9$ i $-K_{\PP^2}$ jest
obszerny. Zatem $\PP^2$ jest powierzchnią del Pezzo
stopnia dziewięć. Dla $\PP^1\times\PP^1$
mamy $-K=2F_1+2F_2$, co jest obszerne;
stopień wynosi osiem. Obie są powierzchniami
del Pezzo, mimo że nie są izomorficzne.
\end{przyklad}

\begin{stwierdzenie}[Stopień powierzchni otrzymanej z $\PP^2$]
\label{st:ga6-del-pezzo-degree}
Niech $\pi:S_r\to\PP^2$ będzie złożeniem
$r$ rozdmuchań w \emph{różnych} punktach
$p_1,\ldots,p_r$ płaszczyzny. Wówczas
\[
 K_{S_r}\sim-3H+E_1+\cdots+E_r,
 \qquad K_{S_r}^2=9-r.
\]
W sytuacji $r=6$ i punktów w położeniu ogólnym
(brak trzech współliniowych i brak stożkowej
przechodzącej przez wszystkie sześć punktów)
dywizor $-K_{S_6}=3H-\sum E_i$ jest nawet
\emph{bardzo obszerny}: jego układ zupełny
wyznacza zanurzenie $S_6\hookrightarrow\PP^3$
jako gładką powierzchnię kubiczną.
\end{stwierdzenie}
\begin{proof}
Formułę na $K_{S_r}$ otrzymujemy, stosując kolejno
stwierdzenie~\ref{st:ga6-Kblowup}. Ponieważ
$H^2=1$, $H\cdot E_i=0$, $E_i^2=-1$
oraz $E_i\cdot E_j=0$ dla $i\ne j$, dostajemy
$K^2=9-r$. Założenie o różnych punktach zapewnia,
że dzielniki wyjątkowe nie przecinają się.

Dla ostatniego zdania potrzeba osobnego twierdzenia
o \emph{bardzo obszerności} układu kubik
płaskich przechodzących przez sześć punktów
w położeniu ogólnym. Nie ukrywamy go pod formułą
$K^2=3$. Można jednak od razu wyjaśnić wymiar
układu: przestrzeń wszystkich jednorodnych form
stopnia trzy w trzech zmiennych ma wymiar
$\binom{3+2}{2}=10$. Sześć punktów nakłada
sześć niezależnych warunków liniowych: dla każdego
$p_i$ istnieje stożkowa przechodząca przez pięć
pozostałych, lecz nie przez $p_i$, a jej iloczyn
z prostą nieprzechodzącą przez $p_i$ rozdziela
wartość w $p_i$ od wartości w pozostałych punktach.
Otrzymujemy zatem czterowymiarową przestrzeń
sekcji $H^0(S_6,\mathcal O_{S_6}(-K))$.
Twierdzenie o bardzo obszerności tej konkretnej
wiązki zapewnia, że odpowiadający morfizm do
$\PP^3$ jest zanurzeniem, a nie tylko odwzorowaniem
wymiernym. Obraz ma stopień
$(-K)^2=3$, a więc jest powierzchnią kubiczną.
\end{proof}

\begin{uwaga}[Jaki wynik przyjęliśmy?]
Twierdzenie o bardzo obszerności $-K$ dla
sześciu punktów w położeniu ogólnym można udowodnić,
sprawdzając, że kubiki przez te punkty rozdzielają
punkty i kierunki styczne na rozdmuchaniu.
Pełny dowód wymagałby systematycznych narzędzi
do badania liniowych układów krzywych, więc
\emph{tę jedną część} przyjęliśmy jako wynik
zewnętrzny. Wszystkie dalsze obliczenia klas,
przecięć i liczby prostych wykonamy jawnie.
\end{uwaga}

\section{Dokładnie 27 prostych na gładkiej powierzchni kubicznej}
\label{sec:ga6-27}

\begin{twierdzenie}[Klasyczny wynik o 27 prostych]
\label{tw:ga6-27-lines}
Niech $k$ będzie algebraicznie domkniętym ciałem
charakterystyki zero, a
$S=\V_+(F)\subseteq\PP^3_k$ gładką powierzchnią
kubiczną, tzn. $F$ jest jednorodnym wielomianem
stopnia trzy. Wtedy $S$ zawiera dokładnie
$27$ prostych projektowych.
\end{twierdzenie}

Wyjaśnimy najpierw jedyny głęboki składnik dowodu.
Klasyczne twierdzenie o klasyfikacji powierzchni
del Pezzo stopnia trzy mówi, że \emph{każda} gładka
powierzchnia kubiczna nad $k$ powstaje jako
rozdmuchanie $\PP^2_k$ w sześciu punktach w położeniu
ogólnym, a jej zanurzenie kubiczne jest określone
pełnym układem antykanonicznym. Wynik ten przyjmujemy;
nie wynika on jedynie z istnienia $27$ krzywych
na pewnym konkretnym przykładzie.

\begin{proof}[Dowód twierdzenia~\ref{tw:ga6-27-lines}, przy użyciu powyższej klasyfikacji]
Niech $\pi:S\to\PP^2$ będzie rozdmuchaniem
sześciu punktów $p_1,\ldots,p_6$. Oznaczmy
przez $H$ pullback prostej, a przez $E_i$
krzywe wyjątkowe. Rozważmy antykanoniczny
dywizor $L=-K_S=3H-\sum_{i=1}^6E_i$.
Ponieważ $L$ definiuje zanurzenie $S\hookrightarrow\PP^3$,
krzywa całkowita $C\subseteq S$ jest prostą
w $\PP^3$ wtedy i tylko wtedy, gdy stopień
$\mathcal O_S(L)|_C$, czyli $L\cdot C$, wynosi jeden
(w kierunku odwrotnym korzystamy z faktu, że
nierozkładalna krzywa projektowa stopnia jeden
jest prostą).

\emph{Krok 1: każda prosta ma samoprzecięcie $-1$.}
Jeśli $C$ jest prostą, to $C\cong\PP^1$,
$(-K_S)\cdot C=1$. Ze wzoru adjunkcji
$-2=C^2+K_S\cdot C=C^2-1$, więc $C^2=-1$.

\emph{Krok 2: sześć krzywych wyjątkowych.}
Dla każdej $E_i$ mamy $E_i^2=-1$ i
$L\cdot E_i=1$. Ich obrazy są prostymi;
daje to sześć różnych prostych.

\emph{Krok 3: kandydaci niebędący krzywymi wyjątkowymi.}
Jeśli $C$ nie jest jedną z $E_i$, jej obraz przez
$\pi$ jest nierozkładalną krzywą płaską stopnia
$a\ge1$. Niech $b_i\ge0$ będą krotnościami
tej krzywej w punktach $p_i$. Wtedy
$[C]=aH-\sum_{i=1}^6b_iE_i$. Z warunków
$L\cdot C=1$ oraz $C^2=-1$ otrzymujemy układ
\[
 3a-\sum_{i=1}^6b_i=1,\qquad
 a^2-\sum_{i=1}^6b_i^2=-1.
\]
Czyli $\sum b_i=3a-1$ oraz
$\sum b_i^2=a^2+1$. Z nierówności Cauchy'ego--Schwarza
\[
 (3a-1)^2=\left(\sum b_i\right)^2
 \le6\sum b_i^2=6(a^2+1)
\]
wynika $3a^2-6a-5\le0$. Ponieważ $a\ge1$
jest liczbą całkowitą, musi zachodzić $a=1$ albo $a=2$.

\emph{Krok 4: przypadek $a=1$.}
Mamy $\sum b_i=2$ oraz $\sum b_i^2=2$.
Dla $b_i\in\Z_{\ge0}$ równość
$\sum b_i(b_i-1)=0$ wymusza, aby każde
$b_i$ było zerem albo jedynką. Dokładnie dwie
liczby $b_i$ są równe jeden. Jest zatem
$\binom{6}{2}=15$ możliwych klas
$H-E_i-E_j$, $i<j$. Każdą z nich reprezentuje
transformata ścisła jedynej prostej płaskiej
przechodzącej przez $p_i,p_j$. Założenie
ogólnego położenia wyklucza trzecie $p_l$
na tej prostej. Wszystkie $15$ są różne.

\emph{Krok 5: przypadek $a=2$.}
Teraz $\sum b_i=5$ i $\sum b_i^2=5$,
więc znowu pięć liczb $b_i$ wynosi jeden,
a szósta wynosi zero. Otrzymujemy sześć klas
$2H-\sum_{j\ne i}E_j$, $1\le i\le6$.
Przez pięć punktów w położeniu ogólnym przechodzi
dokładnie jedna stożkowa. Jest ona gładka:
stożkowa rozkładalna byłaby sumą dwóch prostych,
z których przynajmniej jedna zawierałaby trzy
z pięciu punktów. Nie przechodzi przez szósty
punkt na mocy ogólnego położenia. Jej transformata
ścisła ma podaną klasę i jest prostą w kubicznym
zanurzeniu.

\emph{Krok 6: kompletność.}
Każda prosta należy do jednej z kategorii z
kroków 2--5; każda podana krzywa istnieje,
ma przecięcie z $L$ równe jeden i jest inna
od pozostałych. Łącznie jest więc
$6+15+6=27$ prostych.
\end{proof}

\begin{center}
\begin{minipage}{\linewidth}
\centering
\begin{tikzpicture}[>=Stealth,every node/.style={font=\small}]
 \node[draw,rounded corners,align=center,text width=3.7cm,inner sep=9pt] (a) at (0,0)
 {\textbf{6} wyjątkowych\\$E_1,\ldots,E_6$};
 \node[draw,rounded corners,align=center,text width=3.7cm,inner sep=9pt] (b) at (5.0,0)
 {\textbf{15} prostych przez pary\\$H-E_i-E_j$};
 \node[draw,rounded corners,align=center,text width=3.9cm,inner sep=9pt] (c) at (10.1,0)
 {\textbf{6} stożkowych przez piątki\\$2H-\sum_{j\ne i}E_j$};
 \draw[->] (a.east)--(b.west);
 \draw[->] (b.east)--(c.west);
 \node[align=center] at (5.1,-1.3)
  {$6+15+6=27$ różnych krzywych $(-1)$\\obrazowanych na proste przez $|-K_S|$};
\end{tikzpicture}

{\small Rysunek 6.3. Trzy typy prostych na gładkiej powierzchni kubicznej. Strzałki porządkują klasy, a nie oznaczają odwzorowań między krzywymi.}\par
\end{minipage}
\end{center}

\begin{przyklad}[27 prostych na konkretnej powierzchni Fermata]
\label{prz:ga6-fermat}
Rozważmy powierzchnię
$S_F=\V_+(X_0^3+X_1^3+X_2^3+X_3^3)\subseteq\PP^3_\C$.
Jest gładka, bo wspólne znikanie pochodnych
$3X_0^2,\ldots,3X_3^2$ wymuszałoby znikanie
wszystkich współrzędnych, co w przestrzeni rzutowej
jest niemożliwe. Podzielmy zbiór indeksów
$\{0,1,2,3\}$ na dwie nieuporządkowane pary
$\{i,j\}\sqcup\{l,m\}$; są dokładnie trzy takie
podziały. Wybierzmy dowolne $\zeta,\eta\in\C$
spełniające $\zeta^3=\eta^3=-1$ i nałóżmy
równości liniowe $X_i=\zeta X_j$,
$X_l=\eta X_m$. Wyznaczają one prostą projektową,
a jej każdy punkt należy do $S_F$, ponieważ
$(\zeta^3+1)X_j^3+(\eta^3+1)X_m^3=0$.
Z różnych par współczynników w ustalonym podziale wynikają różne proste;
różne podziały również nie dają tej samej prostej, gdyż
punkty z dwiema zerowymi współrzędnymi należą do niej tylko
wtedy, gdy obie niezerowe współrzędne tworzą jedną z wybranych par.
Mamy $3$ podziały i $3\cdot3=9$ wyborów pary
$(\zeta,\eta)$, łącznie $27$ różnych prostych.
Poprzednie twierdzenie zapewnia, że nie ma już innych.
To przykład, w którym słynny wynik można najpierw
\emph{zobaczyć w równaniach}, a potem zrozumieć,
dlaczego musi zachodzić dla każdej gładkiej powierzchni
kubicznej, a nie tylko dla tej symetrycznej.
\end{przyklad}

\begin{uwaga}[Dlaczego położenie ogólne jest potrzebne?]
Jeśli trzy rozdmuchiwane punkty leżą na jednej prostej,
jej transformata ścisła ma klasę $H-E_1-E_2-E_3$
i kwadrat $1-3=-2$. Ponadto
$(-K)\cdot(H-E_1-E_2-E_3)=3-3=0$,
co wyklucza obszerność $-K$ na tej powierzchni
(z obszerności wynika dodatni stopień na każdej
krzywej całkowitej). W szczególności nie można
bez sprawdzenia ogólnego położenia stosować
wyliczenia $27$ prostych.
\end{uwaga}

\section{Powierzchnie K3: przykład o innej geometrii}
\label{sec:ga6-k3}

\begin{definicja}[Powierzchnia K3]
Gładką, spójną, rzutową powierzchnię $S$ nad $\C$
nazywamy \emph{powierzchnią K3}, gdy
$\omega_S\cong\mathcal O_S$ oraz
$H^1(S,\mathcal O_S)=0$. Tutaj
$\omega_S=\det\Omega_S^1$ jest wiązką kanoniczną,
a $H^1$ oznacza pierwszą kohomologię snopa strukturalnego.
\end{definicja}

\begin{twierdzenie}[Gładka powierzchnia stopnia czwartego jest K3]
\label{tw:ga6-quartic-k3}
Niech $S=\V_+(F)\subseteq\PP^3_\C$, gdzie $F$
jest jednorodnym wielomianem stopnia cztery, a $S$
jest gładka. Wtedy $S$ jest powierzchnią K3.
\end{twierdzenie}
\begin{proof}
Z adjunkcji dla gładkiej hiperpowierzchni
(udowadnianej tak samo przez ciąg konormalny jak
w twierdzeniu~\ref{tw:ga6-adjunction}, z wyznacznikami
wiązek odpowiedniego rzędu) wynika
$\omega_S\cong(\omega_{\PP^3}\otimes
\mathcal O_{\PP^3}(4))|_S$.
Ponieważ $\omega_{\PP^3}\cong\mathcal O_{\PP^3}(-4)$,
otrzymujemy $\omega_S\cong\mathcal O_S$.

Aby obliczyć pierwszą kohomologię strukturalną,
skorzystajmy z dokładnego ciągu snopów
\[
0\longrightarrow\mathcal O_{\PP^3}(-4)
\xrightarrow{\cdot F}\mathcal O_{\PP^3}
\longrightarrow\mathcal O_S\longrightarrow0.
\]
Odpowiedni fragment długiego ciągu kohomologii ma postać
$H^1(\PP^3,\mathcal O_{\PP^3})\to
H^1(S,\mathcal O_S)\to
H^2(\PP^3,\mathcal O_{\PP^3}(-4))$.
Skrajne grupy są zerowe: ogólne obliczenie kohomologii
wiązek $\mathcal O_{\PP^n}(m)$ daje
$H^i(\PP^n,\mathcal O(m))=0$ dla $0<i<n$,
niezależnie od całkowitego $m$. W szczególności
$H^1(S,\mathcal O_S)=0$.
Pozostałe założenia z definicji (spójność i
rzutowość) są spełnione: gładka hiperpowierzchnia
rzutowa dodatniego stopnia w $\PP^3$ jest spójna;
można to również zobaczyć z
$H^0(S,\mathcal O_S)=\C$ otrzymanego z powyższego
ciągu. Tym samym $S$ jest powierzchnią K3.
\end{proof}

\begin{przyklad}[K3 nie może zawierać krzywej $(-1)$]
\label{prz:ga6-k3-no-minus1}
Jeśli $C\cong\PP^1$ jest gładką krzywą na powierzchni K3,
to $K_S\sim0$, więc ze wzoru adjunkcji
$-2=C^2+K_S\cdot C=C^2$.
Każda gładka krzywa wymierna ma zatem
samoprzecięcie $-2$, a nie $-1$. Wynika stąd,
że powierzchnia K3 jest minimalna. Nie oznacza to,
że nie może zawierać żadnych gładkich krzywych
wymiernych: na przykład proste leżące na niektórych
gładkich kwartykach są krzywymi $(-2)$.
\end{przyklad}

\begin{stwierdzenie}[Genus przekroju hiperpłaszczyzną K3]
\label{st:ga6-quartic-plane-section}
Niech $S\subseteq\PP^3_\C$ będzie gładką kwartyką,
a $C=S\cap H$ jej gładkim przekrojem płaszczyzną
$H\cong\PP^2$. Wtedy $g(C)=3$.
\end{stwierdzenie}
\begin{proof}
Na powierzchni $S$ klasa takiego przekroju
jest oznaczana przez $h=\mathcal O_S(1)$.
Jej samoprzecięcie jest równe stopniowi $S$,
czyli $h^2=4$: przecięcie dwóch ogólnych
płaszczyzn z $S$ składa się z czterech punktów
liczonych z krotnościami (twierdzenie Bézouta).
Ponieważ $K_S\sim0$, adjunkcja daje
$2g(C)-2=(K_S+h)\cdot h=4$.
Stąd $g(C)=3$. Jest to zarazem wzór na rodzaj
gładkiej płaskiej krzywej stopnia cztery z rozdziału 4.
\end{proof}

\begin{uwaga}[Wymierna powierzchnia a K3]
Powierzchnia wymierna ma $H^0(S,\omega_S)=0$,
a K3 ma $H^0(S,\omega_S)\cong\C$.
Wymiar $p_g(S):=h^0(S,\omega_S)$ jest
niezmiennikiem biwymiernym gładkich powierzchni
rzutowych (jest to osobne twierdzenie, wynikające
z zachowania regularnych form dwuwymiarowych
pod rozdmuchaniami i rozszerzania przez zbiory
kowymiaru dwa). Zatem żadna powierzchnia K3
nie może być wymierna.
To nie jest tylko różnica w liczbie punktów:
dwie gładkie powierzchnie mogą być obie rzutowe
nad $\C$, a ich geometria biwymierna może być
całkowicie odmienna.
\end{uwaga}

\section{Niezmienniki powierzchni i mapa dalszej teorii}
\label{sec:ga6-invariants}

\begin{definicja}[Niezmienniki kohomologiczne]
Dla gładkiej spójnej rzutowej powierzchni $S$ nad
$\C$ definiujemy:
\[
 q(S)=\dim_{\C}H^1(S,\mathcal O_S),\qquad
 p_g(S)=\dim_{\C}H^0(S,\omega_S),\qquad
 \chi(\mathcal O_S)=\sum_{i=0}^2(-1)^i
 \dim_{\C}H^i(S,\mathcal O_S).
\]
Liczbę $q(S)$ nazywamy \emph{nieregularnością},
a $p_g(S)$ \emph{rodzajem geometrycznym} powierzchni.
Wszystkie występujące grupy są skończenie wymiarowe,
co wynika z twierdzenia o skończoności kohomologii
snopów koherentnych na schemacie właściwym.
\end{definicja}

\begin{stwierdzenie}[Wzór na charakterystykę holomorficzną]
\label{st:ga6-chi}
Dla opisanej powierzchni $S$ zachodzi
$\chi(\mathcal O_S)=1-q(S)+p_g(S)$.
\end{stwierdzenie}
\begin{proof}
Ponieważ $S$ jest spójna, właściwa i zredukowana
nad algebraicznie domkniętym ciałem, jej globalne
funkcje regularne są stałe, zatem
$h^0(S,\mathcal O_S)=1$. Z dualności Serre'a
dla gładkiej rzutowej powierzchni wynika
$H^2(S,\mathcal O_S)^\vee\cong H^0(S,\omega_S)$,
więc $h^2(S,\mathcal O_S)=p_g(S)$.
Podstawienie do definicji $\chi$ daje tezę.
\end{proof}

\begin{uwaga}[Klasy Cherna i wzór Noether]
Dwa klasyczne niezmienniki liczbowe powierzchni
to $K_S^2$ i stopień drugiej klasy Cherna
$c_2(S)=\deg c_2(T_S)$. Dla gładkiej rzutowej
powierzchni nad $\C$ \emph{wzór Noether} głosi
$12\chi(\mathcal O_S)=K_S^2+c_2(S)$.
Ponadto przez twierdzenie Gaussa--Bonneta--Cherna
$c_2(S)$ jest topologiczną charakterystyką Eulera
$S(\C)$. Na tym etapie oba wyniki są zewnętrzne.
Pełne definicje klas Cherna i dowody odpowiednich
formuł należą do rozdziału 7, a relacja z topologią
będzie tematem rozdziału 10. Nie korzystamy z nich
w dowodach o $27$ prostych ani o kwartykach K3.
\end{uwaga}

\begin{definicja}[Wymiar Kodairy]
Niech $S$ będzie gładką rzutową powierzchnią nad $\C$.
Dla $m\ge1$ oznaczmy
$P_m(S)=h^0(S,\omega_S^{\otimes m})$.
Jeżeli wszystkie $P_m(S)$ są zerowe, kładziemy
$\kappa(S)=-\infty$. W przeciwnym razie
$\kappa(S)$ jest największą liczbą całkowitą
$\kappa\in\{0,1,2\}$, dla której istnieje stała
$c>0$ i nieskończony podciąg dodatnich liczb
$m_j\to\infty$ taki, że
$P_{m_j}(S)\ge c\,m_j^{\kappa}$.
Fakt, że w tej definicji wystarczają wykładniki
$0,1,2$, a otrzymana liczba zgadza się z wymiarem obrazu
plurikanonicznego dla dostatecznie podzielnych $m$,
jest głębszym twierdzeniem (teoria wymiaru Iitaki),
które tutaj przyjmujemy.
Wymiar Kodairy mierzy zatem tempo wzrostu
liczby niezależnych sekcji potęg wiązki kanonicznej.
\end{definicja}

\begin{przyklad}[Wymiar Kodairy w dwóch skrajnych sytuacjach]
Dla $\PP^2$ wiązka kanoniczna to $\mathcal O(-3)$,
więc $H^0(\PP^2,\omega^{\otimes m})=0$
dla każdego $m>0$; stąd $\kappa(\PP^2)=-\infty$.
Dla K3 wiązka kanoniczna jest trywialna,
więc $P_m(S)=1$ dla każdego $m\ge1$.
Zatem $\kappa(S)=0$. To wyjaśnia, dlaczego
powierzchnie wymierne i K3 trafiają do różnych
części klasyfikacji Enriquesa--Kodairy.
\end{przyklad}

\begin{uwaga}[Jak odczytywać plan następnych rozdziałów?]
Rozdział 6 wykorzystał lokalne obliczenia rozdmuchań,
ale ograniczył ogólną teorię przecięć do postaci
podstawowych twierdzeń, z wyraźnym wskazaniem miejsc,
w których cytujemy wyniki zewnętrzne.
\begin{enumerate}[label=(\roman*)]
\item W rozdziale 7 zbudujemy cykle, grupy Chowa,
      teorię przecięć i klasy Cherna oraz wrócimy
      do dowodów wzorów liczbowych z tego rozdziału.
\item W rozdziale 8 omówimy geometrię biwymierną,
      modele minimalne, geometrię stożków krzywych
      i kontrakcje.
\item Rozdziały 9 i 10 połączą geometrię algebraiczną
      z analizą zespoloną i topologią: GAGA, Kähler,
      Hodge, Lefschetz i niezmienniki topologiczne.
\end{enumerate}
\end{uwaga}

% Literatura kontrolna i jawne zewnetrzne zaleznosci:
% Hartshorne, Algebraic Geometry, rozdzial V (surfaces).
% Beauville, Complex Algebraic Surfaces.
% Stacks Project, tag 0CBN: SNC versus NC, tag 01OF: blowups.
% M. Reid, Undergraduate Algebraic Geometry (cubic surfaces).
% R. Vakil, lectures on 27 lines on cubic surfaces.



% ===============================================================
% ROZDZIAL 7. TEORIA PRZECIEC I KLASY CHERNA — ITERACJA 9
% Zalozenie domyslne: k algebraicznie domkniete, char(k)=0.
% Wszystkie etykiety maja prefiks ga7:.
% ===============================================================

\chapter{Teoria przecięć, grupy Chowa i klasy Cherna}
\label{ch:ga7-chow-chern}

\section{Motywacja: dlaczego krzywe musimy liczyć z krotnościami?}
\label{sec:ga7-motivation}

W rozdziale~\ref{ch:ga6-powierzchnie} spotkaliśmy dwie pozornie różne
operacje: przecinanie krzywych i obliczanie stopnia wiązki liniowej.
Są one przejawami tej samej teorii. Jeżeli $C,D$ są krzywymi na
gładkiej powierzchni $S$, a $C\cap D$ jest skończony, chcemy
uzyskać liczbę $C\cdot D$, która uwzględnia styczność,
pozostaje niezmiennicza przy zastępowaniu dzielnika innym
liniowo równoważnym i daje się obliczać bez rysowania
wszystkich punktów przecięcia.

\begin{przyklad}[Jedno przecięcie, lecz liczba większa od jedności]
Dla krzywych $C=\V(y)$ i $D=\V(y-x^3)$ w $\A^2_k$ jest
$C\cap D=\{(0,0)\}$, ale
$k[x,y]_{(x,y)}/(y,y-x^3)\cong k[x]_{(x)}/(x^3)$.
Otrzymujemy bazę $1,x,x^2$ i lokalną krotność $3$.
Dwa przypadkowe rysunki tych krzywych nie wykażą,
że w liczbie przecięć należy uwzględnić właśnie trzy.
Wynika to z algebry lokalnego pierścienia.
\end{przyklad}

Założeniem podstawowym tego rozdziału jest, że $k$ jest
ciałem algebraicznie domkniętym charakterystyki zero.
Tam, gdzie jest to potrzebne, $X$ będzie gładką,
quasi-rzutową rozmaitością algebraiczną nad $k$,
a $S$ gładką, spójną, rzutową powierzchnią nad $k$.
Część konstrukcji zachowuje sens dla schematów osobliwych,
lecz wówczas iloczyn cykli wymaga narzędzi bardziej ogólnych
niż zwykłe przecięcia podrozmaitości.

\section{Cykle algebraiczne i równoważność wymierna}
\label{sec:ga7-cycles}

\begin{definicja}[Cykle wymiaru $r$]
\label{def:ga7-cycles}
Niech $X$ będzie schematem skończonego typu nad $k$, a $r\ge0$.
\emph{Cyklem algebraicznym wymiaru $r$} na $X$ nazywamy formalną
sumę $\alpha=\sum_{i=1}^m n_i[V_i]$, gdzie $n_i\in\Z$,
zaś $V_i\subseteq X$ są całkowitymi domkniętymi podschematami
wymiaru $r$ (każdy występuje co najwyżej raz).
Grupę takich cykli oznaczamy $Z_r(X)$.
Przymiotnik \emph{całkowity} oznacza nierozkładalny i zredukowany.
Cyklem zerowym nazywamy element $Z_0(X)$.
\end{definicja}

\begin{uwaga}[Dlaczego współczynniki mogą być ujemne?]
Punkt i krzywa są obiektami geometrycznymi, ale chcemy także
odejmować cykle, tak jak odejmujemy dzielniki.
Cykle efektywne mają współczynniki nieujemne.
Przejście do grupy pozwala pisać relacje, takie jak
$[p]-[q]$, i sprowadzać problem geometryczny do algebry abelowej.
Dla powierzchni całkowitej $S$ grupa $Z_1(S)$
jest po prostu grupą dzielników Weila wprowadzoną w
definicji~\ref{def:ga6-divisor-group}.
\end{uwaga}

\begin{definicja}[Dzielnik funkcji wymiernej na krzywej]
\label{def:ga7-rational-divisor}
Niech $W\subseteq X$ będzie całkowitym, domkniętym
podschematem wymiaru $r+1$, a $\nu:\widetilde W\to W$
jego normalizacją. Dla $f\in k(W)^\times$ definiujemy
$r$-cykl $\operatorname{div}_W(f)$ przez
\[
\operatorname{div}_W(f):=\nu_*
\bigl(\operatorname{div}_{\widetilde W}(f)\bigr).
\]
Ponieważ $\widetilde W$ jest normalna, jej lokalne
pierścienie w punktach kowymiaru jeden są pierścieniami
dyskretnego wartościowania. Współczynniki dzielnika
$\operatorname{div}_{\widetilde W}(f)$ to odpowiednie rzędy
zer i biegunów. Dla skończonego morfizmu $\nu$ obraz
składnika $[\widetilde V]$ ma współczynnik
$[\kappa(\widetilde V):\kappa(\nu(\widetilde V))]$ przy
$[\nu(\widetilde V)]$, gdzie $\widetilde V$ jest
całkowitym podschematem kowymiaru jeden.
W szczególności dla całkowitej krzywej $W$ mamy
\[
\operatorname{div}_W(f)=
\sum_{q\in\widetilde W_{\mathrm{cl}}}
[\kappa(q):\kappa(\nu(q))]\operatorname{ord}_q(f)[\nu(q)].
\]
W rozpatrywanej sytuacji nad algebraicznie domkniętym
$k$ stopnie ciał reszt są równe jeden. Jeśli jednak
$W$ jest osobliwa, nie należy przypisywać jej punktom
osobliwym rzędów wartościowania bez przejścia
do normalizacji; kilka punktów normalizacji może
mieć ten sam obraz w $W$.
\end{definicja}

\begin{przyklad}[Dwa punkty normalizacji nad jednym węzłem]
\label{prz:ga7-nodal-divisor}
Niech $k$ będzie ciałem algebraicznie domkniętym
charakterystyki zero i
$W=\V(y^2-x^2(x+1))\subseteq\A_k^2$.
Jest to krzywa całkowita z węzłem $p=(0,0)$.
Jej normalizację opisuje morfizm
$\nu:\A_k^1\to W$ zadany przez
$x=t^2-1$ oraz $y=t(t^2-1)$.
Równość $t=y/x$ w ciele funkcji pokazuje,
że morfizm jest biwymierny, a równanie
$t^2=x+1$ — że jest skończony.
Nad węzłem $p$ leżą dwa punkty $t=1$ i $t=-1$.
Ponieważ funkcja $x=t^2-1=(t-1)(t+1)$
ma na normalizacji dwa proste zera,
otrzymujemy $\operatorname{div}_W(x)=2[p]$.
Ten współczynnik pochodzi z sumowania rzędów
w dwóch punktach normalizacji, a nie z pojedynczego
wartościowania lokalnego pierścienia węzła.
\end{przyklad}

\begin{definicja}[Grupa Chowa]
\label{def:ga7-chow}
Niech $R_r(X)$ będzie podgrupą $Z_r(X)$ generowaną przez
wszystkie cykle $\operatorname{div}_W(f)$, gdzie
$W\subseteq X$ jest całkowitym domkniętym podschematem
wymiaru $r+1$ i $f\in k(W)^\times$.
\emph{Grupą Chowa cykli wymiaru $r$} nazywamy
$\operatorname{CH}_r(X)=Z_r(X)/R_r(X)$.
Dwa cykle reprezentujące ten sam element $\operatorname{CH}_r(X)$
nazywamy \emph{wymiernie równoważnymi}.
Jeżeli $X$ jest czystego wymiaru $n$, to piszemy
$\operatorname{CH}^i(X)=\operatorname{CH}_{n-i}(X)$.
Dla $r<0$ przyjmujemy $\operatorname{CH}_r(X)=0$.
\end{definicja}

\begin{przyklad}[Punkty na prostej rzutowej]
\label{prz:ga7-p1-points}
Na $\PP^1_k$ wprowadźmy współrzędną wymierną $t$ na
mapie afinicznej, a punkt w nieskończoności oznaczmy
$\infty$. Funkcja $t-a$ ma proste zero w punkcie $a\in k$
i prosty biegun w $\infty$, zatem
$\operatorname{div}(t-a)=[a]-[\infty]$.
W $\operatorname{CH}_0(\PP^1)$ wszystkie punkty zamknięte
mają tę samą klasę. Z drugiej strony każda funkcja wymierna
na gładkiej krzywej właściwej ma tyle samo zer co biegunów,
licząc z krotnościami. Stąd stopień
$\deg(\sum n_i[p_i])=\sum n_i$ nie zmienia się przy
równoważności wymiernej. Każdy cykl zerowy jest wymiernie
równoważny $d[\infty]$, a jego stopień wynosi $d$;
otrzymujemy izomorfizm $\operatorname{CH}_0(\PP^1)\cong\Z$.
\end{przyklad}

\begin{przyklad}[Dlaczego zwartość geometryczna ma znaczenie?]
\label{prz:ga7-a1-chow0}
Na $\A^1_k$ funkcja $t-a$ ma dzielnik $[a]$:
nie ma tu punktu w nieskończoności, w którym należałoby
zapisać biegun. Wobec tego $[a]=0$ w
$\operatorname{CH}_0(\A^1)$ i $\operatorname{CH}_0(\A^1)=0$.
Nie oznacza to, że punkt $a$ nie istnieje. Oznacza jedynie,
że grupa Chowa pamięta również, czy rozpatrujemy przestrzeń
właściwą, czy otwartą.
\end{przyklad}

\begin{stwierdzenie}[Stopień cyklu zerowego na przestrzeni właściwej]
\label{st:ga7-degree}
Niech $X$ będzie właściwym schematem skończonego typu nad
algebraicznie domkniętym ciałem $k$. Wtedy przyporządkowanie
$\deg(\sum n_p[p])=\sum n_p$ określa homomorfizm
$\deg:\operatorname{CH}_0(X)\to\Z$.
\end{stwierdzenie}
\begin{proof}
Trzeba sprawdzić, że stopień znika na generatorach
$R_0(X)$. Taki generator jest dzielnikiem funkcji
wymiernej na całkowitej krzywej właściwej $W\subseteq X$.
Przejdźmy do normalizacji $\widetilde W$.
Jest ona gładką właściwą krzywą, a $f$ wyznacza
odwzorowanie wymierne $\widetilde W\to\PP^1$, które
na gładkiej właściwej krzywej przedłuża się do morfizmu.
Dla funkcji niestałej stopień włókna nad $0$
i nad $\infty$ jest ten sam, równy stopniowi morfizmu;
stąd $\deg\operatorname{div}(f)=0$. Dla funkcji stałej
dzielnik jest zerowy. Ponieważ przedłużenie morfizmu
na właściwej gładkiej krzywej oraz równość stopni włókien
zostały już wykorzystane przy dzielnikach krzywych,
nie potrzeba tu nowego założenia.
\end{proof}

\begin{uwaga}[Gdzie ukryta jest trudność?]
O ile definicja $\operatorname{CH}_r(X)$ wymaga tylko
cykli i dzielników funkcji wymiernych, o tyle iloczyn dwóch
klas cykli nie jest prostą operacją na formalnych sumach:
podrozmaitości mogą przecinać się w zbyt dużym wymiarze.
Dla gładkich przestrzeni algebraicznych problem rozwiązuje
\emph{lemat o przesuwaniu cykli} i konstrukcja przecięcia
z uwzględnieniem krotności. Wyniki te jawnie wskażemy.
\end{uwaga}

\section{Od cykli do przecięć: morfizmy i iloczyn}
\label{sec:ga7-intersection}

\begin{definicja}[Przecięcie właściwe]
\label{def:ga7-proper-intersection}
Niech $X$ będzie gładką całkowitą rozmaitością wymiaru $n$,
a $V,W\subseteq X$ całkowitymi domkniętymi podrozmaitościami
o kowymiarach $a,b$. Mówimy, że $V,W$ \emph{przecinają się
właściwie}, gdy każda składowa nierozkładalna $V\cap W$
ma kowymiar $a+b$ w $X$; gdy $a+b>n$, oznacza to, że
przecięcie jest puste.
Nie należy mylić tego pojęcia z \emph{właściwością morfizmu}.
\end{definicja}

\begin{przyklad}[Wspólna składowa a brak przecięcia właściwego]
Na $\PP^2$ krzywe $\V_+(XY)$ i $\V_+(XZ)$
mają wspólną prostą $X=0$. Ich przecięcie zbiorowe
nie jest zerowymiarowe, choć stopnie obu równań wynoszą dwa.
Nie wolno więc sumować lokalnych długości w punktach,
aby dostać cztery: powstałaby suma po nieskończonej liczbie
punktów. W pierścieniu Chowa obie krzywe mają klasy $2h$,
a iloczyn ich klas to $4h^2$, gdzie $h$ jest klasą prostej.
Dostęp do takiego iloczynu wymaga właśnie równoważności
wymiernej i zastąpienia dzielników odpowiednimi przedstawicielami.
\end{przyklad}

\begin{twierdzenie}[Iloczyn na gładkiej rozmaitości — wynik konstrukcyjny]
\label{tw:ga7-chow-ring}
Niech $X$ będzie gładką quasi-rzutową rozmaitością
wymiaru $n$ nad ciałem $k$. Wówczas grupy Chowa mają kanoniczną, przemienną,
łączną i stopniowaną strukturę pierścienia z jedynką na
$\operatorname{CH}^*(X)=\bigoplus_{i=0}^n\operatorname{CH}^i(X)$,
w której $\alpha\cdot\beta\in\operatorname{CH}^{i+j}(X)$
dla $\alpha\in\operatorname{CH}^i(X)$,
$\beta\in\operatorname{CH}^j(X)$, a $1=[X]$.
Jeżeli całkowite podrozmaitości $V,W$ przecinają się
właściwie, ich iloczyn jest sumą klas składowych
przecięcia z odpowiednimi dodatnimi krotnościami
lokalnymi. Dla dwóch dzielników na gładkiej powierzchni,
które nie mają wspólnej składowej, są to długości
$\operatorname{length}(\mathcal O_{X,p}/(f,g))$ z
rozdziału~\ref{ch:ga6-powierzchnie}.
\end{twierdzenie}

\begin{uwaga}[Dokładny status dowodu]
\label{uw:ga7-moving-lemma}
Ogólne twierdzenie~\ref{tw:ga7-chow-ring} wykorzystuje
lemat o przesuwaniu cykli, teorię przecięć \emph{Gysina}
i dowód niezależności od wyboru przedstawicieli.
Samo wymaganie zgodności z przecięciami właściwymi nie jest
tutaj traktowane jako niezależny dowód istnienia i jednoznaczności
całej struktury pierścienia; chodzi o kanoniczny iloczyn
konstruowany w tej teorii.
Nie ma krótkiego dowodu wyłącznie z aksjomatów pierścienia.
Traktujemy jego istnienie i własności funktorialne jako
\emph{zewnętrzny wynik konstrukcyjny}. Natomiast wzory dla
dzielników na powierzchniach poniżej zostaną wyprowadzone
ze stopni wiązek na krzywych. Dzięki temu
rachunki z rozdziału 6 nie opierają się na niewyjaśnionej
operacji mnożenia symboli.
\end{uwaga}

\begin{stwierdzenie}[Przecięcie z dzielnikiem przez stopień]
\label{st:ga7-divisor-degree}
Niech $k$ będzie ciałem algebraicznie domkniętym
charakterystyki zero, $S$ gładką rzutową powierzchnią
nad $k$, $C\subseteq S$ całkowitą krzywą,
$D$ dzielnikiem na $S$, a $\nu:\widetilde C\to C$
normalizacją. Wtedy liczba przecięcia $D\cdot C$ jest równa
$\deg\bigl(\nu^*(\mathcal O_S(D)|_C)\bigr)$.
Jest addytywna względem $D$, nie zmienia się po
zastąpieniu $D$ dzielnikiem liniowo równoważnym i
dla $D=C$ i $C$ gładkiej daje $C^2=\deg N_{C/S}$.
\end{stwierdzenie}
\begin{proof}
Równość z definicji ogólnego iloczynu stanowi
zastosowanie pierwszej klasy Cherna do cyklu $[C]$;
wyprowadzimy jej definicję w sekcji~\ref{sec:ga7-chern}.
Wyjaśnijmy już teraz, dlaczego pozostałe własności są
konieczne. Dla dzielników $D,D'$ mamy
$\mathcal O_S(D+D')\cong\mathcal O_S(D)\otimes\mathcal O_S(D')$.
Ponieważ stopień iloczynu tensorowego wiązek na gładkiej
właściwej krzywej jest sumą stopni, otrzymujemy
$(D+D')\cdot C=D\cdot C+D'\cdot C$.
Jeżeli $D-D'=\operatorname{div}(f)$, wiązki
$\mathcal O_S(D),\mathcal O_S(D')$ są izomorficzne,
więc ich stopnie na $\widetilde C$ są równe.
W przypadku $C$ gładkiej z definicji normalnej
$N_{C/S}\cong\mathcal O_S(C)|_C$,
a więc $C^2=\deg N_{C/S}$.
Identyfikacja \emph{tego} stopnia z sumą lokalnych
krotności przecięć dla różnych krzywych jest
istotnym twierdzeniem porównawczym teorii przecięć;
nie wynika z samej addytywności stopnia.
\end{proof}

\begin{przyklad}[Długości zamiast odgadywania krotności]
Niech $m\ge1$, $C=\V(y)$ i $D=\V(y-x^m)$ w $\A^2_k$.
W początku oba równania generują ideał $(y,x^m)$.
Domknięcia rzutowe tych krzywych mają równania
$Y=0$ oraz $YZ^{m-1}-X^m=0$, stopni odpowiednio
$1$ i $m$. Twierdzenie Bézouta daje łączną liczbę
przecięć równą $m$.
Sprawdźmy, że w nieskończoności nie ma przecięć.
Dla $m\ge2$ na $Z=0$ drugie równanie daje $X=0$,
zatem jedynym punktem drugiej krzywej w nieskończoności
jest $[0:1:0]$, który nie należy do prostej $Y=0$.
Dla $m=1$ drugie równanie ma postać $Y-X=0$
i jego punkt w nieskończoności to $[1:1:0]$;
on również nie należy do prostej $Y=0$.
Zatem cały wkład do liczby przecięć pochodzi
z początku $(0,0)$ i jest równy $m$.
\end{przyklad}

\begin{center}
\begin{minipage}{\linewidth}
\centering
\begin{tikzpicture}[>=Stealth,every node/.style={font=\small}]
 \node[draw,rounded corners,text width=3.0cm,align=center,minimum height=1.35cm] (a) at (-3.65,0)
 {Dwie krzywe bez wspólnej składowej\\$C\cap D$ skończony};
 \node[draw,rounded corners,text width=3.0cm,align=center,minimum height=1.35cm] (b) at (3.65,0)
 {Suma długości lokalnych\\$\sum_p i_p(C,D)$};
 \draw[->,thick] (a.east)--node[above,font=\footnotesize] {rachunek lokalny} (b.west);
 \node[draw,rounded corners,text width=5.4cm,align=center,minimum height=0.95cm] (c) at (0,-2.0)
 {Globalna klasa $[C]\cdot[D]$ w $\operatorname{CH}^2(S)$};
 \draw[->,thick] (a.south)--(c.north west);
 \draw[->,thick] (b.south)--(c.north east);
\end{tikzpicture}

{\small Rysunek 7.1. Dwie drogi do liczby przecięcia: lokalne pierścienie i globalne klasy cykli.}
\end{minipage}
\end{center}

\section{Podstawowe grupy Chowa: przestrzeń rzutowa}
\label{sec:ga7-projective-chow}

\begin{twierdzenie}[Pierścień Chowa przestrzeni rzutowej]
\label{tw:ga7-chow-pn}
Niech $k$ będzie ciałem algebraicznie domkniętym,
a $n\ge0$ liczbą całkowitą. Dla $n\ge1$ oznaczmy
przez $h=[\PP^{n-1}]\in\operatorname{CH}^1(\PP^n_k)$
klasę hiperpłaszczyzny. Wówczas
\[
\operatorname{CH}^*(\PP^n_k)\cong\Z[h]/(h^{n+1}).
\]
Jeżeli $n=0$, lewą stroną jest $\Z$, a zapis po prawej
oznacza $\Z[h]/(h)\cong\Z$, z formalnym $h=0$.
Dla $n\ge1$ klasa $h^j$ jest klasą liniowej podprzestrzeni
kowymiaru $j$; we wszystkich przypadkach $\deg(h^n)=1$
(przy $n=0$ jest to stopień $[\PP^0_k]$).
\end{twierdzenie}
\begin{proof}
Dla $n=0$ mamy $\PP^0_k=\Spec k$ oraz
$\operatorname{CH}^0(\PP^0_k)=\Z[\PP^0_k]$,
więc twierdzenie zachodzi z $h=0$.
Załóżmy dalej, że $n\ge1$.
Przyjmujemy dwa standardowe fakty o grupach Chowa:
\emph{ciąg lokalizacji} dla domkniętego zanurzenia
$Z\hookrightarrow X$ z otwartym dopełnieniem $U$,
$\operatorname{CH}_r(Z)\to\operatorname{CH}_r(X)
\to\operatorname{CH}_r(U)\to0$, oraz
\emph{niezmienniczość homotopijną}, z której wynika
$\operatorname{CH}_r(\A^n)=0$ dla $r<n$ i
$\operatorname{CH}_n(\A^n)=\Z[\A^n]$.
Dla $X=\PP^n$, $Z=\PP^{n-1}$ będącej hiperpłaszczyzną,
$U=\A^n$ dostajemy, dla $r<n$, surjekcję
$\operatorname{CH}_r(\PP^{n-1})\to\operatorname{CH}_r(\PP^n)$.
Indukcja po $n$ pokazuje, że $\operatorname{CH}_r(\PP^n)$
jest generowana przez klasę liniowej $\PP^r$.
Dla $r=n$ generatorem jest $[\PP^n]$.
Nie może być niezerowego stosunku $m[\PP^r]=0$:
po przecięciu z $r$ ogólnymi hiperpłaszczyznami otrzymujemy
$m$ razy klasę punktu, o stopniu $m$.
Homomorfizm stopnia z
stwierdzenia~\ref{st:ga7-degree} pokazuje, że $m=0$.
Zatem wszystkie te grupy są izomorficzne z $\Z$.
Iloczyn $j$ ogólnych hiperpłaszczyzn jest liniową
podprzestrzenią kowymiaru $j$, zatem jej klasą jest $h^j$.
Przy $j=n+1$ wymiar cyklu stałby się ujemny, więc
$h^{n+1}=0$. Nie ma innych relacji, gdyż
$1,h,\ldots,h^n$ generują niezależne grupy
w różnych stopniach.
\end{proof}

\begin{przyklad}[Bézout jako mnożenie klas]
\label{prz:ga7-bezout}
Niech $d\ge1$, a $F\ne0$ będzie jednorodnym
wielomianem stopnia $d$. Wówczas $F$ definiuje
w $\PP^n$ efektywny dzielnik klasy $dh$.
Jeżeli $F_1,\ldots,F_n$ są niezerowymi jednorodnymi
wielomianami dodatnich stopni $d_1,\ldots,d_n$
i ich wspólny zbiór zer jest zerowymiarowy (bez składowych
wymiaru dodatniego), to jego cykl przecięcia ma klasę
$(d_1\cdots d_n)h^n$. Stopień tego cyklu jest
$d_1\cdots d_n$. Dla dwóch krzywych płaskich daje to
$de$, a dla trzech powierzchni w $\PP^3$ --- $d_1d_2d_3$.
Poszczególne punkty przecięcia mogą mieć różne krotności,
lecz ich suma jest ustalona.
\end{przyklad}

\begin{uwaga}[Co zostało dowiedzione samodzielnie?]
Po przyjęciu ciągu lokalizacji, niezmienniczości
homotopijnej i konstrukcji iloczynu dla gładkich
rozmaitości dowód pierścienia Chowa $\PP^n$ jest
już rachunkiem elementarnym. Nie twierdzimy, że
owe trzy twierdzenia funktorialne zostały dowiedzione
od podstaw w tym skrypcie.
\end{uwaga}

\section{Klasy Cherna: algebraiczny opis skręcenia wiązki}
\label{sec:ga7-chern}

Wiązki liniowe i ich dzielniki mierzyły już, ile zer
ma przekrój na krzywej. Dla wiązki wektorowej rzędu
większego niż jeden pojedynczy dzielnik nie wystarcza.
Klasy Cherna są wielowymiarowym uogólnieniem stopnia:
przypisują wiązce $\mathcal E$ kolejne klasy
$c_i(\mathcal E)$ w pierścieniu Chowa.

\begin{definicja}[Pierwsza klasa Cherna]
\label{def:ga7-c1}
Niech $X$ będzie gładką quasi-rzutową rozmaitością,
a $L$ wiązką liniową na $X$. Jeśli $L\cong\mathcal O_X(D)$
dla dzielnika Cartiera $D$, definiujemy
$c_1(L)=[D]\in\operatorname{CH}^1(X)$.
Jest to poprawna definicja, ponieważ dzielniki liniowo
równoważne reprezentują tę samą klasę, a na gładkiej
rozmaitości $\operatorname{Pic}(X)\cong\operatorname{CH}^1(X)$.
Dla wiązki trywialnej $c_1(\mathcal O_X)=0$.
\end{definicja}

\begin{stwierdzenie}[Pierwsza klasa Cherna i iloczyn tensorowy]
\label{st:ga7-c1-add}
Niech $X$ będzie gładką quasi-rzutową rozmaitością,
a $L,M$ wiązkami liniowymi na $X$. Wtedy
$c_1(L\otimes M)=c_1(L)+c_1(M)$ oraz
$c_1(L^\vee)=-c_1(L)$.
\end{stwierdzenie}
\begin{proof}
Wybierzmy dzielniki $D,E$ z $L\cong\mathcal O_X(D)$
i $M\cong\mathcal O_X(E)$. Z konstrukcji snopa dzielnika
wynika $L\otimes M\cong\mathcal O_X(D+E)$, a
$L^\vee\cong\mathcal O_X(-D)$.
Przejście do klas cykli daje obie równości.
\end{proof}

\begin{twierdzenie}[Formuła wiązki rzutowej — wynik zewnętrzny]
\label{tw:ga7-projective-bundle}
Niech $X$ będzie gładką quasi-rzutową rozmaitością,
$\mathcal E$ wiązką wektorową rzędu $r\ge1$, a
$\pi:\PP_X(\mathcal E)\to X$ przestrzenią
jednowymiarowych ilorazów włókien $\mathcal E$.
Oznaczmy $\xi=c_1(\mathcal O_{\PP(\mathcal E)}(1))$.
Wówczas $\operatorname{CH}^*(\PP_X(\mathcal E))$ jest
wolnym modułem nad $\operatorname{CH}^*(X)$ o bazie
$1,\xi,\ldots,\xi^{r-1}$; działanie współczynników
jest dane przez $\pi^*$.
\end{twierdzenie}

\begin{definicja}[Wyższe klasy Cherna]
\label{def:ga7-c-classes}
Niech $X$ będzie gładką quasi-rzutową rozmaitością nad $k$,
$\mathcal E$ wiązką wektorową rzędu $r\ge1$,
$\pi:\PP_X(\mathcal E)\to X$ projektywizacją
jednowymiarowych ilorazów, a
$\xi=c_1(\mathcal O_{\PP(\mathcal E)}(1))$.
Z formuły wiązki rzutowej~\ref{tw:ga7-projective-bundle}
istnieją jednoznaczne klasy
$c_i(\mathcal E)\in\operatorname{CH}^i(X)$, $1\le i\le r$,
takie, że
\[
\xi^r-\pi^*c_1(\mathcal E)\xi^{r-1}
+\pi^*c_2(\mathcal E)\xi^{r-2}
-\cdots+(-1)^r\pi^*c_r(\mathcal E)=0.
\]
Definiujemy także $c_0(\mathcal E)=1$ i
$c_i(\mathcal E)=0$ dla $i>r$.
\emph{Całkowitą klasą Cherna} nazywamy
$c(\mathcal E)=\sum_{i=0}^r c_i(\mathcal E)$.
\end{definicja}

\begin{uwaga}[Skąd znaki w definicji?]
Przy naszej konwencji ilorazów, jeśli $\mathcal E=L$
ma rząd jeden, to $\PP_X(L)=X$ i
$\mathcal O_{\PP(L)}(1)\cong L$. Równanie staje się
$\xi-c_1(L)=0$, zgodnie z wcześniejszą definicją.
Dla sumy dwóch wiązek liniowych $L\oplus M$ relacja
przyjmuje postać
$\xi^2-(c_1(L)+c_1(M))\xi+c_1(L)c_1(M)=0$.
Różne konwencje projektywizacji zmieniają postać
równań; tutaj będziemy konsekwentnie używać ilorazów.
\end{uwaga}

\begin{twierdzenie}[Własności klas Cherna]
\label{tw:ga7-chern-axioms}
Niech $X$ będzie gładką quasi-rzutową rozmaitością.
Dla wiązek wektorowych $\mathcal E,\mathcal F$ na $X$
zachodzą następujące własności:
\begin{enumerate}[label=(\roman*)]
\item $c_1(\mathcal E)=c_1(\det\mathcal E)$.
\item Jeśli $0\to\mathcal E'\to\mathcal E\to\mathcal E''\to0$
      jest dokładnym ciągiem wiązek wektorowych, to
      $c(\mathcal E)=c(\mathcal E')c(\mathcal E'')$
      (\emph{wzór Whitneya}).
\item Dla morfizmu $f:Y\to X$ między gładkimi
      quasi-rzutowymi rozmaitościami zachodzi
      $c_i(f^*\mathcal E)=f^*c_i(\mathcal E)$.
\item Jeżeli $\mathcal E=L_1\oplus\cdots\oplus L_r$
      jest sumą wiązek liniowych, to
      $c(\mathcal E)=\prod_{j=1}^r(1+c_1(L_j))$.
\end{enumerate}
\end{twierdzenie}
\begin{proof}[Dowód z wykorzystaniem zasady rozszczepienia]
Najpierw wyjaśnimy zasadę rozszczepienia. Dla wiązki
$\mathcal E$ można skonstruować kolejne projektywizacje
wiązek, prowadzące do morfizmu $p:Y\to X$, dla którego
$p^*:\operatorname{CH}^*(X)\to\operatorname{CH}^*(Y)$
jest iniektywne, a $p^*\mathcal E$ posiada filtrację
przez podwiązki, której ilorazy są wiązkami liniowymi
$L_1,\ldots,L_r$. Iniektywność na każdym etapie
wynika z formuły wiązki rzutowej: pullback
identyfikuje pierścień z podmodułem współczynnika
przy $\xi^0$ w swobodnej bazie $1,\xi,\ldots,\xi^{r-1}$.
Pełna konstrukcja flag przez kolejne uniwersalne ilorazy
stanowi standardową część teorii wiązek rzutowych.

Dla sumy wiązek liniowych definicja przez relację
projektową daje
$c(\mathcal E)=\prod_j(1+x_j)$,
gdzie $x_j=c_1(L_j)$; dla filtracji z ilorazami liniowymi
obowiązuje ten sam wzór. Jego wykazanie to podstawowy
etap konstrukcyjny teorii klas Cherna, który
przyjmujemy wraz z formułą wiązki rzutowej.
Należy podkreślić, że sama iniektywność odwzorowania indukowanego
nie dowodzi jeszcze tej tożsamości: jej uzasadnienie
wymaga zgodności klas Cherna z filtracjami wiązek.
Wyrażenia $c_i$ stają się więc elementarnymi funkcjami
symetrycznymi zmiennych $x_j$, zwanych
\emph{formalnymi pierwiastkami Cherna}.

Jeżeli $0\to\mathcal E'\to\mathcal E\to\mathcal E''\to0$
jest dokładny, wybierzmy wspólną przestrzeń rozszczepienia
dla $\mathcal E'$ i $\mathcal E''$, z iniektywnym
odwzorowaniem indukowanym na grupach Chowa. Filtracja
$\mathcal E'$ i przeciwobrazy podwiązek z filtracji
$\mathcal E''$ dają filtrację $\mathcal E$,
której ilorazy są łącznie ilorazami obu wiązek.
W tej przestrzeni mamy zatem
$c(\mathcal E)=c(\mathcal E')c(\mathcal E'')$;
iniektywność odwzorowania indukowanego daje (ii) na $X$.
Przypadek rozszczepionego ciągu daje (iv).
Ponieważ wyznacznik ma pierwszą klasę Cherna równą
sumie pierwszych klas ilorazów, otrzymujemy również (i).

Własność (iii) można odczytać bezpośrednio z definicji:
projektywizacja i tautologiczna wiązka ilorazowa
są zgodne ze zmianą bazy, a pullback zachowuje relację
wielomianową definiującą klasy $c_i$. Jednoznaczność
współczynników zapewniona przez formułę wiązki rzutowej
daje $c_i(f^*\mathcal E)=f^*c_i(\mathcal E)$.
W ten sposób dowód redukuje się do \emph{dwóch} jawnych
wyników konstrukcyjnych: formuły wiązki rzutowej
i zasady rozszczepienia, zamiast przyjmować osobno
każdą z czterech własności.
\end{proof}

\begin{przyklad}[Wiązka sumy dwóch liniowych na $\PP^2$]
\label{prz:ga7-chern-sum}
Niech $h=c_1(\mathcal O_{\PP^2}(1))$,
$\mathcal E=\mathcal O_{\PP^2}(2)\oplus\mathcal O_{\PP^2}(3)$.
Ponieważ $c(\mathcal O(d))=1+dh$, wzór Whitneya daje
$c(\mathcal E)=(1+2h)(1+3h)=1+5h+6h^2$.
Zatem $c_1(\mathcal E)=5h$ i $c_2(\mathcal E)=6h^2$.
Liczba $\deg c_2(\mathcal E)=6$ ma również
interpretację jako liczba zer ogólnego przekroju
tej wiązki, liczonych z krotnościami (pod warunkiem,
że zbiór zer ma oczekiwany wymiar zero).
\end{przyklad}

\begin{stwierdzenie}[Zera przekroju i najwyższa klasa Cherna]
\label{st:ga7-zero-section}
Niech $X$ będzie gładką rzutową rozmaitością wymiaru $n$,
$\mathcal E$ wiązką wektorową rzędu $n$ na $X$,
a $s$ jej regularnym przekrojem, którego schemat zer
$Z(s)$ jest zerowymiarowy. Wtedy
$[Z(s)]=c_n(\mathcal E)\cap[X]$ w $\operatorname{CH}_0(X)$,
a zatem długość $Z(s)$ wynosi $\deg c_n(\mathcal E)$.
\end{stwierdzenie}
\begin{uwaga}[Dlaczego warunek regularności jest potrzebny?]
Przekrój może zerować się na całej krzywej albo
być identycznie zerowy. Wówczas nie można po prostu
liczyć jego zer jako izolowanych punktów. Twierdzenie
o klasie najwyższej wymaga przekroju regularnego,
tj. takiego, którego schemat zer ma kowymiar równy
rzędowi wiązki (lub pustego zbioru zer, gdy klasa znika).
Jego dowód w pełnej ogólności korzysta z teorii
przecięć regularnych zanurzeń; oznaczamy go jako
\emph{zewnętrzny}. Konkretny przykład poniżej
możemy jednak obliczyć bezpośrednio.
\end{uwaga}

\section{Wiązki styczne przestrzeni rzutowych}
\label{sec:ga7-tangent}

\begin{twierdzenie}[Ciąg Eulera i klasy Cherna $\PP^n$]
\label{tw:ga7-euler}
Niech $k$ będzie ciałem algebraicznie domkniętym
charakterystyki zero, $n\ge1$ i
$h=c_1(\mathcal O_{\PP^n_k}(1))$.
Wówczas istnieje dokładny ciąg wiązek
\[
0\longrightarrow\mathcal O_{\PP^n}
\longrightarrow\mathcal O_{\PP^n}(1)^{\oplus(n+1)}
\longrightarrow T_{\PP^n}\longrightarrow0,
\]
a całkowita klasa Cherna wiązki stycznej jest równa
$c(T_{\PP^n})=(1+h)^{n+1}$ w
$\Z[h]/(h^{n+1})$. W szczególności
$c_i(T_{\PP^n})=\binom{n+1}{i}h^i$ dla $0\le i\le n$.
\end{twierdzenie}
\begin{proof}
Oznaczmy przez $V=k^{n+1}$ przestrzeń, której proste
parametryzuje $\PP^n=\PP(V)$ w konwencji prostych.
Wybór niezerowego wektora $v\in\ell$ reprezentującego
$[\ell]$ daje opis
$T_{[\ell]}\PP(V)\cong\operatorname{Hom}(\ell,V/\ell)$:
deformacja prostej jest odwzorowaniem z jej kierunku
do kierunków poprzecznych. Na całej $\PP(V)$ otrzymujemy
ciąg wiązek
$0\to\mathcal O\to\mathcal O(1)\otimes V
\to T_{\PP(V)}\to0$, przy czym pierwszy morfizm
wysyła lokalnie stałą $1$ na sekcję złożoną z
współrzędnych jednorodnych. Jego obraz jest podwiązką,
a iloraz ma włókno $\operatorname{Hom}(\ell,V/\ell)$.
To daje ciąg Eulera. Ze wzoru Whitneya,
$c(\mathcal O)=1$ i $c(\mathcal O(1))=1+h$
otrzymujemy $c(T_{\PP^n})=(1+h)^{n+1}$.
Ponieważ $h^{n+1}=0$, rozwinięcie dwumianowe
kończy się w stopniu $n$.
\end{proof}

\begin{przyklad}[Trzy zera pola wektorowego na $\PP^2$]
\label{prz:ga7-vector-field}
Z twierdzenia~\ref{tw:ga7-euler} mamy
$c_2(T_{\PP^2})=3h^2$, a więc jego stopień wynosi $3$.
Rozważmy algebraiczne działanie grupy multiplikatywnej
$\mathbb G_m$ na $\PP^2_k$ dane przez
$\lambda\cdot[x:y:z]=[x:\lambda y:\lambda^2 z]$.
Różniczkując to działanie w jedynce $\lambda=1$,
otrzymujemy globalne algebraiczne pole wektorowe.
Nad $k=\C$ tę samą konstrukcję można zapisać
analitycznie, podstawiając $\lambda=e^t$
i różniczkując względem $t$ w zerze.
Pole znika dokładnie w trzech punktach współrzędnych
$[1:0:0]$, $[0:1:0]$, $[0:0:1]$.
W pobliżu pierwszego punktu na mapie $x\ne0$
współrzędne $(u=y/x,v=z/x)$ spełniają
$\dot u=u$, $\dot v=2v$; zatem zero jest proste.
Na pozostałych dwóch mapach otrzymujemy odpowiednio
wagi $(-1,1)$ oraz $(-2,-1)$ i znowu proste zera.
Łączna długość schematu zer wynosi $1+1+1=3$,
co dokładnie zgadza się z $\deg c_2(T_{\PP^2})$.
\end{przyklad}

\begin{center}
\begin{minipage}{\linewidth}
\centering
\begin{tikzpicture}[>=Stealth,scale=0.95,every node/.style={font=\small}]
 \coordinate (A) at (-2.8,0); \coordinate (B) at (2.8,0); \coordinate (C) at (0,2.95);
 \draw[thick] (A)--(B)--(C)--cycle;
 \fill[blue!70!black] (A) circle (2.6pt);
 \fill[blue!70!black] (B) circle (2.6pt);
 \fill[blue!70!black] (C) circle (2.6pt);
 \node[below left=4pt of A,font=\footnotesize] {$[1:0:0]$};
 \node[below right=4pt of B,font=\footnotesize] {$[0:1:0]$};
 \node[above=7pt of C,font=\footnotesize] {$[0:0:1]$};
 \node[align=center,font=\footnotesize] at (0,0.87) {trzy zera};
\end{tikzpicture}

{\small Rysunek 7.2. Schematyczne położenie trzech zer algebraicznego pola wektorowego na $\PP^2$; trójkąt przedstawia jedynie współrzędne wierzchołkowe, nie całą rzeczywistą geometrię $\PP^2(\C)$.}
\end{minipage}
\end{center}

\section{Gładkie hiperpowierzchnie i ich niezmienniki}
\label{sec:ga7-hypersurfaces}

\begin{twierdzenie}[Klasy Cherna hiperpowierzchni]
\label{tw:ga7-hypersurface-chern}
Niech $X=\V_+(F)\subseteq\PP^n_k$ będzie gładką
hiperpowierzchnią stopnia $d\ge1$, gdzie $n\ge2$.
Oznaczmy $h=c_1(\mathcal O_X(1))$.
W pierścieniu $\operatorname{CH}^*(X)$ zachodzi
\[
 c(T_X)=\frac{(1+h)^{n+1}}{1+dh},
\]
przy czym prawą stronę rozumiemy jako szereg formalny
obcięty w stopniu większym niż $\dim X=n-1$.
W szczególności $c_1(T_X)=(n+1-d)h$.
\end{twierdzenie}
\begin{proof}
Ponieważ $X$ jest gładka, jest regularną hiperpowierzchnią,
a wiązka normalna jej zanurzenia w $\PP^n$ jest
$N_{X/\PP^n}\cong\mathcal O_X(d)$; wynika to
z tego, że ideał $X$ jest lokalnie generowany przez $F$,
zaś jego snop to $\mathcal O_{\PP^n}(-d)$.
Mamy więc dokładny ciąg wiązek stycznych
$0\to T_X\to T_{\PP^n}|_X\to\mathcal O_X(d)\to0$.
Ze wzoru Whitneya otrzymujemy
$c(T_{\PP^n}|_X)=c(T_X)c(\mathcal O_X(d))$.
Naturalność klas Cherna i ciąg Eulera dają
$c(T_{\PP^n}|_X)=(1+h)^{n+1}$, natomiast
$c(\mathcal O_X(d))=1+dh$.
Ponieważ element $dh$ ma dodatni stopień, $1+dh$
ma odwrotność w skończenie stopniowanym pierścieniu
Chowa $X$. Dzieląc, otrzymujemy wzór.
\end{proof}

\begin{wniosek}[Gładkie powierzchnie stopnia $d$ w $\PP^3$]
\label{wn:ga7-degree-d-surface}
Niech $k$ będzie ciałem algebraicznie domkniętym
charakterystyki zero, $S\subseteq\PP^3_k$ gładką
powierzchnią stopnia $d\ge1$, $H$ dzielnikiem
hiperpłaszczyznowym na $S$ i $h=[H]=c_1(\mathcal O_S(1))$.
Dla dowolnego dzielnika kanonicznego $K_S$ zachodzą
\[
 K_S\sim(d-4)H,\qquad K_S^2=d(d-4)^2,\qquad
 c_2(T_S)=(d^2-4d+6)h^2,
\]
a $\deg c_2(T_S)=d(d^2-4d+6)$.
\end{wniosek}
\begin{proof}
Dla $n=3$ rozwijamy iloraz
$(1+h)^4/(1+dh)$ do drugiego stopnia:
$(1+4h+6h^2)(1-dh+d^2h^2)
=1+(4-d)h+(d^2-4d+6)h^2$.
Ponadto z dokładnego ciągu wiązek normalnych
$0\to T_S\to T_{\PP^3}|_S\to\mathcal O_S(d)\to0$
i tożsamości $\det T_{\PP^3}\cong\mathcal O_{\PP^3}(4)$
wynika po wzięciu wyznaczników
$\det T_S\cong\mathcal O_S(4-d)$, a więc
$\omega_S\cong\mathcal O_S(d-4)$.
Jest to izomorfizm wiązek liniowych, dlatego
$K_S\sim(d-4)H$, nie tylko równość ich klas w grupie Chowa.
Ponieważ $S$ ma stopień $d$, dwa ogólne przekroje
hiperpłaszczyzną przecinają się na $S$ w cyklu zerowym
stopnia $d$, tzn. $\deg(h^2)=d$.
Zatem $K_S^2=(d-4)^2\deg(h^2)=d(d-4)^2$,
a stopień drugiej klasy Cherna jest równy
$d(d^2-4d+6)$.
\end{proof}

\begin{przyklad}[Cztery powierzchnie, cztery różne drugie klasy]
\label{prz:ga7-surface-invariants}
Dla $d=1,2,3,4$ wzory dają odpowiednio:
\[
\begin{array}{c|c|c|c}
 d & \text{powierzchnia} & K_S^2 & \deg c_2(T_S)\\ \hline
 1 & \PP^2 & 9 & 3\\
 2 & \text{gładka kwadryka} & 8 & 4\\
 3 & \text{gładka powierzchnia kubiczna} & 3 & 9\\
 4 & \text{gładka kwartyka K3} & 0 & 24
\end{array}
\]
Wartość $24$ dla kwartyki K3 nie jest osobnym aksjomatem
definicji K3: wynika z obliczenia wiązki stycznej.
\end{przyklad}

\section{Riemann--Roch na powierzchniach: dowód z twierdzenia dla krzywych}
\label{sec:ga7-rr-surface}

Tutaj możemy udowodnić ważne twierdzenie, nie przyjmując
od razu całej teorii klas charakterystycznych. Punktem wyjścia
będzie Riemann--Roch dla \emph{gładkiej krzywej}, którego
dowodowe podstawy zostały omówione w rozdziale 4.

\begin{definicja}[Charakterystyka Eulera snopa koherentnego]
\label{def:ga7-chi}
Niech $S$ będzie gładką rzutową powierzchnią, a
$\mathcal F$ snopem koherentnym na $S$. Definiujemy
$\chi(S,\mathcal F)=\sum_{i=0}^{2}(-1)^i
\dim_k H^i(S,\mathcal F)$.
Wymiary są skończone dzięki twierdzeniu o skończoności
kohomologii koherentnej na schemacie właściwym.
Gdy $\mathcal F=\mathcal O_S(D)$, piszemy
$\chi(\mathcal O_S(D))$.
\end{definicja}

\begin{lemat}[Addytywność charakterystyki Eulera]
\label{lem:ga7-chi-exact}
Niech $0\to\mathcal F'\to\mathcal F\to\mathcal F''\to0$
będzie dokładnym ciągiem snopów koherentnych na
gładkiej rzutowej powierzchni $S$. Wtedy
$\chi(\mathcal F)=\chi(\mathcal F')+\chi(\mathcal F'')$.
\end{lemat}
\begin{proof}
Ciąg krótki indukuje długi ciąg kohomologii
$H^0(\mathcal F')\to H^0(\mathcal F)\to\cdots
\to H^2(\mathcal F'')\to0$.
Jest to skończony ciąg dokładny skończenie wymiarowych
przestrzeni wektorowych. Suma naprzemienna ich wymiarów
wynosi zero: jeśli $V_i$ są kolejnymi wyrazami, to
$\dim V_i=\dim\operatorname{im}(V_{i-1}\to V_i)+
\dim\operatorname{im}(V_i\to V_{i+1})$, a po zsumowaniu
naprzemiennym wszystkie wyrazy obrazu znoszą się parami.
Zgrupowanie wyrazów według trzech snopów daje tezę.
\end{proof}

\begin{lemat}[Przyrost charakterystyki po dodaniu krzywej]
\label{lem:ga7-rr-increment}
Niech $k$ będzie ciałem algebraicznie domkniętym charakterystyki
zero, $S$ gładką rzutową powierzchnią nad $k$,
$K_S$ dowolnym dzielnikiem kanonicznym na $S$,
$D$ dowolnym dzielnikiem, a $C\subseteq S$ gładką
całkowitą krzywą. Wtedy
\[
\chi(\mathcal O_S(D+C))-\chi(\mathcal O_S(D))
= D\cdot C+\frac{C^2-K_S\cdot C}{2}.
\]
\end{lemat}
\begin{proof}
Krzywa $C$ jest efektywnym dzielnikiem Cartiera.
Pomnożenie ciągu
$0\to\mathcal O_S(-C)\to\mathcal O_S\to\mathcal O_C\to0$
przez wiązkę liniową $\mathcal O_S(D+C)$ jest dokładne,
ponieważ iloczyn tensorowy z wiązką lokalnie wolną
zachowuje dokładność. Otrzymujemy
\[
0\longrightarrow\mathcal O_S(D)
\longrightarrow\mathcal O_S(D+C)
\longrightarrow\mathcal O_C(D+C)\longrightarrow0.
\]
Lemat~\ref{lem:ga7-chi-exact} daje
$\chi(\mathcal O_S(D+C))-\chi(\mathcal O_S(D))
=\chi(\mathcal O_C(D+C))$; tu korzystamy także z tego,
że kohomologie snopa $i_*L$ na $S$ są kohomologiami
$L$ na $C$, gdzie $i:C\hookrightarrow S$ jest zanurzeniem.
Riemann--Roch dla gładkiej właściwej krzywej
$C$ mówi, że dla wiązki $L$ na $C$ zachodzi
$\chi(C,L)=\deg L+1-g(C)$. Po podstawieniu
$L=\mathcal O_C(D+C)$ otrzymujemy
$\chi= D\cdot C+C^2+1-g(C)$.
Wzór adjunkcji z twierdzenia~\ref{tw:ga6-adjunction}
daje $1-g(C)=-\frac12(C^2+K_S\cdot C)$.
Łącząc te równości, uzyskujemy tezę.
\end{proof}

\begin{twierdzenie}[Riemann--Roch dla dzielników na powierzchni]
\label{tw:ga7-rr-surfaces}
Niech $S$ będzie gładką spójną rzutową powierzchnią
nad algebraicznie domkniętym ciałem charakterystyki zero,
$K_S$ jej dzielnikiem kanonicznym, a $D$ dowolnym dzielnikiem.
Wtedy
\[
\chi(\mathcal O_S(D))=\chi(\mathcal O_S)
+\frac{D\cdot(D-K_S)}{2}.
\]
\end{twierdzenie}
\begin{proof}
Zdefiniujmy funkcję na grupie dzielników
\[
Q(D)=\chi(\mathcal O_S(D))-\chi(\mathcal O_S)
-\frac{D^2-K_S\cdot D}{2}.
\]
Z lematu~\ref{lem:ga7-rr-increment} wynika, że
dla każdej gładkiej całkowitej krzywej $C$ mamy
$Q(D+C)=Q(D)$, ponieważ
\[
\frac{(D+C)^2-K_S\cdot(D+C)}{2}
-\frac{D^2-K_S\cdot D}{2}
=D\cdot C+\frac{C^2-K_S\cdot C}{2}.
\]
W szczególności, podstawiając $D-C$ za $D$,
dostajemy również $Q(D-C)=Q(D)$.

Pozostaje sprawdzić, że każdy dzielnik jest liniowo
równoważny różnicy dwóch gładkich całkowitych krzywych.
Wybierzmy bardzo obszerny dzielnik $H$.
Dla dostatecznie dużego $m$ zarówno $mH$,
jak i $D+mH$ są bardzo obszerne (jest to znany
wynik Serre'a o obszerności; odróżniamy go od
samego twierdzenia o znikaniu kohomologii).
Na mocy twierdzenia Bertiniego, nad ciałem
algebraicznie domkniętym charakterystyki zero
ogólny element każdego z tych liniowych układów
jest gładką całkowitą krzywą, oznaczmy je
$C_0\sim mH$, $C_1\sim D+mH$.
Stąd $D\sim C_1-C_0$.

Wielkość $Q$ zależy tylko od klasy liniowej,
ponieważ zarówno $\mathcal O_S(D)$, jak i liczby
przecięć nie zmieniają się przy takiej równoważności.
Dzięki dwukrotnemu zastosowaniu równości
$Q(D\pm C)=Q(D)$ mamy
$Q(D)=Q(C_1-C_0)=Q(0)=0$.
Otrzymujemy żądany wzór. Wykorzystane rezultaty
zewnętrzne to Riemann--Roch na krzywej,
Bertini i dostateczna obszerność $D+mH$;
nie używaliśmy Hirzebrucha--Riemanna--Rocha.
\end{proof}

\begin{wniosek}[Parzystość wyrażenia przecięciowego]
\label{wn:ga7-parity}
Niech $S$ będzie gładką rzutową powierzchnią
nad algebraicznie domkniętym ciałem charakterystyki zero,
a $K_S$ jej dowolnym dzielnikiem kanonicznym.
Dla każdego dzielnika $D$ na $S$ liczba
$D\cdot(D-K_S)$ jest parzysta.
\end{wniosek}
\begin{proof}
Z twierdzenia~\ref{tw:ga7-rr-surfaces} wynika,
że połowa tej liczby jest różnicą dwóch
charakterystyk Eulera, a więc liczbą całkowitą.
\end{proof}

\begin{przyklad}[Riemann--Roch na $\PP^2$ dla dowolnego $d\in\Z$]
\label{prz:ga7-rr-p2}
Na $S=\PP^2$ mamy $K=-3H$, $H^2=1$,
$\chi(\mathcal O_S)=1$. Dla $D=dH$ otrzymujemy
$\chi(\mathcal O_{\PP^2}(d))=1+\frac12d(d+3)
=\frac{(d+1)(d+2)}2$.
W szczególności dla $d\ge0$ jest to
$\binom{d+2}{2}$, liczba jednorodnych jednomianów
stopnia $d$ w trzech zmiennych.
Dla $d<0$ nie można utożsamiać tej charakterystyki
z wymiarem $H^0$: na przykład
$\chi(\mathcal O(-3))=1$, choć $H^0(\mathcal O(-3))=0$.
Brakująca jedynka pochodzi z $H^2(\mathcal O(-3))$,
co wynika z dualności Serre'a.
\end{przyklad}

\begin{przyklad}[Riemann--Roch na $\PP^1\times\PP^1$]
Niech $S=\PP^1\times\PP^1$ i
$D=aF_1+bF_2$. Wówczas $D^2=2ab$, a
$K=-2F_1-2F_2$, zatem
$D\cdot(D-K)=2ab+2a+2b$.
Ponieważ $\chi(\mathcal O_S)=1$,
otrzymujemy
$\chi(\mathcal O_S(a,b))=(a+1)(b+1)$
dla \emph{wszystkich} liczb całkowitych $a,b$.
Dla $a,b\ge0$ wynika to także bezpośrednio z liczby
bihomogenicznych jednomianów bitypu $(a,b)$.
\end{przyklad}

\section{Hirzebruch--Riemann--Roch i wzór Noether}
\label{sec:ga7-noether}

Rozszerzenie Riemanna--Rocha z samych dzielników
na dowolne wiązki wektorowe wymaga \emph{charakteru Cherna}
i \emph{klasy Todda}. Te obiekty wygodnie definiować
w pierścieniu $\operatorname{CH}^*(X)\otimes_\Z\Q$:
niektóre współczynniki są ułamkowe, lecz po wzięciu
stopnia otrzymujemy liczbę całkowitą.

\begin{definicja}[Charakter Cherna i klasa Todda]
\label{def:ga7-ch-td}
Niech $X$ będzie gładką quasi-rzutową rozmaitością,
$\mathcal E$ wiązką wektorową rzędu $r$.
Po zastosowaniu zasady rozszczepienia niech
$x_1,\ldots,x_r$ oznaczają \emph{formalne pierwiastki Cherna},
tj. pierwsze klasy Cherna kolejnych ilorazów liniowych
w odpowiedniej filtracji pullbacku $\mathcal E$.
Definiujemy
\[
\operatorname{ch}(\mathcal E)=\sum_{i=1}^r e^{x_i},
\qquad
\operatorname{td}(\mathcal E)=\prod_{i=1}^r
\frac{x_i}{1-e^{-x_i}},
\]
przy czym szeregi rozwijamy formalnie i obcinamy
powyżej wymiaru $X$. Dla $x_i=0$ drugi czynnik
rozumiemy jako szereg o wyrazie stałym $1$.
Ponieważ są to wyrażenia symetryczne w $x_i$,
można je zapisać przez klasy Cherna, a więc
zależność od wyboru rozszczepienia znika.
\end{definicja}

\begin{lemat}[Rozwinięcia do stopnia dwa]
\label{lem:ga7-ch-td-surface}
Niech $X$ będzie gładką powierzchnią,
$\mathcal E$ wiązką wektorową rzędu $r$,
a $T_X$ jej wiązką styczną. Wtedy
\[
\begin{aligned}
\operatorname{ch}(\mathcal E)&=r+c_1(\mathcal E)
+\tfrac12(c_1(\mathcal E)^2-2c_2(\mathcal E)),\\
\operatorname{td}(T_X)&=1+\tfrac12 c_1(T_X)
+\tfrac1{12}(c_1(T_X)^2+c_2(T_X)).
\end{aligned}
\]
\end{lemat}
\begin{proof}
Z rozwinięcia $e^x=1+x+x^2/2+\cdots$ otrzymujemy
$\sum e^{x_i}=r+\sum x_i+\frac12\sum x_i^2$.
Podstawiamy
$c_1=\sum x_i$, $c_2=\sum_{i<j}x_ix_j$ i
$\sum x_i^2=c_1^2-2c_2$.
Ponadto zwykły rachunek formalnych szeregów daje
$x/(1-e^{-x})=1+x/2+x^2/12+\cdots$.
Na powierzchni $T_X$ ma rząd dwa, więc mnożąc
czynniki dla jego dwóch pierwiastków
$x_1,x_2$, otrzymujemy wyraz stopnia dwa
$(x_1^2+x_2^2)/12+x_1x_2/4
=(c_1^2-2c_2)/12+c_2/4=(c_1^2+c_2)/12$.
To dowodzi obu wzorów.
\end{proof}

\begin{twierdzenie}[Hirzebruch--Riemann--Roch — twierdzenie zewnętrzne]
\label{tw:ga7-hrr}
Niech $X$ będzie gładką właściwą rozmaitością nad
ciałem $k$ charakterystyki zero, a $\mathcal E$
wiązką wektorową na $X$. Wtedy
\[
\chi(X,\mathcal E)=\deg\bigl(
\operatorname{ch}(\mathcal E)\operatorname{td}(T_X)
\cap[X]\bigr)_{0},
\]
gdzie indeks $0$ oznacza składową cyklu wymiaru zero.
\end{twierdzenie}

\begin{uwaga}[Dlaczego nie przedstawiamy tego jako prostego wniosku?]
Pełny dowód tego twierdzenia wykorzystuje
teorię $K$-grup, deformację do stożka normalnego
albo teorię indeksową. Jest istotnie
głębszy niż dowód wzoru Riemanna--Rocha dla dzielników
na powierzchni, który podaliśmy wyżej.
Przyjmujemy go jawnie jako \emph{wynik zewnętrzny},
a wszystkie poniższe wyprowadzenia są już rachunkowe.
\end{uwaga}

\begin{twierdzenie}[Wzór Noether i uogólniony Riemann--Roch dla powierzchni]
\label{tw:ga7-noether}
Niech $S$ będzie gładką spójną rzutową powierzchnią
nad ciałem algebraicznie domkniętym charakterystyki zero.
Wówczas
\[
12\chi(\mathcal O_S)=K_S^2+\deg c_2(T_S).
\]
Dla wiązki wektorowej $\mathcal E$ rzędu $r$ na $S$
zachodzi ponadto
\[
\chi(\mathcal E)=r\chi(\mathcal O_S)
+\frac12\deg\bigl(c_1(\mathcal E)
(c_1(\mathcal E)-c_1(\omega_S))\bigr)
-\deg c_2(\mathcal E).
\]
\end{twierdzenie}
\begin{proof}
Zastosujmy twierdzenie~\ref{tw:ga7-hrr} do
$\mathcal E=\mathcal O_S$. Wtedy
$\operatorname{ch}(\mathcal O_S)=1$,
a składową stopnia dwa klasy Todda jest
$(c_1(T_S)^2+c_2(T_S))/12$.
Ponieważ $c_1(T_S)=-c_1(\omega_S)$, jej kwadrat
ma stopień $K_S^2$. Otrzymujemy wzór Noether.

Dla ogólnej $\mathcal E$ mnożymy rozwinięcia
z lematu~\ref{lem:ga7-ch-td-surface}. Składowa
stopnia dwa jest równa
\[
\frac r{12}(c_1(T_S)^2+c_2(T_S))
+\frac12c_1(\mathcal E)c_1(T_S)
+\frac12(c_1(\mathcal E)^2-2c_2(\mathcal E)).
\]
Po wzięciu stopnia pierwszy składnik jest równy
$r\chi(\mathcal O_S)$ na mocy dopiero co
wyprowadzonego wzoru, a
$c_1(T_S)=-c_1(\omega_S)$. Pozostałe wyrazy
są iloczynami klas w $\operatorname{CH}^*(S)$ i dają
drugą formułę, ponieważ $c_1(\omega_S)=[K_S]$.
Dla dzielników używamy więc liczby przecięcia,
a dla wiązek ich pierwszych klas Cherna.
Dla wiązki liniowej $\mathcal O_S(D)$ ma ona
$r=1$, $c_1=D$ i $c_2=0$, czyli dokładnie
wcześniej udowodniony wzór powierzchniowy.
\end{proof}

\begin{przyklad}[Charakterystyka holomorficzna powierzchni stopnia $d$]
Niech $S\subseteq\PP^3$ będzie gładką powierzchnią
stopnia $d$. Ze wzorów dla hiperpowierzchni oraz
wzoru Noether otrzymujemy
\[
\chi(\mathcal O_S)=\frac{d(d-4)^2+d(d^2-4d+6)}{12}
=\frac{d(d^2-6d+11)}6.
\]
Dla $d=1,2,3,4$ daje to odpowiednio $1,1,1,2$.
Z drugiej strony dla kwartyki wiemy z rozdziału 6,
że $q=0$, $p_g=1$, więc $\chi=1-q+p_g=2$.
Otrzymaliśmy niezależne sprawdzenie liczby $24$
we wzorze $12\chi=K^2+c_2$.
\end{przyklad}

\begin{przyklad}[Jak zmienia się $c_2$ po rozdmuchaniu?]
\label{prz:ga7-blowup-c2}
Niech $k$ będzie ciałem algebraicznie domkniętym
charakterystyki zero, $S$ gładką rzutową powierzchnią
nad $k$, a $p\in S$ punktem zamkniętym.
Niech $\pi:\widetilde S=\operatorname{Bl}_pS\to S$
będzie jej rozdmuchaniem w $p$.
W rozdziale 6 udowodniliśmy
$K_{\widetilde S}^2=K_S^2-1$.
Standardowy wynik o kohomologii rozdmuchania mówi, że
$\pi_*\mathcal O_{\widetilde S}=\mathcal O_S$ oraz
$R^i\pi_*\mathcal O_{\widetilde S}=0$ dla $i>0$.
Z ciągu spektralnego Leraya wynika wtedy
$\chi(\mathcal O_{\widetilde S})=\chi(\mathcal O_S)$.
Wzór Noether dla obu powierzchni daje
$\deg c_2(T_{\widetilde S})=\deg c_2(T_S)+1$.
Dla rozdmuchania $\PP^2$ w sześciu punktach otrzymujemy
$K^2=9-6=3$ i $\deg c_2=3+6=9$, zgodnie
z niezależnym rachunkiem dla gładkiej powierzchni kubicznej.
W tym argumencie znikanie wyższych obrazów prostych
jest jawną przesłanką zewnętrzną, a nie pominiętym krokiem.
\end{przyklad}

\section{Podsumowanie: dwa poziomy teorii przecięć}
\label{sec:ga7-summary}

W przypadku powierzchni wiele wzorów liczbowych można
wyprowadzić z normalnej wiązki krzywej, adjunkcji
i Riemanna--Rocha na krzywych. To jest \emph{poziom
konkretny}: pozwolił udowodnić Riemanna--Rocha dla
wszystkich dzielników na gładkiej powierzchni.

W wyższych wymiarach konieczne są pełne grupy Chowa,
klasy Cherna, wzory projekcji, przecięcia Gysina
i twierdzenia o przesuwaniu cykli. To jest
\emph{poziom funktorialny}. Zbudowaliśmy jego definicje,
przeprowadziliśmy liczne rachunki i wskazaliśmy,
które twierdzenia konstrukcyjne nadal są przyjęte
z zewnątrz. W ten sposób teoria z rozdziału 6
otrzymała precyzyjne podstawy, a wzór Noether
stał się konkretną konsekwencją
Hirzebrucha--Riemanna--Rocha.

\begin{przyklad}[Wielomian Hilberta powierzchni jako kontrola wzoru Riemanna--Rocha]
\label{prz:ga7-hilbert-polynomial}
Niech $S\subseteq\PP^3_k$ będzie gładką powierzchnią
stopnia $d\ge1$ nad ciałem algebraicznie domkniętym
charakterystyki zero, a $H$ dzielnikiem hiperpłaszczyznowym
na $S$. Oznaczmy $h=[H]=c_1(\mathcal O_S(1))$.
Dla dowolnego $m\in\Z$ wzór Riemanna--Rocha
i równości $K_S\sim(d-4)H$, $\deg(h^2)=d$,
gdzie $H$ jest dzielnikiem hiperpłaszczyznowym
oraz $h=[H]\in\operatorname{CH}^1(S)$, dają
\[
\chi(\mathcal O_S(m))
=\frac{d(d^2-6d+11)}6
 +\frac d2\,m(m+4-d).
\]
Rachunek ten można niezależnie sprawdzić za pomocą
dokładnego ciągu snopów na $\PP^3$:
\[
0\longrightarrow\mathcal O_{\PP^3}(m-d)
\longrightarrow\mathcal O_{\PP^3}(m)
\longrightarrow\mathcal O_S(m)\longrightarrow0.
\]
Ponieważ
$\chi(\mathcal O_{\PP^3}(t))=\binom{t+3}{3}$
dla każdego $t\in\Z$, z addytywności charakterystyki
Eulera otrzymujemy również
\[
\chi(\mathcal O_S(m))
=\binom{m+3}{3}-\binom{m-d+3}{3}.
\]
W tym wzorze symbol Newtona oznacza wielomian
$\binom{t}{3}=t(t-1)(t-2)/6$, także dla ujemnego
$t$. Rozwinięcie obu wyrażeń daje tę samą funkcję
kwadratową zmiennej $m$. Dla $m=1$ gładka
powierzchnia kubiczna ($d=3$) oraz gładka kwartyka
($d=4$) mają obie $\chi(\mathcal O_S(1))=4$,
choć ich charakterystyki strukturalne wynoszą
odpowiednio $1$ i $2$. Pojedyncza wartość
wielomianu Hilberta nie rozstrzyga więc o typie
biwymiernym powierzchni.
\end{przyklad}

\begin{uwaga}[Przygotowanie do geometrii biwymiernej]
W rozdziale 8 liczby $C^2$, $K_S\cdot C$ i krzywe
z ujemnym samoprzecięciem będą kontrolować możliwe
kontrakcje. Program modeli minimalnych w wymiarze dwa
korzysta szczególnie z tego, że przecięcia z dzielnikiem
kanonicznym rozpoznają krzywe wyjątkowe. Teoria
klas Cherna dostarczy z kolei niezmienników powierzchni,
które pozwalają odróżniać typy biwymierne.
\end{uwaga}

% Źródła matematyczne użyte do weryfikacji:
% Stacks Project: Chow Homology and Chern Classes (tagi 02U4, 02TZ, 02UF),
% Grothendieck-Riemann-Roch (odpowiedni rozdział Stacks Project).
% Fulton, Intersection Theory; Hartshorne, Algebraic Geometry V.
% Wykorzystane zewnętrznie: moving lemma, lokalizacja i homotopia CH,
% projective bundle formula, Whitney/naturalność, sekcje regularne,
% HRR, Bertini i obszerność, wyższe obrazy proste rozdmuchania.

% =============================================================
% ROZDZIAL 8: GEOMETRIA BIWYMIERNA I MODELE MINIMALNE
% =============================================================

\chapter{Geometria biwymierna i modele minimalne powierzchni}
\label{ch:ga8-biwymierna}

\section{Dlaczego izomorfizm bywa zbyt silnym pojęciem?}
\label{sec:ga8-motivation}

Przy badaniu krzywych i powierzchni wielokrotnie zastępowaliśmy
punkt osobliwy nowymi punktami lub kierunkami. Rozdmuchanie
$\pi:\widetilde S\to S$ gładkiej powierzchni w punkcie
nie jest izomorfizmem: włókno nad środkiem ma postać $\PP^1$.
Mimo to poza środkiem $\pi$ jest izomorfizmem. Taka zmiana
jest \emph{biwymierna}. W tym rozdziale badamy jej najbardziej
konkretne konsekwencje: odwzorowania wymierne, przekształcenie
Cremony, rozdmuchania i kontrakcje krzywych.
W końcowej części rozdziału przejdziemy od tych
konkretnych przykładów do dwuwymiarowego programu
modeli minimalnych. Nie będziemy rozwijać jego
wyżej wymiarowej wersji.

Przez cały rozdział $k$ jest algebraicznie domkniętym ciałem
charakterystyki zero. Powierzchnie są gładkie, całkowite i
rzutowe nad $k$, chyba że w przykładzie zaznaczono inaczej.

\begin{definicja}[Odwzorowanie wymierne i biwymierne]
\label{def:ga8-rational}
Niech $X,Y$ będą całkowitymi rozmaitościami algebraicznymi
nad $k$. \emph{Odwzorowaniem wymiernym} $\varphi:X\dashrightarrow Y$
nazywamy klasę par $(U,\varphi_U)$, w której $U\subseteq X$
jest niepustym otwartym podzbiorem Zariskiego, a
$\varphi_U:U\to Y$ jest morfizmem; dwie pary są równoważne,
gdy ich morfizmy zgadzają się na pewnym niepustym
otwartym podzbiorze wspólnej dziedziny. Odwzorowanie jest
\emph{dominujące}, gdy jego obraz jest gęsty w $Y$.
Nazywamy je \emph{biwymiernym}, jeżeli istnieją niepuste
otwarte $U\subseteq X$, $V\subseteq Y$ takie, że
$\varphi_U:U\to V$ jest izomorfizmem.
\end{definicja}

\begin{uwaga}[Dlaczego wystarcza gęste otoczenie?]
Na rozmaitości całkowitej każdy niepusty otwarty
podzbiór jest gęsty. Odwzorowanie wymierne ignoruje
więc podzbiór, na którym nie potrafimy podać morfizmu,
ale pamięta działanie na części gęstej. Aby mówić o
\emph{izomorfizmie}, trzeba natomiast opisać wszystkie punkty.
Domena, na której dana reprezentacja jest morfizmem,
nie musi być największą domeną określoności.
\end{uwaga}

\begin{twierdzenie}[Kryterium przez ciała funkcji]
\label{tw:ga8-function-fields}
Niech $X,Y$ będą całkowitymi rozmaitościami nad $k$.
Każde dominujące odwzorowanie wymierne
$\varphi:X\dashrightarrow Y$ wyznacza zanurzenie ciał
$\varphi^*:k(Y)\hookrightarrow k(X)$.
Odwzorowanie $\varphi$ jest biwymierne wtedy i tylko
wtedy, gdy $\varphi^*$ jest izomorfizmem ciał nad $k$.
\end{twierdzenie}
\begin{proof}
Dominujący morfizm na niepustym otwartym podzbiorze
przenosi funkcję wymierną na jej złożenie z $\varphi$.
Jeśli funkcja znika po złożeniu, jest zerowa na gęstym
obrazie, zatem jest zerowa w $k(Y)$. To daje injektywność.
Jeśli $\varphi$ jest izomorfizmem otwartych podzbiorów,
wyznacza odwrotne izomorfizmy ich ciał funkcji, które
są dokładnie $k(X)$ i $k(Y)$.
W drugą stronę wybierzmy afiniczne niepuste otwarte
$U\subseteq X$ i $V\subseteq Y$. Generatory algebry
współrzędnych $k[V]$ jako elementy $k(X)$ są ilorazami
funkcji regularnych na $U$; na pewnym mniejszym
niepustym otwartym $U'$ ich mianowniki są niezerowe.
To definiuje morfizm wymierny $U'\to Y$.
W taki sam sposób izomorfizm odwrotny ciał funkcji
definiuje odwzorowanie wymierne $Y\dashrightarrow X$.
Złożenia tych odwzorowań mają na poziomie ciał funkcji
identyczność. Ponieważ $X$ i $Y$ są rozdzielone,
odwzorowania te zgadzają się z identycznościami na
odpowiednich niepustych otwartych podzbiorach.
Po ich zmniejszeniu uzyskujemy wzajemnie odwrotne
izomorfizmy otwartych podzbiorów.
\end{proof}

\begin{przyklad}[Rozdmuchanie nie zmienia ciała funkcji]
Dla $\pi:\operatorname{Bl}_0\A^2\to\A^2$ opisanej
w rozdziale~\ref{ch:ga5-rozdmuchania} mamy izomorfizm
po usunięciu punktu $0$ i dzielnika wyjątkowego.
Zatem $k(\operatorname{Bl}_0\A^2)=k(x,y)$.
Obecność całej prostej rzutowej nad jednym punktem
nie zmienia ciała funkcji wymiernych powierzchni.
\end{przyklad}

\section{Kwadratowe przekształcenie Cremony: obliczenie od początku}
\label{sec:ga8-cremona}

\begin{definicja}[Standardowe przekształcenie Cremony]
\label{def:ga8-cremona}
Odwzorowanie wymierne
\[
\sigma:\PP^2_k\dashrightarrow\PP^2_k,\qquad
[x:y:z]\longmapsto[yz:xz:xy]
\]
nazywamy standardowym kwadratowym przekształceniem Cremony.
\end{definicja}

\begin{stwierdzenie}[Punkty bazowe i odwrotność]
\label{st:ga8-cremona-basic}
Niech $k$ będzie ciałem, a
$\sigma:\PP^2_k\dashrightarrow\PP^2_k$ odwzorowaniem wymiernym
$[x:y:z]\mapsto[yz:xz:xy]$.
Odwzorowanie $\sigma$ jest nieokreślone dokładnie w punktach
$p_1=[1:0:0]$, $p_2=[0:1:0]$ i $p_3=[0:0:1]$.
Na otwartym zbiorze $T=\{xyz\ne0\}\subseteq\PP^2$
wyznacza izomorfizm $T\to T$, dla którego
$\sigma^{-1}=\sigma$. Na każdej prostej współrzędnych
poza punktami bazowymi jest stałe i przyjmuje
odpowiedni punkt współrzędnych.
\end{stwierdzenie}
\begin{proof}
Wszystkie trzy współrzędne $yz,xz,xy$ są zerowe
wtedy i tylko wtedy, gdy co najmniej dwie spośród
$x,y,z$ są równe zeru. Są to dokładnie trzy wskazane
punkty. Jeśli $xyz\ne0$, to
\[
\sigma^2([x:y:z])
=[(xz)(xy):(yz)(xy):(yz)(xz)]
=[x^2yz:xy^2z:xyz^2]=[x:y:z].
\]
Na prostej $x=0$ przy $yz\ne0$ obraz to $[1:0:0]$.
Dla pozostałych dwóch prostych otrzymujemy analogicznie
punkty $p_2,p_3$. Ostatni rachunek pokazuje również,
że punkty leżące na tych prostych (poza bazowymi)
nie mogą należeć do otwartego zbioru, na którym
$\sigma$ jest izomorfizmem na obraz otwarty.
\end{proof}

\begin{uwaga}[Po co trzy rozdmuchania?]
Nie da się rozszerzyć wzoru $[yz:xz:xy]$ do morfizmu
całej $\PP^2$ tylko przez \emph{nadanie} mu wartości
w trzech punktach bazowych: różne kierunki dojścia
do danego punktu prowadzą do różnych granicznych
wartości w przestrzeni rzutowej. Rozdmuchanie koduje
te kierunki jako punkty dzielnika wyjątkowego.
\end{uwaga}

\begin{stwierdzenie}[Rozwiązanie punktu bazowego w chartach]
\label{st:ga8-local-resolution}
Niech $p_1=[1:0:0]$. Na charcie $x\ne0$ z
współrzędnymi $(u=y/x,v=z/x)$ przekształcenie Cremony
ma postać $[uv:v:u]$. Po rozdmuchaniu $(u,v)=(0,0)$
na obu chartach jest morfizmem.
\end{stwierdzenie}
\begin{proof}
W pierwszym charcie rozdmuchania kładziemy $v=tu$.
Dla $u\ne0$ mamy
$[uv:v:u]=[tu^2:tu:u]=[tu:t:1]$.
Prawa strona ma sens także dla $u=0$ i opisuje
na dzielniku wyjątkowym punkty $[0:t:1]$.
W drugim charcie kładziemy $u=sv$. Wtedy
$[uv:v:u]=[sv^2:v:sv]=[sv:1:s]$,
co również rozszerza się na $v=0$.
Na przecięciu chartów (gdzie $s=1/t$)
oba wzory reprezentują ten sam punkt projektowy.
Po rozdmuchaniu $p_2$ i $p_3$ wykonujemy identyczny
rachunek z permutacją współrzędnych.
\end{proof}

\begin{center}
\begin{tikzpicture}[>=Stealth,scale=1.0]
 \coordinate (A) at (0,2.2); \coordinate (B) at (-2.45,-1.25);
 \coordinate (C) at (2.45,-1.25);
 \draw[thick] (A)--(B)--(C)--cycle;
 \foreach \P in {A,B,C}{\fill[blue!70!black] (\P) circle (2.6pt);}
 \node[above=5pt,font=\footnotesize] at (A) {$p_3$};
 \node[below left=3pt,font=\footnotesize] at (B) {$p_1$};
 \node[below right=3pt,font=\footnotesize] at (C) {$p_2$};
 \node[align=center,font=\footnotesize,text width=5cm] at (0,-2.15)
 {Trzy punkty bazowe i trzy\\ proste współrzędnych};
\end{tikzpicture}

{\small Rysunek 8.1. Schemat położenia punktów bazowych
standardowego przekształcenia Cremony. Krawędzie trójkąta
oznaczają trzy proste kontraktowane przez odwzorowanie.}\par
\end{center}

\section{Forma przecięć na rozdmuchaniu trzech punktów}
\label{sec:ga8-classes}

Niech $\pi:X\to\PP^2$ będzie rozdmuchaniem
trzech różnych niewspółliniowych punktów $p_1,p_2,p_3$.
Oznaczmy przez $H$ klasę przeciwobrazu prostej,
a przez $E_1,E_2,E_3$ klasy dzielników wyjątkowych.
W rozdziale~\ref{ch:ga6-powierzchnie} wyprowadziliśmy
\[
H^2=1,\qquad H\cdot E_i=0,\qquad
E_i^2=-1,\qquad E_i\cdot E_j=0\quad(i\ne j).
\]
Kanonika ma klasę $K_X=-3H+E_1+E_2+E_3$.

\begin{stwierdzenie}[Transformaty prostych przez pary punktów]
\label{st:ga8-hexagon}
Dla $1\le i<j\le3$ oznaczmy przez $L_{ij}$
transformatę ścisłą prostej łączącej $p_i,p_j$.
Wówczas $[L_{ij}]=H-E_i-E_j$, a
$L_{ij}^2=-1$ oraz $K_X\cdot L_{ij}=-1$.
Ponadto $L_{ij}\cdot E_i=L_{ij}\cdot E_j=1$,
$L_{ij}\cdot E_k=0$ dla $k\notin\{i,j\}$,
oraz pary różnych $L_{ij}$ są rozłączne.
\end{stwierdzenie}
\begin{proof}
Własność transformacji ścisłej z rozdziału 6 daje
$[L_{ij}]=H-E_i-E_j$, ponieważ prosta przechodzi
przez $p_i,p_j$ z krotnością jeden i przez trzeci punkt
nie przechodzi. Formę przecięć obliczamy algebraicznie:
$L_{ij}^2=1-1-1=-1$,
$K_X\cdot L_{ij}=-3+1+1=-1$.
Tak samo $L_{ij}\cdot E_i= -E_i^2=1$.
Dla różnych par, na przykład $(i,j)=(1,2)$ i
$(i,j)=(1,3)$, dostajemy
$(H-E_1-E_2)\cdot(H-E_1-E_3)=1+E_1^2=0$.
Obie krzywe są całkowite, różne i efektywne, więc
nieujemne lokalne krotności przecięć wymuszają
rozłączność.
\end{proof}

\begin{center}
\begin{tikzpicture}[>=Stealth,scale=0.95]
\foreach \i/\a/\txt in {
0/90/$E_1$,1/30/$L_{12}$,2/-30/$E_2$,
3/-90/$L_{23}$,4/-150/$E_3$,5/-210/$L_{13}$}{
 \coordinate (P\i) at (\a:2.05cm);
 \fill[blue!65!black] (P\i) circle (2.4pt);
 \node[font=\footnotesize] at (\a:2.72cm) {\txt};
}
\draw[thick] (P0)--(P1)--(P2)--(P3)--(P4)--(P5)--cycle;
\node[align=center,font=\footnotesize] at (0,0)
 {przecięcia\\równe $1$\\na krawędziach};
\end{tikzpicture}

{\small Rysunek 8.2. Graf przecięć sześciu krzywych
samoprzecięcia $-1$ na rozdmuchaniu trzech niewspółliniowych
punktów $\PP^2$. Krawędzie odpowiadają iloczynowi $1$,
a pozostałe pary różnych krzywych mają iloczyn $0$.}\par
\end{center}

\begin{stwierdzenie}[Układ liniowy rozwiązujący Cremonę]
\label{st:ga8-cremona-linear}
Niech $X$ będzie jak wyżej, a $H'$ klasą przeciwobrazu
prostej z docelowej $\PP^2$ przez morfizm
$\widetilde\sigma:X\to\PP^2$ uzyskany po rozwiązaniu
punktów bazowych. Wtedy
$H'=2H-E_1-E_2-E_3$, a $H'^2=1$.
Morfizm $\widetilde\sigma$ kontraktuje każdą z krzywych
$L_{12},L_{13},L_{23}$ i odwzorowuje każdą $E_i$
izomorficznie na prostą w docelowej płaszczyźnie.
\end{stwierdzenie}
\begin{proof}
Współrzędne $yz,xz,xy$ są formami jednorodnymi
stopnia dwa, każda znika w trzech punktach bazowych
co najmniej do rzędu jeden. Ich transformaty po
rozdmuchaniu wyznaczają układ $|2H-E_1-E_2-E_3|$.
Rachunek w chartach z dowodu
stwierdzenia~\ref{st:ga8-local-resolution} pokazuje,
że transformaty trzech przekrojów nie zerują się
jednocześnie, zatem układ jest bez punktów bazowych
i definiuje $\widetilde\sigma$.
Stąd $H'=2H-E_1-E_2-E_3$ i $H'^2=4-3=1$.
Dla przykładu $H'\cdot L_{12}=2-1-1=0$,
podobnie dla pozostałych. Możemy też sprawdzić
kontrakcję bezpośrednio: prosta $z=0$ poza punktami
bazowymi przechodzi przez $\sigma$ w punkt $p_3$.
Jej transformata ścisła jest więc kontraktowana.
Na $E_1$ w pierwszym charcie mapa ma postać
$t\mapsto[0:t:1]$, a w drugim $s\mapsto[0:1:s]$;
razem tworzą parametryzację docelowej prostej $x'=0$.
Analogicznie zachowują się $E_2,E_3$.
\end{proof}

\section{Krzywe wyjątkowe i odwracanie rozdmuchania}
\label{sec:ga8-contraction}

\begin{definicja}[Krzywa $(-1)$]
\label{def:ga8-minus-one}
Niech $S$ będzie gładką powierzchnią rzutową nad $k$.
Krzywą $(-1)$ nazywamy gładką całkowitą krzywą
$E\subseteq S$, która jest izomorficzna z $\PP^1_k$
i spełnia $E^2=-1$.
\end{definicja}

\begin{stwierdzenie}[Dzielnik wyjątkowy rozdmuchania]
\label{st:ga8-exceptional-curve}
Jeżeli $\pi:\widetilde S\to S$ jest rozdmuchaniem
gładkiej powierzchni rzutowej w punkcie $p$, to
$E=\pi^{-1}(p)$ jest krzywą $(-1)$.
Ponadto $K_{\widetilde S}\cdot E=-1$.
\end{stwierdzenie}
\begin{proof}
Z konstrukcji rozdmuchania $E\cong\PP^1$, a z
rozdziału~\ref{ch:ga6-powierzchnie} $E^2=-1$.
Korzystając ze wzoru $K_{\widetilde S}=\pi^*K_S+E$
oraz z równości $\pi^*K_S\cdot E=0$, otrzymujemy
$K_{\widetilde S}\cdot E=E^2=-1$.
\end{proof}

\begin{twierdzenie}[Kryterium Castelnuova — wynik zewnętrzny]
\label{tw:ga8-castelnuovo}
Niech $S$ będzie gładką rzutową powierzchnią nad
algebraicznie domkniętym ciałem charakterystyki zero,
a $E\subseteq S$ krzywą $(-1)$. Wtedy istnieje
gładka rzutowa powierzchnia $S'$ oraz morfizm
$\rho:S\to S'$, który jest izomorfizmem poza $E$,
kontraktuje $E$ do jednego gładkiego punktu i jest
izomorficzny z rozdmuchaniem $S'$ w tym punkcie.
\end{twierdzenie}

\begin{uwaga}[Dlaczego nie każda ujemna krzywa znika gładko?]
Krzywa o samoprzecięciu $-2$ nie może być dzielnikiem
wyjątkowym rozdmuchania \emph{jednego gładkiego punktu},
gdyż taki dzielnik ma samoprzecięcie $-1$.
Krzywe $(-2)$ mogą być kontraktowane do \emph{punktów osobliwych}.
Przykładem jest ujemna sekcja powierzchni Hirzebrucha
$\mathbb F_2$, którą można kontraktować do wierzchołka
stożka kwadratowego. Jest to inny rodzaj kontrakcji,
którego nie należy nazywać odwrotnością zwykłego
rozdmuchania gładkiego punktu.
\end{uwaga}

\begin{przyklad}[Dwie odwrotne operacje]
Zacznijmy od $\PP^2$ i rozdmuchajmy punkt $p$.
Otrzymujemy $\mathbb F_1$ z wyjątkową sekcją
$E\cong\PP^1$, $E^2=-1$. Kontrakcja tej sekcji
odtwarza $\PP^2$. Nie wszystkie krzywe $\mathbb F_1$
mają samoprzecięcie $-1$: włókno rzutowania
$\mathbb F_1\to\PP^1$ ma samoprzecięcie $0$.
Ten przykład odróżnia kontrakcję od zwykłego
„zmazywania dowolnej krzywej”.
\end{przyklad}

\section{Niezmienniki biwymierne i liczby zmieniające się po rozdmuchaniu}
\label{sec:ga8-invariants}

\begin{stwierdzenie}[Zmiana liczb przecięć przy jednym rozdmuchaniu]
\label{st:ga8-classes-blowup}
Niech $\pi:\widetilde S\to S$ będzie rozdmuchaniem
w punkcie gładkiej powierzchni rzutowej. Wtedy
$K_{\widetilde S}^2=K_S^2-1$, a w charakterystyce zero
$\chi(\mathcal O_{\widetilde S})=\chi(\mathcal O_S)$.
W szczególności $\deg c_2(T_{\widetilde S})=
\deg c_2(T_S)+1$.
\end{stwierdzenie}
\begin{proof}
Pierwsza równość wynika ze wzoru na kanonikę
$K_{\widetilde S}=\pi^*K_S+E$ oraz relacji
$\pi^*K_S\cdot E=0$, $E^2=-1$.
Niezmienniczość charakterystyki Eulera wynika z
$\pi_*\mathcal O_{\widetilde S}=\mathcal O_S$
i $R^j\pi_*\mathcal O_{\widetilde S}=0$ dla $j>0$;
jest to znany wynik o wyższych obrazach prostych przy
rozdmuchaniu gładkiego punktu, przyjęty w
przykładzie~\ref{prz:ga7-blowup-c2}.
Z wzoru Noether, udowodnionego w rozdziale 7
na podstawie Hirzebrucha--Riemanna--Rocha,
otrzymujemy $12\chi=K^2+\deg c_2$.
Po odjęciu wzorów dla obu powierzchni dostajemy
zmianę stopnia $c_2$ równą $+1$.
\end{proof}

\begin{przyklad}[Liczby po trzech rozdmuchaniach]
Dla $X=\operatorname{Bl}_{p_1,p_2,p_3}\PP^2$
mamy $K_X^2=9-3=6$ i $\deg c_2(T_X)=3+3=6$,
a $\chi(\mathcal O_X)=1$. Sześć krzywych $(-1)$
z poprzedniej sekcji tworzy heksagon przecięć;
liczby $K^2$ i $c_2$ nie pozwalają rozpoznać
ich indywidualnego położenia, ale dobrze
kontrolują kolejne rozdmuchania.
\end{przyklad}

\begin{twierdzenie}[Niezmienniczość $q$ i $p_g$ — wynik standardowy]
\label{tw:ga8-q-pg}
Jeżeli $S,S'$ są gładkimi spójnymi rzutowymi
powierzchniami nad $\C$ i są biwymierne, to
$h^1(S,\mathcal O_S)=h^1(S',\mathcal O_{S'})$ oraz
$h^0(S,K_S)=h^0(S',K_{S'})$.
\end{twierdzenie}
\begin{uwaga}[Skąd to wiadomo?]
Dla rozdmuchania pierwsza równość wynika z zanikania
wyższych obrazów prostych $\mathcal O$.
Druga zachodzi, ponieważ holomorficzne formy
najwyższego stopnia na gładkiej powierzchni
przechodzą jednoznacznie między powierzchnią i jej
rozdmuchaniem; można to sprawdzić w chartach, a
formalny dowód korzysta z przedłużania form.
Przejście od pojedynczego rozdmuchania do dowolnego
odwzorowania biwymiernego używa twierdzenia o
faktoryzacji morfizmów biwymiernych gładkich powierzchni
na rozdmuchania punktów. Nie przeprowadzamy tutaj
pełnego dowodu tego standardowego twierdzenia.
\end{uwaga}

\begin{przyklad}[Powierzchnie wymierne i K3]
Powierzchnie $\PP^2$ i $\mathbb F_1$ są biwymierne,
a mimo to nie są izomorficzne: $\mathbb F_1$ zawiera
krzywą o samoprzecięciu $-1$, natomiast $\PP^2$ nie.
Gładka kwartyka K3 w $\PP^3(\C)$ ma
$p_g=h^0(K)=1$, a $\PP^2$ ma $p_g=0$.
Zatem kwartyka K3 nie jest powierzchnią wymierną.
Wniosek nie wynika z samej różnicy $K^2$, ponieważ
ta liczba \emph{zmienia się} przy rozdmuchaniu.
\end{przyklad}

% =============================================================
% DWUWYMIAROWY PROGRAM MODELI MINIMALNYCH: UZUPELNIENIE
% =============================================================
\section{Równoważność numeryczna i stożek krzywych}
\label{sec:ga8-ne-cone}

\emph{Idea.} Odwzorowanie biwymierne może kontraktować krzywe,
ale nie dowolne. Na powierzchni ich istnienie rozpoznajemy,
porównując liczby przecięć z dzielnikiem kanonicznym.
Zamiast rozpatrywać wszystkie krzywe pojedynczo, grupujemy je
według zachowania względem dywizorów. Powstanie w ten sposób
stożek, którego szczególne kierunki wskazują możliwe kontrakcje.

\begin{definicja}[Równoważność numeryczna]
\label{def:ga8-numerical}
Niech $S$ będzie gładką rzutową powierzchnią nad algebraicznie
domkniętym ciałem $k$ charakterystyki zero. Dwa dywizory
$D,D'$ na $S$ są \emph{numerycznie równoważne}, co zapisujemy
$D\equiv D'$, gdy $(D-D')\cdot C=0$ dla każdej całkowitej
krzywej $C\subseteq S$. Dwa cykle jednowymiarowe $Z,Z'$
są numerycznie równoważne, gdy $D\cdot(Z-Z')=0$ dla
każdego dywizora $D$ na $S$. Oznaczamy przez
$N^1(S)_{\R}$ i $N_1(S)_{\R}$ odpowiednie przestrzenie
ilorazowe po rozciągnięciu skalarów do $\R$.
Ich wspólny wymiar $\rho(S)$ nazywamy \emph{liczbą Picarda}.
Przez liczbę przecięcia otrzymujemy parowanie
$N^1(S)_{\R}\times N_1(S)_{\R}\to\R$.
\end{definicja}

\begin{uwaga}[Co jest twierdzeniem zewnętrznym?]
Skończoność wymiaru obu przestrzeni jest konsekwencją
twierdzenia Nérona--Severiego. Z definicji ilorazów
wynika niezdegenerowanie parowania; utożsamienie obu
wymiarów w kontekście powierzchni można uzyskać także
z symetrii i niezdegenerowania formy przecięć.
Nie utożsamiamy równoważności numerycznej z liniową:
na krzywej eliptycznej dwa różne punkty mają ten sam
stopień, a więc są numerycznie równoważne jako dywizory,
lecz na ogół nie są liniowo równoważne.
\end{uwaga}

\begin{definicja}[Stożek krzywych i dywizor nef]
\label{def:ga8-nef}
Niech $S$ będzie jak wyżej. Przez $\operatorname{NE}(S)$
oznaczamy stożek w $N_1(S)_{\R}$ składający się ze
skończonych sum $\sum_i a_i[C_i]$, gdzie $a_i\ge0$,
a $C_i\subseteq S$ są całkowitymi krzywymi.
Jego domknięcie oznaczamy $\overline{\operatorname{NE}}(S)$.
Dywizor rzeczywisty $D$ na $S$ nazywamy \emph{nef}
(ang. \emph{numerically effective}), gdy $D\cdot C\ge0$
dla każdej całkowitej krzywej $C\subseteq S$.
\emph{Promieniem ekstremalnym} stożka wypukłego
$\Gamma\subseteq N_1(S)_{\R}$ nazywamy promień
$R=\R_{\ge0}v\subseteq\Gamma$, $v\ne0$, taki że
z $u,w\in\Gamma$ oraz $u+w\in R$ wynika $u,w\in R$.
\end{definicja}

\begin{przyklad}[Dlaczego nie wystarcza sprawdzić $D^2$?]
Na $S=\operatorname{Bl}_p\PP^2$ niech $H$ będzie
przeciwobrazem klasy prostej, a $E$ dzielnikiem wyjątkowym.
Wtedy $H^2=1$, $H\cdot E=0$, $E^2=-1$.
Dywizor $D=H+E$ ma $D^2=0$, ale $D\cdot E=-1$,
więc nie jest nef. Natomiast $H$ jest nef, ponieważ
$H=\pi^*\mathcal O_{\PP^2}(1)$ i dla każdej krzywej
$C$ mamy $H\cdot C\ge0$; dla $C=E$ liczba ta wynosi zero.
To pokazuje, dlaczego nieujemność samoprzecięcia
dywizora nie jest kryterium nefności.
\end{przyklad}

\begin{stwierdzenie}[Stożek krzywych na $\mathbb F_n$]
\label{st:ga8-fn-cone}
Niech $n\ge0$, a $\mathbb F_n$ będzie powierzchnią
Hirzebrucha z rozdziału~\ref{ch:ga6-powierzchnie}.
Oznaczmy przez $s$ jej standardową sekcję, a przez $f$
włókno projekcji do $\PP^1$, przy czym
$s^2=-n$, $s\cdot f=1$, $f^2=0$.
Wówczas $\overline{\operatorname{NE}}(\mathbb F_n)$
jest stożkiem generowanym przez $[s]$ i $[f]$.
Ponadto $K_{\mathbb F_n}=-2s-(n+2)f$,
a więc $K\cdot s=n-2$ oraz $K\cdot f=-2$.
\end{stwierdzenie}
\begin{proof}
Używamy znanego opisu $\operatorname{Pic}(\mathbb F_n)
\cong\Z[s]\oplus\Z[f]$ (wynik z konstrukcji
wiązki rzutowej). Niech $C\subseteq\mathbb F_n$
będzie całkowitą krzywą różną od $s$ i od włókna.
W klasie numerycznej piszemy $[C]=a[s]+b[f]$.
Dla przecięcia z ogólnym włóknem otrzymujemy
$a=C\cdot f\ge0$. Ponieważ $C$ i $s$ nie mają
wspólnej składowej, $C\cdot s=b-na\ge0$,
a zatem $b\ge na\ge0$. Dla $n=0$ ten sam argument
daje $a,b\ge0$. Włókno i sekcja są skutecznymi
krzywymi, więc każdy generator należy do stożka.
Stąd całość stożka jest dokładnie ich nieujemną
powłoką, która jest już domknięta.
Wzór kanoniczny wynika z adjunkcji dla włókna
($K\cdot f=-2$) i dla sekcji $s\simeq\PP^1$:
$(K+s)\cdot s=-2$. Gdy $K=as+bf$, pierwsza
równość daje $a=-2$, a druga $(-2s+bf+s)\cdot s
=n+b=-2$, czyli $b=-n-2$.
\end{proof}

\begin{center}
\begin{tikzpicture}[>=Stealth,scale=1]
  \draw[->,gray!75] (-0.4,0)--(5.4,0) node[right,font=\footnotesize] {$[s]$};
  \draw[->,gray!75] (0,-0.3)--(0,3.7) node[above,font=\footnotesize] {$[f]$};
  \fill[blue!12] (0,0)--(4.8,0)--(4.8,3.15)--(0,3.15)--cycle;
  \draw[blue!70!black,very thick] (0,0)--(4.8,0);
  \draw[blue!70!black,very thick] (0,0)--(0,3.2);
  \node[font=\footnotesize,align=center,fill=white,inner sep=3pt]
  at (2.65,2.25) {$\overline{\operatorname{NE}}(\mathbb F_n)$\\
  dodatnie kombinacje $[s]$ i $[f]$};
  \fill (0,0) circle (1.5pt);
\end{tikzpicture}

{\small Rysunek 8.3. Schemat stożka krzywych powierzchni
Hirzebrucha w bazie $([s],[f])$. Rysunek przedstawia
wycinek nieograniczonego stożka; długości wektorów
nie reprezentują liczb przecięć.}\par
\end{center}

\section{Krzywe z ujemnym przecięciem kanonicznym}
\label{sec:ga8-negative-curves}

Poniższy lemat jest jednym z najważniejszych elementarnych
ogniw programu modeli minimalnych dla powierzchni.
Pokazuje, że krzywe o ujemnym samoprzecięciu, które
jednocześnie mają ujemne przecięcie z $K_S$, są
ściśle określonej postaci.

\begin{twierdzenie}[Elementarne rozpoznawanie krzywej $(-1)$]
\label{tw:ga8-negative-implies-minusone}
Niech $S$ będzie gładką rzutową powierzchnią nad
algebraicznie domkniętym ciałem $k$, a $C\subseteq S$
całkowitą krzywą. Jeżeli $C^2<0$ i $K_S\cdot C<0$,
to $C$ jest gładką krzywą wymierną oraz
$C^2=K_S\cdot C=-1$. W szczególności
$C\simeq\PP^1_k$.
\end{twierdzenie}
\begin{proof}
Dla całkowitej, niekoniecznie gładkiej krzywej
na gładkiej powierzchni wzór adjunkcji arytmetycznej
przyjmuje postać
$2p_a(C)-2=C^2+K_S\cdot C$, gdzie
$p_a(C)=1-\chi(\mathcal O_C)$.
Krzywa jest całkowita i rzutowa, więc
$H^0(C,\mathcal O_C)=k$; dlatego
$p_a(C)=h^1(C,\mathcal O_C)\ge0$.
Obie liczby $C^2$ i $K_S\cdot C$ są ujemnymi
liczbami całkowitymi, zatem ich suma jest nie
większa niż $-2$. Wzór adjunkcji daje równocześnie
$C^2+K_S\cdot C=2p_a(C)-2\ge-2$.
Musi więc zachodzić równość $-2$ i każda z obu
liczb równa się $-1$. Ponadto $p_a(C)=0$.
Dla normalizacji $\nu:\widetilde C\to C$ zachodzi
$p_a(C)=g(\widetilde C)+\sum_{p\in C}\delta_p$,
gdzie $\delta_p\ge0$ mierzy defekt normalności.
Oba składniki znikają. Zatem $C$ jest normalna,
a krzywa normalna jest gładka; jej rodzaj wynosi zero.
Nad algebraicznie domkniętym ciałem otrzymujemy
$C\cong\PP^1$.
\end{proof}

\begin{przyklad}[Dlaczego nie wystarcza $K_S\cdot C<0$?]
Na $\PP^2$ każda prosta $L$ ma $L^2=1$ i
$K_{\PP^2}\cdot L=-3$, lecz nie jest wyjątkową
krzywą rozdmuchania. Na $\mathbb F_n$ włókno $f$
spełnia $f^2=0$ i $K\cdot f=-2$; również nie
można go kontraktować do gładkiego punktu.
Dopiero dodatkowe założenie $C^2<0$ pozwala
zastosować poprzednie twierdzenie.
\end{przyklad}

\begin{twierdzenie}[Kryterium kontrakcji Castelnuova — wynik zewnętrzny]
\label{tw:ga8-castelnuovo-mmp}
Niech $S$ będzie gładką rzutową powierzchnią nad
algebraicznie domkniętym ciałem charakterystyki zero.
Jeżeli $E\subset S$ jest krzywą izomorficzną z $\PP^1$
i $E^2=-1$, to istnieje gładka rzutowa powierzchnia
$S'$ i morfizm $\pi:S\to S'$, będący rozdmuchaniem
$S'$ w jednym punkcie, dla którego $E$ jest dokładnie
dzielnikiem wyjątkowym.
\end{twierdzenie}

\begin{uwaga}[Ujemność samoprzecięcia nie daje dowolnej kontrakcji]
Krzywa o samoprzecięciu $-2$ może być kontraktowalna
do \emph{punktu osobliwego}, ale nie spełnia kryterium
Castelnuova kontrakcji do gładkiego punktu.
Na $\mathbb F_2$ sekcja $s$ ma $s^2=-2$,
a $K\cdot s=0$. Nie należy więc interpretować
wszystkich krzywych ujemnych jako kolejnych kroków
programu gładkich modeli minimalnych.
\end{uwaga}

\section{Modele minimalne i wymiar Kodairy}
\label{sec:ga8-kodaira-minimal}

\begin{definicja}[Minimalność gładkiej powierzchni]
\label{def:ga8-minimal}
Gładką rzutową powierzchnię $S$ nazywamy
\emph{minimalną w sensie gładkich powierzchni},
gdy nie zawiera żadnej krzywej $E\simeq\PP^1$
o samoprzecięciu $E^2=-1$. Jeśli $K_S$ jest nef,
to $S$ nazywamy \emph{modelem minimalnym w sensie
kanonicznym}. Te określenia nie są bez dodatkowych
hipotez równoważne. Przymiotnik \emph{względnie minimalny}
rezerwujemy natomiast dla powierzchni z ustalonym morfizmem
do bazy, kiedy minimalność bada się względem tego morfizmu.
\end{definicja}

\begin{stwierdzenie}[Istnienie powierzchni minimalnej w sensie gładkim]
\label{st:ga8-terminate-minusone}
Niech $S$ będzie gładką rzutową powierzchnią nad
algebraicznie domkniętym ciałem charakterystyki zero.
Istnieje skończony ciąg kontrakcji krzywych $(-1)$
prowadzący do gładkiej powierzchni minimalnej w sensie gładkim.
\end{stwierdzenie}
\begin{proof}
Jeżeli $S$ nie zawiera krzywych $(-1)$, to koniec.
W przeciwnym wypadku kryterium Castelnuova dostarcza
morfizmu $S\to S_1$ kontraktującego taką krzywą.
Rozdmuchanie punktu zwiększa liczbę Picarda o jeden:
$\rho(S)=\rho(S_1)+1$; pokazaliśmy już rozkład
klas po rozdmuchaniu w rozdziale 6. Powierzchnia
$S_1$ znowu jest gładka i rzutowa. Powtarzając
konstrukcję, za każdym razem zmniejszamy dodatnią
liczbę całkowitą $\rho$. Nie da się więc wykonać
nieskończenie wielu kroków. Ostatnia powierzchnia
nie ma krzywych $(-1)$.
\end{proof}

\begin{przyklad}[Minimalność nie oznacza jednoznaczności]
Powierzchnia $\PP^2$ jest minimalna w sensie gładkim,
ponieważ każda niezerowa efektywna krzywa ma
klasę $dH$ z $d\ge1$, więc jej samoprzecięcie
to $d^2>0$. Powierzchnia
$\PP^1\times\PP^1$ również nie zawiera krzywych
$(-1)$: dla całkowitej krzywej klasy $aF_1+bF_2$
mamy $a,b\ge0$ i $C^2=2ab\ge0$.
Jednocześnie obie powierzchnie są wymierne, a więc
biwymierne, chociaż nie są izomorficzne:
ich liczby Picarda wynoszą odpowiednio $1$ i $2$.
To pokazuje, że bez dodatkowych hipotez
minimalny w sensie gładkim przedstawiciel klasy
biwymierności nie jest jedyny.
\end{przyklad}

\begin{definicja}[Wielowymiarowe układy kanoniczne i wymiar Kodairy]
\label{def:ga8-kodaira}
Niech $S$ będzie gładką spójną rzutową powierzchnią
nad $\C$. Dla $m\ge1$ definiujemy \emph{$m$-ty
plurigenus} przez $P_m(S)=h^0(S,\omega_S^{\otimes m})$.
Jeżeli $P_m(S)=0$ dla wszystkich $m\ge1$, przyjmujemy
$\kappa(S)=-\infty$. W przeciwnym razie
\emph{wymiarem Kodairy} nazywamy
$\kappa(S)=\max_{m:P_m(S)>0}\dim \varphi_{|mK_S|}(S)$,
gdzie $\varphi_{|mK_S|}$ jest odwzorowaniem wymiernym
danym pełnym układem liniowym $|mK_S|$; gdy układ
ma jeden niezależny przekrój, jego obraz jest punktem.
Dla powierzchni $\kappa(S)\in\{-\infty,0,1,2\}$.
\end{definicja}

\begin{stwierdzenie}[Niezmienniczość plurigenera przy rozdmuchaniu]
\label{st:ga8-plurigenera-blowup}
Niech $\pi:\widetilde S\to S$ będzie rozdmuchaniem
punktu gładkiej rzutowej powierzchni $S$ nad $\C$.
Dla każdego $m\ge1$ zachodzi
$H^0(\widetilde S,\omega_{\widetilde S}^{\otimes m})
\cong H^0(S,\omega_S^{\otimes m})$.
W konsekwencji wszystkie $P_m$ i wymiar Kodairy
są niezmiennikami biwymiernymi gładkich
rzutowych powierzchni.
\end{stwierdzenie}
\begin{proof}
Wiemy, że $K_{\widetilde S}=\pi^*K_S+E$,
gdzie $E$ jest dzielnikiem wyjątkowym. Dla $m\ge1$
wzór projekcji daje
$\pi_*\mathcal O_{\widetilde S}(mK_{\widetilde S})
=\mathcal O_S(mK_S)\otimes\pi_*\mathcal O_{\widetilde S}(mE)$.
Pokażemy, że $\pi_*\mathcal O_{\widetilde S}(mE)
=\mathcal O_S$. Nad $S\setminus\{p\}$ obie wiązki
są identyczne. Sekcja $\mathcal O(mE)$ nad
$\pi^{-1}(U)$ może być utożsamiona z funkcją wymierną,
która jest regularna nad $U\setminus\{p\}$:
warunek na dywizor nie dopuszcza biegunów poza $E$.
Ponieważ gładka powierzchnia $S$ jest normalna,
regularne funkcje przedłużają się przez podzbiór
kowymiaru dwa. Tak otrzymujemy sekcję $\mathcal O_U$.
Odwrotnie każda funkcja regularna na $U$ daje
sekcję $\mathcal O(mE)$. Wzór projekcji i wzięcie
sekcji globalnych kończą dowód równości plurigenusów.
Co więcej, po pomnożeniu pullbacku przekroju $mK_S$
przez kanoniczny przekrój $\mathcal O(mE)$ otrzymujemy
\emph{wszystkie} przekroje $mK_{\widetilde S}$. Dlatego
jeśli $|mK_S|$ jest niepusty, układ $|mK_{\widetilde S}|$
powstaje przez pomnożenie pullbacków wszystkich przekrojów
$|mK_S|$ przez ten sam przekrój $\mathcal O(mE)$.
Na otwartym podzbiorze, na którym $\pi$ jest izomorfizmem,
oba odwzorowania plurikanoniczne zgadzają się po utożsamieniu
dziedzin i mają zatem ten sam wymiar obrazu.
To osobno uzasadnia niezmienniczość $\kappa$ z definicji
przez wymiar obrazów, a nie tylko równość liczb $P_m$.
Przejście od rozdmuchań do dowolnej biwymierności
korzysta z twierdzenia o faktoryzacji odwzorowań
biwymiernych gładkich rzutowych powierzchni,
które przyjmujemy jako wynik zewnętrzny.
\end{proof}

\begin{przyklad}[Wszystkie cztery wartości wymiaru Kodairy]
\label{prz:ga8-kodaira-four}
Rozpatrzmy gładkie krzywe projektowe $C_g,C_h$
rodzajów $g,h\ge2$ i krzywą eliptyczną $E$.
\begin{enumerate}[label=(\roman*)]
\item $\kappa(\PP^2)=-\infty$, ponieważ
$K_{\PP^2}=-3H$ i $H^0(\PP^2,\mathcal O(-3m))=0$
dla każdego $m\ge1$.
\item Jeżeli $S$ jest K3, to $K_S\sim0$ i
$P_m(S)=1$ dla wszystkich $m\ge1$, zatem $\kappa(S)=0$.
\item Dla $S=C_g\times E$ mamy
$\omega_S\simeq\mathrm{pr}_1^*\omega_{C_g}$,
a więc $P_m(S)=h^0(C_g,\omega_{C_g}^{\otimes m})$.
Dla $m\ge2$ wzór Riemanna--Rocha na krzywej
daje $(2m-1)(g-1)$, zaś odwzorowania
plurikanoniczne mają obraz będący krzywą.
Stąd $\kappa(S)=1$.
\item Dla $S=C_g\times C_h$ mamy
$P_m(S)=(2m-1)^2(g-1)(h-1)$ dla $m\ge2$.
Wymiar obrazu odpowiednich odwzorowań wynosi dwa,
więc $\kappa(S)=2$.
\end{enumerate}
W dwóch przykładach z iloczynami użyliśmy równości
$H^0(C\times D,L\boxtimes M)
\cong H^0(C,L)\otimes H^0(D,M)$ oraz faktu,
że dostatecznie wysoka wielokrotność dywizora
kanonicznego na krzywej rodzaju co najmniej dwa
daje odwzorowanie o jednowymiarowym obrazie.
\end{przyklad}

\section{Dwuwymiarowy program modeli minimalnych}
\label{sec:ga8-surface-mmp}

Przechodzimy od procedury kontraktowania wszystkich
krzywych $(-1)$ do kryterium, które samo wskazuje,
\emph{czy} i \emph{jakiego typu} należy wykonać następny krok.
To jest zasadnicza idea programu modeli minimalnych.

\begin{twierdzenie}[Stożek i kontrakcje dla gładkich powierzchni — wynik zewnętrzny]
\label{tw:ga8-surface-contraction}
Niech $S$ będzie gładką rzutową powierzchnią nad
algebraicznie domkniętym ciałem charakterystyki zero.
Część stożka $\overline{\operatorname{NE}}(S)$, na której
$K_S$ jest ujemny, jest kontrolowana przez przeliczalną
rodzinę promieni ekstremalnych generowanych przez klasy
krzywych wymiernych (twierdzenie o stożku).
Dla każdego $K_S$-ujemnego promienia ekstremalnego $R$
istnieje kontrakcja $f:S\to Y$ o spójnych włóknach,
kontraktująca dokładnie te krzywe $C$, dla których
$[C]\in R$. Jeśli $Y$ ma wymiar dwa, $f$ jest
kontrakcją krzywej $(-1)$ do gładkiego punktu.
Jeżeli $\dim Y<2$, to $f$ jest \emph{fibracją Moriego}:
ma dodatni względny wymiar, względną liczbę Picarda
równą jeden i $-K_S$ jest względnie obszerny.
\end{twierdzenie}

\begin{uwaga}[Zakres twierdzenia]
Pełny dowód wykorzystuje twierdzenie o stożku i twierdzenie
o kontrakcjach. Nie wyprowadzamy ich z samego twierdzenia
Castelnuova. W przypadku $\dim Y=1$ ogólne włókno
jest izomorficzne z $\PP^1$; dla $\dim Y=0$ dywizor
$-K_S$ jest obszerny, a $\rho(S)=1$.
To są trzy typy kroków: kontrakcja biwymierna,
fibracja nad krzywą lub kontrakcja do punktu.
\end{uwaga}

\begin{twierdzenie}[Program modeli minimalnych dla gładkich powierzchni]
\label{tw:ga8-surface-mmp}
Dla każdej gładkiej rzutowej powierzchni $S$
nad algebraicznie domkniętym ciałem charakterystyki zero
istnieje skończony ciąg kontrakcji $K$-ujemnych
promieni ekstremalnych. Kończy się on w jednej
i tylko jednej z następujących sytuacji:
\begin{enumerate}[label=(\alph*)]
\item otrzymujemy gładką rzutową powierzchnię $S_{\min}$,
dla której $K_{S_{\min}}$ jest nef;
\item otrzymujemy fibrację Moriego $f:S_r\to Y$,
gdzie $\dim Y\in\{0,1\}$.
\end{enumerate}
Jeśli $\kappa(S)\ge0$, zachodzi przypadek (a).
W przypadku (b) powierzchnia ma $\kappa=-\infty$.
\end{twierdzenie}
\begin{proof}
Korzystamy z zewnętrznego twierdzenia
\ref{tw:ga8-surface-contraction}. Gdy $K$ jest nef,
procedura się kończy. W przeciwnym razie istnieje
$K$-ujemny promień ekstremalny i jego kontrakcja.
Jeśli jest to fibracja Moriego, otrzymujemy (b)
i kończymy. Jeśli kontrakcja jest biwymierna,
odwraca rozdmuchanie punktu, więc nowe $S_1$
jest gładkie i $\rho(S_1)=\rho(S)-1$.
Ograniczoność dodatniej liczby $\rho$ daje skończoność
liczby kroków biwymiernych. Wynikiem jest więc
(a) lub (b). Jeśli zachodzi (b), to ogólne włókno
nad krzywą jest $\PP^1$ albo $-K$ jest obszerny
w przypadku bazy punktowej; w obu przypadkach
nie ma niezerowych przekrojów $mK$ dla $m>0$.
W pierwszym przypadku widać to po restrykcji
do ogólnego włókna, gdzie
$\deg(mK|_{\PP^1})=-2m<0$.
W drugim niezerowy efektywny $mK$ miałby
ujemne przecięcie z obszernym $-K$, co jest
niemożliwe. Niezmienniczość plurigenusów przy
krokach biwymiernych daje $\kappa(S)=-\infty$.
Tym samym gdy $\kappa(S)\ge0$, przypadek (b)
jest wykluczony.
\end{proof}

\begin{center}
\begin{tikzpicture}[>=Stealth,
 every node/.style={font=\footnotesize,align=center},
 nodebox/.style={draw,rounded corners,inner sep=4pt}]
\node[nodebox,text width=37mm] (start) at (0,3.1)
 {gładka powierzchnia $S$};
\node[nodebox,text width=37mm] (decision) at (0,1.65)
 {czy $K_S$ jest nef?};
\node[nodebox,text width=36mm] (nef) at (-3.2,-0.2)
 {tak: model kanonicznie minimalny};
\node[nodebox,text width=41mm] (ray) at (3.25,-0.2)
 {nie: kontrakcja promienia $K$-ujemnego};
\node[nodebox,text width=34mm] (blow) at (0.65,-2.5)
 {biwymierna: $\rho$ maleje; powtórz};
\node[nodebox,text width=34mm] (fibre) at (5.1,-2.5)
 {fibracja Moriego: koniec};
\draw[->] (start)--(decision);
\draw[->] (decision)--node[above left,font=\scriptsize] {tak} (nef);
\draw[->] (decision)--node[above right,font=\scriptsize] {nie} (ray);
\draw[->] (ray)--(blow);
\draw[->] (ray)--(fibre);
\end{tikzpicture}

{\small Rysunek 8.4. Schemat dwuwymiarowego programu modeli
minimalnych. W przypadku kontrakcji biwymiernej
wracamy do testu nefności na nowej powierzchni.}\par
\end{center}

\begin{przyklad}[Trzy możliwe zachowania promienia ekstremalnego]
\label{prz:ga8-extremal-examples}
\begin{enumerate}[label=(\roman*)]
\item Na $\operatorname{Bl}_p\PP^2=\mathbb F_1$
krzywa wyjątkowa $s$ spełnia $s^2=-1$,
$K\cdot s=-1$ i jej kontrakcja prowadzi do $\PP^2$.
\item Na $\mathbb F_2$ włókno $f$ spełnia
$f^2=0$, $K\cdot f=-2$. Kontrakcja odpowiada
projekcji $\mathbb F_2\to\PP^1$; nie jest
biwymierna i nie usuwa sekcji $s$ o $s^2=-2$.
\item Na $\PP^2$ promień klasy prostej $H$
ma $H^2=1$, $K\cdot H=-3$.
Kontrakcja prowadzi do punktu, a nie do innej
powierzchni. Jest to przypadek Fano o $\rho=1$.
\end{enumerate}
W szczególności nie należy utożsamiać każdego
$K$-ujemnego promienia ze znikającą krzywą $(-1)$.
\end{przyklad}

\begin{przyklad}[Dwie różne kontrakcje na jednej powierzchni]
\label{prz:ga8-f1-two-rays}
Niech $S=\mathbb F_1\simeq\operatorname{Bl}_p\PP^2$, a $s$
niech oznacza krzywą wyjątkową oraz $f$ włókno
naturalnego rzutowania $S\to\PP^1$. Formuły
ze stwierdzenia~\ref{st:ga8-fn-cone} dla $n=1$ dają
$K_S=-2s-3f$, $K_S\cdot s=-1$ i $K_S\cdot f=-2$.
Oba promienie stożka krzywych są więc $K$-ujemne.
Kontrakcja promienia $[s]$ jest biwymierna i daje
$S\to\PP^2$, czyli odwraca rozdmuchanie punktu.
Kontrakcja promienia $[f]$ ma włókna $\PP^1$
i daje morfizm $S\to\PP^1$; jest fibracją Moriego,
a nie morfizmem biwymiernym. Liczba Picarda włókien
względem bazy wynosi jeden, a $-K_S\cdot f=2>0$.
Dlatego wybór promienia ekstremalnego naprawdę wpływa
na typ kroku i wynik programu, gdy $\kappa=-\infty$.
\end{przyklad}

\begin{twierdzenie}[Jednoznaczność dla nieujemnego wymiaru Kodairy — wynik zewnętrzny]
\label{tw:ga8-minimal-unique}
Niech $S,S'$ będą gładkimi rzutowymi powierzchniami
nad $\C$ takimi, że $\kappa(S)\ge0$ i $S,S'$
są biwymierne. Jeśli $K_S$ i $K_{S'}$ są nef,
to $S\cong S'$. Zatem każda klasa biwymierności
powierzchni o $\kappa\ge0$ ma, z dokładnością do
izomorfizmu, jeden gładki model kanonicznie minimalny.
\end{twierdzenie}
\begin{uwaga}[Dlaczego przypadek wymierny jest inny?]
Dla powierzchni o $\kappa=-\infty$ nie żądamy
kanonicznej nefności, bo $\PP^2$ ma $K=-3H$,
a wszystkie powierzchnie Hirzebrucha mają
$K\cdot f=-2$. Mogą jednak nie mieć krzywych
$(-1)$; otrzymane powierzchnie minimalne w sensie gładkim
nie muszą być izomorficzne. To wyjaśnia przykład
$\PP^2$ i $\PP^1\times\PP^1$.
\end{uwaga}

\section{Nefność, rozkład Zariskiego i granice teorii}
\label{sec:ga8-zariski-decomposition}

\begin{twierdzenie}[Rozkład Zariskiego na powierzchni — wynik zewnętrzny]
\label{tw:ga8-zariski-decomposition}
Niech $S$ będzie gładką rzutową powierzchnią nad
algebraicznie domkniętym ciałem charakterystyki zero,
a $D$ efektywnym dywizorem o współczynnikach wymiernych.
Istnieje jednoznaczny rozkład $D=P+N$, gdzie
$P,N$ są dywizorami o współczynnikach wymiernych,
$P$ jest nef, $N$ jest efektywny, a jeśli $N\ne0$,
to macierz przecięć składowych $N$ jest ujemnie
określona oraz $P\cdot N_i=0$ dla każdej składowej
nierozkładalnej $N_i$ dzielnika $N$.
\end{twierdzenie}

\begin{przyklad}[Rozkład Zariskiego obliczony do końca]
Niech $\pi:\widetilde S\to\PP^2$ będzie rozdmuchaniem
punktu, $H=\pi^*\mathcal O_{\PP^2}(1)$ i $E$
krzywą wyjątkową. Dla efektywnego dywizora
$D=H+E$ kładziemy $P=H$ i $N=E$.
Wówczas $H$ jest nef, $E$ jest efektywny,
$E^2=-1$ i $H\cdot E=0$. Spełnione są wszystkie
warunki twierdzenia, więc jest to jedyny rozkład.
Obserwujemy przy tym $D\cdot E=-1$,
choć dodatnia część $P$ nie ma ujemnych przecięć.
Nie oznacza to, że dodanie $E$ zwiększa liczbę
niezależnych przekrojów: część ujemna jest
właśnie stałą częścią odpowiednio wysokich
wielokrotności układu liniowego $|D|$.
\end{przyklad}

W \emph{wyższych wymiarach} nie wystarcza
kontraktowanie pojedynczych krzywych $(-1)$;
pojawiają się m.in. małe kontrakcje, przekształcenia
\emph{flip} oraz modele z osobliwościami.
Nie omawiamy ich tutaj. Naszym celem było
przeprowadzenie dwuwymiarowego algorytmu na
podstawie przecięć oraz pokazanie, które składniki
wymagają głębokich twierdzeń, a które wynikają
z elementarnego rachunku na powierzchni.

% Literatura do uzupelnienia MMP:
% J. Kollar, S. Mori, Birational Geometry of Algebraic Varieties;
% M. Reid, Chapters on Algebraic Surfaces;
% S. Frengley, Algebraic Surfaces (uniwersyteckie notatki online);
% Stacks Project, contractions and minimal models of surfaces.


\section{Podsumowanie geometrii biwymiernej i modeli minimalnych}
\label{sec:ga8-conclusion}

Przekształcenie Cremony dostarczyło kompletnego,
obliczalnego przykładu odwzorowania wymiernego,
które ma trzy punkty bazowe i kontraktuje trzy proste.
Rozdmuchanie tych punktów zamienia odwzorowanie
w morfizm, a graf przecięć sześciu krzywych
$(-1)$ pokazuje jego geometrię lokalną i globalną.
Wzory przecięć i klasy Cherna pozwalają śledzić
przy tym liczby geometryczne.

Zbudowaliśmy również język równoważności numerycznej,
stożka krzywych i nefności; otrzymaliśmy kryterium
krzywych $(-1)$ oraz dwuwymiarowy algorytm kontrakcji.
Wyraźnie wydzieliliśmy głębokie twierdzenia o stożku,
kontrakcjach i jednoznaczności modeli minimalnych od
rachunków wykonanych w skrypcie. W szczególności
nie utożsamiamy minimalności gładkiej powierzchni
z nefnością $K$ w przypadku $\kappa=-\infty$.

% Odniesienia kontrolne:
% J. Blanc, Introduction to Birational Geometry of Surfaces;
% I. Dolgachev, Classical Algebraic Geometry (Cremona transformations);
% Encyklopedia matematyki: Cremona transformations.


\chapter{Geometria algebraiczna a analiza zespolona: analityfikacja, Chow, GAGA i Kähler}
\label{ch:ga9-analiza-zespolona}

\section{Dwie geometrie tych samych równań}
\label{sec:ga9-motywacja}

Przez kilka poprzednich rozdziałów badaliśmy przestrzenie opisane równaniami
wielomianowymi. Gdy ciało bazowe jest równe $\C$, możemy tę samą przestrzeń
oglądać za pomocą funkcji holomorficznych. Wówczas pojawia się pytanie:
\emph{jak wiele informacji traci się, zastępując wielomiany funkcjami
holomorficznymi?} Odpowiedź na przestrzeni afinicznej jest zupełnie inna niż
na przestrzeni rzutowej. Ta różnica tłumaczy sens twierdzeń Chowa i GAGA.

W całym rozdziale wszystkie schematy algebraiczne są skończonego typu nad
$\C$. Gdy mówimy o rozmaitości algebraicznej bez dodatkowego określenia,
zakładamy, że jest zredukowana i całkowita. W twierdzeniach dotyczących
rozmaitości zespolonych wskazujemy osobno warunki gładkości oraz zwartości.
Nie utożsamiamy przestrzeni topologicznej z jej snopem strukturalnym.

\begin{przyklad}[Pierścień funkcji regularnych i holomorficznych]
Na afinicznej prostej algebraicznej $\A^1_\C$ funkcje regularne na całej
przestrzeni są wielomianami: $\mathcal O_{\rm alg}(\A^1)=\C[z]$.
Na zespolonej płaszczyźnie $\C$ funkcje holomorficzne tworzą znacznie
większy pierścień $\mathcal O_{\rm hol}(\C)$; należy do niego $e^z$.
Nie istnieje wielomian $P\in\C[z]$ taki, że $P(z)=e^z$ dla każdego
$z\in\C$. Istotnie, po ograniczeniu do dodatniej półosi rzeczywistej
otrzymalibyśmy $P(x)=e^x$, lecz $e^x/x^m\to+\infty$ dla każdego
$m\ge0$, podczas gdy $|P(x)|/x^m$ pozostaje ograniczone, jeśli
$m\ge\deg P$. Jest to przykład różnicy między funkcjami regularnymi
i holomorficznymi, a nie tylko różnicy między wybranymi wzorami.
\end{przyklad}

\begin{uwaga}[Ten sam zbiór punktów nie oznacza tej samej kategorii]
Punkty zespolone $\A^1(\C)$ i punkty analityczne $\C$ można utożsamić,
lecz w topologii Zariskiego zbiory otwarte są dopełnieniami zbiorów
skończonych, podczas gdy w topologii euklidesowej otwartych zbiorów jest
znacznie więcej. Co ważniejsze, różnią się też odpowiednie snopy
funkcji. Analityfikacja jest więc konstrukcją na przestrzeniach
\emph{z ich strukturą}, nie tylko zmianą nazwy punktów.
\end{uwaga}

\section{Zbiory analityczne i analityfikacja odmian}
\label{sec:ga9-analityczne}

\begin{definicja}[Zbiór analityczny]
Niech $M$ będzie rozmaitością zespoloną o wymiarze zespolonym $n$,
a $A\subseteq M$ podzbiorem. Mówimy, że $A$ jest
\emph{zespolonym zbiorem analitycznym}, jeżeli dla każdego $p\in M$
istnieje otoczenie otwarte $U$ punktu $p$ i funkcje holomorficzne
$f_1,\ldots,f_r\in\mathcal O_M(U)$ takie, że
$A\cap U=\{x\in U:f_1(x)=\cdots=f_r(x)=0\}$.
Jeśli lokalnie rozpatrujemy dodatkowo snop ideałów funkcji znikających,
otrzymujemy odpowiednią zredukowaną przestrzeń analityczną.
\end{definicja}

\begin{uwaga}[Lokalność nie oznacza globalnego układu równań]
Definicja dopuszcza różne równania na różnych otoczeniach $U$.
W ogólności zbiór analityczny w rozmaitości zespolonej nie musi być
zbiorem wspólnych zer skończonej rodziny globalnych funkcji
holomorficznych. Na przykład na zwartej spójnej rozmaitości zespolonej
wszystkie globalne funkcje holomorficzne są stałe, lecz mogą istnieć
nietrywialne dzielniki analityczne. Snopowość opisu jest zatem
niezbędna.
\end{uwaga}

\begin{definicja}[Analityfikacja]
Niech $X$ będzie schematem skończonego typu nad $\C$. Jego
\emph{analityfikacją} nazywamy zespoloną przestrzeń analityczną
$X^{\rm an}$, której punkty odpowiadają punktom $X(\C)$,
wraz z kanonicznym odwzorowaniem $a:X^{\rm an}\to X$
do przestrzeni Zariskiego i morfizmem snopów
$a^{-1}\mathcal O_X\to\mathcal O_{X^{\rm an}}$.
Jej lokalny model dla
$X=\operatorname{Spec}\C[z_1,\ldots,z_n]/I$ otrzymujemy jako wspólny
zbiór zer generatorów $I$ w $\C^n$, wyposażony w iloraz snopu
funkcji holomorficznych przez snop ideałów generowany przez $I$.
Modele afiniczne skleja się przy użyciu holomorficznych przejść
pochodzących od morfizmów algebraicznych. Dla gładkiego $X$ jest to
zwykła rozmaitość zespolona.
\end{definicja}

\begin{uwaga}[Nilpotenty i osobliwości]
Jeżeli $I=(w^2)\subset\C[z,w]$, to analityfikacja schematu
$\operatorname{Spec}\C[z,w]/(w^2)$ ma ten sam zbiór punktów co prosta
$w=0$, lecz jej snop strukturalny nadal zawiera niezerowy nilpotent
$w$. Analityfikacja schematów nie polega więc na wyrzucaniu informacji
o krotnościach. W przypadku $I=(w^2-z^3)$ otrzymujemy ten sam
analityczny kiełek ostrza, którego normalizację opisaliśmy w
rozdziałach~\ref{ch:ga5-rozdmuchania} i~\ref{ch:ga6-powierzchnie}.
\end{uwaga}

\begin{stwierdzenie}[Wielomiany definiują zbiory analityczne]
Niech $I\subseteq\C[z_1,\ldots,z_n]$ będzie ideałem. Wówczas
$V(I)\subseteq\C^n$, rozpatrywany w zwykłej topologii, jest
zamkniętym zbiorem analitycznym. Jeśli $I$ jest generowany przez $f_1,\ldots,f_r$, to te same
wielomiany definiują analitycznie \emph{zbiór punktów} $V(I)$.
Aby otrzymać analityfikację schematu z ewentualnymi nilpotentami,
należy ponadto zachować iloraz snopu holomorficznego przez snop
ideałów generowany przez $f_1,\ldots,f_r$.
\end{stwierdzenie}
\begin{proof}
Każdy wielomian jest funkcją holomorficzną, a wspólny zbiór zer
funkcji ciągłych jest domknięty. Ponieważ $\C[z_1,\ldots,z_n]$ jest
pierścieniem noetherowskim (twierdzenie Hilberta o bazie z
rozdziału~\ref{ch:pierscienie}), ideał $I$ ma skończony układ
generatorów. Definicję zbioru analitycznego spełnia więc nawet
jedno otoczenie $U=\C^n$ z tym właśnie układem równań.
\end{proof}

\section{Dlaczego w przestrzeni afinicznej odwrotność zawodzi?}
\label{sec:ga9-wykladnicza}

\begin{przyklad}[Analityczna, lecz niealgebraiczna krzywa]
\label{prz:ga9-exp}
Rozważmy $\Gamma=\{(z,w)\in\C^2:w=e^z\}$.
Jest to zamknięty zbiór analityczny, ponieważ $w-e^z$ jest
funkcją holomorficzną, a różniczka $\mathrm d(w-e^z)$ ma współczynnik
przy $\mathrm dw$ równy $1$; zatem $\Gamma$ jest nawet gładką
podrozmaitością zespoloną wymiaru jeden. Pokażemy, że nie jest
zbiorem algebraicznym.

Załóżmy, że niezerowy wielomian $P(z,w)=\sum_{j=0}^m p_j(z)w^j$
znika na $\Gamma$, gdzie $p_m\ne0$. Wtedy dla każdego $z\in\C$
mamy $\sum_{j=0}^m p_j(z)e^{jz}=0$. Wstawmy $z=x\in\R$
i podzielmy przez $e^{mx}$. Otrzymujemy
\[
 p_m(x)=-\sum_{j<m}p_j(x)e^{-(m-j)x}.
\]
Prawa strona dąży do zera, gdy $x\to+\infty$, gdyż każda
funkcja wykładnicza $e^{-cx}$, $c>0$, dominuje wielomianowy
wzrost $|p_j(x)|$. Oznaczałoby to $p_m(x)\to0$, co jest niemożliwe
dla niezerowego wielomianu. Zatem żaden niezerowy wielomian
nie znika na całej $\Gamma$. W szczególności $\Gamma$ nie może
być algebraiczną krzywą w $\A^2_\C$.
\end{przyklad}

\begin{center}
\begin{tikzpicture}[>=Stealth,scale=.93]
  \draw[->] (-3.2,0)--(3.15,0) node[right,font=\footnotesize] {$x$};
  \draw[->] (0,-.35)--(0,3.4) node[above,font=\footnotesize] {$w$};
  \draw[blue!75!black,very thick,domain=-3.1:1.18,samples=80]
    plot (\x,{.65*exp(\x)});
  \node[fill=white,inner sep=2pt,font=\footnotesize,anchor=west]
    at (.72,2.98) {$w=e^x$};
  \node[font=\scriptsize,align=center] at (-1.65,2.65)
    {rzeczywisty przekrój\\zespolonej krzywej $\Gamma$};
\end{tikzpicture}

{\small Rysunek 9.1. Rzeczywisty ślad grafu funkcji wykładniczej;
rysunek jest wyłącznie przekrojem pomocniczym, a nie modelem
pełnej krzywej zespolonej w $\C^2$.}\par
\end{center}

\begin{uwaga}[Dlaczego twierdzenie Chowa nie daje sprzeczności?]
Graf funkcji wykładniczej jest zamknięty w afinicznej przestrzeni
$\C^2$, ale jego obraz w $\PP^2(\C)$ nie jest domkniętym zbiorem
analitycznym. Twierdzenie Chowa będzie dotyczyło tylko zbiorów
\emph{zamkniętych analitycznie w przestrzeni rzutowej}.
Nie wystarcza sama algebraiczność przestrzeni otaczającej.
\end{uwaga}

\section{Lokalne pierścienie holomorficzne i przygotowanie Weierstrassa}
\label{sec:ga9-lokalne}

\begin{definicja}[Pierścień kiełków holomorficznych]
Niech $p\in\C^n$. Przez $\mathcal O_{\C^n,p}$ oznaczamy pierścień
klas par $(U,f)$, gdzie $p\in U\subseteq\C^n$ jest otwarty,
$f\in\mathcal O(U)$, a dwie pary są równoważne, gdy funkcje
zgadzają się na pewnym otoczeniu $p$. Jest to pierścień lokalny
o ideale maksymalnym $\mathfrak m_p=\{[f]:f(p)=0\}$.
Po przesunięciu $p$ do zera identyfikujemy go z pierścieniem
zbieżnych szeregów potęgowych $\C\{z_1,\ldots,z_n\}$.
\end{definicja}

\begin{proof}
Klasy dodajemy i mnożymy przez wybór wspólnego, odpowiednio
zmniejszonego otoczenia. Jeśli $f(p)\ne0$, ciągłość daje
otoczenie, na którym $f$ nie ma zer, więc $1/f$ jest tam
holomorficzna i $[f]$ jest jednostką. Jeśli $f(p)=0$,
to $[f]$ nie może być jednostką: wartość w $p$ iloczynu
$[f][g]$ jest równa zero. Elementy nieodwracalne tworzą
zatem dokładnie ideał $\mathfrak m_p$, który jest maksymalny,
bo homomorfizm wartościowania $[f]\mapsto f(p)$ ma jako
jądro $\mathfrak m_p$ i jako obraz $\C$. Jedyność ideału
maksymalnego wynika stąd, że każdy właściwy ideał składa się
z elementów nieodwracalnych. Identyfikacja z pierścieniem
szeregów zbieżnych wynika z lokalnego rozwinięcia Taylora
funkcji holomorficznej.
\end{proof}

\begin{twierdzenie}[Przygotowanie Weierstrassa — wynik zewnętrzny]
\label{tw:ga9-weierstrass}
Niech $n\ge1$ będzie liczbą całkowitą,
$f\in\C\{z_1,\ldots,z_n\}$ i załóżmy, że
$f(0,\ldots,0,z_n)$ ma w zerze rząd dokładnie $d\ge1$.
Wtedy istnieje jednostka $u\in\C\{z_1,\ldots,z_n\}^\times$
oraz jednoznaczny wielomian wyróżniony
\[
 W=z_n^d+a_1(z')z_n^{d-1}+\cdots+a_d(z'),
 \qquad a_j\in\C\{z'\},\quad a_j(0)=0,
\]
gdzie $z'=(z_1,\ldots,z_{n-1})$, taki że $f=uW$.
\end{twierdzenie}

\begin{uwaga}[Znaczenie twierdzenia]
Choć $f$ może mieć nieskończenie wiele niezerowych wyrazów
szeregu Taylora, po pomnożeniu przez jednostkę można traktować
jego zera lokalnie jak zera \emph{wielomianu w jednej zmiennej}
o współczynnikach holomorficznych w pozostałych zmiennych.
Jest to punkt wyjścia do finitości i lokalnej teorii zbiorów
analitycznych. Dowód przygotowania Weierstrassa jest ważnym,
samodzielnym rezultatem analizy zespolonej; tutaj podajemy jego
precyzyjną treść, nie zastępując dowodu nieuzasadnionym
rachunkiem formalnym.
\end{uwaga}

\begin{przyklad}[Ostrze w pierścieniu zbieżnych szeregów]
Funkcja $f(z,w)=w^2-z^3$ jest już wielomianem wyróżnionym
stopnia dwa względem $w$. Stąd lokalny pierścień analityczny
ostrza w zerze jest $\C\{z,w\}/(w^2-z^3)$.
Parametryzacja $z=t^2$, $w=t^3$ daje naturalny homomorfizm
$\C\{z,w\}/(w^2-z^3)\to\C\{t\}$, którego obrazem jest
$\C\{t^2,t^3\}$. Po dołączeniu funkcji $t$ otrzymujemy
normalizację w sensie analitycznym. Ta sama parametryzacja
działa algebraicznie, lecz pierścienie $\C[t]$ i $\C\{t\}$
nie są tym samym obiektem.
\end{przyklad}

\section{Snopy koherentne: właściwy język porównania}
\label{sec:ga9-koherentne}

\begin{definicja}[Koherentny snop analityczny]
Niech $X$ będzie zespoloną przestrzenią analityczną,
a $\mathcal F$ snopem modułów nad $\mathcal O_X$.
Nazywamy $\mathcal F$ \emph{koherentnym}, jeżeli lokalnie
jest skończenie prezentowany i jądra morfizmów z lokalnie
wolnych snopów skończonego rzędu do $\mathcal F$ są lokalnie
skończenie generowane. Równoważnie, w standardowej teorii
snopów koherentnych nad zespolonymi przestrzeniami
analitycznymi, $\mathcal F$ jest lokalnie ilorazem
$\mathcal O_X^{\oplus r}$ przez koherentny snop relacji.
\end{definicja}

\begin{uwaga}[Dlaczego potrzebny jest warunek na relacje?]
Przy przejściu do obrazu lub jądra morfizmu snopów trzeba
kontrolować nie tylko liczbę generatorów, ale również ich
relacje. Fundamentalne twierdzenie Oki o koherencji snopu
$\mathcal O_{\C^n}$ oraz jego konsekwencje zapewniają
stabilność koherentności przy operacjach występujących
w geometrii analitycznej. Są to wyniki zewnętrzne.
\end{uwaga}

\begin{przyklad}[Snop ideałów podprzestrzeni analitycznej]
Jeśli $A\subseteq X$ jest zamkniętym zredukowanym
podzbiorem analitycznym, przez $\mathcal I_A$ oznaczamy
snop kiełków funkcji holomorficznych zerujących się na $A$.
Wtedy $\mathcal O_A=\mathcal O_X/\mathcal I_A$,
a ciąg
$0\to\mathcal I_A\to\mathcal O_X\to\mathcal O_A\to0$
jest dokładny. Koherentność $\mathcal I_A$ jest
konsekwencją twierdzeń Oki--Cartana; nie wynika wyłącznie
z definicji analityczności przez lokalne równania.
\end{przyklad}

\section{Twierdzenie GAGA: co dokładnie porównujemy?}
\label{sec:ga9-gaga}

\begin{twierdzenie}[Serre, GAGA — wynik zewnętrzny]
\label{tw:ga9-gaga}
Niech $X$ będzie schematem właściwym i skończonego typu nad $\C$.
Analityfikacja wyznacza równoważność kategorii
\[
 \operatorname{Coh}(X)\ \simeq\ \operatorname{Coh}(X^{\rm an}),
 \qquad \mathcal F\longmapsto\mathcal F^{\rm an},
\]
między koherentnymi snopami algebraicznymi i koherentnymi
snopami analitycznymi. Ponadto dla każdego koherentnego
snopu algebraicznego $\mathcal F$ oraz $q\ge0$ naturalny
homomorfizm
$H^q(X,\mathcal F)\to H^q(X^{\rm an},\mathcal F^{\rm an})$
jest izomorfizmem. W szczególności każdy koherentny
snop analityczny na $X^{\rm an}$ ma, z dokładnością do
izomorfizmu, jednoznaczny algebraiczny odpowiednik.
\end{twierdzenie}

\begin{uwaga}[Co znaczy równoważność, a czego ona nie mówi?]
Twierdzenie obejmuje nie tylko istnienie algebraicznego
snopu, lecz także pełną wierność: wszystkie morfizmy
koherentnych snopów analitycznych między analityfikacjami
pochodzą dokładnie z morfizmów algebraicznych.
Nie oznacza to równości pierścieni funkcji holomorficznych
i regularnych na dowolnym \emph{afinicznym} otwartym
podzbiorze $X^{\rm an}$: warunek właściwości jest
istotny. Nie oznacza też, że bezpośrednio porównujemy
wszystkie snopy zbiorów, tylko snopy koherentne.
\end{uwaga}

\begin{center}
\begin{tikzpicture}[>=Stealth,
 every node/.style={font=\footnotesize,align=center},
 nodebox/.style={draw,rounded corners,inner sep=6pt,text width=46mm}]
 \node[nodebox] (a) at (-3.2,1.3)
   {snopy koherentne\\algebraiczne na $X$};
 \node[nodebox] (b) at (3.2,1.3)
   {snopy koherentne\\analityczne na $X^{\rm an}$};
 \node[nodebox] (c) at (-3.2,-1.05)
   {$H^q(X,\mathcal F)$};
 \node[nodebox] (d) at (3.2,-1.05)
   {$H^q(X^{\rm an},\mathcal F^{\rm an})$};
 \draw[<->,thick] (a)--node[above,font=\scriptsize] {GAGA} (b);
 \draw[->] (a)--node[left,font=\scriptsize] {kohomologia} (c);
 \draw[->] (b)--node[right,font=\scriptsize] {kohomologia} (d);
 \draw[<->,thick] (c)--(d);
\end{tikzpicture}

{\small Rysunek 9.2. Dwa porównania twierdzenia GAGA.
Górna strzałka oznacza równoważność kategorii, dolna
izomorfizm grup; nie są to identyfikacje definicji.}\par
\end{center}

\begin{stwierdzenie}[Porównanie wiązek liniowych]
Niech $X$ będzie schematem właściwym nad $\C$.
Analityfikacja wyznacza izomorfizm grup
$\operatorname{Pic}(X)\simeq\operatorname{Pic}(X^{\rm an})$.
W szczególności na zespolonej rozmaitości rzutowej
każda holomorficzna wiązka liniowa pochodzi z algebraicznej.
\end{stwierdzenie}
\begin{proof}
Wiązka liniowa to koherentny snop lokalnie wolny rzędu jeden.
GAGA daje algebraiczny snop koherentny $\mathcal L$ odpowiadający
analitycznej wiązce liniowej. Lokalna wolność $\mathcal L$
może być testowana po analityfikacji, ponieważ homomorfizmy
lokalnych pierścieni algebraicznych do holomorficznych
pierścieni lokalnych są wiernie płaskie; jest to standardowa
lokalna część twierdzenia porównawczego. Stąd $\mathcal L$
jest algebraiczną wiązką liniową. Iloczyn tensorowy
i wiązka odwrotna są zachowywane przez analityfikację,
co daje izomorfizm grup Picarda.
\end{proof}

\begin{uwaga}[Gdzie ukrywa się trudność dowodu GAGA?]
Wśród głębokich składników pełnego dowodu znajdują się
analityczna skończoność kohomologii snopów koherentnych
na przestrzeniach zwartych oraz algebraiczne twierdzenia
Serre'a o snopach koherentnych na przestrzeniach rzutowych.
W tym skrypcie udowodnimy ich najważniejsze konsekwencje
na $\PP^1$ bez odwoływania się do pełnej teorii; samą
równoważność GAGA przyjmujemy jako twierdzenie zewnętrzne.
\end{uwaga}

\section{Chow: dlaczego zamknięta geometria analityczna staje się algebraiczna?}
\label{sec:ga9-chow}

\begin{twierdzenie}[Chow]
\label{tw:ga9-chow}
Niech $n\ge0$ będzie liczbą całkowitą, a
$A\subseteq\PP^n(\C)$ podzbiorem zamkniętym
w zwykłej topologii zespolonej, który
jest zespolonym zbiorem analitycznym. Wtedy istnieją
jednorodne wielomiany $F_1,\ldots,F_r\in\C[Z_0,\ldots,Z_n]$
takie, że
$A=\{[Z]\in\PP^n(\C):F_1(Z)=\cdots=F_r(Z)=0\}$.
Zatem $A$ jest rzutowym zbiorem algebraicznym.
\end{twierdzenie}
\begin{proof}[Dowód jako konsekwencja GAGA]
Wprowadzimy jawnie wyniki użyte w przejściu.
Zbiór $A$ ma koherentny analityczny snop ideałów
$\mathcal I_A\subseteq\mathcal O_{\PP^n{}^{\rm an}}$
(twierdzenie Oki--Cartana). Z twierdzenia GAGA
\ref{tw:ga9-gaga} wynika istnienie koherentnego
algebraicznego snopu $\mathcal I$ na $\PP^n_\C$
oraz izomorfizmu $\mathcal I^{\rm an}\simeq\mathcal I_A$.
Ponieważ analityfikacja jest w GAGA w pełni wierna,
inkluzja $\mathcal I_A\hookrightarrow\mathcal O_{\PP^n{}^{\rm an}}$
pochodzi z morfizmu algebraicznego
$\mathcal I\to\mathcal O_{\PP^n}$. Jest to iniekcja:
można to sprawdzić po wiernie płaskim przejściu do
lokalnych pierścieni analitycznych. Tak więc
$\mathcal I$ jest koherentnym snopem ideałów.
Wyznacza on zamknięty podschemat
$Z\hookrightarrow\PP^n_\C$.

Pozostaje wykazać, że algebraiczny podschemat rzutowy
jest opisany przez skończenie wiele jednorodnych
wielomianów. Na mocy twierdzeń Serre'a o snopach
na przestrzeni rzutowej dla dostatecznie dużego $d$
snop $\mathcal I(d)$ jest generowany przez skończenie
wiele sekcji globalnych. Każdą z nich można
interpretować jako jednorodny wielomian stopnia $d$;
gdy potrzebne są różne stopnie, wybieramy skończony
zestaw takich równań we wszystkich mapach współrzędnych
(z noetherowskości i koherencji). Powstający
jednorodny ideał ma jako związany snop ideałów
$\mathcal I$. W szczególności jego zbiór punktów
zespolonych jest dokładnie $A$. To kończy dowód.
\end{proof}

\begin{uwaga}[Status dowodu]
Dowód jest pełnym \emph{wyprowadzeniem Chow z GAGA,
koherencji analitycznej i twierdzeń Serre'a o przestrzeni
rzutowej}. Nie jest elementarnym dowodem Chowa od podstaw:
wymienione trzy rezultaty przyjmujemy jako wyniki zewnętrzne.
Takie rozgraniczenie jest ważne, by nie ukrywać trudności
we frazie „z teorii snopów wynika”.
\end{uwaga}

\begin{wniosek}[Holomorficzne odwzorowania projektowe są algebraiczne]
\label{wn:ga9-hol-map}
Niech $X,Y$ będą zespolonymi rzutowymi rozmaitościami
algebraicznymi (niekoniecznie gładkimi). Każde
odwzorowanie holomorficzne
$f:X^{\rm an}\to Y^{\rm an}$ jest analityfikacją
jednego morfizmu algebraicznego $X\to Y$.
\end{wniosek}
\begin{proof}
Graf $\Gamma_f\subseteq X^{\rm an}\times Y^{\rm an}$
jest zamkniętym zbiorem analitycznym, ponieważ
$Y^{\rm an}$ jest rozdzielona, a lokalnie graf
definiują równania holomorficzne odwzorowania. Ponieważ
$X\times Y$ jest rzutowy (zanurzenie Segrego),
twierdzenie Chowa pokazuje, że $\Gamma_f$ jest
algebraicznym podzbiorem $X\times Y$. Projekcja
$\Gamma_f\to X$ jest izomorfizmem po analityfikacji.
Algebraiczny morfizm właściwy, którego analityfikacja
jest izomorfizmem, jest izomorfizmem (konsekwencja
lokalnej wiernej płaskości i GAGA). Dlatego
projekcja algebraiczna $\Gamma_f\to X$ jest
izomorfizmem. Złożenie jej odwrotności z drugą
projekcją do $Y$ daje poszukiwany morfizm.
Jednoznaczność wynika z pełnej wierności porównania morfizmów
w sytuacji właściwej (lub, dla zredukowanych odmian, z faktu,
że morfizmy do rozdzielonego celu są wyznaczone przez swoje
wartości na gęstym zbiorze punktów domkniętych).
\end{proof}

\begin{przyklad}[Dlaczego nie wolno pominąć rzutowości?]
Holomorficzne odwzorowanie $\C\to\C^*$, $z\mapsto e^z$,
nie jest morfizmem algebraicznym. Morfizm algebraiczny
$\A^1_\C\to\mathbb G_m$ odpowiada jednostce w $\C[z]$,
a jedynymi jednostkami tego pierścienia są niezerowe
stałe. Zatem w tym przypadku wszystkie morfizmy
algebraiczne do $\mathbb G_m$ są stałe, mimo istnienia wielu
niestacjonarnych odwzorowań holomorficznych.
\end{przyklad}

\section[Rachunek porównawczy na P1: kohomologia w obu językach]{Rachunek porównawczy na $\PP^1$: kohomologia w obu językach}
\label{sec:ga9-p1-cech}

Rozpatrzmy standardowe pokrycie $\PP^1=U_0\cup U_1$,
gdzie $U_0\cong\A^1$ ma współrzędną $z=Z_1/Z_0$,
a $U_1\cong\A^1$ ma współrzędną $w=Z_0/Z_1=1/z$.
Na przecięciu dziedzin obu map współrzędnych $z\ne0$.
Niech $e_0,e_1$ będą standardowymi lokalnymi bazami
wiązki $\mathcal O(n)$, przy konwencji $e_1=z^n e_0$.

\begin{stwierdzenie}[Kohomologia algebraiczna $\mathcal O(n)$]
\label{st:ga9-p1-coh}
Niech $n\in\Z$. Dla wiązki liniowej $\mathcal O_{\PP^1_\C}(n)$
na prostej rzutowej nad $\C$ mamy
\[
 h^0(\PP^1,\mathcal O(n))=\max\{n+1,0\},\qquad
 h^1(\PP^1,\mathcal O(n))=\max\{-n-1,0\}.
\]
Ponadto $H^q(\PP^1,\mathcal O(n))=0$ dla $q\ge2$.
\end{stwierdzenie}
\begin{proof}
Część algebraiczna kompleksu Čecha dla pokrycia
$U_0,U_1$ ma postać
\[
 \C[z]\oplus\C[z^{-1}]
 \ \xrightarrow{\ \delta\ }\ \C[z,z^{-1}],\qquad
 \delta(a,b)=a-z^n b.
\]
Znak nie ma wpływu na iloraz, ale określamy go,
aby ustalić konwencję trywializacji. Sekcja globalna
to para $(a,b)$ z $a=z^n b$. Jeśli $n\ge0$,
możemy przyjąć $b=\sum_{j=0}^n b_jz^{-j}$,
co daje wielomiany $a$ stopnia najwyżej $n$;
wymiar przestrzeni sekcji wynosi $n+1$. Gdy $n<0$,
nie ma niezerowej pary spełniającej równanie.

Grupa $H^1$ jest ilorazem
$\C[z,z^{-1}]/(\C[z]+z^n\C[z^{-1}])$.
W mianowniku znajdują się wszystkie potęgi $z^j$
dla $j\ge0$ oraz wszystkie potęgi $z^j$ dla $j\le n$.
Gdy $n\ge-1$, nie pozostaje żadna potęga. Gdy
$n\le-2$, bazę ilorazu tworzą klasy
$z^{n+1},z^{n+2},\ldots,z^{-1}$, których jest $-n-1$.
Ponieważ $U_0,U_1,U_0\cap U_1$ są afiniczne,
kohomologia koherentna pokrycia jest liczona
przez kompleks Čecha. Wymiar $\PP^1$ wynosi jeden,
więc dla snopów koherentnych wyższe grupy $H^q$
znikają; można to również odczytać z dwuczłonowego
kompleksu Čecha i acykliczności tego pokrycia.
\end{proof}

\begin{stwierdzenie}[Rachunek analityczny daje ten sam wynik]
\label{st:ga9-p1-analytic}
Niech $\mathcal O(n)^{\rm an}$ będzie analityfikacją
wiązki $\mathcal O(n)$ na sferze Riemanna
$\PP^1(\C)$. Grupy $H^0$ i $H^1$ tej wiązki
mają wymiary podane w
stwierdzeniu~\ref{st:ga9-p1-coh}.
\end{stwierdzenie}
\begin{proof}
W dziedzinach map analitycznych $U_0^{\rm an}=\C$ oraz
$U_1^{\rm an}=\C$ sekcje lokalne są już funkcjami
całkowitymi, a na przecięciu $\C^*$ -- funkcjami
holomorficznymi. Dla $H^0$ relacja
$a(z)=z^n b(1/z)$ wraz z rozwinięciami Taylora
w zerze i nieskończoności wymusza dokładnie
te same wielomiany co w dowodzie algebraicznym.

Każdą holomorficzną funkcję $g$ na $\C^*$ można
rozwinąć w szereg Laurenta zbieżny na
$0<|z|<\infty$. Część o wykładnikach nieujemnych
przedłuża się do funkcji całkowitej $a(z)$.
Część o wykładnikach $j\le n$ ma postać
$z^n b(1/z)$ z funkcją całkowitą $b$.
Pozostają tylko wykładniki $n<j<0$, jeśli takie
istnieją, a jest ich $-n-1$ dla $n\le-2$.
Zatem iloraz analitycznych $1$-kokykli przez
kobrzegi ma tę samą bazę co iloraz algebraiczny.
Aby utożsamić ten iloraz z analityczną grupą
$H^1$, wykorzystujemy acykliczność Steinowskich
otwartych zbiorów $U_0^{\rm an}$, $U_1^{\rm an}$
i ich przecięcia, które są Steinowskie (twierdzenie Cartana B).
Nie jest to niejawne zastosowanie samego GAGA.
\end{proof}

\begin{przyklad}[Dualność na sferze jako konkretny rachunek]
Dla $n=-3$ otrzymujemy $H^0(\PP^1,\mathcal O(-3))=0$
i $H^1(\PP^1,\mathcal O(-3))$ o bazie reprezentowanej
przez $z^{-2},z^{-1}$. Z drugiej strony
$H^0(\PP^1,\mathcal O(1))$ ma bazę $1,z$.
Parowanie residuowe z rozdziału~\ref{ch:krzywe-snopy}
łączy te przestrzenie w sposób niezdegenerowany;
odpowiada temu izomorfizm dualności Serre'a
$H^1(\mathcal O(-3))^*\cong H^0(\mathcal O(1))$.
Ten przykład pokazuje jednocześnie, jak wyniki GAGA
łączą algebraiczną teorię dzielników z klasycznymi
całkami po okręgach na sferze Riemanna.
\end{przyklad}

\section{Metryki hermitowskie i geometria Kählerowska}
\label{sec:ga9-kahler}

Na rozmaitości zespolonej $M$ mamy rozkład zespolonej
wiązki kostycznej na podwiązki generowane lokalnie przez
$\mathrm dz_1,\ldots,\mathrm dz_n$ i przez
$\mathrm d\bar z_1,\ldots,\mathrm d\bar z_n$.
Formę będącą kombinacją iloczynów zewnętrznych $p$ symboli
$\mathrm dz_j$ oraz $q$ symboli $\mathrm d\bar z_k$
nazywamy formą typu $(p,q)$. Zewnętrzna różniczka
$\mathrm d$ rozkłada się na $\partial$ i $\bar\partial$:
$\partial$ podnosi pierwszy stopień, a $\bar\partial$ drugi.
Na przykład $\bar\partial f=\sum_j
(\partial f/\partial\bar z_j)\,\mathrm d\bar z_j$;
funkcja klasy $C^1$ jest holomorficzna dokładnie wtedy,
gdy $\bar\partial f=0$. To wyjaśnia oznaczenia używane
w konstrukcji formy Fubiniego--Study'ego.

\begin{definicja}[Forma Kählerowska]
Niech $M$ będzie rozmaitością zespoloną wymiaru $n$
z operatorem struktury zespolonej $J:TM\to TM$.
Metryka Riemannowska $g$ jest \emph{hermitowska}, gdy
$g(Ju,Jv)=g(u,v)$ dla wszystkich rzeczywistych
wektorów stycznych $u,v$. Jej związana forma
$\omega(u,v)=g(Ju,v)$ jest rzeczywistą formą różniczkową
stopnia dwa typu $(1,1)$. Metrykę (lub formę)
nazywamy \emph{Kählerowską}, gdy dodatkowo
$\mathrm d\omega=0$. Rozmaitość $M$ nazywamy
\emph{Kählerowską}, jeśli dopuszcza taką metrykę.
\end{definicja}

\begin{przyklad}[Przestrzeń zespolona jako wzorzec]
Na $\C^n$ ze współrzędnymi $z_j=x_j+iy_j$
standardowa metryka euklidesowa ma formę
$\omega_0=\sum_{j=1}^n\mathrm dx_j\wedge\mathrm dy_j
=\frac{i}{2}\sum_{j=1}^n\mathrm dz_j\wedge\mathrm d\bar z_j$.
Jej współczynniki są stałe, więc $\mathrm d\omega_0=0$.
Każda kompleksowa podrozmaitość otwartego podzbioru
$\C^n$ dziedziczy przez restrykcję metrykę Kählerowską.
\end{przyklad}

\begin{twierdzenie}[Forma Fubiniego--Study'ego]
\label{tw:ga9-fs}
Niech $n\ge1$ będzie liczbą całkowitą i ustalmy standardowy iloczyn
hermitowski na $\C^{n+1}$. Na $\PP^n(\C)$ istnieje
wyznaczona przez niego dodatnia,
zamknięta rzeczywista $(1,1)$-forma
$\omega_{\rm FS}$, która na mapie $Z_0\ne0$
o współrzędnych $z_j=Z_j/Z_0$ jest dana wzorem
\[
 \omega_{\rm FS}=\frac{i}{2\pi}\partial\bar\partial
 \log(1+|z_1|^2+\cdots+|z_n|^2).
\]
Jej klasa de Rhama jest obrazem całkowitej klasy
$c_1(\mathcal O_{\PP^n}(1))$ w $H^2(\PP^n(\C),\R)$.
\end{twierdzenie}
\begin{proof}
Na dowolnym mapie rzutowym wybieramy reprezentanta
$Z=(Z_0,\ldots,Z_n)$ i stosujemy potencjał
$\log\sum_j|Z_j|^2$ po ustawieniu jednej współrzędnej
równej $1$. Dwie takie trywializacje różnią się
pomnożeniem reprezentanta przez niezerową funkcję
holomorficzną $h$, więc ich potencjały różnią się
o $\log|h|^2$. Ponieważ lokalnie
$\partial\bar\partial\log|h|^2=0$, otrzymane formy
sklejają się w globalną formę $\omega_{\rm FS}$.
Zamkniętość wynika z $\partial^2=\bar\partial^2=0$
i $\partial\bar\partial+\bar\partial\partial=0$.

Aby sprawdzić dodatniość, połóżmy $A=1+\|z\|^2$.
Macierz współczynników formy (z pominięciem dodatniego
czynnika $1/\pi$) ma postać
\[
 \frac{\partial^2\log A}{\partial z_j\partial\bar z_k}
 =\frac{A\delta_{jk}-\bar z_j z_k}{A^2}.
\]
Dla $v\ne0$ odpowiadająca jej forma hermitowska
jest proporcjonalna do
$\big(A\|v\|^2-|\sum_j\bar z_jv_j|^2\big)/A^2$.
Z nierówności Cauchy'ego--Schwarza dostajemy dolne
oszacowanie przez $\|v\|^2/A^2>0$.

Na końcu podajemy uzasadnienie identyfikacji klasy.
Tautologiczna wiązka $\mathcal O(-1)$ ma w lokalnej
bazie $(1,z_1,\ldots,z_n)$ metrykę hermitowską
o kwadracie normy $A$. Dualna wiązka $\mathcal O(1)$
ma kwadrat normy lokalnego dualnego przekroju równy $A^{-1}$.
Wzór Cherna--Weila
$c_1(L,h)=-(i/2\pi)\partial\bar\partial\log\|e\|^2_h$
dla lokalnej bazy $e$ daje dla $\mathcal O(1)$
właśnie $\omega_{\rm FS}$. Fakt, że forma krzywizny
reprezentuje topologiczną klasę Cherna, jest
standardowym twierdzeniem Cherna--Weila, przyjętym
tutaj jako wynik zewnętrzny. Wszystkie pozostałe
rachunki wykonaliśmy jawnie.
\end{proof}

\begin{przyklad}[Normalizacja całki na $\PP^1$]
\label{prz:ga9-fs-integral}
Dla $n=1$ oraz $z=x+iy$ mamy
$\omega_{\rm FS}=\frac1\pi\frac{\mathrm dx\wedge\mathrm dy}
{(1+|z|^2)^2}$. Punkt w nieskończoności nie wnosi
dodatkowego składnika do całki po zwartej sferze,
bo jest zbiorem miary zero. Zatem we współrzędnych
biegunowych
\[
 \int_{\PP^1(\C)}\omega_{\rm FS}
 =\frac1\pi\int_0^{2\pi}\int_0^\infty
 \frac{r}{(1+r^2)^2}\,\mathrm dr\,\mathrm d\theta
 =2\int_0^\infty\frac{r}{(1+r^2)^2}\,\mathrm dr=1.
\]
W ten sposób sprawdziliśmy normalizację:
$\int_{\PP^1}\omega_{\rm FS}=\deg\mathcal O(1)=1$.
Ten rachunek wiąże rzutowy stopień algebraiczny
z całką po rzeczywistej powierzchni wymiaru dwa.
\end{przyklad}

\begin{wniosek}[Każda gładka rozmaitość rzutowa jest Kählerowska]
Jeśli $X\hookrightarrow\PP^N_\C$ jest gładką, zamkniętą
podrozmaitością algebraiczną, to $X^{\rm an}$ jest
zwartą rozmaitością Kählerowską; jej formą
Kählerowską może być restrykcja $\omega_{\rm FS}$.
\end{wniosek}
\begin{proof}
Zanurzenie algebraiczne indukuje holomorficzne
zanurzenie analityczne. Forma Kählerowska na
$\PP^N(\C)$ jest zamknięta i dodatnia na
każdej zespolonej podprzestrzeni stycznej, więc
jej restrykcja do $X^{\rm an}$ zachowuje obie własności.
Zwartość wynika z domkniętości $X^{\rm an}$
w zwartej przestrzeni $\PP^N(\C)$.
\end{proof}

\section{Kiedy rozmaitość Kählerowska jest algebraiczna?}
\label{sec:ga9-kodaira}

\begin{twierdzenie}[Twierdzenie Kodairy o zanurzeniu — wynik zewnętrzny]
\label{tw:ga9-kodaira}
Niech $M$ będzie zwartą rozmaitością zespoloną.
Wtedy $M$ dopuszcza holomorficzne zanurzenie jako
zamknięta podrozmaitość pewnej $\PP^N(\C)$ wtedy
i tylko wtedy, gdy istnieje holomorficzna wiązka
liniowa $L\to M$ z dodatnio określoną rzeczywistą
formą krzywizny $(1,1)$ pewnej metryki hermitowskiej.
Taka wiązka nazywa się \emph{dodatnią}.
\end{twierdzenie}

\begin{uwaga}[Dlaczego wynik nie jest tautologią?]
Z jednej strony restrykcja $\mathcal O_{\PP^N}(1)$
na zamkniętą podrozmaitość projektową jest dodatnia.
W drugą stronę potrzebujemy wykazać istnienie tylu
sekcji $L^{\otimes m}$, by rozdzielały wszystkie
punkty oraz kierunki styczne. To jest istotna treść
twierdzenia Kodairy, oparta na teorii analitycznej
$\bar\partial$ i odpowiednich twierdzeniach o znikaniu.
Nie wyprowadzamy jej tutaj wyłącznie z dodatniości
krzywizny, nie dowodząc tych wyników pośrednich.
\end{uwaga}

\begin{wniosek}[Kryterium całkowitej klasy Kählerowskiej]
\label{wn:ga9-integral-kahler}
Zwarta rozmaitość zespolona $M$ jest projektowa
w sensie holomorficznego zanurzenia do przestrzeni
rzutowej wtedy i tylko wtedy, gdy ma formę Kählerowską
$\omega$, której klasa de Rhama $[\omega]$ należy
do obrazu $H^2(M,\Z)\to H^2(M,\R)$.
\end{wniosek}
\begin{proof}
Jeśli $M$ jest projektowa, restrykcja znormalizowanej
formy Fubiniego--Study'ego jest Kählerowska, a jej
klasa jest całkowita, bo jest pierwszą klasą Cherna
wiązki $\mathcal O(1)|_M$.
W drugą stronę twierdzenie o wykładniczym ciągu snopów,
połączone z faktem, że rzeczywista klasa $[\omega]$
jest typu $(1,1)$, daje holomorficzną wiązkę liniową
$L$ z $c_1(L)=[\omega]$ w kohomologii rzeczywistej.
Dzięki lematowi $\partial\bar\partial$ na zwartej
rozmaitości Kählerowskiej można dobrać metrykę
hermitowską na $L$, której znormalizowana krzywizna
jest dokładnie $\omega$. Jest więc ona dodatnia,
a twierdzenie Kodairy daje zanurzenie rzutowe.
W bardziej szczegółowym rozumowaniu całkowitą klasę w\,~$H^2(M,\Z)$
(reprezentującą $[\omega]$ po przejściu do $H^2(M,\R)$)
uznajemy za klasę typu $(1,1)$. Twierdzenie Lefschetza o klasach
$(1,1)$, wyprowadzone z ciągu wykładniczego i rozkładu Hodge'a,
dostarcza wtedy $L$. Lemat $\partial\bar\partial$ zmienia
metrykę na $L$ tak, by jej forma Cherna była równa $\omega$.
Oba te kroki opierają się na twierdzeniach zewnętrznych,
które dokładniej omówimy w następnym rozdziale.
\end{proof}

\begin{przyklad}[Zwarte tory zespolone: Kähler nie oznacza rzutowości]
Niech $\Lambda\subset\C^n$ będzie kratą rzędu $2n$,
a $T=\C^n/\Lambda$ zespolonym torusem. Standardowa
forma $\omega_0$ jest niezmiennicza względem translacji,
więc przechodzi na dodatnią zamkniętą formę na $T$:
każdy torus zespolony jest Kählerowski.
Jednak dla $n\ge2$ istnieją tory zespolone bez
jakiejkolwiek całkowitej klasy Kählerowskiej,
a zatem nieprojektowe. Na bardzo ogólnym torusie
zespolonym o wymiarze dwa zachodzi nawet
$H^{1,1}(T)\cap H^2(T,\Z)=\{0\}$ (po przejściu do
kohomologii bez torsji). Istnienie takich struktur
można wykazać argumentem z przeliczalnością warunków
na macierz okresów; szczegółowy dowód tego istnienia
pozostawiamy jako wynik zewnętrzny.
\end{przyklad}

\section{Hodge, Dolbeault i niezmienniki topologiczne}
\label{sec:ga9-hodge}

\begin{twierdzenie}[Dolbeault — wynik zewnętrzny]
Niech $M$ będzie zwartą rozmaitością zespoloną.
Jeśli $\Omega^p_M$ oznacza snop holomorficznych
$p$-form, to dla każdego $p,q\ge0$ istnieje
naturalny izomorfizm
$H^q(M,\Omega^p_M)\simeq H^{p,q}_{\bar\partial}(M)$,
gdzie prawa strona jest kohomologią kompleksu
gładkich form typu $(p,\bullet)$ z operatorem $\bar\partial$.
\end{twierdzenie}

\begin{twierdzenie}[Rozkład Hodge'a dla zwartych rozmaitości Kählerowskich]
\label{tw:ga9-hodge}
Niech $M$ będzie zwartą rozmaitością Kählerowską.
Wtedy dla każdego $k\ge0$ zachodzi naturalny rozkład
\[
 H^k_{\rm dR}(M;\C)
 =\bigoplus_{p+q=k}H^{p,q}_{\bar\partial}(M),
 \qquad \overline{H^{p,q}}=H^{q,p}.
\]
W szczególności, jeśli $h^{p,q}=\dim_\C H^{p,q}_{\bar\partial}(M)$,
to $b_k=\sum_{p+q=k}h^{p,q}$ i $h^{p,q}=h^{q,p}$.
\end{twierdzenie}

\begin{uwaga}[O co trzeba uzupełnić dowód?]
Kluczowe są tożsamości Kählerowskie, z których
wynika zależność $\Delta_{\rm d}=2\Delta_{\bar\partial}$
(zgodnie ze standardową normalizacją laplasjanów).
Dzięki temu harmoniczna forma de Rhama rozkłada się
na składowe typu $(p,q)$, również harmoniczne.
Pełny dowód tożsamości Kählerowskich i teorii
harmonicznych reprezentantów korzysta z analizy
eliptycznej; w tym rozdziale przyjmujemy go
jako wynik zewnętrzny, a poniżej dowodzimy jego
konkretnych konsekwencji.
\end{uwaga}

\begin{wniosek}[Nieparzyste liczby Bettiego są parzyste]
\label{wn:ga9-odd-betti}
Jeśli $M$ jest zwartą rozmaitością Kählerowską,
to dla każdego $r\ge0$ liczba Bettiego $b_{2r+1}(M)$
jest parzysta. W szczególności $b_1(M)=2h^{1,0}(M)$.
\end{wniosek}
\begin{proof}
Z rozkładu Hodge'a
$b_{2r+1}=\sum_{p+q=2r+1}h^{p,q}$.
Ponieważ $p+q$ jest nieparzyste, nie może być $p=q$.
Składniki możemy zatem pogrupować w rozłączne pary
$(p,q)$ i $(q,p)$ o tych samych wymiarach. Każda para
wnosi liczbę parzystą. Dla $k=1$ pozostają wyłącznie
$(1,0)$ oraz $(0,1)$, co daje drugi wzór.
\end{proof}

\begin{przyklad}[Dlaczego nie wszystkie zwarte rozmaitości zespolone są Kählerowskie?]
Klasyczna powierzchnia Hopfa otrzymywana jako iloraz
$\C^2\setminus\{0\}$ przez działanie
$(z,w)\mapsto(2z,2w)$ jest zwartą rozmaitością zespoloną.
Jako rozmaitość gładka jest dyfeomorficzna z
$S^1\times S^3$, więc jej pierwsza liczba Bettiego
wynosi $b_1=1$. Na mocy
wniosku~\ref{wn:ga9-odd-betti} nie może być Kählerowska.
W szczególności nie jest projektowa. Zwartość
i zespolona struktura nie wymuszają więc algebraiczności.
Identyfikację topologiczną z $S^1\times S^3$
można zobaczyć, zapisując $\C^2\setminus\{0\}
\cong S^3\times\R$ przez promień $\|z\|$ i kierunek
$z/\|z\|$; mnożenie przez $2$ przesuwa współrzędną
$\log\|z\|$ o $\log2$.
\end{przyklad}

\begin{przyklad}[Hodge na krzywej eliptycznej]
Niech $E$ będzie gładką rzutową krzywą zespoloną
rodzaju jeden. Z rozdziału~\ref{ch:krzywe-snopy}
wiemy, że $h^0(E,\Omega_E^1)=1$, a dualność
Serre'a daje $h^1(E,\mathcal O_E)=1$.
Wymiary Hodge'a wynoszą
$h^{0,0}=h^{1,1}=1$, $h^{1,0}=h^{0,1}=1$.
Stąd $b_0=1$, $b_1=2$, $b_2=1$, a
$\chi_{\rm top}(E)=b_0-b_1+b_2=0$.
Warto porównać ten wynik z zerową charakterystyką
Eulera torusa rzeczywistego $S^1\times S^1$.
\end{przyklad}

\begin{twierdzenie}[Twarde Lefschetza — wynik zewnętrzny]
\label{tw:ga9-hard-lefschetz}
Niech $M$ będzie zwartą rozmaitością Kählerowską
o wymiarze zespolonym $n$, a $L$ operatorem
$\alpha\mapsto[\omega]\smile\alpha$ w kohomologii.
Dla każdego $0\le k\le n$ mnożenie przez
$[\omega]^{n-k}$ indukuje izomorfizm
$L^{n-k}:H^k(M;\R)\to H^{2n-k}(M;\R)$.
\end{twierdzenie}

\begin{wniosek}[Ograniczenia topologii gładkich odmian rzutowych]
Niech $X$ będzie gładką rzutową rozmaitością nad $\C$
o wymiarze zespolonym $n$. Jej nieparzyste liczby
Bettiego są parzyste, a dla $k<n$ zachodzi
$b_k(X)\le b_{k+2}(X)$. Ponadto klasy podrozmaitości
algebraicznych kodujące przecięcia są zgodne z
mnożeniem przez klasę Kählerowską.
\end{wniosek}
\begin{proof}
Pierwszy wniosek otrzymujemy z rozkładu Hodge'a.
Aby wykazać nierówność, użyjmy twierdzenia
\ref{tw:ga9-hard-lefschetz}. Jeśli $k<n$, a
$[\omega]\smile\alpha=0$ dla klasy
$\alpha\in H^k(M;\R)$, to tym bardziej
$[\omega]^{n-k}\smile\alpha=0$. Twarde Lefschetza
pokazuje, że $\alpha=0$. Mnożenie przez $[\omega]$
jest więc iniekcją $H^k\hookrightarrow H^{k+2}$.
Zgodność dla klas podrozmaitości wynika z
interpretacji $c_1(\mathcal O(1))$ jako klasy
hiperpłaszczyzny i z formuły Cherna--Weila;
jej ogólny aspekt enumeratywny jest częścią
twierdzeń porównawczych teorii przecięć i
kohomologii, przyjętych tutaj zewnętrznie.
\end{proof}

\section{Rachunek łączący stopień, całkę i rodzaj krzywej}
\label{sec:ga9-degree-genus}

\begin{stwierdzenie}[Stopień krzywej przez całkę]
Niech $C\hookrightarrow\PP^2_\C$ będzie gładką,
spójną, rzutową krzywą płaską stopnia $d\ge1$.
Wówczas
$\int_{C^{\rm an}}\omega_{\rm FS}=d$,
gdzie całkujemy restrykcję znormalizowanej formy
Fubiniego--Study'ego. Ponadto
$2g(C)-2=d(d-3)$.
\end{stwierdzenie}
\begin{proof}
Pierwszą klasą Cherna wiązki $\mathcal O_C(1)$
jest restrykcja klasy hiperpłaszczyzny $h$.
Przez utożsamienie jej z $[\omega_{\rm FS}]$
oraz definicję stopnia wiązki liniowej otrzymujemy
\[
 \int_{C^{\rm an}}\omega_{\rm FS}
 =\langle c_1(\mathcal O_C(1)),[C]\rangle
 =\deg\mathcal O_C(1)=d.
\]
W ostatnim kroku przecięcie z ogólną prostą w
$\PP^2$ ma $d$ punktów, liczonych z krotnościami,
co wynika z Bézouta.

Ze wzoru adjunkcji na gładkiej krzywej mamy
$\omega_C\cong
(\omega_{\PP^2}\otimes\mathcal O_{\PP^2}(C))|_C
\cong\mathcal O_C(d-3)$,
ponieważ $\omega_{\PP^2}\cong\mathcal O_{\PP^2}(-3)$.
Stopień wiązki kanonicznej $C$ wynosi $2g(C)-2$,
a $\deg\mathcal O_C(d-3)=d(d-3)$. Po przyrównaniu
stopni dostajemy drugi wzór.
\end{proof}

\begin{przyklad}[Trzy stopnie, trzy topologie]
Dla $d=1,2,3$ gładkie spójne krzywe płaskie mają
odpowiednio rodzaj $g=0,0,1$. W każdym przypadku
całka restrykcji $\omega_{\rm FS}$ wynosi jednak
odpowiednio $1,2,3$. Stopień zanurzenia i rodzaj
krzywej mierzą różne aspekty geometrii: pierwszy
zależy od wiązki $\mathcal O_C(1)$, drugi od
wewnętrznej kanoniki $K_C$. Z tego powodu tę samą
krzywą abstrakcyjną można zanurzać w przestrzeniach
rzutowych na różne sposoby, uzyskując inne stopnie.
\end{przyklad}

\section{Stein, Cartan i granica między afinicznością a analitycznością}
\label{sec:ga9-stein}

\begin{definicja}[Rozmaitość Steinowska]
Zespoloną rozmaitość $M$ nazywamy \emph{Steinowską},
gdy jest holomorficznie wypukła i holomorficznie
separowalna: dla każdego zwartego $K\subset M$
jego otoczka holomorficzna
$\widehat K_{\mathcal O(M)}=
 \{x\in M:|f(x)|\le\sup_{y\in K}|f(y)|\ \forall f\in\mathcal O(M)\}$
jest zwarta, a dowolne dwa różne punkty $x,y\in M$
można rozdzielić globalną funkcją holomorficzną
$f\in\mathcal O(M)$ z $f(x)\ne f(y)$.
\end{definicja}

\begin{twierdzenie}[Afiniczne odmiany nad $\C$ są Steinowskie — wynik standardowy]
\label{tw:ga9-affine-stein}
Jeśli $X$ jest afiniczną rozmaitością algebraiczną
nad $\C$, to $X^{\rm an}$ jest zespoloną
przestrzenią Steinowską (w przypadku gładkim:
rozmaitością Steinowską).
\end{twierdzenie}
\begin{proof}[Uzasadnienie przez zanurzenie do przestrzeni afinicznej]
Ponieważ $X$ jest afiniczna i skończonego typu,
można ją przedstawić jako domknięty podschemat
$\A^N_\C$. Zatem jej analityfikacja jest
zamkniętą przestrzenią analityczną w $\C^N$.
Przestrzeń $\C^N$ jest Steinowska, a zamknięte
podprzestrzenie analityczne przestrzeni Steinowskich
są Steinowskie. Ostatni krok jest standardowym
twierdzeniem analitycznym i nie wynika tylko
z samej definicji wypukłości.
\end{proof}

\begin{twierdzenie}[Cartan B — wynik zewnętrzny]
Jeżeli $X$ jest przestrzenią Steinowską,
a $\mathcal F$ koherentnym analitycznym snopem
$\mathcal O_X$-modułów, to
$H^q(X,\mathcal F)=0$ dla wszystkich $q\ge1$.
\end{twierdzenie}

\begin{uwaga}[Algebraiczna i analityczna afiniczność to nie to samo]
Dla afinicznego schematu $X$ oraz quasi-koherentnego
snopu $\mathcal F$ zachodzi algebraiczne
$H^q(X,\mathcal F)=0$ dla $q>0$. Dla $X^{\rm an}$
dostajemy analityczny odpowiednik przez Cartana B.
Zgodność dwóch twierdzeń w tej sytuacji nie oznacza
jednak GAGA: już na $\A^1$ grupy globalnych funkcji
regularnych i holomorficznych są różne.
W szczególności analityczność grafu funkcji $e^z$
nie czyni go algebraicznym podzbiorem $\A^2$.
Nie wolno też wnioskować bez dodatkowych założeń,
że każda rozmaitość algebraiczna, której analityfikacja
jest Steinowska, musi być afiniczna.
\end{uwaga}

\section{Synteza: kiedy teoria algebraiczna i analityczna się zgadzają?}
\label{sec:ga9-podsumowanie}

Najbardziej użyteczne rozróżnienie można zebrać
w czterech przypadkach. Na przestrzeni afinicznej
algebraiczne funkcje regularne są zwykle tylko częścią
funkcji holomorficznych. Na przestrzeni rzutowej
zamknięte zbiory analityczne są algebraiczne dzięki
Chow. Dla przestrzeni właściwych nad $\C$ snopy
koherentne i ich kohomologia są porównywalne dzięki
GAGA. Dla gładkich rozmaitości rzutowych geometria
Fubiniego--Study'ego udostępnia dodatkowo narzędzia
Kählera, Hodge'a i Lefschetza.

\begin{uwaga}[Co jest dowiedzione, a co przyjęte?]
W tym rozdziale od podstaw wykazaliśmy niealgebraiczność
grafu wykładniczego, lokalność pierścienia kiełków,
rachunek Čecha na $\PP^1$, dodatniość formy
Fubiniego--Study'ego, jej całkę na $\PP^1$,
parzystość nieparzystych liczb Bettiego przy
założeniu twierdzenia Hodge'a oraz związki
stopnia krzywej z całkami i kanoniką.
Chow został wyprowadzony z jawnie wskazanych
silniejszych wyników: GAGA, koherencji Oki--Cartana
i twierdzeń Serre'a. Zewnętrzne pozostają m.in.:
przygotowanie Weierstrassa, GAGA, Cartan B,
Cherna--Weil, tożsamości Kählerowskie,
Kodaira, Hodge i twarde Lefschetza.
\end{uwaga}

% Literatura podstawowa i źródła kontroli:
% J.-P. Serre, Géométrie algébrique et géométrie analytique (GAGA), 1956.
% P. Griffiths, J. Harris, Principles of Algebraic Geometry.
% C. Voisin, Hodge Theory and Complex Algebraic Geometry I.
% R. O. Wells, Differential Analysis on Complex Manifolds.
% The Stacks Project, Complex Analytic Spaces and GAGA:
% https://stacks.math.columbia.edu/


% ROZDZIAL 10. GEOMETRIA ALGEBRAICZNA A TOPOLOGIA
% Zaleznosci: rozdzialy 2, 4, 6, 7 oraz 9.

\chapter{Geometria algebraiczna a topologia: od cykli do twierdzeń Lefschetza}
\label{ch:ga10-topologia}

\providecommand{\atlasref}[2]{#2}


\section{Dwa sposoby patrzenia na tę samą przestrzeń}
\label{sec:ga10-punkt-wyjscia}

Równanie wielomianowe wyznacza przestrzeń algebraiczną, ale nad $\C$
można badać również jej zwykłą topologię. Dla gładkiej odmiany
rzutowej $X$ nad $\C$ analityfikacja $X^{\rm an}$ jest zwartą,
zorientowaną rozmaitością gładką wymiaru rzeczywistego $2n$,
gdzie $n=\dim_{\C}X$. Orientację wyznaczają zespolone układy
współrzędnych: baza $(v_1,Jv_1,\ldots,v_n,Jv_n)$ jest dodatnio
zorientowana. To właśnie orientacja umożliwia całkowanie form
najwyższego stopnia i definiowanie liczb przecięć.

Chcemy odpowiedzieć na trzy pytania: jak z równań wielomianowych
obliczać grupy kohomologii; dlaczego klasy podrozmaitości tworzą
szczególne klasy kohomologiczne; jakie topologie są niemożliwe
dla gładkich odmian rzutowych?

W całym rozdziale $H^k(M;R)$ oznacza kohomologię singularną
przestrzeni topologicznej $M$ ze współczynnikami w pierścieniu $R$.
Jeśli piszemy $H^k_{\rm dR}(M)$, chodzi o kohomologię kompleksu
gładkich form różniczkowych o rzeczywistych współczynnikach.
Wymiary $b_k(M)=\dim_{\Q}H^k(M;\Q)$ nazywamy \emph{liczbami Bettiego}.
Dla zwartej rozmaitości są one skończone. Zapis $H^q(X,\mathcal F)$
nadal oznacza \emph{kohomologię snopa}; nie utożsamiamy tych dwóch
pojęć bez użycia odpowiedniego twierdzenia porównawczego.

\begin{przyklad}[Krzywa algebraiczna jako powierzchnia topologiczna]
Gładka spójna krzywa rzutowa $C$ nad $\C$ jest zwartą, spójną,
zorientowaną powierzchnią rzeczywistą. Jej rodzaj $g$ oznacza liczbę
uchwytów, a jej liczby Bettiego wynoszą $b_0=b_2=1$ i $b_1=2g$.
Dlatego $\chi_{\rm top}(C)=2-2g$; wyjaśnimy poniżej, jak ta
równość wynika zarówno z topologii, jak i z rachunku dzielników.
\end{przyklad}

\section{Kohomologia de Rhama, dualność Poincarégo i cykle algebraiczne}
\label{sec:ga10-de-rham}

\noindent\emph{Powiązania tematyczne w Bibliotece matematycznej:} \atlasref{TOP-09::thm:de-rham}{TOP-09: twierdzenie de Rhama}; \atlasref{TOP-11::thm:poincare-dualnosc-16}{TOP-11: dualność Poincarégo}.


\begin{definicja}[Kohomologia de Rhama]
Niech $M$ będzie gładką rozmaitością rzeczywistą.
Dla $k\ge0$ przez $\Omega^k(M)$ oznaczamy przestrzeń
rzeczywistych gładkich $k$-form, a przez $\mathrm d\colon\Omega^k(M)
\to\Omega^{k+1}(M)$ różniczkę zewnętrzną. Ponieważ $\mathrm d^2=0$,
definiujemy
\[
 H^k_{\rm dR}(M)=\frac{\ker(\mathrm d\colon\Omega^k(M)\to
 \Omega^{k+1}(M))}{\operatorname{im}(\mathrm d\colon
 \Omega^{k-1}(M)\to\Omega^k(M))},
\]
z umową $\Omega^{-1}(M)=0$. Formy z licznika nazywamy
\emph{zamkniętymi}, a formy z mianownika \emph{dokładnymi}.
\end{definicja}

\begin{uwaga}[Dlaczego iloraz ma sens geometryczny?]
Formy dokładne mają zerową całkę po każdym cyklu bez brzegu:
na mocy twierdzenia Stokesa $\int_Z\mathrm d\beta
=\int_{\partial Z}\beta=0$. Wobec tego całka po cyklu
nie rozróżnia reprezentantów jednej klasy de Rhama.
Twierdzenie de Rhama (wynik zewnętrzny) stwierdza,
że całkowanie po cyklach, po utożsamieniu przestrzeni
funkcjonałów na homologii z kohomologią singularną
ze współczynnikami rzeczywistymi, wyznacza naturalny
izomorfizm $H^k_{\rm dR}(M)\simeq H^k(M;\R)$.
Jest to pomost między rachunkiem różniczkowym a topologią.
\end{uwaga}

\begin{twierdzenie}[Dualność Poincarégo — wynik zewnętrzny]
\label{tw:ga10-pd}
Niech $M$ będzie zwartą, spójną i zorientowaną rozmaitością
gładką wymiaru rzeczywistego $m$, bez brzegu.
Dla każdego $0\le k\le m$ parowanie
\[
 H^k(M;\R)\times H^{m-k}(M;\R)\longrightarrow\R,
 \qquad ([\alpha],[\beta])\longmapsto\int_M\alpha\wedge\beta
\]
jest niezdegenerowane. W szczególności $b_k(M)=b_{m-k}(M)$.
\end{twierdzenie}

\begin{proof}[Dowód ostatniej równości z niezdegenerowania]
Formy gładkie reprezentują klasy rzeczywiste przez twierdzenie
de Rhama. Ponieważ obie przestrzenie kohomologii są skończenie
wymiarowe, niezdegenerowane parowanie identyfikuje
$H^k(M;\R)$ z przestrzenią dualną $H^{m-k}(M;\R)^*$.
Przestrzeń i jej dualna mają ten sam wymiar.
Główna część twierdzenia, czyli konstrukcja dualności,
pozostaje wynikiem topologii algebraicznej przyjętym zewnętrznie.
\end{proof}

\begin{definicja}[Klasa dualna do podrozmaitości]
Niech $M$ będzie zwartą, spójną, zorientowaną rozmaitością
gładką wymiaru $m$ bez brzegu, $p$ liczbą całkowitą
z $0\le2p\le m$, a $Z\subseteq M$ zwartą, zorientowaną,
gładką podrozmaitością wymiaru $m-2p$.
\emph{Klasą dualną Poincarégo} do $Z$ nazywamy
jedyną klasę $\operatorname{PD}[Z]\in H^{2p}(M;\R)$ taką, że
$\int_Z\eta=\int_M\operatorname{PD}[Z]\wedge\eta$
dla każdej zamkniętej formy $\eta$ stopnia $m-2p$.
Istnienie i jednoznaczność wynikają z dualności Poincarégo.
\end{definicja}

\begin{uwaga}[W przypadku osobliwym]
Cykle algebraiczne mogą być osobliwe, więc nie zawsze są
podrozmaitościami. Dla domkniętej pododmiany zespolonej czystego wymiaru
$p$ w zwartej rozmaitości zespolonej istnieje jednak
kanoniczna klasa fundamentalna w homologii singularnej
wymiaru $2p$. Konstrukcja w ogólnym przypadku
wymaga teorii cykli lub homologii Borel–Moore'a; w dalszych
obliczeniach wystarczą gładkie pododmiany i liniowość cykli.
\end{uwaga}

\begin{stwierdzenie}[Intersekcja jako całka]
\label{st:ga10-intersection}
Niech $S$ będzie gładką spójną rzutową powierzchnią nad $\C$,
a $C,D\subset S$ dwiema gładkimi krzywymi przecinającymi się
transwersalnie w skończenie wielu punktach. Wtedy
\[
 C\cdot D=\#(C\cap D)
 =\int_{S^{\rm an}}\operatorname{PD}[C]\wedge\operatorname{PD}[D].
\]
\end{stwierdzenie}
\begin{proof}
W lokalnych współrzędnych holomorficznych wokół każdego
punktu przecięcia krzywe są dwiema transwersalnymi
zespolonymi prostymi. Orientacje zespolone wymuszają dodatni
lokalny znak przecięcia, czyli $+1$. Suma lokalnych znaków
równa się algebraicznej liczbie przecięć, którą dla dwóch
transwersalnych krzywych obliczaliśmy już w
rozdziale~\ref{ch:ga6-powierzchnie}. Z twierdzenia o zgodności
iloczynu kohomologicznego i przecięcia geometrycznego
(zewnętrzny wynik o dualności Poincarégo i klasach przecięć) ta suma jest
równa całce z iloczynu klas dualnych.
\end{proof}

\begin{przyklad}[Dwie proste w $\PP^2$]
Dla dwóch różnych prostych $L_1,L_2\subset\PP^2$ mamy
$L_1\cdot L_2=1$ na mocy Bézouta. Jeśli
$h=\operatorname{PD}[L_1]=\operatorname{PD}[L_2]$,
to $\int_{\PP^2}h^2=1$. Z jednej liczby otrzymujemy
podstawowy rachunek w pierścieniu kohomologii $\PP^2$.
\end{przyklad}

\section{Topologia przestrzeni rzutowej i krzywych}
\label{sec:ga10-projektywne}

\begin{twierdzenie}[Kohomologia przestrzeni rzutowej — wynik standardowy]
\label{tw:ga10-cp-cohom}
Dla $n\ge0$ zachodzi izomorfizm stopniowanych pierścieni
\[
 H^*(\PP^n(\C);\Z)\simeq\Z[h]/(h^{n+1}),
 \qquad \deg h=2,
\]
gdzie $h=c_1(\mathcal O_{\PP^n}(1))$ i
$\int_{\PP^n(\C)}h^n=1$.
\end{twierdzenie}
\begin{proof}[Idea dowodu przez strukturę komórkową]
Dla $n=0$ mamy $\PP^0(\C)$ będącą jednym punktem,
więc $H^*(\PP^0(\C);\Z)=\Z$, a $h=0$ i
$\int_{\PP^0(\C)}h^0=\int_{\PP^0(\C)}1=1$.
Dalej zakładamy $n\ge1$.
Przestrzeń $\PP^n(\C)$ otrzymujemy z
$\PP^{n-1}(\C)$, doklejając jedną rzeczywistą komórkę
wymiaru $2n$: mapa $Z_n\ne0$ jest zespoloną przestrzenią
afiniczną, a jego dopełnienie to $\PP^{n-1}(\C)$.
Zatem $\PP^n(\C)$ ma po jednej komórce
w wymiarach $0,2,\ldots,2n$ i nie ma komórek
w nieparzystych wymiarach. W komórkowym kompleksie
łańcuchowym wszystkie różniczki są zerowe,
ponieważ przynajmniej jedna z sąsiednich grup jest zerowa.
Stąd $H^{2j}(\PP^n;\Z)\cong\Z$ dla $0\le j\le n$,
a pozostałe grupy dodatnich stopni znikają.
Identyfikacja mnożenia wynika z tego, że hiperpłaszczyzna
reprezentuje $h$, a przecięcie $j$ hiperpłaszczyzn w
położeniu ogólnym jest liniową przestrzenią rzutową
kowymiaru $j$; dla $j=n$ przecięcie jest jednym punktem.
Tym samym $h^j$ generuje $H^{2j}$ i $h^n$ ma całkę równą $1$.
Użyliśmy standardowej zgodności dualności Poincarégo
z iloczynem przecięć z poprzedniej sekcji.
\end{proof}

\begin{stwierdzenie}[Rodzaj krzywej z kohomologii]
Niech $C$ będzie gładką spójną krzywą rzutową nad $\C$
o rodzaju $g$. Wtedy $\chi_{\rm top}(C)=2-2g$ oraz
$\dim_{\C}H^0(C,\Omega^1_C)=g$.
\end{stwierdzenie}
\begin{proof}
Pierwszy wzór wynika z klasyfikacji zwartych zorientowanych
powierzchni (twierdzenie topologiczne): powierzchnia rodzaju $g$
ma rozkład CW z jednym wierzchołkiem, $2g$ krawędziami
i jedną komórką dwuwymiarową. Dlatego jej charakterystyka
Eulera wynosi $1-2g+1=2-2g$.
Drugi wzór otrzymujemy przez identyfikację
$H^0(C,\Omega^1_C)$ ze składnikiem Hodge'a $H^{1,0}$.
Dla krzywej o rodzaju $g$ liczba $b_1$ wynosi $2g$,
a symetria Hodge'a daje $b_1=h^{1,0}+h^{0,1}=2h^{1,0}$.
W szczególności $h^{1,0}=g$, zgodnie z Riemannem--Rochem
z rozdziału~\ref{ch:krzywe-snopy}.
\end{proof}

\begin{center}
\begin{tikzpicture}[>=Stealth,every node/.style={font=\footnotesize}]
 \node[draw=blue!55!black,fill=blue!5,rounded corners=5pt,
       minimum height=18mm,text width=43mm,align=left,inner sep=6pt]
       (curve) at (-3.25,0)
       {Krzywa $C$ rodzaju $g$\\[2pt]
        $b_1=2g$\\[2pt]
        $\chi_{\rm top}=2-2g$};
 \node[draw=green!45!black,fill=green!6,rounded corners=5pt,
       minimum height=18mm,text width=43mm,align=left,inner sep=6pt]
       (forms) at (3.25,0)
       {Formy holomorficzne\\[2pt]
        $h^{1,0}=g$\\[2pt]
        $h^{0,1}=g$};
 \draw[->,thick] (curve.east)--node[above,font=\scriptsize]{Hodge}(forms.west);
\end{tikzpicture}

{\small Rysunek 10.1. Rzeczywista topologia i holomorficzne różniczki
krzywej rzutowej niosą tę samą informację o rodzaju.}\par
\end{center}

\section{Słabe twierdzenie Lefschetza: co zachowuje przekrój hiperpłaszczyzną?}
\label{sec:ga10-weak}

\noindent\emph{Powiązania tematyczne w Bibliotece matematycznej:} \atlasref{TOP-28::ch:xvii-topologia-odmian}{TOP-28: słabe twierdzenie Lefschetza}.


% Moduł wspólny, wykorzystywany w geometrii i topologii algebraicznej.
% Dwa historyczne identyfikatory odsyłaczy odpowiadają jednemu twierdzeniu.
% Guard zapobiega podwójnemu składowi w książce łączącej oba rozdziały.
\ifcsname biblioteka@slabelefschetz@loaded\endcsname
Twierdzenie wraz z dowodem zostało już sformułowane wcześniej
(w twierdzeniu~\ref{tw:ga10-weak}).
\else
\expandafter\gdef\csname biblioteka@slabelefschetz@loaded\endcsname{}
\begin{twierdzenie}[Słabe twierdzenie Lefschetza]
\label{tw:ga10-weak}\label{thm:xvii-weak-lefschetz}
Niech $X\hookrightarrow\mathbb P^N_{\mathbb C}$ będzie gładką,
spójną odmianą projektową nad $\mathbb C$ o wymiarze zespolonym
$n\ge 2$. Niech $H\subset\mathbb P^N_{\mathbb C}$ będzie hiperpłaszczyzną,
a $Y=X\cap H$ niepustym gładkim przekrojem, z zanurzeniem
$i\colon Y\hookrightarrow X$. Wówczas homomorfizm kohomologii singularnej
\[
 i^*\colon H^q(X^{\mathrm{an}};\mathbb Z)
 \longrightarrow H^q(Y^{\mathrm{an}};\mathbb Z)
\]
jest izomorfizmem dla $q<n-1$ oraz monomorfizmem dla $q=n-1$.
Wersja kohomologiczna z dowolnymi współczynnikami w ciele
wynika z tej samej konstrukcji. W szczególności $Y$ jest spójna. Dla odpowiednio ustalonych
punktów bazowych indukowane odwzorowanie
$\pi_j(Y^{\mathrm{an}})\to\pi_j(X^{\mathrm{an}})$
jest izomorfizmem dla $1\le j<n-1$ oraz epimorfizmem dla $j=n-1$.
\end{twierdzenie}
\begin{proof}[Redukcja do geometrycznego twierdzenia Lefschetza]
Korzystamy z następującego \emph{głębokiego wyniku zewnętrznego}:
przestrzeń $X^{\mathrm{an}}$ ma typ homotopijny przestrzeni
otrzymanej z $Y^{\mathrm{an}}$ przez dołączenie komórek
rzeczywistego wymiaru co najmniej $n$. Dowód tego wyniku
wymaga zespolonej teorii Morse'a i funkcji plurisubharmonicznych
(lub odpowiedniej teorii znikania). Nie wyprowadzamy go z samej
gładkości $X$.

Względny kompleks komórkowy pary $(X^{\mathrm{an}},Y^{\mathrm{an}})$
nie ma komórek w stopniach mniejszych niż $n$, więc
$H_q(X^{\mathrm{an}},Y^{\mathrm{an}};\mathbb Z)=0$ dla $q<n$.
Względna aproksymacja komórkowa daje również
$\pi_j(X^{\mathrm{an}},Y^{\mathrm{an}})=0$ dla $j<n$.
Długi ciąg dokładny homotopii pary zapewnia stąd podane
izomorfizmy i epimorfizm. Należy przy tym pamiętać, że
w niskich stopniach grupy homotopii wymagają właściwej
interpretacji z ustalonymi punktami bazowymi.

Analogicznie $H^q(X^{\mathrm{an}},Y^{\mathrm{an}};\mathbb Z)=0$
dla $q<n$. W ciągu dokładnym kohomologii pary
\[
 \cdots\longrightarrow H^q(X,Y;\mathbb Z)
 \longrightarrow H^q(X;\mathbb Z)
 \xrightarrow{i^*} H^q(Y;\mathbb Z)
 \longrightarrow H^{q+1}(X,Y;\mathbb Z)\longrightarrow\cdots
\]
(przez $X,Y$ w tym ciągu rozumiemy przestrzenie analityczne)
dla $q<n-1$ znikają obie grupy względne sąsiadujące z $i^*$,
a dla $q=n-1$ znika grupa z lewej. Otrzymujemy odpowiednio
izomorfizm i monomorfizm. Ponieważ $n\ge2$, dla $q=0$
otrzymujemy izomorfizm $H^0(X;\mathbb Z)\simeq H^0(Y;\mathbb Z)$,
a spójność $X$ daje spójność $Y$.
\end{proof}
\fi


\begin{uwaga}[Skąd biorą się granice $n-1$?]
Jedna droga do dowodu prowadzi przez teorię Morse'a:
w odpowiednio dobranej funkcji na dopełnieniu hiperpłaszczyzny
pojawiają się uchwyty rzeczywistego wymiaru $n$.
Oznacza to, że dołączenie takich komórek nie zmienia
niskich grup homotopii. Jest to jedynie idea dowodu:
uzasadnienie istnienia właściwej funkcji Morse'a,
a następnie konstrukcja odpowiedniego rozkładu CW,
są zasadniczą treścią twierdzenia zewnętrznego.
\end{uwaga}

\begin{stwierdzenie}[Gładkie powierzchnie w $\PP^3$ są po prostu spójne]
\label{st:ga10-simplyconn}
Jeżeli $S\subset\PP^3_\C$ jest gładką zespoloną
hiperpowierzchnią stopnia $d\ge1$, to
$\pi_1(S^{\rm an})=0$ i $b_1(S)=b_3(S)=0$.
\end{stwierdzenie}
\begin{proof}
Stosujemy wersję homotopijną
Lefschetza~\ref{tw:ga10-weak} do pary
$X=\PP^3_\C$, $Y=S$, dla $n=3$.
Otrzymujemy izomorfizm $\pi_1(S)\simeq\pi_1(\PP^3)$.
Ta ostatnia grupa jest trywialna: $\PP^3$ ma strukturę
CW z jedną komórką w każdym parzystym wymiarze od $0$ do $6$,
a więc jej jednoszkielet jest punktem, co daje trywialność
$\pi_1$. Zatem $H_1(S;\Z)=0$, więc $b_1(S)=0$.
Ponieważ $S$ jest zwartą zorientowaną rozmaitością
rzeczywistą wymiaru cztery, dualność Poincarégo daje
$b_3(S)=b_1(S)=0$.
\end{proof}

\begin{uwaga}[Słabe a twarde Lefschetza]
Słabe twierdzenie porównuje \emph{dwie różne przestrzenie}:
$Y$ i $X$. Twarde twierdzenie Lefschetza z
rozdziału~\ref{ch:ga9-analiza-zespolona} dotyczy natomiast
mnożenia przez klasę Kählerowską \emph{na tej samej przestrzeni}.
Nie wolno z jednego wyniku bezpośrednio wywnioskować drugiego.
\end{uwaga}

\section{Klasy Cherna i topologiczna charakterystyka Eulera}
\label{sec:ga10-chern-euler}

\begin{twierdzenie}[Chern--Gauss--Bonnet — wynik zewnętrzny]
\label{tw:ga10-cgb}
Jeśli $M$ jest zwartą rozmaitością zespoloną wymiaru $n$,
to dla jej holomorficznej wiązki stycznej $T_M$ zachodzi
\[
 \chi_{\rm top}(M)=\int_M c_n(T_M).
\]
Prawa strona oznacza ocenę całkowitej klasy Cherna
$c_n(T_M)\in H^{2n}(M;\Z)$ na klasie fundamentalnej.
\end{twierdzenie}

\begin{uwaga}[Dlaczego ta tożsamość jest zaskakująca?]
Lewa strona jest niezmiennikiem topologicznym,
liczonym np. z rozkładu komórkowego. Prawa strona
pochodzi z lokalnej geometrii wiązki stycznej.
Przejście między nimi jest szczególnym przypadkiem
uogólnionego twierdzenia Gaussa--Bonneta.
Sam rachunek klas Cherna wykonamy już algebraicznie.
\end{uwaga}

\begin{twierdzenie}[Charakterystyka Eulera gładkiej powierzchni stopnia $d$]
\label{tw:ga10-surface-euler}
Niech $S\subset\PP^3_\C$ będzie gładką powierzchnią
określoną jednym jednorodnym równaniem stopnia $d\ge1$,
a $K_S$ dowolnym dzielnikiem kanonicznym na $S$.
Wtedy
\[
 \chi_{\rm top}(S)=d(d^2-4d+6),\qquad
 K_S^2=d(d-4)^2.
\]
\end{twierdzenie}
\begin{proof}
Niech $h=c_1(\mathcal O_{\PP^3}(1))$ i oznaczmy przez ten sam
symbol jego ograniczenie do $S$.
Z poprzedniego rozdziału o klasach Cherna znamy ciąg dokładny
wiązki stycznej i normalnej
$0\to T_S\to T_{\PP^3}|_S\to\mathcal O_S(d)\to0$.
Z multiplikatywności całkowitej klasy Cherna i z ciągu Eulera
otrzymujemy
\[
 c(T_S)=\frac{(1+h)^4}{1+dh}
 =1+(4-d)h+(d^2-4d+6)h^2,
\]
przy czym pomijamy składniki stopnia wyższego niż $4$,
ponieważ $S$ ma zespolony wymiar dwa.
Na mocy Chern--Gauss--Bonnet
$\chi_{\rm top}(S)=\int_S c_2(T_S)$.
Ze wzoru na stopień hiperpowierzchni
$\int_S h^2=d$ (przecięcie z dwiema ogólnymi
hiperpłaszczyznami) dostajemy pierwszy wzór.

Ponadto przez adjunkcję
$\omega_S\cong\mathcal O_S(d-4)$, więc
$c_1(\omega_S)=(d-4)h$ i
$K_S^2=\int_S (d-4)^2h^2=d(d-4)^2$.
\end{proof}

\begin{przyklad}[Pięć stopni i pięć różnych niezmienników]
Dla $d=1,2,3,4,5$ otrzymujemy kolejno
$\chi_{\rm top}(S)=3,4,9,24,55$.
Wartości $K_S^2$ wynoszą odpowiednio $9,8,3,0,5$.
Gdy $d=4$, kanonika jest trywialna, a mimo to
$\chi_{\rm top}(S)=24$; różnica między tymi niezmiennikami
pokazuje, że znikanie $c_1(T_S)$ nie oznacza znikania
$c_2(T_S)$.
\end{przyklad}

\begin{stwierdzenie}[Liczby Bettiego tych powierzchni]
\label{st:ga10-betti}
Niech $S\subseteq\PP^3_\C$ będzie gładką powierzchnią
stopnia $d\ge1$, określoną jednorodnym wielomianem.
Wówczas
\[
 b_0=b_4=1,\qquad b_1=b_3=0,\qquad
 b_2=d(d^2-4d+6)-2.
\]
\end{stwierdzenie}
\begin{proof}
Spójność daje $b_0=1$, a dualność Poincarégo $b_4=1$.
Stwierdzenie~\ref{st:ga10-simplyconn} daje $b_1=b_3=0$.
Z definicji charakterystyki Eulera
$\chi_{\rm top}(S)=b_0-b_1+b_2-b_3+b_4=2+b_2$,
więc wynik otrzymujemy przez podstawienie
wzoru z twierdzenia~\ref{tw:ga10-surface-euler}.
\end{proof}

\begin{przyklad}[Sześcienna powierzchnia i kwartyk K3]
Gładka powierzchnia kubiczna ma $b_2=7$,
co odpowiada $\PP^2$ rozdmuchanej w sześciu punktach.
Gładka powierzchnia stopnia czwartego ma $b_2=22$.
Dla niej $K_S\cong\mathcal O_S$, a słabe Lefschetza
zapewnia $\pi_1(S)=0$; otrzymujemy więc powierzchnię K3.
Sam niezmiennik $b_2=22$ nie identyfikuje jeszcze całej
struktury zespolonej ani izomorfizmu algebraicznego.
\end{przyklad}

\section{Rozdmuchanie punktu: co dzieje się z topologią?}
\label{sec:ga10-blowup}

\noindent\emph{Powiązania tematyczne w Bibliotece matematycznej:} \atlasref{TOP-28::thm:xvii-blowup-general}{TOP-28: kohomologia rozdmuchania}.


% Wspólna część topologiczna: kohomologia, forma przecięć,
% liczby Bettiego, charakterystyka Eulera i sygnatura.
% Dwie historyczne etykiety identyfikują jedno twierdzenie.
\ifcsname biblioteka@rozdmuchanie@kohomologia@loaded\endcsname
Kohomologię rozdmuchania punktu obliczono wcześniej
w twierdzeniu~\ref{thm:xvii-blowup}.
\else
\expandafter\gdef\csname biblioteka@rozdmuchanie@kohomologia@loaded\endcsname{}
\begin{twierdzenie}[Kohomologia i forma przecięć rozdmuchania punktu]
\label{tw:ga10-blowup}\label{thm:xvii-blowup}
Niech $S$ będzie gładką, zwartą, spójną zespoloną powierzchnią,
$p\in S$, $\pi:\widetilde S=\operatorname{Bl}_pS\to S$ rozdmuchaniem,
$E=\pi^{-1}(p)$ oraz $e=\operatorname{PD}_{\widetilde S}[E]\in
H^2(\widetilde S;\mathbb Z)$ (symbol $\operatorname{PD}$ oznacza dualność
Poincarégo). Wówczas
\[
H^2(\widetilde S;\mathbb Z)=\pi^*H^2(S;\mathbb Z)\oplus\mathbb Z e.
\]
Dla każdej $\alpha\in H^2(S;\mathbb Z)$ zachodzi
$\langle \pi^*\alpha\smile e,[\widetilde S]\rangle=0$ i
$\langle e^2,[\widetilde S]\rangle=-1$.
W innych stopniach $\pi^*:H^k(S;\mathbb Z)\to
H^k(\widetilde S;\mathbb Z)$ jest izomorfizmem.
Przez $\sigma(M)$ rozumiemy sygnaturę rzeczywistej formy
przecięć na $H^2(M;\mathbb R)$; $b_2$ oznacza drugą liczbę
Bettiego, a $\chi_{\mathrm{top}}$ charakterystykę Eulera.
W szczególności
\[
 b_2(\widetilde S)=b_2(S)+1,\qquad
 \chi_{\mathrm{top}}(\widetilde S)=\chi_{\mathrm{top}}(S)+1,
 \qquad \sigma(\widetilde S)=\sigma(S)-1.
\]
\end{twierdzenie}
\begin{proof}
Wybierzmy małą zamkniętą kulę $B\subset S$ o środku $p$,
$U=S\setminus\operatorname{int}B$, i niech
$\widetilde B=\pi^{-1}(B)$. Przestrzeń $\widetilde B$ jest
otoczeniem dyskowym $E\simeq\mathbb P^1(\mathbb C)\simeq S^2$
i deformacyjnie retraktuje się na $E$. Brzegi $B$ oraz
$\widetilde B$ są identyfikowane przez $\pi$ ze sferą $S^3$.
Zatem $S=U\cup B$ i $\widetilde S=U\cup\widetilde B$
mają na przecięciach przestrzeń homotopijnie równoważną $S^3$.

Z ciągów Mayera--Vietorisa (ze współczynnikami $\mathbb Z$)
i równości $H^1(S^3)=H^2(S^3)=0$ dostajemy
\[
 H^2(S;\mathbb Z)\simeq H^2(U;\mathbb Z),\qquad
 H^2(\widetilde S;\mathbb Z)
 \simeq H^2(U;\mathbb Z)\oplus H^2(E;\mathbb Z).
\]
Pierwsza składowa ostatniego izomorfizmu jest obrazem $\pi^*$,
ponieważ $\pi|_{\widetilde B}$ ma obraz w ściągalnej kuli $B$.
Ograniczenie $e$ do $E$ jest klasą
$c_1(N_{E/\widetilde S})=-\eta$, gdzie
$\eta$ jest dodatnim generatorem $H^2(E;\mathbb Z)$;
wynika to z lematu~\ref{lem:xvii-exceptional}.
Dlatego $e$ generuje dodatkowy składnik całkowity.

Dla jasności sprawdzamy także stopień trzeci.
Z dualności Poincarégo--Lefschetza lub z długiego ciągu pary
$(U,\partial U)$ wynika, że homomorfizm łączący
$H^3(\partial U;\mathbb Z)\to H^4(U,\partial U;\mathbb Z)$
jest izomorfizmem: generator sfery granicznej przechodzi
na względną klasę orientacji spójnej rozmaitości $U$.
Stąd $H^3(U;\mathbb Z)\to H^3(\partial U;\mathbb Z)$
jest zerowy. Porównanie obu ciągów Mayera--Vietorisa
pokazuje, że $\pi^*$ jest izomorfizmem również w stopniu trzecim.
Stopnie zero i jeden porównujemy bezpośrednio z ciągów,
a w stopniu czwartym $\pi$ ma stopień $1$, więc
$\pi^*:H^4(S;\mathbb Z)\to H^4(\widetilde S;\mathbb Z)$
jest izomorfizmem (obie grupy są izomorficzne z $\mathbb Z$).

Ortogonalność wynika ze wzoru projekcji i tego, że
$\pi|_E$ jest odwzorowaniem stałym:
\[
 \langle\pi^*\alpha\smile e,[\widetilde S]\rangle
 =\langle (\pi|_E)^*\alpha,[E]\rangle=0.
\]
Samoprzecięcie $e^2=-1$ daje wspomniany lemat.
Ponieważ $\pi$ ma stopień $1$, forma przecięć na
$\pi^*H^2(S;\mathbb R)$ jest identyczna z formą na $S$.
Otrzymujemy więc rozkład ortogonalny
$Q_{\widetilde S}\cong Q_S\oplus(-1)$ na wolnych częściach
kohomologii. To daje $\sigma(\widetilde S)=\sigma(S)-1$.
Zmiana wyłącznie drugiej liczby Bettiego powoduje
$\chi_{\mathrm{top}}(\widetilde S)=\chi_{\mathrm{top}}(S)+1$.
\end{proof}
\fi


\begin{wniosek}[Kwadrat klasy kanonicznej po rozdmuchaniu]
\label{wn:ga10-blowup-canonical}
Niech $S$ będzie gładką spójną rzutową powierzchnią nad $\C$,
$p\in S(\C)$, $\pi:\widetilde S\to S$ rozdmuchaniem punktu,
a $K_S,K_{\widetilde S}$ dowolnymi dzielnikami kanonicznymi
odpowiednich powierzchni. Wówczas
$K_{\widetilde S}^2=K_S^2-1$.
\end{wniosek}
\begin{proof}
Algebraiczny wzór $K_{\widetilde S}\sim\pi^*K_S+E$ jest
udowodniony w stwierdzeniu~\ref{st:ga6-Kblowup}.
Ponieważ $E^2=-1$ i $\pi^*K_S\cdot E=0$, mamy
$K_{\widetilde S}^2=(\pi^*K_S+E)^2=K_S^2-1$.
\end{proof}

\begin{uwaga}[Drugi dowód zmiany charakterystyki Eulera]
Można też skorzystać z addytywności charakterystyki Eulera
z kohomologii o zwartych nośnikach dla rozwarstwienia algebraicznego.
Rozdmuchanie jest izomorfizmem poza punktem $p$, a nad nim zastępuje
punkt krzywą $E\simeq\mathbb P^1(\mathbb C)$, czyli sferą $S^2$.
Dla zespolonych odmian algebraicznych charakterystyka Eulera
z kohomologii o zwartych nośnikach zgadza się ze zwykłą.
Zatem
$\chi_{\rm top}(\widetilde S)=\chi_{\rm top}(S\setminus\{p\})
+\chi_{\rm top}(E)=\chi_{\rm top}(S)-1+2$.
Jest to alternatywny rachunek względem dowodu przez
Mayera--Vietorisa z twierdzenia~\ref{thm:xvii-blowup}.
\end{uwaga}


\begin{uwaga}[Inna postać twierdzenia]
Rozdmuchanie odpowiada różniczkowo-topologicznie sumie
spójnej $S^{\rm an}\#\overline{\PP^2(\C)}$,
gdzie kreska oznacza odwrócenie orientacji.
To standardowy wynik topologii czterowymiarowej,
nie utożsamienie rozdmuchania z algebraiczną sumą dwóch
odmian. Nowy składnik formy przecięć ma znak ujemny,
co jest algebraicznym śladem odwrócenia orientacji.
\end{uwaga}

\begin{przyklad}[Sześć rozdmuchań]
Dla $S=\PP^2$ wiemy, że $b_2(S)=1$ i
$\chi_{\rm top}(S)=3$. Po rozdmuchaniu sześciu
różnych punktów w położeniu ogólnym otrzymujemy
$b_2=7$, $\chi_{\rm top}=9$, $K^2=9-6=3$.
Są to dokładnie liczby gładkiej powierzchni kubicznej.
Uzyskujemy więc topologiczną kontrolę wcześniejszego
rachunku 27 prostych: cała geometria przecięć rozwija się
w siedmiowymiarowej przestrzeni klas dzielników.
\end{przyklad}

\begin{center}
\begin{tikzpicture}[>=Stealth,every node/.style={font=\footnotesize}]
 \draw[rounded corners,draw=blue!60!black,fill=blue!4] (-4.5,-1.2) rectangle (-1.65,1.0);
 \node at (-3.1,.62) {$S$};
 \draw[fill=blue!65] (-3.2,-.16) circle (1.8pt);
 \node[anchor=west,font=\scriptsize] at (-3.08,-.16) {$p$};
 \draw[->,thick] (-1.42,-.09)--(.02,-.09)
  node[midway,above,font=\scriptsize] {$\operatorname{Bl}_p$};
 \draw[rounded corners,draw=blue!60!black,fill=blue!4] (.3,-1.2) rectangle (4.45,1.0);
 \node at (2.4,.65) {$\widetilde S$};
 \draw[thick,draw=red!70!black] (1.76,-.32) ellipse (.58 and .30);
 \node[anchor=west,font=\scriptsize] at (2.55,-.13) {$E\simeq\PP^1$};
 \node[anchor=west,font=\scriptsize] at (2.55,-.51) {$E^2=-1$};
 \node[font=\scriptsize] at (0,-1.62) {$b_2(\widetilde S)=b_2(S)+1$};
\end{tikzpicture}

{\small Rysunek 10.2. Rozdmuchanie zastępuje punkt krzywą wyjątkową.
Schemat nie przedstawia lokalnego kształtu rzeczywistej czterowymiarowej
rozmaitości; wyróżnia jedynie zmianę danych topologicznych.}\par
\end{center}

\section{Twarde Lefschetza, formy pierwotne i twierdzenie Hodge'a o indeksie}
\label{sec:ga10-hard-index}

\noindent\emph{Powiązania tematyczne w Bibliotece matematycznej:} \atlasref{TOP-28::thm:xvii-hard-lefschetz}{TOP-28: silne twierdzenie Lefschetza}.


W rozdziale~\ref{ch:ga9-analiza-zespolona} poznaliśmy
twarde twierdzenie Lefschetza. Teraz wyprowadzimy z niego
ograniczenia formy przecięć na powierzchni.

\begin{definicja}[Operator Lefschetza i kohomologia pierwotna]
Niech $M$ będzie zwartą rozmaitością Kählerowską
o wymiarze zespolonym $n$ z formą Kählerowską $\omega$.
Oznaczamy $L(\alpha)=[\omega]\smile\alpha$.
Dla $0\le k\le n$ definiujemy
\[
 P^k(M)=\ker\bigl(L^{n-k+1}\colon
 H^k(M;\C)\to H^{2n-k+2}(M;\C)\bigr).
\]
Elementy $P^k(M)$ nazywamy \emph{klasami pierwotnymi}.
Dla powierzchni $S$ i $k=2$ dostajemy po prostu
$P^2(S)=\{\alpha\in H^2(S;\C):[\omega]\smile\alpha=0\}$.
\end{definicja}

\begin{twierdzenie}[Rozkład Lefschetza — konsekwencja twardego Lefschetza]
Niech $M$ będzie zwartą rozmaitością Kählerowską wymiaru
zespolonego $n$ z formą $\omega$ i niech
$L(\alpha)=[\omega]\smile\alpha$. Dla $0\le j\le n$
niech $P^j(M)=\ker(L^{n-j+1}:H^j(M;\C)\to
H^{2n-j+2}(M;\C))$. Wówczas dla każdego $0\le k\le n$
zachodzi rozkład
$H^k(M;\C)=\bigoplus_{r\ge0}L^rP^{k-2r}(M)$,
z pominięciem składników o ujemnych wykładnikach.
\end{twierdzenie}
\begin{proof}[Szkic argumentu z jawnym wynikiem zewnętrznym]
Twarde Lefschetza zapewnia izomorfizmy
$L^{n-j}\colon H^j\to H^{2n-j}$.
Z nich, przez rozkład na jądro odpowiedniej potęgi $L$
i obraz operatora $L$ z poprzedniego stopnia,
otrzymuje się sumę $H^k=P^k\oplus LH^{k-2}$.
Indukcja po $k$ daje zapis w tezie. Pełne uzasadnienie
prostopadłości i jednoznaczności tego rozkładu jest
standardowym wnioskiem z teorii reprezentacji $\mathfrak{sl}_2$
operatorów Lefschetza. Ponieważ nie wprowadzaliśmy tej
reprezentacji, tę część traktujemy jako przyjętą
konsekwencję twardego twierdzenia Lefschetza.
\end{proof}

\begin{twierdzenie}[Hodge'a o indeksie — wynik zewnętrzny]
\label{tw:ga10-index}
Niech $S$ będzie gładką spójną rzutową powierzchnią
nad $\C$, niech $H$ będzie dzielnikiem obszernym, a
$\rho(S)=\dim_{\R}N^1(S)_{\R}$ jej liczbą Picarda.
Wówczas forma przecięć na rzeczywistej przestrzeni klas
dzielników modulo równoważność numeryczną
$N^1(S)_{\R}$ ma sygnaturę $(1,\rho(S)-1)$.
Równoważnie, jeśli $D\in N^1(S)_{\R}$,
$D\cdot H=0$ i $D\ne0$, to $D^2<0$.
\end{twierdzenie}

\begin{uwaga}[Skąd bierze się ujemny znak?]
Klasa dzielnika ma typ Hodge'a $(1,1)$.
Jeżeli jest prostopadła do dodatniej klasy $[H]$, to jest
klasą pierwotną stopnia dwa. Relacje dwuliniowe
Hodge'a--Riemanna, przyjęte tutaj jako wynik zewnętrzny,
mówią, że na rzeczywistych pierwotnych klasach $(1,1)$
forma $\alpha\mapsto\int_S\alpha\wedge\alpha$ jest ujemnie
określona. Jedyny dodatni kierunek jest reprezentowany
przez $[H]$, gdyż $H^2>0$. Związek tego rozumowania
z klasami algebraicznymi wyjaśnia twierdzenie Lefschetza
o klasach $(1,1)$ z kolejnej sekcji.
\end{uwaga}

\begin{stwierdzenie}[Nierówność Hodge'a dla dzielników]
\label{st:ga10-hodge-ineq}
Niech $S$ będzie gładką spójną rzutową powierzchnią
nad $\C$, $H$ dzielnikiem obszernym na $S$,
a $N^1(S)_{\R}$ rzeczywistą przestrzenią klas dzielników
modulo równoważność numeryczną. Dla każdego
$D\in N^1(S)_{\R}$ zachodzi
\[
 (D\cdot H)^2\ge D^2H^2,
\]
a równość występuje dokładnie wtedy, gdy $D$ jest
rzeczywistą wielokrotnością $H$ w $N^1(S)_{\R}$.
\end{stwierdzenie}
\begin{proof}
Połóżmy $a=(D\cdot H)/H^2$ i $D_0=D-aH$.
Z konstrukcji $D_0\cdot H=0$.
Twierdzenie Hodge'a o indeksie daje $D_0^2\le0$,
a równość tylko dla $D_0=0$. Po rozwinięciu
$D_0^2=D^2-(D\cdot H)^2/H^2$ otrzymujemy tezę.
\end{proof}

\begin{przyklad}[Dlaczego forma przecięć nie jest dodatnio określona?]
Na rozdmuchaniu $S=\operatorname{Bl}_p\PP^2$
przestrzeń $N^1(S)_{\R}$ ma bazę $H,E$,
gdzie $H$ to klasa przeciwobrazu prostej i $E$ jest dzielnikiem
wyjątkowym. Macierz formy przecięć wynosi
$\left(\begin{smallmatrix}1&0\\0&-1\end{smallmatrix}\right)$.
W szczególności $E^2=-1$. Sama efektywność krzywej
nie implikuje więc nieujemnego samoprzecięcia.
Nie wolno jednak w powyższym twierdzeniu podstawiać
$H$ jako klasę przeciwobrazu prostej: taki $H$ jest nef, ale
nie jest obszerny, gdyż $H\cdot E=0$.
Aby uzyskać klasę obszerną, należy wziąć np.
$A=2H-E$, dla której $A^2=3$ oraz $A\cdot E=1$.
\end{przyklad}

\section{Twierdzenie Lefschetza o klasach \texorpdfstring{$(1,1)$}{(1,1)} i problem Hodge'a}
\label{sec:ga10-11}

\begin{definicja}[Grupy Picarda i Nérona--Severiego]
Niech $X$ będzie gładką spójną rzutową odmianą nad $\C$.
Grupa $\operatorname{Pic}(X)$ to grupa klas izomorfizmu
wiązek liniowych, z działaniem indukowanym przez iloczyn tensorowy.
Podgrupa $\operatorname{Pic}^0(X)$ składa się z klas
wiązek algebraicznie równoważnych wiązce trywialnej;
grupę ilorazową
$\operatorname{NS}(X)=\operatorname{Pic}(X)/\operatorname{Pic}^0(X)$
nazywamy grupą Nérona--Severiego.
Jej rząd części wolnej to liczba Picarda $\rho(X)$.
\end{definicja}

\begin{twierdzenie}[Lefschetz $(1,1)$ — wynik zewnętrzny]
\label{tw:ga10-11}
Niech $X$ będzie gładką spójną rzutową odmianą nad $\C$.
Wówczas obraz homomorfizmu
$c_1\colon\operatorname{Pic}(X)\to H^2(X^{\rm an};\Z)$
jest dokładnie zbiorem tych klas całkowitych,
których obraz w $H^2(X^{\rm an};\C)$ należy do
$H^{1,1}(X)$, to znaczy
\[
 \operatorname{im}c_1
 =\{\alpha\in H^2(X^{\rm an};\Z):
 \alpha_{\C}\in H^{1,1}(X)\}.
\]
\end{twierdzenie}

\begin{proof}[Wyjaśnienie przez ciąg wykładniczy]
W topologii analitycznej mamy dokładny ciąg snopów
$0\to\Z_X\to\mathcal O_X\xrightarrow{\exp(2\pi i\,\cdot)}
\mathcal O_X^{\times}\to1$.
Pochodzący z niego fragment długiego ciągu kohomologii
ma postać
\[
 H^1(X,\mathcal O_X^{\times})\xrightarrow{c_1}
 H^2(X;\Z)\longrightarrow H^2(X,\mathcal O_X).
\]
Pierwsza grupa jest analityczną grupą Picarda,
a przez GAGA również algebraiczną grupą Picarda.
Z dokładności wynika, że obraz $c_1$ jest jądrem
strzałki do $H^2(X,\mathcal O_X)$.
Na zwartej rozmaitości Kählerowskiej
$H^2(X,\mathcal O_X)\simeq H^{0,2}(X)$ przez Dolbeaulta.
Ponieważ dla całkowitej klasy jej składowe typu $(2,0)$
i $(0,2)$ są ze sobą sprzężone, zerowanie składowej
$(0,2)$ jest równoważne temu, że klasa ma czysty typ $(1,1)$.
To dowodzi tezy po przyjęciu dokładności ciągu wykładniczego,
GAGA i rozkładu Hodge'a, wyjaśniając zarazem użyte wcześniej
w kryterium Kodairy twierdzenie o klasach $(1,1)$.
\end{proof}

\begin{wniosek}[Ograniczenie liczby Picarda]
\label{wn:ga10-rho}
Dla gładkiej spójnej rzutowej powierzchni $S$ nad $\C$
zachodzi $1\le\rho(S)\le h^{1,1}(S)$.
\end{wniosek}
\begin{proof}
Klasa obszerna, na przykład $c_1(\mathcal O_S(1))$,
jest całkowita i ma niezerowy kwadrat, więc
$\rho(S)\ge1$.
Twierdzenie~\ref{tw:ga10-11} zanurza rzeczywistą
przestrzeń klas dzielników modulo numeryczną równoważność
w rzeczywistą część $H^{1,1}(S)$.
Jej wymiar rzeczywisty jest równy $h^{1,1}(S)$,
co daje nierówność z góry.
\end{proof}

\begin{przyklad}[K3: dwie różne liczby opisujące kohomologię]
Dla gładkiej kwartykowej powierzchni K3 mamy
$b_2=22$. Pokażemy dalej, że $h^{1,1}=20$.
Jednak liczba Picarda może być znacznie mniejsza:
dla bardzo ogólnej gładkiej kwartykowej powierzchni
w $\PP^3$ zachodzi $\rho=1$ (twierdzenie
Noethera--Lefschetza, wynik zewnętrzny).
A zatem istnieje wiele klas typu $(1,1)$ nad $\C$,
które nie są klasami wymiernych kombinacji dzielników.
To nie przeczy twierdzeniu Lefschetza $(1,1)$:
warunek \emph{całkowitości} jest niezbędny.
\end{przyklad}

\begin{definicja}[Cykl algebraiczny i klasa Hodge'a]
Niech $X$ będzie gładką spójną rzutową odmianą nad $\C$
o wymiarze $n$, a $p$ liczbą całkowitą z $0\le p\le n$.
Dla cyklu algebraicznego $Z$ kowymiaru $p$ jego obraz
$\operatorname{cl}(Z)\in H^{2p}(X^{\rm an};\Z)$
nazywamy klasą cyklu. \emph{Wymierną klasą Hodge'a}
stopnia $2p$ nazywamy element
$\alpha\in H^{2p}(X;\Q)$, którego obraz zespolony
należy do $H^{p,p}(X)$.
\end{definicja}

\begin{uwaga}[Hipoteza Hodge'a — czego nie możemy twierdzić?]
Każdy cykl algebraiczny kowymiaru $p$ daje klasę
w $H^{p,p}(X)$; wynika to ze zgodności orientacji
zespolonej ze stopniami form. Odwrotne stwierdzenie,
że każda wymierna klasa Hodge'a jest wymierną kombinacją
klas cykli algebraicznych, jest \emph{hipotezą Hodge'a}
i nie jest znane w ogólności. Dla $p=1$ jest twierdzeniem
Lefschetza~\ref{tw:ga10-11}. Nie należy mylić
rozwiązania dla dzielników z rozwiązaniem
w dowolnym kowymiarze.
\end{uwaga}

\section{Diament Hodge'a i pełne rachunki dla hiperpowierzchni}
\label{sec:ga10-hodge-numbers}

Dla zwartej powierzchni Kählerowskiej $S$ zapisujemy
$h^{p,q}(S)=\dim_{\C}H^q(S,\Omega^p_S)$.
Symetria sprzężenia daje $h^{p,q}=h^{q,p}$,
a dualność Serre'a i Poincarégo prowadzi do
$h^{p,q}=h^{2-p,2-q}$. Zwyczajowo przedstawia się te
liczby w \emph{diamencie Hodge'a}.

\begin{stwierdzenie}[Diament Hodge'a gładkiej powierzchni stopnia $d$]
\label{st:ga10-hodge-hypersurf}
Niech $S\subset\PP^3_\C$ będzie gładką powierzchnią
stopnia $d\ge1$. Niech $p_g=h^0(S,\omega_S)$.
Wtedy $q=h^1(S,\mathcal O_S)=0$ oraz
\[
 p_g=\begin{cases}0,&1\le d\le3,\\
 \binom{d-1}{3},&d\ge4.\end{cases}
\]
Ponadto $h^{1,1}(S)=b_2(S)-2p_g$, gdzie
$b_2(S)$ jest dane wzorem ze
stwierdzenia~\ref{st:ga10-betti}.
\end{stwierdzenie}
\begin{proof}
Z $b_1(S)=0$ i symetrii Hodge'a mamy
$0=b_1=h^{1,0}+h^{0,1}=2h^{0,1}$.
Twierdzenie Dolbeaulta identyfikuje
$h^{0,1}$ z $\dim H^1(S,\mathcal O_S)=q$,
więc $q=0$.

Przez adjunkcję $\omega_S\simeq\mathcal O_S(d-4)$.
Jeśli $d<4$, wiązka $\mathcal O_S(d-4)$ nie ma
niezerowych przekrojów, ponieważ ma ujemny stopień
na każdej ogólnej sekcji hiperpłaszczyznowej.
Gdy $d\ge4$, korzystamy z ciągu dokładnego
\[
 0\to\mathcal O_{\PP^3}(-4)\longrightarrow
 \mathcal O_{\PP^3}(d-4)\longrightarrow
 \mathcal O_S(d-4)\to0.
\]
Stosujemy standardowe twierdzenie o kohomologii wiązek
$\mathcal O(m)$ na przestrzeni rzutowej (w tym miejscu
przyjęte zewnętrznie w wymiarze trzy), które daje
$H^1(\PP^3,\mathcal O(-4))=0$. Wtedy dokładność
długiego ciągu kohomologii pokazuje, że każdy przekrój
kanoniki jest restrykcją jednorodnego wielomianu
stopnia $d-4$. Przekrojów tych jest
$\binom{(d-4)+3}{3}=\binom{d-1}{3}$.

Wreszcie dla stopnia dwa rozkład Hodge'a ma postać
$b_2=h^{2,0}+h^{1,1}+h^{0,2}$.
Z symetrii Hodge'a oraz identyfikacji
$h^{2,0}=h^0(S,K_S)=p_g$ dostajemy
$h^{1,1}=b_2-2p_g$.
\end{proof}

\begin{przyklad}[Diament K3]
Dla gładkiej kwartykowej powierzchni $S$ w $\PP^3$
mamy $p_g=1$, $q=0$, $b_2=22$ i $h^{1,1}=20$.
Jej diament Hodge'a to
\[
 \begin{array}{ccccc}
 &&1&&\\
 &0&&0&\\
 1&&20&&1\\
 &0&&0&\\
 &&1&&
 \end{array}
\]
(w kolejnych wierszach odczytujemy stopnie $p+q$).
Wyjaśnia to, dlaczego $22$ niezależne klasy w $H^2$
nie są równoważne $22$ niezależnym klasom dzielników:
klasy typu $(2,0)$ i $(0,2)$ również wnoszą wkład do $b_2$.
\end{przyklad}

\begin{center}
\small
\begin{tabular}{c|r|r|r|r|r}
$d$ & $\chi_{\rm top}$ & $b_2$ & $p_g$ & $h^{1,1}$ & $K_S^2$\\\hline
$1$ & $3$ & $1$ & $0$ & $1$ & $9$\\
$2$ & $4$ & $2$ & $0$ & $2$ & $8$\\
$3$ & $9$ & $7$ & $0$ & $7$ & $3$\\
$4$ & $24$ & $22$ & $1$ & $20$ & $0$\\
$5$ & $55$ & $53$ & $4$ & $45$ & $5$
\end{tabular}

\smallskip
{\small Tabela 10.1. Niezmienniki gładkich powierzchni
stopnia $d$ w $\PP^3$; wszystkie kolumny wynikają
z formuł udowodnionych w tym rozdziale przy wskazanych
twierdzeniach zewnętrznych.}\par
\end{center}

\begin{stwierdzenie}[Kontrola wzoru Noethera]
Dla gładkiej powierzchni $S\subset\PP^3_\C$ stopnia $d\ge1$
i dowolnego dzielnika kanonicznego $K_S$ zachodzi
$12\chi(S,\mathcal O_S)=K_S^2+\chi_{\rm top}(S)$.
\end{stwierdzenie}
\begin{proof}
Z obliczonych wyżej liczb Hodge'a wynika
$\chi(S,\mathcal O_S)=h^0-h^1+h^2=1+p_g$;
od strony algebraicznej $h^2(\mathcal O_S)=h^0(K_S)$
z dualności Serre'a. Dla $d\ge4$ dostajemy
$12(1+\binom{d-1}{3})$.
Po podstawieniu wzorów na $K_S^2$ i
$\chi_{\rm top}(S)$ prawa strona jest równa
$d(d-4)^2+d(d^2-4d+6)=2d(d^2-6d+11)$.
Dla $d\ge4$ rozwinięcie dwumianu daje
$12+2(d-1)(d-2)(d-3)=2d(d^2-6d+11)$.
Dla $d=1,2,3$ obie strony wynoszą $12$.
Jest to bezpośrednia kontrola szczególnego przypadku
wzoru Noethera, a nie niezależny dowód pełnego
Hirzebrucha--Riemanna--Rocha.
\end{proof}

\section{Osobliwość w rodzinie i cykl zanikający}
\label{sec:ga10-vanishing}

\noindent\emph{Powiązania tematyczne w Bibliotece matematycznej:} \atlasref{TOP-25::def:nearby-vanishing-xiv}{TOP-25: cykle pobliskie i zanikające}; \atlasref{TOP-25::thm:picard-lefschetz-xiv}{TOP-25: formuła Picarda--Lefschetza}.


Topologia gładkiej krzywej może zmienić się, kiedy
w rodzinie pojawi się punkt osobliwy. Jest to ważna
motywacja teorii osobliwości oraz monodromii.

\begin{przyklad}[Lokalne wygładzenie węzła]
\label{prz:ga10-node}
Rozważmy w $\C^2$ rodzinę $C_t=\{(u,v):uv=t\}$,
dla parametrów $t\in\C$ bliskich zeru.
Dla $t=0$ otrzymujemy dwie osie, które przecinają się
transwersalnie w punkcie $(0,0)$. Dla $t\ne0$
równość $v=t/u$ identyfikuje $C_t$ z $\C^*$.
Jeżeli ograniczymy się do małej kuli wokół zera,
zbiór $C_t$ dla odpowiednio małego $t\ne0$ jest
analitycznym pierścieniem (ma typ homotopii okręgu).
Gdy $t\to0$, ten okrąg może skurczyć się do punktu
w osobliwym włóknie. Nazywamy go \emph{cyklem zanikającym}.
\end{przyklad}

\begin{proof}[Dlaczego pojawia się pierścień?]
Wybierzmy promień $\varepsilon>0$ i ograniczmy $C_t$
do kuli $B_\varepsilon=\{|u|^2+|v|^2<\varepsilon^2\}$.
Dla $t\ne0$ podstawienie $v=t/u$ prowadzi do
$|u|^2+|t|^2/|u|^2<\varepsilon^2$.
Przy $0<2|t|<\varepsilon^2$ nierówność ma postać
$r_-<|u|<r_+$ dla pewnych $0<r_-<r_+$.
Istotnie, po podstawieniu $s=|u|^2$ otrzymujemy
nierówność kwadratową $s^2-\varepsilon^2s+|t|^2<0$
o dwóch dodatnich pierwiastkach.
Jest to obszar pierścieniowy w płaszczyźnie $u$,
a więc odkształca się do okręgu $|u|=\sqrt{|t|}$.
Ten okrąg jest cyklem zanikającym.
\end{proof}

\begin{center}
\begin{tikzpicture}[>=Stealth,every node/.style={font=\footnotesize}]
 \node at (-2.95,1.48) {$uv=0$};
 \draw[->,gray] (-4.25,0)--(-1.55,0);
 \draw[->,gray] (-2.95,-1.18)--(-2.95,1.18);
 \draw[very thick,blue!75!black] (-4.05,0)--(-1.85,0);
 \draw[very thick,blue!75!black] (-2.95,-1.03)--(-2.95,1.04);
 \fill[red!70!black] (-2.95,0) circle (2.5pt);
 \draw[->,thick] (-1.16,.06)--(.98,.06)
  node[midway,above,font=\scriptsize] {wygładzenie};
 \node at (2.72,1.48) {$uv=t\ne0$};
 \draw[blue!75!black,thick] (2.72,0) ellipse (1.18 and .67);
 \draw[blue!75!black,dashed] (2.72,0) ellipse (.49 and .31);
 \draw[->,red!70!black,thick] (2.77,-.29) arc (-80:135:.42 and .30);
 \node[font=\scriptsize,align=center] at (2.72,-1.17)
  {cykl zanikający\\(typ homotopii $S^1$)};
\end{tikzpicture}

{\small Rysunek 10.3. Schemat topologiczny węzła i jego wygładzenia.
Prawy rysunek jest modelem pierścienia, a nie rzeczywistym wykresem
pełnej krzywej zespolonej $uv=t$ w czterech wymiarach rzeczywistych.}\par
\end{center}

\begin{przyklad}[Rzutowa krzywa węzłowa i normalizacja]
Rozważmy krzywą kubiczną
$C_0=\{Y^2Z=X^2(X+Z)\}\subset\PP^2_\C$.
W punkcie $[0:0:1]$ ma ona węzeł: w mapie $Z=1$
równanie brzmi $y^2=x^2(x+1)$, a jego częścią
jednorodną najniższego stopnia jest $y^2-x^2$,
czyli iloczyn dwóch różnych form liniowych.
Krzywa jest nierozkładalna: w mapie afinicznej wielomian
$y^2-x^2(x+1)$ nie rozkłada się w $\C[x,y]$, ponieważ
$x+1$ nie jest kwadratem w $\C[x]$, a jako wielomian
moniczny stopnia dwa w $y$ mógłby rozłożyć się tylko
po istnieniu pierwiastka w $\C[x]$.
Jej normalizacja ma jednak rodzaj zero. Istotnie, na otwartym podzbiorze
$x\ne0$ możemy wprowadzić $s=y/x$; wtedy
$x=s^2-1$ oraz $y=s(s^2-1)$.
Jednorodnie zapisujemy tę parametryzację jako morfizm
$\nu\colon\PP^1\to C_0$ dany wzorem
$\nu([S:T])=[T(S^2-T^2):S(S^2-T^2):T^3]$.
Wzór ma sens także dla $T=0$, nie ma wspólnych zer wszystkich
trzech współrzędnych i spełnia równanie kubiki.
Jest to morfizm biwymierny z gładkiej krzywej $\PP^1$;
ponieważ jest skończony (jest właściwy i ma skończone włókna),
stanowi normalizację $C_0$ na mocy własności uniwersalnej
normalizacji krzywych. Użyliśmy tu standardowych faktów
o skończoności właściwego morfizmu quasi-skończonego.
Morfizm rozdziela dwie gałęzie węzła:
$s=1$ i $s=-1$ mają ten sam obraz $[0:0:1]$.
Zatem topologicznie $C_0$ otrzymujemy z $S^2$
przez sklejenie dwóch różnych punktów, a więc
$\chi_{\rm top}(C_0)=2-1=1$.

Dla małego ogólnego $t\ne0$ równanie
$Y^2Z=X^2(X+Z)+tZ^3$ definiuje gładką
krzywą kubiczną, czyli krzywą rodzaju jeden.
Jej charakterystyka Eulera jest równa zero.
Zmiana $\chi_{\rm top}$ z $1$ na $0$ jest
obrazem lokalnej zmiany w węźle: w miejsce
skrzyżowania dwóch gałęzi pojawia się pierścień.
\end{przyklad}

\begin{uwaga}[Monodromia i wzór Picarda--Lefschetza]
Przechodząc z parametrem $t$ raz wokół zera,
możemy przenosić cykle gładkich włókien.
Otrzymany automorfizm homologii jest \emph{monodromią}.
Dla pojedynczego zwyczajnego punktu podwójnego
w odpowiednim wymiarze opisuje go wzór Picarda--Lefschetza:
transformacja różni się od identyczności o składnik
wyznaczony przez przecięcie z cyklem zanikającym
(z dokładnym znakiem zależnym od wymiaru i przyjętych
konwencji orientacji). Pełna konstrukcja monodromii
wymaga teorii fibracji Milnora i pozostaje poza
zakresem bieżącego rozdziału. Lokalny rachunek $uv=t$
pokazuje jednak geometryczną ideę stojącą za tym wzorem.
\end{uwaga}

\section{Jakie ograniczenia topologiczne są szczególnie użyteczne?}
\label{sec:ga10-konsekwencje}

\begin{stwierdzenie}[Trzy testy dla gładkiej rozmaitości rzutowej]
Niech $X$ będzie gładką spójną rzutową odmianą nad $\C$
o wymiarze zespolonym $n$. Wtedy:
\begin{enumerate}[label=\arabic*)]
\item jej nieparzyste liczby Bettiego są parzyste;
\item dla każdego $k<n$ zachodzi $b_k(X)\le b_{k+2}(X)$;
\item jeżeli $n=2$, to forma przecięć na $N^1(X)_{\R}$
ma dokładnie jeden dodatni kierunek, a pozostałe
$\rho(X)-1$ kierunków są ujemne.
\end{enumerate}
\end{stwierdzenie}
\begin{proof}
Punkt pierwszy wyprowadziliśmy z symetrii Hodge'a
we wniosku~\ref{wn:ga9-odd-betti}. Drugi wynika z
iniektywności mnożenia przez klasę Kählerowską
w stopniach $k<n$, co zostało uzasadnione z twardego
Lefschetza w rozdziale~\ref{ch:ga9-analiza-zespolona}.
Trzeci jest twierdzeniem Hodge'a o indeksie
\ref{tw:ga10-index}. Każde założenie rzutowości
jest tu istotne co najmniej jako gwarancja struktury
Kählerowskiej oraz odpowiednich własności klas dzielników.
\end{proof}

\begin{przyklad}[Powierzchnia Hopfa jako kontrprzykład]
Powierzchnia Hopfa z rozdziału~\ref{ch:ga9-analiza-zespolona}
ma $b_1=1$. Nie może więc być Kählerowska,
a w szczególności nie może być gładką rzutową
powierzchnią nad $\C$. Jest to ograniczenie
czysto topologiczne, wykrywające brak zanurzenia
rzutowego bez rozwiązywania równań wielomianowych.
\end{przyklad}

\begin{uwaga}[Ostrożność: nie każdy niezmiennik wyznacza odmianę]
Dwie powierzchnie o jednakowych liczbach Bettiego,
charakterystyce Eulera i liczbach Cherna nie muszą być
izomorficzne algebraicznie ani biholomorficzne.
Nawet pełna topologia różniczkowa nie determinuje
w ogólności struktury zespolonej. Dla K3 istnieją
wielowymiarowe rodziny nieizomorficznych struktur,
a zmieniająca się część geometrii jest widoczna
w okresach form holomorficznych i w położeniu krat
klas algebraicznych wewnątrz $H^2$.
\end{uwaga}

\section{Synteza i mapa zależności między twierdzeniami}
\label{sec:ga10-koniec}

W tym rozdziale powiązaliśmy trzy poziomy opisu.
Równania wielomianowe dostarczają wiązek, klas Cherna
oraz cykli; analiza Kählerowska dostarcza form
harmonicznych, rozkładu Hodge'a i operatorów Lefschetza;
topologia dostarcza klas fundamentalnych, kohomologii
singularnej, charakterystyki Eulera i form przecięć.

\begin{center}
\begin{tikzpicture}[>=Stealth,every node/.style={font=\footnotesize,align=center},
  bx/.style={draw,rounded corners=3pt,text width=40mm,minimum height=14mm,inner sep=4pt}]
 \node[bx,fill=blue!4] (alg) at (-3.6,.9)
   {geometria algebraiczna\\cykle, dzielniki, klasy Cherna};
 \node[bx,fill=green!4] (an) at (3.5,.9)
   {analiza zespolona\\Dolbeault, Hodge, Kähler};
 \node[bx,fill=orange!5] (top) at (0,-1.45)
   {topologia\\kohomologia, przecięcia, Euler};
 \draw[<->,thick] (alg)--node[above,font=\scriptsize] {GAGA, $(1,1)$} (an);
 \draw[->,thick] (alg)--node[below left,font=\scriptsize] {klasa cyklu} (top);
 \draw[->,thick] (an)--node[below right,font=\scriptsize] {de Rham, Lefschetz} (top);
\end{tikzpicture}

{\small Rysunek 10.4. Powiązania trzech języków geometrii.
Strzałki oznaczają określone konstrukcje i twierdzenia
porównawcze, a nie równość odpowiednich kategorii.}\par
\end{center}

\begin{uwaga}[Granice samowystarczalności]
Dowiedliśmy szczegółowo formuł dla $\chi_{\rm top}$,
$K^2$, liczb Bettiego i części liczb Hodge'a
gładkich powierzchni stopnia $d$ w $\PP^3$,
pokazaliśmy zmianę tych niezmienników przy rozdmuchaniu,
wyprowadziliśmy nierówność Hodge'a dla dzielników,
wyjaśniliśmy ciąg wykładniczy oraz obliczyliśmy
homotopię lokalnego wygładzenia węzła.
Jako twierdzenia zewnętrzne pozostają: dualność Poincarégo,
twierdzenie de Rhama, słabe i twarde Lefschetza,
twierdzenie Hodge'a o indeksie i relacje Hodge'a--Riemanna,
Chern--Gauss--Bonnet, GAGA, Dolbeault i ogólne twierdzenia
porównawcze o klasach cykli. Zostały one użyte z jawnymi
założeniami, a nie ukryte w rachunkowych przeskokach.
\end{uwaga}

% Podstawowe pozycje do dalszej lektury i weryfikacji:
% W. Fulton, Intersection Theory.
% C. Voisin, Hodge Theory and Complex Algebraic Geometry I.
% P. Griffiths, J. Harris, Principles of Algebraic Geometry.
% R. Hartshorne, Algebraic Geometry.
% A. Hatcher, Algebraic Topology (duality, cohomology, Euler).
% Stacks Project: https://stacks.math.columbia.edu/


\end{document}
